\beginnumbering
\eledsection*{Adhyāya 1}
\pstart \ledsidenote{\passref{4.1.1}}\skt{अथातो द्विव्रणीयं चिकित्सितं व्याख्यास्यामः~।}\pend

\pstart \ledsidenote{\passref{4.1.3}}\skt{द्वौ व्रणौ भवतः~। }\edtext{\edtext{\skt{शरी\-र\-स\-मु\-त्थ\-श्चा\-ग\-न्तुश्च}}{\lemma{\skt{शरी\appbrk{}र\appbrk{}स\appbrk{}मु\appbrk{}त्थश्}}\Afootnote{H; \skt{शारीर}\ A}}}{\lemma{\skt{चागन्तुश्च}}\Afootnote{H; \skt{आगन्तुश्च}\ A}}\skt{ तयोः }\edtext{\skt{शरीरसमुत्थो}}{\lemma{\skt{शरीरसमुत्थो}}\Afootnote{H; \skt{शारीरः प}\ A}}\skt{ }\edtext{\skt{वात\-पि\-त्त\-क\-फ\-शो\-णि\-त\-स\-न्नि\-पा\-त\-नि\-मित्त}}{\lemma{\skt{वात\appbrk{}पि\appbrk{}त्त\appbrk{}क\appbrk{}फ\appbrk{}शो\appbrk{}णि\appbrk{}त\appbrk{}स\appbrk{}न्नि\appbrk{}पा\appbrk{}त\appbrk{}नि\appbrk{}मित्त}}\Afootnote{H; \skt{वन\appbrk{}पि\appbrk{}त्त\appbrk{}क\appbrk{}फ\appbrk{}शो\appbrk{}णि\appbrk{}त\appbrk{}स\appbrk{}न्नि\appbrk{}पा\appbrk{}त\appbrk{}नि\appbrk{}मित्तः}\ A}}\skt{ }\edtext{\skt{आगन्तुकः}}{\lemma{\skt{आगन्तुकः}}\Afootnote{H; \skt{आगन्तुरपि}\ A}}\skt{ }\edtext{\skt{पुरु\-ष\-मृ\-ग\-प\-शु\-प\-क्षि\-व्या\-ड\-स\-री\-सृ\-प\-प\-त\-न\-पी\-ड\-न\-प्र\-हा\-रा\-ग्नि\-वि\-षा\-स्थि\-स\-क\-ल\-क\-पा\-ल\-ख\-ड्ग\-च\-क्रेषु}}{\lemma{\skt{पुरु\appbrk{}ष\appbrk{}मृ\appbrk{}ग\appbrk{}प\appbrk{}शु\appbrk{}प\appbrk{}क्षि\appbrk{}व्या\appbrk{}ड\appbrk{}स\appbrk{}री\appbrk{}सृ\appbrk{}प\appbrk{}प\appbrk{}त\appbrk{}न\appbrk{}पी\appbrk{}ड\appbrk{}न\appbrk{}प्र\appbrk{}हा\appbrk{}रा\appbrk{}ग्नि\appbrk{}वि\appbrk{}षा\appbrk{}स्थि\appbrk{}स\appbrk{}क\appbrk{}ल\appbrk{}क\appbrk{}पा\appbrk{}ल\appbrk{}ख\appbrk{}ड्ग\appbrk{}च\appbrk{}क्रेषु}}\Afootnote{H; \skt{पुरु\appbrk{}ष\appbrk{}प\appbrk{}शु\appbrk{}प\appbrk{}क्षि\appbrk{}व्या\appbrk{}ल\appbrk{}स\appbrk{}री\appbrk{}सृ\appbrk{}प\appbrk{}प्र\appbrk{}प\appbrk{}त\appbrk{}न\appbrk{}पी\appbrk{}ड\appbrk{}न\appbrk{}प्र\appbrk{}हा\appbrk{}रा\appbrk{}ग्नि\appbrk{}क्षा\appbrk{}र\appbrk{}वि\appbrk{}ष\appbrk{}ती\appbrk{}क्ष्णौ\appbrk{}ष\appbrk{}ध\appbrk{}श\appbrk{}क\appbrk{}ल\appbrk{}क\appbrk{}पा\appbrk{}ल\appbrk{}शृ\appbrk{}ङ्ग\appbrk{}च्क्रेषु}\ A}}\skt{ }\edtext{\skt{पर\-शु\-श\-क्ति\-कु\-न्ता\-द्या\-यु\-ध\-प्र\-ह\-रे\-णा\-भि\-घा\-त\-नि\-मित्तः}}{\lemma{\skt{पर\appbrk{}शु\appbrk{}श\appbrk{}क्ति\appbrk{}कु\appbrk{}न्ता\appbrk{}द्या\appbrk{}यु\appbrk{}ध\appbrk{}प्र\appbrk{}ह\appbrk{}रे\appbrk{}णा\appbrk{}भि\appbrk{}घा\appbrk{}त\appbrk{}नि\appbrk{}मित्तः}}\Afootnote{H; \skt{पर\appbrk{}शु\appbrk{}श\appbrk{}क्ति\appbrk{}कु\appbrk{}न्ता\appbrk{}द्या\appbrk{}यु\appbrk{}धा\appbrk{}भि\appbrk{}घा\appbrk{}त\appbrk{}नि\appbrk{}मित्तः}\ A}}\skt{~। तत्र तुल्ये व्रणसामान्ये द्विका\-र\-णो\-त्था\-न\-प्र\-यो\-ज\-न\-सा\-म\-र्थ्या\-द्द्वि\-व्र\-णीयः इत्युच्यते~।}\pend

\pstart \ledsidenote{\passref{4.1.4}}\edtext{\skt{सर्व्वस्मिन्नेव}}{\lemma{\skt{सर्व्वस्मिन्नेव}}\Afootnote{H; \skt{सर्वस्मिन्ने}\ A}}\skt{ }\edtext{\skt{चागन्तौ}}{\lemma{\skt{चागन्तौ}}\Afootnote{H; \skt{वागन्तु}\ A}}\skt{ व्रणे तत्कालमेव क्षतोष्मणः~। }\edtext{\edtext{\skt{प्रसृ\-त\-स्यौ\-प\-श\-मा\-र्थ\-म्पि\-त्त\-व\-च्छी\-त\-क्रि\-या\-वि\-भागः}}{\lemma{\skt{प्रसृ\appbrk{}त\appbrk{}स्यौ\appbrk{}प\appbrk{}श\appbrk{}मा\appbrk{}र्थम्}}\Afootnote{H; \skt{प्रसृ\appbrk{}त\appbrk{}स्यो\appbrk{}प\appbrk{}श\appbrk{}मार्थं}\ A}}}{\lemma{\skt{पित्त\appbrk{}व\appbrk{}च्छी\appbrk{}त\appbrk{}क्रि\appbrk{}या\appbrk{}वि\appbrk{}भागः}}\Afootnote{H; \skt{चेति श्लेष्मशोणिताभ्यां, स्फुर\appbrk{}ण\appbrk{}तो\appbrk{}द\appbrk{}दा\appbrk{}ह\appbrk{}धू\appbrk{}मा\appbrk{}य\appbrk{}न\appbrk{}प्रायः पीत\appbrk{}त\appbrk{}नु\appbrk{}र\appbrk{}क्त\appbrk{}स्रावी चेति वात\appbrk{}पि\appbrk{}त्त\appbrk{}शो\appbrk{}णि\appbrk{}तेभ्यः, कण्डूस्फुरण}\ A}}\skt{~। }\edtext{\skt{सन्धानार्थञ्च}}{\lemma{\skt{सन्धानार्थञ्च}}\Afootnote{H; \skt{चु}\ A}}\skt{ }\edtext{\skt{मधु\-घृ\-तो\-प\-योग इति~। एत\-द्वि\-का\-र\-णो\-त्था\-न\-प्र\-यो\-ज\-न\-मु\-त्त\-र\-का\-लन्तु}}{\lemma{\skt{मधु\appbrk{}घृ\appbrk{}तो\appbrk{}प\appbrk{}योग इति~। एत\appbrk{}द्वि\appbrk{}का\appbrk{}र\appbrk{}णो\appbrk{}त्था\appbrk{}न\appbrk{}प्र\appbrk{}यो\appbrk{}ज\appbrk{}न\appbrk{}मु\appbrk{}त्त\appbrk{}र\appbrk{}का\appbrk{}लन्तु}}\Afootnote{H; \skt{मचु\appbrk{}मा\appbrk{}य\appbrk{}मा\appbrk{}न\appbrk{}प्रायः पाण्डु\appbrk{}घ\appbrk{}न\appbrk{}र\appbrk{}क्ता\appbrk{}स्रावी चेति वात\appbrk{}श्ले\appbrk{}ष्म\appbrk{}शो\appbrk{}णि\appbrk{}तेभ्यः, दाह\appbrk{}पा\appbrk{}क\appbrk{}रा\appbrk{}ग\appbrk{}क\appbrk{}ण्डू\appbrk{}प्रायः पाण्डु\appbrk{}घ\appbrk{}न\appbrk{}र\appbrk{}क्ता\appbrk{}स्रावी चेति पित्त\appbrk{}श्ले\appbrk{}ष्म\appbrk{}शो\appbrk{}णि\appbrk{}तेभ्यः, त्रिवि\appbrk{}ध\appbrk{}व\appbrk{}र्ण\appbrk{}वे\appbrk{}द\appbrk{}ना\appbrk{}स्रा\appbrk{}व\appbrk{}वि\appbrk{}शे\appbrk{}षो\appbrk{}पेतः पव}\ A}}\skt{ }\edtext{\skt{दोषो\-प\-प्ल\-व\-वि\-शे\-षा\-द\-व\-स्था\-न्त\-र\-वि\-शेषाः}}{\lemma{\skt{दोषो\appbrk{}प\appbrk{}प्ल\appbrk{}व\appbrk{}वि\appbrk{}शे\appbrk{}षा\appbrk{}द\appbrk{}व\appbrk{}स्था\appbrk{}न्त\appbrk{}र\appbrk{}वि\appbrk{}शेषाः}}\Afootnote{H; \skt{नपित्तकफेभ्यः, निर्द\appbrk{}ह\appbrk{}न\appbrk{}नि\appbrk{}र्म\appbrk{}थ\appbrk{}न\appbrk{}स्फु\appbrk{}र\appbrk{}ण\appbrk{}तो\appbrk{}द\appbrk{}दा\appbrk{}ह\appbrk{}पा\appbrk{}क\appbrk{}रा\appbrk{}ग\appbrk{}क\appbrk{}ण्डू\appbrk{}स्वा\appbrk{}प\appbrk{}ब\appbrk{}हुलो नाना\appbrk{}व\appbrk{}र्ण\appbrk{}वे\appbrk{}द\appbrk{}ना\appbrk{}स्रा\appbrk{}व\appbrk{}वि\appbrk{}शेषो}\ A}}\skt{ }\edtext{\skt{शारी\-र\-व\-त्प्र\-ती\-कारः}}{\lemma{\skt{शारी\appbrk{}र\appbrk{}व\appbrk{}त्प्र\appbrk{}ती\appbrk{}कारः}}\Afootnote{H; \skt{पेतः पव\appbrk{}न\appbrk{}पि\appbrk{}त्त\appbrk{}क\appbrk{}फ\appbrk{}शो\appbrk{}णि\appbrk{}तेभ्यः, जिह्वातलाभो मृदुः स्निग्धः श्लक्ष्णो विगतवेदनः सुव्यवस्थितो निरास्रावश्चेति शुद्धो व्रण इति}\ A}}\skt{~।}\pend

\pstart \ledsidenote{\passref{4.1.5}}\skt{दोषो\-प\-प्ल\-व\-वि\-शेषः पुनः समासतः पञ्च\-द\-श\-प्र\-कारो }\edtext{\skt{भवति}}{\lemma{\skt{भवति}}\Afootnote{H; \emph{om.}\ A}}\skt{~। }\edtext{\skt{यथोक्ते}}{\lemma{\skt{यथोक्ते}}\Afootnote{H; \emph{om.}\ A}}\skt{ }\edtext{\skt{व्रणप्रश्नाध्याये}}{\lemma{\skt{व्रणप्रश्नाध्याये}}\Afootnote{H; \emph{om.}\ A}}\skt{ }\edtext{\edtext{\edtext{\skt{प्रस\-र\-ण\-सा\-म\-र्थ्या\-च्छ\-द्ध\-श्चेति}}{\lemma{\skt{प्रस\appbrk{}र\appbrk{}ण\appbrk{}सा\appbrk{}म\appbrk{}र्थ्याच्}}\Afootnote{H; \skt{प्रस\appbrk{}र\appbrk{}ण\appbrk{}सा\appbrk{}म\appbrk{}र्थ्यात्, यथोक्तो व्रण\appbrk{}प्र\appbrk{}श्ना\appbrk{}धि\appbrk{}कारे}\ A}}}{\lemma{\skt{छद्धश्}}\Afootnote{H; \skt{शुद्धत्}\ A}}}{\lemma{\skt{चेति}}\Afootnote{H; \skt{वात्}\ A}}\skt{ षोडशप्रकार }\edtext{\skt{इत्युच्यते}}{\lemma{\skt{उच्यते}}\Afootnote{H; \skt{एके}\ A}}\skt{~।}\pend

\pstart \ledsidenote{\passref{4.1.6}}\skt{तस्य लक्षणं द्विविधं~। }\edtext{\edtext{\skt{सामा\-न्य\-ल\-क्ष\-ण\-म्वि\-शे\-ष\-ल\-क्ष\-णञ्च}}{\lemma{\skt{सामा\appbrk{}न्य\appbrk{}ल\appbrk{}क्ष\appbrk{}णम्}}\Afootnote{H; \skt{सामान्यं}\ A}}}{\lemma{\skt{विशे\appbrk{}ष\appbrk{}ल\appbrk{}क्ष\appbrk{}णञ्च}}\Afootnote{H; \skt{वैशेषिकं च}\ A}}\skt{~। }\edtext{\skt{तत्र}}{\lemma{\skt{तत्र}}\Afootnote{H; \skt{तत्र सामान्यं रुक्। व्रण गात्र}\ A}}\skt{ }\edtext{\skt{व्रणे}}{\lemma{\skt{व्रणे}}\Afootnote{H; \skt{विचूर्णने, व्रण}\ A}}\skt{ }\edtext{\skt{वृणोतीति}}{\lemma{\skt{वृणोतीति}}\Afootnote{H; \skt{यतीति}\ A}}\skt{ व्रणः~। }\edtext{\skt{विशे\-ष\-ल\-क्ष\-ण\-म्पु\-न\-र्व्वा\-ता\-दि\-स्व\-ल\-क्ष\-ण\-वि\-शेषः}}{\lemma{\skt{विशे\appbrk{}ष\appbrk{}ल\appbrk{}क्ष\appbrk{}ण\appbrk{}म्पु\appbrk{}न\appbrk{}र्व्वा\appbrk{}ता\appbrk{}दि\appbrk{}स्व\appbrk{}ल\appbrk{}क्ष\appbrk{}ण\appbrk{}वि\appbrk{}शेषः}}\Afootnote{H; \skt{विशेषलक्षणं पुन\appbrk{}र्वा\appbrk{}ता\appbrk{}दि\appbrk{}लि\appbrk{}ङ्ग\appbrk{}वि\appbrk{}शेषः}\ A}}\skt{~।}\pend

\pstart \ledsidenote{\passref{4.1.7}}\skt{तत्र }\edtext{\skt{श्यावा\-रु\-णा\-भ\-भ\-स्तनु}}{\lemma{\skt{श्यावा\appbrk{}रु\appbrk{}णा\appbrk{}भ\appbrk{}भ\appbrk{}स्तनु}}\Afootnote{\skt{श्यावा\appbrk{}रु\appbrk{}णा\appbrk{}भ\appbrk{}स्तनुः}\ A; \skt{श्यावा\appbrk{}रु\appbrk{}णा\appbrk{}भ\appbrk{}भ\appbrk{}स्तनु}\ H}}\skt{ }\edtext{\skt{शीत\-पि\-च्छि\-ला\-ल्प\-स्रा\-वी\-रू\-क्ष\-श्च\-ट\-च\-टा\-य\-न\-शीलः}}{\lemma{\skt{शीत\appbrk{}पि\appbrk{}च्छि\appbrk{}ला\appbrk{}ल्प\appbrk{}स्रा\appbrk{}वी\appbrk{}रू\appbrk{}क्ष\appbrk{}श्च\appbrk{}ट\appbrk{}च\appbrk{}टा\appbrk{}य\appbrk{}न\appbrk{}शीलः}}\Afootnote{H; \skt{शीतः पिच्छिलोऽल्पस्रावी रूक्ष\appbrk{}श्च\appbrk{}ट\appbrk{}च\appbrk{}टा\appbrk{}य\appbrk{}न\appbrk{}शीलः}\ A}}\skt{ }\edtext{\edtext{\skt{स्फुर\-णा\-या\-म\-तो\-द\-भे\-द\-वे\-द\-ना\-ब\-हुलो}}{\lemma{\skt{स्फुरणायाम्}}\Afootnote{H; \skt{स्फुरणाया}\ A}}}{\lemma{\skt{अतो\appbrk{}द\appbrk{}भे\appbrk{}द\appbrk{}वे\appbrk{}द\appbrk{}ना\appbrk{}ब\appbrk{}हुलो}}\Afootnote{H; \skt{मतो\appbrk{}द\appbrk{}भे\appbrk{}द\appbrk{}वे\appbrk{}द\appbrk{}ना\appbrk{}ब\appbrk{}हुलो}\ A}}\skt{ निर्म्मांसश्चेति वातात्~॥ क्षिप्रजः पीतनीलाभः किंशु\-को\-द\-का\-भो\-ष्ण\-स्रावी }\edtext{\skt{दाह\-पा\-क\-वि\-कारी}}{\lemma{\skt{दाह\appbrk{}पा\appbrk{}क\appbrk{}वि\appbrk{}कारी}}\Afootnote{H; \skt{दाह\appbrk{}पा\appbrk{}क\appbrk{}रा\appbrk{}ग\appbrk{}वि\appbrk{}का\appbrk{}र\appbrk{}कारी}\ A}}\skt{ }\edtext{\skt{पीत\-पि\-ट\-का\-जु\-ष्ट\-श्चेति}}{\lemma{\skt{पीत\appbrk{}पि\appbrk{}ट\appbrk{}का\appbrk{}जु\appbrk{}ष्टश्}}\Afootnote{H; \skt{पीत\appbrk{}पि\appbrk{}ड\appbrk{}का\appbrk{}जु\appbrk{}ष्टश्}\ A}}\skt{ पित्तात्~॥ }\edtext{\skt{प्रततः}}{\lemma{\skt{प्रततः}}\Afootnote{H; \skt{प्रतत}\ A}}\skt{ }\edtext{\skt{कण्डुबहुलः}}{\lemma{\skt{कण्डुबहुलः}}\Afootnote{H; \skt{चण्ड\appbrk{}क\appbrk{}ण्डू\appbrk{}ब\appbrk{}हुलः}\ A}}\skt{ }\edtext{\skt{स्थूलोष्ठः}}{\lemma{\skt{स्थूलोष्ठः}}\Afootnote{H; \skt{स्थूलौष्ठः}\ A}}\skt{ }\edtext{\skt{स्तव्ध\-सि\-रा\-स्ना\-यु\-जा\-ला\-व\-नद्धः}}{\lemma{\skt{स्तव्ध\appbrk{}सि\appbrk{}रा\appbrk{}स्ना\appbrk{}यु\appbrk{}जा\appbrk{}ला\appbrk{}व\appbrk{}नद्धः}}\Afootnote{H; \skt{स्तब्ध\appbrk{}सि\appbrk{}रा\appbrk{}स्ना\appbrk{}यु\appbrk{}जा\appbrk{}ला\appbrk{}व\appbrk{}ततः}\ A}}\skt{ कठिनः }\edtext{\skt{पाण्ड्वावभासो}}{\lemma{\skt{पाण्ड्वावभासो}}\Afootnote{H; \skt{पाण्ड्ववभासो}\ A}}\skt{ मन्दवेदनः~। शुक्ल\-शी\-त\-सा\-न्द्र\-पि\-च्छि\-ला\-स्रावी गुरुश्चेति कफात्~॥ }\edtext{\skt{प्रवा\-ला\-न\-ल\-नि\-च\-य\-प्र\-काशः}}{\lemma{\skt{प्रवा\appbrk{}ला\appbrk{}न\appbrk{}ल\appbrk{}नि\appbrk{}च\appbrk{}य\appbrk{}प्र\appbrk{}काशः}}\Afootnote{H; \skt{प्रवा\appbrk{}ल\appbrk{}द\appbrk{}ल\appbrk{}नि\appbrk{}च\appbrk{}य\appbrk{}प्र\appbrk{}काशः}\ A}}\skt{ }\edtext{\skt{कृष्ण\-स्फो\-ट\-पि\-ट\-का\-जा\-लो\-प\-चि\-त\-स्तु\-रं\-ग\-स्था\-न\-गन्धी}}{\lemma{\skt{कृष्ण\appbrk{}स्फो\appbrk{}ट\appbrk{}पि\appbrk{}ट\appbrk{}का\appbrk{}जा\appbrk{}लो\appbrk{}प\appbrk{}चि\appbrk{}त\appbrk{}स्तु\appbrk{}रं\appbrk{}ग\appbrk{}स्था\appbrk{}न\appbrk{}गन्धी}}\Afootnote{H; \skt{कृष्ण\appbrk{}स्फो\appbrk{}ट\appbrk{}पि\appbrk{}डा\appbrk{}का\appbrk{}जा\appbrk{}लो\appbrk{}प\appbrk{}चि\appbrk{}त\appbrk{}स्तु\appbrk{}र\appbrk{}ङ्ग\appbrk{}स्था\appbrk{}न\appbrk{}गन्धिः}\ A}}\skt{ सवेदनः~। }\edtext{\skt{प्रधूपनशीलो}}{\lemma{\skt{प्रधूपनशीलो}}\Afootnote{H; \skt{धूमायनशीलो}\ A}}\skt{ रक्तस्रा वी पित्तलिङ्गश्चेति रक्तात्~॥ }\edtext{\skt{प्रतो\-द\-दा\-ह\-धू\-पा\-य\-न\-प्रायः}}{\lemma{\skt{प्रतो\appbrk{}द\appbrk{}दा\appbrk{}ह\appbrk{}धू\appbrk{}पा\appbrk{}य\appbrk{}न\appbrk{}प्रायः}}\Afootnote{H; \skt{तोद\appbrk{}दा\appbrk{}ह\appbrk{}धू\appbrk{}मा\appbrk{}य\appbrk{}न\appbrk{}प्रायः}\ A}}\skt{ पीता\-रु\-णा\-भ\-स्त\-द्व\-र्ण्ण\-स्रावी चेति वातपित्ताभ्यां~॥ कण्डूयनशीलः }\edtext{\skt{सनि\-स्तो\-द\-वा\-त्रूक्षः}}{\lemma{\skt{सनि\appbrk{}स्तो\appbrk{}द\appbrk{}वा\appbrk{}त्रूक्षः}}\Afootnote{H; \skt{सनिस्तोदो रूक्षो}\ A}}\skt{ }\edtext{\skt{खरस्तद्धो}}{\lemma{\skt{खरस्तद्धो}}\Afootnote{H; \skt{गुरुर्दारुणो}\ A}}\skt{ मुहुर्म्मुहुः }\edtext{\skt{शीत\-पि\-च्छ\-ला\-स्रावी}}{\lemma{\skt{शीत\appbrk{}पि\appbrk{}च्छ\appbrk{}ला\appbrk{}स्रावी}}\Afootnote{H; \skt{शीत\appbrk{}पि\appbrk{}च्छि\appbrk{}ला\appbrk{}ल्प\appbrk{}स्रावी}\ A}}\skt{ चेति }\edtext{\skt{वातश्लेष्माभ्यां}}{\lemma{\skt{वातश्लेष्माभ्यां}}\Afootnote{H; \skt{वातश्लेष्मभ्यां}\ A}}\skt{~॥ गुरुः }\edtext{\skt{सदाहोष्णः}}{\lemma{\skt{सदाहोष्णः}}\Afootnote{H; \skt{सदाह उष्णः}\ A}}\skt{ पीतपाण्डुस्रावी चेति }\edtext{\skt{पित्तश्लेष्माभ्यां}}{\lemma{\skt{पित्तश्लेष्माभ्यां}}\Afootnote{H; \skt{पित्तश्लेष्मभ्यां}\ A}}\skt{~॥ }\edtext{\skt{तनु\-रू\-क्ष\-स्तो\-द\-ब\-हुलः}}{\lemma{\skt{तनु\appbrk{}रू\appbrk{}क्ष\appbrk{}स्तो\appbrk{}द\appbrk{}ब\appbrk{}हुलः}}\Afootnote{H; \skt{रूक्ष\appbrk{}स्त\appbrk{}नु\appbrk{}स्तो\appbrk{}द\appbrk{}ब\appbrk{}हुलः}\ A}}\skt{ }\edtext{\skt{प्रसुप्त}}{\lemma{\skt{प्रसुप्त}}\Afootnote{H; \skt{सुप्त}\ A}}\skt{ इव च }\edtext{\skt{कृष्णारुणा}}{\lemma{\skt{कृष्णारुणा}}\Afootnote{H; \skt{रक्तारुणा}\ A}}\skt{ }\edtext{\skt{भस्तद्व्वर्ण्णस्रावी}}{\lemma{\skt{भस्तद्व्वर्ण्णस्रावी}}\Afootnote{H; \skt{भस्तद्वर्णास्रावी}\ A}}\skt{ चेति वातशोणिताभ्यां~॥ घृतमण्डाभो }\edtext{\skt{मीन\-धा\-व\-न\-ग\-न्धी\-मृ\-दु\-र्व्वि\-स\-र्प्पी\-कृ\-ष्णो\-ष्ण\-स्रावी}}{\lemma{\skt{मीन\appbrk{}धा\appbrk{}व\appbrk{}न\appbrk{}ग\appbrk{}न्धी\appbrk{}मृ\appbrk{}दु\appbrk{}र्व्वि\appbrk{}स\appbrk{}र्प्पी\appbrk{}कृ\appbrk{}ष्णो\appbrk{}ष्ण\appbrk{}स्रावी}}\Afootnote{H; \skt{मीन\appbrk{}धा\appbrk{}व\appbrk{}न\appbrk{}तो\appbrk{}य\appbrk{}ग\appbrk{}न्धि\appbrk{}र्मृ\appbrk{}दु\appbrk{}र्वि\appbrk{}स\appbrk{}र्प्यु\appbrk{}ष्ण\appbrk{}कृ\appbrk{}ष्ण\appbrk{}स्रावी}\ A}}\skt{ चेति पित्तशोणिताभ्यां~॥ रक्तो }\edtext{\skt{गुरुस्निग्धः}}{\lemma{\skt{गुरुस्निग्धः}}\Afootnote{H; \skt{गुरुः स्निग्धः}\ A}}\skt{ }\edtext{\skt{कण्डूप्रायः}}{\lemma{\skt{कण्डूप्रायः}}\Afootnote{H; \skt{पिच्छिलः कण्डूप्रायः}\ A}}\skt{ }\edtext{\skt{पिच्छिलः}}{\lemma{\skt{पिच्छिलः}}\Afootnote{H; \skt{स्थिरो}\ A}}\skt{ }\edtext{\skt{सरक्तस्रावी}}{\lemma{\skt{सरक्तस्रावी}}\Afootnote{H; \skt{सर\appbrk{}क्त\appbrk{}पा\appbrk{}ण्डु\appbrk{}स्रावी}\ A}}\skt{ चेति श्लेष्मशोणिताभ्यां~। स्फुर\-ण\-तो\-द\-दा\-ह\-धू\-मा\-य\-न\-प्रायः पीत\-त\-नु\-र\-क्त\-स्रावी चेति वात\-पि\-त्त\-शो\-णि\-तेभ्यः~॥ }\edtext{\skt{कण्डू\-स्फु\-र\-ण\-चु\-मु\-चु\-मा\-य\-मा\-न\-प्रायः}}{\lemma{\skt{कण्डू\appbrk{}स्फु\appbrk{}र\appbrk{}ण\appbrk{}चु\appbrk{}मु\appbrk{}चु\appbrk{}मा\appbrk{}य\appbrk{}मा\appbrk{}न\appbrk{}प्रायः}}\Afootnote{H; \skt{कण्डू\appbrk{}स्फु\appbrk{}र\appbrk{}ण\appbrk{}चु\appbrk{}म\appbrk{}चु\appbrk{}मा\appbrk{}य\appbrk{}मा\appbrk{}न\appbrk{}प्रायः}\ A}}\skt{ }\edtext{\skt{पाण्डु\-त\-नु\-र\-क्त\-स्रावी}}{\lemma{\skt{पाण्डु\appbrk{}त\appbrk{}नु\appbrk{}र\appbrk{}क्त\appbrk{}स्रावी}}\Afootnote{H; \skt{पाण्डु\appbrk{}घ\appbrk{}न\appbrk{}र\appbrk{}क्ता\appbrk{}स्रावी}\ A}}\skt{ चेति वात\-श्ले\-ष्म\-शो\-णि\-तेभ्यः~॥ दाह\-पा\-क\-रा\-ग\-क\-ण्डू\-प्रायः }\edtext{\skt{पाण्डू\-र\-क्त\-पी\-ता\-स्रावी}}{\lemma{\skt{पाण्डू\appbrk{}र\appbrk{}क्त\appbrk{}पी\appbrk{}ता\appbrk{}स्रावी}}\Afootnote{H; \skt{पाण्डु\appbrk{}घ\appbrk{}न\appbrk{}र\appbrk{}क्ता\appbrk{}स्रावी}\ A}}\skt{ चेति पित्त\-श्ले\-ष्म\-शो\-णि\-तेभ्यः~॥ }\edtext{\skt{विवि\-ध\-दो\-ष\-व\-र्ण्ण\-वे\-द\-ना\-स्रा\-व\-वि\-शे\-षो\-पेतः}}{\lemma{\skt{विवि\appbrk{}ध\appbrk{}दो\appbrk{}ष\appbrk{}व\appbrk{}र्ण्ण\appbrk{}वे\appbrk{}द\appbrk{}ना\appbrk{}स्रा\appbrk{}व\appbrk{}वि\appbrk{}शे\appbrk{}षो\appbrk{}पेतः}}\Afootnote{H; \skt{त्रिवि\appbrk{}ध\appbrk{}व\appbrk{}र्ण\appbrk{}वे\appbrk{}द\appbrk{}ना\appbrk{}स्रा\appbrk{}व\appbrk{}वि\appbrk{}शे\appbrk{}षो\appbrk{}पेतः}\ A}}\skt{ पव\-न\-पि\-त्त\-क\-फेभ्यः~॥ }\edtext{\skt{निर्द्द\-र्श\-क\-नि\-र्म्म\-थ\-न\-स्फु\-र\-ण\-तो\-द\-दा\-ह\-पा\-क\-रा\-ग\-क\-ण्डू\-वां\-स्वा\-प\-ब\-हुलो}}{\lemma{\skt{निर्द्द\appbrk{}र्श\appbrk{}क\appbrk{}नि\appbrk{}र्म्म\appbrk{}थ\appbrk{}न\appbrk{}स्फु\appbrk{}र\appbrk{}ण\appbrk{}तो\appbrk{}द\appbrk{}दा\appbrk{}ह\appbrk{}पा\appbrk{}क\appbrk{}रा\appbrk{}ग\appbrk{}क\appbrk{}ण्डू\appbrk{}वां\appbrk{}स्वा\appbrk{}प\appbrk{}ब\appbrk{}हुलो}}\Afootnote{H; \skt{निर्द\appbrk{}ह\appbrk{}न\appbrk{}नि\appbrk{}र्म\appbrk{}थ\appbrk{}न\appbrk{}स्फु\appbrk{}र\appbrk{}ण\appbrk{}तो\appbrk{}द\appbrk{}दा\appbrk{}ह\appbrk{}पा\appbrk{}क\appbrk{}रा\appbrk{}ग\appbrk{}क\appbrk{}ण्डू\appbrk{}स्वा\appbrk{}प\appbrk{}ब\appbrk{}हुलो}\ A}}\skt{ }\edtext{\skt{नाना\-वि\-ध\-व\-र्ण्ण\-वे\-द\-ना\-स्रा\-व\-वि\-शे\-षो\-पेतः}}{\lemma{\skt{नाना\appbrk{}वि\appbrk{}ध\appbrk{}व\appbrk{}र्ण्ण\appbrk{}वे\appbrk{}द\appbrk{}ना\appbrk{}स्रा\appbrk{}व\appbrk{}वि\appbrk{}शे\appbrk{}षो\appbrk{}पेतः}}\Afootnote{H; \skt{नाना\appbrk{}व\appbrk{}र्ण\appbrk{}वे\appbrk{}द\appbrk{}ना\appbrk{}स्रा\appbrk{}व\appbrk{}वि\appbrk{}शे\appbrk{}षो\appbrk{}पेतः}\ A}}\skt{ }\edtext{\skt{पव\-न\-पि\-त्त\-क\-फ\-शो\-णि\-तेभ्य}}{\lemma{\skt{पव\appbrk{}न\appbrk{}पि\appbrk{}त्त\appbrk{}क\appbrk{}फ\appbrk{}शो\appbrk{}णि\appbrk{}तेभ्य}}\Afootnote{H; \skt{पव\appbrk{}न\appbrk{}पि\appbrk{}त्त\appbrk{}क\appbrk{}फ\appbrk{}शो\appbrk{}णि\appbrk{}तेभ्यः}\ A}}\skt{ }\edtext{\skt{इति}}{\lemma{\skt{इति}}\Afootnote{H; \emph{om.}\ A}}\skt{~॥ जिह्वातलाभो }\edtext{\skt{मृदु}}{\lemma{\skt{मृदु}}\Afootnote{H; \skt{मृदुः}\ A}}\skt{ }\edtext{\skt{श्लक्ष्णः}}{\lemma{\skt{श्लक्ष्णः}}\Afootnote{H; \skt{स्निग्धः श्लक्ष्णो}\ A}}\skt{ }\edtext{\skt{स्निग्धो}}{\lemma{\skt{स्निग्धो}}\Afootnote{H; \emph{om.}\ A}}\skt{ विगतवेदनः सुव्यवस्थितो निरास्रावश्चेति शुद्धो }\edtext{\skt{व्रणः}}{\lemma{\skt{व्रणः}}\Afootnote{H; \skt{व्रण इति}\ A}}\skt{~।}\pend

\pstart \ledsidenote{\passref{4.1.8}}\edtext{\skt{तस्यैतस्य}}{\lemma{\skt{तस्यैतस्य}}\Afootnote{H; \skt{तस्य}\ A}}\skt{ व्रणस्य }\edtext{\skt{षष्टिरुपक्रमा}}{\lemma{\skt{षष्टिरुपक्रमा}}\Afootnote{H; \skt{पष्टिरुपक्रमा}\ A}}\skt{ भवन्ति~। }\edtext{\edtext{\skt{तद्य\-था\-प\-त\-र्प\-ण\-मा\-ले\-पनः}}{\lemma{\skt{यथा\appbrk{}प\appbrk{}त\appbrk{}र्प\appbrk{}णम्}}\Afootnote{H; \skt{यथा अपतर्पणम्}\ A}}}{\lemma{\skt{आलेपनः}}\Afootnote{H; \skt{आलेपः}\ A}}\skt{ परिषेकोऽभ्यङ्गः स्वेदो }\edtext{\skt{विम्ला\-य\-न\-मु\-प\-नाहो}}{\lemma{\skt{विम्लायनम्}}\Afootnote{H; \skt{विम्पालनम्}\ A}}\skt{ }\edtext{\skt{विस्रावणं}}{\lemma{\skt{विस्रावणं}}\Afootnote{H; \skt{पाचनं विस्रावणं}\ A}}\skt{ स्नेहो वम\-न\-म्वि\-रे\-चनं~। }\edtext{\skt{च्छेद्यं}}{\lemma{\skt{च्छेद्यं}}\Afootnote{H; \skt{छेदनं}\ A}}\skt{ }\edtext{\skt{भेद्यन्दारणं}}{\lemma{\skt{भेद्यन्दारणं}}\Afootnote{H; \skt{भेदनं दारणं}\ A}}\skt{ }\edtext{\skt{लेख्य}}{\lemma{\skt{लेख्य}}\Afootnote{H; \skt{लेखनम्}\ A}}\skt{ }\edtext{\skt{मेष्यमाहार्यं}}{\lemma{\skt{मेष्यमाहार्यं}}\Afootnote{H; \skt{एष\appbrk{}ण\appbrk{}मा\appbrk{}ह\appbrk{}रणं}\ A}}\skt{ }\edtext{\skt{व्यध्यं}}{\lemma{\skt{व्यध्यं}}\Afootnote{H; \skt{व्यधनं}\ A}}\skt{ }\edtext{\skt{सेव्यं}}{\lemma{\skt{सेव्यं}}\Afootnote{H; \skt{विस्रावणं}\ A}}\skt{ }\edtext{\skt{पीडनं}}{\lemma{\skt{पीडनं}}\Afootnote{H; \skt{सीवनं}\ A}}\skt{ }\edtext{\skt{सन्धानम्विस्राव्यं}}{\lemma{\skt{विस्राव्यं}}\Afootnote{H; \skt{पीडनं}\ A}}\skt{ }\edtext{\skt{शोणि\-त\-स्था\-प\-न\-न्नि\-र्व्वा\-पण}}{\lemma{\skt{शोणि\appbrk{}त\appbrk{}स्था\appbrk{}प\appbrk{}न\appbrk{}न्नि\appbrk{}र्व्वा\appbrk{}पण}}\Afootnote{H; \skt{शोणितास्थापनं निर्वापणम्}\ A}}\skt{ }\edtext{\skt{मूत्कारिका}}{\lemma{\skt{मूत्कारिका}}\Afootnote{H; \skt{उत्कारिका}\ A}}\skt{ कषायो वर्त्तिः कल्कः }\edtext{\skt{सर्प्पि\-स्तै\-ल\-म्पि\-ण्डी\-क\-रणं}}{\lemma{\skt{पिण्डीकरणं}}\Afootnote{H; \emph{om.}\ A}}\skt{ }\edtext{\edtext{\skt{रस\-क्रि\-या\-चू\-र्ण्ण\-म्व्र\-ण\-धू\-प\-न\-मु\-त्सा\-द\-न\-म\-व\-सा\-दनं}}{\lemma{\skt{रस\appbrk{}क्रि\appbrk{}या\appbrk{}चू\appbrk{}र्ण्णम्}}\Afootnote{H; \skt{रस\appbrk{}क्रि\appbrk{}याऽ\appbrk{}व\appbrk{}चू\appbrk{}र्णनं}\ A}}}{\lemma{\skt{व्रण\appbrk{}धू\appbrk{}प\appbrk{}न\appbrk{}मु\appbrk{}त्सा\appbrk{}द\appbrk{}न\appbrk{}म\appbrk{}व\appbrk{}सा\appbrk{}दनं}}\Afootnote{H; \skt{व्रण\appbrk{}धू\appbrk{}प\appbrk{}न\appbrk{}मु\appbrk{}त्सा\appbrk{}द\appbrk{}न\appbrk{}म\appbrk{}व\appbrk{}सा\appbrk{}दनं}\ A}}\skt{ मृदुकर्म्म दारुणकर्म्म क्षार\-क\-र्म्मा\-ग्नि\-कर्म्म कृष्णकर्म्म }\edtext{\skt{पाण्डूकर्म्म}}{\lemma{\skt{पाण्डूकर्म्म}}\Afootnote{H\textsuperscript{pc}; \skt{पाण्डुकर्म}\ A; \skt{पाण्डू}\ H\textsuperscript{ac}}}\skt{ }\edtext{\skt{सम\-व\-र्ण्ण\-लो\-म\-स\-ञ्ज\-ननं}}{\lemma{\skt{सम\appbrk{}व\appbrk{}र्ण्ण\appbrk{}लो\appbrk{}म\appbrk{}स\appbrk{}ञ्ज\appbrk{}ननं}}\Afootnote{H; \skt{प्रतिसारणं रोमसञ्जननं}\ A}}\skt{ }\edtext{\skt{लोमसंहरणं}}{\lemma{\skt{लोमसंहरणं}}\Afootnote{H; \skt{लोमापहरणं}\ A}}\skt{ बस्ति\-क\-र्म्मो\-त्त\-र\-ब\-स्ति\-कर्म्म }\edtext{\skt{वन्धेन}}{\lemma{\skt{वन्धेन}}\Afootnote{H; \skt{बन्धः}\ A}}\skt{ पत्रदानं }\edtext{\skt{क्रिमिघ्नं}}{\lemma{\skt{क्रिमिघ्नं}}\Afootnote{H; \skt{कृमिघ्नं}\ A}}\skt{ वृंहणम्विषघ्नं शिरो\-वि\-रे\-च\-न\-न्नस्यं कवलधारणं }\edtext{\skt{धूपो}}{\lemma{\skt{धूपो}}\Afootnote{H; \skt{धूमो}\ A}}\skt{ }\edtext{\skt{मधु\-स\-र्प्पि\-र्य\-न्त्र\-मा\-हारो}}{\lemma{\skt{मधु\appbrk{}स\appbrk{}र्प्पि\appbrk{}र्य\appbrk{}न्त्रम्}}\Afootnote{H; \skt{मधु\appbrk{}सा\appbrk{}र्पि\appbrk{}र्य\appbrk{}न्त्रम्}\ A}}\skt{ रक्षा\-वि\-धा\-न\-मिति}\pend

\pstart \ledsidenote{\passref{4.1.9}}\skt{तेषु कषायो वर्त्तिः कल्कः सर्प्पिस्तैलं }\edtext{\skt{रस\-क्रि\-या\-चू\-र्ण्ण\-मिति}}{\lemma{\skt{रस\appbrk{}क्रि\appbrk{}या\appbrk{}चू\appbrk{}र्ण्णम्}}\Afootnote{H; \skt{रस\appbrk{}क्रि\appbrk{}याऽ\appbrk{}व\appbrk{}चू\appbrk{}र्ण\appbrk{}नम्}\ A}}\skt{~। शोध\-न\-रो\-प\-णानि~। }\edtext{\edtext{\skt{एवमेते}}{\lemma{\skt{एवम्}}\Afootnote{H; \emph{om.}\ A}}}{\lemma{\skt{एते}}\Afootnote{H; \emph{om.}\ A}}\skt{ }\edtext{\skt{षष्टिरुपक्रमा}}{\lemma{\skt{षष्टिरुपक्रमा}}\Afootnote{H; \emph{om.}\ A}}\skt{ }\edtext{\skt{भवन्ति}}{\lemma{\skt{भवन्ति}}\Afootnote{H; \emph{om.}\ A}}\skt{~॥ तेष्वष्टौ शस्त्रकृत्याः }\edtext{\skt{शोणितस्थापनं}}{\lemma{\skt{शोणितस्थापनं}}\Afootnote{H; \skt{शोणितास्थापनं}\ A}}\skt{ }\edtext{\edtext{\skt{क्षारा\-ग्नि\-र्य\-न्त्र\-मा\-हारो}}{\lemma{\skt{क्षाराग्निर्यन्त्रम्}}\Afootnote{H; \skt{क्षारोऽग्निर्यन्त्रम्}\ A}}}{\lemma{\skt{आहारो}}\Afootnote{H; \skt{अहारो}\ A}}\skt{ }\edtext{\edtext{\skt{रक्षा\-वि\-धा\-न\-मि\-त्युक्ताः}}{\lemma{\skt{इत्य्}}\Afootnote{H; \skt{बन्धविधानं}\ A}}}{\lemma{\skt{उक्ताः}}\Afootnote{H; \skt{चोक्तानि}\ A}}\skt{~। }\edtext{\edtext{\skt{स्नेह\-स्वे\-द\-व\-म\-न\-वि\-रे\-च\-न\-व\-स्त्यु\-त्त\-र\-ब\-स्ति\-वि\-धा\-न\-शि\-रो\-वि\-रे\-च\-न\-र\-क्षा\-ग्नि\-न\-स्य\-धू\-म\-न\-क\-व\-ल\-वि\-धा\-ना\-न्यु\-त्त\-रत्र}}{\lemma{\skt{स्नेह\appbrk{}स्वे\appbrk{}द\appbrk{}व\appbrk{}म\appbrk{}न\appbrk{}वि\appbrk{}रे\appbrk{}च\appbrk{}न\appbrk{}व\appbrk{}स्त्यु\appbrk{}त्त\appbrk{}र\appbrk{}ब\appbrk{}स्ति\appbrk{}वि\appbrk{}धा\appbrk{}न\appbrk{}शि\appbrk{}रो\appbrk{}वि\appbrk{}रे\appbrk{}च\appbrk{}न\appbrk{}र\appbrk{}क्षा\appbrk{}ग्नि\appbrk{}न\appbrk{}स्य\appbrk{}धू\appbrk{}म\appbrk{}न\appbrk{}क\appbrk{}व\appbrk{}ल\appbrk{}वि\appbrk{}धा\appbrk{}नान्य्}}\Afootnote{H; \skt{स्नेह\appbrk{}स्वे\appbrk{}द\appbrk{}न\appbrk{}व\appbrk{}म\appbrk{}न\appbrk{}वि\appbrk{}रे\appbrk{}च\appbrk{}न\appbrk{}ब\appbrk{}स्त्यु\appbrk{}त्त\appbrk{}र\appbrk{}ब\appbrk{}स्ति\appbrk{}शि\appbrk{}रो\appbrk{}वि\appbrk{}रे\appbrk{}च\appbrk{}न\appbrk{}न\appbrk{}स्य\appbrk{}धू\appbrk{}म\appbrk{}क\appbrk{}व\appbrk{}ल\appbrk{}घा\appbrk{}र\appbrk{}णान्}\ A}}}{\lemma{\skt{उत्तरत्र}}\Afootnote{H; \skt{यन्यत्र}\ A}}\skt{ वक्ष्यामः~। }\edtext{\skt{यद\-न्य\-द\-व\-सि\-ष्ट\-मु\-प\-क्र\-म\-जातं}}{\lemma{\skt{यद\appbrk{}न्य\appbrk{}द\appbrk{}व\appbrk{}सि\appbrk{}ष्टम्}}\Afootnote{H; \skt{यद\appbrk{}न्य\appbrk{}द\appbrk{}व\appbrk{}शि\appbrk{}ष्टम्}\ A}}\skt{ }\edtext{\skt{तदत्र}}{\lemma{\skt{तदत्र}}\Afootnote{\skt{तदिह}\ A; \skt{जातदत्र}\ H\textsuperscript{ac}; \skt{तदत्र}\ H\textsuperscript{pc}}}\skt{ वक्ष्यते~।}\pend

\pstart \ledsidenote{\passref{4.1.10}}\edtext{\skt{तत्र}}{\lemma{\skt{तत्र}}\Afootnote{H; \emph{om.}\ A}}\skt{ षड्विधः }\edtext{\skt{प्रादिष्टः}}{\lemma{\skt{प्रादिष्टः}}\Afootnote{H; \skt{प्रागुपदिष्टः}\ A}}\skt{ }\edtext{\skt{शोफ\-स्त\-स्यै\-का\-द\-शो\-प\-क्रमा}}{\lemma{\skt{शोफ\appbrk{}स्त\appbrk{}स्यै\appbrk{}का\appbrk{}द\appbrk{}शो\appbrk{}प\appbrk{}क्रमा}}\Afootnote{H; \skt{शोफः, तस्यै\appbrk{}का\appbrk{}द\appbrk{}शो\appbrk{}प\appbrk{}क्रमा}\ A}}\skt{ भवन्ति~। अप\-त\-र्प्प\-णा\-दयो }\edtext{\skt{विरेचनान्ताः}}{\lemma{\skt{विरेचनान्ताः}}\Afootnote{H; \skt{विरेचनान्ताः; ते च}\ A}}\skt{~। विशेषेण }\edtext{\skt{व्रणप्रतीकारे}}{\lemma{\skt{व्रणप्रतीकारे}}\Afootnote{H; \skt{शोथप्रतीकारे}\ A}}\skt{ वर्त्तन्ते }\edtext{\skt{व्रण\-भा\-व\-मु\-प\-ग\-तस्य न}}{\lemma{\skt{उपगतस्य न}}\Afootnote{H; \skt{आपन्नस्य च न}\ A}}\skt{ विरुध्यन्ते शेषास्तु प्रायेण व्रण\-प्र\-ती\-का\-र\-हे\-तव एव~।}\pend

\pstart \ledsidenote{\passref{4.1.11}}\skt{अपतर्प्पण }\edtext{\skt{आद्य}}{\lemma{\skt{आद्य}}\Afootnote{H; \skt{माद्य}\ A}}\skt{ उपक्रम एष सर्व्वशोफानां }\edtext{\skt{सामान्यतः}}{\lemma{\skt{सामान्यतः}}\Afootnote{H; \skt{सामान्यः}\ A}}\skt{ प्रधानतमश्च~॥ }\edtext{\skt{भ}}{\lemma{\skt{भ}}\Afootnote{H; \emph{om.}\ A}}\skt{~॥}\pend

\pstart\versepar \ledsidenote{\passref{4.1.12}}\edtext{\skt{दोषो\-च्छ्रा\-यो\-प\-शा\-न्त्य\-र्थ\-न्दो\-षा\-न\-द्धस्य}}{\lemma{\skt{दोषो\appbrk{}च्छ्रा\appbrk{}यो\appbrk{}प\appbrk{}शा\appbrk{}न्त्य\appbrk{}र्थन्}}\Afootnote{H; \skt{भवन्ति चात्र दोषो\appbrk{}च्छ्रा\appbrk{}यो\appbrk{}प\appbrk{}शा\appbrk{}न्त्यर्थं}\ A}}\skt{ देहिनः~। }\pend
\pstart\versepar \skt{अवेक्ष्य दोषं प्राणञ्च कार्यं }\edtext{\skt{स्यादपतर्प्पणं}}{\lemma{\skt{स्यादपतर्प्पणं}}\Afootnote{A; \skt{स्यादपतर्प्पणं}\ H}}\skt{~॥}\pend

\pstart\versepar \ledsidenote{\passref{4.1.13}}\edtext{\skt{तद्धि}}{\lemma{\skt{तद्धि}}\Afootnote{H; \skt{ऊर्ध्व}\ A}}\skt{ }\edtext{\skt{मारुतक्षुत्तृष्णा}}{\lemma{\skt{मारुतक्षुत्तृष्णा}}\Afootnote{H; \skt{मारुततृष्णा}\ A}}\skt{ }\edtext{\skt{मुख\-शो\-ष\-स्र\-मा\-न्विते}}{\lemma{\skt{मुख\appbrk{}शो\appbrk{}ष\appbrk{}स्र\appbrk{}मा\appbrk{}न्विते}}\Afootnote{H; \skt{क्षुन्मु\appbrk{}ख\appbrk{}शो\appbrk{}ष\appbrk{}श्र\appbrk{}मा\appbrk{}न्वितैः~। न}\ A}}\skt{~। }\pend
\pstart\versepar \skt{}\edtext{\skt{कार्यं न}}{\lemma{\skt{कार्यं न}}\Afootnote{H; \skt{कार्यं}\ A}}\skt{ }\edtext{\skt{बालवृद्धे वा न}}{\lemma{\skt{बालवृद्धे वा न}}\Afootnote{H; \emph{om.}\ A}}\skt{ }\edtext{\skt{गर्ब्भिण्यां न}}{\lemma{\skt{गर्ब्भिण्यां न}}\Afootnote{H; \skt{गर्भि\appbrk{}णी\appbrk{}वृ\appbrk{}द्ध\appbrk{}बाल}\ A}}\skt{ }\edtext{\skt{दुर्व्वले}}{\lemma{\skt{दुर्व्वले}}\Afootnote{H; \skt{दुर्बलभीरुभिः ध्यामे सरक्ते वेदनावति}\ A}}\skt{~॥}\pend

\pstart\versepar \ledsidenote{\passref{4.1.14}}\skt{शोफे\-षू\-त्थि\-त\-मा\-त्रेषु व्रणे\-षू\-ग्र\-रु\-जेषु च~। }\pend
\pstart\versepar \skt{}\edtext{\edtext{\edtext{\skt{यथो\-क्तै\-रौ\-ष\-धै\-र्ल्लेपं}}{\lemma{\skt{यथोक्तैर्}}\Afootnote{\skt{यथास्वैर्}\ A; \skt{यथोक्तै}\ H}}}{\lemma{\skt{औषधैर्}}\Afootnote{A; \skt{रौषधै}\ H}}}{\lemma{\skt{ल्लेपं}}\Afootnote{A; \skt{र्ल्लेपं}\ H}}\skt{ }\edtext{\skt{प्रत्येकांश्येन}}{\lemma{\skt{प्रत्येकांश्येन}}\Afootnote{H; \skt{प्रत्येकश्येन}\ A}}\skt{ कारयेत्~॥}\pend

\pstart\versepar \ledsidenote{\passref{4.1.15}}\skt{यथा }\edtext{\skt{प्रज्वलिते}}{\lemma{\skt{प्रज्वलिते}}\Afootnote{H; \skt{प्रवज्वलिते}\ A}}\skt{ }\edtext{\skt{वेश्मन्यम्भसा}}{\lemma{\skt{वेश्मन्यम्भसा}}\Afootnote{\skt{वेश्मन्यम्भसा}\ A H}}\skt{ }\edtext{\skt{परिषेचनं}}{\lemma{\skt{परिषेचनं}}\Afootnote{H; \skt{परिबेचनम्}\ A}}\skt{~। }\pend
\pstart\versepar \skt{क्षिप्रं }\edtext{\skt{प्रश\-म\-य\-त्य\-ग्नि\-मे\-व\-मा\-ले\-पनं}}{\lemma{\skt{प्रश\appbrk{}म\appbrk{}य\appbrk{}त्य\appbrk{}ग्नि\appbrk{}मे\appbrk{}वम्}}\Afootnote{\skt{प्रश\appbrk{}म\appbrk{}य\appbrk{}त्य\appbrk{}ग्नि\appbrk{}मे\appbrk{}वम्}\ A H}}\skt{ रुजः~॥}\pend

\pstart\versepar \ledsidenote{\passref{4.1.16}}\skt{प्रह्लादने शोधने च शोफस्य हरणे तथा~। }\pend
\pstart\versepar \skt{उत्सादने }\edtext{\skt{रोपने च}}{\lemma{\skt{रोपने च}}\Afootnote{H; \skt{रोपणे च}\ A}}\skt{ लेपः स्यात्तु तदर्थकृत्~॥}\pend

\pstart \ledsidenote{\passref{4.1.17}}\edtext{\skt{वातशेफे तु}}{\lemma{\skt{वातशेफे तु}}\Afootnote{H; \skt{वातशोफे तु}\ A}}\skt{ वेद\-नो\-प\-श\-मार्थं }\edtext{\edtext{\skt{सर्प्पि\-स्तै\-ल\-धा\-न्या\-म्ल\-मां\-स\-र\-स\-वा\-त\-ह\-रौ\-ष\-ध\-निः\-क्वा\-थै\-र\-शीतैः}}{\lemma{\skt{सर्प्पि\appbrk{}स्तै\appbrk{}ल\appbrk{}धा\appbrk{}न्या\appbrk{}म्ल\appbrk{}मां\appbrk{}स\appbrk{}र\appbrk{}स\appbrk{}वा\appbrk{}त\appbrk{}ह\appbrk{}रौ\appbrk{}ष\appbrk{}ध\appbrk{}निः\appbrk{}क्वा\appbrk{}थैर्}}\Afootnote{A; \skt{सर्प्पि\appbrk{}स्तै\appbrk{}ल\appbrk{}धा\appbrk{}न्या\appbrk{}म्ल\appbrk{}मां\appbrk{}स\appbrk{}र\appbrk{}स\appbrk{}वा\appbrk{}त\appbrk{}ह\appbrk{}रौ\appbrk{}ष\appbrk{}ध\appbrk{}निः\appbrk{}क्वाथै}\ H}}}{\lemma{\skt{अशीतैः}}\Afootnote{A; \skt{रशीतैः}\ H}}\skt{ }\edtext{\skt{परिषेकं}}{\lemma{\skt{परिषेकं}}\Afootnote{H; \skt{परिषेकान्}\ A}}\skt{ कुर्व्वीत~। पित्त\-र\-क्ता\-भि\-घा\-त\-वि\-ष\-नि\-मि\-त्तेषु }\edtext{\skt{क्षीर\-वृ\-क्ष\-क्षी\-र\-श\-र्क्क\-रो\-द\-क\-के\-क्षु\-र\-स\-घृ\-त\-म\-धु\-रौ\-ष\-धनिः }\edtext{\skt{क्वाथैरनुष्णैः}}{\lemma{\skt{अनुष्णैः}}\Afootnote{A; \skt{रनुष्णैः}\ H}}}{\lemma{\skt{क्षीर\appbrk{}वृ\appbrk{}क्ष\appbrk{}क्षी\appbrk{}र\appbrk{}श\appbrk{}र्क्क\appbrk{}रो\appbrk{}द\appbrk{}क\appbrk{}के\appbrk{}क्षु\appbrk{}र\appbrk{}स\appbrk{}घृ\appbrk{}त\appbrk{}म\appbrk{}धु\appbrk{}रौ\appbrk{}ष\appbrk{}धनिः क्वाथैर्}}\Afootnote{\skt{क्षीर\appbrk{}घृ\appbrk{}त\appbrk{}म\appbrk{}धु\appbrk{}श\appbrk{}र्क\appbrk{}रो\appbrk{}द\appbrk{}के\appbrk{}क्षु\appbrk{}र\appbrk{}स\appbrk{}म\appbrk{}धु\appbrk{}रौ\appbrk{}ष\appbrk{}ध\appbrk{}क्षी\appbrk{}र\appbrk{}वृ\appbrk{}क्ष\appbrk{}नि\appbrk{}ष्क्वा\appbrk{}थैर्}\ A; \skt{क्षीर\appbrk{}वृ\appbrk{}क्ष\appbrk{}क्षी\appbrk{}र\appbrk{}श\appbrk{}र्क्क\appbrk{}रो\appbrk{}द\appbrk{}क\appbrk{}के\appbrk{}क्षु\appbrk{}र\appbrk{}स\appbrk{}घृ\appbrk{}त\appbrk{}म\appbrk{}धु\appbrk{}रौ\appbrk{}ष\appbrk{}ध\appbrk{}निः\appbrk{}क्वाथै}\ H}}\skt{~। }\edtext{\skt{श्लेष्मशोफे तु}}{\lemma{\skt{श्लेष्मशोफे तु}}\Afootnote{H; \skt{परि\appbrk{}षे\appbrk{}का\appbrk{}न्कु\appbrk{}र्वीत, श्लेष्मशोफे तु}\ A}}\skt{ }\edtext{\edtext{\skt{तैल\-मू\-त्र\-क्षा\-रो\-द\-क\-सु\-रा\-शु\-क्त\-क\-फ\-घ्नौ\-ष\-ध\-निः\-क्वा\-थै\-र\-शी\-तैर्}}{\lemma{\skt{तैल\appbrk{}मू\appbrk{}त्र\appbrk{}क्षा\appbrk{}रो\appbrk{}द\appbrk{}क\appbrk{}सु\appbrk{}रा\appbrk{}शु\appbrk{}क्त\appbrk{}क\appbrk{}फ\appbrk{}घ्नौ\appbrk{}ष\appbrk{}ध\appbrk{}निः\appbrk{}क्वा\appbrk{}थैर्}}\Afootnote{A; \skt{तैल\appbrk{}मू\appbrk{}त्र\appbrk{}क्षा\appbrk{}रो\appbrk{}द\appbrk{}क\appbrk{}सु\appbrk{}रा\appbrk{}शु\appbrk{}क्त\appbrk{}क\appbrk{}फ\appbrk{}घ्नौ\appbrk{}ष\appbrk{}ध\appbrk{}निः\appbrk{}क्वाथै}\ H}}}{\lemma{\skt{अशीतैर्}}\Afootnote{\skt{अशीतैः परि\appbrk{}षे\appbrk{}का\appbrk{}न्कु\appbrk{}र्वीत}\ A; \skt{रशीतैर्}\ H}}\pend

\pstart\versepar \ledsidenote{\passref{4.1.18}}\edtext{\skt{यथाम्वुभिः}}{\lemma{\skt{यथाम्वुभिः}}\Afootnote{H; \skt{यथा+अम्बुभिः}\ A}}\skt{ सिच्यमानः }\edtext{\skt{शान्ति\-म\-ग्नि\-र्न्नि\-य\-च्छति}}{\lemma{\skt{शान्ति\appbrk{}म\appbrk{}ग्नि\appbrk{}र्न्नि\appbrk{}य\appbrk{}च्छति}}\Afootnote{\skt{शान्ति\appbrk{}म\appbrk{}ग्नि\appbrk{}र्नि\appbrk{}य\appbrk{}च्छति}\ A H}}\skt{~। }\pend
\pstart\versepar \skt{}\edtext{\skt{दोषाग्निरेवं}}{\lemma{\skt{दोषाग्निरेवं}}\Afootnote{\skt{दोषाग्निरेवं}\ A H}}\skt{ सहसा परिषेकेण शाम्यति~॥}\pend

\pstart \ledsidenote{\passref{4.1.19}}\skt{अभ्यङ्गस्तु }\edtext{\edtext{\skt{दोष\-मा\-लो\-च्यो\-प\-युक्तो}}{\lemma{\skt{दोषम्}}\Afootnote{H; \skt{दोष}\ A}}}{\lemma{\skt{आलोच्योपयुक्तो}}\Afootnote{H; \skt{मालोक्योपयुक्तो}\ A}}\skt{ }\edtext{\skt{दोषोपशमो}}{\lemma{\skt{दोषोपशमो}}\Afootnote{H; \skt{दोषोपशमं}\ A}}\skt{ मृदुतां च करोति~।}\pend

\pstart\versepar \ledsidenote{\passref{4.1.20}}\skt{स्वेद\-वि\-म्ला\-प\-ना\-दीनां क्रियाणां प्राक्स उच्यते~। }\pend
\pstart\versepar \skt{पश्चात्कर्म्मसु }\edtext{\skt{चादिष्टः स तु}}{\lemma{\skt{चादिष्टः स तु}}\Afootnote{H; \skt{चादिष्टः स च}\ A}}\skt{ विस्रावणादिषु~॥}\edtext{}{\lemma{\skt{उच्यते~। पश्}}\Afootnote{H; \skt{उच्यते। पञ्}\ A}}\pend

\pstart\versepar \ledsidenote{\passref{4.1.21}}\skt{रुजा\-व\-ता\-न्दा\-रु\-णा\-ना\-ङ्क\-ठि\-ना\-ना\-न्त\-थैव च~। }\pend
\pstart\versepar \skt{शोफानां स्वेदनं कार्यं ये }\edtext{\skt{चाप्येवम्विधा}}{\lemma{\skt{चाप्येवम्विधा}}\Afootnote{\skt{चाप्येवंविधा}\ A H}}\skt{ व्रणाः~॥}\pend

\pstart\versepar \ledsidenote{\passref{4.1.22}}\skt{स्थिराणां रुजताम्मन्दं कार्यं विम्ला\-प\-न\-म्भ\-वेत्~। }\pend
\pstart\versepar \skt{अभ्यज्य }\edtext{\skt{स्वेदयित्वा च}}{\lemma{\skt{स्वेदयित्वा च}}\Afootnote{H; \skt{स्वेदयित्वा तु}\ A}}\skt{ }\edtext{\skt{वेणुनाड्या}}{\lemma{\skt{वेणुनाड्या}}\Afootnote{H; \skt{वेशुनाड्या}\ A}}\skt{ ततः शनैः~॥}\pend

\pstart\versepar \ledsidenote{\passref{4.1.23}}\edtext{\skt{विम\-र्द्द\-ये\-द्भि\-ष\-क्प्रा\-ज्ञ\-स्त\-ले\-ना\-ङ्गु\-ष्ठ\-केन वा}}{\lemma{\skt{प्राज्ञ\appbrk{}स्त\appbrk{}ले\appbrk{}ना\appbrk{}ङ्गु\appbrk{}ष्ठ\appbrk{}केन वा}}\Afootnote{A; \skt{प्राज्ञ\appbrk{}स्त\appbrk{}ले\appbrk{}ना\appbrk{}ङ्गु\appbrk{}ष्ठ\appbrk{}केन वा}\ H}}\skt{~। }\pend
\pstart\versepar \skt{शोफ\-यो\-रु\-प\-ना\-हन्तु कुर्या\-दा\-म\-वि\-द\-ग्धयोः~॥}\pend

\pstart\versepar \ledsidenote{\passref{4.1.24}}\edtext{\edtext{\skt{प्रशा\-म्य\-त्य\-वि\-द\-ग्धस्तु}}{\lemma{\skt{प्रशाम्यत्य्}}\Afootnote{H; \emph{om.}\ A}}}{\lemma{\skt{तु}}\Afootnote{H; \skt{शमं याति}\ A}}\skt{ विदग्धः पाकमेति च~। }\pend
\pstart\versepar \skt{निवर्त्तते न यः शोफो विरे\-का\-न्तै\-रु\-प\-क्रमैः~॥}\pend

\pstart\versepar \ledsidenote{\passref{4.1.25}}\edtext{\skt{तत्र}}{\lemma{\skt{तत्र}}\Afootnote{H; \skt{तस्य}\ A}}\skt{ }\edtext{\edtext{\skt{संपा\-च\-न\-द्र\-व्या\-ण्या\-हृ\-त्यो\-क्तानि}}{\lemma{\skt{संपा\appbrk{}च\appbrk{}न\appbrk{}द्र\appbrk{}व्याण्य्}}\Afootnote{H; \skt{संपाचनं कुर्यात्स}\ A}}}{\lemma{\skt{आहृत्योक्तानि}}\Afootnote{H; \skt{माहृत्यौषधानि}\ A}}\skt{ }\edtext{\skt{पिव्य च}}{\lemma{\skt{पिव्य च}}\Afootnote{H; \skt{तु}\ A}}\skt{~। }\pend
\pstart\versepar \skt{}\edtext{\edtext{\skt{दधि\-त\-क्र\-सु\-रा\-शु\-क्त\-धा\-न्या\-म्लै\-र्यो\-ज\-ये\-द्भि\-षक्}}{\lemma{\skt{योजयेद्}}\Afootnote{H; \skt{यो}\ A}}}{\lemma{\skt{भिषक्}}\Afootnote{H; \skt{जितानि तु}\ A}}\skt{~॥}\pend

\pstart\versepar \ledsidenote{\passref{4.1.26}}\skt{स्निग्धानि लवणी }\edtext{\skt{कृत्वा}}{\lemma{\skt{कृत्वा}}\Afootnote{H; \skt{कृत्य}\ A}}\skt{ }\edtext{\skt{पचे\-दु\-त्का\-रि\-का\-न्ततः}}{\lemma{\skt{ततः}}\Afootnote{H; \skt{शुभाम्}\ A}}\skt{~। }\pend
\pstart\versepar \skt{सैरण्डपत्रया शोफ\-न्ना\-ह\-ये\-दु\-ष्णया तया~॥}\pend

\pstart\versepar \ledsidenote{\passref{4.1.27}}\skt{हितं संभोजनञ्चापि }\edtext{\skt{पाकाय}}{\lemma{\skt{पाकाय}}\Afootnote{H; \skt{पाकाया}\ A}}\skt{ }\edtext{\skt{विमुखो}}{\lemma{\skt{विमुखो}}\Afootnote{H; \skt{भिमुखो}\ A}}\skt{ यदि~। }\pend
\pstart\versepar \skt{वेद\-नो\-प\-श\-मा\-र्थाय तथा }\edtext{\skt{पाकभयाय}}{\lemma{\skt{पाकभयाय}}\Afootnote{H; \skt{पाकशमाय}\ A}}\skt{ च~॥}\pend

\pstart\versepar \ledsidenote{\passref{4.1.28}}\skt{अचिरोत्पतिते शोफे कुर्या\-च्छो\-णि\-त\-मो\-क्षणं~। }\pend
\pstart\versepar \skt{सशोफे कठिने ध्यामे सरक्ते वेदनावति~॥}\pend

\pstart \ledsidenote{\passref{4.1.29}}\skt{संरव्धे विषमे चापि व्रणे }\edtext{\skt{विश्रावणं}}{\lemma{\skt{विश्रावणं}}\Afootnote{H; \skt{विस्रावणं}\ A}}\skt{ }\edtext{\skt{हितं}}{\lemma{\skt{हितं}}\Afootnote{H; \skt{हितम्~। सविषे च विशेषेण जलौकोभिः पदैस्तथा}\ A}}\skt{~।}\pend

\pstart \ledsidenote{\passref{4.1.30}}\edtext{\skt{सोप\-द्र\-वा\-द्र\-वाणां}}{\lemma{\skt{सोप\appbrk{}द्र\appbrk{}वा\appbrk{}द्र\appbrk{}वाणां}}\Afootnote{H; \skt{वेदनायाः प्रशान्त्यर्थं पाकस्याप्राप्तये तथा~। सोपद्रवाणां}\ A}}\skt{ रूक्षाणाङ्कृशानां व्रणशोषिणां~॥}\pend

\pstart\versepar \ledsidenote{\passref{4.1.31}}\skt{यथा }\edtext{\skt{स्वौषधसंसिद्धं}}{\lemma{\skt{स्वौषधसंसिद्धं}}\Afootnote{H; \skt{स्वमौषधैः सिद्धं}\ A}}\skt{ स्नेह\-पा\-न\-म्वि\-धी\-यते~। }\pend
\pstart\versepar \skt{उत्स\-न्न\-मां\-स\-शोफे तु }\edtext{\skt{कफे}}{\lemma{\skt{कफे}}\Afootnote{H; \skt{कफ}\ A}}\skt{ }\edtext{\skt{दुष्टे}}{\lemma{\skt{दुष्टे}}\Afootnote{H; \skt{जुष्टे}\ A}}\skt{ विशेषतः~॥}\pend

\pstart\versepar \ledsidenote{\passref{4.1.32}}\edtext{\skt{संसृ\-ष्ट\-ध्या\-म\-रु\-धिरे}}{\lemma{\skt{संसृ\appbrk{}ष्ट\appbrk{}ध्या\appbrk{}म\appbrk{}रु\appbrk{}धिरे}}\Afootnote{H; \skt{संक्लि\appbrk{}ष्ट\appbrk{}श्या(\appbrk{}ध्या)\appbrk{}म\appbrk{}रु\appbrk{}धिरे}\ A}}\skt{ व्रणे प्रच्छर्द्दनं हितं~। }\pend
\pstart\versepar \skt{वात\-पि\-त्त\-प्र\-दु\-ष्टेषु दीर्घ\-का\-ला\-नु\-व\-न्धिषु~॥}\pend

\pstart\versepar \ledsidenote{\passref{4.1.33}}\skt{विरे\-च\-न\-म्प्र\-शं\-सन्ति व्रणेषु व्रणकोविदाः~। }\pend
\pstart\versepar \skt{}\edtext{\skt{अपा\-क\-व\-त्सु\-रो\-गेषु}}{\lemma{\skt{अपा\appbrk{}क\appbrk{}व\appbrk{}त्सु\appbrk{}रो\appbrk{}गेषु}}\Afootnote{H; \skt{अपाकेषु तु रोगेषु}\ A}}\skt{ कठिनेषु स्थिरेषु च~॥}\pend

\pstart\versepar \ledsidenote{\passref{4.1.34}}\skt{स्नायुकोथादिषु तथा च्छेदनं }\edtext{\skt{प्राप्तमुच्यते}}{\lemma{\skt{प्राप्तमुच्यते}}\Afootnote{\skt{प्राप्तमुच्यते}\ A H}}\skt{~। }\pend
\pstart\versepar \skt{अन्तः पूयेष्ववक्त्रेषु तथै\-वो\-त्स\-ङ्ग\-व\-त्स्वपि॥}\pend

\pstart\versepar \ledsidenote{\passref{4.1.35}}\skt{गतिमत्सु च रोगेषु }\edtext{\skt{भेद\-न\-म्प्रा\-प्त\-मु\-च्यते}}{\lemma{\skt{भेद\appbrk{}न\appbrk{}म्प्रा\appbrk{}प्त\appbrk{}मु\appbrk{}च्यते}}\Afootnote{H; \skt{भेदनं प्राप्तमुच्यते}\ A}}\skt{~। }\pend
\pstart\versepar \skt{}\edtext{\skt{वाल\-वृ\-द्ध\-क्ष\-त\-क्षी\-ण\-भी\-रूणां}}{\lemma{\skt{वाल\appbrk{}वृ\appbrk{}द्ध\appbrk{}क्ष\appbrk{}त\appbrk{}क्षी\appbrk{}ण\appbrk{}भी\appbrk{}रूणां}}\Afootnote{H; \skt{बाल\appbrk{}वृ\appbrk{}द्धा\appbrk{}स\appbrk{}ह\appbrk{}क्षी\appbrk{}ण\appbrk{}भी\appbrk{}रूणां}\ A}}\skt{ योषितामपि~॥}\pend

\pstart\versepar \ledsidenote{\passref{4.1.36}}\skt{मर्म्मोपरि च जातेषु }\edtext{\skt{रोगेषूक्तञ्च}}{\lemma{\skt{रोगेषूक्तञ्च}}\Afootnote{H; \skt{रोगेषूक्तेषु}\ A}}\skt{ }\edtext{\skt{दारुणं}}{\lemma{\skt{दारुणं}}\Afootnote{H; \skt{दारणम्}\ A}}\skt{~। }\pend
\pstart\versepar \skt{सुपक्वे पिण्डिते शोफे पीड\-नै\-रु\-प\-पी\-डिते~॥}\pend

\pstart\versepar \ledsidenote{\passref{4.1.37}}\edtext{\skt{पाकोद्धतेषु}}{\lemma{\skt{पाकोद्धतेषु}}\Afootnote{H; \skt{पाकोद्वृत्तेषु}\ A}}\skt{ दोषेषु तत्तु कार्यम्विजानता~। }\pend
\pstart\versepar \skt{}\edtext{\skt{सुपि\-ष्टै\-र्दा\-र\-ण\-द्रव्यैः}}{\lemma{\skt{सुपि\appbrk{}ष्टै\appbrk{}र्दा\appbrk{}र\appbrk{}ण\appbrk{}द्रव्यैः}}\Afootnote{H; \skt{सुपि\appbrk{}ष्टै\appbrk{}र्दा\appbrk{}र\appbrk{}ण\appbrk{}द्र\appbrk{}व्यैर्}\ A}}\skt{ }\edtext{\skt{ऋतक्षारण वा}}{\lemma{\skt{ऋतक्षारण वा}}\Afootnote{H; \skt{युक्तैः क्षारेण वा}\ A}}\skt{ पुनः~॥}\pend

\pstart\versepar \ledsidenote{\passref{4.1.38}}\edtext{\edtext{\skt{कठि\-न\-स्थू\-ल\-व\-त्तौ\-ष्ठा\-न्दी\-र्य\-मा\-ना\-न्पुनः}}{\lemma{\skt{कठि\appbrk{}न\appbrk{}स्थू\appbrk{}ल\appbrk{}व\appbrk{}त्तौ\appbrk{}ष्ठान्}}\Afootnote{H; \skt{कठि\appbrk{}ना\appbrk{}न्स्थू\appbrk{}ल\appbrk{}वृ\appbrk{}त्तौ\appbrk{}ष्ठान्}\ A}}}{\lemma{\skt{दीर्यमानान्}}\Afootnote{H; \skt{दीर्यमाणान्}\ A}}\skt{ पुनः~। }\pend
\pstart\versepar \skt{}\edtext{\skt{कठि\-ना\-त्स\-न्न\-मां\-साश्च}}{\lemma{\skt{कठि\appbrk{}ना\appbrk{}त्स\appbrk{}न्न\appbrk{}मां\appbrk{}साश्च}}\Afootnote{H; \skt{कठि\appbrk{}नो\appbrk{}त्स\appbrk{}न्न\appbrk{}मां\appbrk{}सांश्च}\ A}}\skt{ }\edtext{\skt{लेख\-नै\-रा\-च\-रे\-द्भि\-षक्}}{\lemma{\skt{लेख\appbrk{}नै\appbrk{}रा\appbrk{}च\appbrk{}रेद्}}\Afootnote{\skt{लेस्व\appbrk{}ने\appbrk{}ना\appbrk{}च\appbrk{}रेद्}\ A; \skt{लेख\appbrk{}नै\appbrk{}रा\appbrk{}च\appbrk{}रेद्}\ H}}\skt{~॥}\pend

\pstart\versepar \ledsidenote{\passref{4.1.39}}\skt{समं लिखेत्सुलिखितं लिखे\-न्नि\-र\-व\-शे\-षतः~। }\pend
\pstart\versepar \skt{}\edtext{\skt{चर्म्मणान्तु}}{\lemma{\skt{चर्म्मणान्तु}}\Afootnote{H; \skt{वर्त्मनां तु}\ A}}\skt{ प्रमाणेन समं शस्त्रेण }\edtext{\skt{संलिखेत्}}{\lemma{\skt{संलिखेत्}}\Afootnote{H; \skt{निर्लिखेत्}\ A}}\skt{~॥}\pend

\pstart\versepar \ledsidenote{\passref{4.1.40}}\skt{क्षौमं प्लोतं पिचुं फेनं यावशूकं ससैन्धवं~। }\pend
\pstart\versepar \skt{}\edtext{\skt{कर्कशानि च}}{\lemma{\skt{कर्कशानि च}}\Afootnote{H\textsuperscript{pc}; \skt{ककशानि च}\ A; \skt{कर्कशानि}\ H\textsuperscript{ac}}}\skt{ पत्राणि लेखनार्थे प्रदापयेत्~॥}\pend

\pstart\versepar \ledsidenote{\passref{4.1.41}}\edtext{\skt{नाडीव्रणं}}{\lemma{\skt{नाडीव्रणं}}\Afootnote{H; \skt{नाडीव्रणाञ्}\ A}}\skt{ }\edtext{\skt{शल्य\-ग\-र्ब्भ\-मु\-न्मा\-र्ग्यु\-त्सं\-गिनं}}{\lemma{\skt{शल्य\appbrk{}ग\appbrk{}र्ब्भ\appbrk{}मु\appbrk{}न्मा\appbrk{}र्ग्यु\appbrk{}त्सं\appbrk{}गिनं}}\Afootnote{\skt{शल्य\appbrk{}ग\appbrk{}र्भा\appbrk{}नु\appbrk{}न्मा\appbrk{}र्ग्यु\appbrk{}त्स\appbrk{}ङ्गिनः}\ A; \skt{शल्य\appbrk{}ग\appbrk{}र्व्भ\appbrk{}मु\appbrk{}न्मा\appbrk{}र्ग्यु\appbrk{}त्सं\appbrk{}गिनं}\ H}}\skt{ शनैः~। }\pend
\pstart\versepar \skt{करी\-र\-वा\-ला\-ङ्गु\-लि\-भि\-रे\-षण्या }\edtext{\skt{वेशयेद्भिषक्}}{\lemma{\skt{वेशयेद्}}\Afootnote{H; \skt{वैषयेद्}\ A}}\skt{~॥}\pend

\pstart\versepar \ledsidenote{\passref{4.1.42}}\edtext{\skt{नेत्र\-व\-र्म्म\-गु\-दा\-भ्या\-स\-गत्यो}}{\lemma{\skt{नेत्र\appbrk{}व\appbrk{}र्म्म\appbrk{}गु\appbrk{}दा\appbrk{}भ्या\appbrk{}स\appbrk{}गत्यो}}\Afootnote{H; \skt{नेत्र\appbrk{}व\appbrk{}र्त्म\appbrk{}गु\appbrk{}दा\appbrk{}भ्या\appbrk{}स\appbrk{}नाड्यो}\ A}}\skt{ }\edtext{\skt{वक्त्राः}}{\lemma{\skt{वक्त्राः}}\Afootnote{H; \skt{ऽवक्त्राः}\ A}}\skt{ सशोणिताः~। }\pend
\pstart\versepar \skt{}\edtext{\skt{चुच्चूपोतकजैः}}{\lemma{\skt{चुच्चूपोतकजैः}}\Afootnote{H; \skt{चुच्चूपोदकजैः}\ A}}\skt{ श्लक्ष्णैः }\edtext{\skt{करी\-रै\-रे\-ष\-ये\-द्भि\-षक्}}{\lemma{\skt{भिषक्}}\Afootnote{H; \skt{तु ताः}\ A}}\skt{~॥}\pend

\pstart\versepar \ledsidenote{\passref{4.1.43}}\edtext{\skt{संवृ\-ता\-वि\-वृ\-ता\-स्येषु}}{\lemma{\skt{संवृ\appbrk{}ता\appbrk{}वि\appbrk{}वृ\appbrk{}ता\appbrk{}स्येषु}}\Afootnote{H; \skt{संवृ\appbrk{}ता\appbrk{}सं\appbrk{}वृ\appbrk{}ता\appbrk{}स्येषु}\ A}}\skt{ व्रणेषु }\edtext{\skt{मतिमान्भिषक्}}{\lemma{\skt{मतिमान्भिषक्}}\Afootnote{H; \skt{मतिमान्भिषक्}\ A}}\skt{~। }\pend
\pstart\versepar \skt{यथो\-क्त\-मा\-ह\-रे\-च्छल्यं }\edtext{\skt{प्राप्तो\-द्ध\-र\-ण\-ल\-क्षणैः}}{\lemma{\skt{प्राप्तो\appbrk{}द्ध\appbrk{}र\appbrk{}ण\appbrk{}ल\appbrk{}क्षणैः}}\Afootnote{H; \skt{प्राप्तो\appbrk{}द्ध\appbrk{}र\appbrk{}ण\appbrk{}ल\appbrk{}क्ष\appbrk{}णम्}\ A}}\skt{~॥}\pend

\pstart\versepar \ledsidenote{\passref{4.1.44}}\skt{रोगे }\edtext{\skt{व्यधनसाध्येषु}}{\lemma{\skt{व्यधनसाध्येषु}}\Afootnote{H; \skt{व्यधनसाध्ये तु}\ A}}\skt{ }\edtext{\skt{यथादेशं}}{\lemma{\skt{यथादेशं}}\Afootnote{H; \skt{यथोद्देशं}\ A}}\skt{ प्रमाणतः~। }\pend
\pstart\versepar \skt{शस्त्रं निदध्याद्दोषञ्च }\edtext{\skt{स्राव\-ये\-त्की\-र्त्ति\-त\-न्तथा}}{\lemma{\skt{कीर्त्तितन्तथा}}\Afootnote{H; \skt{कीर्तितं यथा}\ A}}\skt{~॥}\pend

\pstart\versepar \ledsidenote{\passref{4.1.45}}\skt{अपाकोपद्रुता ये च }\edtext{\skt{मांस\-स्था\-वि\-वृ\-ताश्च ये}}{\lemma{\skt{मांस\appbrk{}स्था\appbrk{}वि\appbrk{}वृ\appbrk{}ताश्च ये}}\Afootnote{H; \skt{मांसस्था विवृतांश्च ये}\ A}}\skt{~। }\pend
\pstart\versepar \skt{}\edtext{\skt{तेषु}}{\lemma{\skt{तेषु}}\Afootnote{H; \skt{यथोक्तं सीवनं तेषु}\ A}}\skt{ }\edtext{\skt{कार्यं}}{\lemma{\skt{कार्यं}}\Afootnote{H; \skt{कार्}\ A}}\skt{ }\edtext{\skt{यथोद्दिष्टं}}{\lemma{\skt{यथोद्दिष्टं}}\Afootnote{H; \emph{om.}\ A}}\skt{ }\edtext{\skt{सेव्यं}}{\lemma{\skt{सेव्यं}}\Afootnote{H; \skt{यं}\ A}}\skt{ सन्धानमेव च~॥}\pend

\pstart\versepar \ledsidenote{\passref{4.1.46}}\edtext{\skt{पूय\-ग\-र्ब्भा\-न\-नु\-द्वा\-रा\-न्व्र\-णा\-न्म\-र्म्म\-ग\-ता\-नपि}}{\lemma{\skt{पूय\appbrk{}ग\appbrk{}र्ब्भा\appbrk{}न\appbrk{}नु\appbrk{}द्वा\appbrk{}रान्}}\Afootnote{H; \skt{पूय\appbrk{}ग\appbrk{}र्भा\appbrk{}न\appbrk{}णु\appbrk{}द्वा\appbrk{}रान्}\ A}}\skt{~। }\pend
\pstart\versepar \skt{यथोक्तैः पीडनद्रव्यैः }\edtext{\skt{सम\-न्ता\-त्प्र\-ति\-पी\-ड\-येत्}}{\lemma{\skt{सम\appbrk{}न्ता\appbrk{}त्प्र\appbrk{}ति\appbrk{}पी\appbrk{}ड\appbrk{}येत्}}\Afootnote{H; \skt{सम\appbrk{}न्ता\appbrk{}त्प\appbrk{}रि\appbrk{}पी\appbrk{}ड\appbrk{}येत्}\ A}}\skt{~॥}\pend

\pstart\versepar \ledsidenote{\passref{4.1.47}}\skt{शुष्य\-मा\-ण\-मु\-पे\-क्षेत प्रदे\-ह\-म्पी\-ड\-न\-म्प्रति~। }\pend
\pstart\versepar \skt{न }\edtext{\skt{चापि}}{\lemma{\skt{चापि}}\Afootnote{H; \skt{चाभि}\ A}}\skt{ मुख\-मा\-लि\-म्पे\-त्त\-था\-दोषः प्रसिच्यते~॥}\pend

\pstart\versepar \ledsidenote{\passref{4.1.48}}\skt{तैस्तै\-र्न्नि\-मि\-त्तै\-र्ब्ब\-हुधा शोणिते प्रस्रुते भृशं~। }\pend
\pstart\versepar \skt{कार्यं यथोक्तं वैद्येन }\edtext{\skt{शोणि\-त\-स्था\-प\-न\-म्भ\-वेत्}}{\lemma{\skt{शोणि\appbrk{}त\appbrk{}स्था\appbrk{}प\appbrk{}न\appbrk{}म्भ\appbrk{}वेत्}}\Afootnote{H; \skt{शोणितास्थापनं भवेत्}\ A}}\skt{~॥}\pend

\pstart\versepar \ledsidenote{\passref{4.1.49}}\skt{दाह\-पा\-क\-ज्व\-र\-वतां व्रणा\-ना\-म्पि\-त्त\-को\-पतः~। }\pend
\pstart\versepar \skt{रक्तेन चाभिभूतानां कार्य\-न्नि\-र्व्वा\-प\-ण\-म्भ\-वेत्~॥}\pend

\pstart\versepar \ledsidenote{\passref{4.1.50}}\skt{यथोक्तैः शीतलद्रव्यैः क्षीर\-पि\-ष्टै\-र्घृ\-ता\-प्लुतैः~। }\pend
\pstart\versepar \skt{}\edtext{\skt{दिह्येत}}{\lemma{\skt{दिह्येत}}\Afootnote{H; \skt{दिह्याद}\ A}}\skt{ }\edtext{\skt{वहलां}}{\lemma{\skt{वहलां}}\Afootnote{H; \skt{बहलान्}\ A}}\skt{ }\edtext{\skt{सेकां}}{\lemma{\skt{सेकां}}\Afootnote{H; \skt{सेकान्}\ A}}\skt{ }\edtext{\skt{सुशी\-ता\-ञ्चा\-व\-चा\-र\-येत्}}{\lemma{\skt{सुशीताञ्}}\Afootnote{H; \skt{सुशीतांश्}\ A}}\skt{~॥}\pend

\pstart\versepar \ledsidenote{\passref{4.1.52}}\edtext{\skt{वाताभिभूते}}{\lemma{\skt{वाताभिभूते}}\Afootnote{H; \skt{वातघ्नवर्गेऽम्}\ A}}\skt{ }\edtext{\skt{वलयाः}}{\lemma{\skt{वलयाः}}\Afootnote{H; \skt{लगणे}\ A}}\skt{ }\edtext{\skt{काकोल्यादीगणे}}{\lemma{\skt{काकोल्यादीगणे}}\Afootnote{H; \skt{काकोल्यादिगणे}\ A}}\skt{ तथा~। }\pend
\pstart\versepar \skt{स्नैहिकेषु च वीजेषु }\edtext{\skt{कुर्यादुत्कारिकां}}{\lemma{\skt{कुर्यादुत्कारिकां}}\Afootnote{H; \skt{पचेदुत्कारिकां}\ A}}\skt{ शुभां~॥}\pend

\pstart\versepar \ledsidenote{\passref{4.1.53}}\edtext{\skt{येषां}}{\lemma{\skt{येषां}}\Afootnote{H; \skt{तेषां च}\ A}}\skt{ स्वेदनङ्कार्यं स्थिरा\-णा\-म्वे\-द\-ना\-वतां~। }\pend
\pstart\versepar \skt{दुर्ग्गन्धानां क्लेदवतां }\edtext{\skt{पिच्छि\-ला\-ना\-न्त\-थैव}}{\lemma{\skt{पिच्छि\appbrk{}ला\appbrk{}ना\appbrk{}न्त\appbrk{}थैव}}\Afootnote{H; \skt{पिच्छिलानां विशेष}\ A}}\skt{ }\edtext{\skt{च}}{\lemma{\skt{च}}\Afootnote{H; \skt{तः}\ A}}\skt{~॥}\pend

\pstart\versepar \ledsidenote{\passref{4.1.54}}\edtext{\skt{कषायं}}{\lemma{\skt{कषायं}}\Afootnote{H; \skt{कषायैः}\ A}}\skt{ शोधनङ्कार्यं शोधनैः प्रागुदीरितैः~। }\pend
\pstart\versepar \skt{}\edtext{\edtext{\skt{अन्तः\-पू\-या\-न\-णु\-मु\-खा\-न्ग\-म्भीरा}}{\lemma{\skt{अन्तः\appbrk{}पू\appbrk{}या\appbrk{}न\appbrk{}णु\appbrk{}मु\appbrk{}खान्}}\Afootnote{H; \skt{अन्तः\appbrk{}श\appbrk{}ल्या\appbrk{}न\appbrk{}णु\appbrk{}मु\appbrk{}खान्}\ A}}}{\lemma{\skt{गम्भीरा}}\Afootnote{H; \skt{गम्भीरान्}\ A}}\skt{ मांससंश्रितान्~॥}\pend

\pstart\versepar \ledsidenote{\passref{4.1.55}}\skt{शोध\-न\-द्र\-व्य\-यु\-क्ता\-भि\-र्व्व\-र्त्तिभिः }\edtext{\skt{समुपायरेत्}}{\lemma{\skt{समुपायरेत्}}\Afootnote{H; \skt{तान्यथाक्रमम्}\ A}}\skt{~। }\pend
\pstart\versepar \skt{}\edtext{\skt{पुति\-मां\-स\-प्र\-ति\-च्छन्नां}}{\lemma{\skt{पुति\appbrk{}मां\appbrk{}स\appbrk{}प्र\appbrk{}ति\appbrk{}च्छन्नां}}\Afootnote{H; \skt{पूति\appbrk{}मां\appbrk{}स\appbrk{}प्र\appbrk{}ति\appbrk{}च्छ\appbrk{}न्नान्}\ A}}\skt{ }\edtext{\skt{महा\-दो\-षा\-नु\-दी\-रितैः}}{\lemma{\skt{महा\appbrk{}दो\appbrk{}षा\appbrk{}नु\appbrk{}दी\appbrk{}रितैः}}\Afootnote{H; \skt{महादोषांश्च शोधयेत्}\ A}}\skt{~॥}\pend

\pstart\versepar \ledsidenote{\passref{4.1.56}}\edtext{\skt{वर्त्तिद्रव्यैः}}{\lemma{\skt{वर्त्तिद्रव्यैः}}\Afootnote{H; \skt{कल्की\appbrk{}कृ\appbrk{}तै\appbrk{}र्य\appbrk{}था\appbrk{}लाभं वर्तिद्रव्यैः}\ A}}\skt{ }\edtext{\skt{कल्क\-कृ\-तै\-र्य\-था\-ला\-भेन}}{\lemma{\skt{कल्क\appbrk{}कृ\appbrk{}तै\appbrk{}र्य\appbrk{}था\appbrk{}ला\appbrk{}भेन}}\Afootnote{H; \emph{om.}\ A}}\skt{ }\edtext{\skt{शोधयेत्}}{\lemma{\skt{शोधयेत्}}\Afootnote{H; \skt{पुरोदितैः}\ A}}\skt{~। }\pend
\pstart\versepar \skt{}\edtext{\skt{पित्तप्रदुष्टां}}{\lemma{\skt{पित्तप्रदुष्टां}}\Afootnote{H; \skt{पित्तप्रदुष्टान्}\ A}}\skt{ }\edtext{\skt{गंभी\-रा\-न्दा\-ह\-पा\-क\-प्र\-पी\-डितां}}{\lemma{\skt{गंभी\appbrk{}रा\appbrk{}न्दा\appbrk{}ह\appbrk{}पा\appbrk{}क\appbrk{}प्र\appbrk{}पी\appbrk{}डितां}}\Afootnote{H; \skt{गम्भी\appbrk{}रा\appbrk{}न्दा\appbrk{}ह\appbrk{}पा\appbrk{}क\appbrk{}प्र\appbrk{}पी\appbrk{}डि\appbrk{}तान्}\ A}}\skt{~॥}\pend

\pstart\versepar \ledsidenote{\passref{4.1.57}}\skt{कार्प्पा\-सी\-फ\-ल\-मि\-श्रेण }\edtext{\skt{युञ्ज्या\-च्छो\-ध\-न\-स\-र्प्पिषा}}{\lemma{\skt{युञ्ज्या\appbrk{}च्छो\appbrk{}ध\appbrk{}न\appbrk{}स\appbrk{}र्प्पिषा}}\Afootnote{H; \skt{जये\appbrk{}च्छो\appbrk{}ध\appbrk{}न\appbrk{}स\appbrk{}र्पिषा}\ A}}\skt{~। }\pend
\pstart\versepar \skt{}\edtext{\skt{उत्स\-न्न\-मा\-न्सा\-न\-स्नि\-ग्धा\-न\-ल्प\-स्रा\-वा\-न्नृ\-णा\-म्व्र\-णान्}}{\lemma{\skt{उत्स\appbrk{}न्न\appbrk{}मा\appbrk{}न्सा\appbrk{}न\appbrk{}स्नि\appbrk{}ग्धा\appbrk{}न\appbrk{}ल्प\appbrk{}स्रा\appbrk{}वा\appbrk{}न्नृ\appbrk{}णा\appbrk{}म्व्र\appbrk{}णान्}}\Afootnote{H; \skt{उत्स\appbrk{}न्न\appbrk{}मां\appbrk{}सा\appbrk{}न\appbrk{}स्नि\appbrk{}ग्धा\appbrk{}न\appbrk{}ल्प\appbrk{}स्रा\appbrk{}वा\appbrk{}न्व्र\appbrk{}णां\appbrk{}स्तथा}\ A}}\skt{~॥}\pend

\pstart\versepar \ledsidenote{\passref{4.1.58}}\skt{सर्ष\-प\-स्ने\-ह\-यु\-क्तेन }\edtext{\skt{धीमान्तैले}}{\lemma{\skt{धीमान्तैले}}\Afootnote{H; \skt{धीमांस्तैलेन}\ A}}\skt{ शोधयेत्~। }\pend
\pstart\versepar \skt{}\edtext{\skt{व्रणानां}}{\lemma{\skt{व्रणानां}}\Afootnote{H; \emph{om.}\ A}}\skt{ }\edtext{\skt{स्थिर\-मां\-सा\-ना\-न्दुः\-शो\-ध्या\-ना\-न्त\-थैव च}}{\lemma{\skt{स्थिर\appbrk{}मां\appbrk{}सा\appbrk{}ना\appbrk{}न्दुः\appbrk{}शो\appbrk{}ध्या\appbrk{}ना\appbrk{}न्त\appbrk{}थैव च}}\Afootnote{H; \skt{तैले\appbrk{}ना\appbrk{}शु\appbrk{}ध्य\appbrk{}मा\appbrk{}नानां शोधनीयां रसक्रियाम्}\ A}}\skt{~।}\pend

\pstart\versepar \ledsidenote{\passref{4.1.59}}\edtext{\skt{रसक्रियां}}{\lemma{\skt{रसक्रियां}}\Afootnote{H; \skt{व्रणानां स्थिर}\ A}}\skt{ }\edtext{\skt{शोध\-नी\-या\-ङ्कु\-र्या\-द्द्र\-व्यै\-रु\-दी\-रितैः}}{\lemma{\skt{शोधनीयाङ्}}\Afootnote{H; \skt{मांसानां}\ A}}\skt{~॥ }\pend
\pstart\versepar \skt{कषाये विधिवत्तेषां कृते }\edtext{\skt{चापि}}{\lemma{\skt{चापि}}\Afootnote{H; \skt{चाधि}\ A}}\skt{ }\edtext{\skt{स्रिते}}{\lemma{\skt{स्रिते}}\Afootnote{H; \skt{श्रयेत्}\ A}}\skt{ }\edtext{\skt{वुधः}}{\lemma{\skt{वुधः}}\Afootnote{H; \skt{पुनः}\ A}}\skt{~।}\pend

\pstart\versepar \ledsidenote{\passref{4.1.60}}\skt{सुराष्ट्रजां }\edtext{\skt{सकाशीसां}}{\lemma{\skt{सकाशीसां}}\Afootnote{H; \skt{सकासीसां}\ A}}\skt{ दद्याच्चापि मनःशिलां~॥ }\pend
\pstart\versepar \skt{}\edtext{\skt{हरि\-ता\-ल\-ञ्च\-म\-ति\-मा\-न्त\-त\-स्ता\-न\-व\-चा\-र\-येत्}}{\lemma{\skt{तत\appbrk{}स्ता\appbrk{}न\appbrk{}व\appbrk{}चा\appbrk{}र\appbrk{}येत्}}\Afootnote{H; \skt{स्तत\appbrk{}स्ता\appbrk{}म\appbrk{}व\appbrk{}चा\appbrk{}र\appbrk{}येत्}\ A}}\skt{~।}\pend

\pstart\versepar \ledsidenote{\passref{4.1.61}}\edtext{\skt{मातु\-लु\-ङ्ग\-र\-सो\-पेताः}}{\lemma{\skt{मातु\appbrk{}लु\appbrk{}ङ्ग\appbrk{}र\appbrk{}सो\appbrk{}पेताः}}\Afootnote{H; \skt{मातु\appbrk{}लु\appbrk{}ङ्ग\appbrk{}र\appbrk{}सो\appbrk{}पेतां}\ A}}\skt{ }\edtext{\skt{सक्षौद्राः}}{\lemma{\skt{सक्षौद्राः}}\Afootnote{H; \skt{सक्षौद्रा}\ A}}\skt{ }\edtext{\skt{सुप्रमर्द्दिताः}}{\lemma{\skt{सुप्रमर्द्दिताः}}\Afootnote{H; \skt{मतिमर्दिताम्}\ A}}\skt{~॥ }\pend
\pstart\versepar \skt{व्रणेषु }\edtext{\skt{दत्ता}}{\lemma{\skt{दत्ता}}\Afootnote{H; \skt{दत्त्वा तां}\ A}}\skt{ }\edtext{\skt{तिष्ठेत्तु}}{\lemma{\skt{तिष्ठेत्तु}}\Afootnote{H; \skt{तिष्ठेत्त्}\ A}}\skt{ }\edtext{\skt{त्रीम्त्रीश्च}}{\lemma{\skt{त्रीम्त्रीश्च}}\Afootnote{H; \skt{रींस्त्रींश्च}\ A}}\skt{ }\edtext{\skt{दिवसान्परान्}}{\lemma{\skt{परान्}}\Afootnote{H; \skt{परम्}\ A}}\skt{~।}\pend

\pstart\versepar \ledsidenote{\passref{4.1.62}}\edtext{\skt{मेदो\-प\-सृ\-ष्टा\-न्दु\-र्ग्गन्धां}}{\lemma{\skt{मेदो\appbrk{}प\appbrk{}सृ\appbrk{}ष्टा\appbrk{}न्दु\appbrk{}र्ग्गन्धां}}\Afootnote{H; \skt{मेदो\appbrk{}जु\appbrk{}ष्टा\appbrk{}न\appbrk{}ग\appbrk{}म्भी\appbrk{}रा\appbrk{}न्दु\appbrk{}र्ग\appbrk{}न्धांश्}\ A}}\skt{ }\edtext{\skt{गम्भी\-रा\-ञ्चू\-र्ण्ण\-शो\-धनैः}}{\lemma{\skt{गम्भी\appbrk{}रा\appbrk{}ञ्चू\appbrk{}र्ण्ण\appbrk{}शो\appbrk{}धनैः}}\Afootnote{H; \skt{चूर्णशोधनैः}\ A}}\skt{~। }\pend
\pstart\versepar \skt{}\edtext{\skt{उपा\-च\-रे\-द्भि\-ष\-क्प्राज्ञः}}{\lemma{\skt{उपा\appbrk{}च\appbrk{}रे\appbrk{}द्भि\appbrk{}ष\appbrk{}क्प्राज्ञः}}\Afootnote{H; \skt{उपा\appbrk{}च\appbrk{}रे\appbrk{}त्भि\appbrk{}ष\appbrk{}क्प्राज्ञः}\ A}}\skt{ }\edtext{\skt{शूक्ष्मैः}}{\lemma{\skt{शूक्ष्मैः}}\Afootnote{H; \skt{श्लक्ष्णैः}\ A}}\skt{ }\edtext{\skt{शोधनबस्तिजैः}}{\lemma{\skt{शोधनबस्तिजैः}}\Afootnote{H; \skt{शोधनवर्तिजैः}\ A}}\skt{~।}\pend

\pstart\versepar \ledsidenote{\passref{4.1.63}}\edtext{\skt{शुद्ध\-ल\-क्ष\-ण\-सं\-युक्ते}}{\lemma{\skt{शुद्ध\appbrk{}ल\appbrk{}क्ष\appbrk{}ण\appbrk{}सं\appbrk{}युक्ते}}\Afootnote{H; \skt{शुद्ध\appbrk{}ल\appbrk{}क्ष\appbrk{}ण\appbrk{}यु\appbrk{}क्तानां}\ A}}\skt{ कषायं रोपणं हितं~॥ }\pend
\pstart\versepar \skt{तत्र कार्यं }\edtext{\skt{यथो\-द्दि\-ष्टै\-र्द्द्र\-व्यै\-र्व्वै\-द्येन}}{\lemma{\skt{यथो\appbrk{}द्दि\appbrk{}ष्टै\appbrk{}र्द्द्र\appbrk{}व्यैर्}}\Afootnote{H; \skt{यथो\appbrk{}द्दि\appbrk{}ष्टै\appbrk{}र्द्र\appbrk{}व्यैर्}\ A}}\skt{ जानता~।}\pend

\pstart\versepar \ledsidenote{\passref{4.1.64}}\skt{अवेदनानां शुद्धानां गम्भी\-रा\-णा\-न्त\-थैव च~॥ }\pend
\pstart\versepar \skt{}\edtext{\skt{हिताः}}{\lemma{\skt{हिताः}}\Afootnote{H; \skt{हिता}\ A}}\skt{ }\edtext{\skt{संरोपणावर्त्त्य}}{\lemma{\skt{संरोपणावर्त्त्य}}\Afootnote{H; \skt{रोपणवर्}\ A}}\skt{ }\edtext{\skt{वत्यङ्गैरोपणैः}}{\lemma{\skt{वत्यङ्गैरोपणैः}}\Afootnote{H; \skt{त्यङ्ग}\ A}}\skt{ }\edtext{\skt{कृताः}}{\lemma{\skt{कृताः}}\Afootnote{H; \skt{कृता रोपणवर्तयः}\ A}}\skt{~।}\pend

\pstart\versepar \ledsidenote{\passref{4.1.65}}\skt{अपे\-त\-पू\-ति\-मां\-सानां मांस\-स्था\-ना\-म\-रो\-हतां~॥ }\pend
\pstart\versepar \skt{}\edtext{\skt{कल्कस्तुरोपणः}}{\lemma{\skt{कल्कस्तुरोपणः}}\Afootnote{H; \skt{कल्कः संरोहणः}\ A}}\skt{ कार्यस्तिलजो }\edtext{\skt{मधुसंयुतः}}{\lemma{\skt{मधुसंयुतः}}\Afootnote{H; \skt{मधुस्ंयुतः}\ A}}\skt{~।}\pend

\pstart\versepar \ledsidenote{\passref{4.1.66}}\edtext{\skt{स माधु\-र्या\-त्त\-थै\-वो\-ष्ण्या\-त्स्ने\-हा\-च्चा\-नि\-ल\-ना\-शनः}}{\lemma{\skt{स माधु\appbrk{}र्या\appbrk{}त्त\appbrk{}थै\appbrk{}वो\appbrk{}ष्ण्या\appbrk{}त्स्ने\appbrk{}हा\appbrk{}च्चा\appbrk{}नि\appbrk{}ल\appbrk{}ना\appbrk{}शनः}}\Afootnote{H; \skt{स माधु\appbrk{}र्या\appbrk{}त्त\appbrk{}थौ\appbrk{}ष्ण्याच्च स्नेहा\appbrk{}च्चा\appbrk{}नि\appbrk{}ल\appbrk{}ना\appbrk{}शनः}\ A}}\skt{~॥ }\pend
\pstart\versepar \skt{}\edtext{\skt{कषा\-य\-भा\-वा\-त्मा\-धु\-र्या\-त्ति\-क्त\-त्वा\-च्चापि}}{\lemma{\skt{कषा\appbrk{}य\appbrk{}भा\appbrk{}वा\appbrk{}त्मा\appbrk{}धु\appbrk{}र्या\appbrk{}त्ति\appbrk{}क्त\appbrk{}त्वाच्}}\Afootnote{H; \skt{कषा\appbrk{}य\appbrk{}भा\appbrk{}वा\appbrk{}न्मा\appbrk{}धु\appbrk{}र्या\appbrk{}त्ति\appbrk{}क्त\appbrk{}त्वाच्}\ A}}\skt{ }\edtext{\skt{पित्तहा}}{\lemma{\skt{पित्तहा}}\Afootnote{H; \skt{पित्तहृत्}\ A}}\skt{~।}\pend

\pstart\versepar \ledsidenote{\passref{4.1.67}}\skt{औष्ण्या\-त्क\-षा\-य\-भा\-वाच्च तिक्तत्वाच्च कफे हितः~॥ }\pend
\pstart\versepar \skt{शोध\-ये\-द्रो\-प\-ये\-च्चापि युक्तः शोधनरोपणैः~।}\pend

\pstart\versepar \ledsidenote{\passref{4.1.68}}\skt{निम्व\-प\-त्र\-म\-धु\-भ्यान्तु युक्तः संशोधनः स्मृतः~॥ }\pend
\pstart\versepar \skt{पूर्व्वाभ्यां सर्प्पिषा चापि }\edtext{\skt{युक्त\-श्चा\-प्यु\-प\-रो\-पणे}}{\lemma{\skt{चाप्युपरोपणे}}\Afootnote{H; \skt{चाप्युपरोपणः}\ A}}\skt{~।}\pend

\pstart\versepar \ledsidenote{\passref{4.1.69}}\skt{तिल\-व\-द्य\-व\-क\-ल्कन्तु }\edtext{\edtext{\skt{केचि\-दा\-हु\-र्म्म\-नी\-षिणः}}{\lemma{\skt{केचिदाहुर्}}\Afootnote{A; \skt{केचिदाहु}\ H}}}{\lemma{\skt{म्मनीषिणः}}\Afootnote{A; \skt{र्म्मनीषिणः}\ H}}\skt{~॥ }\pend
\pstart\versepar \skt{}\edtext{\skt{अविदग्धं}}{\lemma{\skt{अविदग्धं}}\Afootnote{H; \emph{om.}\ A}}\skt{ }\edtext{\edtext{\edtext{\skt{प्रश\-म\-ये\-द्वि\-द\-ग्ध\-म्वापि}}{\lemma{\skt{प्रशमयेद्}}\Afootnote{H; \skt{शमयेद}\ A}}}{\lemma{\skt{विदग्धम्}}\Afootnote{H; \skt{विदग्धं च}\ A}}}{\lemma{\skt{वापि}}\Afootnote{H; \skt{विदग्धमपि}\ A}}\skt{ पाचयेत्~।}\pend

\pstart\versepar \ledsidenote{\passref{4.1.70}}\edtext{\skt{पक्वञ्च}}{\lemma{\skt{पक्वञ्च}}\Afootnote{H; \skt{पक्वं}\ A}}\skt{ }\edtext{\skt{भिन्द्याद्भिन्नश्च}}{\lemma{\skt{भिन्द्याद्भिन्नश्च}}\Afootnote{H; \skt{भिनत्ति भिन्नं च}\ A}}\skt{ }\edtext{\skt{शोध\-ये\-द्रो\-प\-यीत च}}{\lemma{\skt{रोपयीत च}}\Afootnote{H; \skt{रोपयेत्तथा}\ A}}\skt{~॥ }\pend
\pstart\versepar \skt{पित्त\-र\-क्त\-वि\-षा\-ग\-न्तू\-न्ग\-म्भी\-रा\-नपि च व्रणान्~।}\pend

\pstart\versepar \ledsidenote{\passref{4.1.71}}\skt{रोप\-ये\-द्रो\-प\-णी\-येन क्षीरसिद्धेन सर्प्पिषा~॥ }\pend
\pstart\versepar \skt{}\edtext{\skt{कफ\-वा\-ता\-र्द्दि\-ता\-नान्तु}}{\lemma{\skt{कफ\appbrk{}वा\appbrk{}ता\appbrk{}र्द्दि\appbrk{}ता\appbrk{}नान्तु}}\Afootnote{H; \skt{कफ\appbrk{}वा\appbrk{}ता\appbrk{}भि\appbrk{}भू\appbrk{}तानां}\ A}}\skt{ }\edtext{\skt{व्रणा\-ना\-न्तै\-ल\-रो\-पणं}}{\lemma{\skt{व्रणा\appbrk{}ना\appbrk{}न्तै\appbrk{}ल\appbrk{}रो\appbrk{}पणं}}\Afootnote{H; \skt{व्रणानां मतिमान्भिषक्}\ A}}\skt{~।}\pend

\pstart\versepar \ledsidenote{\passref{4.1.72}}\edtext{\edtext{\skt{कुर्यात्तत्र}}{\lemma{\skt{कुर्यात्}}\Afootnote{H; \skt{कारयेद्}\ A}}}{\lemma{\skt{तत्र}}\Afootnote{H; \skt{रोपणं तैलं भेषजैस्तद्}\ A}}\skt{ }\edtext{\skt{यथोक्तेषु}}{\lemma{\skt{यथोक्तेषु}}\Afootnote{H; \skt{यथो}\ A}}\skt{ }\edtext{\skt{द्रव्येषु}}{\lemma{\skt{द्रव्येषु}}\Afootnote{H; \emph{om.}\ A}}\skt{ }\edtext{\skt{मतिमान्भिषक्}}{\lemma{\skt{मतिमान्}}\Afootnote{H; \emph{om.}\ A}}\skt{~॥ }\pend
\pstart\versepar \skt{अव्य\-ध्या\-ना\-ञ्च\-ल\-स्थानां शुद्धानाञ्च प्रदुष्यतां~।}\edtext{}{\lemma{\skt{भिषक्~॥ अव्य\appbrk{}ध्या\appbrk{}ना\appbrk{}ञ्च\appbrk{}ल\appbrk{}स्थानां}}\Afootnote{H; \skt{दितैः~। अबन्ध्यानां चलस्थानां}\ A}}\pend

\pstart\versepar \ledsidenote{\passref{4.1.73}}\skt{द्विहरिद्रायुतां }\edtext{\skt{कुर्याद्रोपणीयै}}{\lemma{\skt{रोपणीयै}}\Afootnote{H; \skt{रोपणार्थां}\ A}}\skt{ }\edtext{\skt{रसक्रिया}}{\lemma{\skt{रसक्रिया}}\Afootnote{H; \skt{रसक्रियाम्}\ A}}\skt{~॥ }\pend
\pstart\versepar \skt{समानां स्थिरमांसानां }\edtext{\skt{त्वक्स्था\-ना\-ञ्चू\-र्ण्ण\-रो\-पणं}}{\lemma{\skt{त्वक्स्था\appbrk{}ना\appbrk{}ञ्चू\appbrk{}र्ण्ण\appbrk{}रो\appbrk{}पणं}}\Afootnote{H; \skt{त्वक्स्थानां रोपणं भिषक्}\ A}}\skt{~।}\pend

\pstart\versepar \ledsidenote{\passref{4.1.74}}\edtext{\edtext{\skt{प्राप्त\-म्वि\-द\-ध्या\-त्म\-ति\-मा\-न्प्राक्स}}{\lemma{\skt{प्राप्तम्विदध्यात्}}\Afootnote{H; \skt{चूर्णं विदध्यान्}\ A}}}{\lemma{\skt{प्राक्स}}\Afootnote{H; \skt{प्राक्स्}\ A}}\skt{ }\edtext{\skt{उक्तो}}{\lemma{\skt{उक्तो}}\Afootnote{H; \skt{थानोक्तो}\ A}}\skt{ विधिर्यथा~॥ }\pend
\pstart\versepar \skt{शोधनो रोपणश्चैव }\edtext{\edtext{\skt{विधिर्यो}}{\lemma{\skt{विधिर्}}\Afootnote{H; \skt{विधि}\ A}}\skt{ यम्प्रकीर्त्तितः}}{\lemma{\skt{यो यम्}}\Afootnote{H; \skt{र्योऽयं}\ A}}\skt{~।}\pend

\pstart\versepar \ledsidenote{\passref{4.1.75}}\skt{सर्व्वव्रणानां सामान्ये नोक्तो }\edtext{\skt{दोषविशेषतः}}{\lemma{\skt{दोषविशेषतः}}\Afootnote{H; \skt{दोषाविशेषतः}\ A}}\skt{~॥ }\pend
\pstart\versepar \skt{एष आग\-म\-सि\-द्ध\-त्वा\-त्त\-थैव फलदर्शनात्~।}\pend

\pstart\versepar \ledsidenote{\passref{4.1.76}}\skt{मन्त्र\-व\-त्सं\-प्र\-यो\-क्तव्यो न मीमांस्यः कथंचन~॥ }\pend
\pstart\versepar \skt{स्ववुद्ध्या }\edtext{\skt{विभजेद्युक्त्या}}{\lemma{\skt{विभजेद्युक्त्या}}\Afootnote{H; \skt{चापि विभजेत्}\ A}}\skt{ कषायादिषु }\edtext{\skt{योगवित्}}{\lemma{\skt{योगवित्}}\Afootnote{H; \skt{सप्तसु}\ A}}\skt{~।}\pend

\pstart\versepar \ledsidenote{\passref{4.1.77}}\edtext{\skt{दोषो\-च्छ्रा\-ये\-णौ\-ष\-धानि}}{\lemma{\skt{दोषो\appbrk{}च्छ्रा\appbrk{}ये\appbrk{}णौ\appbrk{}ष\appbrk{}धानि}}\Afootnote{H; \skt{भेषजानि}\ A}}\skt{ }\edtext{\skt{यान्युक्तानि}}{\lemma{\skt{यान्युक्तानि}}\Afootnote{H; \skt{यथायोगं यान्युक्तानि}\ A}}\skt{ पुरा मया~॥ }\pend
\pstart\versepar \skt{आद्ये द्वे पञ्चमूल्यौ तु गणो यश्चानिलापहः~।}\pend

\pstart\versepar \ledsidenote{\passref{4.1.78}}\skt{स वातदुष्टे }\edtext{\skt{देयस्तु}}{\lemma{\skt{देयस्तु}}\Afootnote{H; \skt{दातव्यः}\ A}}\skt{ कषायादिषु सप्तसु~॥ }\pend
\pstart\versepar \skt{}\edtext{\skt{न्यग्रोधादिगणो यस्तु}}{\lemma{\skt{न्यग्रोधादिगणो यस्तु}}\Afootnote{H; \skt{न्यग्रोधादिर्गणो यस्तु}\ A}}\skt{ काकोल्यादिश्च यो }\edtext{\skt{गणः}}{\lemma{\skt{गणः}}\Afootnote{H; \skt{स्मृतः}\ A}}\skt{~।}\pend

\pstart\versepar \ledsidenote{\passref{4.1.79}}\skt{तौ पित्तदुष्टे दातव्यौ कषायादिषु सप्तसु~। }\pend
\pstart\versepar \skt{}\edtext{\skt{आरग्वधादिश्च}}{\lemma{\skt{आरग्वधादिश्च}}\Afootnote{H; \skt{आरग्वधादिस्तु}\ A}}\skt{ गणो }\edtext{\skt{यश्चोष्णः}}{\lemma{\skt{यश्चोष्णः}}\Afootnote{H; \skt{यश्चोष्णः}\ A}}\skt{ परिकीर्त्तितः~।}\pend

\pstart\versepar \ledsidenote{\passref{4.1.80}}\skt{तौ देयौ कफदुष्टे तु संसृष्टे }\edtext{\skt{संयुतागणाः}}{\lemma{\skt{संयुतागणाः}}\Afootnote{H; \skt{संयता गणाः}\ A}}\skt{~॥ }\pend
\pstart\versepar \skt{}\edtext{\skt{वातो\-प\-सृ\-ष्टा\-नु\-ग्र\-रु\-जा\-न्सा\-स्रा\-वा\-नपि}}{\lemma{\skt{वातो\appbrk{}प\appbrk{}सृ\appbrk{}ष्टा\appbrk{}नु\appbrk{}ग्र\appbrk{}रु\appbrk{}जान्}}\Afootnote{H; \skt{वाता\appbrk{}त्म\appbrk{}का\appbrk{}नु\appbrk{}ग\appbrk{}र\appbrk{}रु\appbrk{}जान्}\ A}}\skt{ च व्रणान्~।}\pend

\pstart\versepar \ledsidenote{\passref{4.1.81}}\edtext{\skt{सक्षौ\-म\-य\-व\-स\-र्प्पि\-र्ब्भि\-र्द्धू\-प\-ना\-ङ्गैस्तु}}{\lemma{\skt{सक्षौ\appbrk{}म\appbrk{}य\appbrk{}व\appbrk{}स\appbrk{}र्प्पि\appbrk{}र्ब्भि\appbrk{}र्द्धू\appbrk{}प\appbrk{}ना\appbrk{}ङ्गैस्तु}}\Afootnote{H; \skt{सक्षौ\appbrk{}म\appbrk{}य\appbrk{}व\appbrk{}स\appbrk{}र्पि\appbrk{}र्भि\appbrk{}र्धू\appbrk{}प\appbrk{}ना\appbrk{}ङ्गैश्च}\ A}}\skt{ }\edtext{\skt{पूरयेत्}}{\lemma{\skt{पूरयेत्}}\Afootnote{H; \skt{धूपयेत्}\ A}}\skt{~। }\pend
\pstart\versepar \skt{परि\-शु\-ष्का\-ल्प\-मां\-सानां गम्भी\-रा\-णा\-न्त\-थैव च~।}\pend

\pstart\versepar \ledsidenote{\passref{4.1.82}}\skt{कुर्या\-दु\-त्सा\-द\-नी\-यानि सर्प्पीं\-ष्या\-ले\-प\-नानि च~॥ }\pend
\pstart\versepar \skt{}\edtext{\skt{मांसादानाञ्च}}{\lemma{\skt{मांसादानाञ्च}}\Afootnote{H; \skt{मांसाशिनां च}\ A}}\skt{ मांसानि भक्ष\-ये\-द्वि\-धि\-व\-न्नरः~।}\pend

\pstart\versepar \ledsidenote{\passref{4.1.83}}\edtext{\skt{विशु\-द्ध\-म\-न\-स\-स्तस्य}}{\lemma{\skt{विशु\appbrk{}द्ध\appbrk{}म\appbrk{}न\appbrk{}स\appbrk{}स्तस्य}}\Afootnote{H; \skt{विशु\appbrk{}द्ध\appbrk{}म\appbrk{}न\appbrk{}स\appbrk{}स्तस्य}\ A}}\skt{ मांसम्मांसेन वर्द्धते~॥ }\pend
\pstart\versepar \skt{उत्स\-न्न\-मृ\-दु\-मां\-सानां व्रणा\-ना\-म\-व\-सा\-दनं~।}\pend

\pstart\versepar \ledsidenote{\passref{4.1.84}}\edtext{\edtext{\skt{कुर्या\-द्द्र\-व्यै\-र्य\-क्तैस्तु}}{\lemma{\skt{द्रव्यैर्}}\Afootnote{A; \skt{द्रव्यै}\ H}}}{\lemma{\skt{यक्तैस्तु}}\Afootnote{\skt{यथोद्दिष्टैश्}\ A; \skt{र्यक्ताक्तैस्तु}\ H\textsuperscript{ac}; \skt{र्यक्तैस्तु}\ H\textsuperscript{pc}}}\skt{ }\edtext{\skt{चूर्ण्णितैर्म्मधुना}}{\lemma{\skt{चूर्ण्णितैर्म्मधुना}}\Afootnote{H; \skt{चूर्णितैर्मधुना}\ A}}\skt{ सह~॥ }\pend
\pstart\versepar \skt{}\edtext{\skt{कठि\-ना\-ना\-म\-मां\-सानां}}{\lemma{\skt{कठि\appbrk{}ना\appbrk{}ना\appbrk{}म\appbrk{}मां\appbrk{}सानां}}\Afootnote{H; \skt{कठि\appbrk{}ना\appbrk{}ना\appbrk{}म\appbrk{}मां\appbrk{}सानां}\ A}}\skt{ दुष्टानां }\edtext{\skt{मातरिश्वना}}{\lemma{\skt{मातरिश्वना}}\Afootnote{H; \skt{मतिरिश्वना}\ A}}\skt{~।}\pend

\pstart\versepar \ledsidenote{\passref{4.1.85}}\skt{मृद्वी क्रिया विधातव्या शोणितञ्चापि मोक्षयेत्~॥ }\pend
\pstart\versepar \skt{}\edtext{\skt{वात\-घ्नौ\-ष\-ध\-सं\-युक्तैः}}{\lemma{\skt{वात\appbrk{}घ्नौ\appbrk{}ष\appbrk{}ध\appbrk{}सं\appbrk{}युक्तैः}}\Afootnote{H; \skt{वात\appbrk{}घ्नौ\appbrk{}ष\appbrk{}ध\appbrk{}सं\appbrk{}यु\appbrk{}क्तान्}\ A}}\skt{ }\edtext{\skt{स्नेहस्वेदौ च}}{\lemma{\skt{स्नेहस्वेदौ च}}\Afootnote{H; \skt{स्नेहान्सेकांश्च}\ A}}\skt{ }\edtext{\skt{कारयेत्}}{\lemma{\skt{कारयेत्}}\Afootnote{H; \skt{पूरयेत्~। कासीसं रोचनां तुत्थं हरितालं मनःशिलाम्}\ A}}\skt{~।}\pend

\pstart\versepar \ledsidenote{\passref{4.1.86}}\skt{मृदु\-त्व\-मा\-शु\-रो\-हञ्च }\edtext{\skt{गाढ\-व\-न्ध\-ङ्क\-रोति हि}}{\lemma{\skt{गाढ\appbrk{}व\appbrk{}न्ध\appbrk{}ङ्क\appbrk{}रोति हि}}\Afootnote{H; \skt{गाढो बन्धः करोति हि}\ A}}\skt{~॥ }\pend
\pstart\versepar \skt{}\edtext{\skt{व्रणेषु}}{\lemma{\skt{व्रणेषु}}\Afootnote{H; \skt{वर्णेषु}\ A}}\skt{ मृदुमांसेषु दारुणीकरणं हितं~। }\pend
\pstart\versepar \skt{धव\-प्रि\-यं\-ग्व\-शो\-कानां रोहिण्याश्च }\edtext{\edtext{\skt{त्वचस्तथा}}{\lemma{\skt{त्वचस्}}\Afootnote{H; \skt{त्वच}\ A}}}{\lemma{\skt{तथा}}\Afootnote{H; \skt{स्तथा}\ A}}\skt{~॥}\pend

\pstart\versepar \ledsidenote{\passref{4.1.87}}\edtext{\skt{त्रिफ\-ला\-पु\-ष्प\-का\-सीसे}}{\lemma{\skt{त्रिफ\appbrk{}ला\appbrk{}पु\appbrk{}ष्प\appbrk{}का\appbrk{}सीसे}}\Afootnote{H; \skt{त्रिफ\appbrk{}ला\appbrk{}धा\appbrk{}त\appbrk{}की\appbrk{}पुष्प}\ A}}\skt{ }\edtext{\skt{तथासर्ज्जरसं}}{\lemma{\skt{तथासर्ज्जरसं}}\Afootnote{H; \skt{रोध्र\appbrk{}स\appbrk{}र्ज\appbrk{}र\appbrk{}सान्स}\ A}}\skt{ }\edtext{\skt{समं}}{\lemma{\skt{समं}}\Afootnote{H; \skt{मान्}\ A}}\skt{~। }\pend
\pstart\versepar \skt{कृत्वा सूक्ष्माणि चूर्ण्णानि व्रणं }\edtext{\skt{तैरे\-व\-चू\-र्ण्ण\-येत्}}{\lemma{\skt{तैरे\appbrk{}व\appbrk{}चू\appbrk{}र्ण्ण\appbrk{}येत्}}\Afootnote{H; \skt{तैर\appbrk{}व\appbrk{}चू\appbrk{}र्ण\appbrk{}येत्}\ A}}\skt{~॥}\pend

\pstart\versepar \ledsidenote{\passref{4.1.88}}\skt{उत्स\-न्न\-मा\-न्सा\-न्क\-ठि\-ना\-न्क\-ण्डू\-यु\-क्ता\-श्चि\-रो\-त्थि\-तान्~। }\pend
\pstart\versepar \skt{}\edtext{\skt{दुःशोधनीयाश्च}}{\lemma{\skt{दुःशोधनीयाश्च}}\Afootnote{H; \skt{तथैव खलु दुःशोध्}\ A}}\skt{ }\edtext{\skt{तथा\-शो\-ध\-ये\-त्क्षा\-र\-क\-र्म्मणा}}{\lemma{\skt{तथाशोधयेत्}}\Afootnote{H; \skt{याञ्शोधयेत्}\ A}}\skt{~॥}\pend

\pstart\versepar \ledsidenote{\passref{4.1.89}}\edtext{\skt{स्रव\-तोऽ\-श्म\-रि\-जा\-न्मूत्रं}}{\lemma{\skt{ऽश्मरिजान्मूत्रं}}\Afootnote{H; \skt{ऽश्मभवान्मूत्रं}\ A}}\skt{ येचान्ये रक्तवाहिनः~। }\pend
\pstart\versepar \skt{}\edtext{\skt{संधी}}{\lemma{\skt{संधी}}\Afootnote{H; \emph{om.}\ A}}\skt{ }\edtext{\skt{निःशे\-ष\-त\-श्च्छि\-न्ना\-न्सा\-ध\-ये\-द\-ग्नि\-क\-र्म्मणा}}{\lemma{\skt{निःशे\appbrk{}ष\appbrk{}त\appbrk{}श्च्छि\appbrk{}न्नान्}}\Afootnote{H; \skt{निःशे\appbrk{}ष\appbrk{}च्छि\appbrk{}न्न\appbrk{}स\appbrk{}न्धींश्च}\ A}}\skt{~॥}\pend

\pstart \ledsidenote{\passref{4.1.90}}\skt{दुरूढत्वात्तु शुक्लानां कृष्णकर्म्म }\edtext{\skt{हितम्भवेत्}}{\lemma{\skt{भवेत्}}\Afootnote{H; \skt{भवेत्। भल्ला\appbrk{}त\appbrk{}का\appbrk{}न्वा\appbrk{}स\appbrk{}येत्तु क्षीरे प्राङ्मू\appbrk{}त्र\appbrk{}भा\appbrk{}वि\appbrk{}तान्। ततो द्विधा च्छेदयित्वा लौहे कुम्भे निधापयेत्}\ A}}\skt{~।}\pend

\pstart \ledsidenote{\passref{4.1.93}}\edtext{\skt{भल्ला\-त\-क\-वि\-धा\-नेन सा च}}{\lemma{\skt{भल्ला\appbrk{}त\appbrk{}क\appbrk{}वि\appbrk{}धा\appbrk{}नेन सा च}}\Afootnote{H; \skt{तैलेनानेन संसृष्टं शुक्ल\appbrk{}मा\appbrk{}ले\appbrk{}प\appbrk{}ये\appbrk{}द्व्र\appbrk{}णम्। भल्ला\appbrk{}त\appbrk{}क\appbrk{}वि\appbrk{}धा\appbrk{}नेन सार}\ A}}\skt{ }\edtext{\edtext{\skt{स्नेहा\-न्प्र\-यो\-ज\-येत्}}{\lemma{\skt{स्नेहान्}}\Afootnote{H; \skt{स्नेहांस्}\ A}}}{\lemma{\skt{प्रयोजयेत्}}\Afootnote{H; \skt{तु कारयेत्}\ A}}\skt{~॥}\pend

\pstart\versepar \ledsidenote{\passref{4.1.94}}\skt{ये च केचित्फलस्नेहा विधानन्तेषु }\edtext{\skt{वक्ष्यते}}{\lemma{\skt{वक्ष्यते}}\Afootnote{H; \skt{पूर्ववत्।}\ A}}\skt{~॥ }\pend
\pstart\versepar \skt{दुरूढत्वात्तु कृष्णानां }\edtext{\skt{हिता}}{\lemma{\skt{हिता}}\Afootnote{H; \emph{om.}\ A}}\skt{ }\edtext{\skt{पाण्डुक्रिया}}{\lemma{\skt{पाण्डुक्रिया}}\Afootnote{H; \skt{पाण्डुकर्म हितं}\ A}}\skt{ भवेत्~।}\pend

\pstart\versepar \ledsidenote{\passref{4.1.95}}\edtext{\skt{सप्तरात्रस्थितं}}{\lemma{\skt{सप्तरात्रस्थितं}}\Afootnote{H; \skt{सप्तरात्रं स्थितं}\ A}}\skt{ क्षीरे }\edtext{\skt{वार्क्करे}}{\lemma{\skt{वार्क्करे}}\Afootnote{H; \skt{छागले}\ A}}\skt{ रोहिणीफलं~॥ }\pend
\pstart\versepar \skt{तेनैव }\edtext{\skt{पिष्ट्वा}}{\lemma{\skt{पिष्ट्वा}}\Afootnote{H; \skt{पिष्टं}\ A}}\skt{ सुश्लक्ष्णं }\edtext{\skt{कृष्णम्पाण्डूकरं}}{\lemma{\skt{कृष्णम्पाण्डूकरं}}\Afootnote{H; \skt{सवर्णकर}\ A}}\skt{ }\edtext{\skt{स्मृतं}}{\lemma{\skt{स्मृतं}}\Afootnote{H; \skt{णं हितम्}\ A}}\skt{~।}\pend

\pstart\versepar \ledsidenote{\passref{4.1.96}}\skt{नवं कपालिकाचूर्ण्णं }\edtext{\skt{वैदलं}}{\lemma{\skt{वैदलं}}\Afootnote{H; \skt{वैदुलं}\ A}}\skt{ }\edtext{\edtext{\skt{सर्ज्जमेव}}{\lemma{\skt{सर्ज्जम्}}\Afootnote{H; \skt{सर्ज}\ A}}\skt{ च}}{\lemma{\skt{एव च}}\Afootnote{H; \skt{नाम च}\ A}}\skt{~॥ }\pend
\pstart\versepar \skt{}\edtext{\skt{कासीसमधुकं}}{\lemma{\skt{कासीसमधुकं}}\Afootnote{H; \skt{कासीसं मधुकं}\ A}}\skt{ चैव }\edtext{\skt{क्षौद्राक्तं}}{\lemma{\skt{क्षौद्राक्तं}}\Afootnote{H; \skt{क्षौद्रयुक्तं}\ A}}\skt{ }\edtext{\skt{लेपने}}{\lemma{\skt{लेपने}}\Afootnote{H; \skt{प्रलेपये}\ A}}\skt{ }\edtext{\skt{हितं}}{\lemma{\skt{हितं}}\Afootnote{H; \skt{त्}\ A}}\skt{~।}\pend

\pstart\versepar \ledsidenote{\passref{4.1.97}}\edtext{\skt{कुक्कु\-टा\-ण्ड\-क\-पा\-लानि}}{\lemma{\skt{कुक्कु\appbrk{}टा\appbrk{}ण्ड\appbrk{}क\appbrk{}पा\appbrk{}लानि}}\Afootnote{H; \emph{om.}\ A}}\skt{ }\edtext{\skt{कतकं}}{\lemma{\skt{कतकं}}\Afootnote{H; \emph{om.}\ A}}\skt{ }\edtext{\skt{मधुरन्तथा}}{\lemma{\skt{मधुरन्तथा}}\Afootnote{H; \emph{om.}\ A}}\skt{~॥ }\pend
\pstart\versepar \skt{}\edtext{\skt{तथा}}{\lemma{\skt{तथा}}\Afootnote{H; \skt{कपि}\ A}}\skt{ }\edtext{\skt{समुद्रमण्डूकी}}{\lemma{\skt{समुद्रमण्डूकी}}\Afootnote{H; \skt{त्थिमुद्धृते मां}\ A}}\skt{ }\edtext{\edtext{\skt{मशिचूर्ण्णञ्च}}{\lemma{\skt{मशिचूर्ण्णञ्}}\Afootnote{H; \skt{से मूत्रेणा}\ A}}}{\lemma{\skt{च}}\Afootnote{H; \skt{जेन}\ A}}\skt{ }\edtext{\skt{दापयेत्}}{\lemma{\skt{दापयेत्}}\Afootnote{H; \skt{पूरयेत्। कासीसं रोचनां तुत्थं हरितालं मनःशिलाम्}\ A}}\skt{~।}\pend

\pstart\versepar \ledsidenote{\passref{4.1.98}}\edtext{\skt{कपि\-त्था\-न्मा\-न्स\-मु\-द्धृत्य}}{\lemma{\skt{कपि\appbrk{}त्था\appbrk{}न्मा\appbrk{}न्स\appbrk{}मु\appbrk{}द्धृत्य}}\Afootnote{H; \skt{वेणुनिर्लेखनं चापि प्रपु\appbrk{}न्ना\appbrk{}ड\appbrk{}र\appbrk{}साञ्}\ A}}\skt{ }\edtext{\skt{मूत्रेणाजस्य}}{\lemma{\skt{मूत्रेणाजस्य}}\Afootnote{H; \skt{जनम्~। अध\appbrk{}स्ता\appbrk{}द\appbrk{}र्जु\appbrk{}नस्यै}\ A}}\skt{ }\edtext{\skt{पूरयेत्}}{\lemma{\skt{पूरयेत्}}\Afootnote{H; \skt{तन्मासं भूमौ निधापयेत्}\ A}}\skt{~। }\pend
\pstart\versepar \skt{}\edtext{\skt{कासीसं}}{\lemma{\skt{कासीसं}}\Afootnote{H; \emph{om.}\ A}}\skt{ }\edtext{\skt{रोचना}}{\lemma{\skt{रोचना}}\Afootnote{H; \emph{om.}\ A}}\skt{ }\edtext{\skt{चैव}}{\lemma{\skt{चैव}}\Afootnote{H; \emph{om.}\ A}}\skt{ }\edtext{\skt{हरितालम्मनः}}{\lemma{\skt{हरितालम्मनः}}\Afootnote{H; \emph{om.}\ A}}\skt{ }\edtext{\skt{शिला}}{\lemma{\skt{शिला}}\Afootnote{H; \emph{om.}\ A}}\skt{~॥}\pend

\pstart\versepar \ledsidenote{\passref{4.1.99}}\edtext{\skt{वेणु\-नि\-र्ल्ले\-ख\-न\-ञ्चैव}}{\lemma{\skt{वेणु\appbrk{}नि\appbrk{}र्ल्ले\appbrk{}ख\appbrk{}न\appbrk{}ञ्चैव}}\Afootnote{H; \emph{om.}\ A}}\skt{ }\edtext{\skt{प्रपुनाटञ्च}}{\lemma{\skt{प्रपुनाटञ्च}}\Afootnote{H; \emph{om.}\ A}}\skt{ दापयेत्~। }\pend
\pstart\versepar \skt{अध\-स्ता\-द\-र्ज्जु\-न\-स्यै\-त\-न्मांसं }\edtext{\skt{भूमौ निधापयेत्}}{\lemma{\skt{भूमौ निधापयेत्}}\Afootnote{H; \skt{दूर्ध्वं ततस्तेन कृष्ण\appbrk{}मा\appbrk{}ले\appbrk{}प\appbrk{}ये\appbrk{}द्व्र\appbrk{}णम्~। कुक्कु\appbrk{}टा\appbrk{}ण्ड\appbrk{}क\appbrk{}पा\appbrk{}लानि कतकं मधुकं समम्}\ A}}\skt{~॥}\edtext{}{\lemma{\skt{दापयेत्~। अध\appbrk{}स्ता\appbrk{}द\appbrk{}र्ज्जु\appbrk{}न\appbrk{}स्यै\appbrk{}त\appbrk{}न्मांसं}}\Afootnote{H; \skt{मासा}\ A}}\pend

\pstart \ledsidenote{\passref{4.1.100}}\edtext{\skt{मासा\-दू\-र्द्ध\-न्त\-त\-स्तेन}}{\lemma{\skt{मासा\appbrk{}दू\appbrk{}र्द्ध\appbrk{}न्त\appbrk{}त\appbrk{}स्तेन}}\Afootnote{H; \skt{तथा समुद्रमण्डूकी मणिचूर्णं च}\ A}}\skt{ }\edtext{\edtext{\edtext{\skt{कृष्ण\-मा\-ले\-प\-ये\-द्व्रणः}}{\lemma{\skt{कृष्णम्}}\Afootnote{H; \emph{om.}\ A}}}{\lemma{\skt{आलेपयेद्}}\Afootnote{H; \skt{दापयेत्}\ A}}}{\lemma{\skt{व्रणः}}\Afootnote{H; \skt{गुटिका मूत्रपिष्टास्ता व्रणानां प्रतिसारणम्}\ A}}\skt{~॥}\pend

\pstart\versepar \ledsidenote{\passref{4.1.101}}\skt{हस्तिदन्तमसीं कृत्वा मुख्यञ्चैव रसाञ्जनं~। }\pend
\pstart\versepar \skt{}\edtext{\skt{रोमाण्यनेन}}{\lemma{\skt{रोमाण्यनेन}}\Afootnote{H; \skt{रोमाण्येतेन}\ A}}\skt{ जायन्ते }\edtext{\skt{तूर्ण्णं}}{\lemma{\skt{तूर्ण्णं}}\Afootnote{H; \skt{लेपात्}\ A}}\skt{ पाणितलेष्वपि~॥}\pend

\pstart\versepar \ledsidenote{\passref{4.1.102}}\edtext{\skt{चतु\-ष्प\-दा\-ना\-न्त्व\-ग्रो\-म\-न\-ख\-शृ\-ङ्गा\-स्थि\-भ\-स्मना}}{\lemma{\skt{चतु\appbrk{}ष्प\appbrk{}दा\appbrk{}ना\appbrk{}न्त्व\appbrk{}ग्रो\appbrk{}म\appbrk{}न\appbrk{}ख\appbrk{}शृ\appbrk{}ङ्गा\appbrk{}स्थि\appbrk{}भ\appbrk{}स्मना}}\Afootnote{H; \skt{चतुष्पदानां त्वग्रो\appbrk{}म\appbrk{}खु\appbrk{}र\appbrk{}शृ\appbrk{}ङ्गा\appbrk{}स्थि\appbrk{}भ\appbrk{}स्मना}\ A}}\skt{~। }\pend
\pstart\versepar \skt{तैलाक्ता चूर्ण्णिता }\edtext{\skt{भूमि\-र्भ\-वे\-द्रो\-म\-रुहा}}{\lemma{\skt{भूमि\appbrk{}र्भ\appbrk{}वे\appbrk{}द्रो\appbrk{}म\appbrk{}रुहा}}\Afootnote{H; \skt{भूमि\appbrk{}र्भ\appbrk{}वे\appbrk{}द्रो\appbrk{}म\appbrk{}वती}\ A}}\skt{ पुनः~॥}\pend

\pstart\versepar \ledsidenote{\passref{4.1.103}}\skt{कासी\-स\-न्न\-क्त\-मा\-लस्य }\edtext{\skt{पल्लवाश्चापि}}{\lemma{\skt{चापि}}\Afootnote{H; \skt{चैव}\ A}}\skt{ संहरेत्~। }\pend
\pstart\versepar \skt{कपि\-त्थ\-र\-स\-पि\-ष्टानि रोम\-स\-ञ्ज\-न\-न\-म्परं~॥}\pend

\pstart\versepar \ledsidenote{\passref{4.1.104}}\skt{रोमाकीर्ण्णो व्रणो यस्तु न सम्य\-गु\-प\-रो\-हति~। }\pend
\pstart\versepar \skt{क्षुरकर्त्तरि }\edtext{\skt{सन्दंशैस्तस्य}}{\lemma{\skt{सन्दंशैस्}}\Afootnote{H; \skt{सन्दंशौस्}\ A}}\skt{ }\edtext{\skt{लोमानि}}{\lemma{\skt{लोमानि}}\Afootnote{H; \skt{रोमाणि}\ A}}\skt{ निर्हरेत्~॥}\pend

\pstart\versepar \ledsidenote{\passref{4.1.105}}\edtext{\edtext{\skt{शंखचूर्ण्णञ्च}}{\lemma{\skt{शंखचूर्ण्णञ्}}\Afootnote{H; \skt{शङ्खचूर्ण}\ A}}}{\lemma{\skt{च}}\Afootnote{H; \skt{स्य}\ A}}\skt{ }\edtext{\skt{भागौष्टौ}}{\lemma{\skt{भागौष्टौ}}\Afootnote{H; \skt{भागौ द्वौ}\ A}}\skt{ हरितालं च भागिकं~। }\pend
\pstart\versepar \skt{शुक्तेन सह }\edtext{\skt{पिष्टन्तु}}{\lemma{\skt{पिष्टन्तु}}\Afootnote{H; \skt{पिष्टानि}\ A}}\skt{ }\edtext{\skt{रोम\-शा\-त\-न\-मु\-त्तमं}}{\lemma{\skt{रोमशातनम्}}\Afootnote{H; \skt{लोमशातनम्}\ A}}\skt{~॥}\pend

\pstart\versepar \ledsidenote{\passref{4.1.106}}\skt{तैलं भल्लातकस्याथ स्नुही\-क्षी\-र\-न्त\-थैव च~। }\pend
\pstart\versepar \skt{}\edtext{\skt{विमृद्यैकत्र}}{\lemma{\skt{विमृद्यैकत्र}}\Afootnote{H; \skt{प्रगृह्यैकत्र}\ A}}\skt{ }\edtext{\skt{दातव्यं}}{\lemma{\skt{दातव्यं}}\Afootnote{H; \skt{मतिमान्}\ A}}\skt{ रोम\-शा\-त\-न\-मु\-त्तमं~॥}\pend

\pstart\versepar \ledsidenote{\passref{4.1.109}}\skt{वातदुष्टो व्रणो }\edtext{\skt{यस्तु}}{\lemma{\skt{यस्तु}}\Afootnote{H; \skt{यस्तु}\ A}}\skt{ }\edtext{\skt{रूक्ष\-स्या\-त्य\-र्थ\-वे\-दनः}}{\lemma{\skt{रूक्ष\appbrk{}स्या\appbrk{}त्य\appbrk{}र्थ\appbrk{}वे\appbrk{}दनः}}\Afootnote{H; \skt{रूक्ष\appbrk{}श्चा\appbrk{}त्य\appbrk{}र्थ\appbrk{}वे\appbrk{}दनः}\ A}}\skt{~। }\pend
\pstart\versepar \skt{अधःकाये विशेषेण तत्र }\edtext{\edtext{\skt{बस्तिर्विधीयते}}{\lemma{\skt{बस्तिर्}}\Afootnote{A; \skt{वस्ति}\ H}}}{\lemma{\skt{विधीयते}}\Afootnote{A; \skt{र्विधीयते}\ H}}\skt{~॥}\pend

\pstart\versepar \ledsidenote{\passref{4.1.110}}\skt{मूत्राघाते मूत्रदोषे शुक्रदोषे }\edtext{\skt{तथैव च}}{\lemma{\skt{तथैव च}}\Afootnote{H; \skt{अश्मरीव्रणे}\ A}}\skt{~। }\pend
\pstart\versepar \skt{}\edtext{\skt{तथै\-वा\-र्त्त\-व\-दो\-षेषु}}{\lemma{\skt{तथै\appbrk{}वा\appbrk{}र्त्त\appbrk{}व\appbrk{}दो\appbrk{}षेषु}}\Afootnote{H; \skt{तथै\appbrk{}वा\appbrk{}र्त\appbrk{}व\appbrk{}दोषे च}\ A}}\skt{ }\edtext{\edtext{\edtext{\skt{बस्तिरप्युत्तरो}}{\lemma{\skt{बस्तिर्}}\Afootnote{A; \skt{वस्ति}\ H}}}{\lemma{\skt{अप्य्}}\Afootnote{A; \skt{रप्}\ H}}}{\lemma{\skt{उत्तरो}}\Afootnote{A; \skt{युत्तरो}\ H}}\skt{ हितः~॥}\pend

\pstart\versepar \ledsidenote{\passref{4.1.111}}\edtext{\skt{व्रणः}}{\lemma{\skt{व्रणः}}\Afootnote{H; \skt{यस्माच्छ्}\ A}}\skt{ }\edtext{\skt{शुध्यभिवन्धेन}}{\lemma{\skt{शुध्यभिवन्धेन}}\Afootnote{H; \skt{उध्यति बन्धेन}\ A}}\skt{ }\edtext{\skt{मृदु\-त्व\-ञ्चो\-प\-ग\-च्छति}}{\lemma{\skt{मृदु\appbrk{}त्व\appbrk{}ञ्चो\appbrk{}प\appbrk{}ग\appbrk{}च्छति}}\Afootnote{H; \skt{व्रणो याति च मा}\ A}}\skt{~। }\pend
\pstart\versepar \skt{}\edtext{\skt{रोहत्यपि}}{\lemma{\skt{रोहत्य्}}\Afootnote{H; \skt{र्दवम्~। रोहत्य्}\ A}}\skt{ च }\edtext{\skt{निःस\-ङ्गा\-त्त\-स्मा\-द्वन्धः}}{\lemma{\skt{निःसङ्गात्तस्माद्}}\Afootnote{H; \skt{निःशङ्कस्तस्माद्}\ A}}\skt{ }\edtext{\skt{प्रशस्यते}}{\lemma{\skt{प्रशस्यते}}\Afootnote{H; \skt{विधीयते}\ A}}\skt{~॥}\pend

\pstart\versepar \ledsidenote{\passref{4.1.112}}\skt{स्थिरा\-णा\-म\-ल्प\-मां\-सानां रौक्ष्या\-द\-नु\-प\-रो\-हतां~। }\pend
\pstart\versepar \skt{पत्र\-दा\-न\-म्भ\-वे\-त्कार्यं यथादोषं }\edtext{\skt{यथर्त्तु}}{\lemma{\skt{यथर्त्तु}}\Afootnote{H; \skt{यथर्तु}\ A}}\skt{ च~॥}\pend

\pstart\versepar \ledsidenote{\passref{4.1.113}}\edtext{\skt{एर\-ण्ड\-भू\-र्ज्ज\-पू\-ती\-क\-का\-श्म\-र्या\-णाञ्च}}{\lemma{\skt{एर\appbrk{}ण्ड\appbrk{}भू\appbrk{}र्ज्ज\appbrk{}पू\appbrk{}ती\appbrk{}क\appbrk{}का\appbrk{}श्म\appbrk{}र्या\appbrk{}णाञ्च}}\Afootnote{H; \skt{एर\appbrk{}ण्ड\appbrk{}भू\appbrk{}र्ज\appbrk{}पू\appbrk{}ती\appbrk{}क\appbrk{}ह\appbrk{}रि\appbrk{}द्राणां तु}\ A}}\skt{ }\edtext{\skt{वातजः}}{\lemma{\skt{वातजः}}\Afootnote{H; \skt{वातजे}\ A}}\skt{~। }\pend
\pstart\versepar \skt{पत्रमाश्ववलं यच्च }\edtext{\skt{हरि\-द्रा\-प\-त्र\-मेव}}{\lemma{\skt{हरिद्रापत्रम्}}\Afootnote{H; \skt{काश्मरीपत्रम्}\ A}}\skt{ च~॥}\pend

\pstart\versepar \ledsidenote{\passref{4.1.114}}\skt{पत्राणि }\edtext{\skt{क्षीर\-वृ\-क्षा\-णा\-मौ\-द\-का\-ना\-न्त\-थैव च}}{\lemma{\skt{क्षीर\appbrk{}वृ\appbrk{}क्षा\appbrk{}णा\appbrk{}मौ\appbrk{}द\appbrk{}का\appbrk{}ना\appbrk{}न्त\appbrk{}थैव च}}\Afootnote{H; \skt{क्षीर\appbrk{}वृ\appbrk{}क्षा\appbrk{}णा\appbrk{}मौ\appbrk{}द\appbrk{}कानि तथैव च}\ A}}\skt{~। }\pend
\pstart\versepar \skt{दूषिते रक्त\-पि\-त्ता\-भ्या\-म्व्रणे दद्याद्विचक्षणः~॥}\pend

\pstart\versepar \ledsidenote{\passref{4.1.115}}\edtext{\skt{पाठामूर्व्वाश्च}}{\lemma{\skt{पाठामूर्व्वाश्च}}\Afootnote{H; \skt{पाठामूर्वागु}\ A}}\skt{ }\edtext{\skt{गन्धानां}}{\lemma{\skt{गन्धानां}}\Afootnote{H; \skt{डूचीनां}\ A}}\skt{ काक\-मा\-ची\-ह\-रि\-द्रयोः~। }\pend
\pstart\versepar \skt{}\edtext{\skt{पत्रं}}{\lemma{\skt{पत्रं}}\Afootnote{H; \skt{पत्रं च}\ A}}\skt{ }\edtext{\skt{शुकनसायाश्च}}{\lemma{\skt{शुकनसायाश्च}}\Afootnote{H; \skt{शुकनासाया}\ A}}\skt{ योजयेत्कफजे व्रणे~॥}\pend

\pstart\versepar \ledsidenote{\passref{4.1.116}}\edtext{\skt{अक\-र्क्क\-श\-म\-वि\-च्छि\-न्न\-म\-जी\-र्ण्ण\-म्वि\-ष\-व\-र्ज्जितं}}{\lemma{\skt{अक\appbrk{}र्क्क\appbrk{}श\appbrk{}म\appbrk{}वि\appbrk{}च्छि\appbrk{}न्न\appbrk{}म\appbrk{}जी\appbrk{}र्ण्ण\appbrk{}म्वि\appbrk{}ष\appbrk{}व\appbrk{}र्ज्जितं}}\Afootnote{H; \skt{अक\appbrk{}र्क\appbrk{}श\appbrk{}म\appbrk{}वि\appbrk{}च्छि\appbrk{}न्न\appbrk{}म\appbrk{}जीर्णं सुकुमारकम्}\ A}}\skt{~। }\pend
\pstart\versepar \skt{अजन्तु जग्धम्मृदु च पत्रं }\edtext{\skt{गुणवदुच्यते}}{\lemma{\skt{गुणवदुच्यते}}\Afootnote{\skt{गुणवदुच्यते}\ A H}}\skt{~॥}\pend

\pstart\versepar \ledsidenote{\passref{4.1.117}}\edtext{\skt{स्नेह\-मौ\-ष\-ध\-सा\-रम्वा}}{\lemma{\skt{स्नेह\appbrk{}मौ\appbrk{}ष\appbrk{}ध\appbrk{}सा\appbrk{}रम्वा}}\Afootnote{\skt{स्नेह\appbrk{}मौ\appbrk{}ष\appbrk{}ध\appbrk{}सारं च}\ A; \skt{स्नेह\appbrk{}मौ\appbrk{}ष\appbrk{}ध\appbrk{}सा\appbrk{}रम्वा}\ H}}\skt{ }\edtext{\skt{पट्ट\-प\-त्रा\-न्त\-री\-कृतं}}{\lemma{\skt{पट्ट\appbrk{}प\appbrk{}त्रा\appbrk{}न्त\appbrk{}री\appbrk{}कृतं}}\Afootnote{H; \skt{पट्टः पत्रान्तरीकृतः}\ A}}\skt{~। }\pend
\pstart\versepar \skt{}\edtext{\skt{तद्ददाति}}{\lemma{\skt{तद्ददाति}}\Afootnote{\skt{गादत्ते}\ A; \skt{तद्ददाति}\ H}}\skt{ }\edtext{\edtext{\skt{यत\-स्त\-स्मा\-दौ\-ष\-धो\-परि}}{\lemma{\skt{यतस्}}\Afootnote{\skt{यत्ततः}\ A; \skt{य}\ H}}}{\lemma{\skt{तस्मा\appbrk{}दौ\appbrk{}ष\appbrk{}धो\appbrk{}परि}}\Afootnote{\skt{पत्रं लेपस्योपरि}\ A; \skt{तस्त\appbrk{}स्मा\appbrk{}दौ\appbrk{}ष\appbrk{}धो\appbrk{}परि}\ H}}\skt{ दापयेत्~॥}\pend

\pstart\versepar \ledsidenote{\passref{4.1.118}}\edtext{\skt{शीतो\-ष्ण\-ज\-न\-ना\-र्थाय}}{\lemma{\skt{शीतो\appbrk{}ष्ण\appbrk{}ज\appbrk{}न\appbrk{}ना\appbrk{}र्थाय}}\Afootnote{H; \skt{शैत्यौ\appbrk{}ष्ण्य\appbrk{}ज\appbrk{}न\appbrk{}ना\appbrk{}र्थाय}\ A}}\skt{ स्नेह\-सं\-ग्र\-ह\-णाय च~। }\pend
\pstart\versepar \skt{दत्तौषधेषु दातव्यम्पत्रं वैद्येन जानता~॥}\pend

\pstart\versepar \ledsidenote{\passref{4.1.119}}\skt{मक्षि\-का\-व्र\-ण\-मा\-गत्य }\edtext{\skt{निक्षिपन्ति}}{\lemma{\skt{निक्षिपन्ति}}\Afootnote{H; \skt{निःक्षिपन्ति}\ A}}\skt{ यदा }\edtext{\skt{क्रिमिं}}{\lemma{\skt{क्रिमिं}}\Afootnote{H; \skt{कृमीन्}\ A}}\skt{~। }\pend
\pstart\versepar \skt{}\edtext{\skt{श्वयथुर्ब्भक्षिते तैस्तु}}{\lemma{\skt{श्वयथुर्ब्भक्षिते तैस्तु}}\Afootnote{H; \skt{श्वयथुर्भक्षिते तैस्तु}\ A}}\skt{ जायते }\edtext{\skt{भृशदारुणं}}{\lemma{\skt{भृशदारुणं}}\Afootnote{H; \skt{भृशदारुणः}\ A}}\skt{~॥}\pend

\pstart\versepar \ledsidenote{\passref{4.1.120}}\skt{तीव्रा रुजो विचित्राश्च }\edtext{\skt{रक्तस्रावं}}{\lemma{\skt{रक्तस्रावं}}\Afootnote{H; \skt{रक्तास्रावश्}\ A}}\skt{ }\edtext{\skt{स्रवेदपि}}{\lemma{\skt{स्रवेदपि}}\Afootnote{H; \skt{च}\ A}}\skt{~। }\pend
\pstart\versepar \skt{}\edtext{\skt{स्वसु\-र\-सा\-दि\-हि\-त\-स्तत्र}}{\lemma{\skt{स्वसु\appbrk{}र\appbrk{}सा\appbrk{}दि\appbrk{}हि\appbrk{}तस्}}\Afootnote{H\textsuperscript{pc}; \skt{जायते~। सुर\appbrk{}सा\appbrk{}दि\appbrk{}र्हि\appbrk{}तस्}\ A; \skt{स्वर\appbrk{}सा\appbrk{}दि\appbrk{}हि\appbrk{}तस्}\ H\textsuperscript{ac}}}\skt{ धावने पूरणे तथा~॥}\pend

\pstart\versepar \ledsidenote{\passref{4.1.121}}\skt{सप्त\-प\-र्ण्ण\-क\-र\-ञ्जार्क्क निम्व\-रा\-जा\-द\-न\-त्वचः~। }\pend
\pstart\versepar \skt{}\edtext{\skt{गोमूत्रपिष्टा च}}{\lemma{\skt{गोमूत्रपिष्टा च}}\Afootnote{H; \skt{हिता गोमूत्रपिष्टाश्च}\ A}}\skt{ }\edtext{\skt{हिताः}}{\lemma{\skt{हिताः}}\Afootnote{H; \emph{om.}\ A}}\skt{ सेकः क्षारोदकेन वा~॥}\pend

\pstart\versepar \ledsidenote{\passref{4.1.122}}\edtext{\skt{व्रणं}}{\lemma{\skt{व्रणं}}\Afootnote{H; \emph{om.}\ A}}\skt{ }\edtext{\skt{प्रच्छादयित्वा च}}{\lemma{\skt{प्रच्छादयित्वा च}}\Afootnote{H; \skt{प्रच्छाद्य}\ A}}\skt{ }\edtext{\edtext{\skt{मांस\-पे\-ष्या\-म्वि\-नि\-र्ह\-रेत्}}{\lemma{\skt{मांसपेष्याम्}}\Afootnote{H; \skt{मांसपेश्या}\ A}}}{\lemma{\skt{विनिर्हरेत्}}\Afootnote{H; \skt{वा कृमी\appbrk{}न\appbrk{}प\appbrk{}ह\appbrk{}रे\appbrk{}द्व्र\appbrk{}णात्}\ A}}\skt{~। }\pend
\pstart\versepar \skt{}\edtext{\skt{विंशति}}{\lemma{\skt{विंशति}}\Afootnote{H; \skt{विंशतिं}\ A}}\skt{ }\edtext{\skt{क्रिमिजातीस्तु}}{\lemma{\skt{क्रिमिजातीस्तु}}\Afootnote{H; \skt{कृमिजातीस्तु}\ A}}\skt{ }\edtext{\edtext{\skt{वक्ष्याम्युपरि}}{\lemma{\skt{वक्ष्याम्य्}}\Afootnote{H; \skt{वक्स्याम्}\ A}}}{\lemma{\skt{उपरि}}\Afootnote{H; \skt{युपरि}\ A}}\skt{ भागशः~॥}\pend

\pstart\versepar \ledsidenote{\passref{4.1.123}}\edtext{\skt{दीर्घ\-का\-तु\-रा\-णान्तु}}{\lemma{\skt{दीर्घ\appbrk{}का\appbrk{}तु\appbrk{}रा\appbrk{}णान्तु}}\Afootnote{\skt{दीर्घ\appbrk{}का\appbrk{}ला\appbrk{}तु\appbrk{}राणां तु}\ A; \skt{दीर्घ\appbrk{}का…\appbrk{}तु\appbrk{}रा\appbrk{}णान्तु}\ H}}\skt{ }\edtext{\skt{कृषा\-णा\-म्व्र\-ण\-शो\-षिणां}}{\lemma{\skt{कृषाणाम्}}\Afootnote{H; \skt{कृशानां}\ A}}\skt{~। }\pend
\pstart\versepar \skt{वृंहणीयो विधिः }\edtext{\skt{कार्यः }\edtext{\skt{सोऽप्य\-ग्नि\-म्प\-रि\-र\-क्षता}}{\lemma{\skt{ऽप्यग्निम्}}\Afootnote{H; \skt{याग्निं}\ A}}}{\lemma{\skt{कार्यः सो}}\Afootnote{H; \skt{सर्वः का}\ A}}\skt{~॥}\pend

\pstart\versepar \ledsidenote{\passref{4.1.124}}\edtext{\skt{विषदुष्टस्य}}{\lemma{\skt{विषदुष्टस्य}}\Afootnote{H; \skt{विषजुष्टस्य}\ A}}\skt{ }\edtext{\skt{विज्ञानं}}{\lemma{\skt{विज्ञानं}}\Afootnote{H; \skt{विज्ञानां}\ A}}\skt{ विष\-नि\-श्च\-य\-मेव च~। }\pend
\pstart\versepar \skt{}\edtext{\skt{चिकित्सितन्तु}}{\lemma{\skt{चिकित्सितन्तु}}\Afootnote{H; \skt{चिकित्सितं च}\ A}}\skt{ वक्ष्यामि }\edtext{\skt{तङ्कल्पेषु}}{\lemma{\skt{तङ्कल्पेषु}}\Afootnote{H; \skt{कल्पेषु}\ A}}\skt{ }\edtext{\skt{विभागशः}}{\lemma{\skt{विभागशः}}\Afootnote{H; \skt{प्रतिभागशः}\ A}}\skt{~॥}\pend

\pstart\versepar \ledsidenote{\passref{4.1.125}}\skt{कण्डूमन्तः }\edtext{\skt{सशोफश्च}}{\lemma{\skt{सशोफश्च}}\Afootnote{H; \skt{सशोफाश्च ये च}\ A}}\skt{ }\edtext{\skt{जत्रूणोपरि ये}}{\lemma{\skt{जत्रूणोपरि ये}}\Afootnote{H; \skt{जत्रूपरि}\ A}}\skt{ व्रणाः~। }\pend
\pstart\versepar \skt{शिरोविरेचनं तेषु }\edtext{\skt{विधेयं}}{\lemma{\skt{विधेयं}}\Afootnote{H; \skt{विदध्यात्कुशलो}\ A}}\skt{ }\edtext{\skt{वक्ष्यते}}{\lemma{\skt{वक्ष्यते}}\Afootnote{H; \skt{भिषक्}\ A}}\skt{ }\edtext{\skt{तथा}}{\lemma{\skt{तथा}}\Afootnote{H; \emph{om.}\ A}}\skt{~॥}\pend

\pstart\versepar \ledsidenote{\passref{4.1.126}}\skt{रुजा\-व\-न्तोऽ\-नि\-ला\-विष्टा रूक्षा ये चोर्द्ध्वजत्रुजाः~। }\pend
\pstart\versepar \skt{व्रणेषु तेषु कर्त्तव्यं नस्यं }\edtext{\skt{स्नेहेन}}{\lemma{\skt{स्नेहेन}}\Afootnote{H; \skt{वैद्येन}\ A}}\skt{ जानता~॥}\pend

\pstart\versepar \ledsidenote{\passref{4.1.127}}\edtext{\edtext{\skt{दोष\-म्प्र\-च्या\-व\-णा\-र्थाय}}{\lemma{\skt{दोषम्}}\Afootnote{H; \skt{दोष}\ A}}}{\lemma{\skt{प्रच्यावणार्थाय}}\Afootnote{H; \skt{प्रच्यावनार्थाय}\ A}}\skt{ रुजा\-दा\-ह\-क्ष\-याय च~। }\pend
\pstart\versepar \skt{जिह्वा\-द\-न्त\-स\-मु\-त्थस्य हर\-णा\-र्थ\-म्म\-लस्य च~॥}\pend

\pstart\versepar \ledsidenote{\passref{4.1.128}}\skt{शोधनो रोपणश्चैव व्रणस्य मुखजस्य वै~। }\pend
\pstart\versepar \skt{उष्णो वा यदि वा शीतः कवलग्रह इष्यते~॥}\pend

\pstart\versepar \ledsidenote{\passref{4.1.129}}\edtext{\edtext{\edtext{\skt{ऊर्द्ध्व\-ज\-त्रु\-ग\-ता\-न्त्रा\-गा\-न्व्र\-णान्वा}}{\lemma{\skt{ऊर्द्ध्व\appbrk{}ज\appbrk{}त्रु\appbrk{}ग\appbrk{}तान्}}\Afootnote{H; \skt{ऊर्ध्वजत्रुगता}\ A}}}{\lemma{\skt{त्रागान्}}\Afootnote{H; \skt{न्रोगान्}\ A}}}{\lemma{\skt{व्रणान्वा}}\Afootnote{H; \skt{व्रणांश्च}\ A}}\skt{ कफवातजान्~। }\pend
\pstart\versepar \skt{}\edtext{\skt{शोफा\-स्रा\-व\-रु\-जा\-यु\-क्ता\-न्धू\-म\-पा\-नै\-रु\-पा\-च\-रेत्}}{\lemma{\skt{शोफा\appbrk{}स्रा\appbrk{}व\appbrk{}रु\appbrk{}जा\appbrk{}यु\appbrk{}क्तान्}}\Afootnote{H; \skt{शोफ\appbrk{}स्रा\appbrk{}व\appbrk{}रु\appbrk{}जा\appbrk{}यु\appbrk{}क्तान्}\ A}}\skt{~॥}\pend

\pstart\versepar \ledsidenote{\passref{4.1.130}}\edtext{\skt{क्षतोष्मणां}}{\lemma{\skt{क्षतोष्मणां}}\Afootnote{H; \skt{क्षतोष्मणो}\ A}}\skt{ निग्रहार्थं सन्धा\-ना\-र्थ\-न्त\-थैव च~। }\pend
\pstart\versepar \skt{}\edtext{\skt{सद्यो\-व्र\-णे\-ष्वा\-य\-तेषु}}{\lemma{\skt{सद्यो\appbrk{}व्र\appbrk{}णे\appbrk{}ष्वा\appbrk{}य\appbrk{}तेषु}}\Afootnote{H; \skt{सद्यो\appbrk{}व्र\appbrk{}णे\appbrk{}ष्वा\appbrk{}य\appbrk{}तेषु}\ A}}\skt{ }\edtext{\skt{मधु\-स\-र्प्पि\-र्व्वि\-धी\-यते}}{\lemma{\skt{मधु\appbrk{}स\appbrk{}र्प्पि\appbrk{}र्व्वि\appbrk{}धी\appbrk{}यते}}\Afootnote{H; \skt{क्षौद्र\appbrk{}स\appbrk{}र्पि\appbrk{}र्वि\appbrk{}धी\appbrk{}यते}\ A}}\skt{~॥}\pend

\pstart\versepar \ledsidenote{\passref{4.1.131}}\edtext{\skt{अव\-गा\-ढ\-स्त्व\-णु\-मुखा ये}}{\lemma{\skt{अव\appbrk{}गा\appbrk{}ढ\appbrk{}स्त्व\appbrk{}णु\appbrk{}मुखा ये}}\Afootnote{H; \skt{अव\appbrk{}गा\appbrk{}ढा\appbrk{}स्त्व\appbrk{}णु\appbrk{}मुखा ये}\ A}}\skt{ व्रणाः शल्यपीडिताः~। }\pend
\pstart\versepar \skt{निवृ\-त्त\-ह\-स्तो\-द्ध\-रणा यन्त्रं तेषु विधीयते~॥}\pend

\pstart\versepar \ledsidenote{\passref{4.1.132}}\skt{लघुमात्रो लघुश्चैव }\edtext{\skt{स्निग्धाष्णो}}{\lemma{\skt{स्निग्धाष्णो}}\Afootnote{H; \skt{स्निग्ध उष्णो}\ A}}\skt{ }\edtext{\skt{वह्निदीपनं}}{\lemma{\skt{वह्निदीपनं}}\Afootnote{H; \skt{ऽग्निदीपनः}\ A}}\skt{~। }\pend
\pstart\versepar \skt{सर्व्वव्रणिभ्यो देयस्तु }\edtext{\skt{सदा\-हा\-र\-म्वि\-जा\-नता}}{\lemma{\skt{सदाहारम्}}\Afootnote{H; \skt{सदाऽऽआहारो}\ A}}\skt{~॥}\pend

\pstart\versepar \ledsidenote{\passref{4.1.133}}\skt{निशाचरेभ्यो रक्ष्यस्तु नित्यमेव क्षतातुरः~। }\pend
\pstart\versepar \skt{रक्षा\-वि\-धा\-नै\-रु\-द्दिष्टैः }\edtext{\skt{सदा}}{\lemma{\skt{सदा}}\Afootnote{H; \skt{यमैः सनि\appbrk{}य\appbrk{}मै\appbrk{}स्त\appbrk{}तहा}\ A}}\skt{ }\edtext{\skt{चापि}}{\lemma{\skt{चापि}}\Afootnote{H; \emph{om.}\ A}}\skt{ }\edtext{\skt{प्रवक्ष्यते}}{\lemma{\skt{प्रवक्ष्यते}}\Afootnote{H; \emph{om.}\ A}}\skt{~॥}\pend

\pstart\versepar \ledsidenote{\passref{4.1.134}}\edtext{\skt{षण्मू\-लो\-ष्ट\-प\-रि\-ग्राही}}{\lemma{\skt{षण्मू\appbrk{}लो\appbrk{}ष्ट\appbrk{}प\appbrk{}रि\appbrk{}ग्राही}}\Afootnote{H; \skt{षण्मू\appbrk{}लोऽ\appbrk{}ष्ट\appbrk{}प\appbrk{}रि\appbrk{}ग्राही}\ A}}\skt{ पञ्च\-ल\-क्ष\-ण\-ल\-क्षितः~। }\pend
\pstart\versepar \skt{}\edtext{\skt{षष्ट्याविधानैः}}{\lemma{\skt{षष्ट्याविधानैः}}\Afootnote{H; \skt{षष्ट्या चिधानैर्}\ A}}\skt{ }\edtext{\skt{प्राप्नोति}}{\lemma{\skt{प्राप्नोति}}\Afootnote{H; \skt{निर्दि}\ A}}\skt{ }\edtext{\skt{चतुर्ब्भिः}}{\lemma{\skt{चतुर्ब्भिः}}\Afootnote{H; \skt{ष्टैश्चतुर्भिः}\ A}}\skt{ साध्यते व्रणः~॥}\pend

\pstart\versepar \ledsidenote{\passref{4.1.135}}\edtext{\skt{अल्पग्रन्थं}}{\lemma{\skt{अल्पग्रन्थं}}\Afootnote{H; \emph{om.}\ A}}\skt{ }\edtext{\skt{मृगयता}}{\lemma{\skt{मृगयता}}\Afootnote{H; \emph{om.}\ A}}\skt{ }\edtext{\skt{योगो }\edtext{\skt{योऽल्पौषधः}}{\lemma{\skt{ऽल्पौषधः}}\Afootnote{H; \skt{ऽल्पौषध}\ A}}}{\lemma{\skt{योगो यो}}\Afootnote{H; \skt{यो}\ A}}\skt{ कृतः~। }\pend
\pstart\versepar \skt{}\edtext{\skt{द्रव्याणां तु}}{\lemma{\skt{द्रव्याणां तु}}\Afootnote{H; \skt{योगो बहु\appbrk{}ग्र\appbrk{}न्थ\appbrk{}भ\appbrk{}या\appbrk{}न्मया~। द्रव्याणां त}\ A}}\skt{ }\edtext{\skt{समानानान्तत्र}}{\lemma{\skt{समानानान्तत्र}}\Afootnote{H; \skt{त्समानानां तत्रा}\ A}}\skt{ }\edtext{\skt{वायो न}}{\lemma{\skt{वायो न}}\Afootnote{H; \skt{वापो न}\ A}}\skt{ दुष्यति~॥}\pend

\pstart\versepar \ledsidenote{\passref{4.1.136}}\skt{प्रसंगाभिहितो यो वा बहु\-दु\-र्ल्ल\-भ\-भे\-षजः~। }\pend
\pstart\versepar \skt{}\edtext{\skt{यथोपपत्या}}{\lemma{\skt{यथोपपत्या}}\Afootnote{H; \skt{यथोपपत्ति}\ A}}\skt{ तत्रापि कार्यमेव चिकित्सितं~॥}\pend

\pstart\versepar \ledsidenote{\passref{4.1.140}}\edtext{\skt{व्रण\-क्रि\-या\-स्वि\-मा\-स्वेव}}{\lemma{\skt{व्रण\appbrk{}क्रि\appbrk{}या\appbrk{}स्वि\appbrk{}मा\appbrk{}स्वेव}}\Afootnote{H; \skt{व्रणक्रियास्वेव}\ A}}\skt{ }\edtext{\skt{समासोक्तास्तु}}{\lemma{\skt{समासोक्तास्तु}}\Afootnote{H; \skt{मासु व्यासेनोक्तास्व्}\ A}}\skt{ }\edtext{\skt{या}}{\lemma{\skt{या}}\Afootnote{H; \skt{अपि}\ A}}\skt{ }\edtext{\skt{क्रिया}}{\lemma{\skt{क्रिया}}\Afootnote{H; \skt{क्रियाम्}\ A}}\skt{~। }\pend
\pstart\versepar \skt{}\edtext{\edtext{\edtext{\skt{भूयोऽ\-भि\-धा\-स्या\-म्यु\-परि}}{\lemma{\skt{भूयो}}\Afootnote{A; \skt{भूयो}\ H}}}{\lemma{\skt{ऽभिधास्याम्य्}}\Afootnote{\skt{ऽप्}\ A; \skt{भिधास्याम्य्}\ H}}}{\lemma{\skt{उपरि}}\Afootnote{H; \skt{युपरि}\ A}}\skt{ }\edtext{\skt{तथा}}{\lemma{\skt{तथा}}\Afootnote{H; \skt{वक्ष्यामि}\ A}}\skt{ }\edtext{\skt{सद्योव्रणेषु ता}}{\lemma{\skt{सद्योव्रणेषु ता}}\Afootnote{H; \skt{सद्योव्रणचिकि}\ A}}\skt{ }\edtext{\skt{इति}}{\lemma{\skt{इति}}\Afootnote{H; \skt{त्सिते}\ A}}\skt{~॥}\pend

\pstart \ledsidenote{\passref{4.1.trailer}}\edtext{\skt{॥ द्विव्रणीयं}}{\lemma{\skt{॥ द्विव्रणीयं}}\Afootnote{H; \skt{इति सुश्रु\appbrk{}त\appbrk{}सं\appbrk{}हि\appbrk{}तायां चिकित्सास्थाने द्विव्रणीय}\ A}}\skt{ चिकित्सितं }\edtext{\skt{प्रथमः}}{\lemma{\skt{प्रथमः}}\Afootnote{H; \skt{नाम प्रथमोऽध्यायः}\ A}}\skt{~॥}\pend

\endnumbering
\beginnumbering
\eledsection*{Adhyāya 2}
\pstart \ledsidenote{\passref{4.2.1}}\skt{अथातः सद्यो\-व्र\-ण\-चि\-कि\-त्सितं व्याख्यास्यामः~॥}\pend

\pstart\versepar \ledsidenote{\passref{4.2.3}}\skt{धन्व\-न्त\-रि\-र्द्ध\-र्म्म\-भृ\-ता\-म्व\-रिष्ठो }\edtext{\edtext{\skt{वाग्विदाम्वरः}}{\lemma{\skt{वाग्विदाम्}}\Afootnote{H; \skt{वाग्विशा}\ A}}}{\lemma{\skt{वरः}}\Afootnote{H; \skt{रदः}\ A}}\skt{~। }\pend
\pstart\versepar \skt{}\edtext{\edtext{\skt{विश्वा\-मि\-त्रा\-त्म\-ज\-मृषिं}}{\lemma{\skt{विश्वा\appbrk{}मि\appbrk{}त्रा\appbrk{}त्म\appbrk{}जम्}}\Afootnote{H; \skt{विश्वामित्र}\ A}}}{\lemma{\skt{ऋषिं}}\Afootnote{H; \skt{सुतं}\ A}}\skt{ शिष्यं }\edtext{\skt{सुश्रु\-त\-म\-न्व\-शात्}}{\lemma{\skt{सुश्रुतम्}}\Afootnote{H; \skt{ऋषिं सुश्रुतम्}\ A}}\skt{~॥}\pend

\pstart\versepar \ledsidenote{\passref{4.2.4}}\skt{नानाधारामुखैः शस्त्रै\-र्न्ना\-ना\-स्था\-ननि पातितैः~। }\pend
\pstart\versepar \skt{}\edtext{\skt{भवंति}}{\lemma{\skt{भवंति}}\Afootnote{H; \emph{om.}\ A}}\skt{ }\edtext{\skt{नानाकृतयो}}{\lemma{\skt{नानाकृतयो}}\Afootnote{H; \skt{नानारूपा}\ A}}\skt{ }\edtext{\edtext{\skt{व्रणास्तास्त्व}}{\lemma{\skt{व्रणास्तास्}}\Afootnote{H; \skt{व्रणा}\ A}}\skt{ निबोध मे}}{\lemma{\skt{त्व निबोध मे}}\Afootnote{H; \skt{ये स्युस्तेषां वक्ष्यामि लक्षणम्}\ A}}\skt{~॥}\pend

\pstart\versepar \ledsidenote{\passref{4.2.5}}\skt{आय\-ता\-श्च\-तु\-र\-स्राश्च त्र्यस्रा }\edtext{\skt{मुण्डलिनस्तथा}}{\lemma{\skt{मुण्डलिनस्}}\Afootnote{H; \skt{मण्डलिनस्}\ A}}\skt{~। }\pend
\pstart\versepar \skt{अर्द्ध\-च\-न्द्र\-प्र\-ती\-काशा विशालाः कुटिलास्तथा~॥}\pend

\pstart\versepar \ledsidenote{\passref{4.2.6}}\skt{शरा\-व\-नि\-म्न\-म\-ध्याश्च }\edtext{\skt{यव\-म\-ध्या\-स्त\-था\-परे}}{\lemma{\skt{तथापरे}}\Afootnote{H; \skt{तथाऽपरे}\ A}}\skt{~। }\pend
\pstart\versepar \skt{एव\-म्प्र\-का\-रा\-कृ\-तयो भवन्त्यागन्तवो व्रणाः~॥}\pend

\pstart\versepar \ledsidenote{\passref{4.2.7}}\skt{दोषजा वा स्वयम्भिन्ना न तु वैद्यनिमित्तजाः~। }\pend
\pstart\versepar \skt{}\edtext{\skt{व्रणाकृतिज्ञो हि}}{\lemma{\skt{व्रणाकृतिज्ञो हि}}\Afootnote{H; \skt{भिष\appbrk{}ग्व्र\appbrk{}णा\appbrk{}कृ\appbrk{}तिज्ञो हि}\ A}}\skt{ }\edtext{\skt{भिषक्न}}{\lemma{\skt{भिषक्न}}\Afootnote{H; \skt{न}\ A}}\skt{ }\edtext{\skt{मोह\-मु\-प\-ग\-च्छति}}{\lemma{\skt{उपगच्छति}}\Afootnote{H; \skt{अधिगच्छति}\ A}}\skt{~॥}\pend

\pstart\versepar \ledsidenote{\passref{4.2.8}}\edtext{\edtext{\skt{दुर्द्दर्शनेष्वपि}}{\lemma{\skt{दुर्द्दर्शनेष्व्}}\Afootnote{H; \skt{भृशं दुर्दर्शरूपेषु}\ A}}}{\lemma{\skt{अपि}}\Afootnote{H; \emph{om.}\ A}}\skt{ }\edtext{\skt{भृशं}}{\lemma{\skt{भृशं}}\Afootnote{H; \emph{om.}\ A}}\skt{ व्रणेषु विकृतेष्वपि~। }\pend
\pstart\versepar \skt{}\edtext{\skt{अने\-का\-कृ\-ति\-रा\-गन्तुः}}{\lemma{\skt{अनेकाकृतिर्}}\Afootnote{H; \skt{अनन्ताकृतिर्}\ A}}\skt{ स भिषग्भिः पुरातनैः~॥}\pend

\pstart\versepar \ledsidenote{\passref{4.2.9}}\skt{समासतो लक्षणतः षड्विधः परिकीर्त्तितः~। }\pend
\pstart\versepar \skt{छिन्नम्भिन्नन्तथा विद्धं }\edtext{\skt{क्षत\-म्पि\-चि\-त\-मेव}}{\lemma{\skt{पिचितम्}}\Afootnote{H; \skt{पिच्चितम्}\ A}}\skt{ च~॥}\pend

\pstart\versepar \ledsidenote{\passref{4.2.10}}\skt{घृष्टमाहुस्तथा षष्ठ\-न्ते\-षा\-म्व\-क्ष्यामि लक्षणं~॥ }\pend
\pstart\versepar \skt{}\edtext{\skt{तिर्यक्चीनो}}{\lemma{\skt{तिर्यक्चीनो}}\Afootnote{H; \skt{तिरश्चीन}\ A}}\skt{ }\edtext{\edtext{\skt{मृदुर्व्वापि}}{\lemma{\skt{मृदुर्}}\Afootnote{H; \skt{ऋजु}\ A}}\skt{ सो}}{\lemma{\skt{व्वापि सो}}\Afootnote{H; \skt{र्वाऽपि यो}\ A}}\skt{ }\edtext{\skt{व्रणस्त्वायतो}}{\lemma{\skt{त्वायतो}}\Afootnote{H; \skt{चायतो}\ A}}\skt{ भवेत्~।}\pend

\pstart\versepar \ledsidenote{\passref{4.2.11}}\edtext{\skt{गात्रपाती न वा}}{\lemma{\skt{गात्रपाती न वा}}\Afootnote{H; \skt{गात्रस्य पातनं चा}\ A}}\skt{ }\edtext{\skt{तद्धि}}{\lemma{\skt{तद्धि}}\Afootnote{H; \skt{पि}\ A}}\skt{ }\edtext{\edtext{\skt{छिन्न\-ल\-क्ष\-ण\-मु\-च्यते}}{\lemma{\skt{छिन्नलक्षणम्}}\Afootnote{H; \skt{छिन्न\appbrk{}मि\appbrk{}त्यु\appbrk{}प\appbrk{}दि\appbrk{}श्यते~। कुन्त}\ A}}}{\lemma{\skt{उच्यते}}\Afootnote{H; \emph{om.}\ A}}\skt{~॥ }\pend
\pstart\versepar \skt{}\edtext{\skt{शक्तिकुन्तेषु }\edtext{\edtext{\skt{खड्गा\-दि\-वि\-षा\-णै\-रा\-श\-यो\-र्हतः}}{\lemma{\skt{आशयोर्}}\Afootnote{H; \skt{आशयः}\ A}}}{\lemma{\skt{हतः}}\Afootnote{H; \emph{om.}\ A}}}{\lemma{\skt{शक्तिकुन्तेषु खड्गा\appbrk{}दि\appbrk{}वि\appbrk{}षा\appbrk{}णैर्}}\Afootnote{H\textsuperscript{pc}; \skt{शक्त्यृ\appbrk{}ष्टि\appbrk{}ख\appbrk{}ङ्गा\appbrk{}ग्र\appbrk{}वि\appbrk{}षा\appbrk{}णा\appbrk{}दि\appbrk{}भिर्}\ A; \skt{शक्ति\appbrk{}कु\appbrk{}न्ते\appbrk{}ख\appbrk{}ड्गा\appbrk{}दि\appbrk{}वि\appbrk{}षा\appbrk{}णैर्}\ H\textsuperscript{ac}}}\skt{~।}\pend

\pstart\versepar \ledsidenote{\passref{4.2.12}}\edtext{\edtext{\skt{यत्किञ्चित्स्रवते}}{\lemma{\skt{यत्}}\Afootnote{H; \skt{हतः}\ A}}}{\lemma{\skt{स्रवते}}\Afootnote{H; \skt{स्रवेत्}\ A}}\skt{ तद्धि भिन्न\-ल\-क्ष\-ण\-मु\-च्यते~॥ }\pend
\pstart\versepar \skt{स्थाना\-न्या\-मा\-ग्नि\-प\-क्वानां मूत्रस्य रुधिरस्य च~।}\pend

\pstart\versepar \ledsidenote{\passref{4.2.13}}\skt{हृदुण्डुकः फुप्फुसश्च }\edtext{\skt{कोष्ठ\-मि\-त्य\-भि\-धी\-यते}}{\lemma{\skt{कोष्ठम्}}\Afootnote{H; \skt{कोष्ठ}\ A}}\skt{~॥ }\pend
\pstart\versepar \skt{तस्मिन्भिन्ने रक्तपूर्ण्णे ज्वरो दाहश्च जायते~।}\pend

\pstart\versepar \ledsidenote{\passref{4.2.14}}\skt{मूत्र\-मा\-र्ग्ग\-गु\-दा\-स्येभ्यो रक्तं घ्राणाच्च गच्छति~॥ }\pend
\pstart\versepar \skt{}\edtext{\edtext{\skt{मूर्च्छा\-ध्मा\-न\-तृ\-षा\-श्वा\-स\-स्त्व\-भ\-क्त\-च्छन्द}}{\lemma{\skt{मूर्च्छा\appbrk{}ध्मा\appbrk{}न\appbrk{}तृ\appbrk{}षा\appbrk{}श्वा\appbrk{}सस्}}\Afootnote{H; \skt{मूर्च्छा\appbrk{}श्वा\appbrk{}स\appbrk{}तृ\appbrk{}डा\appbrk{}ध्मान}\ A}}}{\lemma{\skt{त्वभक्तच्छन्द}}\Afootnote{H\textsuperscript{pc}; \skt{मभक्तच्छन्द}\ A; \skt{त्वभक्तच्छन्न}\ H\textsuperscript{ac}}}\skt{ एव च~।}\pend

\pstart\versepar \ledsidenote{\passref{4.2.15}}\edtext{\skt{विण्मू\-त्र\-वा\-त\-स\-ङ्गञ्च}}{\lemma{\skt{विण्मू\appbrk{}त्र\appbrk{}वा\appbrk{}त\appbrk{}स\appbrk{}ङ्गञ्च}}\Afootnote{H; \skt{विण्मू\appbrk{}त्र\appbrk{}वा\appbrk{}त\appbrk{}स\appbrk{}ङ्गश्च}\ A}}\skt{ स्वेदास्रावो क्षिरक्तता~॥ }\pend
\pstart\versepar \skt{लोह\-ग\-न्धि\-त्व\-मा\-स्यस्य }\edtext{\skt{गात्र\-दौ\-र्ग्ग\-न्ध\-मेव}}{\lemma{\skt{गात्रदौर्ग्गन्धम्}}\Afootnote{H; \skt{गात्रदौर्गन्ध्यम्}\ A}}\skt{ च~।}\pend

\pstart\versepar \ledsidenote{\passref{4.2.16}}\skt{हृच्छूलं पार्श्वयोश्चापि विशेषं चात्र मे शृणु~॥ }\pend
\pstart\versepar \skt{आमाशयस्थे रुधिरे रुधिरं }\edtext{\edtext{\skt{छर्द्दयत्यपि}}{\lemma{\skt{छर्द्दयत्य्}}\Afootnote{H; \skt{संगृह्य समुत्थाप्य महाबलाः~। भवत्य्}\ A}}}{\lemma{\skt{अपि}}\Afootnote{H; \skt{अन्तःप्रवेशस्तु यथा निर्धुनुयुस्तथा}\ A}}\skt{~।}\pend

\pstart\versepar \ledsidenote{\passref{4.2.17}}\edtext{\edtext{\skt{आध्मा\-न\-म\-ति\-मात्रं}}{\lemma{\skt{आध्मानम्}}\Afootnote{H; \skt{आध्मान}\ A}}\skt{ च}}{\lemma{\skt{अतिमात्रं च}}\Afootnote{H; \skt{मतिमात्रं च}\ A}}\skt{ शूलं च भृशदारुणम्~॥ }\pend
\pstart\versepar \skt{}\edtext{\skt{पक्वासयगते}}{\lemma{\skt{पक्वासयगते}}\Afootnote{H; \skt{पक्वाशयगते}\ A}}\skt{ चापि }\edtext{\skt{रुजा}}{\lemma{\skt{रुजा}}\Afootnote{H; \skt{रुजो}\ A}}\skt{ गौरवमेव च~।}\pend

\pstart \ledsidenote{\passref{4.2.18}}\edtext{\skt{अधःकाये}}{\lemma{\skt{अधःकाये}}\Afootnote{H; \emph{om.}\ A}}\skt{ }\edtext{\skt{विशेषेण}}{\lemma{\skt{विशेषेण}}\Afootnote{H; \emph{om.}\ A}}\skt{ शीतता }\edtext{\skt{च}}{\lemma{\skt{च}}\Afootnote{H; \skt{चाप्यधो नामेः खेभ्यो रक्तस्य चागमः। अभि\appbrk{}न्ने+\appbrk{}अ\appbrk{}प्या\appbrk{}श\appbrk{}ये+\appbrk{}अ\appbrk{}न्त्राणां खैः सूक्ष्मैरन्त्र}\ A}}\skt{ }\edtext{\skt{भवन्ति}}{\lemma{\skt{भवन्ति}}\Afootnote{H; \skt{पूरणम्}\ A}}\skt{ }\edtext{\skt{हि}}{\lemma{\skt{हि}}\Afootnote{H; \emph{om.}\ A}}\skt{~॥}\pend

\pstart \ledsidenote{\passref{4.2.19}}\edtext{\skt{सूक्ष्मास्यं}}{\lemma{\skt{सूक्ष्मास्यं}}\Afootnote{H; \skt{पिहितास्ये घटे यद्वल्लक्ष्यते तस्य गौरवम्। सूक्ष्मास्य}\ A}}\skt{ शल्याभिहतं }\edtext{\edtext{\skt{यद\-ग\-न्धा\-स\-य\-म्विना}}{\lemma{\skt{यद\appbrk{}ग\appbrk{}न्धा\appbrk{}स\appbrk{}यम्}}\Afootnote{H; \skt{यदङ्गं त्वाशया}\ A}}}{\lemma{\skt{विना}}\Afootnote{H; \skt{द्विना}\ A}}\skt{~।}\pend

\pstart\versepar \ledsidenote{\passref{4.2.20}}\skt{उत्तु\-ण्डि\-त\-न्नि\-र्ग\-तम्वा तद्विद्धमिति निर्द्दिशेत्~॥ }\pend
\pstart\versepar \skt{}\edtext{\edtext{\skt{नाति\-च्छि\-न्न\-न्ना\-ति\-वि\-द्ध\-मु\-भ\-यो\-र्ल्ल\-क्ष\-णा\-न्वि\-तम्}}{\lemma{\skt{नातिविद्धम्}}\Afootnote{H; \skt{नातिभिन्न}\ A}}}{\lemma{\skt{उभयोर्}}\Afootnote{H; \skt{मुभयोर्}\ A}}\skt{~।}\pend

\pstart\versepar \ledsidenote{\passref{4.2.21}}\edtext{\skt{विष\-म\-व्र\-ण\-संगो यं}}{\lemma{\skt{विष\appbrk{}म\appbrk{}व्र\appbrk{}ण\appbrk{}संगो यं}}\Afootnote{H\textsuperscript{pc}; \skt{विषमं व्रणमङ्गे य}\ A; \skt{विष\appbrk{}म\appbrk{}व्र\appbrk{}ण\appbrk{}मङ्गे यन्}\ H\textsuperscript{ac}}}\skt{ }\pend
\pstart\versepar \skt{}\edtext{\skt{प्रहा\-र\-पी\-ड\-ना\-भ्यान्तु}}{\lemma{\skt{प्रहा\appbrk{}र\appbrk{}पी\appbrk{}ड\appbrk{}ना\appbrk{}भ्यान्तु}}\Afootnote{\skt{त्तत्क्षतं त्वभिनिर्दिशेत्। प्रहा\appbrk{}र\appbrk{}पी\appbrk{}ड\appbrk{}नाभ्यां तु}\ A; \skt{तत्क्ष\appbrk{}त\appbrk{}न्त्व\appbrk{}भि\appbrk{}नि\appbrk{}र्द्दि\appbrk{}शेत्~॥ प्रहा\appbrk{}र\appbrk{}पी\appbrk{}ड\appbrk{}ना\appbrk{}भ्यान्तु}\ H}}\skt{ यदङ्गम्पृथुतां गतं~।}\pend

\pstart\versepar \ledsidenote{\passref{4.2.22}}\skt{सास्थि }\edtext{\edtext{\skt{तत्पि\-च्चि\-त\-म्वि\-द्या\-त्म\-ज्जा\-र\-क्त\-प\-रि\-प्लुतं}}{\lemma{\skt{विद्यात्}}\Afootnote{H; \skt{विद्यान्}\ A}}}{\lemma{\skt{मज्जा\appbrk{}र\appbrk{}क्त\appbrk{}प\appbrk{}रि\appbrk{}प्लुतं}}\Afootnote{H; \skt{मज्ज\appbrk{}र\appbrk{}क्त\appbrk{}प\appbrk{}रिप्लु}\ A}}\skt{~॥ }\pend
\pstart\versepar \skt{}\edtext{\skt{घर्ष\-णा\-द\-भि\-घा\-ताद्वा}}{\lemma{\skt{घर्ष\appbrk{}णा\appbrk{}द\appbrk{}भि\appbrk{}घा\appbrk{}ताद्वा}}\Afootnote{H; \emph{om.}\ A}}\skt{ }\edtext{\edtext{\skt{यद\-ङ्ग\-म्वि\-ग\-त\-त्वचं}}{\lemma{\skt{यदङ्गम्}}\Afootnote{H; \skt{तम्}\ A}}}{\lemma{\skt{विगतत्वचं}}\Afootnote{H; \skt{विग\appbrk{}त\appbrk{}त्व\appbrk{}ग्य\appbrk{}दङ्गं हि संघ\appbrk{}र्षा\appbrk{}द\appbrk{}न्य\appbrk{}था+\appbrk{}अपि वा}\ A}}\skt{~।}\pend

\pstart\versepar \ledsidenote{\passref{4.2.23}}\edtext{\edtext{\skt{ऊषा\-स्रा\-वा\-न्वि\-त\-न्तन्तु}}{\lemma{\skt{ऊषा\appbrk{}स्रा\appbrk{}वा\appbrk{}न्वि\appbrk{}त\appbrk{}न्तन्}}\Afootnote{H; \skt{उषास्रावान्वितं त}\ A}}}{\lemma{\skt{तु}}\Afootnote{H; \skt{त्तु}\ A}}\skt{ }\edtext{\edtext{\edtext{\skt{घृष्ट\-मि\-त्य\-भि\-धी\-यते}}{\lemma{\skt{घृष्टम्}}\Afootnote{H; \skt{घृष्ट}\ A}}}{\lemma{\skt{इत्य्}}\Afootnote{H; \skt{मित्य्}\ A}}}{\lemma{\skt{अभिधीयते}}\Afootnote{H; \skt{उपदिश्यते}\ A}}\skt{~॥ }\pend
\pstart\versepar \skt{छिन्ने भिन्ने तथा विद्धे क्षते }\edtext{\skt{वासृगतिस्रवे}}{\lemma{\skt{वासृगतिस्रवे}}\Afootnote{H; \skt{वा+अ\appbrk{}सृ\appbrk{}ग\appbrk{}ति\appbrk{}स्र\appbrk{}वेत्}\ A}}\skt{~।}\pend

\pstart \ledsidenote{\passref{4.2.24}}\edtext{\skt{रक्तक्षयादत्र}}{\lemma{\skt{रक्तक्षयादत्र}}\Afootnote{H; \skt{रक्त\appbrk{}क्ष\appbrk{}या\appbrk{}द्रु\appbrk{}ज\appbrk{}स्तत्र}\ A; \emph{om.}\ K}}\skt{ }\edtext{\skt{रुजाः}}{\lemma{\skt{रुजाः}}\Afootnote{H; \emph{om.}\ A K}}\skt{ }\edtext{\skt{करोति}}{\lemma{\skt{करोति}}\Afootnote{A H; \emph{om.}\ K}}\skt{ }\edtext{\skt{पवनो}}{\lemma{\skt{पवनो}}\Afootnote{A H; \emph{om.}\ K}}\skt{ }\edtext{\skt{भृशं}}{\lemma{\skt{भृशं}}\Afootnote{A H; \emph{om.}\ K}}\skt{~॥ }\edtext{\skt{स्नेहपानं}}{\lemma{\skt{स्नेहपानं}}\Afootnote{A H; \emph{om.}\ K}}\skt{ हितन्तत्र तत्सेको }\edtext{\edtext{\skt{विहितस्तथा}}{\lemma{\skt{विहितस्}}\Afootnote{K H; \skt{विहित}\ A}}}{\lemma{\skt{तथा}}\Afootnote{K H; \skt{स्तथा}\ A}}\skt{~।}\pend

\pstart\versepar \ledsidenote{\passref{4.2.25}}\skt{वेशवारैः सकृशरैः }\edtext{\edtext{\skt{सुस्नि\-ग्ध\-श्चो\-प\-ना\-हतं}}{\lemma{\skt{सुस्निग्धश्}}\Afootnote{K; \skt{सुस्निग्धैश्}\ A; \skt{सुस्निग्धैश्}\ H}}}{\lemma{\skt{चोपनाहतं}}\Afootnote{K; \skt{चोपनाहनम्}\ A H}}\skt{~। }\pend
\pstart\versepar \skt{धान्यस्वेदांश्च कुर्वीत }\edtext{\skt{स्निग्धा\-न्या\-ले\-प\-नानि}}{\lemma{\skt{स्निग्धा\appbrk{}न्या\appbrk{}ले\appbrk{}प\appbrk{}नानि}}\Afootnote{A H; \skt{स्निग्धा\appbrk{}न्या\appbrk{}ले\appbrk{}प\appbrk{}नानि}\ K}}\skt{ च~॥}\pend

\pstart\versepar \ledsidenote{\passref{4.2.26}}\skt{वात\-घ्नौ\-ष\-ध\-सि\-द्धैश्च }\edtext{\edtext{\skt{स्नेहै\-र्ब\-स्ति\-र्वि\-धी\-यते}}{\lemma{\skt{बस्तिर्}}\Afootnote{K H; \skt{बस्ति}\ A}}}{\lemma{\skt{विधीयते}}\Afootnote{K H; \skt{र्विधीयते}\ A}}\skt{~॥ }\pend
\pstart\versepar \skt{पिच्चिते च विघृष्टे च }\edtext{\skt{नातिस्रवति}}{\lemma{\skt{नातिस्रवति}}\Afootnote{K H; \skt{नातिस्रवपि}\ A}}\skt{ शोणितं~।}\pend

\pstart\versepar \ledsidenote{\passref{4.2.27}}\edtext{\skt{तस्मिन्न}}{\lemma{\skt{तस्मिन्न}}\Afootnote{K H; \skt{अ}\ A}}\skt{ गछति भृशं }\edtext{\skt{दाहः}}{\lemma{\skt{दाहः}}\Afootnote{K H; \skt{तस्मिन्दाहः}\ A}}\skt{ पाकश्च जायते~। }\pend
\pstart\versepar \skt{}\edtext{\skt{तस्योष्मणो}}{\lemma{\skt{तस्योष्मणो}}\Afootnote{K; \skt{तत्रोष्मणो}\ A; \skt{तस्योष्मनो}\ H}}\skt{ निग्रहार्थं }\edtext{\skt{दाहपाकभया य}}{\lemma{\skt{दाहपाकभया य}}\Afootnote{K H; \skt{तथा दाहप्रपाकयोः}\ A}}\skt{ }\edtext{\skt{च}}{\lemma{\skt{च}}\Afootnote{K H; \emph{om.}\ A}}\skt{~।}\pend

\pstart\versepar \ledsidenote{\passref{4.2.28}}\skt{शीतमालेपनं }\edtext{\skt{कार्य\-म्प\-रि\-षे\-क\-श्शी\-तलः}}{\lemma{\skt{परिषेकश्}}\Afootnote{K; \skt{परिषेकश्च}\ A H}}\skt{~। }\pend
\pstart\versepar \skt{षट्स्वेतेषु यथोक्तेषु छिन्नादिषु समासतः~।}\pend

\pstart\versepar \ledsidenote{\passref{4.2.29}}\skt{ज्ञेयं }\edtext{\skt{समर्पितं}}{\lemma{\skt{समर्पितं}}\Afootnote{A K; \skt{ससर्प्पिषं}\ H}}\skt{ सर्वं }\edtext{\skt{सद्यो\-व्र\-ण\-चि\-कि\-त्सि\-तम्~॥ {\symfont ❈}}}{\lemma{\skt{सद्यो\appbrk{}व्र\appbrk{}ण\appbrk{}चि\appbrk{}कि\appbrk{}त्सि\appbrk{}तम्~॥ {\symfont ❈}}}\Afootnote{K H; \skt{सद्यो\appbrk{}व्र\appbrk{}ण\appbrk{}चि\appbrk{}कि\appbrk{}त्सि\appbrk{}तम्}\ A}}\skt{~॥ }\pend
\pstart\versepar \skt{अत ऊर्ध्वम्प्रवक्ष्यामि छिन्नानान्तु चिकित्सितं~।}\pend

\pstart\versepar \ledsidenote{\passref{4.2.30}}\skt{ये व्रणा विवृताः }\edtext{\skt{केचिच्छिन्नाः}}{\lemma{\skt{छिन्नाः}}\Afootnote{K H; \skt{इरः}\ A}}\skt{ पार्श्वावलम्वितः~। }\pend
\pstart\versepar \skt{तां सीव्ये\-द्वि\-धि\-नो\-क्तेन }\edtext{\skt{वन्धी\-या\-द्गा\-ढ\-मेव}}{\lemma{\skt{वन्धीयाद्}}\Afootnote{K; \skt{बध्नीयाद्}\ A H}}\skt{ च~।}\edtext{}{\lemma{\skt{पार्श्वावलम्वितः~। तां}}\Afootnote{K; \skt{पार्श्वावलम्बिनः। तान्}\ A; \skt{पार्श्वावलम्बिनः~॥ तां}\ H}}\pend

\pstart\versepar \ledsidenote{\passref{4.2.31}}\skt{कर्ण्णं }\edtext{\skt{स्थानादपोढन्तु}}{\lemma{\skt{स्थानादपोढन्तु}}\Afootnote{K; \skt{स्थानादपहृतं}\ A; \skt{स्थानादपोढन्तु}\ H}}\skt{ }\edtext{\skt{स्थापयुत्वा}}{\lemma{\skt{स्थापयुत्वा}}\Afootnote{K; \skt{स्थापयित्वा}\ A H}}\skt{ यथास्थितं~। }\pend
\pstart\versepar \skt{}\edtext{\skt{सीव्येद्यथोक्तं}}{\lemma{\skt{सीव्येद्}}\Afootnote{A H; \skt{सीव्येद्}\ K}}\skt{ तैलेन }\edtext{\edtext{\skt{स्रोत\-श्चा\-प्य\-भि\-त\-र्प\-येत्}}{\lemma{\skt{चाप्य्}}\Afootnote{H; \skt{चाभि}\ A; \skt{चा}\ K}}}{\lemma{\skt{अभितर्पयेत्}}\Afootnote{H; \skt{प्रतर्पयेत्}\ A; \skt{प्यभितर्पयेत्}\ K}}\skt{~।}\pend

\pstart\versepar \ledsidenote{\passref{4.2.32}}\skt{कृकाटिकान्ते छिन्ने तु }\edtext{\edtext{\skt{गछत्यपि}}{\lemma{\skt{गछत्य्}}\Afootnote{K H; \skt{गच्छत्}\ A}}}{\lemma{\skt{अपि}}\Afootnote{K H; \skt{यपि}\ A}}\skt{ समीरणे~। }\pend
\pstart\versepar \skt{सम्यङ्निवेश्य वध्नी\-या\-त्सी\-व्ये\-च्चापि निरन्तरम्~।}\pend

\pstart\versepar \ledsidenote{\passref{4.2.33}}\edtext{\skt{आजेन}}{\lemma{\skt{आजेन}}\Afootnote{A K; \skt{आज्येन}\ H}}\skt{ }\edtext{\skt{सर्पिषा}}{\lemma{\skt{सर्पिषा}}\Afootnote{K H; \skt{सपिषा}\ A}}\skt{ }\edtext{\skt{चैवं}}{\lemma{\skt{चैवं}}\Afootnote{A K; \skt{चैव}\ H}}\skt{ }\edtext{\skt{परिषेको}}{\lemma{\skt{परिषेको}}\Afootnote{K H; \skt{परिषेकं}\ A}}\skt{ }\edtext{\skt{भिधीयते}}{\lemma{\skt{भिधीयते}}\Afootnote{K; \skt{तु कारये}\ A; \skt{विधीयते}\ H}}\skt{~। }\pend
\pstart\versepar \skt{}\edtext{\skt{उत्तानोन्नं}}{\lemma{\skt{उत्तानोन्नं}}\Afootnote{K H; \skt{त्। उत्तानोऽन्नं}\ A}}\skt{ }\edtext{\skt{सम\-श्र्नी\-या\-च्छ\-यीत}}{\lemma{\skt{समश्र्नीयाच्}}\Afootnote{K; \skt{समश्नीयाच्छ्}\ A H}}\skt{ च सुयन्त्रितः~।}\pend

\pstart\versepar \ledsidenote{\passref{4.2.34}}\skt{शाखासु }\edtext{\skt{तिर्यक्सतितां}}{\lemma{\skt{तिर्यक्सतितां}}\Afootnote{K; \skt{पतितांस्तिर्यक्}\ A; \skt{तिर्यक्षतिताम्}\ H}}\skt{ }\edtext{\skt{प्रहारां}}{\lemma{\skt{प्रहारां}}\Afootnote{K; \skt{प्रहारान्}\ A; \skt{प्रहराम्}\ H}}\skt{ }\edtext{\skt{विवृतां}}{\lemma{\skt{विवृतां}}\Afootnote{K H; \skt{विवृतान्}\ A}}\skt{ }\edtext{\skt{भृसं}}{\lemma{\skt{भृसं}}\Afootnote{K H; \skt{भृशम्}\ A}}\skt{~। }\pend
\pstart\versepar \skt{}\edtext{\skt{सीव्ये\-त्स\-म्य\-ङ्नि\-वे\-श्याशुः}}{\lemma{\skt{निवेश्याशुः}}\Afootnote{K; \skt{निवेश्याशु}\ A H}}\skt{ }\edtext{\skt{सन्ध्य\-स्थी\-न्य\-नु\-पू\-र्वशः}}{\lemma{\skt{सन्ध्य\appbrk{}स्थी\appbrk{}न्य\appbrk{}नु\appbrk{}पू\appbrk{}र्वशः}}\Afootnote{A; \skt{सन्ध्य\appbrk{}स्थी\appbrk{}न्य\appbrk{}नु\appbrk{}पू\appbrk{}र्वशः}\ K H}}\skt{~।}\pend

\pstart\versepar \ledsidenote{\passref{4.2.35}}\edtext{\skt{वध्या\-वे\-ल्लि\-त\-के\-नाशुः}}{\lemma{\skt{वध्या\appbrk{}वे\appbrk{}ल्लि\appbrk{}त\appbrk{}के\appbrk{}नाशुः}}\Afootnote{K; \skt{बद्ध्वा वेल्लितकेनाशु}\ A; \skt{बध्वा वेलितकेनाशु}\ H}}\skt{ }\edtext{\skt{वस्तान्तेन}}{\lemma{\skt{वस्तान्तेन}}\Afootnote{K; \skt{ततस्तैलेन}\ A; \skt{वस्त्रान्तेन}\ H}}\skt{ }\edtext{\skt{घनेन वा}}{\lemma{\skt{घनेन वा}}\Afootnote{K H; \skt{सेचयेत्}\ A}}\skt{~। }\pend
\pstart\versepar \skt{}\edtext{\skt{चर्मणा गोः}}{\lemma{\skt{चर्मणा गोः}}\Afootnote{K H; \skt{चर्मणा गो}\ A}}\skt{ }\edtext{\skt{फणो}}{\lemma{\skt{फणो}}\Afootnote{K H; \skt{फणा}\ A}}\skt{ वन्धः }\edtext{\skt{कार्यश्वा}}{\lemma{\skt{कार्यश्वा}}\Afootnote{K; \skt{कार्यो}\ A; \skt{कार्यश्चां}\ H}}\skt{ }\edtext{\skt{सहते}}{\lemma{\skt{सहते}}\Afootnote{K; \emph{om.}\ A; \skt{से हते}\ H}}\skt{ }\edtext{\skt{सिना}}{\lemma{\skt{सिना}}\Afootnote{K H; \skt{यो वा हितो भवेत्}\ A}}\skt{~।}\pend

\pstart\versepar \ledsidenote{\passref{4.2.36}}\skt{पृष्ठे व्रणो यस्य }\edtext{\skt{भवेदुत्तानञ्च}}{\lemma{\skt{भवेदुत्तानञ्च}}\Afootnote{K H; \skt{भवेदुत्तानं शा}\ A}}\skt{ }\edtext{\skt{शयीत च}}{\lemma{\skt{शयीत च}}\Afootnote{K H; \skt{ययेत्तु तम्}\ A}}\skt{~। }\pend
\pstart\versepar \skt{}\edtext{\skt{अतोन्यथा}}{\lemma{\skt{अतोन्यथा}}\Afootnote{\skt{अतोऽन्यथा}\ A; \skt{अतो न्यथा}\ K H}}\skt{ }\edtext{\skt{चोरसिजे}}{\lemma{\skt{चोरसिजे}}\Afootnote{A K; \skt{चारसिजे}\ H}}\skt{ }\edtext{\skt{शयीत}}{\lemma{\skt{शयीत}}\Afootnote{K H; \skt{शाययेत्}\ A}}\skt{ }\edtext{\skt{पुरुषो}}{\lemma{\skt{पुरुषो}}\Afootnote{K H; \skt{पुरुषं}\ A}}\skt{ व्रणे~।}\pend

\pstart\versepar \ledsidenote{\passref{4.2.37}}\edtext{\skt{छिन्नां}}{\lemma{\skt{छिन्नां}}\Afootnote{A K; \skt{छिन्ना}\ H}}\skt{ निःशेषतः }\edtext{\edtext{\skt{शाखान्दग्ध्वा}}{\lemma{\skt{शाखान्}}\Afootnote{A K; \skt{शाखा}\ H}}}{\lemma{\skt{दग्ध्वा}}\Afootnote{A K; \skt{दग्धा}\ H}}\skt{ तैलेन }\edtext{\skt{बुद्धिमां}}{\lemma{\skt{बुद्धिमां}}\Afootnote{K; \skt{बुद्धिमान्}\ A H}}\skt{~। }\pend
\pstart\versepar \skt{}\edtext{\skt{वध्नी\-या\-त्को\-श\-व\-न्धेन}}{\lemma{\skt{कोशवन्धेन}}\Afootnote{A K; \skt{कोषबन्धेन}\ H}}\skt{ }\edtext{\skt{प्राप्तं}}{\lemma{\skt{प्राप्तं}}\Afootnote{A K; \skt{प्राप्त}\ H}}\skt{ }\edtext{\skt{कुर्याच्च}}{\lemma{\skt{कुर्याच्च}}\Afootnote{K H; \skt{कार्यं च}\ A}}\skt{ रोपणं~॥}\pend

\pstart\versepar \ledsidenote{\passref{4.2.38}}\skt{चन्दनं पद्मकं }\edtext{\edtext{\skt{लोध्रमुत्पलादि}}{\lemma{\skt{लोध्रम्}}\Afootnote{K H; \skt{रोध्र}\ A}}}{\lemma{\skt{उत्पलादि}}\Afootnote{K H; \skt{मुत्पलानि}\ A}}\skt{ प्रियंगवः~॥ }\pend
\pstart\versepar \skt{हरिद्रा मधुकं चैव }\edtext{\edtext{\skt{पय\-श्चा\-ष्ट\-म\-मेव}}{\lemma{\skt{चाष्टमम्}}\Afootnote{K H; \skt{स्यादत्र चाष्टमम्}\ A}}}{\lemma{\skt{एव}}\Afootnote{K H; \emph{om.}\ A}}\skt{ }\edtext{\skt{तु}}{\lemma{\skt{तु}}\Afootnote{K H; \emph{om.}\ A}}\skt{~।}\pend

\pstart\versepar \ledsidenote{\passref{4.2.39}}\edtext{\edtext{\edtext{\skt{तैल\-मे\-भि\-र्वि\-पक्कं}}{\lemma{\skt{तैलम्}}\Afootnote{K H; \skt{तैल}\ A}}}{\lemma{\skt{एभिर्}}\Afootnote{K H; \skt{मेभिर्}\ A}}}{\lemma{\skt{विपक्कं}}\Afootnote{K; \skt{विपक्वं}\ A H}}\skt{ }\edtext{\skt{स्यात्प्रधानं}}{\lemma{\skt{स्यात्}}\Afootnote{K H; \skt{तु}\ A}}\skt{ व्रणरोपणं~। }\pend
\pstart\versepar \skt{चन्दनं कर्कटाख्या च सहे }\edtext{\skt{मांसाह्वयामृता}}{\lemma{\skt{मांसाह्वयामृता}}\Afootnote{K; \skt{मांस्या\appbrk{}ह्व\appbrk{}या+\appbrk{}अ\appbrk{}मृता}\ A; \skt{मासाह्वयामृताः}\ H}}\skt{~।}\pend

\pstart\versepar \ledsidenote{\passref{4.2.40}}\skt{हरेणवो मृणालञ्च त्रिफला }\edtext{\skt{पद्मकोत्पलैः}}{\lemma{\skt{पद्मकोत्पलैः}}\Afootnote{K H; \skt{पद्मकोत्पले}\ A}}\skt{~। }\pend
\pstart\versepar \skt{}\edtext{\skt{त्रयो\-द\-शां\-ग\-त्रि\-वृ\-त\-मे\-तद्वा}}{\lemma{\skt{त्रयो\appbrk{}द\appbrk{}शां\appbrk{}ग\appbrk{}त्रि\appbrk{}वृ\appbrk{}त\appbrk{}मे\appbrk{}तद्वा}}\Afootnote{K; \skt{त्रयोदशाङ्गं त्रिवृतमेतद्वा}\ A; \skt{त्रयो\appbrk{}द\appbrk{}शा\appbrk{}ङ्ग\appbrk{}न्त्रि\appbrk{}वृ\appbrk{}त\appbrk{}मे\appbrk{}तद्वा}\ H}}\skt{ }\edtext{\edtext{\skt{पयसामान्वितं}}{\lemma{\skt{पयसाम्}}\Afootnote{K H; \skt{पयसा+}\ A}}}{\lemma{\skt{आन्वितं}}\Afootnote{K; \skt{अन्वितम्}\ A; \skt{वितं}\ H}}\skt{~।}\pend

\pstart\versepar \ledsidenote{\passref{4.2.41}}\skt{तैलं विपक्वं }\edtext{\skt{सेकार्यहितन्तु}}{\lemma{\skt{सेकार्यहितन्तु}}\Afootnote{K; \skt{सेकार्थे हितं तु}\ A; \skt{सेकाय हितन्तद्}\ H}}\skt{ }\edtext{\skt{व्रणरोपने~॥ {\symfont ❈}}}{\lemma{\skt{व्रणरोपने~॥ {\symfont ❈}}}\Afootnote{K; \skt{व्रणरोपणे}\ A H}}\skt{~॥ }\pend
\pstart\versepar \skt{अत ऊर्ध्वं प्रवक्ष्यामि भिन्नानान्तु चिकित्सितं~।}\pend

\pstart\versepar \ledsidenote{\passref{4.2.42}}\edtext{\skt{भिन्न\-ने\-त्र\-म\-क\-र्म\-ण्य\-म\-भिन्नं}}{\lemma{\skt{भिन्न\appbrk{}ने\appbrk{}त्र\appbrk{}म\appbrk{}क\appbrk{}र्म\appbrk{}ण्य\appbrk{}म\appbrk{}भिन्नं}}\Afootnote{K; \skt{भिन्नं नेत्र\appbrk{}म\appbrk{}क\appbrk{}र्म\appbrk{}ण्य\appbrk{}म\appbrk{}भिन्नं}\ A; \skt{भिन्न\appbrk{}ने\appbrk{}त्र\appbrk{}म\appbrk{}क\appbrk{}र्म्म\appbrk{}ण्य\appbrk{}म\appbrk{}भिन्नं}\ H}}\skt{ }\edtext{\skt{लम्वते तु यः}}{\lemma{\skt{लम्वते तु यः}}\Afootnote{K H; \skt{लम्बते तु यत्}\ A}}\skt{~। }\pend
\pstart\versepar \skt{तन्निवेश्य }\edtext{\skt{यथा\-स्था\-न\-म\-व्या\-वि\-द्ध\-सिरं}}{\lemma{\skt{यथा\appbrk{}स्था\appbrk{}न\appbrk{}म\appbrk{}व्या\appbrk{}वि\appbrk{}द्ध\appbrk{}सिरं}}\Afootnote{A K; \skt{यथा\appbrk{}स्था\appbrk{}न\appbrk{}म\appbrk{}व्या\appbrk{}वि\appbrk{}द्ध\appbrk{}सिरां}\ H}}\skt{ शनैः~।}\pend

\pstart\versepar \ledsidenote{\passref{4.2.43}}\edtext{\skt{पीडयेन्पाणिना}}{\lemma{\skt{पीडयेन्}}\Afootnote{K; \skt{पीडयेत्}\ A H}}\skt{ }\edtext{\skt{सम्म्य\-क्प\-द्म\-प\-त्रा\-न्त\-रेण}}{\lemma{\skt{सम्म्यक्}}\Afootnote{K; \skt{सम्यक्}\ A H}}\skt{ तु~। }\pend
\pstart\versepar \skt{}\edtext{\skt{ततोस्य}}{\lemma{\skt{ततोस्य}}\Afootnote{K; \skt{ततोऽस्य}\ A H}}\skt{ तर्पणं कार्य\-न्न\-स्य\-ञ्चा\-नेन सर्पिषा~।}\pend

\pstart\versepar \ledsidenote{\passref{4.2.44}}\edtext{\skt{अजाघृतं}}{\lemma{\skt{अजाघृतं}}\Afootnote{K H; \skt{आजं घृतं}\ A}}\skt{ क्षीरपात्रं मधु\-क\-ञ्चो\-त्प\-लानि च~। }\pend
\pstart\versepar \skt{}\edtext{\skt{जीव\-क\-र्स\-ष\-भकौ}}{\lemma{\skt{जीव\appbrk{}क\appbrk{}र्स\appbrk{}ष\appbrk{}भकौ}}\Afootnote{K; \skt{जीवकर्षभकौ}\ A H}}\skt{ }\edtext{\skt{चापि}}{\lemma{\skt{चापि}}\Afootnote{K H; \skt{चैव}\ A}}\skt{ पिष्ट्वा }\edtext{\edtext{\skt{सर्पि\-र्वि\-पा\-च\-येत्}}{\lemma{\skt{सर्पिर्}}\Afootnote{A H; \skt{सर्पि}\ K}}}{\lemma{\skt{विपाचयेत्}}\Afootnote{A; \skt{र्विपाचयेत्}\ K; \skt{व्विपाचयेत्}\ H}}\skt{~।}\pend

\pstart\versepar \ledsidenote{\passref{4.2.45}}\edtext{\skt{सर्व\-ने\-त्रा\-भि\-घा\-तेषु}}{\lemma{\skt{सर्व\appbrk{}ने\appbrk{}त्रा\appbrk{}भि\appbrk{}घा\appbrk{}तेषु}}\Afootnote{K H; \skt{सर्व\appbrk{}ने\appbrk{}त्रा\appbrk{}भि\appbrk{}घाते तु}\ A}}\skt{ }\edtext{\skt{सर्पि\-रे\-त\-त्प्र\-श\-स्यते}}{\lemma{\skt{सर्पिरेतत्}}\Afootnote{A K; \skt{सर्प्पिरेतत्}\ H}}\skt{~। }\pend
\pstart\versepar \skt{}\edtext{\skt{उदरान्मेदसो}}{\lemma{\skt{उदरान्मेदसो}}\Afootnote{A K; \skt{उदरान्मेदसो}\ H}}\skt{ }\edtext{\skt{वर्त्तिर्न्निर्गता}}{\lemma{\skt{वर्त्तिर्न्}}\Afootnote{K; \skt{वर्तिर्}\ A; \skt{वर्त्तिर्}\ H}}\skt{ यस्य देहिनः~।}\pend

\pstart\versepar \ledsidenote{\passref{4.2.46}}\edtext{\skt{तां कषा\-या\-व\-की\-र्ण्णन्तु}}{\lemma{\skt{तां कषा\appbrk{}या\appbrk{}व\appbrk{}की\appbrk{}र्ण्णन्तु}}\Afootnote{K H; \skt{कषा\appbrk{}य\appbrk{}भ\appbrk{}स्म\appbrk{}मृ\appbrk{}त्कीर्णां}\ A}}\skt{ वध्वा सूत्रेण }\edtext{\skt{सूत्रवित्}}{\lemma{\skt{सूत्रवित्}}\Afootnote{K H; \skt{सूत्रवित्। अग्नितप्तेन शस्त्}\ A}}\skt{~। }\pend
\pstart\versepar \skt{}\edtext{\edtext{\edtext{\skt{सम्य\-ग्छि\-न्द्या\-द्व्र\-णेन}}{\lemma{\skt{सम्यग्}}\Afootnote{K; \skt{रेण छिन्द्या\appbrk{}न्म\appbrk{}धु\appbrk{}स\appbrk{}मा\appbrk{}यु\appbrk{}तम्}\ A; \skt{सम्यक्}\ H}}}{\lemma{\skt{छिन्द्याद्}}\Afootnote{K; \emph{om.}\ A; \skt{छिन्द्याद्}\ H}}}{\lemma{\skt{व्रणेन}}\Afootnote{K; \emph{om.}\ A; \skt{व्रणन्}\ H}}\skt{ }\edtext{\skt{तच्च}}{\lemma{\skt{तच्च}}\Afootnote{K H; \emph{om.}\ A}}\skt{ }\edtext{\skt{क्षौद्रेण}}{\lemma{\skt{क्षौद्रेण}}\Afootnote{K; \emph{om.}\ A; \skt{क्षौद्रेणाभ्यज्य}\ H}}\skt{ }\edtext{\skt{वेष्टयेत्}}{\lemma{\skt{वेष्टयेत्}}\Afootnote{K H; \emph{om.}\ A}}\skt{~।}\pend

\pstart\versepar \ledsidenote{\passref{4.2.49}}\edtext{\edtext{\skt{अच्छि\-न्द्य\-मा\-ना\-म\-रणं}}{\lemma{\skt{अच्छिन्द्यमानाम्}}\Afootnote{K; \skt{आटोपो मरणं वा स्याच्छूलो वा+अच्छिद्यमा}\ A; \skt{अच्छिद्यमाना}\ H}}}{\lemma{\skt{अरणं}}\Afootnote{K; \emph{om.}\ A; \skt{मरणं}\ H}}\skt{ }\edtext{\edtext{\skt{कुर्या\-दा\-टो\-प\-मेव}}{\lemma{\skt{कुर्यादाटोपम्}}\Afootnote{K H; \emph{om.}\ A}}\skt{ वा}}{\lemma{\skt{एव वा}}\Afootnote{K H; \skt{नया}\ A}}\skt{~। }\pend
\pstart\versepar \skt{}\edtext{\skt{मेदाग्रन्थौ तु}}{\lemma{\skt{मेदाग्रन्थौ तु}}\Afootnote{K; \skt{मेदोग्रन्थौ तु}\ A H}}\skt{ यत्तैलं वक्ष्यते }\edtext{\skt{तत्तु}}{\lemma{\skt{तत्तु}}\Afootnote{K; \skt{तच्च}\ A; \skt{तन्तु}\ H}}\skt{ योजयेत्~।}\pend

\pstart\versepar \ledsidenote{\passref{4.2.50}}\edtext{\skt{त्वचोतीत्य}}{\lemma{\skt{त्वचोतीत्य}}\Afootnote{K; \skt{त्वचोऽतीत्य}\ A H}}\skt{ सिरादीनि भित्वा वा परिहृत्य वा~। }\pend
\pstart\versepar \skt{}\edtext{\skt{काष्ठे}}{\lemma{\skt{काष्ठे}}\Afootnote{K H; \skt{कोष्ठे}\ A}}\skt{ प्रतिष्ठितं शल्यं कुर्या\-दु\-क्ता\-नु\-प\-द्र\-वान्~।}\pend

\pstart\versepar \ledsidenote{\passref{4.2.51}}\edtext{\skt{तत्रा\-न्त\-र्लो\-हि\-त\-म्पा\-ण्डु\-शी\-त\-पा\-द\-क\-रा\-ननं}}{\lemma{\skt{पाण्डु\appbrk{}शी\appbrk{}त\appbrk{}पा\appbrk{}द\appbrk{}क\appbrk{}रा\appbrk{}ननं}}\Afootnote{K H; \skt{पाण्डुं शीत\appbrk{}पा\appbrk{}द\appbrk{}क\appbrk{}रा\appbrk{}न\appbrk{}नम्}\ A}}\skt{~। }\pend
\pstart\versepar \skt{शीतोच्छ्वासं }\edtext{\skt{रक्त\-ने\-त्र\-मा\-न\-द्धञ्च}}{\lemma{\skt{रक्त\appbrk{}ने\appbrk{}त्र\appbrk{}मा\appbrk{}न\appbrk{}द्धञ्च}}\Afootnote{A K; \skt{रक्त\appbrk{}ने\appbrk{}त्र\appbrk{}मा\appbrk{}न\appbrk{}द्धञ्च}\ H}}\skt{ विवर्जयेत्~।}\pend

\pstart\versepar \ledsidenote{\passref{4.2.52}}\skt{आमाशयस्थे रुधिरे }\edtext{\edtext{\skt{वम\-न\-म्प\-थ्य\-मु\-च्यते}}{\lemma{\skt{पथ्यम्}}\Afootnote{K H; \skt{पथ्य}\ A}}}{\lemma{\skt{उच्यते}}\Afootnote{K H; \skt{मुच्यते}\ A}}\skt{~। }\pend
\pstart\versepar \skt{}\edtext{\skt{पक्वाशयथे}}{\lemma{\skt{पक्वाशयथे}}\Afootnote{K H; \skt{पक्वाशयस्थे}\ A}}\skt{ देयञ्च }\edtext{\edtext{\skt{विरे\-च\-न\-म\-सं\-शयं}}{\lemma{\skt{विरेचनम्}}\Afootnote{K H; \skt{विरेचन}\ A}}}{\lemma{\skt{असंशयं}}\Afootnote{K H; \skt{मसंशयम्}\ A}}\skt{~।}\pend

\pstart\versepar \ledsidenote{\passref{4.2.53}}\edtext{\skt{आस्थापञ्च निः}}{\lemma{\skt{आस्थापञ्च निः}}\Afootnote{K; \skt{आस्थापनं च निः}\ A H}}\skt{ }\edtext{\skt{स्नेहः}}{\lemma{\skt{स्नेहः}}\Afootnote{K H; \skt{स्हेहं}\ A}}\skt{ }\edtext{\edtext{\skt{कार्य\-मु\-ष्ण\-वि\-शो\-धनैः}}{\lemma{\skt{कार्यम्}}\Afootnote{K H; \skt{कार्य}\ A}}}{\lemma{\skt{उष्णविशोधनैः}}\Afootnote{K; \skt{मुष्णैर्विशोधनैः}\ A; \skt{उष्णम्विशोधनैः}\ H}}\skt{~। }\pend
\pstart\versepar \skt{यव\-को\-ल\-कु\-ल\-त्थानां निः स्नेहेन रसेन च~।}\pend

\pstart\versepar \ledsidenote{\passref{4.2.54}}\skt{भुञ्जीतान्नं }\edtext{\skt{यवागुं वा}}{\lemma{\skt{यवागुं वा}}\Afootnote{K; \skt{यवागूं वा}\ A H}}\skt{ }\edtext{\skt{पिवे\-त्से\-न्ध\-व\-सं\-युतां}}{\lemma{\skt{सेन्धवसंयुतां}}\Afootnote{K; \skt{सैन्ध\appbrk{}व\appbrk{}सं\appbrk{}यु\appbrk{}ताम्}\ A H}}\skt{~। }\pend
\pstart\versepar \skt{अतिनिः स्रुतरक्तो वा भिन्नकोष्ठः }\edtext{\skt{पिवेद्सृक्}}{\lemma{\skt{पिवेद्सृक्}}\Afootnote{K; \skt{पिबेदसृक्}\ A H}}\skt{~।}\pend

\pstart\versepar \ledsidenote{\passref{4.2.55}}\edtext{\skt{स्वमा\-र्प्र\-ति\-प\-न्नास्तु}}{\lemma{\skt{स्वमा\appbrk{}र्प्र\appbrk{}ति\appbrk{}प\appbrk{}न्नास्तु}}\Afootnote{K; \skt{स्वमा\appbrk{}र्ग\appbrk{}प्र\appbrk{}ति\appbrk{}प\appbrk{}न्नास्तु}\ A H}}\skt{ यस्य विण्मूत्रमारुताः~। }\pend
\pstart\versepar \skt{}\edtext{\skt{निरूपद्रवः}}{\lemma{\skt{निरूपद्रवः}}\Afootnote{K; \skt{व्युपद्रवः}\ A; \skt{निरुपद्रवः}\ H}}\skt{ }\edtext{\skt{सम्भिन्नेपि}}{\lemma{\skt{सम्भिन्नेपि}}\Afootnote{K; \skt{स भिन्ने+अपि}\ A; \skt{सम्भिन्ने}\ H}}\skt{ कोष्ठे जीवति मानवः~।}\pend

\pstart\versepar \ledsidenote{\passref{4.2.56}}\skt{अभिन्नमन्त्रं निः }\edtext{\skt{क्रान्तां}}{\lemma{\skt{क्रान्तां}}\Afootnote{K H; \skt{क्रान्तं}\ A}}\skt{ प्रवेश्यं नान्यथा भवेत्~। }\pend
\pstart\versepar \skt{अलङ्गले }\edtext{\skt{सिरोग्रस्तं}}{\lemma{\skt{सिरोग्रस्तं}}\Afootnote{K H; \skt{काशिरोग्रस्तं}\ A}}\skt{ }\edtext{\edtext{\skt{तदप्यके}}{\lemma{\skt{तदप्य्}}\Afootnote{H; \skt{तदप्}\ A K}}}{\lemma{\skt{अके}}\Afootnote{\skt{येके}\ A; \skt{यके}\ K; \skt{एके}\ H}}\skt{ वदन्ति }\edtext{\skt{हि}}{\lemma{\skt{हि}}\Afootnote{K H; \skt{तु}\ A}}\skt{~।}\edtext{}{\lemma{\skt{भवेत्~। अलङ्गले}}\Afootnote{K; \skt{भवेत्। पिपीलि}\ A; \skt{भवेत्~॥ अलङ्गल}\ H}}\pend

\pstart\versepar \ledsidenote{\passref{4.2.57}}\skt{प्रक्षाल्य पयसा दिग्धं }\edtext{\skt{तृण\-शो\-णि\-त\-पां\-सुभिः}}{\lemma{\skt{तृण\appbrk{}शो\appbrk{}णि\appbrk{}त\appbrk{}पां\appbrk{}सुभिः}}\Afootnote{K H; \skt{तृण्शो\appbrk{}णि\appbrk{}त\appbrk{}पां\appbrk{}शुभिः}\ A}}\skt{~॥ }\pend
\pstart\versepar \skt{}\edtext{\skt{प्रवे\-श\-ये\-त्कृ\-त्त\-नखो}}{\lemma{\skt{कृत्तनखो}}\Afootnote{A K; \skt{क्लिप्तनखो}\ H}}\skt{ घृतेनाक्तं शनैः शनैः~॥}\pend

\pstart\versepar \ledsidenote{\passref{4.2.58}}\skt{प्रवे\-श\-ये\-त्क्षी\-र\-सिक्तं }\edtext{\skt{शुष्कमन्त्रं}}{\lemma{\skt{शुष्कमन्त्रं}}\Afootnote{K H; \skt{शुष्कमत्रं}\ A}}\skt{ }\edtext{\skt{घृताप्लतं}}{\lemma{\skt{घृताप्लतं}}\Afootnote{K; \skt{घृताप्लुतम्}\ A H}}\skt{~। }\pend
\pstart\versepar \skt{}\edtext{\skt{अंगु\-ल्या\-भि\-मृ\-शे\-त्कण्ठं}}{\lemma{\skt{अंगु\appbrk{}ल्या\appbrk{}भि\appbrk{}मृ\appbrk{}शेत्}}\Afootnote{K H; \skt{अङ्गु\appbrk{}ल्या+\appbrk{}अ\appbrk{}भि\appbrk{}मृ\appbrk{}शेत्}\ A}}\skt{ }\edtext{\skt{जले\-नो\-द्वे\-ज\-ये\-दपि}}{\lemma{\skt{जले\appbrk{}नो\appbrk{}द्वे\appbrk{}ज\appbrk{}ये\appbrk{}दपि}}\Afootnote{A H; \skt{जले\appbrk{}नो\appbrk{}द्वे\appbrk{}ज\appbrk{}ये\appbrk{}दपि}\ K}}\skt{~।}\pend

\pstart\versepar \ledsidenote{\passref{4.2.60}}\edtext{\skt{तथान्त्राणि}}{\lemma{\skt{तथान्त्राणि}}\Afootnote{K H; \skt{तथाऽन्त्राणि}\ A}}\skt{ }\edtext{\skt{विशन्त्यन्तः}}{\lemma{\skt{विशन्त्यन्तः}}\Afootnote{A K; \skt{विशत्यन्तः}\ H}}\skt{ स्वां कलाम्पीडयन्ति च~। }\pend
\pstart\versepar \skt{}\edtext{\skt{व्रण\-सौ\-द्ब\-हु\-त्वाद्वा}}{\lemma{\skt{व्रण\appbrk{}सौ\appbrk{}द्ब\appbrk{}हु\appbrk{}त्वाद्वा}}\Afootnote{K; \skt{व्रणा\appbrk{}ल्प\appbrk{}त्वा\appbrk{}द्ब\appbrk{}हु\appbrk{}त्वाद्वा}\ A; \skt{व्रण\appbrk{}सौ\appbrk{}क्ष्म्या\appbrk{}द्ब\appbrk{}हु\appbrk{}त्वाद्वा}\ H}}\skt{ }\edtext{\skt{दुः प्रवेशन्तु}}{\lemma{\skt{दुः प्रवेशन्तु}}\Afootnote{K H; \skt{दुष्प्रवेशं}\ A}}\skt{ }\edtext{\edtext{\skt{यद्भवेत्}}{\lemma{\skt{यद्}}\Afootnote{K H; \emph{om.}\ A}}}{\lemma{\skt{भवेत्}}\Afootnote{K H; \skt{भवेत्तु यत्}\ A}}\skt{~।}\pend

\pstart\versepar \ledsidenote{\passref{4.2.61}}\edtext{\edtext{\skt{तद्विपाद्य}}{\lemma{\skt{तद्}}\Afootnote{K H; \skt{त}\ A}}}{\lemma{\skt{विपाद्य}}\Afootnote{K H; \skt{दापाठ्य}\ A}}\skt{ प्रमाणेन }\edtext{\skt{भिष\-ग\-न्त\-म्प्र\-वे\-श\-येत्}}{\lemma{\skt{भिष\appbrk{}ग\appbrk{}न्त\appbrk{}म्प्र\appbrk{}वे\appbrk{}श\appbrk{}येत्}}\Afootnote{K; \skt{भिषगन्त्रं पर्वेशयेत्}\ A; \skt{भिष\appbrk{}ग\appbrk{}न्त्र\appbrk{}म्प्र\appbrk{}वे\appbrk{}श\appbrk{}येत्}\ H}}\skt{~। }\pend
\pstart\versepar \skt{}\edtext{\skt{यथा\-स्था\-न\-नि\-विष्टे च}}{\lemma{\skt{यथा\appbrk{}स्था\appbrk{}न\appbrk{}नि\appbrk{}विष्टे च}}\Afootnote{K; \skt{यथास्थानं निवेष्टे च}\ A; \skt{यथा\appbrk{}स्था\appbrk{}न\appbrk{}न्नि\appbrk{}विष्}\ H\textsuperscript{ac}; \skt{यथा\appbrk{}स्था\appbrk{}न\appbrk{}न्वि\appbrk{}शिष्टे च}\ H\textsuperscript{pc}}}\skt{ }\edtext{\skt{व्रणं}}{\lemma{\skt{व्रणं}}\Afootnote{A K; \skt{टे च व्रणं}\ H\textsuperscript{ac}; \skt{व्रणं}\ H\textsuperscript{pc}}}\skt{ }\edtext{\skt{सीव्यमतन्द्रितः}}{\lemma{\skt{सीव्यमतन्द्रितः}}\Afootnote{\skt{सीव्येदतन्द्रियः}\ A; \skt{सीव्यमतन्द्रितः}\ K H}}\skt{~।}\pend

\pstart\versepar \ledsidenote{\passref{4.2.62}}\edtext{\skt{स्थाना\-द\-पे\-त\-मा\-दत्ते}}{\lemma{\skt{स्थाना\appbrk{}द\appbrk{}पे\appbrk{}त\appbrk{}मा\appbrk{}दत्ते}}\Afootnote{A; \skt{स्थाना\appbrk{}द\appbrk{}पे\appbrk{}त\appbrk{}मा\appbrk{}दत्ते}\ K H}}\skt{ }\edtext{\skt{प्राणां}}{\lemma{\skt{प्राणां}}\Afootnote{K H; \skt{प्राणान्}\ A}}\skt{ गुपितमेव वा~। }\pend
\pstart\versepar \skt{}\edtext{\skt{वेष्टयित्वा च}}{\lemma{\skt{वेष्टयित्वा च}}\Afootnote{K H; \skt{वेष्टयित्वा तु}\ A}}\skt{ पट्टेन घृत\-से\-क\-म्प्र\-दा\-प\-येत्~।}\pend

\pstart\versepar \ledsidenote{\passref{4.2.63}}\edtext{\skt{पयः}}{\lemma{\skt{पयः}}\Afootnote{K H; \skt{घृतं}\ A}}\skt{ पिवेत्सुखोष्णञ्च चित्रा\-तै\-ल\-स\-म\-न्वितं~। }\pend
\pstart\versepar \skt{}\edtext{\skt{मृदु\-क्रि\-या\-र्थ\-श\-कृतो}}{\lemma{\skt{मृदु\appbrk{}क्रि\appbrk{}या\appbrk{}र्थ\appbrk{}श\appbrk{}कृतो}}\Afootnote{K; \skt{मृदुक्रियार्थं शकृतो}\ A; \skt{मृदुक्रियार्थं सकृतो}\ H}}\skt{ वायोश्चाधः प्रवृत्तये~।}\pend

\pstart\versepar \ledsidenote{\passref{4.2.64}}\edtext{\edtext{\skt{ततस्तैलमिदं}}{\lemma{\skt{ततस्तैलम्}}\Afootnote{\skt{ततस्तैल}\ A; \skt{ततस्तैलम्}\ K H}}}{\lemma{\skt{इदं}}\Afootnote{K H; \skt{मिदं}\ A}}\skt{ कुर्या\-द्रो\-प\-णा\-र्थ\-ञ्चि\-कि\-त्सकः~। }\pend
\pstart\versepar \skt{}\edtext{\edtext{\skt{त्वचो\-श्व\-क\-र्ण\-ध\-व\-यो\-र्पा\-त\-की\-मे\-ष\-शृ\-ङ्गयोः}}{\lemma{\skt{त्वचो\appbrk{}श्व\appbrk{}क\appbrk{}र्ण\appbrk{}ध\appbrk{}व\appbrk{}योर्}}\Afootnote{K; \skt{त्वचोऽ\appbrk{}श्व\appbrk{}क\appbrk{}र्ण\appbrk{}ध\appbrk{}व\appbrk{}योर्}\ A H}}}{\lemma{\skt{पात\appbrk{}की\appbrk{}मे\appbrk{}ष\appbrk{}शृ\appbrk{}ङ्गयोः}}\Afootnote{K; \skt{मोच\appbrk{}की\appbrk{}मे\appbrk{}ष\appbrk{}शृ\appbrk{}ङ्गयोः}\ A; \skt{द्धात\appbrk{}की\appbrk{}मे\appbrk{}ष\appbrk{}शृ\appbrk{}ङ्गयोः}\ H}}\skt{~।}\pend

\pstart\versepar \ledsidenote{\passref{4.2.65}}\edtext{\skt{शल्ल\-क्य\-र्ज\-न\-यो\-श्चापि}}{\lemma{\skt{शल्ल\appbrk{}क्य\appbrk{}र्ज\appbrk{}न\appbrk{}योश्}}\Afootnote{K; \skt{शल्ल\appbrk{}क्य\appbrk{}र्जु\appbrk{}न\appbrk{}योश्}\ A; \skt{शल्ल\appbrk{}क्या\appbrk{}र्ज्जु\appbrk{}न\appbrk{}योश्}\ H}}\skt{ }\edtext{\skt{विदा\-र्या\-क्षी\-रि\-णा\-न्तथा}}{\lemma{\skt{विदा\appbrk{}र्या\appbrk{}क्षी\appbrk{}रि\appbrk{}णान्}}\Afootnote{K H; \skt{विदार्याः क्षीरिणां}\ A}}\skt{~। }\pend
\pstart\versepar \skt{वलामूलानि चाहृत्य }\edtext{\edtext{\skt{तैल\-मे\-त\-द्वि\-पा\-च\-येत्}}{\lemma{\skt{तैलम्}}\Afootnote{K H; \skt{तैल}\ A}}}{\lemma{\skt{एतद्}}\Afootnote{K H; \skt{मेतैर्}\ A}}\skt{~।}\pend

\pstart\versepar \ledsidenote{\passref{4.2.66}}\skt{व्रणं संरोपयेत्तेन }\edtext{\skt{रक्ते\-त्स\-म्व\-त्स\-रञ्च}}{\lemma{\skt{रक्ते\appbrk{}त्स\appbrk{}म्व\appbrk{}त्स\appbrk{}रञ्च}}\Afootnote{K; \skt{वर्षमात्रं यते}\ A; \skt{रक्षे\appbrk{}त्स\appbrk{}म्व\appbrk{}त्स\appbrk{}रञ्च}\ H}}\skt{ }\edtext{\skt{तत्}}{\lemma{\skt{तत्}}\Afootnote{K H; \skt{त च}\ A}}\skt{~॥ }\pend
\pstart\versepar \skt{पादौ निरस्तमुष्कस्य जलेन }\edtext{\skt{प्रओक्ष्य}}{\lemma{\skt{प्रओक्ष्य}}\Afootnote{K; \skt{प्रोक्ष्य}\ A H}}\skt{ चाक्षिणी~।}\pend

\pstart\versepar \ledsidenote{\passref{4.2.67}}\skt{प्रवेश्य }\edtext{\skt{मुष्कौ}}{\lemma{\skt{मुष्कौ}}\Afootnote{K H; \skt{तुन्}\ A}}\skt{ }\edtext{\skt{सूत्रेण तु}}{\lemma{\skt{सूत्रेण तु}}\Afootnote{K; \emph{om.}\ A; \skt{सूत्रेण तू}\ H}}\skt{ }\edtext{\skt{नो}}{\lemma{\skt{नो}}\Afootnote{K; \skt{नसेवन्या मुष्कौ}\ A; \skt{नोसेवन}\ H}}\skt{ }\edtext{\skt{सेवनसेवितं}}{\lemma{\skt{सेवनसेवितं}}\Afootnote{K; \skt{सीव्येत्ततः परम्}\ A; \skt{सेवितं}\ H}}\skt{~। }\pend
\pstart\versepar \skt{}\edtext{\skt{कार्यञ्च गोः}}{\lemma{\skt{कार्यञ्च गोः}}\Afootnote{K; \skt{कार्यो गो}\ A; \skt{कार्यश्च गोः}\ H}}\skt{ }\edtext{\skt{फलावन्धः}}{\lemma{\skt{फलावन्धः}}\Afootnote{K; \skt{फणिकाबन्धः}\ A; \skt{फणो बन्धः}\ H}}\skt{ }\edtext{\skt{कट्यामावेश्य}}{\lemma{\skt{कट्यामावेश्य}}\Afootnote{\skt{कट्यामावेश्य}\ A K H}}\skt{ यन्त्रकं~।}\pend

\pstart\versepar \ledsidenote{\passref{4.2.68}}\skt{न कुर्या\-त्स्ने\-ह\-से\-कञ्च तेन }\edtext{\skt{प्रक्लिद्यते}}{\lemma{\skt{प्रक्लिद्यते}}\Afootnote{K H; \skt{क्लिद्यति हि}\ A}}\skt{ व्रणः~। }\pend
\pstart\versepar \skt{}\edtext{\skt{काला\-नु\-सा\-री\-म\-गुरूं}}{\lemma{\skt{काला\appbrk{}नु\appbrk{}सा\appbrk{}री\appbrk{}म\appbrk{}गुरूं}}\Afootnote{\skt{काला\appbrk{}नु\appbrk{}सा\appbrk{}र्या\appbrk{}गुर्}\ A; \skt{कालानुसा}\ K; \skt{काला\appbrk{}नु\appbrk{}सा\appbrk{}री\appbrk{}म\appbrk{}गुरुं}\ H}}\skt{ }\edtext{\skt{चन्द\-न\-न्दे\-व\-दारू च}}{\lemma{\skt{चन्द\appbrk{}न\appbrk{}न्दे\appbrk{}व\appbrk{}दारू च}}\Afootnote{H; \skt{वेला\appbrk{}जा\appbrk{}ती\appbrk{}च\appbrk{}न्द\appbrk{}न\appbrk{}प\appbrk{}द्मकैः}\ A; \skt{रीमगुरूं चन्द\appbrk{}न\appbrk{}न्दे\appbrk{}व\appbrk{}दारू च}\ K}}\skt{~।}\pend

\pstart\versepar \ledsidenote{\passref{4.2.69}}\edtext{\skt{मनः}}{\lemma{\skt{मनः}}\Afootnote{K H; \emph{om.}\ A}}\skt{ }\edtext{\skt{शिलाले}}{\lemma{\skt{शिलाले}}\Afootnote{K H; \skt{शिलादा}\ A}}\skt{ }\edtext{\skt{वा हृत्य}}{\lemma{\skt{वा हृत्य}}\Afootnote{K; \skt{र्व्यमृतातुत्थैस्}\ A; \skt{चाहृत्य}\ H}}\skt{ तैलं कुर्वीत रोपणं~। }\pend
\pstart\versepar \skt{}\edtext{\skt{शीर्षाद्यपहते}}{\lemma{\skt{शीर्षाद्यपहते}}\Afootnote{\skt{शिर\appbrk{}सोऽ\appbrk{}अ\appbrk{}प\appbrk{}हृते}\ A; \skt{शीर्षाद्यपहते}\ K; \skt{शीर्षाद्व्यपहृते}\ H}}\skt{ शल्ये }\edtext{\skt{वाल\-व\-र्ति\-म्प्र\-वे\-श\-येत्}}{\lemma{\skt{प्रवेशयेत्}}\Afootnote{K H; \skt{निवेशयेत्}\ A}}\skt{~।}\pend

\pstart\versepar \ledsidenote{\passref{4.2.70}}\edtext{\skt{वाल\-व\-र्त्या\-म\-द\-त्तायां}}{\lemma{\skt{वाल\appbrk{}व\appbrk{}र्त्या\appbrk{}म\appbrk{}द\appbrk{}त्तायां}}\Afootnote{K; \skt{बाल\appbrk{}व\appbrk{}र्त्या\appbrk{}म\appbrk{}द\appbrk{}त्तायां}\ A; \skt{बालवर्त्ता मदन्ताया}\ H\textsuperscript{ac}; \skt{बालव र्त्त्यामदन्ताया}\ H\textsuperscript{pc}}}\skt{ }\edtext{\skt{मस्तु\-लुं\-ग\-व्र\-णा\-त्स्र\-वेत्}}{\lemma{\skt{मस्तु\appbrk{}लुं\appbrk{}ग\appbrk{}व्र\appbrk{}णात्}}\Afootnote{K\textsuperscript{pc}; \skt{मस्तुलुङ्गं व्रणात्}\ A; \skt{मस्तुलुं व्रणात्}\ K\textsuperscript{ac}; \skt{मस्त\appbrk{}लुं\appbrk{}ग\appbrk{}व्र\appbrk{}णान्}\ H}}\skt{~। }\pend
\pstart\versepar \skt{हन्यादेनन्ततो वायु\-स्त\-स्मा\-दे\-व\-मु\-पा\-च\-रेत्~।}\pend

\pstart\versepar \ledsidenote{\passref{4.2.71}}\skt{व्रणे }\edtext{\skt{रूहति}}{\lemma{\skt{रूहति}}\Afootnote{K; \skt{रोहति}\ A; \skt{रुह्यति}\ H}}\skt{ चैकैकं }\edtext{\edtext{\skt{शनै\-र्वा\-ल\-मु\-प\-क्षि\-पेत्}}{\lemma{\skt{वालम्}}\Afootnote{K H; \skt{बाल}\ A}}}{\lemma{\skt{उपक्षिपेत्}}\Afootnote{K H; \skt{मपक्षिपेत्}\ A}}\skt{~। }\pend
\pstart\versepar \skt{}\edtext{\skt{गात्राद्यपहृते}}{\lemma{\skt{गात्राद्यपहृते}}\Afootnote{K; \skt{गात्रादपहृते}\ A; \skt{गात्राद्व्यपहृते}\ H}}\skt{ }\edtext{\skt{तस्मिं}}{\lemma{\skt{तस्मिं}}\Afootnote{K; \skt{ऽन्यस्मात्}\ A; \skt{तस्मिन्}\ H}}\skt{ स्नेह\-व\-र्ति\-म्प्र\-वे\-श\-येत्~।}\pend

\pstart\versepar \ledsidenote{\passref{4.2.72}}\edtext{\skt{निः}}{\lemma{\skt{निः}}\Afootnote{K H; \skt{कृते निः}\ A}}\skt{ शोणिते चापि }\edtext{\skt{कृतविधिः}}{\lemma{\skt{कृतविधिः}}\Afootnote{K; \skt{विधिः}\ A; \skt{कृते विधिः}\ H}}\skt{ }\edtext{\skt{सद्यक्षते}}{\lemma{\skt{सद्यक्षते}}\Afootnote{K; \skt{सद्यःक्षते}\ A H}}\skt{ हितः~। }\pend
\pstart\versepar \skt{}\edtext{\skt{दूरावगाढाः}}{\lemma{\skt{दूरावगाढाः}}\Afootnote{A K; \skt{दूरावगाढः}\ H}}\skt{ }\edtext{\skt{सूक्ष्मा}}{\lemma{\skt{सूक्ष्मा}}\Afootnote{K; \skt{सूक्ष्माः}\ A H}}\skt{ }\edtext{\edtext{\skt{स्युर्ये}}{\lemma{\skt{स्युर्}}\Afootnote{K H; \skt{स्यु}\ A}}}{\lemma{\skt{ये}}\Afootnote{K H; \skt{र्ये}\ A}}\skt{ }\edtext{\edtext{\skt{व्रणांस्तां}}{\lemma{\skt{व्रणांस्}}\Afootnote{K H; \skt{व्रणा}\ A}}}{\lemma{\skt{तां}}\Afootnote{K H; \skt{स्तान्}\ A}}\skt{ }\edtext{\skt{विशोणितां}}{\lemma{\skt{विशोणितां}}\Afootnote{K H; \skt{विशोणितान्}\ A}}\skt{~।}\pend

\pstart\versepar \ledsidenote{\passref{4.2.73}}\skt{कृत्वा }\edtext{\skt{नेत्रेण}}{\lemma{\skt{नेत्रेण}}\Afootnote{K H; \emph{om.}\ A}}\skt{ }\edtext{\skt{सूक्ष्मेण}}{\lemma{\skt{सूक्ष्मेण}}\Afootnote{K H; \skt{सूक्ष्मेण नेत्रेण}\ A}}\skt{ चक्रतैलेन तर्पयेत्~। }\pend
\pstart\versepar \skt{}\edtext{\skt{पद्मां}}{\lemma{\skt{पद्मां}}\Afootnote{K H; \emph{om.}\ A}}\skt{ समङ्गां रजनीं }\edtext{\skt{त्रिवर्गं}}{\lemma{\skt{त्रिवर्गं}}\Afootnote{K H; \skt{पद्मां त्रिवर्गं}\ A}}\skt{ }\edtext{\edtext{\skt{तुत्थमेव}}{\lemma{\skt{तुत्थम्}}\Afootnote{K H; \skt{तुत्थ}\ A}}}{\lemma{\skt{एव}}\Afootnote{K H; \skt{मेव}\ A}}\skt{ च~।}\pend

\pstart \ledsidenote{\passref{4.2.74}}\edtext{\skt{संहृत्य}}{\lemma{\skt{संहृत्य}}\Afootnote{K H; \skt{विडङ्गं कटुकां पथ्यां गुडूचीं सकरञ्जिकाम्~। संहृत्य}\ A}}\skt{ विपचेत्काले तैलं }\edtext{\edtext{\skt{रोपणमुत्तमं}}{\lemma{\skt{रोपणम्}}\Afootnote{K H; \skt{रोपण}\ A}}}{\lemma{\skt{उत्तमं}}\Afootnote{K H; \skt{मुत्तमम्}\ A}}\skt{~।}\pend

\pstart \ledsidenote{\passref{4.2.76}}\edtext{\skt{क्षते}}{\lemma{\skt{क्षते}}\Afootnote{K H; \skt{पक्वं सद्योव्रणेषूक्तं तैलं रोप\appbrk{}ण\appbrk{}मु\appbrk{}त्त\appbrk{}मम्~। क्षते}\ A}}\skt{ क्षतविधिः कार्यः पिच्चिते }\edtext{\skt{भग्नवद्विधिः}}{\lemma{\skt{भग्नवद्}}\Afootnote{K H; \skt{भग्नवद्}\ A}}\skt{~।}\pend

\pstart\versepar \ledsidenote{\passref{4.2.77}}\skt{घृष्टे }\edtext{\skt{रुजां}}{\lemma{\skt{रुजां}}\Afootnote{K H; \skt{रुजो}\ A}}\skt{ }\edtext{\skt{निगृ\-ह्या\-शु\-स्चू\-र्णै\-रु\-प\-च\-रे\-द्व्रणं}}{\lemma{\skt{निगृह्याशुस्}}\Afootnote{K H; \skt{निगृह्याशु}\ A}}\skt{~। }\pend
\pstart\versepar \skt{विश्लिष्टदेहं पतितं मथितं हतमेव }\edtext{\skt{वा}}{\lemma{\skt{वा}}\Afootnote{K H; \skt{च}\ A}}\skt{~।}\pend

\pstart\versepar \ledsidenote{\passref{4.2.78}}\edtext{\edtext{\skt{वास\-ये\-त्तै\-ल\-पू\-र्ण्णायां}}{\lemma{\skt{वासयेत्}}\Afootnote{A H; \skt{वास}\ K}}}{\lemma{\skt{तैलपूर्ण्णायां}}\Afootnote{A H; \emph{om.}\ K}}\skt{ }\edtext{\skt{द्रोण्यां}}{\lemma{\skt{द्रोण्यां}}\Afootnote{A H; \emph{om.}\ K}}\skt{ }\edtext{\skt{मांसरसाशनं}}{\lemma{\skt{मांसरसाशनं}}\Afootnote{A H; \emph{om.}\ K}}\skt{~॥}\footnoteA{MS K breaks off at this point and resumes at 4.19.22.}\skt{ }\pend
\pstart\versepar \skt{}\edtext{\skt{एष}}{\lemma{\skt{एष}}\Afootnote{H; \skt{अयम्}\ A; \emph{om.}\ K}}\skt{ }\edtext{\skt{एव}}{\lemma{\skt{एव}}\Afootnote{A H; \emph{om.}\ K}}\skt{ }\edtext{\skt{विधिः}}{\lemma{\skt{विधिः}}\Afootnote{A H; \emph{om.}\ K}}\skt{ }\edtext{\skt{कार्यः}}{\lemma{\skt{कार्यः}}\Afootnote{A H; \emph{om.}\ K}}\skt{ }\edtext{\skt{क्षीणे}}{\lemma{\skt{क्षीणे}}\Afootnote{A H; \emph{om.}\ K}}\skt{ }\edtext{\skt{मर्म्महते}}{\lemma{\skt{मर्म्महते}}\Afootnote{A H; \emph{om.}\ K}}\skt{ }\edtext{\skt{तथा}}{\lemma{\skt{तथा}}\Afootnote{A H; \emph{om.}\ K}}\skt{~।}\pend

\pstart \ledsidenote{\passref{4.2.79}}\skt{रोपणे }\edtext{\skt{चैव}}{\lemma{\skt{चैव}}\Afootnote{H; \skt{सप}\ A}}\skt{ }\edtext{\skt{सेके च}}{\lemma{\skt{सेके च}}\Afootnote{H; \skt{रीषेके}\ A}}\skt{ पाने च व्रणिनां सदा~॥ तैलं घृतं }\edtext{\skt{चावचार्यं}}{\lemma{\skt{चावचार्यं}}\Afootnote{H; \skt{वा संयोज्यं}\ A}}\skt{ शरी\-र\-र्तू\-न\-वेक्ष्य हि~।}\pend

\pstart \ledsidenote{\passref{4.2.80}}\skt{घृतानि यानि वक्ष्यामि }\edtext{\skt{यन्नतः}}{\lemma{\skt{यन्नतः}}\Afootnote{H; \skt{यत्नतः}\ A}}\skt{ पित्तविद्रधौ~॥ }\edtext{\skt{सद्यो}}{\lemma{\skt{सद्यो}}\Afootnote{H; \skt{सद्यो}\ A}}\skt{ व्रणेषु }\edtext{\skt{दातव्यन्तानि}}{\lemma{\skt{दातव्यन्}}\Afootnote{H; \skt{देयानि}\ A}}\skt{ वैद्येन जानता~॥}\pend

\pstart \ledsidenote{\passref{4.2.81}}\edtext{\skt{सद्यः\-क्ष\-त\-म्व्र\-ण\-म्वैद्यः}}{\lemma{\skt{सद्यःक्षतम्}}\Afootnote{H; \skt{सद्यःक्षत}\ A}}\skt{ सशू\-ल\-म्प\-रि\-षे\-च\-येत्~। सर्प्पिषा नातिशीतेन बलातैलेन वा पुनः~॥}\pend

\pstart \ledsidenote{\passref{4.2.82}}\edtext{\skt{पद्मां}}{\lemma{\skt{पद्मां}}\Afootnote{H; \emph{om.}\ A}}\skt{ समंगां }\edtext{\skt{रज\-नी\-म्प\-थ्या\-न्तुत्थं}}{\lemma{\skt{पथ्यान्}}\Afootnote{H; \skt{पद्मां पथ्यां}\ A}}\skt{ सुवर्च्चलां~। }\edtext{\skt{पद्मकम्मधुकं}}{\lemma{\skt{मधुकं}}\Afootnote{H; \skt{रोध्रमधुकं}\ A}}\skt{ }\edtext{\skt{लोध्रं}}{\lemma{\skt{लोध्रं}}\Afootnote{H; \emph{om.}\ A}}\skt{ विडंगानि }\edtext{\skt{हरेणवः}}{\lemma{\skt{हरेणवः}}\Afootnote{H; \skt{हरेणुकाम्}\ A}}\skt{~॥}\pend

\pstart \ledsidenote{\passref{4.2.83}}\edtext{\skt{तालीसपत्त्रं}}{\lemma{\skt{तालीसपत्त्रं}}\Afootnote{H; \skt{तालीशपत्रं}\ A}}\skt{ नलदं }\edtext{\skt{चन्द\-न\-म्प\-द्म\-के\-सरं}}{\lemma{\skt{पद्मकेसरं}}\Afootnote{H; \skt{पद्मकेशरम्}\ A}}\skt{~। मञ्जि\-ष्ठो\-शी\-र\-ला\-क्षाश्च क्षीरिणाञ्चापि पल्लवान्~॥}\pend

\pstart\versepar \ledsidenote{\passref{4.2.84}}\edtext{\skt{पियालबीजं}}{\lemma{\skt{पियालबीजं}}\Afootnote{H; \skt{प्रियालबीजं}\ A}}\skt{ }\edtext{\edtext{\skt{तिन्दु\-क्या\-स्त\-रु\-णानि}}{\lemma{\skt{तिन्दुक्यास्}}\Afootnote{H; \skt{तिन्दुक्या}\ A}}}{\lemma{\skt{तरुणानि}}\Afootnote{H; \skt{स्तरुणानि}\ A}}\skt{ फलानि च~॥ }\pend
\pstart\versepar \skt{}\edtext{\skt{यथालाभेन}}{\lemma{\skt{यथालाभेन}}\Afootnote{H; \skt{यथालाभं स}\ A}}\skt{ }\edtext{\skt{संहृत्य}}{\lemma{\skt{संहृत्य}}\Afootnote{H; \skt{माहृत्य}\ A}}\skt{ }\edtext{\edtext{\skt{तैलन्धीरो}}{\lemma{\skt{तैलन्}}\Afootnote{H; \skt{तैल}\ A}}}{\lemma{\skt{धीरो}}\Afootnote{H; \skt{मेभिर्}\ A}}\skt{ विपाचयेत्~।}\pend

\pstart\versepar \ledsidenote{\passref{4.2.85}}\skt{सद्योव्रणानां }\edtext{\edtext{\skt{सर्व्वे\-षा\-म\-दु\-ष्टा\-नाञ्च}}{\lemma{\skt{सर्व्वेषाम्}}\Afootnote{H; \skt{सर्वेषा}\ A}}}{\lemma{\skt{अदुष्टानाञ्च}}\Afootnote{H; \skt{मदुष्टानां तु}\ A}}\skt{ रोपणम्~॥ }\pend
\pstart\versepar \skt{कषायमधुराः शीताः क्रियाः स्निग्धाश्च योजयेत्~।}\pend

\pstart\versepar \ledsidenote{\passref{4.2.86}}\skt{सद्योव्रणानां }\edtext{\skt{सप्ताहं }\edtext{\edtext{\edtext{\skt{पश्चा\-त्पू\-र्व्वो\-क्त\-मा\-च\-रेत्}}{\lemma{\skt{चात्}}\Afootnote{H; \skt{श्चात्}\ A}}}{\lemma{\skt{पूर्व्वोक्तम्}}\Afootnote{H; \skt{पूर्वोक्त}\ A}}}{\lemma{\skt{आचरेत्}}\Afootnote{H; \skt{माचरेत्}\ A}}}{\lemma{\skt{सप्ताहं पश्}}\Afootnote{H; \skt{सप्ताहं प}\ A}}\skt{~। }\pend
\pstart\versepar \skt{दुष्टव्रणेषु }\edtext{\edtext{\skt{कर्त्त\-व्य\-मू\-र्द्ध\-ञ्चा\-धश्च}}{\lemma{\skt{कर्त्तव्यम्}}\Afootnote{H; \skt{कर्तव्य}\ A}}}{\lemma{\skt{ऊर्द्धञ्}}\Afootnote{H; \skt{मूर्ध्वं}\ A}}\skt{ शोधनम्~॥}\pend

\pstart\versepar \ledsidenote{\passref{4.2.87}}\edtext{\skt{विशो\-ष\-ण\-न्त\-था\-हारः}}{\lemma{\skt{तथाहारः}}\Afootnote{H; \skt{तथा++आहारः}\ A}}\skt{ शोणितस्य च मोक्षणं~। }\pend
\pstart\versepar \skt{}\edtext{\skt{कषा\-य\-मा\-र\-ग्व\-धा\-दि\-सु\-र\-सा\-द्यञ्च}}{\lemma{\skt{कषा\appbrk{}य\appbrk{}मा\appbrk{}र\appbrk{}ग्व\appbrk{}धा\appbrk{}दि\appbrk{}सु\appbrk{}र\appbrk{}सा\appbrk{}द्यञ्च}}\Afootnote{H; \skt{कषायं राजवृक्षादौ सुरसादौ च}\ A}}\skt{ }\edtext{\skt{शोधनं}}{\lemma{\skt{शोधनं}}\Afootnote{H; \skt{धावनम्}\ A}}\skt{~॥}\pend

\pstart\versepar \ledsidenote{\passref{4.2.88}}\edtext{\edtext{\skt{कषा\-य\-यो\-रे\-त\-योस्तु}}{\lemma{\skt{कषाययोर्}}\Afootnote{H; \skt{तयोरेव कषाये}\ A}}}{\lemma{\skt{एतयोस्तु}}\Afootnote{H; \skt{ण}\ A}}\skt{ तैलं }\edtext{\edtext{\skt{शोधनमिष्यते}}{\lemma{\skt{शोधनम्}}\Afootnote{H; \skt{शोधन}\ A}}}{\lemma{\skt{इष्यते}}\Afootnote{H; \skt{मिष्यते}\ A}}\skt{~। }\pend
\pstart\versepar \skt{क्षारकल्पेन वा तैलं क्षारद्रव्येषु साधितं~॥}\pend

\pstart\versepar \ledsidenote{\passref{4.2.89}}\skt{द्रवन्ती }\edtext{\skt{चिरिबिल्वञ्च}}{\lemma{\skt{चिरिबिल्वञ्च}}\Afootnote{H; \skt{चिरबिल्वश्च}\ A}}\skt{ दन्ती }\edtext{\edtext{\skt{चित्त्रकमेव}}{\lemma{\skt{चित्त्रकम्}}\Afootnote{H; \skt{चित्रक}\ A}}\skt{ च}}{\lemma{\skt{एव च}}\Afootnote{H; \skt{मेव च}\ A}}\skt{~। }\pend
\pstart\versepar \skt{पृथ्वीका निम्बपत्त्राणि }\edtext{\skt{काशीसं}}{\lemma{\skt{काशीसं}}\Afootnote{H; \skt{कासीसं}\ A}}\skt{ }\edtext{\edtext{\skt{तुत्थमेव}}{\lemma{\skt{तुत्थम्}}\Afootnote{H; \skt{तुत्थ}\ A}}}{\lemma{\skt{एव}}\Afootnote{H; \skt{मेव}\ A}}\skt{ च~॥}\pend

\pstart\versepar \ledsidenote{\passref{4.2.90}}\edtext{\skt{त्रिवृता}}{\lemma{\skt{त्रिवृता}}\Afootnote{H; \skt{त्रिवृत्ते}\ A}}\skt{ }\edtext{\skt{नीलिनी}}{\lemma{\skt{नीलिनी}}\Afootnote{H; \skt{जोवती नीली}\ A}}\skt{ }\edtext{\skt{श्यामा}}{\lemma{\skt{श्यामा}}\Afootnote{H; \emph{om.}\ A}}\skt{ }\edtext{\skt{हरिद्रा}}{\lemma{\skt{हरिद्रा}}\Afootnote{H; \skt{हरिद्रे}\ A}}\skt{ सैन्धवन्तिला~। }\pend
\pstart\versepar \skt{}\edtext{\skt{नीपः}}{\lemma{\skt{नीपः}}\Afootnote{H; \skt{भूमि}\ A}}\skt{ }\edtext{\skt{कदम्बः}}{\lemma{\skt{कदम्बः}}\Afootnote{A H\textsuperscript{pc}; \skt{कदम्ब}\ H\textsuperscript{ac}}}\skt{ }\edtext{\skt{सुवहां}}{\lemma{\skt{सुवहां}}\Afootnote{H; \skt{सुवहा}\ A}}\skt{ }\edtext{\skt{शुकाख्यं}}{\lemma{\skt{शुकाख्यं}}\Afootnote{H; \skt{शुकाख्या}\ A}}\skt{ }\edtext{\skt{लाङ्गलाह्वयं}}{\lemma{\skt{लाङ्गलाह्वयं}}\Afootnote{H; \skt{लाङ्गलाह्वया}\ A}}\skt{~॥}\pend

\pstart\versepar \ledsidenote{\passref{4.2.91}}\edtext{\skt{मृगादनी}}{\lemma{\skt{मृगादनी}}\Afootnote{H; \emph{om.}\ A}}\skt{ }\edtext{\skt{मदयन्ती}}{\lemma{\skt{मदयन्ती}}\Afootnote{H; \emph{om.}\ A}}\skt{ }\edtext{\skt{नेपाली}}{\lemma{\skt{नेपाली}}\Afootnote{H; \skt{नैपाली}\ A}}\skt{ }\edtext{\skt{जालिनी}}{\lemma{\skt{जालिनी}}\Afootnote{H; \skt{जालिनी चैव मदयन्ती मृगादनी}\ A}}\skt{ }\edtext{\skt{सुरा}}{\lemma{\skt{सुरा}}\Afootnote{H; \skt{सुधा}\ A}}\skt{~। }\pend
\pstart\versepar \skt{}\edtext{\skt{स्नुहा\-चा\-र्क्क\-वि\-ड\-ङ्गा\-नि\-ह\-रि\-ता\-ल\-क\-र\-ञ्जिका}}{\lemma{\skt{स्नुहा\appbrk{}चा\appbrk{}र्क्क\appbrk{}वि\appbrk{}ड\appbrk{}ङ्गा\appbrk{}नि\appbrk{}ह\appbrk{}रि\appbrk{}ता\appbrk{}ल\appbrk{}क\appbrk{}र\appbrk{}ञ्जिका}}\Afootnote{H; \skt{मूर्वा\appbrk{}र्क\appbrk{}की\appbrk{}टा\appbrk{}रि\appbrk{}ह\appbrk{}रि\appbrk{}ता\appbrk{}ल\appbrk{}क\appbrk{}र\appbrk{}ञ्जिकाः}\ A}}\skt{~॥}\pend

\pstart\versepar \ledsidenote{\passref{4.2.92}}\edtext{\skt{यथोपपत्या}}{\lemma{\skt{यथोपपत्या}}\Afootnote{H; \skt{यथोपपत्ति}\ A}}\skt{ }\edtext{\skt{कर्त्त\-व्य\-न्तै\-ल\-मे\-तेषु}}{\lemma{\skt{एतेषु}}\Afootnote{H; \skt{एतैस्तु}\ A}}\skt{ शोधनम्~। }\pend
\pstart\versepar \skt{घृतम्वा }\edtext{\skt{यदि}}{\lemma{\skt{यदि}}\Afootnote{H; \skt{यदि वा}\ A}}\skt{ प्राप्तं कल्काः }\edtext{\edtext{\skt{संशोधनास्तथा}}{\lemma{\skt{संशोधनास्}}\Afootnote{H; \skt{संशोधना}\ A}}}{\lemma{\skt{तथा}}\Afootnote{H; \skt{स्तथा}\ A}}\skt{~॥}\pend

\pstart \ledsidenote{\passref{4.2.94}}\edtext{\skt{दुष्टव्रणविधिः}}{\lemma{\skt{दुष्टव्रणविधिः}}\Afootnote{H; \skt{कफजे तिल\appbrk{}ते\appbrk{}जो\appbrk{}ह्वा\appbrk{}द\appbrk{}न्ती\appbrk{}स्व\appbrk{}र्जि\appbrk{}क\appbrk{}चि\appbrk{}त्रकाः~। दुष्टव्रणविधिः}\ A}}\skt{ कार्यः }\edtext{\edtext{\skt{कुष्ठ\-मे\-ह\-व्र\-णे\-ष्वपि}}{\lemma{\skt{कुष्ठ\appbrk{}मे\appbrk{}ह\appbrk{}व्र\appbrk{}णेष्व्}}\Afootnote{H; \skt{मेह\appbrk{}कु\appbrk{}ष्ठ\appbrk{}व्र\appbrk{}णेष्}\ A}}}{\lemma{\skt{अपि}}\Afootnote{H; \skt{वपि}\ A}}\skt{~।}\pend

\pstart\versepar \ledsidenote{\passref{4.2.95}}\skt{षड्विधः प्राक्प्रदिष्टो यः सद्यो\-व्र\-ण\-वि\-नि\-श्चयः~॥ }\pend
\pstart\versepar \skt{नातः शक्यम्परं }\edtext{\edtext{\skt{वक्तुमपि}}{\lemma{\skt{वक्तुम्}}\Afootnote{H; \skt{वक्तु}\ A}}}{\lemma{\skt{अपि}}\Afootnote{H; \skt{मपि}\ A}}\skt{ }\edtext{\skt{पण्डितमानिना}}{\lemma{\skt{पण्डितमानिना}}\Afootnote{H; \skt{निश्चितवादिभिः}\ A}}\skt{~।}\pend

\pstart\versepar \ledsidenote{\passref{4.2.96}}\edtext{\skt{उप\-स\-र्ग्गै\-र्न्नि\-पा\-तैस्तु तन्तु}}{\lemma{\skt{न्निपातैस्तु तन्तु}}\Afootnote{H; \skt{निपातैश्च तत्तु}\ A}}\skt{ }\edtext{\skt{पण्डितवादिनः}}{\lemma{\skt{पण्डितवादिनः}}\Afootnote{H; \skt{पण्डितमानिनः}\ A}}\skt{~॥ }\pend
\pstart\versepar \skt{}\edtext{\skt{संयोज्य}}{\lemma{\skt{संयोज्य}}\Afootnote{H; \skt{केचित्संयोज्य}\ A}}\skt{ }\edtext{\skt{केचिद्भाषन्ते}}{\lemma{\skt{केचिद्}}\Afootnote{H; \emph{om.}\ A}}\skt{ बहुधा मानगर्व्विताः~।}\pend

\pstart\versepar \ledsidenote{\passref{4.2.97}}\skt{बहु तद्भाषितन्तेषां }\edtext{\edtext{\skt{षट्स्वेतेषु}}{\lemma{\skt{षट्स्व्}}\Afootnote{H; \skt{षट्स्}\ A}}}{\lemma{\skt{एतेषु}}\Afootnote{H; \skt{वेष्वेवा}\ A}}\skt{ }\edtext{\skt{निरुध्यते}}{\lemma{\skt{निरुध्यते}}\Afootnote{H; \skt{वतिष्ठते}\ A}}\skt{~॥ }\pend
\pstart\versepar \skt{विशेषा इव सामान्ये }\edtext{\edtext{\edtext{\edtext{\skt{तस्मा\-त्ष\-ट्त्व\-म्प\-र\-म्म\-त\-मिति}}{\lemma{\skt{तस्मात्}}\Afootnote{H; \emph{om.}\ A}}}{\lemma{\skt{षट्त्वम्}}\Afootnote{H; \skt{षट्त्वं तु}\ A}}}{\lemma{\skt{परम्}}\Afootnote{H; \skt{परमं}\ A}}}{\lemma{\skt{इति}}\Afootnote{H; \emph{om.}\ A}}\skt{~॥}\pend

\pstart \ledsidenote{\passref{4.2.trailer}}\edtext{\skt{चिकित्सा}}{\lemma{\skt{चिकित्सा}}\Afootnote{H; \skt{इति सुश्रु\appbrk{}त\appbrk{}सं\appbrk{}हि\appbrk{}तायां चिकित्सास्थाने सद्यो\appbrk{}च्र\appbrk{}ण\appbrk{}चि\appbrk{}कि\appbrk{}त्सितं}\ A}}\skt{ }\edtext{\skt{द्वि}}{\lemma{\skt{द्वि}}\Afootnote{H; \skt{नाम द्वितीयोऽअध्यायः}\ A}}\skt{~॥}\pend

\endnumbering
\beginnumbering
\eledsection*{Adhyāya 3}
\pstart \ledsidenote{\passref{4.3.1}}\skt{अथातो }\edtext{\skt{भग्न\-चि\-कि\-त्सि\-त\-म्व्या\-ख्या\-स्यामः}}{\lemma{\skt{भग्न\appbrk{}चि\appbrk{}कि\appbrk{}त्सि\appbrk{}तम्}}\Afootnote{H; \skt{भग्नानां चिकित्सितं}\ A}}\skt{~॥}\pend

\pstart\versepar \ledsidenote{\passref{4.3.3}}\skt{अल्पा\-शि\-नोऽ\-अ\-ना\-त्म\-वतो }\edtext{\skt{जन्तो\-र्वा\-ता\-त्म\-कस्य वा}}{\lemma{\skt{जन्तो\appbrk{}र्वा\appbrk{}ता\appbrk{}त्म\appbrk{}कस्य वा}}\Afootnote{H; \skt{जन्तो\appbrk{}र्वा\appbrk{}ता\appbrk{}त्म\appbrk{}कस्य च}\ A}}\skt{~। }\pend
\pstart\versepar \skt{उपद्रवैर्वा जुष्टस्य }\edtext{\skt{भग्नः}}{\lemma{\skt{भग्नः}}\Afootnote{H; \skt{भग्नं}\ A}}\skt{ कृच्छ्रेण सिध्यति~॥}\pend

\pstart\versepar \ledsidenote{\passref{4.3.4}}\skt{लवणङ्कटुकं क्षारमम्लं मैथुनमातपं~। }\pend
\pstart\versepar \skt{व्यायामञ्च न सेवेत भग्नो }\edtext{\edtext{\skt{रूक्षान्नमेव}}{\lemma{\skt{रूक्षान्नम्}}\Afootnote{H; \skt{रूक्षान्न}\ A}}}{\lemma{\skt{एव}}\Afootnote{H; \skt{मेव}\ A}}\skt{ च~॥}\pend

\pstart\versepar \ledsidenote{\passref{4.3.5}}\edtext{\skt{मान्सम्मांसरसः}}{\lemma{\skt{मान्सम्}}\Afootnote{H; \skt{शालिर्}\ A}}\skt{ क्षीरं सर्प्पिर्यूषः सतीनजः~। }\pend
\pstart\versepar \skt{बृंह\-ण\-ञ्चा\-न्न\-पानं स्याद्देयं भग्नाय जानता~॥}\pend

\pstart\versepar \ledsidenote{\passref{4.3.6}}\edtext{\skt{पला\-शो\-दु\-म्ब\-रा\-श्व\-त्थ\-क\-द\-म्ब\-नि\-चु\-ल\-त्वचः}}{\lemma{\skt{पला\appbrk{}शो\appbrk{}दु\appbrk{}म्ब\appbrk{}रा\appbrk{}श्व\appbrk{}त्थ\appbrk{}क\appbrk{}द\appbrk{}म्ब\appbrk{}नि\appbrk{}चु\appbrk{}ल\appbrk{}त्वचः}}\Afootnote{H; \skt{मधू\appbrk{}को\appbrk{}दु\appbrk{}म्ब\appbrk{}रा\appbrk{}श्व\appbrk{}त्थ\appbrk{}प\appbrk{}ला\appbrk{}श\appbrk{}क\appbrk{}कु\appbrk{}भ\appbrk{}त्वचः}\ A}}\skt{~। }\pend
\pstart\versepar \skt{}\edtext{\skt{वंश\-स\-र्ज्जा\-र्जु\-ना\-दीनां}}{\lemma{\skt{वंश\appbrk{}स\appbrk{}र्ज्जा\appbrk{}र्जु\appbrk{}ना\appbrk{}दीनां}}\Afootnote{H; \skt{वंश\appbrk{}स\appbrk{}र्ज\appbrk{}व\appbrk{}टानां}\ A}}\skt{ }\edtext{\edtext{\skt{कुप्य\-र्थ\-मु\-प\-सं\-ह\-रेत्}}{\lemma{\skt{कुप्यर्थम्}}\Afootnote{H; \skt{च कुशार्थ}\ A}}}{\lemma{\skt{उपसंहरेत्}}\Afootnote{H; \skt{मुपसंहरेत्}\ A}}\skt{~॥}\pend

\pstart\versepar \ledsidenote{\passref{4.3.7}}\edtext{\skt{आलेपनार्थे}}{\lemma{\skt{आलेपनार्थे}}\Afootnote{H; \skt{आलेपनार्थं}\ A}}\skt{ }\edtext{\skt{मञ्जिष्ठा}}{\lemma{\skt{मञ्जिष्ठा}}\Afootnote{H; \skt{मञ्जिष्ठां}\ A}}\skt{ }\edtext{\skt{मधु\-क\-ञ्चा\-म्ल\-पे\-षिमं}}{\lemma{\skt{चाम्लपेषिमं}}\Afootnote{H; \skt{रक्तचन्दनम्}\ A}}\skt{~। }\pend
\pstart\versepar \skt{शत\-धौ\-त\-घृ\-तो\-न्मिश्रं शालिपिष्टञ्च }\edtext{\skt{लेपनं}}{\lemma{\skt{लेपनं}}\Afootnote{H; \skt{संहरेत्}\ A}}\skt{~॥}\pend

\pstart \ledsidenote{\passref{4.3.8ab}}\edtext{\edtext{\skt{सप्त\-रा\-त्रा\-त्स\-प्त\-रा\-त्रा\-त्सौ\-म्ये\-ष्वृ\-तुषु}}{\lemma{\skt{सप्तरात्रात्}}\Afootnote{H; \skt{सप्ताहादथ}\ A}}}{\lemma{\skt{सप्तरात्रात्}}\Afootnote{H; \skt{सप्ताहात्}\ A}}\skt{ बन्धनं~।}\pend

\pstart \ledsidenote{\passref{4.3.8.add1}}\skt{कर्त्तव्यञ्च त्रिरात्रं स्यात्तथा }\edtext{\skt{ग्रीष्मेषु}}{\lemma{\skt{ग्रीष्मेषु}}\Afootnote{H\textsuperscript{pc}; \skt{ग्नेये षु}\ H\textsuperscript{ac}}}\skt{ जानता~॥}\pend

\pstart \ledsidenote{\passref{4.3.8.add2}}\skt{काले तु समशीतोष्णे पञ्च\-रा\-त्रा\-न्नि\-ब\-न्ध\-येत्~।}\pend

\pstart \ledsidenote{\passref{4.3.9}}\edtext{\skt{तत्रातिशिथिलं}}{\lemma{\skt{तत्रातिशिथिलं}}\Afootnote{H; \skt{आग्नेयेषु त्र्यहा तु कुर्या\appbrk{}द्भ\appbrk{}ग्न\appbrk{}दो\appbrk{}ष\appbrk{}व\appbrk{}शेन वा~। तत्रातिशिथिलं}\ A}}\skt{ बद्धे }\edtext{\skt{सन्धिः}}{\lemma{\skt{सन्धिः}}\Afootnote{H; \skt{सन्धि}\ A}}\skt{ स्थैर्यन्न जायते~॥}\pend

\pstart\versepar \ledsidenote{\passref{4.3.10}}\skt{गाढेनापि त्वगादीनां शोफो }\edtext{\skt{रुक्पाकमेव}}{\lemma{\skt{पाकम्}}\Afootnote{H; \skt{पाक}\ A}}\skt{ च~। }\pend
\pstart\versepar \skt{तस्मा\-त्सा\-धा\-र\-ण\-म्बन्धं भग्ने शंसन्ति तद्विदः~॥}\pend

\pstart\versepar \ledsidenote{\passref{4.3.11}}\skt{न्यग्रो\-धा\-दि\-क\-षा\-यन्तु सुशीतं परिषेचने~। }\pend
\pstart\versepar \skt{}\edtext{\skt{पञ्च\-मू\-ली\-क\-षा\-यन्तु}}{\lemma{\skt{पञ्च\appbrk{}मू\appbrk{}ली\appbrk{}क\appbrk{}षा\appbrk{}यन्तु}}\Afootnote{H; \skt{पञ्च\appbrk{}मू\appbrk{}ली\appbrk{}वि\appbrk{}पक्वं तु}\ A}}\skt{ }\edtext{\edtext{\skt{कुर्यात्क्षीरं}}{\lemma{\skt{कुर्यात्}}\Afootnote{H; \skt{क्षीरं कुर्यात्}\ A}}}{\lemma{\skt{क्षीरं}}\Afootnote{H; \emph{om.}\ A}}\skt{ सवेदने~॥}\pend

\pstart\versepar \ledsidenote{\passref{4.3.12}}\skt{सुखो\-ष्ण\-म\-व\-चा\-र्यम्वा चक्र\-तै\-ल\-म्वि\-जा\-नता~। }\pend
\pstart\versepar \skt{}\edtext{\skt{गृष्टिक्षीरं}}{\lemma{\skt{गृष्टिक्षीरं}}\Afootnote{H; \emph{om.}\ A}}\skt{ }\edtext{\edtext{\skt{सस\-र्प्पि\-ष्क\-म्म\-धु\-रौ\-ष\-ध\-सा\-धितं}}{\lemma{\skt{ससर्प्पिष्कम्}}\Afootnote{H; \skt{विभज्य कालं दोषं}\ A}}}{\lemma{\skt{मधु\appbrk{}रौ\appbrk{}ष\appbrk{}ध\appbrk{}सा\appbrk{}धितं}}\Afootnote{H; \skt{च दोष\appbrk{}घ्नौ\appbrk{}ष\appbrk{}ध\appbrk{}सं\appbrk{}यु\appbrk{}तम्}\ A}}\skt{~॥}\pend

\pstart\versepar \ledsidenote{\passref{4.3.14}}\skt{शीतलं लाक्षया }\edtext{\skt{युक्तम्प्रातर्भ्भग्नः}}{\lemma{\skt{प्रातर्भ्भग्नः}}\Afootnote{H; \skt{प्रातर्भग्नः}\ A}}\skt{ पिबेन्नरः~। }\pend
\pstart\versepar \skt{सव्रणस्य तु भग्नस्य व्रणं सर्प्पिर्म्मधूत्तरैः~॥}\pend

\pstart\versepar \ledsidenote{\passref{4.3.15}}\skt{प्रतिसार्य }\edtext{\edtext{\skt{कषायैस्तु}}{\lemma{\skt{कषायैस्}}\Afootnote{H; \skt{कषायै}\ A}}}{\lemma{\skt{तु}}\Afootnote{H; \skt{स्तु}\ A}}\skt{ शेष\-म्भ\-ग्न\-व\-दा\-च\-रेत्~। }\pend
\pstart\versepar \skt{प्रथमे वयसि त्वेवं भग्नं सुक\-र\-मा\-दि\-शेत्~॥}\pend

\pstart\versepar \ledsidenote{\passref{4.3.16}}\skt{अल्पदोषस्य }\edtext{\edtext{\skt{जन्तोस्तु}}{\lemma{\skt{जन्तोस्}}\Afootnote{H; \skt{जन्तो}\ A}}}{\lemma{\skt{तु}}\Afootnote{H; \skt{स्तु}\ A}}\skt{ काले च शिशिरात्मके~। }\pend
\pstart\versepar \skt{प्रथमे वयसि त्वेवं मासात्सन्धिः स्थिरो भवेत्~॥}\pend

\pstart\versepar \ledsidenote{\passref{4.3.17}}\skt{मध्यमे }\edtext{\skt{द्विगु\-णा\-त्का\-ला\-दु\-त्तमे}}{\lemma{\skt{कालादुत्तमे}}\Afootnote{H; \skt{कालादुत्तरे}\ A}}\skt{ }\edtext{\skt{त्रिगुणात्स्मृताः}}{\lemma{\skt{स्मृताः}}\Afootnote{H; \skt{स्मृतः}\ A}}\skt{~। }\pend
\pstart\versepar \skt{}\edtext{\edtext{\skt{अव\-ना\-मि\-त\-मु\-न्न\-ह्ये\-दु\-न्न\-त\-ञ्चा\-व\-पी\-ड\-ञ्चा\-व\-येत्}}{\lemma{\skt{अव\appbrk{}ना\appbrk{}मि\appbrk{}त\appbrk{}मु\appbrk{}न्न\appbrk{}ह्ये\appbrk{}दु\appbrk{}न्न\appbrk{}तञ्}}\Afootnote{H\textsuperscript{pc}; \skt{अव\appbrk{}ना\appbrk{}मि\appbrk{}त\appbrk{}मु\appbrk{}न्न\appbrk{}ह्ये\appbrk{}दु\appbrk{}न्नतं}\ A; \skt{अव\appbrk{}ना\appbrk{}मि\appbrk{}त\appbrk{}मु\appbrk{}ह्ये\appbrk{}दु\appbrk{}न्न\appbrk{}तञ्}\ H\textsuperscript{ac}}}}{\lemma{\skt{चाव\appbrk{}पी\appbrk{}ड\appbrk{}ञ्चा\appbrk{}व\appbrk{}येत्}}\Afootnote{H\textsuperscript{ac}; \skt{चवपीडयेत्}\ A; \skt{चावपीडयेत्}\ H\textsuperscript{pc}}}\skt{~॥}\pend

\pstart\versepar \ledsidenote{\passref{4.3.18}}\edtext{\skt{आञ्छे\-द\-ति\-क्षि\-प्त\-मधो}}{\lemma{\skt{आञ्छे\appbrk{}द\appbrk{}ति\appbrk{}क्षि\appbrk{}प्त\appbrk{}मधो}}\Afootnote{H; \skt{आञ्छे\appbrk{}द\appbrk{}ति\appbrk{}क्षि\appbrk{}प्त\appbrk{}मधो}\ A}}\skt{ गतञ्चोपरि वर्त्तयेत्~। }\pend
\pstart\versepar \skt{}\edtext{\skt{आञ्छ\-नो\-त्पी\-ड\-नो\-न्नाम}}{\lemma{\skt{पीडनोन्नाम}}\Afootnote{H; \skt{पीडनैश्चैव}\ A}}\skt{ }\edtext{\skt{संक्षो\-भै\-र्ब्ब\-न्ध\-नै\-स्तथा}}{\lemma{\skt{संक्षो\appbrk{}भै\appbrk{}र्ब्ब\appbrk{}न्ध\appbrk{}नैस्}}\Afootnote{H; \skt{सङ्क्षे\appbrk{}पै\appbrk{}र्ब\appbrk{}न्ध\appbrk{}नैस्}\ A}}\skt{~॥}\edtext{}{\lemma{\skt{वर्त्तयेत्~। आञ्छनोत्}}\Afootnote{H; \skt{वर्तयेत्~। आञ्छनैः}\ A}}\pend

\pstart\versepar \ledsidenote{\passref{4.3.19}}\edtext{\skt{सन्धीन्शरीरे}}{\lemma{\skt{सन्धीन्शरीरे}}\Afootnote{H; \skt{सन्धीञ्छरीरे}\ A}}\skt{ }\edtext{\skt{सर्व्वांस्तु}}{\lemma{\skt{तु}}\Afootnote{H; \skt{स्तु}\ A}}\skt{ चला\-न\-प्य\-च\-ला\-नपि~। }\pend
\pstart\versepar \skt{एतैस्तु स्थापनोपायैः }\edtext{\skt{स्थाप\-ये\-त्म\-ति\-मा\-न्भि\-षक्}}{\lemma{\skt{स्थापयेत्}}\Afootnote{H; \skt{स्थापयेन्}\ A}}\skt{~॥}\pend

\pstart\versepar \ledsidenote{\passref{4.3.20}}\edtext{\skt{उत्पिष्टमथ}}{\lemma{\skt{उत्पिष्टमथ}}\Afootnote{H; \skt{उत्पिष्टमथ}\ A}}\skt{ विश्लिष्टं सन्धिं वैद्यो न घट्टयेत्~। }\pend
\pstart\versepar \skt{तस्य }\edtext{\edtext{\skt{शीता\-त्प\-री\-षे\-का\-म्प्र\-दे\-हां\-श्चा\-व\-चा\-र\-येत्}}{\lemma{\skt{शीतात्}}\Afootnote{H; \skt{शीतान्}\ A}}}{\lemma{\skt{परीषेकाम्}}\Afootnote{H; \skt{परीषेकान्}\ A}}\skt{~॥}\pend

\pstart \ledsidenote{\passref{4.3.22}}\edtext{\skt{प्रत्यङ्गभग्नस्य}}{\lemma{\skt{प्रत्यङ्गभग्नस्य}}\Afootnote{H; \skt{पट्टोपरि कुशान्दत्त्वा यथा\appbrk{}व\appbrk{}द्व\appbrk{}न्ध\appbrk{}मा\appbrk{}च\appbrk{}रेत्~। प्रत्यङ्गभग्नस्य}\ A}}\skt{ }\edtext{\edtext{\skt{विधिमत}}{\lemma{\skt{विधिम्}}\Afootnote{H; \skt{विधि}\ A}}}{\lemma{\skt{अत}}\Afootnote{H; \skt{रत}\ A}}\skt{ }\edtext{\skt{ऊर्द्ध्वम्प्रवक्ष्महे}}{\lemma{\skt{प्रवक्ष्महे}}\Afootnote{H; \skt{प्रवक्ष्यते}\ A}}\skt{~।}\pend

\pstart\versepar \ledsidenote{\passref{4.3.23}}\skt{नखसन्धिं समुत्पिष्टं }\edtext{\skt{रक्षा\-नु\-ग\-त\-मा\-रया}}{\lemma{\skt{रक्षानुगतम्}}\Afootnote{H; \skt{रक्तानुगतम्}\ A}}\skt{~॥ }\pend
\pstart\versepar \skt{अवमथ्य }\edtext{\skt{श्रूते}}{\lemma{\skt{श्रूते}}\Afootnote{H; \skt{स्रुते}\ A}}\skt{ रक्ते शालिपिष्टेन लेपयेत्~॥}\pend

\pstart\versepar \ledsidenote{\passref{4.3.24}}\edtext{\skt{भग्ना वा}}{\lemma{\skt{भग्ना वा}}\Afootnote{H; \skt{भग्नां वा}\ A}}\skt{ }\edtext{\edtext{\skt{सन्धिमुक्ता}}{\lemma{\skt{सन्धिम्}}\Afootnote{H; \skt{सन्धि}\ A}}\skt{ वा}}{\lemma{\skt{उक्ता वा}}\Afootnote{H; \skt{मुक्तां वा}\ A}}\skt{ }\edtext{\skt{स्थाप\-यि\-त्वा\-ङ्गु\-ली\-म्भि\-षक्}}{\lemma{\skt{स्थाप\appbrk{}यि\appbrk{}त्वा\appbrk{}ङ्गु\appbrk{}लीम्}}\Afootnote{H; \skt{स्थाप\appbrk{}यि\appbrk{}त्वा+\appbrk{}अ\appbrk{}ङ्गुलीं}\ A}}\skt{~॥ }\pend
\pstart\versepar \skt{अनुनावेष्ट्य पट्टेन }\edtext{\edtext{\skt{घृत\-मे\-क\-म्प्र\-दा\-प\-येत्}}{\lemma{\skt{घृतम्}}\Afootnote{H; \skt{घृत}\ A}}}{\lemma{\skt{एकम्}}\Afootnote{H; \skt{सेकं}\ A}}\skt{~।}\edtext{}{\lemma{\skt{भिषक्~॥ अनुनावेष्ट्य}}\Afootnote{H; \skt{समाम्~। अणुना++आवेष्ट्य}\ A}}\pend

\pstart\versepar \ledsidenote{\passref{4.3.25}}\skt{अभ्यज्य }\edtext{\skt{सर्प्पिषा}}{\lemma{\skt{सर्प्पिषा}}\Afootnote{H; \skt{सपिषा}\ A}}\skt{ }\edtext{\skt{पादम्भग्नन्दन्वा}}{\lemma{\skt{भग्नन्दन्वा}}\Afootnote{H; \skt{तलभग्नं}\ A}}\skt{ }\edtext{\edtext{\skt{कुशाम्भिष्षक्}}{\lemma{\skt{कुशाम्}}\Afootnote{H; \skt{कुशोत्तरम्}\ A}}}{\lemma{\skt{भिष्षक्}}\Afootnote{H; \emph{om.}\ A}}\skt{~॥ }\pend
\pstart\versepar \skt{वस्त्रपट्टेन }\edtext{\skt{बध्नीयान्न तु}}{\lemma{\skt{बध्नीयान्न तु}}\Afootnote{H; \skt{बध्नीयान्न च}\ A}}\skt{ }\edtext{\skt{व्याया\-म\-मा\-च\-रेत्}}{\lemma{\skt{व्याया\appbrk{}म\appbrk{}मा\appbrk{}च\appbrk{}रेत्}}\Afootnote{H; \skt{व्याया\appbrk{}म\appbrk{}मा\appbrk{}च\appbrk{}रेत्}\ A}}\skt{~।}\pend

\pstart\versepar \ledsidenote{\passref{4.3.26}}\edtext{\skt{अभ्य\-ज्य\-स\-म\-ये\-जंघां}}{\lemma{\skt{अभ्य\appbrk{}ज्य\appbrk{}स\appbrk{}म\appbrk{}ये\appbrk{}जंघां}}\Afootnote{H; \skt{अभ्य\appbrk{}ज्या\appbrk{}या\appbrk{}म\appbrk{}ये\appbrk{}ज्जङ्}\ A}}\skt{ }\edtext{\edtext{\skt{घृते\-ना\-क्तां\-स्सं\-स्थितां}}{\lemma{\skt{घृतेनाक्तांस्}}\Afootnote{H; \skt{घामू}\ A}}}{\lemma{\skt{संस्थितां}}\Afootnote{H; \skt{रुं च सुसमाहितः}\ A}}\skt{~॥ }\pend
\pstart\versepar \skt{दत्वा वृक्षत्वचः }\edtext{\skt{शीता\-म्व\-स्त्र\-प\-ट्टेन}}{\lemma{\skt{शीताम्}}\Afootnote{H; \skt{शीता}\ A}}\skt{ वेष्टयेत्~।}\pend

\pstart\versepar \ledsidenote{\passref{4.3.27}}\edtext{\skt{चक्रयोगेन}}{\lemma{\skt{चक्रयोगेन}}\Afootnote{H; \skt{मति\appbrk{}मां\appbrk{}श्च\appbrk{}क्र\appbrk{}यो\appbrk{}गेन}\ A}}\skt{ }\edtext{\skt{मति\-मा\-ना\-ञ्छे\-दू\-र्द्ध्वस्थि}}{\lemma{\skt{मति\appbrk{}मा\appbrk{}ना\appbrk{}ञ्छे\appbrk{}दू\appbrk{}र्द्ध्वस्थि}}\Afootnote{H; \skt{ह्याञ्छेदूर्वस्थि}\ A}}\skt{ निर्ग्गतं~। }\pend
\pstart\versepar \skt{स्फुटितं }\edtext{\skt{पिच्चितम्वापि}}{\lemma{\skt{वापि}}\Afootnote{H; \skt{चापि}\ A}}\skt{ बध्नी\-या\-त्पू\-र्व्व\-व\-द्भि\-षक्~।}\pend

\pstart\versepar \ledsidenote{\passref{4.3.28}}\edtext{\skt{आञ्छे\-दू\-र्द्ध्व\-म\-ध\-श्चापि}}{\lemma{\skt{चापि}}\Afootnote{H; \skt{वा+अपि}\ A}}\skt{ कटिभग्नन्तु मानवं~॥ }\pend
\pstart\versepar \skt{ततः स्थानस्थिते सन्धौ बस्तिभिः समुपाचरेत्~॥}\pend

\pstart\versepar \ledsidenote{\passref{4.3.29}}\edtext{\edtext{\skt{पर्शुकास्वथ}}{\lemma{\skt{पर्शुकास्व्}}\Afootnote{H; \skt{पर्शुकास्}\ A}}}{\lemma{\skt{अथ}}\Afootnote{H; \skt{वथ}\ A}}\skt{ भग्नासु घृताभ्यक्तस्य }\edtext{\skt{पृष्ठतः}}{\lemma{\skt{पृष्ठतः}}\Afootnote{H; \skt{तिष्ठतः}\ A}}\skt{~॥ }\pend
\pstart\versepar \skt{}\edtext{\skt{दक्षिणामथवा}}{\lemma{\skt{दक्षिणाम्}}\Afootnote{H; \skt{दक्षिणास्व्}\ A}}\skt{ }\edtext{\edtext{\skt{वामा\-म\-नु\-मृ\-ज्ज्या\-नु\-ब\-न्धनीं}}{\lemma{\skt{वामाम्}}\Afootnote{H; \skt{वामास्व्}\ A}}}{\lemma{\skt{अनु\appbrk{}मृ\appbrk{}ज्ज्या\appbrk{}नु\appbrk{}ब\appbrk{}न्धनीं}}\Afootnote{H; \skt{अनुमृज्य निबन्धनीः}\ A}}\skt{~।}\pend

\pstart\versepar \ledsidenote{\passref{4.3.30}}\skt{ततः कवलिकान्दत्वा }\edtext{\skt{वेष्ट\-ये\-त्सु\-स\-मा\-हितं}}{\lemma{\skt{सुसमाहितं}}\Afootnote{H; \skt{सुसमाहितः}\ A}}\skt{~॥ }\pend
\pstart\versepar \skt{तैलपूर्ण्णे कटाहे वा }\edtext{\edtext{\skt{द्रोण्या\-म्वा\-प्य\-वे\-गा\-ह\-येत्}}{\lemma{\skt{वाप्य्}}\Afootnote{H; \skt{वा शा}\ A}}}{\lemma{\skt{अवेगाहयेत्}}\Afootnote{H; \skt{ययेन्नरम्}\ A}}\skt{~।}\pend

\pstart\versepar \ledsidenote{\passref{4.3.31}}\edtext{\edtext{\skt{मुश्ले\-नो\-त्क्षि\-पे\-त्क\-क्ष\-म\-क्ष\-स\-सन्धौ}}{\lemma{\skt{मुश्लेनोत्क्षिपेत्}}\Afootnote{H; \skt{मुस\appbrk{}ले\appbrk{}नो\appbrk{}त्क्षि\appbrk{}पेत्}\ A}}}{\lemma{\skt{कक्ष\appbrk{}म\appbrk{}क्ष\appbrk{}स\appbrk{}सन्धौ}}\Afootnote{H; \skt{कक्षामंससन्धौ}\ A}}\skt{ }\edtext{\skt{विसङ्गते}}{\lemma{\skt{विसङ्गते}}\Afootnote{H; \skt{विसंहते}\ A}}\skt{~॥ }\pend
\pstart\versepar \skt{स्थानस्थितञ्च बध्नी\-या\-त्स्व\-स्ति\-केन विचक्षणः~।}\pend

\pstart\versepar \ledsidenote{\passref{4.3.32}}\edtext{\skt{कौर्प्यरन्तु}}{\lemma{\skt{कौर्प्यरन्तु}}\Afootnote{H; \skt{कौर्परं तु}\ A}}\skt{ }\edtext{\skt{तथा}}{\lemma{\skt{तथा}}\Afootnote{A H\textsuperscript{pc}; \skt{था}\ H\textsuperscript{ac}}}\skt{ सन्धि\-म\-ङ्गु\-ष्ठे\-ना\-नु\-मा\-र्ज्ज\-येत्~॥ }\pend
\pstart\versepar \skt{अनुमृज्य }\edtext{\skt{तथा}}{\lemma{\skt{तथा}}\Afootnote{H; \skt{ततः}\ A}}\skt{ }\edtext{\edtext{\skt{सन्धि\-म्पी\-ड\-ये\-त्कु\-र्प्य\-रा\-च्च्यूतं}}{\lemma{\skt{कुर्प्यराच्}}\Afootnote{H; \skt{कूर्पराच्}\ A}}}{\lemma{\skt{च्यूतं}}\Afootnote{H; \skt{च्युतम्}\ A}}\skt{~।}\pend

\pstart\versepar \ledsidenote{\passref{4.3.33}}\edtext{\skt{प्रसार्य}}{\lemma{\skt{प्रसार्य}}\Afootnote{H; \skt{प्रसार्या}\ A}}\skt{ कुञ्चयेच्चैनं स्नेहसेकञ्च }\edtext{\skt{दापयेत्}}{\lemma{\skt{दापयेत्}}\Afootnote{H; \skt{दापये}\ A}}\skt{~॥ }\pend
\pstart\versepar \skt{}\edtext{\skt{मणिबन्धे च}}{\lemma{\skt{मणिबन्धे च}}\Afootnote{H; \skt{त्~। एवं}\ A}}\skt{ }\edtext{\skt{जानौ च}}{\lemma{\skt{जानौ च}}\Afootnote{H; \skt{जानुनि}\ A}}\skt{ गुल्फे }\edtext{\skt{चैव}}{\lemma{\skt{चैव}}\Afootnote{H; \skt{च म}\ A}}\skt{ }\edtext{\skt{समाचरेत्}}{\lemma{\skt{समाचरेत्}}\Afootnote{H; \skt{णिबन्धे च कारयेत्}\ A}}\skt{~।}\pend

\pstart\versepar \ledsidenote{\passref{4.3.34}}\skt{उभे तले समे कृत्वा तलभग्नस्य देहिनः~॥ }\pend
\pstart\versepar \skt{बध्नी\-या\-दा\-म\-तै\-लेन परिषेकञ्च }\edtext{\skt{दापयेत्}}{\lemma{\skt{दापयेत्}}\Afootnote{H; \skt{कारयेत्}\ A}}\skt{~॥}\pend

\pstart\versepar \ledsidenote{\passref{4.3.35}}\skt{मृत्पि\-ण्ड\-न्धा\-र\-ये\-त्पूर्व्वं }\edtext{\edtext{\skt{लवणाण्डमतः}}{\lemma{\skt{लवणाण्डम्}}\Afootnote{H; \skt{लवणं च}\ A}}}{\lemma{\skt{अतः}}\Afootnote{H; \skt{ततः}\ A}}\skt{ परं~॥ }\pend
\pstart\versepar \skt{}\edtext{\skt{हस्तेन}}{\lemma{\skt{हस्तेन}}\Afootnote{H; \skt{हस्ते}\ A}}\skt{ जातबले चापि }\edtext{\skt{कार्य\-म्पा\-षा\-ण\-धा\-रणं}}{\lemma{\skt{कार्यम्}}\Afootnote{H; \skt{कुर्यात्}\ A}}\skt{~।}\pend

\pstart\versepar \ledsidenote{\passref{4.3.36}}\edtext{\edtext{\edtext{\edtext{\skt{सन्न\-मु\-न्ना\-म\-ये\-त्स्नि\-न्न\-म\-क्ष\-क\-म्मु\-श\-लेन}}{\lemma{\skt{उन्नामयेत्}}\Afootnote{H; \skt{उन्नमयेत्}\ A}}}{\lemma{\skt{स्निन्नम्}}\Afootnote{H; \skt{स्विन्न}\ A}}}{\lemma{\skt{अक्षकम्}}\Afootnote{H; \skt{मक्षकं}\ A}}\skt{ तु}}{\lemma{\skt{मुशलेन तु}}\Afootnote{H; \skt{मुसलेन तु}\ A}}\skt{~॥ }\pend
\pstart\versepar \skt{तथोन्नतं }\edtext{\skt{पीडयीत}}{\lemma{\skt{पीडयीत}}\Afootnote{H; \skt{पीडयेच्च}\ A}}\skt{ }\edtext{\edtext{\edtext{\skt{बध्नी\-या\-द्गा\-ढ\-मेव}}{\lemma{\skt{बध्नीयाद्}}\Afootnote{H; \skt{बध्नीयाद्}\ A}}}{\lemma{\skt{गाढम्}}\Afootnote{H; \skt{गाढ}\ A}}}{\lemma{\skt{एव}}\Afootnote{H; \skt{मेव}\ A}}\skt{ च~।}\pend

\pstart\versepar \ledsidenote{\passref{4.3.37}}\edtext{\skt{ऊरूवच्चापि}}{\lemma{\skt{ऊरूवच्}}\Afootnote{H; \skt{ऊरुवच्}\ A}}\skt{ }\edtext{\skt{कर्त्त\-व्य\-म्बा\-हु\-भग्ने}}{\lemma{\skt{बाहुभग्ने}}\Afootnote{H; \skt{बाहुभग्न}\ A}}\skt{ चिकित्सितं~॥ }\pend
\pstart\versepar \skt{ग्रीवायान्तु }\edtext{\skt{विवृ\-त्ता\-या\-म्प्र\-वि\-ष्टा\-या\-म\-थापि}}{\lemma{\skt{प्रवि\appbrk{}ष्टा\appbrk{}या\appbrk{}म\appbrk{}थापि}}\Afootnote{H; \skt{प्रवि\appbrk{}ष्टा\appbrk{}या\appbrk{}म\appbrk{}धोऽ\appbrk{}अपि}\ A}}\skt{ वा~।}\pend

\pstart\versepar \ledsidenote{\passref{4.3.38}}\edtext{\skt{आवट्यामथ}}{\lemma{\skt{आवट्यामथ}}\Afootnote{H; \skt{अवटावथ}\ A}}\skt{ हन्वोश्च }\edtext{\skt{प्रगृ\-ह्यो\-न्ना\-म\-ये\-न्नरं}}{\lemma{\skt{प्रगृ\appbrk{}ह्यो\appbrk{}न्ना\appbrk{}म\appbrk{}येन्}}\Afootnote{H; \skt{प्रगृ\appbrk{}ह्यो\appbrk{}न्न\appbrk{}म\appbrk{}येन्}\ A}}\skt{~। }\pend
\pstart\versepar \skt{ततः }\edtext{\edtext{\skt{कुशीसमङ्कृत्वा}}{\lemma{\skt{कुशीसमङ्}}\Afootnote{H; \skt{कुशां समं}\ A}}}{\lemma{\skt{कृत्वा}}\Afootnote{H; \skt{दत्त्वा}\ A}}\skt{ वस्त्रपट्टेन वेष्टयेत्~।}\pend

\pstart\versepar \ledsidenote{\passref{4.3.39}}\skt{उत्तानं शाययेच्चैनं }\edtext{\skt{सप्तरात्रं}}{\lemma{\skt{सप्तरात्रं}}\Afootnote{H; \skt{सप्तरात्र}\ A}}\skt{ }\edtext{\skt{सुयन्त्रितः}}{\lemma{\skt{सुयन्त्रितः}}\Afootnote{H; \skt{मतन्द्रियः}\ A}}\skt{~॥ }\pend
\pstart\versepar \skt{}\edtext{\skt{हन्वस्थीनि}}{\lemma{\skt{हन्वस्थीनि}}\Afootnote{H; \skt{हन्वस्थिनी}\ A}}\skt{ समानीय हनुसन्धौ }\edtext{\skt{विसङ्गते}}{\lemma{\skt{विसङ्गते}}\Afootnote{H; \skt{विसंहते}\ A}}\skt{~।}\pend

\pstart\versepar \ledsidenote{\passref{4.3.40}}\skt{स्वेदयित्वा }\edtext{\skt{स्थिते}}{\lemma{\skt{स्थिते}}\Afootnote{H; \skt{स्थितेन्}\ A}}\skt{ }\edtext{\skt{चापि}}{\lemma{\skt{चापि}}\Afootnote{H; \skt{सम्यक्}\ A}}\skt{ }\edtext{\skt{पञ्चा\-ङ्गी\-म्वि\-भ\-जे\-द्भि\-षक्}}{\lemma{\skt{विभजेद्}}\Afootnote{H; \skt{वितरेद्}\ A}}\skt{~॥ }\pend
\pstart\versepar \skt{}\edtext{\skt{घृत\-म्वा\-त\-घ्न\-म\-धुरैः}}{\lemma{\skt{घृतम्}}\Afootnote{H; \emph{om.}\ A}}\skt{ }\edtext{\edtext{\skt{सिद्धन्नस्तञ्च}}{\lemma{\skt{सिद्धन्}}\Afootnote{H; \skt{सर्पिः सिद्धं}\ A}}}{\lemma{\skt{नस्तञ्च}}\Afootnote{H; \skt{नस्ये च}\ A}}\skt{ पूजितं~।}\pend

\pstart \ledsidenote{\passref{4.3.41}}\edtext{\edtext{\edtext{\skt{दन्ता\-न\-भ\-ग्ना\-न्च\-लि\-ता\-न्स\-र\-क्ता\-नेव}}{\lemma{\skt{दन्तानभग्नान्}}\Afootnote{H; \skt{अभग्नांश्}\ A}}}{\lemma{\skt{चलितान्}}\Afootnote{H; \skt{चलितान्दन्तान्}\ A}}}{\lemma{\skt{सरक्तानेव}}\Afootnote{H; \skt{सरक्तानव}\ A}}\skt{ }\edtext{\skt{पीडयेत्}}{\lemma{\skt{पीडयेत्}}\Afootnote{H; \skt{पीडयेत्~। तरुणस्य मनुष्यस्य शीतै\appbrk{}रा\appbrk{}ले\appbrk{}प\appbrk{}ये\appbrk{}द्वहिः}\ A}}\skt{~॥}\pend

\pstart \ledsidenote{\passref{4.3.43ab}}\edtext{\skt{वयस्थस्य}}{\lemma{\skt{वयस्थस्य}}\Afootnote{H; \skt{जीर्णस्य तु}\ A}}\skt{ }\edtext{\skt{मनुष्यस्य}}{\lemma{\skt{मनुष्यस्य}}\Afootnote{H; \skt{मनुष्य}\ A}}\skt{ }\edtext{\skt{जीर्ण्णस्य तु}}{\lemma{\skt{जीर्ण्णस्य तु}}\Afootnote{H; \skt{स्य}\ A}}\skt{ }\edtext{\skt{विवर्ज्जयेत्}}{\lemma{\skt{विवर्ज्जयेत्}}\Afootnote{H; \skt{वर्ज\appbrk{}ये\appbrk{}च्च\appbrk{}लि\appbrk{}ता\appbrk{}न्द्वि\appbrk{}जान्}\ A}}\skt{~।}\pend

\pstart \ledsidenote{\passref{4.3.42ab}}\edtext{\skt{संसिञ्च्य}}{\lemma{\skt{संसिञ्च्य}}\Afootnote{H; \skt{सिक्त्वा+अम्}\ A}}\skt{ }\edtext{\skt{वारिभिः}}{\lemma{\skt{वारिभिः}}\Afootnote{H; \skt{बुभिस्ततः}\ A}}\skt{ शीतैः सन्धा\-नी\-यै\-रु\-पा\-च\-रेत्~॥}\pend

\pstart \ledsidenote{\passref{4.3.43cd}}\skt{नासां सन्नाम्विवृत्तां वा ऋज्वीं कृत्वा शलाकया~।}\pend

\pstart\versepar \ledsidenote{\passref{4.3.44}}\edtext{\skt{प्रत्ये\-क\-न\-स्त\-तो\-नाडीं}}{\lemma{\skt{प्रत्ये\appbrk{}क\appbrk{}न\appbrk{}स्त\appbrk{}तो\appbrk{}नाडीं}}\Afootnote{H; \skt{पृथ\appbrk{}ङ्न\appbrk{}सि\appbrk{}क\appbrk{}यो\appbrk{}र्नाड्यौ}\ A}}\skt{ }\edtext{\skt{द्विमुखीं}}{\lemma{\skt{द्विमुखीं}}\Afootnote{H; \skt{द्विमुख्यौ}\ A}}\skt{ सम्प्रवेशयेत्~॥ }\pend
\pstart\versepar \skt{ततः पट्टेन संवेष्ट्य }\edtext{\skt{घृत\-से\-का\-म्प्र\-दा\-प\-येत्}}{\lemma{\skt{घृतसेकाम्}}\Afootnote{H; \skt{घृतसेकं}\ A}}\skt{~।}\pend

\pstart\versepar \ledsidenote{\passref{4.3.45}}\skt{भग्नं }\edtext{\skt{कर्ण्णं}}{\lemma{\skt{कर्ण्णं}}\Afootnote{H; \skt{कर्णं तु}\ A}}\skt{ }\edtext{\skt{घृताभ्यक्तं}}{\lemma{\skt{घृताभ्यक्तं}}\Afootnote{H; \skt{बध्नीयात्समं कृत्वा घृत}\ A}}\skt{ }\edtext{\skt{स्थाप्य}}{\lemma{\skt{स्थाप्य}}\Afootnote{H; \skt{प्लु}\ A}}\skt{ }\edtext{\skt{बन्धेन}}{\lemma{\skt{बन्धेन}}\Afootnote{H; \skt{तम्}\ A}}\skt{ }\edtext{\skt{योजयेत्}}{\lemma{\skt{योजयेत्}}\Afootnote{H; \emph{om.}\ A}}\skt{~॥ }\pend
\pstart\versepar \skt{सद्यः\-क्ष\-त\-वि\-धा\-नञ्च }\edtext{\skt{नयने}}{\lemma{\skt{नयने}}\Afootnote{H; \emph{om.}\ A}}\skt{ }\edtext{\skt{निर्ग्गते}}{\lemma{\skt{निर्ग्गते}}\Afootnote{H; \skt{त}\ A}}\skt{ }\edtext{\skt{हितं}}{\lemma{\skt{हितं}}\Afootnote{H; \skt{तः पञ्चा\appbrk{}त्स\appbrk{}मा\appbrk{}च\appbrk{}रेत्}\ A}}\skt{~।}\pend

\pstart\versepar \ledsidenote{\passref{4.3.46}}\skt{मस्तुलुङ्गाद्विना भिन्ने कपाले }\edtext{\skt{मधुसर्प्पिषा}}{\lemma{\skt{मधुसर्प्पिषा}}\Afootnote{H; \skt{मधुसर्पिषी}\ A}}\skt{~॥ }\pend
\pstart\versepar \skt{दत्त्वा }\edtext{\edtext{\skt{बन्धेत्ततः}}{\lemma{\skt{बन्धेत्}}\Afootnote{H; \skt{ततो निब\appbrk{}ध्नी\appbrk{}या\appbrk{}त्स\appbrk{}प्ताहं च पिबेद्}\ A}}}{\lemma{\skt{ततः}}\Afootnote{H; \skt{घृतम्}\ A}}\skt{ }\edtext{\skt{सर्प्पिः}}{\lemma{\skt{सर्प्पिः}}\Afootnote{H; \emph{om.}\ A}}\skt{ }\edtext{\edtext{\edtext{\skt{सप्त्रा\-त्र\-म्पि\-बे\-न्नरः}}{\lemma{\skt{सप्त्रात्रम्}}\Afootnote{H; \emph{om.}\ A}}}{\lemma{\skt{पिबेन्}}\Afootnote{H; \emph{om.}\ A}}}{\lemma{\skt{नरः}}\Afootnote{H; \emph{om.}\ A}}\skt{~॥}\pend

\pstart\versepar \ledsidenote{\passref{4.3.47}}\skt{पत\-ना\-द\-भि\-घा\-ताद्वा }\edtext{\skt{सूनमङ्गं}}{\lemma{\skt{सूनमङ्गं}}\Afootnote{H; \skt{शूनमङ्गं}\ A}}\skt{ यदक्षतं~। }\pend
\pstart\versepar \skt{}\edtext{\skt{शीतां}}{\lemma{\skt{शीतां}}\Afootnote{H; \skt{शीता}\ A}}\skt{ }\edtext{\edtext{\skt{सेकाम्प्रदेहांश्च}}{\lemma{\skt{सेकाम्}}\Afootnote{H; \skt{न्}\ A}}}{\lemma{\skt{प्रदेहांश्च}}\Afootnote{H; \skt{प्रदेहान्सेकांश्च}\ A}}\skt{ भिष\-क्त\-स्या\-व\-चा\-र\-येत्~॥}\pend

\pstart\versepar \ledsidenote{\passref{4.3.48}}\edtext{\edtext{\skt{कटि\-जं\-घो\-रु\-भ\-ग्ना\-ना\-ङ्क\-वा\-ट\-श\-यनं}}{\lemma{\skt{कटि\appbrk{}जं\appbrk{}घो\appbrk{}रु\appbrk{}भ\appbrk{}ग्ना\appbrk{}नाङ्}}\Afootnote{H; \skt{अथ जङ्घोरुभग्नानां}\ A}}}{\lemma{\skt{कवाटशयनं}}\Afootnote{H; \skt{कपाटशयनं}\ A}}\skt{ हितं~। }\pend
\pstart\versepar \skt{}\edtext{\skt{यन्त्रणार्थे च}}{\lemma{\skt{यन्त्रणार्थे च}}\Afootnote{H; \skt{कीलका बन्धनार्थं च}\ A}}\skt{ }\edtext{\skt{कीलाः}}{\lemma{\skt{कीलाः}}\Afootnote{H; \emph{om.}\ A}}\skt{ }\edtext{\edtext{\skt{स्युस्तत्र}}{\lemma{\skt{स्युस्}}\Afootnote{H; \emph{om.}\ A}}}{\lemma{\skt{तत्र}}\Afootnote{H; \skt{पञ्च}\ A}}\skt{ कार्या }\edtext{\skt{निबन्धनाः}}{\lemma{\skt{निबन्धनाः}}\Afootnote{\skt{विजानता}\ A; \skt{निबन्धना}\ H}}\skt{~॥}\pend

\pstart\versepar \ledsidenote{\passref{4.3.49}}\skt{यथा न चलनन्तस्य भग्नस्य }\edtext{\skt{प्रयतेत्तथा}}{\lemma{\skt{प्रयतेत्}}\Afootnote{H; \skt{क्रियते}\ A}}\skt{~। }\pend
\pstart\versepar \skt{सन्धेरुभयतो द्वौ द्वौ तले }\edtext{\skt{चैकस्तथा}}{\lemma{\skt{तथा}}\Afootnote{H; \skt{च की}\ A}}\skt{ }\edtext{\skt{भवेत्}}{\lemma{\skt{भवेत्}}\Afootnote{H; \skt{लकः}\ A}}\skt{~॥}\pend

\pstart\versepar \ledsidenote{\passref{4.3.50}}\edtext{\skt{श्रोण्या वा}}{\lemma{\skt{श्रोण्या वा}}\Afootnote{H; \skt{श्रोण्यां वा}\ A}}\skt{ पृष्ठवंशे वा }\edtext{\skt{वक्षः}}{\lemma{\skt{वक्षः}}\Afootnote{H; \skt{वक्षस्य्}\ A}}\skt{ }\edtext{\skt{कक्षाङ्शयोस्तथा}}{\lemma{\skt{कक्षाङ्शयोस्}}\Afootnote{H; \skt{अक्षकयोस्}\ A}}\skt{~। }\pend
\pstart\versepar \skt{भग्न\-स\-न्धि\-वि\-मो\-क्षेषु }\edtext{\skt{विधिमेवं}}{\lemma{\skt{विधिमेवं}}\Afootnote{H; \skt{विधिमेनं}\ A}}\skt{ }\edtext{\skt{समारभते}}{\lemma{\skt{समारभते}}\Afootnote{H; \skt{समाचरेत्}\ A}}\skt{~॥}\pend

\pstart\versepar \ledsidenote{\passref{4.3.51}}\edtext{\skt{सन्धी\-न्चि\-र\-वि\-मु\-क्तांस्तु}}{\lemma{\skt{सन्धीन्}}\Afootnote{H; \skt{सन्धींश्}\ A}}\skt{ }\edtext{\skt{स्निग्धस्विन्ना}}{\lemma{\skt{स्निग्धस्विन्ना}}\Afootnote{H; \skt{स्निग्धान्स्विन्नान्}\ A}}\skt{ मृदूकृतान्~। }\pend
\pstart\versepar \skt{उक्तै\-र्व्वि\-धा\-नै\-र्ब्बुद्ध्या च }\edtext{\skt{यथा}}{\lemma{\skt{यथा}}\Afootnote{H; \skt{सम्यक्}\ A}}\skt{ प्रकृ\-ति\-मा\-न\-येत्~॥}\pend

\pstart\versepar \ledsidenote{\passref{4.3.52}}\skt{काण्डभग्ने प्ररूढे तु }\edtext{\skt{विष\-मो\-ल्ब\-ण\-सं\-स्थिते}}{\lemma{\skt{विष\appbrk{}मो\appbrk{}ल्ब\appbrk{}ण\appbrk{}सं\appbrk{}स्थिते}}\Afootnote{H; \skt{विष\appbrk{}मो\appbrk{}ल्ब\appbrk{}ण\appbrk{}सं\appbrk{}हिते}\ A}}\skt{~। }\pend
\pstart\versepar \skt{}\edtext{\skt{आपोह्य}}{\lemma{\skt{आपोह्य}}\Afootnote{H; \skt{आपोथ्य}\ A}}\skt{ }\edtext{\skt{भङ्गं}}{\lemma{\skt{भङ्गं}}\Afootnote{H; \emph{om.}\ A}}\skt{ }\edtext{\skt{समये}}{\lemma{\skt{समये}}\Afootnote{H; \skt{समयेद्}\ A}}\skt{ }\edtext{\skt{ततो}}{\lemma{\skt{ततो}}\Afootnote{H; \skt{भग्नं ततो}\ A}}\skt{ भग्न\-व\-दा\-च\-रेत्~॥}\pend

\pstart\versepar \ledsidenote{\passref{4.3.54}}\skt{ऊर्द्ध्वकाये तु }\edtext{\skt{भग्ना\-ना\-म्म\-स्ति\-ष्क\-ङ्क\-र्ण्ण\-पू\-रणं}}{\lemma{\skt{मस्तिष्कङ्}}\Afootnote{H; \skt{मस्तिष्क्यं}\ A}}\skt{~। }\pend
\pstart\versepar \skt{}\edtext{\skt{शस्यते}}{\lemma{\skt{शस्यते}}\Afootnote{H; \emph{om.}\ A}}\skt{ }\edtext{\edtext{\skt{नस्यपानञ्च}}{\lemma{\skt{नस्यपानञ्}}\Afootnote{H; \skt{घृतपानं}\ A}}}{\lemma{\skt{च}}\Afootnote{H; \skt{हितं नस्यं}\ A}}\skt{ प्रशा\-खा\-स्व\-नु\-वा\-सनं~॥}\pend

\pstart\versepar \ledsidenote{\passref{4.3.55}}\skt{अत ऊर्द्ध्वम्प्रवक्ष्यामि तैलं भग्नप्रसाधकं~। }\pend
\pstart\versepar \skt{रात्रौ रात्रौ }\edtext{\skt{तिलां}}{\lemma{\skt{तिलां}}\Afootnote{H; \skt{तिलान्}\ A}}\skt{ }\edtext{\skt{कृष्णां}}{\lemma{\skt{कृष्णां}}\Afootnote{H; \skt{कृष्णान्}\ A}}\skt{ वासयेदस्थिरे जले~॥}\pend

\pstart\versepar \ledsidenote{\passref{4.3.56}}\skt{दिवा दिवा शोषयित्वा }\edtext{\skt{क्षीरेण}}{\lemma{\skt{क्षीरेण}}\Afootnote{H; \skt{गवां क्षीरेण}\ A}}\skt{ }\edtext{\skt{सह}}{\lemma{\skt{सह}}\Afootnote{H; \emph{om.}\ A}}\skt{ भावयेत्~। }\pend
\pstart\versepar \skt{}\edtext{\skt{तृतीये}}{\lemma{\skt{तृतीये}}\Afootnote{H; \skt{तृतीयं}\ A}}\skt{ }\edtext{\skt{सप्तरात्रे तु}}{\lemma{\skt{सप्तरात्रे तु}}\Afootnote{H; \skt{सप्तरात्रं तु}\ A}}\skt{ भाव\-ये\-न्म\-धु\-का\-म्बुना~॥}\pend

\pstart\versepar \ledsidenote{\passref{4.3.57}}\skt{ततः क्षीरम्पुनः }\edtext{\skt{पीतां}}{\lemma{\skt{पीतां}}\Afootnote{H; \skt{पीतान्}\ A}}\skt{ }\edtext{\skt{शुष्कां}}{\lemma{\skt{शुष्कां}}\Afootnote{H; \skt{सुशुष्कांश्}\ A}}\skt{ }\edtext{\edtext{\skt{सूक्षा\-म्वि\-चू\-र्ण्ण\-येत्}}{\lemma{\skt{सूक्षाम्}}\Afootnote{H; \emph{om.}\ A}}}{\lemma{\skt{विचूर्ण्णयेत्}}\Afootnote{H; \skt{चूर्णयेद्}\ A}}\skt{~। }\pend
\pstart\versepar \skt{}\edtext{\skt{काकोल्यादि}}{\lemma{\skt{काकोल्यादि}}\Afootnote{H; \skt{भिषक्~। काकोल्यादिं}\ A}}\skt{ }\edtext{\skt{स्वदम्ष्ट्राम्मञ्जिष्ठां}}{\lemma{\skt{स्वदम्ष्ट्राम्}}\Afootnote{H; \skt{सयष्ट्याह्वं}\ A}}\skt{ सारिवान्तथा~॥}\pend

\pstart\versepar \ledsidenote{\passref{4.3.58}}\skt{कुष्ठं }\edtext{\skt{सर्ज्जरसम्मांसी}}{\lemma{\skt{मांसी}}\Afootnote{H; \skt{मांसीं}\ A}}\skt{ सुरदारु सचन्दनं~। }\pend
\pstart\versepar \skt{}\edtext{\skt{शतपुष्पञ्च}}{\lemma{\skt{शतपुष्पञ्च}}\Afootnote{H; \skt{शतपुष्पां च}\ A}}\skt{ सञ्चूर्ण्ण्य तिलचूर्ण्णेन योजयेत्~॥}\pend

\pstart\versepar \ledsidenote{\passref{4.3.59}}\edtext{\skt{पीडनार्थे च}}{\lemma{\skt{पीडनार्थे च}}\Afootnote{H; \skt{पीडनार्थं च}\ A}}\skt{ कर्त्तव्यं }\edtext{\skt{सर्व्व\-ग\-न्ध\-स्रि\-त\-म्पयः}}{\lemma{\skt{सर्व्व\appbrk{}ग\appbrk{}न्ध\appbrk{}स्रि\appbrk{}तम्}}\Afootnote{H; \skt{सर्वगन्धशृतं}\ A}}\skt{~। }\pend
\pstart\versepar \skt{चतुर्ग्गुणेन पयसा तत्तैलं }\edtext{\edtext{\skt{योजयेत्पुनः}}{\lemma{\skt{योजयेत्}}\Afootnote{H; \skt{विपचेद्}\ A}}}{\lemma{\skt{पुनः}}\Afootnote{H; \skt{भिषक्}\ A}}\skt{~॥}\pend

\pstart\versepar \ledsidenote{\passref{4.3.60}}\skt{एला\-मं\-शु\-म\-ती\-म्पत्रं जीव\-क\-न्त\-ग\-र\-न्तथा~। }\pend
\pstart\versepar \skt{}\edtext{\skt{लोध्रं}}{\lemma{\skt{लोध्रं}}\Afootnote{H; \skt{रोध्रं}\ A}}\skt{ }\edtext{\skt{प्रपोण्डरीकञ्च}}{\lemma{\skt{प्रपोण्डरीकञ्च}}\Afootnote{H; \skt{प्रपौण्डरीकं च}\ A}}\skt{ तथा }\edtext{\skt{काला\-नु\-सा\-रि\-वान्}}{\lemma{\skt{काला\appbrk{}नु\appbrk{}सा\appbrk{}रि\appbrk{}वान्}}\Afootnote{H; \skt{काला\appbrk{}नु\appbrk{}सा\appbrk{}रि(\appbrk{}वा)\appbrk{}णम्}\ A}}\skt{~॥}\pend

\pstart\versepar \ledsidenote{\passref{4.3.61}}\edtext{\skt{शैलेयकं}}{\lemma{\skt{शैलेयकं}}\Afootnote{H; \skt{सैरेयकं}\ A}}\skt{ }\edtext{\edtext{\skt{क्षीर\-शु\-क्ला\-म\-नन्ता}}{\lemma{\skt{क्षीरशुक्लाम्}}\Afootnote{H; \skt{क्षीरशुक्ला}\ A}}}{\lemma{\skt{अनन्ता}}\Afootnote{H; \skt{मनन्तां}\ A}}\skt{ }\edtext{\skt{समधूलिकान्}}{\lemma{\skt{समधूलिकान्}}\Afootnote{H; \skt{समधूलिकाम्}\ A}}\skt{~। }\pend
\pstart\versepar \skt{पिष्ट्वा शृङ्गाटकञ्चैव }\edtext{\skt{प्रागु\-क्ता\-न्यौ\-ष\-धानि}}{\lemma{\skt{प्रागुक्तान्य्}}\Afootnote{H; \skt{पूर्वोक्तान्य्}\ A}}\skt{ च~॥}\pend

\pstart\versepar \ledsidenote{\passref{4.3.62}}\skt{एभि\-स्त\-द्वि\-प\-चे\-त्तैलं }\edtext{\skt{शास्त्र\-वि\-न्मृ\-दु\-ना\-ग्निना}}{\lemma{\skt{शास्त्र\appbrk{}वि\appbrk{}न्मृ\appbrk{}दु\appbrk{}ना\appbrk{}ग्निना}}\Afootnote{H; \skt{शास्त्र\appbrk{}वि\appbrk{}न्मृ\appbrk{}दु\appbrk{}नाऽ\appbrk{}ग्निना}\ A}}\skt{~। }\pend
\pstart\versepar \skt{एतत्तैलं सदा पथ्यं भग्नानां सर्व्वकर्म्मसु~॥}\pend

\pstart\versepar \ledsidenote{\passref{4.3.63}}\skt{आक्षेपके पक्षघाते तालुशोषे }\edtext{\skt{तथार्द्दिते}}{\lemma{\skt{तथार्द्दिते}}\Afootnote{H; \skt{तथाऽर्दिते}\ A}}\skt{~। }\pend
\pstart\versepar \skt{मन्यास्तम्भे शिरोरोगे }\edtext{\skt{कर्ण्णमूले}}{\lemma{\skt{कर्ण्णमूले}}\Afootnote{H; \skt{कर्णशूले}\ A}}\skt{ हनुग्रहे~॥}\pend

\pstart\versepar \ledsidenote{\passref{4.3.64}}\skt{बाधिर्ये तिमिरे चैव }\edtext{\skt{येषु}}{\lemma{\skt{येषु}}\Afootnote{H; \skt{ये च}\ A}}\skt{ स्त्रीषु क्षयङ्गताः~। }\pend
\pstart\versepar \skt{पथ्यम्पाने }\edtext{\skt{तथाभ्यङ्गे}}{\lemma{\skt{तथाभ्यङ्गे}}\Afootnote{H; \skt{तथाऽभ्यङ्गे}\ A}}\skt{ नस्ये बस्तिषु भोजने~॥}\pend

\pstart\versepar \ledsidenote{\passref{4.3.65}}\skt{ग्रीवास्कन्धोरसां }\edtext{\skt{वृद्धि\-र\-ने\-नै\-वो\-प\-जा\-यते}}{\lemma{\skt{वृद्धि\appbrk{}र\appbrk{}ने\appbrk{}नै\appbrk{}वो\appbrk{}प\appbrk{}जा\appbrk{}यते}}\Afootnote{H; \skt{वृद्धि\appbrk{}र\appbrk{}मु\appbrk{}नै\appbrk{}वो\appbrk{}प\appbrk{}जा\appbrk{}यते}\ A}}\skt{~। }\pend
\pstart\versepar \skt{}\edtext{\skt{मुखं}}{\lemma{\skt{मुखं}}\Afootnote{H; \skt{मुखं च}\ A}}\skt{ }\edtext{\skt{पद्मप्रतिसमं}}{\lemma{\skt{पद्मप्रतिसमं}}\Afootnote{H; \skt{पद्मप्रतिमं}\ A}}\skt{ }\edtext{\skt{सुस\-ग्र\-न्धि\-स\-मी\-रणे}}{\lemma{\skt{सुस\appbrk{}ग्र\appbrk{}न्धि\appbrk{}स\appbrk{}मी\appbrk{}रणे}}\Afootnote{H; \skt{ससु\appbrk{}ग\appbrk{}न्धि\appbrk{}स\appbrk{}मी\appbrk{}र\appbrk{}णम्}\ A}}\skt{~॥}\pend

\pstart\versepar \ledsidenote{\passref{4.3.66}}\skt{गन्धतैलमिदं नाम्ना सर्व्व\-वा\-त\-वि\-का\-र\-नुत्~। }\pend
\pstart\versepar \skt{राजा\-र्ह\-मे\-त\-त्क\-र्त्तव्यं राज्ञामेव विचक्षणैः~॥}\pend

\pstart\versepar \ledsidenote{\passref{4.3.67}}\edtext{\skt{त्रपु\-षा\-क्ष\-पि\-या\-लानां}}{\lemma{\skt{त्रपु\appbrk{}षा\appbrk{}क्ष\appbrk{}पि\appbrk{}या\appbrk{}लानां}}\Afootnote{H; \skt{त्रपु\appbrk{}षा\appbrk{}क्ष्प्रि\appbrk{}या\appbrk{}लानां}\ A}}\skt{ तैलानि मधुरैः सह~। }\pend
\pstart\versepar \skt{वसान्दत्त्वा }\edtext{\skt{यथालाभे}}{\lemma{\skt{यथालाभे}}\Afootnote{H; \skt{यथालाभं}\ A}}\skt{ क्षीरे दशगुणे पचेत्~॥}\pend

\pstart\versepar \ledsidenote{\passref{4.3.68}}\edtext{\edtext{\edtext{\skt{स्नेहो\-त्त\-म\-मि\-द\-ञ्चात्र}}{\lemma{\skt{स्नेहोत्तमम्}}\Afootnote{H; \skt{स्नेहोत्तम}\ A}}}{\lemma{\skt{इदञ्}}\Afootnote{H; \skt{मिदं}\ A}}}{\lemma{\skt{चात्र}}\Afootnote{H; \skt{चाशु}\ A}}\skt{ कुर्या\-द्भ\-ग्न\-प्र\-सा\-धनं~। }\pend
\pstart\versepar \skt{पाना\-भ्य\-ञ्ज\-न\-न\-स्येषु बस्तिकर्म्मणि }\edtext{\skt{सेवने}}{\lemma{\skt{सेवने}}\Afootnote{H; \skt{सेचने}\ A}}\skt{~॥}\pend

\pstart\versepar \ledsidenote{\passref{4.3.69}}\skt{भग्नन्नैति यथा\-पा\-क\-म्प्र\-य\-तेत तथा भिषक्~। }\pend
\pstart\versepar \skt{}\edtext{\skt{पक्वं}}{\lemma{\skt{पक्वं}}\Afootnote{H; \skt{पक्व}\ A}}\skt{ }\edtext{\edtext{\skt{हिंस्या\-त्सि\-रा\-स्ना\-यु\-न्तद्धि}}{\lemma{\skt{हिंस्यात्}}\Afootnote{H; \skt{मांस}\ A}}}{\lemma{\skt{सिरास्नायुन्तद्धि}}\Afootnote{H; \skt{सिरास्नायु तद्धि}\ A}}\skt{ कृच्छ्रेण सिद्ध्यति~॥}\pend

\pstart\versepar \ledsidenote{\passref{4.3.70}}\edtext{\skt{भग्न\-स\-न्धि\-म\-ना\-वि\-द्ध\-म\-ही\-नां\-ग\-म\-नु\-ल्बणं}}{\lemma{\skt{भग्न\appbrk{}स\appbrk{}न्धि\appbrk{}म\appbrk{}ना\appbrk{}वि\appbrk{}द्ध\appbrk{}म\appbrk{}ही\appbrk{}नां\appbrk{}ग\appbrk{}म\appbrk{}नु\appbrk{}ल्बणं}}\Afootnote{H; \skt{भग्नं सन्धि\appbrk{}म\appbrk{}ना\appbrk{}वि\appbrk{}द्ध\appbrk{}म\appbrk{}ही\appbrk{}ना\appbrk{}ङ्ग\appbrk{}म\appbrk{}नु\appbrk{}ल्ब\appbrk{}णम्}\ A}}\skt{~। }\pend
\pstart\versepar \skt{}\edtext{\skt{सुख\-चे\-ष्टा\-प्र\-सा\-रञ्च}}{\lemma{\skt{सुख\appbrk{}चे\appbrk{}ष्टा\appbrk{}प्र\appbrk{}सा\appbrk{}रञ्च}}\Afootnote{H; \skt{सुख\appbrk{}चे\appbrk{}ष्टा\appbrk{}प्र\appbrk{}चारं च}\ A}}\skt{ }\edtext{\skt{सम्य\-क्सा\-धि\-त\-मा\-दि\-शे\-दिति}}{\lemma{\skt{सम्य\appbrk{}क्सा\appbrk{}धि\appbrk{}त\appbrk{}मा\appbrk{}दि\appbrk{}शे\appbrk{}दिति}}\Afootnote{H; \skt{संहितं सम्यगादिशेत्}\ A}}\skt{~॥}\pend

\pstart \ledsidenote{\passref{4.3.trailer}}\edtext{\skt{चिकित्सायां}}{\lemma{\skt{चिकित्सायां}}\Afootnote{H; \skt{इति सुश्रु\appbrk{}त\appbrk{}सं\appbrk{}हि\appbrk{}तायां चिकित्सास्थाने भग्नचिकित्सितं}\ A}}\skt{ }\edtext{\skt{तृतीयोऽध्यायः}}{\lemma{\skt{तृतीयो}}\Afootnote{H; \skt{नाम तृतीयो}\ A}}\skt{~॥~॥}\pend

\endnumbering
\beginnumbering
\eledsection*{Adhyāya 4}
\pstart \ledsidenote{\passref{4.4.1}}\skt{अथातो वात\-व्या\-धि\-चि\-कि\-त्सितं व्याख्यास्यामः~॥}\pend

\pstart\versepar \ledsidenote{\passref{4.4.3}}\skt{आमाशयगते }\edtext{\skt{वायौ}}{\lemma{\skt{वायौ}}\Afootnote{H; \skt{वाते}\ A}}\skt{ }\edtext{\skt{च्छर्द्दिताय}}{\lemma{\skt{च्छर्द्दिताय}}\Afootnote{H; \skt{च्छर्दयित्वा}\ A}}\skt{ यथाक्रमं~। }\pend
\pstart\versepar \skt{देयः षड्धरणो योगः सप्तरात्रं सुखाम्बुना~॥~॥}\pend

\pstart\versepar \ledsidenote{\passref{4.4.4}}\edtext{\skt{चित्रकेन्द्रयवे}}{\lemma{\skt{चित्रकेन्द्रयवे}}\Afootnote{A; \skt{चित्रकेन्द्रयवा}\ H}}\skt{ पाठा }\edtext{\skt{कटू\-का\-ति\-वि\-षा\-भया}}{\lemma{\skt{कटू\appbrk{}का\appbrk{}ति\appbrk{}वि\appbrk{}षा\appbrk{}भया}}\Afootnote{H; \skt{कटु\appbrk{}काऽ\appbrk{}ति\appbrk{}वि\appbrk{}षाऽ\appbrk{}भया}\ A}}\skt{~।}\footnoteA{We emend \emph{citrakendrayavā} to the dual, with the vulgate, to make the sentence construe.}\skt{ }\pend
\pstart\versepar \skt{}\edtext{\skt{महा\-व्या\-धि\-प्र\-श\-मनो}}{\lemma{\skt{महा\appbrk{}व्या\appbrk{}धि\appbrk{}प्र\appbrk{}श\appbrk{}मनो}}\Afootnote{H; \skt{वात\appbrk{}व्या\appbrk{}धि\appbrk{}प्र\appbrk{}श\appbrk{}मनो}\ A}}\skt{ योगः षड्धरणः स्मृतः~॥}\pend

\pstart\versepar \ledsidenote{\passref{4.4.5}}\skt{पक्वाशयगते चापि देयं स्नेहविरेचनं~। }\pend
\pstart\versepar \skt{बस्तयः शोधनीयाश्च प्राशाश्च लवणोत्तराः~॥}\pend

\pstart\versepar \ledsidenote{\passref{4.4.6}}\skt{कार्यो बस्तिगते चापि }\edtext{\skt{विधिर्बस्तिश्च}}{\lemma{\skt{बस्तिश्च}}\Afootnote{H; \skt{बस्तिवि}\ A}}\skt{ शोधनः~। }\pend
\pstart\versepar \skt{}\edtext{\skt{श्रोत्रादिषु}}{\lemma{\skt{श्रोत्रादिषु}}\Afootnote{A; \skt{श्रोतादिषु}\ H}}\skt{ प्रकुपिते कार्यश्चानिलहा क्रमः~॥}\footnoteA{We emend \emph{śrotādi-} to \emph{śrotrādi-} as attested in the vulgate.}\pend

\pstart\versepar \ledsidenote{\passref{4.4.7}}\edtext{\skt{स्नेहा\-भ्य\-ङ्गो\-प\-ना\-हांश्च}}{\lemma{\skt{स्नेहा\appbrk{}भ्य\appbrk{}ङ्गो\appbrk{}प\appbrk{}ना\appbrk{}हांश्च}}\Afootnote{H; \skt{स्नेहा\appbrk{}भ्य\appbrk{}ङ्गो\appbrk{}प\appbrk{}ना\appbrk{}हाश्च}\ A}}\skt{ मर्द्द\-ना\-ले\-प\-नानि च~। }\pend
\pstart\versepar \skt{}\edtext{\skt{त्वङ्मांसरुधिरे}}{\lemma{\skt{त्वङ्मांसरुधिरे}}\Afootnote{H; \skt{त्वङ्मांसासृक्सिरा}\ A}}\skt{ }\edtext{\skt{प्राप्ते}}{\lemma{\skt{प्राप्ते}}\Afootnote{A H\textsuperscript{pc}; \skt{प्रोक्ते}\ H\textsuperscript{ac}}}\skt{ कुर्या\-च्चा\-सृ\-ग्वि\-मो\-क्ष\-णम्~॥}\pend

\pstart\versepar \ledsidenote{\passref{4.4.8}}\skt{स्नेहो\-प\-ना\-हा\-ग्नि\-क\-र्म्म\-व\-न्ध\-नो\-न्म\-र्द्द\-नानि च~। }\pend
\pstart\versepar \skt{स्नायु\-स\-न्ध्य\-स्थि\-स\-म्प्राप्ते }\edtext{\skt{कुर्याद्वाते}}{\lemma{\skt{वाते}}\Afootnote{H; \skt{वाया}\ A}}\skt{ }\edtext{\skt{विचक्षणः}}{\lemma{\skt{विचक्षणः}}\Afootnote{H; \skt{वतन्द्रितः}\ A}}\skt{~॥}\pend

\pstart\versepar \ledsidenote{\passref{4.4.9}}\edtext{\edtext{\skt{निगूढेऽस्थिगते}}{\lemma{\skt{निगूढे}}\Afootnote{H; \skt{निरुद्धे}\ A}}}{\lemma{\skt{ऽस्थिगते}}\Afootnote{H; \skt{ऽस्थनि वा}\ A}}\skt{ वायौ पाणिमन्थेन दारिते~। }\pend
\pstart\versepar \skt{}\edtext{\skt{नाडीन्दत्वास्थनि}}{\lemma{\skt{दत्वास्थनि}}\Afootnote{\skt{दत्त्वाऽस्थनि}\ A; \skt{दत्वास्थीनि}\ H}}\skt{ }\edtext{\skt{भिष\-गा\-चू\-षे\-त्प\-व\-न\-म्बली}}{\lemma{\skt{भिषगाचूषेत्}}\Afootnote{H; \skt{भिषक्चूषयेत्}\ A}}\skt{~॥}\footnoteA{We emend \emph{asthīni} to the locative \emph{asthani} with the vulgate, to make the sentence construe.}\pend

\pstart\versepar \ledsidenote{\passref{4.4.10}}\skt{शुक्रप्राप्तेऽनिले }\edtext{\edtext{\skt{कुर्या\-च्छु\-क्र\-दो\-ष\-चि\-कि\-त्सि\-तम्}}{\lemma{\skt{कुर्याच्}}\Afootnote{H; \skt{कार्यं}\ A}}}{\lemma{\skt{छुक्र\appbrk{}दो\appbrk{}ष\appbrk{}चि\appbrk{}कि\appbrk{}त्सि\appbrk{}तम्}}\Afootnote{H; \skt{शुक्र\appbrk{}दो\appbrk{}ष\appbrk{}चि\appbrk{}कि\appbrk{}त्सि\appbrk{}तम्}\ A}}\skt{~। }\pend
\pstart\versepar \skt{}\edtext{\skt{अव\-गा\-ह\-कु\-टी\-क\-र्षू\-प्र\-स्त\-रा\-भ्य\-ङ्ग\-ब\-स्तिभिः}}{\lemma{\skt{अव\appbrk{}गा\appbrk{}ह\appbrk{}कु\appbrk{}टी\appbrk{}क\appbrk{}र्षू\appbrk{}प्र\appbrk{}स्त\appbrk{}रा\appbrk{}भ्य\appbrk{}ङ्ग\appbrk{}ब\appbrk{}स्तिभिः}}\Afootnote{A; \skt{अव\appbrk{}गा\appbrk{}ह\appbrk{}कु\appbrk{}टी\appbrk{}क\appbrk{}र्ष\appbrk{}प्र\appbrk{}स्त\appbrk{}रा\appbrk{}भ्य\appbrk{}ङ्ग\appbrk{}ब\appbrk{}स्तिनि}\ H}}\skt{~॥}\footnoteA{We emend \emph{karṣa} to \emph{karṣū} with the vulgate, and we emend \emph{-bastini} to \emph{bastibhiḥ}, again with the vulgate, for a meaningful text. The latter error may reflect a scribal misreading of \emph{ni} for \emph{bhi(ḥ)} at an earlier stage of transmission, since \emph{na} and \emph{bha} are similar in old Nepalese scripts.}\pend

\pstart\versepar \ledsidenote{\passref{4.4.11}}\edtext{\skt{जयेत्सर्व्वाङ्गगं}}{\lemma{\skt{सर्व्वाङ्गगं}}\Afootnote{H; \skt{सर्वाङ्गजं}\ A}}\skt{ वातं सिरामोक्षैश्च बुद्धिमान्~। }\pend
\pstart\versepar \skt{एकाङ्गगम्वा }\edtext{\edtext{\skt{मति\-मा\-ञ्शृ\-ङ्गै\-श्चा\-व\-स्थि\-त\-ञ्ज\-येत्}}{\lemma{\skt{मतिमाञ्}}\Afootnote{\skt{मतिमाञ्}\ A; \skt{मतिमान्}\ H}}}{\lemma{\skt{शृङ्गैश्}}\Afootnote{H; \skt{छृङ्गैश्}\ A}}\skt{~॥}\edtext{}{\lemma{\skt{बुद्धिमान्~। एकाङ्गगम्वा}}\Afootnote{H; \skt{बुद्धिमान्~। एकाङ्गगं च}\ A}}\pend

\pstart\versepar \ledsidenote{\passref{4.4.12}}\edtext{\skt{बला\-स\-पि\-त्त\-र\-क्तैश्च}}{\lemma{\skt{बला\appbrk{}स\appbrk{}पि\appbrk{}त्त\appbrk{}र\appbrk{}क्तैश्च}}\Afootnote{H; \skt{बला\appbrk{}स\appbrk{}पि\appbrk{}त्त\appbrk{}र\appbrk{}क्तैस्तु}\ A}}\skt{ संसृ\-ष्ट\-म\-वि\-रो\-धिभिः~। }\pend
\pstart\versepar \skt{}\edtext{\skt{सुप्तवाते}}{\lemma{\skt{सुप्तवाते}}\Afootnote{H; \skt{सुप्तिवाते}\ A}}\skt{ त्वसृङ्मोक्षं }\edtext{\skt{कुर्याच्च}}{\lemma{\skt{कुर्याच्च}}\Afootnote{H; \skt{कुर्यात्तु}\ A}}\skt{ बहुशो भिषक्~॥}\pend

\pstart\versepar \ledsidenote{\passref{4.4.13}}\edtext{\skt{लिह्याच्च}}{\lemma{\skt{लिह्याच्च}}\Afootnote{H; \skt{दिह्याच्च}\ A}}\skt{ }\edtext{\skt{लव\-णा\-गा\-र\-धू\-मै\-स्तै\-ल\-स\-मा\-युतैः}}{\lemma{\skt{तैलसमायुतैः}}\Afootnote{H; \skt{तैलसमन्वितैः}\ A}}\skt{~। }\pend
\pstart\versepar \skt{}\edtext{\skt{पञ्चमूलीशृतं}}{\lemma{\skt{पञ्चमूलीशृतं}}\Afootnote{A; \skt{पञ्चमूलीश्रितं}\ H}}\skt{ क्षीरं फलाम्लो रस एव च~॥}\pend

\pstart\versepar \ledsidenote{\passref{4.4.14}}\skt{सुस्निग्धो }\edtext{\skt{धान्ययूषो वा}}{\lemma{\skt{धान्ययूषो वा}}\Afootnote{H; \skt{धन्ययूषो वा}\ A}}\skt{ }\edtext{\edtext{\skt{वातनुद्भोजनं}}{\lemma{\skt{वातनुद्}}\Afootnote{H; \skt{हितो वात}\ A}}}{\lemma{\skt{भोजनं}}\Afootnote{H; \emph{om.}\ A}}\skt{ हितं~। }\pend
\pstart\versepar \skt{काकोल्यादिः सवातघ्नः सर्व्वा\-म्ल\-द्र\-व्य\-सं\-युतः~॥।}\footnoteA{We emend to \emph{kākolyādiḥ} with the vulgate and against witness H.}\edtext{}{\lemma{\skt{हितं~। काकोल्यादिः}}\Afootnote{\skt{विकारिणाम्~। काकोल्यादिः}\ A; \skt{हितं~। काकोल्यादि}\ H}}\pend

\pstart\versepar \ledsidenote{\passref{4.4.15}}\edtext{\skt{सानू\-पो\-द\-क\-मां\-सस्तु}}{\lemma{\skt{सानू\appbrk{}पो\appbrk{}द\appbrk{}क\appbrk{}मां\appbrk{}सस्तु}}\Afootnote{H; \skt{सानू\appbrk{}पौ\appbrk{}द\appbrk{}क\appbrk{}मां\appbrk{}सस्तु}\ A}}\skt{ }\edtext{\skt{सर्व्व\-स्ने\-ह\-स\-मा\-युतः}}{\lemma{\skt{सर्व्व\appbrk{}स्ने\appbrk{}ह\appbrk{}स\appbrk{}मा\appbrk{}युतः}}\Afootnote{\skt{सर्व\appbrk{}स्ने\appbrk{}ह\appbrk{}स\appbrk{}म\appbrk{}न्वितः}\ A; \skt{सर्व्व\appbrk{}स्ने\appbrk{}ह\appbrk{}स\appbrk{}मा\appbrk{}युतं}\ H}}\skt{~।}\footnoteA{We emend to \emph{-samāyutaḥ} with the vulgate and against witness H.}\skt{ }\pend
\pstart\versepar \skt{सुखोष्णः स्पष्टलवणः }\edtext{\skt{शाल्वलः}}{\lemma{\skt{शाल्वलः}}\Afootnote{H; \skt{साल्वणः}\ A}}\skt{ परिकीर्त्तितः~॥}\footnoteA{In witness H, the two \emph{la} characters in \emph{salvalaḥ} are differently written.}\pend

\pstart\versepar \ledsidenote{\passref{4.4.16}}\skt{तेनोपनाहं कुर्व्वीत सर्व्वदा वातरोगिणां~। }\pend
\pstart\versepar \skt{कुञ्च्यमानं रुजार्त्तम्वा }\edtext{\skt{गात्र\-स्त\-ब्ध\-म\-थापि}}{\lemma{\skt{गात्रस्तब्धम्}}\Afootnote{H; \skt{गात्रं स्तब्धम्}\ A}}\skt{ वा~॥}\pend

\pstart\versepar \ledsidenote{\passref{4.4.17}}\edtext{\skt{गाढम्पट्टैस्तु}}{\lemma{\skt{पट्टैस्तु}}\Afootnote{H; \skt{पट्टैर्नि}\ A}}\skt{ }\edtext{\skt{बध्नी\-या\-द्व\-ल्क\-क\-र्प्पा\-स\-को\-र्ण्णकैः}}{\lemma{\skt{वल्क\appbrk{}क\appbrk{}र्प्पा\appbrk{}स\appbrk{}को\appbrk{}र्ण्णकैः}}\Afootnote{H; \skt{क्षौम\appbrk{}का\appbrk{}र्पा\appbrk{}सि\appbrk{}कौ\appbrk{}र्णिकैः}\ A}}\skt{~। }\pend
\pstart\versepar \skt{}\edtext{\skt{बिडा\-ल\-न\-कु\-लो\-द्राणां}}{\lemma{\skt{बिडा\appbrk{}ल\appbrk{}न\appbrk{}कु\appbrk{}लो\appbrk{}द्राणां}}\Afootnote{H; \skt{बिडा\appbrk{}ल\appbrk{}न\appbrk{}कु\appbrk{}लो\appbrk{}न्द्राणां}\ A}}\skt{ चर्म्म\-गो\-ण्या\-म्मृ\-गस्य वा~॥}\pend

\pstart\versepar \ledsidenote{\passref{4.4.18}}\skt{प्रवेशयेद्वा स्वभ्यक्तं }\edtext{\skt{शाल्व\-ले\-नो\-प\-ना\-हितं}}{\lemma{\skt{शाल्व\appbrk{}ले\appbrk{}नो\appbrk{}प\appbrk{}ना\appbrk{}हितं}}\Afootnote{H; \skt{साल्व\appbrk{}णे\appbrk{}नो\appbrk{}प\appbrk{}ना\appbrk{}हि\appbrk{}तम्}\ A}}\skt{~। }\pend
\pstart\versepar \skt{}\edtext{\skt{वक्ष\-स्त्रि\-क\-स्क\-न्ध\-ग\-त\-म्वा\-यु\-म्म\-न्या\-ग\-त\-न्तथा}}{\lemma{\skt{वक्ष\appbrk{}स्त्रि\appbrk{}क\appbrk{}स्क\appbrk{}न्ध\appbrk{}ग\appbrk{}तम्}}\Afootnote{\skt{स्कन्ध\appbrk{}व\appbrk{}क्ष\appbrk{}स्त्रि\appbrk{}क\appbrk{}प्राप्तं}\ A; \skt{वक्ष\appbrk{}स्त्रि\appbrk{}क\appbrk{}स्क\appbrk{}न्धि\appbrk{}ग\appbrk{}तम्}\ H}}\skt{~॥}\pend

\pstart\versepar \ledsidenote{\passref{4.4.19}}\skt{वमनं हन्ति नस्यञ्च कुशलेन प्रयोजितं~। }\pend
\pstart\versepar \skt{शिरोगतं शिरोबस्तिर्हन्ति }\edtext{\skt{चासृग्विमोक्षणं}}{\lemma{\skt{चासृग्विमोक्षणं}}\Afootnote{H; \skt{वासृ\appbrk{}ग्वि\appbrk{}मो\appbrk{}क्ष\appbrk{}णम्}\ A}}\skt{~॥}\pend

\pstart\versepar \ledsidenote{\passref{4.4.20}}\skt{स्नेहं मात्रासहस्रन्तु धारयेत्तत्र योगतः~। }\pend
\pstart\versepar \skt{}\edtext{\edtext{\edtext{\skt{सर्व्वा\-ङ्ग\-ग\-त\-मे\-का\-ङ्ग\-स्थि\-त\-म्वापि}}{\lemma{\skt{सर्व्वाङ्गगतम्}}\Afootnote{H; \skt{सर्वाङ्गत}\ A}}}{\lemma{\skt{एकाङ्गस्थितम्}}\Afootnote{H; \skt{मेकाङ्गस्थितं}\ A}}}{\lemma{\skt{वापि}}\Afootnote{H; \skt{वाऽपि}\ A}}\skt{ समीरणम्~॥}\pend

\pstart\versepar \ledsidenote{\passref{4.4.21}}\skt{रुणद्धि केवलो }\edtext{\skt{बस्तिस्तस्य}}{\lemma{\skt{तस्य}}\Afootnote{H; \skt{वायु}\ A}}\skt{ }\edtext{\skt{वेगमिवानिलः}}{\lemma{\skt{इवानिलः}}\Afootnote{H; \skt{इवाचलः}\ A}}\skt{~। }\pend
\pstart\versepar \skt{}\edtext{\skt{स्नेहा}}{\lemma{\skt{स्नेहा}}\Afootnote{H; \skt{स्नेह}\ A}}\skt{ }\edtext{\skt{स्वेदस्तथाभ्यङ्गो}}{\lemma{\skt{तथाभ्यङ्गो}}\Afootnote{H; \skt{तथाऽभ्यङ्गो}\ A}}\skt{ }\edtext{\skt{बस्ति\-स्ने\-ह\-वि\-रे\-चनं}}{\lemma{\skt{बस्ति\appbrk{}स्ने\appbrk{}ह\appbrk{}वि\appbrk{}रे\appbrk{}चनं}}\Afootnote{H; \skt{बस्तिः स्नेह\appbrk{}वि\appbrk{}रे\appbrk{}च\appbrk{}नम्}\ A}}\skt{~॥}\pend

\pstart\versepar \ledsidenote{\passref{4.4.22}}\skt{शिरोबस्तिः शिरः स्नेहो }\edtext{\skt{धूमं}}{\lemma{\skt{धूमं}}\Afootnote{H; \skt{धूमः}\ A}}\skt{ }\edtext{\skt{स्नैहिकमेव}}{\lemma{\skt{स्नैहिकम्}}\Afootnote{H; \skt{स्नैहिक}\ A}}\skt{ च~। }\pend
\pstart\versepar \skt{सुखोष्णः }\edtext{\skt{कवलो}}{\lemma{\skt{कवलो}}\Afootnote{H; \skt{स्नेहगण्डूषो}\ A}}\skt{ नस्यं }\edtext{\skt{कल्कं}}{\lemma{\skt{कल्कं}}\Afootnote{H; \emph{om.}\ A}}\skt{ स्नैहिकमेव च~॥}\pend

\pstart\versepar \ledsidenote{\passref{4.4.23}}\edtext{\skt{क्षीराणि}}{\lemma{\skt{क्षीराणि}}\Afootnote{H; \skt{रसाः क्षीराणि}\ A}}\skt{ मांसानि }\edtext{\skt{रसाः}}{\lemma{\skt{रसाः}}\Afootnote{H; \emph{om.}\ A}}\skt{ }\edtext{\edtext{\skt{स्नेहान्वितञ्चअ}}{\lemma{\skt{स्नेहान्वितञ्}}\Afootnote{\skt{स्नेहाः स्नेहान्वितं}\ A H}}\skt{ यत्}}{\lemma{\skt{चअ यत्}}\Afootnote{\skt{च यत्}\ A H}}\skt{~। }\pend
\pstart\versepar \skt{भोजनानि फलाम्लानि स्निग्धानि लवणानि च~॥}\pend

\pstart\versepar \ledsidenote{\passref{4.4.24}}\skt{सुखोष्णाश्च }\edtext{\skt{परिसेकास्तथा}}{\lemma{\skt{परिसेकास्}}\Afootnote{H; \skt{परीषेकास्}\ A}}\skt{ सम्वाहनानि च~। }\pend
\pstart\versepar \skt{कुङ्कु\-मा\-गु\-रु\-प\-त्राणि कुष्ठै\-ला\-त\-ग\-राणि च~॥}\pend

\pstart\versepar \ledsidenote{\passref{4.4.25}}\edtext{\skt{कौषे\-यो\-र्ण्ण\-क\-रो\-माणि}}{\lemma{\skt{कौषे\appbrk{}यो\appbrk{}र्ण्ण\appbrk{}क\appbrk{}रो\appbrk{}माणि}}\Afootnote{H; \skt{कौशे\appbrk{}यौ\appbrk{}र्णि\appbrk{}क\appbrk{}रौ\appbrk{}माणि}\ A}}\skt{ कार्पासानि }\edtext{\skt{मृदूनि च}}{\lemma{\skt{मृदूनि च}}\Afootnote{H; \skt{गुरूणि च}\ A}}\skt{~। }\pend
\pstart\versepar \skt{}\edtext{\skt{निवा\-ता\-त\-प\-यु\-क्तानि}}{\lemma{\skt{निवा\appbrk{}ता\appbrk{}त\appbrk{}प\appbrk{}यु\appbrk{}क्तानि}}\Afootnote{A H\textsuperscript{pc}; \skt{निवा\appbrk{}ता\appbrk{}त\appbrk{}प\appbrk{}युक्ता}\ H\textsuperscript{ac}}}\skt{ तथा }\edtext{\skt{गर्ब्भगृहाणि}}{\lemma{\skt{गर्ब्भगृहाणि}}\Afootnote{H; \skt{गर्भगृहाणि}\ A}}\skt{ च~॥}\pend

\pstart\versepar \ledsidenote{\passref{4.4.26}}\skt{मृद्वी }\edtext{\skt{शय्याग्निसन्तापो}}{\lemma{\skt{शय्याग्निसन्तापो}}\Afootnote{H; \skt{शय्याऽग्निसंतापो}\ A}}\skt{ ब्रह्म\-च\-र्य\-न्त\-थैव च~। }\pend
\pstart\versepar \skt{}\edtext{\edtext{\skt{समा\-से\-नै\-व\-मा\-दीनि}}{\lemma{\skt{समासेनैवम्}}\Afootnote{H; \skt{समासेनैव}\ A}}}{\lemma{\skt{आदीनि}}\Afootnote{H; \skt{मादीनि}\ A}}\skt{ योज्या\-न्य\-नि\-ल\-रो\-गिषु~॥}\pend

\pstart \ledsidenote{\passref{4.4.27}}\edtext{\edtext{\skt{त्रिवृ\-द्द\-न्ती\-श\-ङ्खि\-नी\-सु\-व\-र्ण्ण\-क्षी\-री\-त्रि\-फ\-ला\-वि\-ड\-ङ्गा\-ना\-म\-क्ष\-समाः}}{\lemma{\skt{त्रिवृ\appbrk{}द्द\appbrk{}न्ती\appbrk{}श\appbrk{}ङ्खि\appbrk{}नी\appbrk{}सु\appbrk{}व\appbrk{}र्ण्ण\appbrk{}क्षी\appbrk{}री\appbrk{}त्रि\appbrk{}फ\appbrk{}ला\appbrk{}वि\appbrk{}ड\appbrk{}ङ्गा\appbrk{}नाम्}}\Afootnote{\skt{त्रिवृ\appbrk{}द्द\appbrk{}न्ती\appbrk{}सु\appbrk{}व\appbrk{}र्ण\appbrk{}क्षी\appbrk{}री\appbrk{}स\appbrk{}प्त\appbrk{}ला\appbrk{}श\appbrk{}ङ्खि\appbrk{}नी\appbrk{}त्रि\appbrk{}फ\appbrk{}ला\appbrk{}वि\appbrk{}ड\appbrk{}ङ्गा\appbrk{}नाम्}\ A; \skt{तृवृ\appbrk{}द्द\appbrk{}न्ती\appbrk{}श\appbrk{}ङ्खि\appbrk{}नी\appbrk{}सु\appbrk{}व\appbrk{}र्ण्ण\appbrk{}क्षी\appbrk{}री\appbrk{}त्रि\appbrk{}फ\appbrk{}ला\appbrk{}वि\appbrk{}ड\appbrk{}ङ्गा\appbrk{}नाम्}\ H}}}{\lemma{\skt{अक्षसमाः}}\Afootnote{H; \skt{अक्षसमाः}\ A}}\skt{ }\edtext{\skt{कल्का}}{\lemma{\skt{कल्का}}\Afootnote{H; \skt{भागाः}\ A}}\skt{ }\edtext{\skt{बिल्व\-मा\-त्र\-स्ति\-ल्व\-क\-मू\-ल\-क\-म्पि\-ल्ल्य\-यो\-स्त्रि\-फ\-ला\-र\-स\-द\-धि\-पात्रे}}{\lemma{\skt{तिल्व\appbrk{}क\appbrk{}मू\appbrk{}ल\appbrk{}क\appbrk{}म्पि\appbrk{}ल्ल्य\appbrk{}योस्}}\Afootnote{H; \skt{कल्क\appbrk{}स्ति\appbrk{}ल्व\appbrk{}क\appbrk{}मू\appbrk{}ल\appbrk{}क\appbrk{}म्पि\appbrk{}ल्ल\appbrk{}कयोः}\ A}}\skt{ द्वे द्वे }\edtext{\skt{घृत\-पा\-त्र\-मे\-क\-न्त\-दै\-क\-ध्य\-म्प्र\-ति\-सं\-सृज्य}}{\lemma{\skt{प्रतिसंसृज्य}}\Afootnote{H; \skt{संसृज्य}\ A}}\skt{ }\edtext{\skt{विप\-चे\-त्ति\-ल्व\-क\-स\-र्प्पि\-रे\-त\-त्स्ने\-ह\-वि\-रे\-च\-न\-मु\-प\-दि\-शन्ति}}{\lemma{\skt{स्नेह\appbrk{}वि\appbrk{}रे\appbrk{}च\appbrk{}नम्}}\Afootnote{A H\textsuperscript{pc}; \skt{स्नेह\appbrk{}वि\appbrk{}रे\appbrk{}चि\appbrk{}नम्}\ H\textsuperscript{ac}}}\skt{ }\edtext{\skt{वातरोगेषु}}{\lemma{\skt{वातरोगेषु}}\Afootnote{H; \skt{वातरोगिषु}\ A}}\skt{~। }\edtext{\edtext{\skt{तिल्व\-क\-वि\-धि\-र्चा\-शो\-क\-र\-म्य\-क\-यो\-र्द्र\-ष्टव्यः}}{\lemma{\skt{तिल्वकविधिर्}}\Afootnote{A; \skt{तिल्वकविधि}\ H}}}{\lemma{\skt{चाशो\appbrk{}क\appbrk{}र\appbrk{}म्य\appbrk{}क\appbrk{}योर्}}\Afootnote{\skt{एवा\appbrk{}शो\appbrk{}क\appbrk{}र\appbrk{}म्य\appbrk{}क\appbrk{}योर्}\ A; \skt{ना चाशो\appbrk{}क\appbrk{}र\appbrk{}म्य\appbrk{}क\appbrk{}योर्}\ H}}\skt{~॥}\footnoteA{We read \emph{tilvakavidhir}with the vulgate for grammar and sense. The syntax of witness H does not have grammatical concord.}\pend

\pstart \ledsidenote{\passref{4.4.28}}\edtext{\edtext{\skt{तिल\-पी\-ड\-नो\-प\-क\-र\-ण\-का\-ष्ठा\-न्या\-हृ\-त्या\-न\-ल्प\-का\-ल\-तै\-ल\-प\-रि\-पी\-डि\-तानि}}{\lemma{\skt{तिल\appbrk{}पी\appbrk{}ड\appbrk{}नो\appbrk{}प\appbrk{}क\appbrk{}र\appbrk{}ण\appbrk{}का\appbrk{}ष्ठान्य्}}\Afootnote{H; \skt{तिल\appbrk{}प\appbrk{}रि\appbrk{}पी\appbrk{}ड\appbrk{}नो\appbrk{}प\appbrk{}क\appbrk{}र\appbrk{}ण\appbrk{}का\appbrk{}ष्ठान्}\ A}}}{\lemma{\skt{आहृ\appbrk{}त्या\appbrk{}न\appbrk{}ल्प\appbrk{}का\appbrk{}ल\appbrk{}तै\appbrk{}ल\appbrk{}प\appbrk{}रि\appbrk{}पी\appbrk{}डि\appbrk{}तानि}}\Afootnote{H; \skt{याहृ\appbrk{}त्या\appbrk{}न\appbrk{}ल्प\appbrk{}कालं तैल\appbrk{}प\appbrk{}रि\appbrk{}पी\appbrk{}तान्य्}\ A}}\skt{ अणूनि }\edtext{\skt{च्छेद\-यि\-त्वा\-व\-क्षुद्य}}{\lemma{\skt{च्छेद\appbrk{}यि\appbrk{}त्वा\appbrk{}व\appbrk{}क्षुद्य}}\Afootnote{H; \skt{खण्डशः कल्प\appbrk{}यि\appbrk{}त्वाऽ\appbrk{}व\appbrk{}क्षुद्}\ A}}\skt{ }\edtext{\skt{प्रक्षिप्य}}{\lemma{\skt{प्रक्षिप्य}}\Afootnote{H; \skt{य}\ A}}\skt{ महति कटाहे }\edtext{\skt{पानीयेनाप्लाव्य}}{\lemma{\skt{पानीयेनाप्लाव्य}}\Afootnote{H; \skt{पानी\appbrk{}ये\appbrk{}ना\appbrk{}भि\appbrk{}प्लाव्य}\ A}}\skt{ क्वाथयेत्~। }\edtext{\skt{ततस्तं}}{\lemma{\skt{ततस्तं}}\Afootnote{H; \skt{ततः}\ A}}\skt{ }\edtext{\skt{स्नेह\-म\-म्बु\-पृ\-ष्था\-त्स\-र\-क\-पा\-ण्यो\-र\-न्य\-त\-रे\-णा\-दाय}}{\lemma{\skt{अम्बुपृष्थात्}}\Afootnote{H; \skt{अम्बु\appbrk{}पृ\appbrk{}ष्ठा\appbrk{}द्य\appbrk{}दु\appbrk{}देति तत्}\ A}}\skt{ }\edtext{\skt{वात\-घ्नौ\-ष\-ध\-द\-त्त\-प्र\-ती\-वापं}}{\lemma{\skt{वात\appbrk{}घ्नौ\appbrk{}ष\appbrk{}ध\appbrk{}द\appbrk{}त्त\appbrk{}प्र\appbrk{}ती\appbrk{}वापं}}\Afootnote{\skt{वात\appbrk{}घ्नौ\appbrk{}ष\appbrk{}ध\appbrk{}प्र\appbrk{}ती\appbrk{}वापं}\ A; \skt{वात\appbrk{}घ्नौ\appbrk{}ष\appbrk{}ध\appbrk{}द\appbrk{}न्त\appbrk{}प्र\appbrk{}ती\appbrk{}वायं}\ H}}\skt{ स्नेह\-पा\-क\-क\-ल्पेन }\edtext{\skt{विप\-चे\-दे\-त\-द\-णु\-तै\-ल\-मु\-प\-दि\-शन्ति}}{\lemma{\skt{विप\appbrk{}चे\appbrk{}दे\appbrk{}त\appbrk{}द\appbrk{}णु\appbrk{}तै\appbrk{}लम्}}\Afootnote{A; \skt{विप\appbrk{}चे\appbrk{}दे\appbrk{}त\appbrk{}द\appbrk{}नु\appbrk{}तै\appbrk{}लम्}\ H}}\skt{ }\edtext{\skt{वातरोरोगेषु}}{\lemma{\skt{वातरोरोगेषु}}\Afootnote{\skt{वातरोगिषु}\ A; \skt{वातरोगेषु}\ H}}\skt{~।}\footnoteA{We emend \emph{-danta-} in witness H to \emph{datta} since this is the reading of the parallel passage in the Aṣṭāṅgasaṅgraha (cikitsāsthāna 23.18). But the meaning of the compound remains obscure.}\skt{ }\edtext{\skt{अणु\-तै\-ल\-द्र\-व्येभ्यो}}{\lemma{\skt{अणु\appbrk{}तै\appbrk{}ल\appbrk{}द्र\appbrk{}व्येभ्यो}}\Afootnote{\skt{अणु\appbrk{}भ्य\appbrk{}स्तै\appbrk{}ल\appbrk{}द्र\appbrk{}व्येभ्यो}\ A; \skt{अनु\appbrk{}तै\appbrk{}ल\appbrk{}द्र\appbrk{}व्येभ्यो}\ H}}\skt{ निष्पाद्यत }\edtext{\skt{इत्यणुतैलं}}{\lemma{\skt{अणुतैलं}}\Afootnote{A; \skt{अनुतैलं}\ H}}\skt{~।}\pend

\pstart \ledsidenote{\passref{4.4.29}}\edtext{\skt{अथवा}}{\lemma{\skt{अथवा}}\Afootnote{H; \skt{अथ}\ A}}\skt{ }\edtext{\edtext{\edtext{\skt{महा\-प\-ञ्च\-मू\-ली\-का\-ष्ठै\-र्ब्ब\-हु\-भि\-र\-प\-द\-ह्या\-व\-नि\-प्र\-दे\-श\-मे\-क\-रा\-त्र\-मु\-प\-शा\-न्तेऽ\-ग्ना\-व\-पोह्य}}{\lemma{\skt{महा\appbrk{}प\appbrk{}ञ्च\appbrk{}मू\appbrk{}ली\appbrk{}का\appbrk{}ष्ठैर्}}\Afootnote{H; \skt{महा\appbrk{}प\appbrk{}ञ्च\appbrk{}मू\appbrk{}ल\appbrk{}का\appbrk{}ष्ठैर्}\ A}}}{\lemma{\skt{अप\appbrk{}द\appbrk{}ह्या\appbrk{}व\appbrk{}नि\appbrk{}प्र\appbrk{}दे\appbrk{}शम्}}\Afootnote{H\textsuperscript{pc}; \skt{अव\appbrk{}द\appbrk{}ह्या\appbrk{}व\appbrk{}नि\appbrk{}प्र\appbrk{}दे\appbrk{}शम्}\ A; \skt{अप\appbrk{}द\appbrk{}ह्या\appbrk{}व\appbrk{}नि\appbrk{}प्रदे}\ H\textsuperscript{ac}}}}{\lemma{\skt{एकरात्रम्}}\Afootnote{H\textsuperscript{pc}; \skt{असि\appbrk{}त\appbrk{}मु\appbrk{}षि\appbrk{}त\appbrk{}मे\appbrk{}क\appbrk{}रा\appbrk{}त्रम्}\ A; \skt{शेमेकरात्रम्}\ H\textsuperscript{ac}}}\skt{ }\edtext{\skt{भस्म\-नि\-र्व्वृ\-ता\-म्भू\-मि\-म्वि\-दा\-रि\-ग\-न्धा\-दि\-सि\-द्धेन}}{\lemma{\skt{भस्मनिर्व्वृताम्}}\Afootnote{H; \skt{भस्म निवृत्तां}\ A}}\skt{ }\edtext{\skt{तैल\-घ\-ट\-श\-ते\-ना\-व\-सि\-च्यै\-क\-रा\-त्र\-म\-व\-स्थाप्य}}{\lemma{\skt{तैल\appbrk{}घ\appbrk{}ट\appbrk{}श\appbrk{}ते\appbrk{}ना\appbrk{}व\appbrk{}सि\appbrk{}च्यै\appbrk{}क\appbrk{}रा\appbrk{}त्रम्}}\Afootnote{H; \skt{तैल\appbrk{}घ\appbrk{}ट\appbrk{}श\appbrk{}तेन तुल्य\appbrk{}प\appbrk{}य\appbrk{}साऽ\appbrk{}भि\appbrk{}षि\appbrk{}च्यै\appbrk{}क\appbrk{}रा\appbrk{}त्रम्}\ A}}\skt{ ततो }\edtext{\edtext{\skt{यावान्पांशुः}}{\lemma{\skt{यावान्}}\Afootnote{\skt{यावती}\ A; \skt{यावन्}\ H}}}{\lemma{\skt{पांशुः}}\Afootnote{H; \skt{मृत्तिका}\ A}}\skt{ }\edtext{\skt{स्निग्धः}}{\lemma{\skt{स्निग्धः}}\Afootnote{H; \skt{स्निग्धा}\ A}}\skt{ }\edtext{\edtext{\skt{स्यात्ता\-वा\-न्ग्रा\-ह\-यि\-त्वो\-दके}}{\lemma{\skt{तावान्}}\Afootnote{H; \skt{ताम्}\ A}}}{\lemma{\skt{ग्राहयित्वोदके}}\Afootnote{H; \skt{आदा\appbrk{}यो\appbrk{}ष्णो\appbrk{}द\appbrk{}केन}\ A}}\skt{ महति }\edtext{\skt{कटा\-हेऽ\-भ्या\-सि\-ञ्चेत्}}{\lemma{\skt{ऽभ्यासिञ्चेत्}}\Afootnote{A; \skt{भ्यःसिंचेत्}\ H}}\skt{~। तत्र }\edtext{\edtext{\skt{यत्तै\-ल\-मु\-त्ति\-ष्ठे\-त्त\-त्पा\-णि\-भ्या\-मा\-दाय}}{\lemma{\skt{उत्तिष्ठेत्}}\Afootnote{A; \skt{उत्तिष्ठते}\ H}}}{\lemma{\skt{आदाय}}\Afootnote{H; \skt{पर्यादाय}\ A}}\skt{ स्वनु\-गु\-प्त\-न्नि\-द\-ध्यात्~। }\edtext{\skt{तत्तै\-ल\-म्वा\-त\-ह\-रौ\-ष\-ध\-क्वा\-थ\-मां\-स\-र\-स\-क्षी\-रा\-म्ल\-भा\-ग\-स\-ह\-स्रेण}}{\lemma{\skt{तत्}}\Afootnote{H; \skt{ततस्}\ A}}\skt{ सह\-स्र\-पा\-क\-म्वि\-प\-चेत्~॥ यावता कालेन शक्नु\-या\-त्प\-क्तु\-म्प्र\-ति\-वा\-प\-श्चात्र }\edtext{\skt{हैमवता}}{\lemma{\skt{हैमवता}}\Afootnote{A; \skt{हेमवता}\ H}}\skt{ }\edtext{\skt{दक्षि\-णा\-प\-थ\-गा\-श्चा\-श्व\-गन्धा}}{\lemma{\skt{चाश्वगन्धा}}\Afootnote{H; \skt{च गन्धा}\ A}}\skt{ वातघ्नानि च~। तस्मिन्सिद्ध्यति }\edtext{\edtext{\edtext{\edtext{\skt{शङ्खा\-ना\-ध्मा\-प\-ये\-च्छ\-त्र\-न्धा\-र\-ये\-द्दु\-न्दु\-भी\-स्ता\-ड\-ये\-त्वा\-ल\-व्य\-ज\-नैश्च}}{\lemma{\skt{शङ्खा\appbrk{}ना\appbrk{}ध्मा\appbrk{}प\appbrk{}येच्}}\Afootnote{H; \skt{शङ्खा\appbrk{}ना\appbrk{}ध्मा\appbrk{}प\appbrk{}ये\appbrk{}द्दु\appbrk{}न्दु\appbrk{}भी\appbrk{}ना\appbrk{}घा\appbrk{}त\appbrk{}येच्}\ A}}}{\lemma{\skt{धारयेद्}}\Afootnote{H; \skt{धार}\ A}}}{\lemma{\skt{दुन्दुभीस्}}\Afootnote{H; \emph{om.}\ A}}}{\lemma{\skt{ताडयेत्}}\Afootnote{H; \skt{येद्}\ A}}\skt{ }\edtext{\skt{व्यजेत्}}{\lemma{\skt{व्यजेत्}}\Afootnote{H; \skt{वीज\appbrk{}ये\appbrk{}द्ब्रा\appbrk{}ह्म\appbrk{}ण\appbrk{}स\appbrk{}हस्रं भोजयेत्}\ A}}\skt{~। }\edtext{\skt{ततः}}{\lemma{\skt{ततः}}\Afootnote{H; \skt{तत्}\ A}}\skt{ साधु\-सि\-द्ध\-म\-व\-तार्य }\edtext{\skt{सौवर्णे}}{\lemma{\skt{सौवर्णे}}\Afootnote{A; \skt{सौवर्ण्ण}\ H}}\skt{ राजते }\edtext{\skt{वा}}{\lemma{\skt{वा}}\Afootnote{H; \skt{मृन्मये वा पात्रे स्व}\ A}}\skt{ }\edtext{\skt{कुम्भे}}{\lemma{\skt{कुम्भे}}\Afootnote{H; \skt{नुगु}\ A}}\skt{ }\edtext{\skt{निदध्यात्}}{\lemma{\skt{निदध्यात्}}\Afootnote{H; \skt{प्तं निदध्यात्}\ A}}\skt{~। }\edtext{\skt{तदे\-त\-त्स\-ह\-स्र\-पा\-क\-म\-प्र\-ति\-वा\-र\-वीयं}}{\lemma{\skt{अप्र\appbrk{}ति\appbrk{}वा\appbrk{}र\appbrk{}वीयं}}\Afootnote{\skt{अप्र\appbrk{}ति\appbrk{}वा\appbrk{}र\appbrk{}वीर्यं}\ A; \skt{अप्र\appbrk{}ति\appbrk{}वा\appbrk{}र\appbrk{}वीय}\ H}}\skt{ राजार्हन्तैलम्~। }\edtext{\skt{एव\-म्भा\-ग\-श\-त\-पक्वं}}{\lemma{\skt{भागशतपक्वं}}\Afootnote{\skt{भाग\appbrk{}श\appbrk{}त\appbrk{}वि\appbrk{}पक्वं}\ A; \skt{भागश}\ H\textsuperscript{ac}; \skt{भाग\appbrk{}श\appbrk{}स\appbrk{}ह\appbrk{}स्र\appbrk{}पक्वं}\ H\textsuperscript{pc}}}\skt{ }\edtext{\edtext{\skt{शतपाकमिति}}{\lemma{\skt{शतपाकम्}}\Afootnote{A; \emph{om.}\ H\textsuperscript{ac}; \skt{सहस्रपाकम्}\ H\textsuperscript{pc}}}\skt{ शतं }\edtext{\skt{पाकमिति}}{\lemma{\skt{इति}}\Afootnote{H; \emph{om.}\ A}}}{\lemma{\skt{इति शतं पाकम्}}\Afootnote{\emph{om.}\ A; \skt{तपक्वं शतपाकम्}\ H\textsuperscript{ac}; \skt{इतितपक्वं शतपाकम्}\ H\textsuperscript{pc}}}\skt{~॥}\pend

\pstart \ledsidenote{\passref{4.4.30}}\edtext{\edtext{\skt{गन्ध\-र्व्व\-ह\-स्ता\-ट\-रू\-ष\-क\-मु\-ष्क\-क\-न\-क्त\-मा\-ल\-पू\-ती\-क\-चि\-त्र\-का\-णा\-म्प\-त्त्रा\-ण्या\-हृ\-त्या\-र्द्द्राणि}}{\lemma{\skt{गन्ध\appbrk{}र्व्व\appbrk{}ह\appbrk{}स्ता\appbrk{}ट\appbrk{}रू\appbrk{}ष\appbrk{}क\appbrk{}मु\appbrk{}ष्क\appbrk{}क\appbrk{}न\appbrk{}क्त\appbrk{}मा\appbrk{}ल\appbrk{}पू\appbrk{}ती\appbrk{}क\appbrk{}चि\appbrk{}त्र\appbrk{}का\appbrk{}णाम्}}\Afootnote{\skt{गन्ध\appbrk{}र्व\appbrk{}ह\appbrk{}स्त\appbrk{}मु\appbrk{}ष्क\appbrk{}क\appbrk{}न\appbrk{}क्त\appbrk{}मा\appbrk{}ला\appbrk{}ट\appbrk{}रू\appbrk{}ष\appbrk{}क\appbrk{}पू\appbrk{}ती\appbrk{}का\appbrk{}र\appbrk{}ग्व\appbrk{}ध\appbrk{}चि\appbrk{}त्र\appbrk{}का\appbrk{}दीनां}\ A; \skt{गन्ध\appbrk{}र्व्व\appbrk{}ह\appbrk{}स्ता\appbrk{}ट\appbrk{}रू\appbrk{}ष\appbrk{}क\appbrk{}मु\appbrk{}ष्क\appbrk{}क\appbrk{}न\appbrk{}क्त\appbrk{}मा\appbrk{}ल\appbrk{}पू\appbrk{}ती\appbrk{}क\appbrk{}चि\appbrk{}त्र\appbrk{}का\appbrk{}नाम्}\ H}}}{\lemma{\skt{आहृत्यार्द्द्राणि}}\Afootnote{H; \skt{आर्द्राणि}\ A}}\skt{ लवणेन सहोदूखले }\edtext{\skt{संक्षुद्य}}{\lemma{\skt{संक्षुद्य}}\Afootnote{H; \skt{ऽवक्षुद्य}\ A}}\skt{ स्नेहघटे }\edtext{\skt{प्रक्षिप्यालिप्य}}{\lemma{\skt{प्रक्षिप्यालिप्य}}\Afootnote{H; \skt{प्रक्षिप्यावलिप्य}\ A}}\skt{ गोश\-कृ\-द्भि\-र्द्दा\-ह\-येत्~। एत\-त्प\-त्र\-ल\-व\-ण\-मु\-प\-दि\-शन्ति वातरोगेषु~॥}\pend

\pstart \ledsidenote{\passref{4.4.31}}\skt{एवं }\edtext{\skt{स्नुहा\-का\-ण्ड\-वा\-र्त्ता\-कु\-ल\-व\-णा\-न्सं\-क्षुद्य}}{\lemma{\skt{स्नुहा\appbrk{}का\appbrk{}ण्ड\appbrk{}वा\appbrk{}र्त्ता\appbrk{}कु\appbrk{}ल\appbrk{}व\appbrk{}णान्}}\Afootnote{\skt{स्नुही\appbrk{}का\appbrk{}ण्ड\appbrk{}वा\appbrk{}र्ता\appbrk{}कु\appbrk{}शि\appbrk{}ग्रु\appbrk{}ल\appbrk{}व\appbrk{}णानि}\ A; \skt{स्नुहा\appbrk{}का\appbrk{}ण्ड\appbrk{}वा\appbrk{}र्त्ता\appbrk{}कु\appbrk{}ल\appbrk{}वणा}\ H}}\skt{ }\edtext{\skt{पूर्ण्णघटे}}{\lemma{\skt{पूर्ण्णघटे}}\Afootnote{H; \skt{घटं पूरयित्वा}\ A}}\skt{ }\edtext{\skt{सर्प्पि\-स्तै\-ल\-व\-सा\-म\-ज्जानः}}{\lemma{\skt{सर्प्पि\appbrk{}स्तै\appbrk{}ल\appbrk{}व\appbrk{}सा\appbrk{}म\appbrk{}ज्जानः}}\Afootnote{H; \skt{सर्पि\appbrk{}स्तै\appbrk{}ल\appbrk{}व\appbrk{}सा\appbrk{}म\appbrk{}ज्जभिः}\ A}}\skt{ }\edtext{\skt{प्रक्षिप्यालिप्य}}{\lemma{\skt{प्रक्षिप्यालिप्य}}\Afootnote{H; \skt{प्रक्षिप्यावलिप्य}\ A}}\skt{ }\edtext{\edtext{\skt{पूर्व्ववद्दहेत्}}{\lemma{\skt{पूर्व्ववद्}}\Afootnote{H; \skt{गोशकृद्भिर्}\ A}}\skt{~। एत\-त्स्ने\-ह\-ल\-व\-ण\-मु\-प\-दि\-शन्ति}}{\lemma{\skt{दहेत्~। एतत्}}\Afootnote{H; \skt{दाहयेत्; एतत्}\ A}}\skt{ }\edtext{\skt{वातरोगेषु}}{\lemma{\skt{वातरोगेषु}}\Afootnote{H; \skt{वातरोगेषु~।(इति काण्डलवणम्)}\ A}}\skt{~॥}\pend

\pstart \ledsidenote{\passref{4.4.32}}\edtext{\edtext{\skt{गण्डी\-र\-प\-ला\-श\-कु\-ट\-ज\-बि\-ल्वा\-र्क्क\-स्नु\-हा\-पा\-मा\-र्ग\-पा\-ट\-ली\-पा\-रि\-भ\-द्र\-क\-ना\-दे\-यी\-कृ\-ष्ण\-ग\-न्धा\-नी\-प\-नि\-र्द्द\-ह्या\-न्या\-ट\-रू\-ष\-क\-न\-क्त\-मा\-ल\-पू\-ती\-क\-बृ\-ह\-ती\-क\-ण्ट\-का\-रि\-का\-भ\-ल्ला\-त\-का\-शो\-क\-वै\-ज\-य\-न्ती\-त्ये\-त\-देव}}{\lemma{\skt{गण्डी\appbrk{}र\appbrk{}प\appbrk{}ला\appbrk{}श\appbrk{}कु\appbrk{}ट\appbrk{}ज\appbrk{}बि\appbrk{}ल्वा\appbrk{}र्क्क\appbrk{}स्नु\appbrk{}हा\appbrk{}पा\appbrk{}मा\appbrk{}र्ग\appbrk{}पा\appbrk{}ट\appbrk{}ली\appbrk{}पा\appbrk{}रि\appbrk{}भ\appbrk{}द्र\appbrk{}क\appbrk{}ना\appbrk{}दे\appbrk{}यी\appbrk{}कृ\appbrk{}ष्ण\appbrk{}ग\appbrk{}न्धा\appbrk{}नी\appbrk{}प\appbrk{}नि\appbrk{}र्द्द\appbrk{}ह्या\appbrk{}न्या\appbrk{}ट\appbrk{}रू\appbrk{}ष\appbrk{}क\appbrk{}न\appbrk{}क्त\appbrk{}मा\appbrk{}ल\appbrk{}पू\appbrk{}ती\appbrk{}क\appbrk{}बृ\appbrk{}ह\appbrk{}ती\appbrk{}क\appbrk{}ण्ट\appbrk{}का\appbrk{}रि\appbrk{}का\appbrk{}भ\appbrk{}ल्ला\appbrk{}त\appbrk{}का\appbrk{}शो\appbrk{}क\appbrk{}वै\appbrk{}ज\appbrk{}य\appbrk{}न्तीत्य्}}\Afootnote{H; \skt{गण्डी\appbrk{}र\appbrk{}प\appbrk{}ला\appbrk{}श\appbrk{}कु\appbrk{}ट\appbrk{}ज\appbrk{}बि\appbrk{}ल्वा\appbrk{}र्क\appbrk{}स्नु\appbrk{}ह्य\appbrk{}पा\appbrk{}मा\appbrk{}र्ग\appbrk{}पा\appbrk{}ट\appbrk{}ला\appbrk{}पा\appbrk{}रि\appbrk{}भ\appbrk{}द्र\appbrk{}क\appbrk{}ना\appbrk{}दे\appbrk{}यी\appbrk{}कृ\appbrk{}ष्ण\appbrk{}ग\appbrk{}न्धा\appbrk{}नी\appbrk{}प\appbrk{}नि\appbrk{}म्ब\appbrk{}नि\appbrk{}र्द\appbrk{}ह\appbrk{}न्य\appbrk{}ट\appbrk{}रू\appbrk{}ष\appbrk{}क\appbrk{}न\appbrk{}क्त\appbrk{}मा\appbrk{}ल\appbrk{}क\appbrk{}पू\appbrk{}ति\appbrk{}क\appbrk{}बृ\appbrk{}ह\appbrk{}ती\appbrk{}क\appbrk{}ण्ट\appbrk{}का\appbrk{}रि\appbrk{}का\appbrk{}भ\appbrk{}ल्ला\appbrk{}त\appbrk{}के\appbrk{}ङ्गु\appbrk{}दी\appbrk{}वै\appbrk{}ज\appbrk{}य\appbrk{}न्ती\appbrk{}क\appbrk{}द\appbrk{}ली\appbrk{}बा\appbrk{}ष्प\appbrk{}द्व\appbrk{}ये\appbrk{}क्षु\appbrk{}र\appbrk{}के\appbrk{}न्द्र\appbrk{}वा\appbrk{}रु\appbrk{}णी\appbrk{}श्वे\appbrk{}त\appbrk{}मो\appbrk{}क्ष\appbrk{}का\appbrk{}शोका इत्य्}\ A}}}{\lemma{\skt{एतदेव}}\Afootnote{H; \skt{ए}\ A}}\skt{ }\edtext{\skt{विंशतिवर्गं}}{\lemma{\skt{विंशतिवर्गं}}\Afootnote{\skt{वं वर्गं}\ A; \skt{विंशतिवर्ग्ग}\ H}}\skt{ }\edtext{\skt{सफ\-ल\-मू\-ल\-प\-त्र\-शा\-ख\-मा\-र्द्र\-मा\-हृत्य}}{\lemma{\skt{सफ\appbrk{}ल\appbrk{}मू\appbrk{}ल\appbrk{}प\appbrk{}त्र\appbrk{}शा\appbrk{}खम्}}\Afootnote{\skt{समू\appbrk{}ल\appbrk{}प\appbrk{}त्र\appbrk{}शा\appbrk{}खम्}\ A; \skt{सफ\appbrk{}ल\appbrk{}मू\appbrk{}ल\appbrk{}प\appbrk{}त्र\appbrk{}शा\appbrk{}खाम्}\ H}}\skt{ लवणेन सह संसृज्य पूर्व्ववद्दग्ध्वा क्षारकल्पेन परिस्राव्य विपचेत्~।}\footnoteA{The emendations against H to the accusative \emph{viṃśativargaṃ saphalamūlapatraśākham} seem necessary for the grammar.}\skt{ प्रतिवापश्चात्र }\edtext{\skt{पिप्पल्यादिर्}}{\lemma{\skt{पिप्पल्यादिर्}}\Afootnote{H; \skt{हिङ्ग्वादिभिः पिप्पल्यादिभिर्वा}\ A}}\skt{~। }\edtext{\edtext{\edtext{\skt{एत\-त्क\-ल्या\-ण\-क\-न्नाम}}{\lemma{\skt{एतत्}}\Afootnote{H; \skt{इत्येतत्}\ A}}}{\lemma{\skt{कल्याणकन्}}\Afootnote{H; \skt{कल्याण}\ A}}}{\lemma{\skt{नाम}}\Afootnote{H; \skt{क}\ A}}\skt{ }\edtext{\skt{लव\-ण\-म्वा\-त\-रोगे}}{\lemma{\skt{वातरोगे}}\Afootnote{H; \skt{वातरोगगु}\ A}}\skt{ }\edtext{\edtext{\edtext{\skt{प्लीहा\-ग्नि\-सं\-गा\-जी\-र्ण्णा\-रो\-च\-का\-र्शो\-भि\-रु\-प\-द्रु\-ता\-ना\-मु\-प\-दि\-शन्ति}}{\lemma{\skt{प्लीहा\appbrk{}ग्नि\appbrk{}सं\appbrk{}गा\appbrk{}जी\appbrk{}र्ण्णा\appbrk{}रो\appbrk{}च\appbrk{}का\appbrk{}र्शो\appbrk{}भिर्}}\Afootnote{H; \skt{ल्मप्ली\appbrk{}हा\appbrk{}ग्नि\appbrk{}ष\appbrk{}ङ्गा\appbrk{}जी\appbrk{}र्णा\appbrk{}र्शोऽ\appbrk{}रो\appbrk{}च\appbrk{}का\appbrk{}र्तानां कासादिभिः कृमिभिर्}\ A}}}{\lemma{\skt{उपद्रुतानाम्}}\Afootnote{A; \skt{ऊपद्रुतानाम्}\ H}}}{\lemma{\skt{उपदिशन्ति}}\Afootnote{H; \skt{चोपदिशन्ति}\ A}}\skt{ }\edtext{\skt{पानभोजनेषु}}{\lemma{\skt{पानभोजनेषु}}\Afootnote{\skt{पान\appbrk{}भो\appbrk{}ज\appbrk{}ने\appbrk{}ष्व\appbrk{}पीति}\ A; \skt{पान\appbrk{}भो\appbrk{}ज\appbrk{}ने\appbrk{}ष्विति}\ H}}\skt{~॥}\pend

\pstart \ledsidenote{\passref{4.4.trailer}}\edtext{\skt{इति }\edtext{\skt{वात\-व्या\-धि\-चि\-कि\-त्सि\-त\-ञ्च\-तुर्थः}}{\lemma{\skt{चतुर्थः}}\Afootnote{H; \skt{नाम चतुर्थोऽध्यायः}\ A}}}{\lemma{\skt{इति वात\appbrk{}व्या\appbrk{}धि\appbrk{}चि\appbrk{}कि\appbrk{}त्सि\appbrk{}तञ्}}\Afootnote{H; \skt{इति श्रीसु\appbrk{}श्रु\appbrk{}त\appbrk{}सं\appbrk{}हि\appbrk{}तायां चिकित्सास्थाने वात\appbrk{}व्या\appbrk{}धि\appbrk{}चि\appbrk{}कि\appbrk{}त्सितं}\ A}}\skt{~॥~॥}\pend

\endnumbering
\beginnumbering
\eledsection*{Adhyāya 5}
\pstart \ledsidenote{\passref{4.5.1}}\skt{अथातो महा\-वा\-त\-व्या\-धि\-चि\-कि\-त्सि\-त\-म्व्या\-ख्या\-स्यामः~॥}\pend

\pstart \ledsidenote{\passref{4.5.3}}\skt{द्विविधं }\edtext{\skt{वात\-र\-क्त\-मु\-त्ता\-न\-म\-व\-गाढं}}{\lemma{\skt{वातरक्तम्}}\Afootnote{H; \skt{वातशोणितम्}\ A}}\skt{ चेत्येके भाषन्ते~॥}\footnoteA{Emmerick (1984: 101) pointed out that the view being criticized is Caraka Ci.29.19.}\skt{ }\edtext{\skt{तत्तु न}}{\lemma{\skt{तत्तु न}}\Afootnote{A; \skt{तन्तु न}\ H}}\skt{ सम्यक्~। }\edtext{\skt{कस्मात्}}{\lemma{\skt{कस्मात्}}\Afootnote{H; \skt{तद्धि}\ A}}\skt{~। कुष्ठ\-व\-दु\-त्ता\-न\-म्भूत्वा काला\-न्त\-रे\-णा\-व\-गा\-ढी\-भ\-वति~। }\edtext{\edtext{\skt{तस्माद्द्वित्वसंख्या}}{\lemma{\skt{तस्माद्}}\Afootnote{H; \skt{तस्मान्न}\ A}}}{\lemma{\skt{द्वित्वसंख्या}}\Afootnote{H; \skt{द्विविधम्}\ A}}\skt{ }\edtext{\skt{हीयते}}{\lemma{\skt{हीयते}}\Afootnote{H; \emph{om.}\ A}}\skt{~।}\pend

\pstart \ledsidenote{\passref{4.5.4}}\skt{तत्र }\edtext{\skt{बल\-व\-द्वि\-ग्र\-हा\-दि\-भि\-र्व्वि\-शेषैः}}{\lemma{\skt{व्विशेषैः}}\Afootnote{H; \emph{om.}\ A}}\skt{ प्रकुपितस्य वायो\-र्ग्गु\-रू\-ष्णा\-ध्य\-श\-न\-शी\-लस्य }\edtext{\skt{दुष्ट\-शो\-णि\-त\-म्मा\-र्ग\-मा\-वृत्य}}{\lemma{\skt{दुष्टशोणितम्}}\Afootnote{H; \skt{प्रदुष्टं शोणितं}\ A}}\skt{ वातेन सहैकीभूतं }\edtext{\skt{युग\-प\-द्वा\-त\-शो\-णि\-त\-नि\-मित्तां}}{\lemma{\skt{वात\appbrk{}शो\appbrk{}णि\appbrk{}त\appbrk{}नि\appbrk{}मित्तां}}\Afootnote{H; \skt{वात\appbrk{}र\appbrk{}क्त\appbrk{}नि\appbrk{}मित्तां}\ A}}\skt{ }\edtext{\skt{वेद\-ना\-ञ्ज\-न\-यति}}{\lemma{\skt{जनयति}}\Afootnote{H; \skt{जनयती}\ A}}\skt{~। }\edtext{\edtext{\edtext{\edtext{\skt{तद्वा\-त\-शो\-णि\-त\-मि\-त्या\-च\-क्षते}}{\lemma{\skt{तद्}}\Afootnote{H; \skt{ति}\ A}}}{\lemma{\skt{वातशोणितम्}}\Afootnote{H; \skt{वात}\ A}}}{\lemma{\skt{इत्य्}}\Afootnote{H; \emph{om.}\ A}}}{\lemma{\skt{आचक्षते}}\Afootnote{H; \skt{रक्तम्}\ A}}\skt{~। }\edtext{\skt{तत्तु}}{\lemma{\skt{तत्तु}}\Afootnote{A; \skt{तन्तु}\ H}}\skt{ पूर्व्वं हस्त\-पा\-द\-यो\-र\-व\-स्था\-न\-ङ्कृत्वा }\edtext{\skt{पश्चाद्देहे}}{\lemma{\skt{देहे}}\Afootnote{H; \skt{देहं}\ A}}\skt{ }\edtext{\skt{विसर्प्पति}}{\lemma{\skt{विसर्प्पति}}\Afootnote{H; \skt{व्याप्नोति}\ A}}\skt{~। तस्य पूर्व्वरूपाणि }\edtext{\skt{तोद\-दा\-ह\-क\-ण्डू\-शो\-फ\-स्त\-म्भ\-त्व\-क्पा\-रु\-ष्य\-सि\-रा\-स्ना\-यु\-ध\-म\-नी\-स्प\-न्द\-न\-स\-क्थि\-दौ\-र्ब्ब\-ल्यानि}}{\lemma{\skt{तोद\appbrk{}दा\appbrk{}ह\appbrk{}क\appbrk{}ण्डू\appbrk{}शो\appbrk{}फ\appbrk{}स्त\appbrk{}म्भ\appbrk{}त्व\appbrk{}क्पा\appbrk{}रु\appbrk{}ष्य\appbrk{}सि\appbrk{}रा\appbrk{}स्ना\appbrk{}यु\appbrk{}ध\appbrk{}म\appbrk{}नी\appbrk{}स्प\appbrk{}न्द\appbrk{}न\appbrk{}स\appbrk{}क्थि\appbrk{}दौ\appbrk{}र्ब्ब\appbrk{}ल्यानि}}\Afootnote{A; \skt{तोद\appbrk{}दा\appbrk{}ह\appbrk{}क\appbrk{}ण्डू\appbrk{}शो\appbrk{}फ\appbrk{}स्त\appbrk{}म्भ\appbrk{}स्त्व\appbrk{}क्पा\appbrk{}रु\appbrk{}ष्य\appbrk{}सि\appbrk{}रा\appbrk{}स्ना\appbrk{}यु\appbrk{}ध\appbrk{}म\appbrk{}नी\appbrk{}स्प\appbrk{}न्द\appbrk{}न\appbrk{}स\appbrk{}क्थि\appbrk{}दौ\appbrk{}र्ब्ब\appbrk{}ल्यानि}\ H}}\skt{ }\edtext{\skt{तथा}}{\lemma{\skt{तथा}}\Afootnote{H; \emph{om.}\ A}}\skt{ }\edtext{\edtext{\edtext{\skt{श्यावा\-रु\-ण\-र\-क्त\-म\-ण्ड\-लो\-त्प\-त्ति\-र\-क\-स्माच्च}}{\lemma{\skt{श्यावा\appbrk{}रु\appbrk{}ण\appbrk{}र\appbrk{}क्त\appbrk{}म\appbrk{}ण्ड\appbrk{}लो\appbrk{}त्प\appbrk{}त्तिर्}}\Afootnote{H; \skt{श्यावा\appbrk{}रु\appbrk{}ण\appbrk{}म\appbrk{}ण्ड\appbrk{}लो\appbrk{}त्प\appbrk{}त्तिश्}\ A}}}{\lemma{\skt{अकस्माच्}}\Afootnote{\skt{चाकस्मात्}\ A; \skt{अ{\symfont ❈}कस्माच्}\ H}}}{\lemma{\skt{च}}\Afootnote{H; \skt{पाणि}\ A}}\skt{ }\edtext{\edtext{\edtext{\edtext{\skt{पाद\-त\-ला\-ङ्गु\-लि\-गु\-ल्फ\-म\-णि\-ब\-न्धे\-ष्वे\-षा\-मु\-त्प\-त्ति\-र्ब्भ\-वति}}{\lemma{\skt{पाद\appbrk{}त\appbrk{}ला\appbrk{}ङ्गु\appbrk{}लि\appbrk{}गु\appbrk{}ल्फ\appbrk{}म\appbrk{}णि\appbrk{}ब\appbrk{}न्धेष्व्}}\Afootnote{H; \skt{पाद\appbrk{}त\appbrk{}ला\appbrk{}ङ्गु\appbrk{}लि\appbrk{}गु\appbrk{}ल्फ\appbrk{}म\appbrk{}णि\appbrk{}ब\appbrk{}न्ध\appbrk{}प्र\appbrk{}भृ\appbrk{}तिषु}\ A}}}{\lemma{\skt{एषाम्}}\Afootnote{H; \emph{om.}\ A}}}{\lemma{\skt{उत्पत्तिर्}}\Afootnote{H; \emph{om.}\ A}}}{\lemma{\skt{ब्भवति}}\Afootnote{H; \emph{om.}\ A}}\skt{~। }\edtext{\edtext{\skt{तत्रा\-प्र\-ति\-सा\-रि\-णोऽ\-प\-चा\-रि\-णश्च}}{\lemma{\skt{तत्रा\appbrk{}प्र\appbrk{}ति\appbrk{}सा\appbrk{}रिणो}}\Afootnote{\skt{तत्रा\appbrk{}प्र\appbrk{}ति\appbrk{}का\appbrk{}रिणो}\ A; \skt{तत्रा\appbrk{}प्र\appbrk{}ति\appbrk{}सा\appbrk{}रिणो}\ H}}}{\lemma{\skt{ऽपचारिणश्च}}\Afootnote{A; \skt{पचारिणश्च}\ H}}\skt{ रोगः }\edtext{\skt{प्रव्यक्तो}}{\lemma{\skt{प्रव्यक्तो}}\Afootnote{H; \skt{व्यक्त}\ A}}\skt{ }\edtext{\skt{भवति}}{\lemma{\skt{भवति}}\Afootnote{H; \skt{तरः}\ A}}\skt{~। तस्य लक्षणमुक्तं~। }\edtext{\edtext{\skt{तत्रा\-प्य\-प्र\-ति\-का\-रिणो}}{\lemma{\skt{तत्राप्य्}}\Afootnote{H; \skt{तत्रा}\ A}}}{\lemma{\skt{अप्रतिकारिणो}}\Afootnote{H; \skt{प्रतिकारिणो}\ A}}\skt{ वैकल्यम्भवति~॥}\pend

\pstart\versepar \ledsidenote{\passref{4.5.5}}\edtext{\skt{प्रायशः}}{\lemma{\skt{प्रायशः}}\Afootnote{H; \skt{भवति चात्र प्रायशः}\ A}}\skt{ सुकु\-मा\-रा\-णा\-म्मि\-थ्या\-हा\-र\-वि\-हा\-रिणां~। }\pend
\pstart\versepar \skt{स्थूलानां सुखिनाञ्चापि }\edtext{\skt{कुप्यते}}{\lemma{\skt{कुप्यते}}\Afootnote{H; \skt{वातरक्तं प्रकुप्यति}\ A}}\skt{ }\edtext{\skt{वातशोणितम्}}{\lemma{\skt{वातशोणितम्}}\Afootnote{H; \emph{om.}\ A}}\skt{~॥}\pend

\pstart \ledsidenote{\passref{4.5.6}}\skt{तत्र }\edtext{\skt{प्राण\-क्ष\-य\-पि\-पा\-सा\-ज्व\-र\-मू\-र्च्छा\-श्वा\-स\-स्त\-म्भा\-रो\-च\-का\-न\-व\-शी\-र्य\-माणं}}{\lemma{\skt{प्राण\appbrk{}क्ष\appbrk{}य\appbrk{}पि\appbrk{}पा\appbrk{}सा\appbrk{}ज्व\appbrk{}र\appbrk{}मू\appbrk{}र्च्छा\appbrk{}श्वा\appbrk{}स\appbrk{}स्त\appbrk{}म्भा\appbrk{}रो\appbrk{}च\appbrk{}का\appbrk{}न\appbrk{}व\appbrk{}शी\appbrk{}र्य\appbrk{}माणं}}\Afootnote{H; \skt{प्राण\appbrk{}मां\appbrk{}स\appbrk{}क्ष\appbrk{}य\appbrk{}पि\appbrk{}पा\appbrk{}सा\appbrk{}ज्व\appbrk{}र\appbrk{}मू\appbrk{}र्च्छा\appbrk{}श्वा\appbrk{}स\appbrk{}का\appbrk{}स\appbrk{}स्त\appbrk{}म्भा\appbrk{}रो\appbrk{}च\appbrk{}का\appbrk{}वि\appbrk{}पा\appbrk{}क\appbrk{}वि\appbrk{}स\appbrk{}रण}\ A}}\skt{ }\edtext{\skt{सङ्को\-चै\-र\-नु\-प\-द्रु\-त\-म्ब\-ल\-व\-न्त\-मा\-त्म\-व\-न्त\-मु\-प\-क\-र\-ण\-व\-न्त\-ञ्चो\-प\-क्र\-मेत्}}{\lemma{\skt{सङ्कोचैर्}}\Afootnote{H; \skt{संकोचनैर्}\ A}}\skt{~।}\footnoteA{The negative present passive participle \emph{anavaśīryamāṇa} appears to be a hapax legomenon in Ayurvedic literature and rare everywhere (two occurrences in Bāṇa). However, the related term \emph{aviśīryamāṇa} occurs in the Carakasaṃhitā (3.8.122) and \emph{viśīryamāṇa} in the Aṣṭāṅgahṛdayasaṃhitā (2.5.98, 6.37.12) and elsewhere. An emendation to \emph{°arocakair aviśīryamāṇaṃ} might produce a more plausible text.}\pend

\pstart \ledsidenote{\passref{4.5.7}}\edtext{\skt{तत्रादावेव}}{\lemma{\skt{तत्रादाव्}}\Afootnote{H; \skt{तत्र आदाव्}\ A}}\skt{ }\edtext{\skt{बहु\-वा\-त\-रू\-क्ष\-म्ला\-ना\-ङ्गा\-दृते}}{\lemma{\skt{बहु\appbrk{}वा\appbrk{}त\appbrk{}रू\appbrk{}क्ष\appbrk{}म्ला\appbrk{}ना\appbrk{}ङ्गा\appbrk{}दृते}}\Afootnote{A; \skt{बहु\appbrk{}वा\appbrk{}त\appbrk{}रू\appbrk{}क्षा\appbrk{}म्ला\appbrk{}ना\appbrk{}ङ्गा\appbrk{}दृते}\ H}}\skt{ }\edtext{\skt{मार्ग्गावरणकं}}{\lemma{\skt{मार्ग्गावरणकं}}\Afootnote{H; \skt{मार्गावरणाद्}\ A}}\skt{ दुष्ट\-शो\-णि\-त\-म\-स\-कृ\-द\-ल्पा\-ल्प\-म\-व\-सि\-ञ्चे\-द्वा\-त\-को\-प\-भ\-यात्~। ततो वम\-ना\-दि\-भि\-रु\-प\-क्र\-मै\-रु\-प\-पाद्य }\edtext{\skt{संसृ\-ष्ट\-भ\-क्त\-वाते}}{\lemma{\skt{संसृ\appbrk{}ष्ट\appbrk{}भ\appbrk{}क्त\appbrk{}वाते}}\Afootnote{H; \skt{प्रति\appbrk{}सं\appbrk{}सृ\appbrk{}ष्ट\appbrk{}भक्तं वात}\ A}}\skt{ प्रबले पुराणघृतं पाययेत्~। अजाक्षीरं }\edtext{\skt{वार्द्धतैलं}}{\lemma{\skt{वार्द्धतैलं}}\Afootnote{H; \skt{वाऽर्धतैलं}\ A}}\skt{ }\edtext{\skt{मधुकाक्षयुक्तं}}{\lemma{\skt{मधुकाक्षयुक्तं}}\Afootnote{H; \skt{मधु\appbrk{}का\appbrk{}क्ष\appbrk{}य\appbrk{}युक्तं}\ A}}\skt{ शृगा\-ल\-वि\-न्ना\-सि\-द्धम्वा शर्क्क\-रा\-म\-धु\-म\-धुरं~। }\edtext{\skt{शुण्ठी\-क\-शे\-रु\-सि\-द्धम्वा}}{\lemma{\skt{शुण्ठी\appbrk{}क\appbrk{}शे\appbrk{}रु\appbrk{}सि\appbrk{}द्धम्वा}}\Afootnote{H; \skt{शुण्ठी\appbrk{}शृ\appbrk{}ङ्गा\appbrk{}ट\appbrk{}क\appbrk{}क\appbrk{}शे\appbrk{}रु\appbrk{}क\appbrk{}सिद्धं वा}\ A}}\skt{~। }\edtext{\skt{श्यामा\-रा\-स्ना\-सु\-ष\-वी\-शृ\-गा\-ल\-वि\-न्ना\-पी\-लु\-श\-ता\-व\-री\-श्व\-दं\-ष्ट्रा\-द्वि\-प\-ञ्च\-मू\-ल\-क्वा\-था\-ष्ट\-गु\-ण\-सि\-द्धेन वा}}{\lemma{\skt{श्यामा\appbrk{}रा\appbrk{}स्ना\appbrk{}सु\appbrk{}ष\appbrk{}वी\appbrk{}शृ\appbrk{}गा\appbrk{}ल\appbrk{}वि\appbrk{}न्ना\appbrk{}पी\appbrk{}लु\appbrk{}श\appbrk{}ता\appbrk{}व\appbrk{}री\appbrk{}श्व\appbrk{}दं\appbrk{}ष्ट्रा\appbrk{}द्वि\appbrk{}प\appbrk{}ञ्च\appbrk{}मू\appbrk{}ल\appbrk{}क्वा\appbrk{}था\appbrk{}ष्ट\appbrk{}गु\appbrk{}ण\appbrk{}सि\appbrk{}द्धेन वा}}\Afootnote{H; \skt{श्यामा\appbrk{}रा\appbrk{}स्ना\appbrk{}सु\appbrk{}ष\appbrk{}वी\appbrk{}शृ\appbrk{}गा\appbrk{}ल\appbrk{}वि\appbrk{}न्ना\appbrk{}पी\appbrk{}लु\appbrk{}श\appbrk{}ता\appbrk{}व\appbrk{}री\appbrk{}श्व\appbrk{}दं\appbrk{}ष्ट्रा\appbrk{}द्वि\appbrk{}प\appbrk{}ञ्च\appbrk{}मू\appbrk{}ली\appbrk{}सिद्धं वा~। द्विप\appbrk{}ञ्च\appbrk{}मू\appbrk{}ली\appbrk{}क्वा\appbrk{}था\appbrk{}ष्ट\appbrk{}गु\appbrk{}ण\appbrk{}सि\appbrk{}द्धेन}\ A}}\skt{ पयसा }\edtext{\skt{मेष\-शृ\-ङ्गी\-श्व\-दं\-ष्ट्रा\-म\-धु\-ना\-ग\-ब\-ला\-भ\-द्र\-दा\-रु\-व\-चा\-सु\-र\-भि\-क\-ल्क\-प्र\-ती\-वा\-प\-न्तै\-ल\-म्पा\-च\-यित्वा}}{\lemma{\skt{मेष\appbrk{}शृ\appbrk{}ङ्गी\appbrk{}श्व\appbrk{}दं\appbrk{}ष्ट्रा\appbrk{}म\appbrk{}धु\appbrk{}ना\appbrk{}ग\appbrk{}ब\appbrk{}ला\appbrk{}भ\appbrk{}द्र\appbrk{}दा\appbrk{}रु\appbrk{}व\appbrk{}चा\appbrk{}सु\appbrk{}र\appbrk{}भि\appbrk{}क\appbrk{}ल्क\appbrk{}प्र\appbrk{}ती\appbrk{}वा\appbrk{}पन्}}\Afootnote{\skt{मधु\appbrk{}क\appbrk{}मे\appbrk{}ष\appbrk{}शृ\appbrk{}ङ्गी\appbrk{}श्व\appbrk{}दं\appbrk{}ष्ट्रा\appbrk{}स\appbrk{}र\appbrk{}ल\appbrk{}भ\appbrk{}द्र\appbrk{}दा\appbrk{}रु\appbrk{}व\appbrk{}चा\appbrk{}सु\appbrk{}र\appbrk{}भि\appbrk{}क\appbrk{}ल्क\appbrk{}प्र\appbrk{}ती\appbrk{}वापं}\ A; \skt{मेष\appbrk{}शृ\appbrk{}ङ्गी\appbrk{}श्व\appbrk{}दं\appbrk{}ष्ट्रा\appbrk{}म\appbrk{}धु\appbrk{}ना\appbrk{}ग\appbrk{}ब\appbrk{}ला\appbrk{}भ\appbrk{}द्र\appbrk{}दा\appbrk{}रु\appbrk{}व\appbrk{}चा\appbrk{}सु\appbrk{}र\appbrk{}भि\appbrk{}क\appbrk{}ल्क\appbrk{}प्र\appbrk{}ति\appbrk{}वा\appbrk{}पन्}\ H}}\skt{ }\edtext{\skt{पाना\-दि\-षू\-प\-यु\-ञ्ज्यात्}}{\lemma{\skt{पाना\appbrk{}दि\appbrk{}षू\appbrk{}प\appbrk{}यु\appbrk{}ञ्ज्यात्}}\Afootnote{H; \skt{पाना\appbrk{}दि\appbrk{}षू\appbrk{}प\appbrk{}यु\appbrk{}ञ्जीत}\ A}}\skt{~। }\edtext{\skt{शता\-व\-री\-म\-यू\-र\-क\-ध\-व\-क\-म\-धु\-क\-क्षी\-र\-वि\-दा\-री\-ब\-ला\-ति\-ब\-ला\-तृ\-ण\-प\-ञ्च\-मू\-ली\-क्वा\-थ\-सि\-द्धम्वा}}{\lemma{\skt{शता\appbrk{}व\appbrk{}री\appbrk{}म\appbrk{}यू\appbrk{}र\appbrk{}क\appbrk{}ध\appbrk{}व\appbrk{}क\appbrk{}म\appbrk{}धु\appbrk{}क\appbrk{}क्षी\appbrk{}र\appbrk{}वि\appbrk{}दा\appbrk{}री\appbrk{}ब\appbrk{}ला\appbrk{}ति\appbrk{}ब\appbrk{}ला\appbrk{}तृ\appbrk{}ण\appbrk{}प\appbrk{}ञ्च\appbrk{}मू\appbrk{}ली\appbrk{}क्वा\appbrk{}थ\appbrk{}सि\appbrk{}द्धम्वा}}\Afootnote{H; \skt{शता\appbrk{}व\appbrk{}री\appbrk{}म\appbrk{}यू\appbrk{}र\appbrk{}क\appbrk{}कि\appbrk{}णि\appbrk{}ह्य\appbrk{}ज\appbrk{}मो\appbrk{}दा\appbrk{}म\appbrk{}धु\appbrk{}क\appbrk{}क्षी\appbrk{}र\appbrk{}वि\appbrk{}दा\appbrk{}री\appbrk{}ब\appbrk{}ला\appbrk{}ति\appbrk{}ब\appbrk{}ला\appbrk{}तृ\appbrk{}ण\appbrk{}प\appbrk{}ञ्च\appbrk{}मू\appbrk{}ली\appbrk{}क्वा\appbrk{}थ\appbrk{}सिद्धं वा}\ A}}\skt{ काको\-ल्या\-दि\-प्र\-ति\-वापं~। }\edtext{\skt{बलातैलम्वा}}{\lemma{\skt{बलातैलम्वा}}\Afootnote{H; \skt{बलातैलं}\ A}}\skt{ }\edtext{\skt{शतपाकमिति}}{\lemma{\skt{इति}}\Afootnote{H; \skt{वेति}\ A}}\skt{~। वात\-ह\-र\-मू\-ल\-सि\-द्धेन च पयसा }\edtext{\skt{अम्लैर्व्वा}}{\lemma{\skt{अम्लैर्व्वा}}\Afootnote{\skt{परि\appbrk{}षे\appbrk{}च\appbrk{}न\appbrk{}म\appbrk{}म्लैर्वा}\ A; \skt{म्लैर्व्वा}\ H}}\skt{ }\edtext{\skt{परि\-षे\-क\-ङ्कु\-र्व्वीत}}{\lemma{\skt{परिषेकङ्}}\Afootnote{H; \emph{om.}\ A}}\skt{~॥ }\edtext{\skt{तत्र}}{\lemma{\skt{तत्र}}\Afootnote{H; \skt{यव\appbrk{}म\appbrk{}धु\appbrk{}कै\appbrk{}र\appbrk{}ण्ड\appbrk{}ति\appbrk{}ल\appbrk{}व\appbrk{}र्षा\appbrk{}भू\appbrk{}भिर्वा प्रदेहः कार्यः~। तत्र}\ A}}\skt{ चूर्ण्णितेषु यव\-गो\-धू\-म\-ति\-ल\-मु\-द्ग\-मा\-षेषु प्रत्येकशः }\edtext{\skt{काको\-ली\-क्षी\-र\-का\-को\-ली\-जी\-व\-क\-र्ष\-भ\-क\-ब\-ला\-ति\-ब\-ला\-शृ\-गा\-ल\-वि\-न्ना\-मे\-ष\-शृ\-ङ्गी\-पि\-या\-ल\-श\-र्क्क\-रा\-क\-शे\-रु\-का\-सु\-र\-भि\-व\-चा\-क\-ल्क\-मि\-श्रे\-षू\-प\-ना\-हार्थे}}{\lemma{\skt{काको\appbrk{}ली\appbrk{}क्षी\appbrk{}र\appbrk{}का\appbrk{}को\appbrk{}ली\appbrk{}जी\appbrk{}व\appbrk{}क\appbrk{}र्ष\appbrk{}भ\appbrk{}क\appbrk{}ब\appbrk{}ला\appbrk{}ति\appbrk{}ब\appbrk{}ला\appbrk{}शृ\appbrk{}गा\appbrk{}ल\appbrk{}वि\appbrk{}न्ना\appbrk{}मे\appbrk{}ष\appbrk{}शृ\appbrk{}ङ्गी\appbrk{}पि\appbrk{}या\appbrk{}ल\appbrk{}श\appbrk{}र्क्क\appbrk{}रा\appbrk{}क\appbrk{}शे\appbrk{}रु\appbrk{}का\appbrk{}सु\appbrk{}र\appbrk{}भि\appbrk{}व\appbrk{}चा\appbrk{}क\appbrk{}ल्क\appbrk{}मि\appbrk{}श्रे\appbrk{}षू\appbrk{}प\appbrk{}ना\appbrk{}हार्थे}}\Afootnote{H; \skt{काको\appbrk{}ली\appbrk{}क्षी\appbrk{}र\appbrk{}का\appbrk{}को\appbrk{}ली\appbrk{}जी\appbrk{}व\appbrk{}क\appbrk{}र्ष\appbrk{}भ\appbrk{}क\appbrk{}ब\appbrk{}ला\appbrk{}ति\appbrk{}ब\appbrk{}ला\appbrk{}बि\appbrk{}स\appbrk{}मृ\appbrk{}णा\appbrk{}ल\appbrk{}शृ\appbrk{}गा\appbrk{}ल\appbrk{}वि\appbrk{}न्ना\appbrk{}मे\appbrk{}ष\appbrk{}शृ\appbrk{}ङ्गी\appbrk{}प्रि\appbrk{}या\appbrk{}ल\appbrk{}श\appbrk{}र्क\appbrk{}रा\appbrk{}क\appbrk{}शे\appbrk{}रु\appbrk{}क\appbrk{}सु\appbrk{}र\appbrk{}भि\appbrk{}व\appbrk{}चा\appbrk{}क\appbrk{}ल्क\appbrk{}मि\appbrk{}श्रे\appbrk{}षू\appbrk{}प\appbrk{}ना\appbrk{}हार्थं}\ A}}\skt{ }\edtext{\skt{सर्प्पि\-स्तै\-ल\-व\-सा\-म\-ज्जा\-दु\-ग्ध\-सिद्धाः}}{\lemma{\skt{सर्प्पि\appbrk{}स्तै\appbrk{}ल\appbrk{}व\appbrk{}सा\appbrk{}म\appbrk{}ज्जा\appbrk{}दु\appbrk{}ग्ध\appbrk{}सिद्धाः}}\Afootnote{H; \skt{सर्पि\appbrk{}स्तै\appbrk{}ल\appbrk{}व\appbrk{}सा\appbrk{}म\appbrk{}ज्ज\appbrk{}दु\appbrk{}ग्ध\appbrk{}सिद्धाः}\ A}}\skt{ पञ्च पायसा व्याख्याताः~। }\edtext{\skt{स्नैहि\-क\-फ\-ल\-सा\-रो\-त्का\-रिका}}{\lemma{\skt{स्नैहि\appbrk{}क\appbrk{}फ\appbrk{}ल\appbrk{}सा\appbrk{}रो\appbrk{}त्का\appbrk{}रिका}}\Afootnote{H; \skt{स्नैहि\appbrk{}क\appbrk{}फ\appbrk{}ल\appbrk{}सा\appbrk{}रो\appbrk{}त्का\appbrk{}रिका वा}\ A}}\skt{ }\edtext{\skt{नल\-मी\-न\-म\-त्स्य\-पि\-शि\-त\-वै\-श\-वारो वा}}{\lemma{\skt{नल\appbrk{}मी\appbrk{}न\appbrk{}म\appbrk{}त्स्य\appbrk{}पि\appbrk{}शि\appbrk{}त\appbrk{}वै\appbrk{}श\appbrk{}वारो वा}}\Afootnote{\skt{चूर्णितेषु यव\appbrk{}गो\appbrk{}धू\appbrk{}म\appbrk{}ति\appbrk{}ल\appbrk{}मु\appbrk{}द्ग\appbrk{}मा\appbrk{}षेषु मत्स्य\appbrk{}पि\appbrk{}शि\appbrk{}त\appbrk{}वे\appbrk{}श\appbrk{}वारो वा}\ A; \skt{नल\appbrk{}पी\appbrk{}न\appbrk{}म\appbrk{}त्स्य\appbrk{}पि\appbrk{}शि\appbrk{}त\appbrk{}वै\appbrk{}श\appbrk{}वारो वा}\ H}}\skt{~। }\edtext{\skt{बिल्व\-पे\-सि\-का\-त\-ग\-र\-दे\-व\-दा\-रु\-स\-र\-ल\-रा\-स्ना\-ह\-रे\-णु\-कु\-ष्ठ\-श\-त\-पु\-ष्प\-सु\-रा\-द\-धि\-म\-स्तु\-यु\-क्त\-मु\-प\-नाहः}}{\lemma{\skt{बिल्व\appbrk{}पे\appbrk{}सि\appbrk{}का\appbrk{}त\appbrk{}ग\appbrk{}र\appbrk{}दे\appbrk{}व\appbrk{}दा\appbrk{}रु\appbrk{}स\appbrk{}र\appbrk{}ल\appbrk{}रा\appbrk{}स्ना\appbrk{}ह\appbrk{}रे\appbrk{}णु\appbrk{}कु\appbrk{}ष्ठ\appbrk{}श\appbrk{}त\appbrk{}पु\appbrk{}ष्प\appbrk{}सु\appbrk{}रा\appbrk{}द\appbrk{}धि\appbrk{}म\appbrk{}स्तु\appbrk{}यु\appbrk{}क्तम्}}\Afootnote{H; \skt{बल्व\appbrk{}पे\appbrk{}शि\appbrk{}का\appbrk{}त\appbrk{}ग\appbrk{}र\appbrk{}दे\appbrk{}व\appbrk{}दा\appbrk{}रु\appbrk{}स\appbrk{}र\appbrk{}ला\appbrk{}रा\appbrk{}स्ना\appbrk{}ह\appbrk{}रे\appbrk{}णु\appbrk{}कु\appbrk{}ष्ठ\appbrk{}श\appbrk{}त\appbrk{}पु\appbrk{}ष्पै\appbrk{}ला\appbrk{}सु\appbrk{}रा\appbrk{}द\appbrk{}धि\appbrk{}म\appbrk{}स्तु\appbrk{}युक्त}\ A}}\skt{~। }\edtext{\edtext{\skt{मातु\-लु\-ङ्गा\-म्ल\-सै\-न्ध\-व\-घृ\-त\-मि\-श्र\-म\-धु\-शि\-ग्रु\-मू\-ल\-मा\-ले\-पनं}}{\lemma{\skt{मातु\appbrk{}लु\appbrk{}ङ्गा\appbrk{}म्ल\appbrk{}सै\appbrk{}न्ध\appbrk{}व\appbrk{}घृ\appbrk{}त\appbrk{}मि\appbrk{}श्र\appbrk{}म\appbrk{}धु\appbrk{}शि\appbrk{}ग्रु\appbrk{}मू\appbrk{}लम्}}\Afootnote{H; \skt{मातु\appbrk{}लु\appbrk{}ङ्गा\appbrk{}म्ल\appbrk{}सै\appbrk{}न्ध\appbrk{}व\appbrk{}घृ\appbrk{}त\appbrk{}मिश्रं मधु\appbrk{}शि\appbrk{}ग्रु\appbrk{}मू\appbrk{}लम्}\ A}}}{\lemma{\skt{आलेपनं}}\Afootnote{H; \skt{आलेपस्}\ A}}\skt{ तिलकल्को }\edtext{\skt{वेति}}{\lemma{\skt{वेति}}\Afootnote{H; \skt{वेति वातप्रबले}\ A}}\skt{~॥}\footnoteA{The fish-name \emph{nalapīna} in witness H is a hapax legomenon; Ḍalhaṇa reported seeing manuscripts reading \emph{nalamīna} at this locus (Ācārya 1938: 425b).}\pend

\pstart \ledsidenote{\passref{4.5.8}}\skt{पित्तप्रबले }\edtext{\edtext{\skt{द्राक्षा\-रे\-व\-त\-क\-ट्फ\-ल\-प\-य\-स्या\-म\-धु\-क\-च\-न्द\-न\-का\-श्म\-र्य\-क\-षा\-य\-म्म\-धु\-म\-धु\-र\-म्पा\-य\-येत्}}{\lemma{\skt{द्राक्षा\appbrk{}रे\appbrk{}व\appbrk{}त\appbrk{}क\appbrk{}ट्फ\appbrk{}ल\appbrk{}प\appbrk{}य\appbrk{}स्या\appbrk{}म\appbrk{}धु\appbrk{}क\appbrk{}च\appbrk{}न्द\appbrk{}न\appbrk{}का\appbrk{}श्म\appbrk{}र्य\appbrk{}क\appbrk{}षा\appbrk{}यम्}}\Afootnote{A; \skt{द्राक्षा\appbrk{}रे\appbrk{}व\appbrk{}त\appbrk{}क\appbrk{}फ\appbrk{}ल\appbrk{}प\appbrk{}य\appbrk{}स्या\appbrk{}म\appbrk{}धु\appbrk{}क\appbrk{}च\appbrk{}न्द\appbrk{}न\appbrk{}का\appbrk{}श्म\appbrk{}र्य\appbrk{}क\appbrk{}षा\appbrk{}यम्}\ H}}}{\lemma{\skt{मधुमधुरम्}}\Afootnote{H; \skt{शर्क\appbrk{}रा\appbrk{}म\appbrk{}धु\appbrk{}म\appbrk{}धुरं}\ A}}\skt{~। }\edtext{\edtext{\skt{शता\-व\-री\-प\-टो\-ल\-प\-त्त्र\-त्रि\-फ\-ला\-क\-टु\-रो\-हि\-णी\-गु\-डू\-ची\-क\-षा\-यम्वा}}{\lemma{\skt{शता\appbrk{}व\appbrk{}री\appbrk{}प\appbrk{}टो\appbrk{}ल\appbrk{}प\appbrk{}त्त्र\appbrk{}त्रि\appbrk{}फ\appbrk{}ला\appbrk{}क\appbrk{}टु\appbrk{}रो\appbrk{}हि\appbrk{}णी\appbrk{}गु\appbrk{}डू\appbrk{}ची\appbrk{}क\appbrk{}षा\appbrk{}यम्}}\Afootnote{\skt{शता\appbrk{}व\appbrk{}री\appbrk{}म\appbrk{}धु\appbrk{}क\appbrk{}प\appbrk{}टो\appbrk{}ल\appbrk{}त्रि\appbrk{}फ\appbrk{}ला\appbrk{}क\appbrk{}टु\appbrk{}रो\appbrk{}हि\appbrk{}णी\appbrk{}क\appbrk{}षायं, गुडूचीकषायं वा}\ A; \skt{शता\appbrk{}व\appbrk{}री\appbrk{}प\appbrk{}टो\appbrk{}ल\appbrk{}प\appbrk{}त्त्र\appbrk{}त्रि\appbrk{}फ\appbrk{}ला\appbrk{}क\appbrk{}डु\appbrk{}रो\appbrk{}हि\appbrk{}णी\appbrk{}गु\appbrk{}डू\appbrk{}ची\appbrk{}क\appbrk{}षा\appbrk{}यम्}\ H}}}{\lemma{\skt{वा}}\Afootnote{H; \skt{पित्तज्वरहरं वा}\ A}}\skt{~। }\edtext{\skt{मधु\-र\-ति\-क्त\-क\-षा\-य\-सिद्धं}}{\lemma{\skt{मधु\appbrk{}र\appbrk{}ति\appbrk{}क्त\appbrk{}क\appbrk{}षा\appbrk{}य\appbrk{}सिद्धं}}\Afootnote{H; \skt{चन्द\appbrk{}ना\appbrk{}दि\appbrk{}क\appbrk{}षायं शर्क\appbrk{}रा\appbrk{}म\appbrk{}धु\appbrk{}म\appbrk{}धुरं, मधु\appbrk{}र\appbrk{}ति\appbrk{}क्त\appbrk{}क\appbrk{}षा\appbrk{}य\appbrk{}सिद्धं वा}\ A}}\skt{ सर्प्पिः~। }\edtext{\skt{बिस\-मृ\-णा\-ल\-भ\-द्र\-श्रि\-य\-प\-द्म\-क\-क\-षा\-ये\-णा\-जा\-क्षी\-रेण}}{\lemma{\skt{बिस\appbrk{}मृ\appbrk{}णा\appbrk{}ल\appbrk{}भ\appbrk{}द्र\appbrk{}श्रि\appbrk{}य\appbrk{}प\appbrk{}द्म\appbrk{}क\appbrk{}क\appbrk{}षा\appbrk{}ये\appbrk{}णा\appbrk{}जा\appbrk{}क्षी\appbrk{}रेण}}\Afootnote{H; \skt{बिस\appbrk{}मृ\appbrk{}णा\appbrk{}ल\appbrk{}भ\appbrk{}द्र\appbrk{}श्रि\appbrk{}य\appbrk{}प\appbrk{}द्म\appbrk{}क\appbrk{}क\appbrk{}षा\appbrk{}ये\appbrk{}णा\appbrk{}र्ध\appbrk{}क्षी\appbrk{}रेण}\ A}}\skt{ परिषेकः~। }\edtext{\skt{क्षीरे\-क्षु\-र\-स\-म\-धु\-क\-श\-र्क्क\-रा\-त\-ण्डु\-लो\-द\-कैर्व्वा}}{\lemma{\skt{क्षीरे\appbrk{}क्षु\appbrk{}र\appbrk{}स\appbrk{}म\appbrk{}धु\appbrk{}क\appbrk{}श\appbrk{}र्क्क\appbrk{}रा\appbrk{}त\appbrk{}ण्डु\appbrk{}लो\appbrk{}द\appbrk{}कैर्व्वा}}\Afootnote{\skt{क्षीरे\appbrk{}क्षु\appbrk{}र\appbrk{}सै\appbrk{}र्म\appbrk{}धु\appbrk{}क\appbrk{}श\appbrk{}र्क\appbrk{}रा\appbrk{}त\appbrk{}ण्डु\appbrk{}लो\appbrk{}द\appbrk{}कैर्वा}\ A; \skt{क्षीरे\appbrk{}क्षु\appbrk{}र\appbrk{}स\appbrk{}म\appbrk{}धु\appbrk{}श\appbrk{}र्क्क\appbrk{}रा\appbrk{}त\appbrk{}ण्डु\appbrk{}लो\appbrk{}द\appbrk{}कैर्व्वा}\ H}}\skt{~। द्राक्षे\-क्षु\-क\-षा\-य\-मि\-श्रैर्व्वा~। }\edtext{\skt{मस्तु\-धा\-न्या\-म्लै\-र्ज्जी\-व\-नी\-य\-सि\-द्धेन}}{\lemma{\skt{मस्तुधान्याम्लैर्}}\Afootnote{H; \skt{मस्तु\appbrk{}म\appbrk{}द्य\appbrk{}धा\appbrk{}न्याम्लैः}\ A}}\skt{ वा सर्प्पिषा~। }\edtext{\skt{शतधौतेन वा}}{\lemma{\skt{शतधौतेन वा}}\Afootnote{\skt{ऽभ्यङ्गः, शत\appbrk{}धौ\appbrk{}त\appbrk{}घृ\appbrk{}तेन वा, काको\appbrk{}ल्या\appbrk{}दि\appbrk{}क\appbrk{}ल्क\appbrk{}क\appbrk{}षा\appbrk{}य\appbrk{}वि\appbrk{}प\appbrk{}क्वेन वा}\ A; \skt{शतधौते}\ H}}\skt{ }\edtext{\skt{शालिपिष्टकं}}{\lemma{\skt{शालिपिष्टकं}}\Afootnote{\skt{सर्पिषा; शालिषष्टिक}\ A; \skt{न वा शालिपिष्टकं}\ H}}\skt{~। }\edtext{\skt{नल\-व\-ञ्जु\-ल\-ता\-ली\-स\-शृ\-ङ्गा\-ट\-क\-क\-लो\-ड्य\-गौ\-री\-शै\-व\-ल\-प\-द्म\-प्र\-भृ\-ति\-भिर्व्वा}}{\lemma{\skt{नल\appbrk{}व\appbrk{}ञ्जु\appbrk{}ल\appbrk{}ता\appbrk{}ली\appbrk{}स\appbrk{}शृ\appbrk{}ङ्गा\appbrk{}ट\appbrk{}क\appbrk{}क\appbrk{}लो\appbrk{}ड्य\appbrk{}गौ\appbrk{}री\appbrk{}शै\appbrk{}व\appbrk{}ल\appbrk{}प\appbrk{}द्म\appbrk{}प्र\appbrk{}भृ\appbrk{}ति\appbrk{}भिर्व्वा}}\Afootnote{\skt{नल\appbrk{}व\appbrk{}ञ्जु\appbrk{}ल\appbrk{}ता\appbrk{}ली\appbrk{}स\appbrk{}शृ\appbrk{}ङ्गा\appbrk{}ट\appbrk{}क\appbrk{}ग\appbrk{}लो\appbrk{}ड्य\appbrk{}गौ\appbrk{}री\appbrk{}गै\appbrk{}रि\appbrk{}क\appbrk{}शै\appbrk{}व\appbrk{}ल\appbrk{}प\appbrk{}द्म\appbrk{}क\appbrk{}प\appbrk{}द्म\appbrk{}प\appbrk{}त्र\appbrk{}प्र\appbrk{}भृ\appbrk{}ति\appbrk{}भिर्}\ A; \skt{नल\appbrk{}व\appbrk{}ञ्जु\appbrk{}ल\appbrk{}ता\appbrk{}ली\appbrk{}स\appbrk{}शृ\appbrk{}ङ्गा\appbrk{}ट\appbrk{}क\appbrk{}का\appbrk{}लो\appbrk{}द्या\appbrk{}गौ\appbrk{}री\appbrk{}शै\appbrk{}वा\appbrk{}ल\appbrk{}प\appbrk{}द्म\appbrk{}प्र\appbrk{}भृ\appbrk{}ति\appbrk{}भिर्व्वा}\ H}}\skt{ }\edtext{\skt{धान्याम्लपिष्टः}}{\lemma{\skt{धान्याम्लपिष्टः}}\Afootnote{H; \skt{धान्याम्लपिष्टैः}\ A}}\skt{ प्रदेहो }\edtext{\skt{घृतमिश्रः}}{\lemma{\skt{घृतमिश्रः}}\Afootnote{H; \skt{घृतमिश्रः, वात\appbrk{}प्र\appbrk{}ब\appbrk{}लेऽ\appbrk{}प्येष सुखोष्णः प्रदेहः कार्यः}\ A}}\skt{~॥}\footnoteA{We emend witness H's \emph{kālodyā} to the attested lexeme \emph{kaloḍya}.}\pend

\pstart \ledsidenote{\passref{4.5.9}}\skt{रक्त\-प्र\-ब\-लेऽ\-प्येवं बहुशश्च }\edtext{\skt{शोणि\-त\-म\-व\-से\-च\-येत्}}{\lemma{\skt{अवसेचयेत्}}\Afootnote{H; \skt{अवसेचयेत्, शीततमाश्च प्रदेहाः कार्या इति}\ A}}\skt{~॥}\pend

\pstart \ledsidenote{\passref{4.5.10}}\skt{श्लेष्मप्रबले त्वाम\-ल\-क\-ह\-रि\-द्रा\-क\-षा\-य\-म्म\-धु\-म\-धु\-र\-म्पा\-य\-येत्~। त्रिफलाकषायं वा~। }\edtext{\skt{मधु\-क\-शृ\-ङ्ग\-वे\-र\-ह\-री\-त\-की\-ति\-क्त\-रो\-हि\-णी\-क\-ल्कम्वा}}{\lemma{\skt{मधु\appbrk{}क\appbrk{}शृ\appbrk{}ङ्ग\appbrk{}वे\appbrk{}र\appbrk{}ह\appbrk{}री\appbrk{}त\appbrk{}की\appbrk{}ति\appbrk{}क्त\appbrk{}रो\appbrk{}हि\appbrk{}णी\appbrk{}क\appbrk{}ल्कम्वा}}\Afootnote{A; \skt{मधु\appbrk{}क\appbrk{}शृ\appbrk{}ङ्ग\appbrk{}वे\appbrk{}र\appbrk{}ह\appbrk{}री\appbrk{}त\appbrk{}की\appbrk{}ति\appbrk{}क्त\appbrk{}रो\appbrk{}हि\appbrk{}णी\appbrk{}के\appbrk{}ल्कम्वा}\ H}}\skt{ }\edtext{\skt{सक्षौद्रम्}}{\lemma{\skt{सक्षौद्रम्}}\Afootnote{A; \skt{सक्षौद्र}\ H}}\skt{~। }\edtext{\skt{मूत्रेण}}{\lemma{\skt{मूत्रेण}}\Afootnote{H; \skt{मूत्र\appbrk{}तो\appbrk{}य\appbrk{}यो\appbrk{}र\appbrk{}न्य\appbrk{}त\appbrk{}रेण}\ A}}\skt{ }\edtext{\skt{तोयेन}}{\lemma{\skt{तोयेन}}\Afootnote{H; \skt{गुड}\ A}}\skt{ }\edtext{\skt{हरीतकीम्वा}}{\lemma{\skt{हरीतकीम्वा}}\Afootnote{A; \skt{हरीतकीम्}\ H}}\skt{ }\edtext{\skt{पाययेत्}}{\lemma{\skt{पाययेत्}}\Afootnote{H; \skt{भक्षयेत्}\ A}}\skt{~। }\edtext{\skt{तैल\-मू\-त्र\-क्षा\-रो\-द\-क\-सु\-रा\-शु\-क्तैश्च}}{\lemma{\skt{तैल\appbrk{}मू\appbrk{}त्र\appbrk{}क्षा\appbrk{}रो\appbrk{}द\appbrk{}क\appbrk{}सु\appbrk{}रा\appbrk{}शु\appbrk{}क्तैश्च}}\Afootnote{H; \skt{तैल\appbrk{}मू\appbrk{}त्र\appbrk{}क्षा\appbrk{}रो\appbrk{}द\appbrk{}क\appbrk{}सु\appbrk{}रा\appbrk{}शु\appbrk{}क्त\appbrk{}क\appbrk{}फ\appbrk{}घ्नौ\appbrk{}ष\appbrk{}ध\appbrk{}निः\appbrk{}क्वा\appbrk{}थैश्च}\ A}}\skt{ परिषेकः~। }\edtext{\skt{आर\-ग्व\-धा\-दि\-क\-षा\-येण}}{\lemma{\skt{आर\appbrk{}ग्व\appbrk{}धा\appbrk{}दि\appbrk{}क\appbrk{}षा\appbrk{}येण}}\Afootnote{H; \skt{आर\appbrk{}ग्व\appbrk{}धा\appbrk{}दि\appbrk{}क\appbrk{}षा\appbrk{}यैर्वो}\ A}}\skt{ }\edtext{\skt{वा}}{\lemma{\skt{वा}}\Afootnote{H; \skt{ष्णैः}\ A}}\skt{~। }\edtext{\skt{मस्तु\-मू\-त्र\-सु\-रा\-शु\-क्ता\-म\-धु\-क\-सा\-रि\-वा\-प\-द्म\-क\-सि\-द्धेन वा}}{\lemma{\skt{मस्तु\appbrk{}मू\appbrk{}त्र\appbrk{}सु\appbrk{}रा\appbrk{}शु\appbrk{}क्ता\appbrk{}म\appbrk{}धु\appbrk{}क\appbrk{}सा\appbrk{}रि\appbrk{}वा\appbrk{}प\appbrk{}द्म\appbrk{}क\appbrk{}सि\appbrk{}द्धेन वा}}\Afootnote{\skt{मस्तु\appbrk{}मू\appbrk{}त्र\appbrk{}सु\appbrk{}रा\appbrk{}शु\appbrk{}क्त\appbrk{}म\appbrk{}धु\appbrk{}क\appbrk{}सा\appbrk{}रि\appbrk{}वा\appbrk{}प\appbrk{}द्म\appbrk{}क\appbrk{}सिद्धं वा}\ A; \skt{मस्तु\appbrk{}मू\appbrk{}त्र\appbrk{}सु\appbrk{}रा\appbrk{}शु\appbrk{}क्ता\appbrk{}म\appbrk{}धु\appbrk{}क\appbrk{}शा\appbrk{}रि\appbrk{}वा\appbrk{}प\appbrk{}द्म\appbrk{}क\appbrk{}सिद्धे वा}\ H}}\skt{ }\edtext{\skt{सर्प्पिषाभ्यङ्गः}}{\lemma{\skt{सर्प्पिषाभ्यङ्गः}}\Afootnote{H; \skt{घृतमभ्यङ्गः}\ A}}\skt{~। }\edtext{\skt{श्वेत\-स\-र्ष्ष\-प\-कल्कः}}{\lemma{\skt{श्वेत\appbrk{}स\appbrk{}र्ष्ष\appbrk{}प\appbrk{}कल्कः}}\Afootnote{H; \skt{तिल\appbrk{}स\appbrk{}र्ष\appbrk{}पा\appbrk{}त\appbrk{}सी\appbrk{}य\appbrk{}व\appbrk{}चू\appbrk{}र्णानि श्लेष्मा\appbrk{}त\appbrk{}क\appbrk{}क\appbrk{}पि\appbrk{}त्थ\appbrk{}मु\appbrk{}धु\appbrk{}शि\appbrk{}ग्रु\appbrk{}मि\appbrk{}श्राणि क्षार\appbrk{}मू\appbrk{}त्र\appbrk{}पि\appbrk{}ष्टानि प्रदेहः; श्वेत\appbrk{}स\appbrk{}र्ष\appbrk{}प\appbrk{}कल्कः}\ A}}\skt{~। }\edtext{\skt{तिला\-श्व\-ग\-न्धा\-कल्कः}}{\lemma{\skt{तिला\appbrk{}श्व\appbrk{}ग\appbrk{}न्धा\appbrk{}कल्कः}}\Afootnote{H\textsuperscript{ac}; \skt{तिला\appbrk{}श्व\appbrk{}ग\appbrk{}न्धा\appbrk{}कल्कः, प्रियाल}\ A; \skt{तिला\appbrk{}श्व\appbrk{}ग\appbrk{}न्धा\appbrk{}य\appbrk{}वागू कल्कः}\ H\textsuperscript{pc}}}\skt{~।}\footnoteA{A scribe recorded a reading \emph{yavāgū} in the top margin, apparently wanting to replace one of the terms in this line, but it is unclear which.}\skt{ }\edtext{\skt{दारु\-शे\-लु\-क\-पि\-त्थ\-कल्कः}}{\lemma{\skt{दारु\appbrk{}शे\appbrk{}लु\appbrk{}क\appbrk{}पि\appbrk{}त्थ\appbrk{}कल्कः}}\Afootnote{H; \skt{सेलु\appbrk{}क\appbrk{}पि\appbrk{}त्थ\appbrk{}कल्कः}\ A}}\skt{~। }\edtext{\skt{मधु\-शि\-ग्रु\-पु\-न\-र्न्न\-वा\-कल्कः}}{\lemma{\skt{मधु\appbrk{}शि\appbrk{}ग्रु\appbrk{}पु\appbrk{}न\appbrk{}र्न्न\appbrk{}वा\appbrk{}कल्कः}}\Afootnote{A; \skt{मधु\appbrk{}शि\appbrk{}ग्रु\appbrk{}पु\appbrk{}न\appbrk{}र्न्न\appbrk{}व\appbrk{}कल्कः}\ H}}\skt{~। }\edtext{\skt{व्योष\-ति\-क्ता\-पृ\-थ\-क्प\-र्ण्णी\-बृ\-ह\-ती\-क\-ल्कैर्व्वा}}{\lemma{\skt{व्योष\appbrk{}ति\appbrk{}क्ता\appbrk{}पृ\appbrk{}थ\appbrk{}क्प\appbrk{}र्ण्णी\appbrk{}बृ\appbrk{}ह\appbrk{}ती\appbrk{}क\appbrk{}ल्कैर्व्वा}}\Afootnote{H; \skt{व्योष\appbrk{}ति\appbrk{}क्ता\appbrk{}पृ\appbrk{}थ\appbrk{}क्प\appbrk{}र्णी\appbrk{}बृ\appbrk{}ह\appbrk{}ती\appbrk{}कल्क}\ A}}\skt{~। }\edtext{\edtext{\skt{एवमेते}}{\lemma{\skt{एवम्}}\Afootnote{H; \skt{इत्य्}\ A}}}{\lemma{\skt{एते}}\Afootnote{H; \skt{एतेषां}\ A}}\skt{ पञ्च प्रदेहाः }\edtext{\skt{क्षारोदकपिष्टाः}}{\lemma{\skt{क्षारोदकपिष्टाः}}\Afootnote{H; \skt{सुखोष्णाः क्षारोदकपिष्}\ A}}\skt{ }\edtext{\skt{व्याख्याताः}}{\lemma{\skt{व्याख्याताः}}\Afootnote{H; \skt{टाः, शालिपर्णी पृश्निपर्णी बृहत्यौ वा क्षीर\appbrk{}पि\appbrk{}ष्टा\appbrk{}स्त\appbrk{}र्प\appbrk{}ण\appbrk{}मिश्राः}\ A}}\skt{~॥}\pend

\pstart \ledsidenote{\passref{4.5.11}}\skt{संसर्ग्गे सन्निपाते च क्रियापथमुक्तं मिश्रं }\edtext{\skt{कुर्यादिति}}{\lemma{\skt{कुर्यादिति}}\Afootnote{H; \skt{कुर्यात्}\ A}}\skt{~।}\pend

\pstart \ledsidenote{\passref{4.5.12}}\skt{सर्व्वेषु च गुड\-ह\-री\-त\-की\-मा\-से\-वेत~। }\edtext{\skt{पिप्पलीर्वा}}{\lemma{\skt{पिप्पलीर्वा}}\Afootnote{A; \skt{पिप्पल्यो वा}\ H}}\skt{ }\edtext{\skt{क्षीरपिष्टाः}}{\lemma{\skt{क्षीरपिष्टाः}}\Afootnote{H; \skt{क्षीरपिष्टा वारिपिष्टा वा}\ A}}\skt{ }\edtext{\skt{पञ्चाभिवृद्ध्या}}{\lemma{\skt{पञ्चाभिवृद्ध्या}}\Afootnote{H; \skt{पञ्चाभिवृद्ध्या दशाभिवृद्ध्या वा}\ A}}\skt{ }\edtext{\skt{पिबे\-त्क्षी\-रो\-द\-ना\-हारो}}{\lemma{\skt{क्षीरोदनाहारो}}\Afootnote{H; \skt{क्षीरौदनाहारो}\ A}}\skt{ दशरात्रं~। }\edtext{\skt{पञ्चाभिवृद्ध्या}}{\lemma{\skt{पञ्चाभिवृद्ध्या}}\Afootnote{H; \emph{om.}\ A}}\skt{ }\edtext{\skt{भूय\-श्चा\-प\-क\-र्ष\-येत्}}{\lemma{\skt{भूयश्}}\Afootnote{A; \skt{भूयञ्}\ H}}\skt{~। }\edtext{\skt{एवमेव}}{\lemma{\skt{एव}}\Afootnote{H; \emph{om.}\ A}}\skt{ }\edtext{\skt{यावत्पञ्चेति}}{\lemma{\skt{पञ्चेति}}\Afootnote{H; \skt{पञ्च दश वेति}\ A}}\skt{~। }\edtext{\skt{तदे\-त\-त्पि\-प्प\-ली\-व\-र्द्ध\-मा\-न\-क\-मा\-च\-क्षते}}{\lemma{\skt{आचक्षते}}\Afootnote{H; \emph{om.}\ A}}\skt{~॥ }\edtext{\skt{एतद्धि}}{\lemma{\skt{एतद्धि}}\Afootnote{H; \emph{om.}\ A}}\skt{ }\edtext{\skt{वात\-शो\-णि\-त\-वि\-ष\-म\-ज्व\-रा\-रो\-च\-क\-पा\-ण्डु\-रो\-गो\-द\-रा\-र्शः\-श्वा\-स\-का\-स\-शो\-षा\-ग्नि\-सा\-द\-हृ\-द्रो\-गा\-नु\-प\-हन्ति}}{\lemma{\skt{वात\appbrk{}शो\appbrk{}णि\appbrk{}त\appbrk{}वि\appbrk{}ष\appbrk{}म\appbrk{}ज्व\appbrk{}रा\appbrk{}रो\appbrk{}च\appbrk{}क\appbrk{}पा\appbrk{}ण्डु\appbrk{}रो\appbrk{}गो\appbrk{}द\appbrk{}रा\appbrk{}र्शः\appbrk{}श्वा\appbrk{}स\appbrk{}का\appbrk{}स\appbrk{}शो\appbrk{}षा\appbrk{}ग्नि\appbrk{}सा\appbrk{}द\appbrk{}हृ\appbrk{}द्रो\appbrk{}गा\appbrk{}नु\appbrk{}प\appbrk{}हन्ति}}\Afootnote{\skt{वात\appbrk{}शो\appbrk{}णि\appbrk{}त\appbrk{}वि\appbrk{}ष\appbrk{}म\appbrk{}ज्व\appbrk{}रा\appbrk{}रो\appbrk{}च\appbrk{}क\appbrk{}पा\appbrk{}ण्डु\appbrk{}रो\appbrk{}ग\appbrk{}प्ली\appbrk{}हो\appbrk{}द\appbrk{}रा\appbrk{}र्शः\appbrk{}का\appbrk{}स\appbrk{}श्वा\appbrk{}स\appbrk{}शो\appbrk{}फ\appbrk{}शो\appbrk{}षा\appbrk{}ग्नि\appbrk{}सा\appbrk{}द\appbrk{}हृ\appbrk{}द्रो\appbrk{}गो\appbrk{}द\appbrk{}रा\appbrk{}ण्य\appbrk{}प\appbrk{}हन्ति}\ A; \skt{वात\appbrk{}शो\appbrk{}णि\appbrk{}त\appbrk{}वि\appbrk{}ष\appbrk{}म\appbrk{}ज्व\appbrk{}रा\appbrk{}रो\appbrk{}च\appbrk{}क\appbrk{}पा\appbrk{}ण्डु\appbrk{}रो\appbrk{}गो\appbrk{}द\appbrk{}रा\appbrk{}र्शः\appbrk{}श्वा\appbrk{}स\appbrk{}का\appbrk{}स\appbrk{}शो\appbrk{}षा\appbrk{}ग्नि\appbrk{}सा\appbrk{}द\appbrk{}हृद्रो}\ H}}\skt{~। }\edtext{\skt{प्रदेहः}}{\lemma{\skt{प्रदेहः}}\Afootnote{\skt{जीव\appbrk{}नी\appbrk{}य\appbrk{}प्र\appbrk{}ती\appbrk{}वापं सर्पिः पयसा पाच\appbrk{}यि\appbrk{}त्वाऽ\appbrk{}भ्य\appbrk{}ज्यात्; सहासहदे}\ A; \skt{गानुपहन्ति~। प्रदेहः}\ H}}\skt{ }\edtext{\skt{सहा\-च\-न्द\-न\-मू\-र्व्वा\-पि\-या\-ल\-श\-ता\-व\-री\-क\-शे\-रु\-का\-स\-ह\-दे\-वा\-प\-द्म\-क\-म\-धु\-क\-श\-त\-पु\-ष्पा\-कु\-ष्ठानि}}{\lemma{\skt{सहा\appbrk{}च\appbrk{}न्द\appbrk{}न\appbrk{}मू\appbrk{}र्व्वा\appbrk{}पि\appbrk{}या\appbrk{}ल\appbrk{}श\appbrk{}ता\appbrk{}व\appbrk{}री\appbrk{}क\appbrk{}शे\appbrk{}रु\appbrk{}का\appbrk{}स\appbrk{}ह\appbrk{}दे\appbrk{}वा\appbrk{}प\appbrk{}द्म\appbrk{}क\appbrk{}म\appbrk{}धु\appbrk{}क\appbrk{}श\appbrk{}त\appbrk{}पु\appbrk{}ष्पा\appbrk{}कु\appbrk{}ष्ठानि}}\Afootnote{H; \skt{वाच\appbrk{}न्द\appbrk{}न\appbrk{}मू\appbrk{}र्वा\appbrk{}मु\appbrk{}स्ता\appbrk{}प्रि\appbrk{}या\appbrk{}ल\appbrk{}श\appbrk{}ता\appbrk{}व\appbrk{}री\appbrk{}क\appbrk{}शे\appbrk{}रु\appbrk{}प\appbrk{}द्म\appbrk{}क\appbrk{}म\appbrk{}धु\appbrk{}क\appbrk{}श\appbrk{}त\appbrk{}पु\appbrk{}ष्पा\appbrk{}वि\appbrk{}दा\appbrk{}री\appbrk{}कु\appbrk{}ष्ठानि}\ A}}\skt{ }\edtext{\skt{शैले\-य\-का\-ट\-रू\-ष\-क\-ब\-ला\-ति\-ब\-ला\-जी\-व\-न्ती\-कल्को वा}}{\lemma{\skt{शैले\appbrk{}य\appbrk{}का\appbrk{}ट\appbrk{}रू\appbrk{}ष\appbrk{}क\appbrk{}ब\appbrk{}ला\appbrk{}ति\appbrk{}ब\appbrk{}ला\appbrk{}जी\appbrk{}व\appbrk{}न्ती\appbrk{}कल्को वा}}\Afootnote{H; \skt{क्षीरपिष्टः प्रदेहो घृतमण्डयुक्तः, सैरे\appbrk{}य\appbrk{}का\appbrk{}ट\appbrk{}रू\appbrk{}ष\appbrk{}क\appbrk{}ब\appbrk{}ला\appbrk{}ति\appbrk{}ब\appbrk{}ला\appbrk{}जी\appbrk{}व\appbrk{}न्ती\appbrk{}सु\appbrk{}ष\appbrk{}वी\appbrk{}कल्को वा}\ A}}\skt{ }\edtext{\skt{क्षीरयुक्तः}}{\lemma{\skt{क्षीरयुक्तः}}\Afootnote{H; \skt{च्छागक्षीरपिष्टः, गोक्षीरपिष्टः}\ A}}\skt{~। काश्म\-र्य\-म\-धु\-क\-त\-र्प्प\-ण\-कल्को वा }\edtext{\skt{घृतमण्डयुक्तः}}{\lemma{\skt{घृतमण्डयुक्तः}}\Afootnote{H; \emph{om.}\ A}}\skt{~। }\edtext{\skt{मधू\-च्छि\-ष्ट\-शा\-रि\-वा\-स\-र्ज्ज\-र\-स\-म\-धु\-क\-म\-धु\-र\-व\-र्ग्ग\-क्षी\-र\-सि\-द्धम्वा}}{\lemma{\skt{मधू\appbrk{}च्छि\appbrk{}ष्ट\appbrk{}शा\appbrk{}रि\appbrk{}वा\appbrk{}स\appbrk{}र्ज्ज\appbrk{}र\appbrk{}स\appbrk{}म\appbrk{}धु\appbrk{}क\appbrk{}म\appbrk{}धु\appbrk{}र\appbrk{}व\appbrk{}र्ग्ग\appbrk{}क्षी\appbrk{}र\appbrk{}सि\appbrk{}द्धम्वा}}\Afootnote{H; \skt{मधू\appbrk{}च्छि\appbrk{}ष्ट\appbrk{}म\appbrk{}ञ्जि\appbrk{}ष्ठा\appbrk{}स\appbrk{}र्ज\appbrk{}र\appbrk{}स\appbrk{}सा\appbrk{}रि\appbrk{}वा\appbrk{}क्षी\appbrk{}र\appbrk{}सिद्धं}\ A}}\skt{ }\edtext{\edtext{\skt{पिण्डतैलमिति}}{\lemma{\skt{पिण्डतैलम्}}\Afootnote{H; \skt{पिण्ड\appbrk{}तै\appbrk{}ल\appbrk{}म\appbrk{}भ्यङ्गः}\ A}}}{\lemma{\skt{इति}}\Afootnote{H; \skt{सर्वेषु च}\ A}}\skt{~॥ }\edtext{\skt{पुरा\-ण\-घृ\-त\-मा\-म\-ल\-क\-स\-र\-ल\-वि\-पक्वं}}{\lemma{\skt{आम\appbrk{}ल\appbrk{}क\appbrk{}स\appbrk{}र\appbrk{}ल\appbrk{}वि\appbrk{}पक्वं}}\Afootnote{H\textsuperscript{ac}; \skt{आम\appbrk{}ल\appbrk{}क\appbrk{}र\appbrk{}स\appbrk{}वि\appbrk{}पक्वं}\ A; \skt{आम\appbrk{}ल\appbrk{}क\appbrk{}स\appbrk{}र\appbrk{}ल\appbrk{}वि\appbrk{}पक्व}\ H\textsuperscript{pc}}}\skt{ }\edtext{\skt{शर्क्क\-रा\-म\-धु\-म\-धुरं}}{\lemma{\skt{शर्क्क\appbrk{}रा\appbrk{}म\appbrk{}धु\appbrk{}म\appbrk{}धुरं}}\Afootnote{H\textsuperscript{pc}; \skt{वा}\ A; \emph{om.}\ H\textsuperscript{ac}}}\skt{ पानार्थे~। जीवनीयसिद्धं परिषेकार्थे~। }\edtext{\skt{सुष\-वी\-क्वा\-थ\-सिद्धं वा}}{\lemma{\skt{सुष\appbrk{}वी\appbrk{}क्वा\appbrk{}थ\appbrk{}सिद्धं वा}}\Afootnote{H; \skt{काको\appbrk{}ल्या\appbrk{}दि\appbrk{}क्वा\appbrk{}थ\appbrk{}क\appbrk{}ल्क\appbrk{}सिद्धं वा; सुष\appbrk{}वी\appbrk{}क्वा\appbrk{}थ\appbrk{}क\appbrk{}ल्क\appbrk{}सिद्धं वा, कार\appbrk{}वे\appbrk{}ल्ल\appbrk{}क\appbrk{}क्व\appbrk{}थ\appbrk{}मा\appbrk{}त्र\appbrk{}सिद्धं वा}\ A}}\skt{~। बलातैलं }\edtext{\skt{सिद्धं}}{\lemma{\skt{सिद्धं}}\Afootnote{H; \skt{वा}\ A}}\skt{ }\edtext{\edtext{\skt{परि\-षे\-का\-व\-गा\-ह\-ब\-स्ति\-भो\-ज\-ना\-र्थ\-मिति}}{\lemma{\skt{परि\appbrk{}षे\appbrk{}का\appbrk{}व\appbrk{}गा\appbrk{}ह\appbrk{}ब\appbrk{}स्ति\appbrk{}भो\appbrk{}ज\appbrk{}ना\appbrk{}र्थम्}}\Afootnote{H; \skt{परि\appbrk{}षे\appbrk{}का\appbrk{}व\appbrk{}गा\appbrk{}ह\appbrk{}ब\appbrk{}स्ति\appbrk{}भो\appbrk{}जने}\ A}}}{\lemma{\skt{इति}}\Afootnote{H; \skt{षु}\ A}}\skt{~। }\edtext{\edtext{\skt{शालि\-ष\-ष्टि\-क\-य\-व\-गो\-धू\-मा\-न्न\-मु\-प\-से\-वेत}}{\lemma{\skt{शालि\appbrk{}ष\appbrk{}ष्टि\appbrk{}क\appbrk{}य\appbrk{}व\appbrk{}गो\appbrk{}धू\appbrk{}मा\appbrk{}न्नम्}}\Afootnote{H; \skt{शालि\appbrk{}ष\appbrk{}ष्टि\appbrk{}क\appbrk{}य\appbrk{}व\appbrk{}गो\appbrk{}धू\appbrk{}मान्न}\ A}}}{\lemma{\skt{उपसेवेत}}\Afootnote{H; \skt{मनवं भुञ्जीत}\ A}}\skt{ पयसा }\edtext{\skt{जाङ्गलरसेन}}{\lemma{\skt{जाङ्गलरसेन}}\Afootnote{H; \skt{जाङ्गलरसेन वा}\ A}}\skt{ मुद्गयूषेण }\edtext{\skt{वानम्लेन}}{\lemma{\skt{वानम्लेन}}\Afootnote{H; \skt{वाऽनम्लेन}\ A}}\skt{~। }\edtext{\edtext{\skt{शोणि\-त\-मो\-क्ष\-ण\-ञ्चो\-च्छ्रि\-त\-दो\-ष\-ग्र\-ह\-णेन}}{\lemma{\skt{शोणि\appbrk{}त\appbrk{}मो\appbrk{}क्ष\appbrk{}णञ्}}\Afootnote{\skt{शोणितमोक्षं चाभीक्ष्णं}\ A; \skt{शोणितमोक्ष}\ H}}}{\lemma{\skt{चोच्छ्रि\appbrk{}त\appbrk{}दो\appbrk{}ष\appbrk{}ग्र\appbrk{}ह\appbrk{}णेन}}\Afootnote{\skt{कुर्वीत; उच्छितदोषे च}\ A; \skt{णञ्चो\appbrk{}च्छ्रि\appbrk{}त\appbrk{}दो\appbrk{}ष\appbrk{}ग्र\appbrk{}ह\appbrk{}णेन}\ H}}\skt{ }\edtext{\skt{वम\-न\-वि\-रे\-च\-ना\-स्था\-प\-ना\-नु\-वा\-स\-न\-क\-र्म्मा\-से\-वेत}}{\lemma{\skt{वम\appbrk{}न\appbrk{}वि\appbrk{}रे\appbrk{}च\appbrk{}ना\appbrk{}स्था\appbrk{}प\appbrk{}ना\appbrk{}नु\appbrk{}वा\appbrk{}स\appbrk{}न\appbrk{}क\appbrk{}र्म्मा\appbrk{}से\appbrk{}वेत}}\Afootnote{H; \skt{वम\appbrk{}न\appbrk{}वि\appbrk{}रे\appbrk{}च\appbrk{}ना\appbrk{}स्था\appbrk{}प\appbrk{}ना\appbrk{}नु\appbrk{}वा\appbrk{}स\appbrk{}न\appbrk{}कर्म कर्तव्यम्}\ A}}\skt{~॥}\pend

\pstart\versepar \ledsidenote{\passref{4.5.14}}\edtext{\skt{भवन्ति}}{\lemma{\skt{भवन्ति}}\Afootnote{A; \skt{भवति}\ H}}\skt{ चात्र~॥ }\pend
\pstart\versepar \skt{एवमाद्यैः क्रिया\-यो\-गै\-र\-चि\-रो\-त्प\-तितं सुखम्~। }\pend
\pstart\versepar \skt{वातासृक्साध्यते }\edtext{\skt{वैद्यैर्याप्यते च}}{\lemma{\skt{याप्यते च}}\Afootnote{H; \skt{याप्यते तु}\ A}}\skt{ चिरोत्थितं~॥}\pend

\pstart\versepar \ledsidenote{\passref{4.5.15}}\edtext{\skt{अभ्य\-ञ्ज\-न\-प\-री\-षे\-क\-प्र\-दे\-ह\-व्य\-ज\-ना\-निलाः}}{\lemma{\skt{अभ्य\appbrk{}ञ्ज\appbrk{}न\appbrk{}प\appbrk{}री\appbrk{}षे\appbrk{}क\appbrk{}प्र\appbrk{}दे\appbrk{}ह\appbrk{}व्य\appbrk{}ज\appbrk{}ना\appbrk{}निलाः}}\Afootnote{\skt{उप\appbrk{}ना\appbrk{}ह\appbrk{}प\appbrk{}री\appbrk{}षे\appbrk{}क\appbrk{}प्र\appbrk{}दे\appbrk{}हा\appbrk{}भ्य\appbrk{}ञ्ज\appbrk{}नानि च}\ A; \skt{अभ्य\appbrk{}ञ्ज\appbrk{}न\appbrk{}प\appbrk{}रि\appbrk{}षे\appbrk{}क\appbrk{}प्र\appbrk{}दे\appbrk{}ह\appbrk{}व्य\appbrk{}ज\appbrk{}ना\appbrk{}निलाः}\ H}}\skt{ }\pend
\pstart\versepar \skt{शर\-णा\-न्य\-प्र\-वा\-तानि मनोज्ञानि महान्ति च~॥}\pend

\pstart\versepar \ledsidenote{\passref{4.5.16}}\skt{मृदु\-ग\-ण्डो\-प\-धा\-नानि शयनानि सुखानि च~। }\pend
\pstart\versepar \skt{वातरक्ते प्रशस्यन्ते मृदु\-स\-म्वा\-ह\-नानि च~॥}\pend

\pstart\versepar \ledsidenote{\passref{4.5.17}}\skt{व्यायामम्मैथुनं कोप\-मु\-ष्णा\-म्ल\-ल\-व\-णा\-शनं~। }\pend
\pstart\versepar \skt{दिवा\-स्व\-प्न\-म\-भि\-ष्यन्दि गुरु }\edtext{\skt{चान्न\-म्वि\-व\-र्ज्ज\-येद्~॥ इति}}{\lemma{\skt{विवर्ज्जयेद्~॥ इति}}\Afootnote{H; \skt{विवर्जयेत्}\ A}}\skt{~॥}\pend

\pstart \ledsidenote{\passref{4.5.18}}\edtext{\edtext{\edtext{\edtext{\edtext{\edtext{\edtext{\skt{अथा\-प\-ता\-न\-कि\-न\-म\-स्र\-स्ता\-क्ष\-म\-व\-क्र\-भ्रु\-व\-म\-स्त\-ब्धा\-ङ्गु\-ल\-म\-स्वे\-द\-न\-म\-वे\-प\-न\-म\-प्र\-ला\-पि\-न\-म\-ख\-ट्वा\-पा\-ति\-न\-म\-ब\-हि\-रा\-या\-मि\-न\-ञ्चो\-प\-क्र\-मेत्}}{\lemma{\skt{अथा\appbrk{}प\appbrk{}ता\appbrk{}न\appbrk{}कि\appbrk{}नम्}}\Afootnote{\skt{अप\appbrk{}ता\appbrk{}न\appbrk{}कि\appbrk{}नम्}\ A; \skt{अथा\appbrk{}प\appbrk{}ता\appbrk{}कि\appbrk{}नम्}\ H}}}{\lemma{\skt{अस्तब्धाङ्गुलम्}}\Afootnote{H; \skt{अस्तब्धमेढ्रम्}\ A}}}{\lemma{\skt{अवेपनम्}}\Afootnote{A; \skt{अवे}\ H\textsuperscript{ac}; \skt{अवेम्}\ H\textsuperscript{pc}}}}{\lemma{\skt{अप्रलापिनम्}}\Afootnote{A; \skt{पमानम्}\ H\textsuperscript{ac}; \skt{अप्र\appbrk{}ला\appbrk{}पि\appbrk{}न\appbrk{}मिति पाठःपमानम्}\ H\textsuperscript{pc}}}}{\lemma{\skt{अख\appbrk{}ट्वा\appbrk{}पा\appbrk{}ति\appbrk{}नम्}}\Afootnote{A; \skt{अख\appbrk{}ट्वा\appbrk{}या\appbrk{}ति\appbrk{}नम्}\ H}}}{\lemma{\skt{अब\appbrk{}हि\appbrk{}रा\appbrk{}या\appbrk{}मि\appbrk{}नञ्}}\Afootnote{A; \skt{अब\appbrk{}हि\appbrk{}रा\appbrk{}या\appbrk{}मि\appbrk{}नञ्}\ H}}}{\lemma{\skt{चोपक्रमेत्}}\Afootnote{A; \emph{om.}\ H}}\skt{~। तत्र }\edtext{\skt{प्रागेव}}{\lemma{\skt{प्रागेव}}\Afootnote{A; \skt{प्रागैव}\ H}}\skt{ स्नेहाभ्यक्तं }\edtext{\skt{स्विन्न\-श\-री\-र\-म\-व\-पी\-डेन}}{\lemma{\skt{अवपीडेन}}\Afootnote{H; \skt{अवपीडनेन}\ A}}\skt{ }\edtext{\edtext{\skt{तीक्ष्णे\-नो\-प\-क्र\-मे\-त्शि\-रः\-शु\-द्ध्य\-र्थम्}}{\lemma{\skt{तीक्ष्णे\appbrk{}नो\appbrk{}प\appbrk{}क्र\appbrk{}मेत्}}\Afootnote{H; \skt{तीक्ष्णे\appbrk{}नो\appbrk{}प\appbrk{}क्र\appbrk{}मेत}\ A}}}{\lemma{\skt{शिरःशुद्ध्यर्थम्}}\Afootnote{H; \skt{शिरःशुद्ध्यर्थं}\ A}}\skt{~। }\edtext{\edtext{\edtext{\skt{अथा\-न\-न्त\-र\-म्वि\-दा\-रि\-ग\-न्धा\-दि\-क्वा\-थे\-क्षु\-र\-स\-क्षी\-र\-द\-धि\-वि\-पक्वं}}{\lemma{\skt{अथानन्तरम्}}\Afootnote{H; \skt{अनन्तरं}\ A}}}{\lemma{\skt{क्वाथेक्षुरस}}\Afootnote{H; \skt{क्वाथमांसरस}\ A}}}{\lemma{\skt{क्षीर\appbrk{}द\appbrk{}धि\appbrk{}वि\appbrk{}पक्वं}}\Afootnote{H; \skt{क्षीरदधिपक्वं}\ A}}\skt{ सर्प्पि\-र\-च्छ\-म्पा\-य\-येत्~। }\edtext{\skt{तथा}}{\lemma{\skt{तथा}}\Afootnote{H; \skt{तथा हि}\ A}}\skt{ नातिमात्रम्वायुः प्रसरति~। ततो }\edtext{\skt{भद्रदार्वादीनि}}{\lemma{\skt{भद्रदार्वादीनि}}\Afootnote{H; \skt{भद्रदार्वादि}\ A}}\skt{ }\edtext{\edtext{\skt{वातघ्नान्यन्यानि}}{\lemma{\skt{वातघ्नान्य्}}\Afootnote{H; \skt{वातघ्नगणम्}\ A}}}{\lemma{\skt{अन्यानि}}\Afootnote{H; \emph{om.}\ A}}\skt{ }\edtext{\skt{चाङ्गान्याहृत्य}}{\lemma{\skt{चाङ्गान्य्}}\Afootnote{H; \emph{om.}\ A}}\skt{ }\edtext{\edtext{\skt{यव\-को\-ल\-कु\-ल\-त्थांश्च}}{\lemma{\skt{यव\appbrk{}को\appbrk{}ल\appbrk{}कु\appbrk{}ल\appbrk{}त्थांश्}}\Afootnote{H; \skt{सय\appbrk{}व\appbrk{}को\appbrk{}ल\appbrk{}कु\appbrk{}लत्थं}\ A}}}{\lemma{\skt{च}}\Afootnote{H; \emph{om.}\ A}}\skt{ }\edtext{\skt{सानू\-पो\-द\-क\-मां\-सा\-न्ये\-कतः}}{\lemma{\skt{सानू\appbrk{}पो\appbrk{}द\appbrk{}क\appbrk{}मां\appbrk{}सान्य्}}\Afootnote{H; \skt{सानू\appbrk{}पौ\appbrk{}द\appbrk{}क\appbrk{}मांसं पञ्चवर्गम्}\ A}}\skt{ }\edtext{\skt{प्रक्वाथ्य तमादाय}}{\lemma{\skt{प्रक्वाथ्य तम्}}\Afootnote{A; \skt{क्वाथ्यम्}\ H}}\footnoteA{We emend to the vulgate's \emph{ prakvāthya tam ādāya} for sense. The passage as transmitted by the Nepalese witnesses remains puzzling.}\skt{ कषा\-य\-म\-म्ल\-क्षीरैः सहोन्मिश्र्य }\edtext{\skt{सर्प्पि\-स्तै\-ल\-व\-सा\-म\-ज्जाभिः}}{\lemma{\skt{सर्प्पि\appbrk{}स्तै\appbrk{}ल\appbrk{}व\appbrk{}सा\appbrk{}म\appbrk{}ज्जाभिः}}\Afootnote{H; \skt{सर्पि\appbrk{}स्तै\appbrk{}ल\appbrk{}व\appbrk{}सा\appbrk{}म\appbrk{}ज्जभिः}\ A}}\skt{ सह }\edtext{\skt{पचे\-त्म\-धु\-र\-क\-प्र\-ती\-वापं}}{\lemma{\skt{पचेत्}}\Afootnote{H\textsuperscript{pc}; \skt{विपचेन्}\ A; \skt{पचेत्चेत्}\ H\textsuperscript{ac}}}\skt{~। }\edtext{\edtext{\edtext{\skt{तदे\-त\-त्त्रै\-वृ\-त\-म\-प\-ता\-न\-कि\-ना\-म्प\-रि\-षे\-का\-भ्य\-ङ्ग\-पा\-न\-भो\-ज\-ना\-नु\-वा\-स\-न\-न\-स्येषु}}{\lemma{\skt{त्रैवृतम्}}\Afootnote{A; \skt{त्रिवृतम्}\ H}}}{\lemma{\skt{परिषेकाभ्यङ्ग}}\Afootnote{H; \skt{परि\appbrk{}षे\appbrk{}का\appbrk{}व\appbrk{}गा\appbrk{}हा\appbrk{}भ्यङ्ग}\ A}}}{\lemma{\skt{वासननस्येषु}}\Afootnote{A H\textsuperscript{ac}; \skt{वासनस्येषु}\ H\textsuperscript{pc}}}\skt{ विदध्याद्~। यथोक्तैश्च स्वेदविधानैः स्वेदयेत्~। }\edtext{\skt{बलीयसि}}{\lemma{\skt{बलीयसि}}\Afootnote{A H\textsuperscript{pc}; \skt{बलीलीयसि}\ H\textsuperscript{ac}}}\skt{ वाते }\edtext{\edtext{\skt{सुखो\-ष्ण\-प\-रि\-षेके}}{\lemma{\skt{सुखोष्ण}}\Afootnote{H; \skt{सुखो\appbrk{}ष्ण\appbrk{}तु\appbrk{}ष\appbrk{}बुस}\ A}}\skt{ पूर्ण्णे}}{\lemma{\skt{परिषेके पूर्ण्णे}}\Afootnote{\skt{करीषपूर्णे कूपे}\ A; \skt{परिषेकपूर्ण्णे}\ H}}\skt{ }\edtext{\skt{निखनेत्}}{\lemma{\skt{निखनेत्}}\Afootnote{H; \skt{निद\appbrk{}ध्या\appbrk{}दा\appbrk{}मु\appbrk{}खात्}\ A}}\footnoteA{The word \emph{kūpe} "trench" may have been introduced in the vulgate because of the similarity of this treatment with that of the Carakasaṃhitā's "trench sweat" described at Carakasaṃhitā 1.14.59cd--60ad.}\skt{~। तप्तायाम्वा }\edtext{\skt{रथकारचुल्ल्यां}}{\lemma{\skt{रथकारचुल्ल्यां}}\Afootnote{\skt{रथकारचुल्ल्यां तप्तायाम्वा शिला\appbrk{}या\appbrk{}म्सु\appbrk{}रा\appbrk{}प\appbrk{}रि\appbrk{}षि\appbrk{}क्तायां पला\appbrk{}श\appbrk{}द\appbrk{}ल\appbrk{}च्छ\appbrk{}न्नायां शाययेत्}\ A; \skt{रथाकारचुल्ल्यां}\ H}}\skt{ }\edtext{\skt{कृस\-र\-वै\-श\-वा\-र\-पा\-य\-सैर्व्वा}}{\lemma{\skt{कृस\appbrk{}र\appbrk{}वै\appbrk{}श\appbrk{}वा\appbrk{}र\appbrk{}पा\appbrk{}य\appbrk{}सैर्}}\Afootnote{H; \skt{कृश\appbrk{}रा\appbrk{}वे\appbrk{}श\appbrk{}वा\appbrk{}र\appbrk{}पा\appbrk{}य\appbrk{}सैर्}\ A}}\skt{ }\edtext{\edtext{\edtext{\edtext{\edtext{\skt{स्वेद\-ये\-न्मू\-ल\-को\-रु\-वू\-क\-स्फू\-र्ज्ज\-का\-र्क्क\-स\-प्त\-ला\-श\-ङ्खि\-नी\-स्व\-र\-स\-सि\-द्ध\-म\-प\-ता\-न\-कि\-म्प\-रि\-षे\-का\-दिषु}}{\lemma{\skt{स्वेदयेन्}}\Afootnote{H; \skt{स्वेदयेत्}\ A}}}{\lemma{\skt{वूकस्फूर्ज्ज}}\Afootnote{H; \skt{बूस्फूर्जार्ज}\ A}}}{\lemma{\skt{कार्क्क\appbrk{}स\appbrk{}प्त\appbrk{}ला\appbrk{}श\appbrk{}ङ्खि\appbrk{}नी\appbrk{}स्व\appbrk{}र\appbrk{}स\appbrk{}सि\appbrk{}द्धम्}}\Afootnote{H; \skt{कार्क\appbrk{}स\appbrk{}प्त\appbrk{}ला\appbrk{}श\appbrk{}ङ्खि\appbrk{}नी\appbrk{}स्व\appbrk{}र\appbrk{}स\appbrk{}सिद्धं तैलम्}\ A}}}{\lemma{\skt{अपतानकिम्}}\Afootnote{\skt{अपतानकिनां}\ A H}}}{\lemma{\skt{परिषेकादिषु}}\Afootnote{H; \skt{परेषेकादिषू}\ A}}\skt{ }\edtext{\edtext{\skt{योज्य\-म\-भु\-क्त\-व\-ता\-म्पी\-त\-म\-म्ल\-न्दधि}}{\lemma{\skt{योज्यम्}}\Afootnote{H; \skt{पयोज्यम्}\ A}}}{\lemma{\skt{अभुक्तवताम्}}\Afootnote{H; \skt{अभुक्तवता}\ A}}\skt{ }\edtext{\edtext{\skt{मरि\-च\-यु\-क्त\-म\-प\-ता\-न\-क\-म\-प\-हन्ति}}{\lemma{\skt{मरिचयुक्तम्}}\Afootnote{H; \skt{मरि\appbrk{}च\appbrk{}व\appbrk{}चा\appbrk{}यु\appbrk{}क्तम्}\ A}}}{\lemma{\skt{अपहन्ति}}\Afootnote{H; \skt{हन्ति}\ A}}\skt{~। }\edtext{\skt{सर्प्पि\-स्तै\-ल\-व\-सा\-मज्जा}}{\lemma{\skt{सर्प्पि\appbrk{}स्तै\appbrk{}ल\appbrk{}व\appbrk{}सा\appbrk{}मज्जा}}\Afootnote{H; \skt{तैल\appbrk{}स\appbrk{}र्पि\appbrk{}र्व\appbrk{}सा\appbrk{}क्षौ\appbrk{}द्राणि}\ A}}\skt{ वा }\edtext{\skt{एत\-च्छु\-द्ध\-वा\-ता\-प\-ता\-न\-क\-वि\-धि\-मुक्तं}}{\lemma{\skt{छुद्ध\appbrk{}वा\appbrk{}ता\appbrk{}प\appbrk{}ता\appbrk{}न\appbrk{}क\appbrk{}वि\appbrk{}धिम्}}\Afootnote{H; \skt{छुद्ध\appbrk{}वा\appbrk{}ता\appbrk{}प\appbrk{}ता\appbrk{}न\appbrk{}क\appbrk{}वि\appbrk{}धा\appbrk{}नम्}\ A}}\skt{~। संसृष्टे संसृष्टङ्कर्त्तव्यं~। वेगान्तरेषु }\edtext{\skt{चाव\-पी\-डा\-न्द\-द्यात्}}{\lemma{\skt{चावपीडान्}}\Afootnote{H; \skt{चावपीडं}\ A}}\skt{~। }\edtext{\edtext{\edtext{\edtext{\skt{ताम्र\-चू\-ड\-क\-र्क्क\-ट\-क\-कृ\-ष्ण\-म\-त्स्य\-शि\-शु\-मा\-र\-र\-सा\-ञ्चा\-वे\-क्ष्यते}}{\lemma{\skt{ताम्र\appbrk{}चू\appbrk{}ड\appbrk{}क\appbrk{}र्क्क\appbrk{}ट\appbrk{}क\appbrk{}कृ\appbrk{}ष्ण\appbrk{}मत्स्य}}\Afootnote{H; \skt{ताम्र\appbrk{}चू\appbrk{}ड\appbrk{}क\appbrk{}क\appbrk{}र्क\appbrk{}ट\appbrk{}कृ\appbrk{}ष्ण\appbrk{}मत्स्य}\ A}}}{\lemma{\skt{शिशुमार}}\Afootnote{H; \skt{शिशु\appbrk{}मा\appbrk{}र\appbrk{}व\appbrk{}राह}\ A}}}{\lemma{\skt{रसाञ्}}\Afootnote{H; \skt{वसाश्}\ A}}}{\lemma{\skt{चावेक्ष्यते}}\Afootnote{\skt{चासेवेत}\ A; \skt{चावेक्ष्येत}\ H}}\skt{~। }\edtext{\skt{क्षीरानि}}{\lemma{\skt{क्षीरानि}}\Afootnote{H; \skt{क्षीराणि}\ A}}\skt{ }\edtext{\skt{वात\-ह\-र\-सि\-द्धानि}}{\lemma{\skt{वात\appbrk{}ह\appbrk{}र\appbrk{}सि\appbrk{}द्धानि}}\Afootnote{H; \skt{वावा\appbrk{}त\appbrk{}ह\appbrk{}र\appbrk{}सि\appbrk{}द्धानि}\ A}}\skt{~। }\edtext{\skt{यव\-को\-ल\-कु\-ल\-त्थ\-मू\-ल\-क\-द\-धि\-घृ\-त\-तै\-ल\-सिद्धां}}{\lemma{\skt{यव\appbrk{}को\appbrk{}ल\appbrk{}कु\appbrk{}ल\appbrk{}त्थ\appbrk{}मू\appbrk{}ल\appbrk{}क\appbrk{}द\appbrk{}धि\appbrk{}घृ\appbrk{}त\appbrk{}तै\appbrk{}ल\appbrk{}सिद्धां}}\Afootnote{H; \skt{यव\appbrk{}को\appbrk{}ल\appbrk{}कु\appbrk{}ल\appbrk{}त्थ\appbrk{}मू\appbrk{}ल\appbrk{}क\appbrk{}द\appbrk{}धि\appbrk{}घृ\appbrk{}त\appbrk{}तै\appbrk{}ल\appbrk{}सिद्धा वा}\ A}}\skt{ }\edtext{\skt{यवागूं}}{\lemma{\skt{यवागूं}}\Afootnote{H; \skt{यवागूः}\ A}}\skt{~। }\edtext{\edtext{\edtext{\edtext{\edtext{\skt{स्नेह\-वि\-रे\-च\-ना\-स्था\-प\-ना\-नु\-वा\-स\-नै\-श्चे\-न\-न्द\-श\-रा\-त्र\-मा\-ग\-त\-वे\-ग\-मु\-प\-च\-रेत्}}{\lemma{\skt{स्नेह\appbrk{}वि\appbrk{}रे\appbrk{}च\appbrk{}ना\appbrk{}स्था\appbrk{}प\appbrk{}ना\appbrk{}नु\appbrk{}वा\appbrk{}स\appbrk{}नैश्}}\Afootnote{H; \skt{स्नेह\appbrk{}वि\appbrk{}रे\appbrk{}च\appbrk{}ना\appbrk{}स्था\appbrk{}प\appbrk{}ना\appbrk{}नु\appbrk{}वा\appbrk{}स\appbrk{}नश्}\ A}}}{\lemma{\skt{चेनन्}}\Afootnote{\skt{चैनं}\ A H}}}{\lemma{\skt{दशरात्रम्}}\Afootnote{H; \skt{दशरात्}\ A}}}{\lemma{\skt{आगतवेगम्}}\Afootnote{H; \skt{राहृतवेगम्}\ A}}}{\lemma{\skt{उपचरेत्}}\Afootnote{H; \skt{उपक्रमेत}\ A}}\skt{~। }\edtext{\skt{वात\-व्या\-धि\-चि\-कि\-त्सि\-त\-ञ्चा\-वे\-क्ष्येत}}{\lemma{\skt{चावेक्ष्येत}}\Afootnote{H; \skt{चावेक्षेत}\ A}}\skt{~। रक्षाकर्म्म च }\edtext{\skt{कुर्यात्}}{\lemma{\skt{कुर्यात्}}\Afootnote{H; \skt{कुर्यादिति}\ A}}\skt{~।}\pend

\pstart \ledsidenote{\passref{4.5.19}}\skt{पक्षा\-घा\-तो\-प\-द्रु\-त\-म\-म्ला\-न\-गात्रं }\edtext{\skt{सरु\-ज\-मा\-त्म\-व\-न्त\-ञ्चो\-प\-क्र\-मेत्}}{\lemma{\skt{आत्मवन्तञ्}}\Afootnote{H; \skt{आत्म\appbrk{}व\appbrk{}न्त\appbrk{}मु\appbrk{}प\appbrk{}क\appbrk{}र\appbrk{}ण\appbrk{}वन्तं}\ A}}\skt{~। तत्र प्रागेव स्नेह\-स्वे\-दो\-प\-पन्नं~। मृदुना }\edtext{\skt{शोधनेन}}{\lemma{\skt{शोधनेन}}\Afootnote{A; \skt{संशोधनेन}\ H}}\skt{ संशो\-ध्या\-नु\-वा\-स्या\-स्थाप्य च यथा\-का\-ल\-मा\-क्षे\-प\-क\-वि\-धा\-ने\-नो\-प\-च\-रेत्~। }\edtext{\edtext{\skt{वैशे\-षि\-क\-म\-स्ति\-ष्क\-शि\-रो\-ब\-स्ति\-भि\-रिति}}{\lemma{\skt{वैशे\appbrk{}षि\appbrk{}क\appbrk{}म\appbrk{}स्ति\appbrk{}ष्क\appbrk{}शि\appbrk{}रो\appbrk{}ब\appbrk{}स्ति\appbrk{}भिर्}}\Afootnote{H; \skt{वैशेषिकश्चात्र मस्तिष्क्यः शिरोबस्}\ A}}}{\lemma{\skt{इति}}\Afootnote{H; \skt{तिः}\ A}}\skt{~। }\edtext{\skt{अणु\-तै\-ल\-म\-भ्य\-ङ्गार्थे}}{\lemma{\skt{अणुतैलम्}}\Afootnote{A; \skt{अनुतैलम्}\ H}}\skt{~। }\edtext{\skt{साल्व\-ल\-मु\-प\-ना\-हार्थे}}{\lemma{\skt{साल्वलम्}}\Afootnote{H; \skt{साल्वणम्}\ A}}\skt{~। बला\-तै\-ल\-म\-नु\-वा\-स\-नार्थे~। }\edtext{\edtext{\edtext{\skt{एव\-म\-त\-न्द्रि\-त\-स्त्रीं\-श्च\-तुरो}}{\lemma{\skt{एवम्}}\Afootnote{H; \skt{एव}\ A}}}{\lemma{\skt{अतन्द्रितस्}}\Afootnote{H; \skt{मतन्द्रितस्}\ A}}\skt{ वा}}{\lemma{\skt{चतुरो वा}}\Afootnote{A; \skt{चतुरो}\ H}}\skt{ }\edtext{\edtext{\skt{मासा\-ङ्क्रि\-या\-मा\-से\-वेत}}{\lemma{\skt{मासाङ्क्रियाम्}}\Afootnote{H; \skt{मासान्क्रिया}\ A}}}{\lemma{\skt{आसेवेत}}\Afootnote{H; \skt{पथ\appbrk{}मु\appbrk{}प\appbrk{}से\appbrk{}वेत}\ A}}\skt{~॥}\pend

\pstart \ledsidenote{\passref{4.5.20}}\skt{मन्या\-स्त\-म्भेऽ\-प्ये\-त\-देव विधानं, विशेषतो वात\-श्ले\-ष्म\-ह\-रै\-र्नस्यै रूक्ष\-स्वे\-दै\-श्चो\-प\-च\-रेत्~॥}\pend

\pstart \ledsidenote{\passref{4.5.22}}\edtext{\skt{अर्द्दि\-ता\-तु\-र\-म्ब\-ल\-व\-न्त\-मु\-प\-क\-र\-ण\-व\-न्तञ्च}}{\lemma{\skt{बलवन्तम्}}\Afootnote{H; \skt{बल\appbrk{}व\appbrk{}न्त\appbrk{}मा\appbrk{}त्म\appbrk{}व\appbrk{}न्तम्}\ A}}\skt{ वात\-व्या\-धि\-वि\-धा\-ने\-नो\-प\-च\-रेत्~। }\edtext{\skt{वैशेषिकश्च}}{\lemma{\skt{वैशेषिकश्च}}\Afootnote{H; \skt{वैशेषिकैश्च}\ A}}\skt{ }\edtext{\skt{मस्ति\-ष्क\-शि\-रो\-ब\-स्ति\-न\-स्य\-धू\-पो\-प\-ना\-ह\-ना\-डी\-स्वे\-दैस्}}{\lemma{\skt{मस्ति\appbrk{}ष्क\appbrk{}शि\appbrk{}रो\appbrk{}ब\appbrk{}स्ति\appbrk{}न\appbrk{}स्य\appbrk{}धू\appbrk{}पो\appbrk{}प\appbrk{}ना\appbrk{}ह\appbrk{}ना\appbrk{}डी\appbrk{}स्वे\appbrk{}दैस्}}\Afootnote{H; \skt{मस्ति\appbrk{}ष्क्य\appbrk{}शि\appbrk{}रो\appbrk{}ब\appbrk{}स्ति\appbrk{}न\appbrk{}स्य\appbrk{}धू\appbrk{}मो\appbrk{}प\appbrk{}ना\appbrk{}ह\appbrk{}स्ने\appbrk{}ह\appbrk{}ना\appbrk{}डी\appbrk{}स्वे\appbrk{}दा\appbrk{}दिभिः}\ A}}\skt{~। ततः }\edtext{\edtext{\skt{सतृ\-ण\-म\-ह\-त्प\-ञ्च\-मू\-ली\-मा\-हृत्य}}{\lemma{\skt{सतृ\appbrk{}ण\appbrk{}म\appbrk{}ह\appbrk{}त्प\appbrk{}ञ्च\appbrk{}मू\appbrk{}लीम्}}\Afootnote{H; \skt{सतृणं महापञ्चमूलं}\ A}}}{\lemma{\skt{आहृत्य}}\Afootnote{H; \skt{काकोल्यादिं विदा\appbrk{}रि\appbrk{}ग\appbrk{}न्धा\appbrk{}दि\appbrk{}मौ\appbrk{}द\appbrk{}का\appbrk{}नू\appbrk{}प\appbrk{}मांसं तथै\appbrk{}वौ\appbrk{}द\appbrk{}क\appbrk{}क\appbrk{}न्दां\appbrk{}श्चा\appbrk{}हृत्य}\ A}}\skt{ द्विगुणोदके }\edtext{\skt{क्षीरे}}{\lemma{\skt{क्षीरे}}\Afootnote{H; \skt{क्षीरद्रोणे}\ A}}\skt{ }\edtext{\skt{क्वाथ्य\-क्षी\-रा\-व\-शि\-ष्ट\-म\-व\-तार्य}}{\lemma{\skt{क्वाथ्य\appbrk{}क्षी\appbrk{}रा\appbrk{}व\appbrk{}शि\appbrk{}ष्टम्}}\Afootnote{H; \skt{निःक्वाथ्य क्षीरावशिष्टम्}\ A}}\skt{ परिस्राव्य तैल\-प्र\-स्थे\-नो\-न्मिश्र्य }\edtext{\skt{पुन\-र\-ग्ना\-व\-धि\-स्रित्य}}{\lemma{\skt{अधिस्रित्य}}\Afootnote{H; \skt{अधिश्}\ A}}\skt{ }\edtext{\skt{विप\-चे\-त्त\-त\-स्तैलं}}{\lemma{\skt{विपचेत्}}\Afootnote{H; \skt{रगेत्}\ A}}\skt{ क्षीरा\-नु\-ग\-त\-म\-व\-तार्य }\edtext{\skt{शीती\-भू\-त\-म्म\-थ्नी\-यात्}}{\lemma{\skt{मथ्नीयात्}}\Afootnote{H; \skt{अभिमथ्नीयात्}\ A}}\skt{~। }\edtext{\edtext{\edtext{\edtext{\edtext{\skt{तदे\-त\-त्क्षी\-र\-तै\-ल\-म\-र्द्दि\-ता\-ना\-म्पा\-ना\-दि\-षू\-प\-यो\-ज्य\-मिति}}{\lemma{\skt{तदेतत्}}\Afootnote{H; \skt{तत्र यः स्नेह उत्ति\appbrk{}ष्ठे\appbrk{}त्त\appbrk{}मा\appbrk{}दाय मधु\appbrk{}रौ\appbrk{}ष\appbrk{}ध\appbrk{}स\appbrk{}हा\appbrk{}क्षी\appbrk{}र\appbrk{}युक्तं विपचेत्, एतत्}\ A}}}{\lemma{\skt{क्षीरतैलम्}}\Afootnote{H; \skt{क्षीरतैल}\ A}}}{\lemma{\skt{अर्द्दितानाम्}}\Afootnote{H; \skt{मर्दितातुराणां}\ A}}}{\lemma{\skt{पाना\appbrk{}दि\appbrk{}षू\appbrk{}प\appbrk{}यो\appbrk{}ज्यम्}}\Afootnote{H; \skt{पाना\appbrk{}भ्य\appbrk{}ङ्गा\appbrk{}दि\appbrk{}षू\appbrk{}प\appbrk{}योज्यं}\ A}}}{\lemma{\skt{इति}}\Afootnote{H; \skt{तैलहीनं वा क्षीर\appbrk{}स\appbrk{}र्पि\appbrk{}र\appbrk{}क्षि\appbrk{}त\appbrk{}र्प\appbrk{}ण\appbrk{}मिति}\ A}}\skt{~।}\pend

\pstart \ledsidenote{\passref{4.5.23}}\edtext{\skt{गृध्र\-सी\-वि\-श्च\-ञ्ची\-क्रो\-ष्टु\-क\-शी\-र्ष\-प\-ङ्गु\-क\-ला\-य\-ख\-ञ्ज\-वा\-त\-क\-ण्ट\-क\-पा\-द\-दा\-ह\-पा\-द\-हर्ष}}{\lemma{\skt{गृध्र\appbrk{}सी\appbrk{}वि\appbrk{}श्च\appbrk{}ञ्ची\appbrk{}क्रो\appbrk{}ष्टु\appbrk{}क\appbrk{}शी\appbrk{}र्ष\appbrk{}प\appbrk{}ङ्गु\appbrk{}क\appbrk{}ला\appbrk{}य\appbrk{}ख\appbrk{}ञ्ज\appbrk{}वा\appbrk{}त\appbrk{}क\appbrk{}ण्ट\appbrk{}क\appbrk{}पा\appbrk{}द\appbrk{}दा\appbrk{}ह\appbrk{}पा\appbrk{}द\appbrk{}हर्ष}}\Afootnote{H; \skt{घृध्र\appbrk{}सी\appbrk{}वि\appbrk{}श्वा\appbrk{}ची\appbrk{}क्रो\appbrk{}ष्टु\appbrk{}क\appbrk{}शि\appbrk{}रः\appbrk{}ख\appbrk{}ञ्ज\appbrk{}प\appbrk{}ङ्गु\appbrk{}ल\appbrk{}वा\appbrk{}त\appbrk{}क\appbrk{}ण्टा\appbrk{}क\appbrk{}पा\appbrk{}द\appbrk{}दा\appbrk{}ह\appbrk{}पा\appbrk{}द\appbrk{}हर्ष}\ A}}\skt{ }\edtext{\skt{अव\-बा\-हु\-क\-बा\-धि\-र्य\-ध\-म\-नी\-वा\-तेषु}}{\lemma{\skt{अव\appbrk{}बा\appbrk{}हु\appbrk{}क\appbrk{}बा\appbrk{}धि\appbrk{}र्य\appbrk{}ध\appbrk{}म\appbrk{}नी\appbrk{}वा\appbrk{}तेषु}}\Afootnote{H; \skt{वबा\appbrk{}हु\appbrk{}क\appbrk{}बा\appbrk{}धि\appbrk{}र्य\appbrk{}ध\appbrk{}म\appbrk{}नी\appbrk{}ग\appbrk{}त\appbrk{}वा\appbrk{}त\appbrk{}रो\appbrk{}गेषु}\ A}}\skt{ }\edtext{\skt{यथो\-क्त\-सि\-रा\-व्यधं}}{\lemma{\skt{यथो\appbrk{}क्त\appbrk{}सि\appbrk{}रा\appbrk{}व्यधं}}\Afootnote{H; \skt{यथोक्तं यथोद्देशं च सिराव्यधं}\ A}}\skt{ }\edtext{\edtext{\skt{कृत्वा\-न्य\-त्रा\-व\-बा\-हु\-का\-द्वा\-त\-व्या\-धि\-चि\-कि\-त्सि\-त\-म\-वे\-क्ष्येत}}{\lemma{\skt{कृत्वा\appbrk{}न्य\appbrk{}त्रा\appbrk{}व\appbrk{}बा\appbrk{}हु\appbrk{}काद्}}\Afootnote{H; \skt{कुर्यात् , अन्य\appbrk{}त्रा\appbrk{}व\appbrk{}बा\appbrk{}हु\appbrk{}कात्}\ A}}}{\lemma{\skt{अवेक्ष्येत}}\Afootnote{H; \skt{चावेक्षेत}\ A}}\skt{~।}\pend

\pstart \ledsidenote{\passref{4.5.24}}\edtext{\skt{कर्ण्णमूले तु}}{\lemma{\skt{कर्ण्णमूले तु}}\Afootnote{H; \skt{कर्णशूले तु}\ A}}\skt{ }\edtext{\edtext{\skt{शृङ्ग\-वे\-र\-र\-सै\-र्म्म\-धु\-क\-तै\-ल\-सं\-सृष्टः}}{\lemma{\skt{शृङ्ग\appbrk{}वे\appbrk{}र\appbrk{}र\appbrk{}सैर्}}\Afootnote{H; \skt{शृङ्गवेररसं}\ A}}}{\lemma{\skt{म्मधु\appbrk{}क\appbrk{}तै\appbrk{}ल\appbrk{}सं\appbrk{}सृष्टः}}\Afootnote{H; \skt{तैलम्धुसंसृष्टं}\ A}}\skt{ }\edtext{\skt{सैन्धवोपहितः}}{\lemma{\skt{सैन्धवोपहितः}}\Afootnote{H; \skt{सैधवोपहितं}\ A}}\skt{ }\edtext{\skt{सुखोष्णः}}{\lemma{\skt{सुखोष्णः}}\Afootnote{H; \skt{सुखोष्णं}\ A}}\skt{ }\edtext{\edtext{\skt{कर्ण्णयोर्द्देयः}}{\lemma{\skt{कर्ण्णयोर्}}\Afootnote{H; \skt{कर्णे द}\ A}}}{\lemma{\skt{द्देयः}}\Afootnote{H; \skt{द्या}\ A}}\skt{~। }\edtext{\skt{अज\-मू\-त्र\-म\-धु\-क\-तै\-लानि वा}}{\lemma{\skt{अज\appbrk{}मू\appbrk{}त्र\appbrk{}म\appbrk{}धु\appbrk{}क\appbrk{}तै\appbrk{}लानि वा}}\Afootnote{H; \skt{त्, अजा\appbrk{}मू\appbrk{}त्र\appbrk{}म\appbrk{}धु\appbrk{}तै\appbrk{}लानि वा}\ A}}\skt{~। }\edtext{\skt{मातु\-लु\-ङ्ग\-दा\-डि\-म\-ति\-न्ति\-डी\-क\-र\-स\-मू\-त्र\-सि\-द्धम्वा}}{\lemma{\skt{मातु\appbrk{}लु\appbrk{}ङ्ग\appbrk{}दा\appbrk{}डि\appbrk{}म\appbrk{}ति\appbrk{}न्ति\appbrk{}डी\appbrk{}क\appbrk{}र\appbrk{}स\appbrk{}मू\appbrk{}त्र\appbrk{}सि\appbrk{}द्धम्वा}}\Afootnote{H; \skt{मातु\appbrk{}लु\appbrk{}ङ्ग\appbrk{}दा\appbrk{}डि\appbrk{}म\appbrk{}ति\appbrk{}न्ति\appbrk{}डी\appbrk{}क\appbrk{}स्व\appbrk{}र\appbrk{}स\appbrk{}मू\appbrk{}त्र\appbrk{}सिद्धं}\ A}}\skt{ }\edtext{\skt{तैलम्}}{\lemma{\skt{तैलम्}}\Afootnote{A; \skt{तैल}\ H}}\skt{~। }\edtext{\skt{सुक्त\-सु\-रा\-त\-क्र\-मू\-त्र\-सि\-द्धम्वा}}{\lemma{\skt{सुक्त\appbrk{}सु\appbrk{}रा\appbrk{}त\appbrk{}क्र\appbrk{}मू\appbrk{}त्र\appbrk{}सि\appbrk{}द्धम्वा}}\Afootnote{H; \skt{शुक्त\appbrk{}सु\appbrk{}रा\appbrk{}त\appbrk{}क्र\appbrk{}मू\appbrk{}त्र\appbrk{}ल\appbrk{}व\appbrk{}ण\appbrk{}सिद्धं वा}\ A}}\skt{ नाडीस्वेदैश्च स्वेदयेत्~। }\edtext{\skt{वात\-व्या\-धि\-चि\-कि\-त्सा\-ञ्चा\-वे\-क्ष्येत}}{\lemma{\skt{चावेक्ष्येत}}\Afootnote{H; \skt{चावेक्षेत}\ A}}\skt{~॥ भूयश्चोत्तरे वक्ष्यामः~।}\pend

\pstart \ledsidenote{\passref{4.5.25}}\edtext{\skt{तूनीप्रतितून्योः}}{\lemma{\skt{तूनीप्रतितून्योः}}\Afootnote{H; \skt{तूनीप्रतून्योः}\ A}}\skt{ }\edtext{\skt{स्नेहं}}{\lemma{\skt{स्नेहं}}\Afootnote{H; \skt{स्नेह}\ A}}\skt{ लव\-ण\-मु\-ष्णो\-द\-केन पाययेत्~। पिप्प\-ल्या\-दि\-चू\-र्ण्णम्वा~। हिङ्गु\-य\-व\-क्षा\-र\-प्र\-गा\-ढम्वा सर्प्पिः~। }\edtext{\skt{बस्ति\-भि\-श्चो\-प\-क्र\-मेत्}}{\lemma{\skt{चोपक्रमेत्}}\Afootnote{H; \skt{चैन\appbrk{}मु\appbrk{}प\appbrk{}क्र\appbrk{}मेत्}\ A}}\skt{~।}\pend

\pstart \ledsidenote{\passref{4.5.26}}\skt{आध्माने }\edtext{\skt{त्ववतर्प्पणं}}{\lemma{\skt{त्ववतर्प्पणं}}\Afootnote{H; \skt{त्वपर्पण}\ A}}\skt{ }\edtext{\skt{पाणि\-ता\-प\-फ\-ल\-व\-र्त्ति\-क्रि\-या\-दी\-प\-न\-पा\-च\-नानि}}{\lemma{\skt{पाणि\appbrk{}ता\appbrk{}प\appbrk{}फ\appbrk{}ल\appbrk{}व\appbrk{}र्त्ति\appbrk{}क्रि\appbrk{}या\appbrk{}दी\appbrk{}प\appbrk{}न\appbrk{}पा\appbrk{}च\appbrk{}नानि}}\Afootnote{H; \skt{पाणि\appbrk{}ता\appbrk{}प(\appbrk{}दी\appbrk{}प\appbrk{}न\appbrk{}चू\appbrk{}र्ण)\appbrk{}फ\appbrk{}ल\appbrk{}व\appbrk{}र्ति\appbrk{}क्रि\appbrk{}या\appbrk{}पा\appbrk{}च\appbrk{}नी\appbrk{}य\appbrk{}दी\appbrk{}प\appbrk{}नी\appbrk{}य\appbrk{}ब\appbrk{}स्ति\appbrk{}भि\appbrk{}रु\appbrk{}पा\appbrk{}च\appbrk{}रेत्; लङ्घनानन्तरं चान्नकाले धान्य\appbrk{}क\appbrk{}जी\appbrk{}र\appbrk{}का\appbrk{}दि\appbrk{}दी\appbrk{}प\appbrk{}न\appbrk{}सि\appbrk{}द्धा\appbrk{}न्य\appbrk{}न्नानि}\ A}}\skt{ }\edtext{\edtext{\skt{कुर्य\-च्छो\-ध\-न\-ब\-स्त्यु\-प\-यो\-गश्च}}{\lemma{\skt{कुर्यच्}}\Afootnote{H; \emph{om.}\ A}}}{\lemma{\skt{छोध\appbrk{}न\appbrk{}ब\appbrk{}स्त्यु\appbrk{}प\appbrk{}यो\appbrk{}गश्च}}\Afootnote{H; \emph{om.}\ A}}\skt{~। प्रत्याध्माने }\edtext{\skt{छर्द्द\-ना\-प\-त\-र्प्प\-ण\-दी\-प\-नी\-यानि}}{\lemma{\skt{छर्द्द\appbrk{}ना\appbrk{}प\appbrk{}त\appbrk{}र्प्प\appbrk{}ण\appbrk{}दी\appbrk{}प\appbrk{}नी\appbrk{}यानि}}\Afootnote{H; \skt{छर्द\appbrk{}ना\appbrk{}प\appbrk{}त\appbrk{}र्प\appbrk{}ण\appbrk{}दी\appbrk{}प\appbrk{}नानि कुर्यात्}\ A}}\skt{~।}\pend

\pstart \ledsidenote{\passref{4.5.27}}\edtext{\skt{अष्ठी\-ला\-प्र\-त्य\-ष्ठी\-ल\-यो\-र्ग्गु\-ल्मा\-भ्य\-न्त\-र\-वि\-द्र\-धि\-क्रि\-या\-वि\-भाग}}{\lemma{\skt{ग्गुल्मा\appbrk{}भ्य\appbrk{}न्त\appbrk{}र\appbrk{}वि\appbrk{}द्र\appbrk{}धि\appbrk{}क्रि\appbrk{}या\appbrk{}वि\appbrk{}भाग}}\Afootnote{H; \skt{गुल्मा\appbrk{}भ्य\appbrk{}न्त\appbrk{}र\appbrk{}वि\appbrk{}द्र\appbrk{}धि\appbrk{}व\appbrk{}त्क्रि\appbrk{}या\appbrk{}वि\appbrk{}भाग}\ A}}\skt{ इति~॥}\pend

\pstart \ledsidenote{\passref{4.5.28}}\edtext{\skt{हिङ्गु\-त्रि\-क\-टु\-क\-व\-चा\-ज\-मो\-द\-धा\-न्या\-ज\-ग\-न्धा\-दा\-डि\-म\-ति\-न्ति\-डी\-क\-पा\-ठा\-चि\-त्त्र\-क\-च\-व्या\-सै\-न्ध\-व\-वि\-ड\-सौ\-व\-र्च्च\-ल\-य\-व\-क्षा\-र\-सु\-व\-र्च्चि\-का\-पि\-प्प\-ली\-मू\-ला\-म्ल\-वे\-त\-स\-श\-ठी\-पु\-ष्क\-र\-ह\-पु\-षा\-जा\-जी\-पथ्याः}}{\lemma{\skt{हिङ्गु\appbrk{}त्रि\appbrk{}क\appbrk{}टु\appbrk{}क\appbrk{}व\appbrk{}चा\appbrk{}ज\appbrk{}मो\appbrk{}द\appbrk{}धा\appbrk{}न्या\appbrk{}ज\appbrk{}ग\appbrk{}न्धा\appbrk{}दा\appbrk{}डि\appbrk{}म\appbrk{}ति\appbrk{}न्ति\appbrk{}डी\appbrk{}क\appbrk{}पा\appbrk{}ठा\appbrk{}चि\appbrk{}त्त्र\appbrk{}क\appbrk{}च\appbrk{}व्या\appbrk{}सै\appbrk{}न्ध\appbrk{}व\appbrk{}वि\appbrk{}ड\appbrk{}सौ\appbrk{}व\appbrk{}र्च्च\appbrk{}ल\appbrk{}य\appbrk{}व\appbrk{}क्षा\appbrk{}र\appbrk{}सु\appbrk{}व\appbrk{}र्च्चि\appbrk{}का\appbrk{}पि\appbrk{}प्प\appbrk{}ली\appbrk{}मू\appbrk{}ला\appbrk{}म्ल\appbrk{}वे\appbrk{}त\appbrk{}स\appbrk{}श\appbrk{}ठी\appbrk{}पु\appbrk{}ष्क\appbrk{}र\appbrk{}ह\appbrk{}पु\appbrk{}षा\appbrk{}जा\appbrk{}जी\appbrk{}पथ्याः}}\Afootnote{H; \skt{हिङ्गु\appbrk{}त्रि\appbrk{}क\appbrk{}टु\appbrk{}व\appbrk{}चा\appbrk{}ज\appbrk{}मो\appbrk{}दा\appbrk{}धा\appbrk{}न्या\appbrk{}ज\appbrk{}ग\appbrk{}न्धा\appbrk{}दा\appbrk{}डि\appbrk{}म\appbrk{}ति\appbrk{}न्ति\appbrk{}डी\appbrk{}क\appbrk{}पा\appbrk{}ठा\appbrk{}चि\appbrk{}त्र\appbrk{}क\appbrk{}य\appbrk{}व\appbrk{}क्षा\appbrk{}र\appbrk{}सै\appbrk{}न्ध\appbrk{}व\appbrk{}वि\appbrk{}ड\appbrk{}सौ\appbrk{}र्व\appbrk{}च\appbrk{}ल\appbrk{}स्व\appbrk{}र्जि\appbrk{}का\appbrk{}पि\appbrk{}प्प\appbrk{}ली\appbrk{}मू\appbrk{}ला\appbrk{}म्ल\appbrk{}वे\appbrk{}त\appbrk{}स\appbrk{}श\appbrk{}टी\appbrk{}पु\appbrk{}ष्क\appbrk{}र\appbrk{}मू\appbrk{}ल\appbrk{}ह\appbrk{}पु\appbrk{}षा\appbrk{}च\appbrk{}व्या\appbrk{}जा\appbrk{}जी\appbrk{}प\appbrk{}थ्याश्}\ A}}\skt{ }\edtext{\skt{सञ्चूर्ण्ण्य}}{\lemma{\skt{सञ्चूर्ण्ण्य}}\Afootnote{H; \skt{चूर्णयि}\ A}}\skt{ }\edtext{\skt{मातुलुङ्गाम्लेन}}{\lemma{\skt{मातुलुङ्गाम्लेन}}\Afootnote{H; \skt{त्वा मातुलुङ्गालेन}\ A}}\skt{ बहुशः परि\-भा\-व्या\-क्ष\-मात्रा }\edtext{\skt{गुडिकाः}}{\lemma{\skt{गुडिकाः}}\Afootnote{H; \skt{गुटिकाः}\ A}}\skt{ कारयेत्ततः }\edtext{\skt{प्रात\-रे\-कै\-का\-म्वा\-त\-रोगी}}{\lemma{\skt{वातरोगी}}\Afootnote{H; \skt{वातविकारी}\ A}}\skt{ भक्षयेदेष }\edtext{\skt{खलु}}{\lemma{\skt{खलु}}\Afootnote{H; \emph{om.}\ A}}\skt{ योगो }\edtext{\skt{गुल्म\-श्वा\-स\-का\-सा\-रो\-च\-क\-हृ\-द्रो\-गा\-ध्मा\-न\-पा\-र्श्वो\-द\-र\-ब\-स्ति\-शू\-ला\-ना\-ह\-मू\-त्र\-कृ\-च्छ्रा\-र्शः\-प्ली\-हो\-द\-र\-पा\-ण्डु\-रो\-गा\-न\-प\-हन्ति}}{\lemma{\skt{गुल्म\appbrk{}श्वा\appbrk{}स\appbrk{}का\appbrk{}सा\appbrk{}रो\appbrk{}च\appbrk{}क\appbrk{}हृ\appbrk{}द्रो\appbrk{}गा\appbrk{}ध्मा\appbrk{}न\appbrk{}पा\appbrk{}र्श्वो\appbrk{}द\appbrk{}र\appbrk{}ब\appbrk{}स्ति\appbrk{}शू\appbrk{}ला\appbrk{}ना\appbrk{}ह\appbrk{}मू\appbrk{}त्र\appbrk{}कृ\appbrk{}च्छ्रा\appbrk{}र्शः\appbrk{}प्ली\appbrk{}हो\appbrk{}द\appbrk{}र\appbrk{}पा\appbrk{}ण्डु\appbrk{}रो\appbrk{}गा\appbrk{}न\appbrk{}प\appbrk{}हन्ति}}\Afootnote{H; \skt{कास\appbrk{}श्वा\appbrk{}स\appbrk{}गु\appbrk{}ल्मो\appbrk{}द\appbrk{}रा\appbrk{}रो\appbrk{}च\appbrk{}क\appbrk{}हृ\appbrk{}द्रो\appbrk{}गा\appbrk{}ध्मा\appbrk{}न\appbrk{}पा\appbrk{}र्श्वो\appbrk{}द\appbrk{}र\appbrk{}ब\appbrk{}स्ति\appbrk{}शू\appbrk{}ला\appbrk{}ना\appbrk{}ह\appbrk{}मू\appbrk{}त्र\appbrk{}कृ\appbrk{}च्छ्र\appbrk{}प्ली\appbrk{}हा\appbrk{}र्शस्}\ A}}\skt{ }\edtext{\skt{तूनी\-प्र\-ति\-तू\-नीषु}}{\lemma{\skt{तूनी\appbrk{}प्र\appbrk{}ति\appbrk{}तू\appbrk{}नीषु}}\Afootnote{H; \skt{तूनी\appbrk{}प्र\appbrk{}तू\appbrk{}नी\appbrk{}र\appbrk{}प\appbrk{}हन्ति}\ A}}\skt{ }\edtext{\edtext{\skt{चात्य\-र्थ\-मु\-प\-यु\-ज्यते}}{\lemma{\skt{चात्यर्थम्}}\Afootnote{H; \emph{om.}\ A}}}{\lemma{\skt{उपयुज्यते}}\Afootnote{H; \emph{om.}\ A}}\skt{~॥}\pend

\pstart\versepar \ledsidenote{\passref{4.5.29}}\edtext{\skt{भवन्ति}}{\lemma{\skt{भवन्ति}}\Afootnote{A; \skt{भभवन्ति}\ H}}\skt{ चात्र॥ }\pend
\pstart\versepar \skt{केवलो दोषयुक्तो वा }\edtext{\skt{धातू\-न\-नु\-ग\-तोऽ\-निलः}}{\lemma{\skt{धातूननुगतो}}\Afootnote{H; \skt{धातु\appbrk{}भि\appbrk{}र्वा++\appbrk{}आ\appbrk{}वृतो}\ A}}\skt{~। }\pend
\pstart\versepar \skt{विज्ञेयो }\edtext{\skt{लक्षणोन्याश्च}}{\lemma{\skt{लक्षणोन्याश्च}}\Afootnote{H; \skt{लक्षणोहाभ्यां}\ A}}\skt{ चिकि\-त्स्य\-श्चा\-वि\-रो\-धतः~॥}\pend

\pstart\versepar \ledsidenote{\passref{4.5.30}}\skt{रुजावन्तं घनं शीतं }\edtext{\skt{शोफ\-मे\-दो\-यु\-तोऽ\-निलः}}{\lemma{\skt{शोफमेदोयुतो}}\Afootnote{H; \skt{शोफं मेदोयुतो}\ A}}\skt{~। }\pend
\pstart\versepar \skt{करोति यस्य तम्वैद्यः }\edtext{\skt{शोफ\-व\-त्स\-मु\-पा\-च\-रेत्}}{\lemma{\skt{शोफवत्}}\Afootnote{H; \skt{शोथवत्}\ A}}\skt{~॥}\pend

\pstart\versepar \ledsidenote{\passref{4.5.31}}\edtext{\skt{कफमेदोयुतो}}{\lemma{\skt{कफमेदोयुतो}}\Afootnote{H; \skt{कफमेदोवृतो}\ A}}\skt{ }\edtext{\skt{वायुरूरू}}{\lemma{\skt{ऊरू}}\Afootnote{H; \skt{यदोरू}\ A}}\skt{ }\edtext{\skt{संश्रयते}}{\lemma{\skt{संश्रयते}}\Afootnote{H; \skt{प्रतिपद्यते}\ A}}\skt{ }\edtext{\skt{यदा}}{\lemma{\skt{यदा}}\Afootnote{H; \skt{तदा}\ A}}\skt{~। }\pend
\pstart\versepar \skt{}\edtext{\skt{तदा\-ङ्ग\-म\-र्द्द\-स्तै\-मित्यं}}{\lemma{\skt{तदा\appbrk{}ङ्ग\appbrk{}म\appbrk{}र्द्द\appbrk{}स्तै\appbrk{}मित्यं}}\Afootnote{H; \skt{ऽङ्ग\appbrk{}म\appbrk{}र्द\appbrk{}स्तै\appbrk{}मित्य}\ A}}\skt{ रोम\-ह\-र्ष\-रु\-जा\-ज्वरैः~॥}\pend

\pstart\versepar \ledsidenote{\passref{4.5.32}}\skt{निद्रया }\edtext{\skt{चार्द्दितस्तब्धौ}}{\lemma{\skt{चार्द्दितस्तब्धौ}}\Afootnote{H; \skt{चार्दितौ स्तब्धौ}\ A}}\skt{ }\edtext{\skt{शीत\-ला\-व\-प्र\-वे\-पनौ}}{\lemma{\skt{अप्रवेपनौ}}\Afootnote{H; \skt{अप्रचेतनौ}\ A}}\skt{~। }\pend
\pstart\versepar \skt{}\edtext{\skt{गुरु\-का\-व\-स्थि\-रा\-वूरू}}{\lemma{\skt{अस्थिराव्}}\Afootnote{A; \skt{आस्थिराव्}\ H}}\skt{ }\edtext{\edtext{\skt{पारक्याविव}}{\lemma{\skt{पारक्याव्}}\Afootnote{H; \skt{न स्वा}\ A}}}{\lemma{\skt{इव}}\Afootnote{H; \skt{विव च}\ A}}\skt{ }\edtext{\skt{वर्त्तते}}{\lemma{\skt{वर्त्तते}}\Afootnote{H; \skt{मन्यते}\ A}}\skt{~॥}\pend

\pstart\versepar \ledsidenote{\passref{4.5.33}}\edtext{\skt{तदू\-रू\-स्त\-म्भ\-मि\-त्या\-हु\-रा\-ढ्य\-वा\-त\-म\-था\-परे}}{\lemma{\skt{तदूरूस्तम्भम्}}\Afootnote{H; \skt{तमूरुस्तम्भम्}\ A}}\skt{~। }\pend
\pstart\versepar \skt{}\edtext{\skt{सुखाम्बुना}}{\lemma{\skt{सुखाम्बुना}}\Afootnote{H; \skt{स्नेहव}\ A}}\skt{ }\edtext{\skt{पिबेत्तत्र}}{\lemma{\skt{पिबेत्}}\Afootnote{H; \skt{र्जं पिबेत्}\ A}}\skt{ चूर्ण्णं }\edtext{\skt{षण्ढरणन्नरः}}{\lemma{\skt{षण्ढरणन्}}\Afootnote{H; \skt{षड्धरणं}\ A}}\skt{~॥}\pend

\pstart\versepar \ledsidenote{\passref{4.5.34}}\skt{हितमुष्णाम्बुना }\edtext{\skt{तद्व\-त्पि\-प्प\-ल्या\-दि\-गणे}}{\lemma{\skt{पिप्पल्यादिगणे}}\Afootnote{H; \skt{पिप्पल्यादिगणैः}\ A}}\skt{ कृतं~। }\pend
\pstart\versepar \skt{लिह्याद्वा }\edtext{\skt{त्रैफलाचूर्ण्णं}}{\lemma{\skt{त्रैफलाचूर्ण्णं}}\Afootnote{H; \skt{त्रैफलं चूर्णं}\ A}}\skt{ क्षौद्रेण }\edtext{\skt{कटुकायुतं}}{\lemma{\skt{कटुकायुतं}}\Afootnote{H; \skt{कटुकान्वितम्}\ A}}\skt{~॥}\pend

\pstart\versepar \ledsidenote{\passref{4.5.35}}\skt{मूत्रैर्व्वा }\edtext{\skt{गुग्गुलुं}}{\lemma{\skt{गुग्गुलुं}}\Afootnote{H; \skt{गुग्गुलं}\ A}}\skt{ }\edtext{\skt{श्रेष्ठम्पिबेद्वापि}}{\lemma{\skt{वापि}}\Afootnote{H; \skt{वाऽपि}\ A}}\skt{ शिलाजतुं~। }\pend
\pstart\versepar \skt{}\edtext{\skt{सोऽपहन्ति}}{\lemma{\skt{ऽपहन्ति}}\Afootnote{H; \skt{ततो हन्ति}\ A}}\skt{ कफाक्रान्तं समे\-द\-स्क\-म्प्र\-भ\-ञ्जनं~॥}\edtext{}{\lemma{\skt{शिलाजतुं~। सो}}\Afootnote{H; \skt{शिलाजतु}\ A}}\pend

\pstart\versepar \ledsidenote{\passref{4.5.36}}\skt{हृद्रोगमरुचिं }\edtext{\skt{गुल्म\-न्त\-था\-भ्य\-न्त\-र\-वि\-द्रधिं}}{\lemma{\skt{तथा\appbrk{}भ्य\appbrk{}न्त\appbrk{}र\appbrk{}वि\appbrk{}द्रधिं}}\Afootnote{H; \skt{तथाऽ\appbrk{}भ्य\appbrk{}न्त\appbrk{}र\appbrk{}वि\appbrk{}द्र\appbrk{}धिम्}\ A}}\skt{~। }\pend
\pstart\versepar \skt{}\edtext{\skt{सक्षा\-र\-मू\-त्रा\-न्स्वे\-दाश्च}}{\lemma{\skt{सक्षारमूत्रान्}}\Afootnote{H; \skt{सक्षारमूत्र}\ A}}\skt{ रूक्षा\-ण्यु\-त्सा\-द\-नानि च~॥}\pend

\pstart\versepar \ledsidenote{\passref{4.5.37}}\edtext{\skt{कुर्याद्दिह्याच्च}}{\lemma{\skt{दिह्याच्च}}\Afootnote{A; \skt{दिहेच्च}\ H}}\skt{ मूत्राढ्यैः कर\-ञ्ज\-फ\-ल\-स\-र्षपैः~। }\pend
\pstart\versepar \skt{}\edtext{\skt{भोज्या}}{\lemma{\skt{भोज्या}}\Afootnote{H; \skt{भोज्याः}\ A}}\skt{ }\edtext{\skt{पुराणाः}}{\lemma{\skt{पुराणाः}}\Afootnote{H; \skt{पुराण}\ A}}\skt{ }\edtext{\skt{श्यामाकाः}}{\lemma{\skt{श्यामाकाः}}\Afootnote{H; \skt{श्यामाक}\ A}}\skt{ }\edtext{\skt{कोद्र\-वो\-द्दा\-ल\-का\-दयः}}{\lemma{\skt{कोद्र\appbrk{}वो\appbrk{}द्दा\appbrk{}ल\appbrk{}का\appbrk{}दयः}}\Afootnote{H; \skt{कोद्र\appbrk{}वो\appbrk{}हा\appbrk{}ल\appbrk{}शा\appbrk{}लयः}\ A}}\skt{~॥}\pend

\pstart \ledsidenote{\passref{4.5.38cd}}\edtext{\skt{जाङ्ग\-लै\-र\-घृ\-तै\-र्म्मांसैः}}{\lemma{\skt{अघृतैर्}}\Afootnote{A; \skt{अकृतैर्}\ H}}\skt{ शाकै\-श्चा\-ल\-व\-णै\-र्हितैः~।}\pend

\pstart\versepar \ledsidenote{\passref{4.5.39}}\edtext{\skt{परिक्षीणौ}}{\lemma{\skt{परिक्षीणौ}}\Afootnote{H; \emph{om.}\ A}}\skt{ यदा स्यातां }\edtext{\edtext{\skt{भूयि\-ष्ठ\-ङ्क\-फ\-मे\-दसी}}{\lemma{\skt{भूयिष्ठङ्}}\Afootnote{H; \skt{परिक्षीणे भूयिष्ठे}\ A}}}{\lemma{\skt{कफमेदसी}}\Afootnote{A H\textsuperscript{pc}; \skt{कफमेसी}\ H\textsuperscript{ac}}}\skt{~॥ }\pend
\pstart\versepar \skt{तदा स्नेहादिकङ्कर्म्म }\edtext{\edtext{\skt{पुन\-र\-त्रा\-व\-चा\-र\-ये\-दिति}}{\lemma{\skt{पुनर्}}\Afootnote{H; \skt{पुन}\ A}}}{\lemma{\skt{अत्रा\appbrk{}व\appbrk{}चा\appbrk{}र\appbrk{}ये\appbrk{}दिति}}\Afootnote{H; \skt{रत्रा\appbrk{}व\appbrk{}चा\appbrk{}र\appbrk{}येत्}\ A}}\skt{~॥}\pend

\pstart \ledsidenote{\passref{4.5.trailer}}\edtext{\skt{चि}}{\lemma{\skt{चि}}\Afootnote{H; \skt{इति श्रीसु\appbrk{}श्रु\appbrk{}त\appbrk{}सं\appbrk{}हि\appbrk{}तायां चिकित्सास्थाने महा\appbrk{}वा\appbrk{}त\appbrk{}व्या\appbrk{}धि\appbrk{}चि\appbrk{}कि\appbrk{}त्सितं नाम पञ्चमोऽध्यायः}\ A}}\skt{ ५}\pend

\endnumbering
\beginnumbering
\eledsection*{Adhyāya 6}
\pstart \ledsidenote{\passref{4.6.1}}\skt{अथातोऽर्शसां }\edtext{\skt{चिकि\-त्सि\-त\-म्व्या\-ख्याम्याः}}{\lemma{\skt{चिकि\appbrk{}त्सि\appbrk{}त\appbrk{}म्व्या\appbrk{}ख्याम्याः}}\Afootnote{H; \skt{चिकित्सितं व्याख्यास्यामः}\ A}}\skt{~॥}\pend

\pstart \ledsidenote{\passref{4.6.2}}\edtext{\skt{अथ}}{\lemma{\skt{अथ}}\Afootnote{H; \emph{om.}\ A}}\skt{ }\edtext{\skt{खलु}}{\lemma{\skt{खलु}}\Afootnote{H; \emph{om.}\ A}}\skt{ }\edtext{\skt{वत्ससुश्रुतः}}{\lemma{\skt{वत्ससुश्रुतः}}\Afootnote{H; \skt{यथोवाच भग\appbrk{}वा\appbrk{}न्ध\appbrk{}न्व\appbrk{}न्तरिः}\ A}}\skt{~।}\pend

\pstart \ledsidenote{\passref{4.6.3}}\skt{चतुर्व्विधोऽर्शसां साधनोपायः~। तद्यथाभेषजं क्षारोऽग्निः }\edtext{\skt{शस्त्र इ~॥~॥ ति}}{\lemma{\skt{शस्त्र इ~॥~॥ ति}}\Afootnote{H; \skt{शस्त्रमिति}\ A}}\skt{~। }\edtext{\edtext{\skt{तत्रा\-चि\-र\-का\-ल\-जा\-ता\-न्य\-ल्प\-दो\-ष\-लि\-ङ्गो\-प\-द्र\-वानि}}{\lemma{\skt{तत्रा\appbrk{}चि\appbrk{}र\appbrk{}का\appbrk{}ल\appbrk{}जा\appbrk{}तान्य्}}\Afootnote{H; \skt{तत्र अचि\appbrk{}र\appbrk{}का\appbrk{}ल\appbrk{}जा\appbrk{}तान्}\ A}}}{\lemma{\skt{अल्प\appbrk{}दो\appbrk{}ष\appbrk{}लि\appbrk{}ङ्गो\appbrk{}प\appbrk{}द्र\appbrk{}वानि}}\Afootnote{H; \skt{यल्प\appbrk{}दो\appbrk{}ष\appbrk{}लि\appbrk{}ङ्गो\appbrk{}प\appbrk{}द्र\appbrk{}वाणि}\ A}}\skt{ भेषजसाध्यानि~। मृदु\-प्र\-सृ\-ता\-व\-गा\-ढा\-न्यु\-छ्रि\-तानि क्षारेण~। }\edtext{\edtext{\skt{कर्क्क\-श\-स्थि\-र\-पृ\-थु\-क\-ठि\-ना\-न्य\-ग्निना}}{\lemma{\skt{कर्क्क\appbrk{}श\appbrk{}स्थि\appbrk{}र\appbrk{}पृ\appbrk{}थु\appbrk{}क\appbrk{}ठि\appbrk{}नान्य्}}\Afootnote{H; \skt{कर्क\appbrk{}श\appbrk{}स्थि\appbrk{}र\appbrk{}पृ\appbrk{}थु\appbrk{}क\appbrk{}ठि\appbrk{}नान्}\ A}}}{\lemma{\skt{अग्निना}}\Afootnote{H; \skt{यग्निना}\ A}}\skt{~। }\edtext{\edtext{\skt{अत\-नु\-मू\-ला\-न्यु\-छ्रि\-तानि}}{\lemma{\skt{अतनुमूलान्य्}}\Afootnote{H; \skt{तनुमूलन्}\ A}}}{\lemma{\skt{उछ्रितानि}}\Afootnote{H; \skt{युच्छ्रितानि}\ A}}\skt{ क्लेदवन्ति च शस्त्रेण~। तत्र }\edtext{\edtext{\skt{भेष\-ज\-सा\-ध्या\-ना\-म\-र्श\-सा\-म\-दृ\-श्या\-नाञ्च}}{\lemma{\skt{भेष\appbrk{}ज\appbrk{}सा\appbrk{}ध्या\appbrk{}नाम्}}\Afootnote{H; \skt{भेषजसाध्याना}\ A}}}{\lemma{\skt{अर्श\appbrk{}सा\appbrk{}म\appbrk{}दृ\appbrk{}श्या\appbrk{}नाञ्च}}\Afootnote{H; \skt{मर्श\appbrk{}सा\appbrk{}म\appbrk{}दृ\appbrk{}श्यानां च}\ A}}\skt{ भेषजम्भवति~। क्षाराग्नि शस्त्रसाध्यानां }\edtext{\skt{क्षाराग्निशस्त्राणि च}}{\lemma{\skt{क्षाराग्निशस्त्राणि च}}\Afootnote{H; \emph{om.}\ A}}\skt{~। }\edtext{\edtext{\edtext{\skt{तेषा\-म्वि\-धा\-न\-मु\-प\-धा\-र\-यस्व}}{\lemma{\skt{तेषाम्}}\Afootnote{H; \skt{तु}\ A}}}{\lemma{\skt{विधानम्}}\Afootnote{H; \skt{विधा\appbrk{}न\appbrk{}मु\appbrk{}च्य\appbrk{}मा\appbrk{}नम्}\ A}}}{\lemma{\skt{उपधारयस्व}}\Afootnote{H; \skt{उपधारय}\ A}}\skt{~।}\pend

\pstart \ledsidenote{\passref{4.6.4}}\edtext{\edtext{\edtext{\edtext{\skt{तत्र\-व\-ल\-व\-न्त\-म\-र्शो\-भि\-रु\-प\-द्रु\-त\-मु\-प\-स्नि\-ग्ध\-प\-रि\-स्नि\-न्न\-म\-नि\-ल\-वे\-द\-ना\-भि\-व्र्द्धि\-श\-म\-नार्थं}}{\lemma{\skt{तत्र\appbrk{}व\appbrk{}ल\appbrk{}व\appbrk{}न्तम्}}\Afootnote{H; \skt{तत्र बलवन्त}\ A}}}{\lemma{\skt{अर्शोभिर्}}\Afootnote{H; \skt{मातु\appbrk{}र\appbrk{}म\appbrk{}र्शो\appbrk{}भिर्}\ A}}}{\lemma{\skt{उप\appbrk{}स्नि\appbrk{}ग्ध\appbrk{}प\appbrk{}रि\appbrk{}स्नि\appbrk{}न्नम्}}\Afootnote{H; \skt{उपस्निग्धं परिस्विन्नम्}\ A}}}{\lemma{\skt{अनि\appbrk{}ल\appbrk{}वे\appbrk{}द\appbrk{}ना\appbrk{}भि\appbrk{}व्र्द्धि\appbrk{}श\appbrk{}म\appbrk{}नार्थं}}\Afootnote{H; \skt{अनि\appbrk{}ल\appbrk{}वे\appbrk{}द\appbrk{}ना\appbrk{}भि\appbrk{}वृ\appbrk{}द्धि\appbrk{}प्र\appbrk{}श\appbrk{}मार्थं}\ A}}\skt{ }\edtext{\edtext{\skt{स्निग्ध\-म\-न\-ल्प\-म\-न्न\-द्र\-व\-भु\-क्त\-वन्तं}}{\lemma{\skt{अनल्पम्}}\Afootnote{H; \skt{उष्णमल्पम्}\ A}}}{\lemma{\skt{अन्न\appbrk{}द्र\appbrk{}व\appbrk{}भु\appbrk{}क्त\appbrk{}वन्तं}}\Afootnote{H; \skt{अन्नं द्रवप्रायं भुक्तवन्तम्}\ A}}\skt{ }\edtext{\skt{सम्वेश्य}}{\lemma{\skt{सम्वेश्य}}\Afootnote{H; \skt{उपवेश्य}\ A}}\skt{ }\edtext{\skt{शुचौ}}{\lemma{\skt{शुचौ}}\Afootnote{H; \skt{संवृ(भृ)ते शुचौ}\ A}}\skt{ }\edtext{\skt{साधारणे}}{\lemma{\skt{साधारणे}}\Afootnote{H; \skt{देशे साधारणे}\ A}}\skt{ }\edtext{\skt{देशे}}{\lemma{\skt{देशे}}\Afootnote{H; \emph{om.}\ A}}\skt{ व्यभ्रे काले समे फलके }\edtext{\skt{प्रत्या\-दि\-त्य\-गु\-द\-म\-न्य\-स्यो\-त्सङ्गे}}{\lemma{\skt{प्रत्यादित्य्}}\Afootnote{H; \skt{शय्यायां वा प्रत्यादित्य्}\ A}}\skt{ }\edtext{\skt{निषर्ण्णं}}{\lemma{\skt{निषर्ण्णं}}\Afootnote{H; \skt{निषण्ण}\ A}}\skt{ }\edtext{\skt{पूर्व्व\-का\-य\-मु\-त्ता\-न\-मु\-न्न\-त\-क\-टीकं}}{\lemma{\skt{उन्नतकटीकं}}\Afootnote{H; \skt{किञ्चि\appbrk{}दु\appbrk{}न्न\appbrk{}त\appbrk{}क\appbrk{}टिकं}\ A}}\skt{ }\edtext{\skt{यस्त्रसाटकेन}}{\lemma{\skt{यस्त्रसाटकेन}}\Afootnote{H; \skt{वस्त्र\appbrk{}क\appbrk{}म्ब\appbrk{}ल\appbrk{}को\appbrk{}प\appbrk{}विष्टं यन्त्र\appbrk{}ण\appbrk{}शा\appbrk{}ट\appbrk{}केन}\ A}}\skt{ }\edtext{\edtext{\skt{परि\-क्षि\-प्य\-ग्री\-वा\-स\-क्थि\-रु\-परि}}{\lemma{\skt{परि\appbrk{}क्षि\appbrk{}प्य\appbrk{}ग्री\appbrk{}वा\appbrk{}स\appbrk{}क्थिर्}}\Afootnote{H; \skt{परि\appbrk{}क्षि\appbrk{}प्त\appbrk{}ग्री\appbrk{}वा\appbrk{}सक्थिं}\ A}}}{\lemma{\skt{उपरि}}\Afootnote{H; \skt{परि}\ A}}\skt{ }\edtext{\skt{कर्म्मभिः}}{\lemma{\skt{कर्म्मभिः}}\Afootnote{H; \skt{कर्मिभिः}\ A}}\skt{ सुपरिगृहीतं }\edtext{\skt{कारयित्वा}}{\lemma{\skt{कारयित्वा}}\Afootnote{H; \skt{अस्प\appbrk{}न्द\appbrk{}न\appbrk{}श\appbrk{}रीरं कृत्वा}\ A}}\skt{ ततोऽस्मै }\edtext{\skt{घृताभ्यक्तं}}{\lemma{\skt{घृताभ्यक्तं}}\Afootnote{H; \skt{घृता\appbrk{}भ्य\appbrk{}क्त\appbrk{}गु\appbrk{}दाय घृताभ्यक्तं}\ A}}\skt{ यन्त्र\-मृ\-ज्व\-नु\-मुखं }\edtext{\skt{शन्नैः}}{\lemma{\skt{शन्नैः}}\Afootnote{H; \skt{पायौ शनैः शनैः प्रवाहमाणस्य}\ A}}\skt{ प्रणिधाय प्रविष्टे }\edtext{\skt{चार्शोऽ\-भि\-वी\-क्ष्य\-स\-ला\-क\-यो\-त्पीड्य}}{\lemma{\skt{ऽभि\appbrk{}वी\appbrk{}क्ष्य\appbrk{}स\appbrk{}ला\appbrk{}क\appbrk{}यो\appbrk{}त्पीड्य}}\Afootnote{H; \skt{वीक्ष्य, शलाकयोत्पीड्य}\ A}}\skt{ }\edtext{\skt{पिचु\-प्लो\-त\-यो\-र\-न्य\-त\-रेण}}{\lemma{\skt{पिचुप्लोतयोर्}}\Afootnote{H; \skt{पिचुवस्त्रयोर्}\ A}}\skt{ }\edtext{\skt{प्रमृ\-ज्या\-क्षा\-र\-म्पा\-त\-येत्}}{\lemma{\skt{प्रमृज्याक्षारम्}}\Afootnote{H; \skt{प्रमृज्य, क्षारं}\ A}}\skt{~। }\edtext{\skt{पातयित्वा}}{\lemma{\skt{पातयित्वा}}\Afootnote{H; \skt{पातयित्वा च}\ A}}\skt{ पाणिना यन्त्र\-द्वा\-र\-म्पि\-धाय }\edtext{\edtext{\skt{वाक्छ\-त\-मा\-त्र\-मु\-पे\-क्षेत}}{\lemma{\skt{वाक्छतमात्रम्}}\Afootnote{H; \skt{वाक्च्छतमात्र}\ A}}}{\lemma{\skt{उपेक्षेत}}\Afootnote{H; \skt{मुपेक्षेत}\ A}}\skt{~। ततः }\edtext{\skt{प्रमृज्य}}{\lemma{\skt{प्रमृज्य}}\Afootnote{H; \skt{पर्मृज्य}\ A}}\skt{ }\edtext{\skt{क्षार\-ल\-व\-ण\-ञ्चा\-वेक्ष्य}}{\lemma{\skt{क्षारलवणञ्}}\Afootnote{H; \skt{क्षारबलं व्याधिबलं}\ A}}\skt{ }\edtext{\skt{पुनः}}{\lemma{\skt{पुनः}}\Afootnote{H; \skt{पु}\ A}}\skt{ }\edtext{\skt{प्रतिसारयेत्~। अथर्शः}}{\lemma{\skt{प्रतिसारयेत्~। अथर्शः}}\Afootnote{H; \skt{नरालेपयेत्, अथार्शः}\ A}}\skt{ }\edtext{\edtext{\skt{पक्व\-जा\-म्वू\-प्र\-का\-श\-म\-भि\-स\-मीक्ष्यो}}{\lemma{\skt{पक्व\appbrk{}जा\appbrk{}म्वू\appbrk{}प्र\appbrk{}का\appbrk{}शम्}}\Afootnote{H; \skt{पक्व\appbrk{}जा\appbrk{}म्ब\appbrk{}व\appbrk{}प्र\appbrk{}तीका}\ A}}}{\lemma{\skt{अभिसमीक्ष्यो}}\Afootnote{H; \skt{शम}\ A}}\skt{ }\edtext{\skt{वस\-न्न\-मी\-ष\-त्ततः}}{\lemma{\skt{वस\appbrk{}न्न\appbrk{}मी\appbrk{}ष\appbrk{}त्ततः}}\Afootnote{H; \skt{वस\appbrk{}न्न\appbrk{}मी\appbrk{}ष\appbrk{}न्न\appbrk{}त\appbrk{}म\appbrk{}भिस्}\ A}}\skt{~। }\edtext{\skt{अपावर्त्तयेत्}}{\lemma{\skt{अपावर्त्तयेत्}}\Afootnote{H; \skt{अमी\appbrk{}क्ष्यो\appbrk{}पा\appbrk{}व\appbrk{}र्त\appbrk{}येत्, क्षारं प्रक्षालयेद्}\ A}}\skt{~। }\edtext{\edtext{\skt{धान्या\-म्ल\-द\-धि\-म\-स्तु\-शु\-क्ता\-ना\-म\-न्य\-त\-मेन}}{\lemma{\skt{धान्या\appbrk{}म्ल\appbrk{}द\appbrk{}धि\appbrk{}म\appbrk{}स्तु\appbrk{}शु\appbrk{}क्ता\appbrk{}नाम्}}\Afootnote{H; \skt{धान्याम्लेन दधिमस्तुशुक्त}\ A}}}{\lemma{\skt{अन्यतमेन}}\Afootnote{H; \skt{फलाम्लैर्वा}\ A}}\skt{~। ततो }\edtext{\edtext{\skt{यष्टी\-म\-धु\-क\-मि\-श्रे\-ण\-स\-र्प्पिषा}}{\lemma{\skt{यष्टीम्}}\Afootnote{H; \skt{यष्टी}\ A}}}{\lemma{\skt{अधु\appbrk{}क\appbrk{}मि\appbrk{}श्रे\appbrk{}ण\appbrk{}स\appbrk{}र्प्पिषा}}\Afootnote{H; \skt{मधुकमिश्रेण सर्पिषा}\ A}}\skt{ निर्व्वाप्य }\edtext{\edtext{\edtext{\edtext{\skt{यन्त्र\-मु\-प\-नी\-यो\-त्था\-प्या\-तु\-र\-सु\-ष्णो\-द\-केऽ\-व\-गा\-ह्य\-शी\-ता\-भि\-रद्भिः}}{\lemma{\skt{यन्त्रम्}}\Afootnote{H; \skt{यन्त्र}\ A}}}{\lemma{\skt{उपनीयोत्थाप्य्}}\Afootnote{H; \skt{मपनीयोत्थाप्}\ A}}}{\lemma{\skt{आतु\appbrk{}र\appbrk{}सु\appbrk{}ष्णो\appbrk{}दके}}\Afootnote{H; \skt{यातु\appbrk{}र\appbrk{}मु\appbrk{}ष्णो\appbrk{}दको}\ A}}}{\lemma{\skt{ऽव\appbrk{}गा\appbrk{}ह्य\appbrk{}शी\appbrk{}ता\appbrk{}भिर्}}\Afootnote{H; \skt{पविष्टं शीताभिर्}\ A}}\skt{ }\edtext{\skt{परिषेचयेत्~। }\edtext{\skt{अशी\-ता\-भि\-र\-द्भि\-रि\-त्येके}}{\lemma{\skt{अद्भिर्}}\Afootnote{H; \skt{भिर्}\ A}}}{\lemma{\skt{परिषेचयेत्~। अशीताभिर्}}\Afootnote{H; \skt{परिषिञ्चेत्, अशीता}\ A}}\skt{~। ततो }\edtext{\edtext{\edtext{\edtext{\skt{निवा\-त\-म\-गा\-र\-म्प्रे\-वे\-श्या\-चा\-रि\-क\-मु\-प\-दि\-शेत्}}{\lemma{\skt{निवातम्}}\Afootnote{H; \skt{निर्वात}\ A}}}{\lemma{\skt{अगारम्}}\Afootnote{H; \skt{मागारं}\ A}}}{\lemma{\skt{प्रेवे\appbrk{}श्या\appbrk{}चा\appbrk{}रि\appbrk{}कम्}}\Afootnote{H; \skt{प्रवे\appbrk{}श्या\appbrk{}चा\appbrk{}रि\appbrk{}कम्}\ A}}}{\lemma{\skt{उपदिशेत्}}\Afootnote{H; \skt{आदिशेत्}\ A}}\skt{~। }\edtext{\edtext{\skt{साव\-शे\-ष\-म्पु\-न\-र्द्दा\-ह\-येत्}}{\lemma{\skt{पुनर्}}\Afootnote{H; \skt{पुन}\ A}}}{\lemma{\skt{द्दाहयेत्}}\Afootnote{H; \skt{र्दहे}\ A}}\skt{~। }\edtext{\skt{तत्र}}{\lemma{\skt{तत्र}}\Afootnote{H; \skt{त्; एवं}\ A}}\skt{ }\edtext{\edtext{\skt{सप्त\-रा\-त्रा\-च्चै\-कै\-क\-मु\-प\-क्रा\-मेत्}}{\lemma{\skt{सप्त\appbrk{}रा\appbrk{}त्रा\appbrk{}च्चै\appbrk{}कै\appbrk{}कम्}}\Afootnote{\skt{सप्त\appbrk{}रा\appbrk{}त्रा\appbrk{}त्स\appbrk{}प्त\appbrk{}रा\appbrk{}त्रा\appbrk{}दे\appbrk{}कै\appbrk{}कम्}\ A; \skt{सप्त\appbrk{}रा\appbrk{}त्रा\appbrk{}च्चै\appbrk{}कै\appbrk{}कम्}\ H}}}{\lemma{\skt{उपक्रामेत्}}\Afootnote{H; \skt{उपक्रमेत; तत्र बहुषु पूर्व्ंअ दक्षिणं साधयेत्}\ A}}\skt{~। }\edtext{\skt{ततो}}{\lemma{\skt{ततो}}\Afootnote{H; \skt{दक्षिणाद्वामं}\ A}}\skt{ }\edtext{\skt{वामात्ततः}}{\lemma{\skt{ततः}}\Afootnote{H; \emph{om.}\ A}}\skt{ }\edtext{\skt{पृष्ठ\-जा\-ग्र\-ज\-मिति}}{\lemma{\skt{पृष्ठजाग्रजम्}}\Afootnote{H; \skt{पृष्ठजं, ततोऽग्रजम्}\ A}}\skt{~।}\pend

\pstart \ledsidenote{\passref{4.6.5}}\skt{तत्र }\edtext{\skt{वात\-श्ले\-ष्म\-नि\-मि\-त्ता\-न्य\-ग्नि\-क्षा\-राभ्यां}}{\lemma{\skt{वात\appbrk{}श्ले\appbrk{}ष्म\appbrk{}नि\appbrk{}मि\appbrk{}त्ता\appbrk{}न्य\appbrk{}ग्नि\appbrk{}क्षा\appbrk{}राभ्यां}}\Afootnote{\skt{वात\appbrk{}श्ले\appbrk{}ष्म\appbrk{}नि\appbrk{}मि\appbrk{}त्ता\appbrk{}न्य\appbrk{}ग्नि\appbrk{}क्षा\appbrk{}राभ्यां}\ A H}}\skt{ साधयेत्~। क्षारेणैव }\edtext{\skt{पित्त\-र\-क्त\-स\-मु\-त्थानि}}{\lemma{\skt{पित्त\appbrk{}र\appbrk{}क्त\appbrk{}स\appbrk{}मु\appbrk{}त्थानि}}\Afootnote{H; \skt{मृदुना पित्त\appbrk{}र\appbrk{}क्त\appbrk{}नि\appbrk{}मि\appbrk{}त्तानि}\ A}}\skt{~।}\pend

\pstart \ledsidenote{\passref{4.6.6}}\skt{तत्र }\edtext{\edtext{\edtext{\edtext{\edtext{\skt{वाता\-नु\-लो\-म्य\-म\-न्ने\-भि\-रु\-चि\-र\-ग्नि\-दी\-प्ति\-ला\-घ\-व\-म्व\-ल\-व\-र्ण्णो\-त्प\-त्ति\-र्म्म\-न\-स्तु\-ष्टि\-रिति}}{\lemma{\skt{वाता\appbrk{}नु\appbrk{}लो\appbrk{}म्य\appbrk{}म\appbrk{}न्ने\appbrk{}भिर्}}\Afootnote{H; \skt{वाता\appbrk{}नु\appbrk{}लो\appbrk{}म्य\appbrk{}मन्न}\ A}}}{\lemma{\skt{उचि\appbrk{}र\appbrk{}ग्नि\appbrk{}दी\appbrk{}प्ति\appbrk{}ला\appbrk{}घ\appbrk{}वम्}}\Afootnote{H; \skt{रुचि\appbrk{}र\appbrk{}ग्नि\appbrk{}दी\appbrk{}प्ति\appbrk{}र्ला\appbrk{}घवं}\ A}}}{\lemma{\skt{वल\appbrk{}व\appbrk{}र्ण्णो\appbrk{}त्प\appbrk{}त्तिर्}}\Afootnote{H; \skt{बलवर्णोत्पत्ति}\ A}}}{\lemma{\skt{म्मनस्तुष्टिर्}}\Afootnote{\skt{र्मनस्तुष्टि}\ A; \skt{म्मनस्तुष्टिर्}\ H}}}{\lemma{\skt{इति}}\Afootnote{H; \skt{रिति}\ A}}\skt{ सम्य\-ग्द\-ग्ध\-लि\-ङ्गानि~। अतिदग्धे तु }\edtext{\edtext{\skt{गुदा\-व\-र\-ण\-न्दा\-हो\-द्ज्वरः}}{\lemma{\skt{गुदावरणन्}}\Afootnote{H; \skt{गुदावदरणं}\ A}}}{\lemma{\skt{दाहोद्ज्वरः}}\Afootnote{H; \skt{दाहो मूर्च्छा ज्वरः}\ A}}\skt{ पिपासा }\edtext{\edtext{\skt{शोणि\-ता\-ति\-प्र\-वृ\-त्ति\-स्त\-न्नि\-मि\-त्ता\-श्चो\-प\-द्रवा}}{\lemma{\skt{शोणि\appbrk{}ता\appbrk{}ति\appbrk{}प्र\appbrk{}वृ\appbrk{}त्तिस्}}\Afootnote{H; \skt{शोणि\appbrk{}ता\appbrk{}ति\appbrk{}प्र\appbrk{}वृत्ति}\ A}}}{\lemma{\skt{तन्}}\Afootnote{H; \skt{स्तन्}\ A}}\skt{ भवन्ति~। }\edtext{\edtext{\skt{ध्यामा\-ल्प\-ता\-क\-ण्डू\-न\-नि\-ल\-वै\-गु\-ण्यञ्च}}{\lemma{\skt{ध्यामा\appbrk{}ल्प\appbrk{}ता\appbrk{}क\appbrk{}ण्डू\appbrk{}न\appbrk{}नि\appbrk{}ल\appbrk{}वै\appbrk{}गु\appbrk{}ण्यञ्}}\Afootnote{H; \skt{ध्यामाल्पव्रणता कण्डू\appbrk{}र\appbrk{}नि\appbrk{}ल\appbrk{}वै\appbrk{}गु\appbrk{}ण्य\appbrk{}मि\appbrk{}न्दि\appbrk{}र्या\appbrk{}णा\appbrk{}म\appbrk{}प्र\appbrk{}सादो विकारस्य}\ A}}}{\lemma{\skt{च}}\Afootnote{H; \skt{चाशान्ति}\ A}}\skt{ }\edtext{\skt{हीनदग्धे}}{\lemma{\skt{हीनदग्धे}}\Afootnote{H; \skt{र्हीनदग्धे}\ A}}\skt{~।}\pend

\pstart \ledsidenote{\passref{4.6.7}}\edtext{\skt{म हान्ति वा}}{\lemma{\skt{म हान्ति वा}}\Afootnote{H; \skt{महान्ति च}\ A}}\skt{ }\edtext{\skt{प्राण\-व\-त\-स्थि\-त्वा\-द\-हेत्}}{\lemma{\skt{प्राण\appbrk{}व\appbrk{}त\appbrk{}स्थि\appbrk{}त्वा\appbrk{}द\appbrk{}हेत्}}\Afootnote{H; \skt{प्राणवतच्छित्त्वा दहेत्}\ A}}\skt{~। निर्गतानि }\edtext{\skt{चात्य\-र्थ\-न्दो\-ष\-प\-रि\-पू\-र्ण्णा\-ग\-तानि}}{\lemma{\skt{दोष\appbrk{}प\appbrk{}रि\appbrk{}पू\appbrk{}र्ण्णा\appbrk{}ग\appbrk{}तानि}}\Afootnote{H; \skt{दोषपूर्णानि यन्त्राद्विना}\ A}}\skt{ }\edtext{\skt{स्नेहा\-भ्य\-ङ्ग\-स्वे\-दो\-प\-नाह}}{\lemma{\skt{स्नेहा\appbrk{}भ्य\appbrk{}ङ्ग\appbrk{}स्वे\appbrk{}दो\appbrk{}प\appbrk{}नाह}}\Afootnote{H; \skt{स्वेदा\appbrk{}भ्य\appbrk{}ङ्ग\appbrk{}स्ने\appbrk{}हा\appbrk{}व\appbrk{}गा\appbrk{}हो\appbrk{}प\appbrk{}नाह}\ A}}\skt{ }\edtext{\edtext{\skt{विस्रा\-व\-णा\-ले\-प\-नै\-रु\-प\-च\-रेत्}}{\lemma{\skt{विस्रा\appbrk{}व\appbrk{}णा\appbrk{}ले\appbrk{}प\appbrk{}नैर्}}\Afootnote{H; \skt{विस्रा\appbrk{}व\appbrk{}णा\appbrk{}ले\appbrk{}प\appbrk{}क्षा\appbrk{}रा\appbrk{}ग्नि\appbrk{}श\appbrk{}स्त्रैर्}\ A}}}{\lemma{\skt{उपचरेत्}}\Afootnote{H; \skt{उपाचरेत्}\ A}}\skt{~। }\edtext{\skt{प्रवृत्त्रक्तानि वा}}{\lemma{\skt{प्रवृत्त्रक्तानि वा}}\Afootnote{H; \skt{प्रवृत्तरक्तानि}\ A}}\skt{ }\edtext{\skt{पित्त}}{\lemma{\skt{पित्त}}\Afootnote{H; \skt{च}\ A}}\skt{ }\edtext{\skt{रक्तविधानेन}}{\lemma{\skt{रक्तविधानेन}}\Afootnote{H; \skt{रक्त\appbrk{}पि\appbrk{}त्त\appbrk{}वि\appbrk{}धा\appbrk{}नेन}\ A}}\skt{~। }\edtext{\skt{भिन्नवर्च्चांसि}}{\lemma{\skt{भिन्नवर्च्चांसि}}\Afootnote{H; \skt{भिन्नपुरीषाणि}\ A}}\skt{ चाती\-सा\-र\-वि\-धा\-नेन~। वद्धवर्च्चांसि }\edtext{\skt{चोदा\-व\-र्त्त\-वि\-धा\-ने\-नेति}}{\lemma{\skt{चोदा\appbrk{}व\appbrk{}र्त्त\appbrk{}वि\appbrk{}धा\appbrk{}ने\appbrk{}नेति}}\Afootnote{H; \skt{स्नेह\appbrk{}पा\appbrk{}न\appbrk{}वि\appbrk{}धा\appbrk{}ने\appbrk{}नो\appbrk{}दा\appbrk{}व\appbrk{}र्त\appbrk{}वि\appbrk{}धा\appbrk{}नेन वा}\ A}}\skt{~। एष }\edtext{\edtext{\skt{सर्व्व\-स्था\-न\-ग\-ता\-ना\-म\-र्श\-सा\-न्द\-ह\-न\-कल्पः}}{\lemma{\skt{सर्व्व\appbrk{}स्था\appbrk{}न\appbrk{}ग\appbrk{}ता\appbrk{}नाम्}}\Afootnote{H; \skt{सर्व\appbrk{}स्था\appbrk{}न\appbrk{}ग\appbrk{}ताना}\ A}}}{\lemma{\skt{अर्शसान्}}\Afootnote{H; \skt{मर्शसां}\ A}}\skt{~।}\pend

\pstart \ledsidenote{\passref{4.6.8}}\edtext{\skt{आसाद्य}}{\lemma{\skt{आसाद्य}}\Afootnote{H; \skt{आसाद्य च}\ A}}\skt{ दर्व्वीकू }\edtext{\skt{र्च्चस\-ला\-का\-ना\-म\-न्य\-त\-मेन}}{\lemma{\skt{र्च्चसलाकानाम्}}\Afootnote{H; \skt{र्चक\appbrk{}श\appbrk{}ला\appbrk{}का\appbrk{}नाम्}\ A}}\skt{ क्षारम्पातयेत्~। भ्रष्टगुदस्य तुविना }\edtext{\skt{यन्त्रं}}{\lemma{\skt{यन्त्रं}}\Afootnote{H; \skt{यन्त्रेण}\ A}}\skt{ क्षारादिकर्म्म प्रयुञ्जीत~। सर्व्वेषु च शालि\-ष\-ष्टि\-क\-य\-व\-गो\-धू\-मान्नं सर्प्पिः स्निग्ध\-मु\-प\-से\-वेत~। पयसा निम्वयूषेण पटोलयूषेण वा~। }\edtext{\edtext{\skt{यथा\-दो\-षै\-र्व्वा\-स्तू\-क\-त\-ण्डु\-ली\-य\-क\-जी\-व\-न्त्यु\-पो\-द\-काश्च}}{\lemma{\skt{यथादोषैर्}}\Afootnote{H; \skt{यथादोषं शाकैर्}\ A}}}{\lemma{\skt{व्वास्तू\appbrk{}क\appbrk{}त\appbrk{}ण्डु\appbrk{}ली\appbrk{}य\appbrk{}क\appbrk{}जी\appbrk{}व\appbrk{}न्त्यु\appbrk{}पो\appbrk{}द\appbrk{}काश्च}}\Afootnote{H; \skt{वास्तू\appbrk{}क\appbrk{}त\appbrk{}ण्डु\appbrk{}ली\appbrk{}य\appbrk{}क\appbrk{}जी\appbrk{}व\appbrk{}न्त्यु\appbrk{}पो\appbrk{}दि\appbrk{}काश्व}\ A}}\skt{ }\edtext{\edtext{\skt{वल\-वा\-ल\-मू\-ल\-क\-श\-ण\-चु\-ञ्चु\-चि\-ल्ली\-ना\-म\-न्य\-त\-मेन}}{\lemma{\skt{वल\appbrk{}वा\appbrk{}ल\appbrk{}मू\appbrk{}ल\appbrk{}क\appbrk{}श\appbrk{}ण\appbrk{}चु\appbrk{}ञ्चु\appbrk{}चि\appbrk{}ल्ली\appbrk{}नाम्}}\Afootnote{H; \skt{बला\appbrk{}ब\appbrk{}ल\appbrk{}मू\appbrk{}ल\appbrk{}क\appbrk{}पा\appbrk{}ल\appbrk{}ङ्क्य\appbrk{}स\appbrk{}न\appbrk{}चि\appbrk{}ल्ली\appbrk{}चु\appbrk{}च्चू\appbrk{}क\appbrk{}ला\appbrk{}य\appbrk{}व\appbrk{}ल्ली\appbrk{}भिर्}\ A}}}{\lemma{\skt{अन्यतमेन}}\Afootnote{H; \skt{अन्यैर्वा}\ A}}\skt{~। }\edtext{\skt{यच्चान्यदपि}}{\lemma{\skt{यच्चान्यदपि}}\Afootnote{H; \skt{यच्चान्यदपि}\ A}}\skt{ }\edtext{\skt{स्निघ्द\-म\-ग्नि\-दी\-प\-न\-म\-र्शोघ्नं}}{\lemma{\skt{स्निघ्द\appbrk{}म\appbrk{}ग्नि\appbrk{}दी\appbrk{}प\appbrk{}न\appbrk{}म\appbrk{}र्शोघ्नं}}\Afootnote{H; \skt{स्निग्ध\appbrk{}म\appbrk{}ग्नि\appbrk{}दी\appbrk{}प\appbrk{}न\appbrk{}म\appbrk{}र्शोघ्नं}\ A}}\skt{ सृष्ट\-मू\-त्र\-पु\-री\-षञ्च }\edtext{\skt{तदुपसेवेत}}{\lemma{\skt{तदुपसेवेत}}\Afootnote{H; \skt{तदुपसेवेत}\ A}}\skt{~।}\pend

\pstart \ledsidenote{\passref{4.6.9}}\skt{दग्धेषु }\edtext{\skt{चार्श\-स\-स्त्व\-न\-ल\-स\-न्धु\-क्ष\-णार्थं}}{\lemma{\skt{चार्श\appbrk{}स\appbrk{}स्त्व\appbrk{}न\appbrk{}ल\appbrk{}स\appbrk{}न्धु\appbrk{}क्ष\appbrk{}णार्थं}}\Afootnote{H; \skt{चार्शः\appbrk{}स्व\appbrk{}भ्य\appbrk{}क्तोऽ\appbrk{}न\appbrk{}ल\appbrk{}स\appbrk{}न्धु\appbrk{}क्ष\appbrk{}णा\appbrk{}र्थ\appbrk{}म\appbrk{}नि\appbrk{}ल\appbrk{}प्र\appbrk{}को\appbrk{}प\appbrk{}सं\appbrk{}र\appbrk{}क्ष\appbrk{}णार्थं च}\ A}}\skt{ }\edtext{\skt{स्नेहादी}}{\lemma{\skt{स्नेहादी}}\Afootnote{H; \skt{स्नेहादीनां}\ A}}\skt{ }\edtext{\skt{सामान्यतो}}{\lemma{\skt{सामान्यतो}}\Afootnote{H; \skt{सामान्य}\ A}}\skt{ }\edtext{\skt{विशेषतश्च}}{\lemma{\skt{विशेषतश्च}}\Afootnote{H; \skt{तः}\ A}}\skt{ }\edtext{\edtext{\skt{क्रिया\-य\-थ\-मा\-से\-वेत}}{\lemma{\skt{क्रियायथम्}}\Afootnote{H; \skt{क्रियापथम्}\ A}}\skt{~। 𑑛 सर्प्पींषि}}{\lemma{\skt{आसेवेत~। 𑑛}}\Afootnote{H; \skt{उपसेवेत~। विशेषतस्तु वातार्शःसु}\ A}}\skt{ च वात\-ह\-र\-दी\-प\-नी\-य\-सि\-द्धानि }\edtext{\skt{हिंग्वादीनां चूर्ण्णैः}}{\lemma{\skt{हिंग्वादीनां}}\Afootnote{H; \skt{हिङ्ग्वादिभिश्}\ A}}\skt{ प्रतिसंसृज्य पिवेत्~। पित्ता\-र्श\-स्सु\-पृ\-थ\-क्प\-र्ण्ण्या\-दी\-ना\-ङ्क\-षा\-येण }\edtext{\skt{दीप\-नी\-य\-प्र\-ति\-वापं सर्प्पिः}}{\lemma{\skt{दीप\appbrk{}नी\appbrk{}य\appbrk{}प्र\appbrk{}ति\appbrk{}वापं}}\Afootnote{H; \skt{दीप\appbrk{}नी\appbrk{}य\appbrk{}प्र\appbrk{}ती\appbrk{}वापं}\ A}}\skt{ }\edtext{\skt{पिवेत्~।}}{\lemma{\skt{पिवेत्}}\Afootnote{H; \emph{om.}\ A}}\skt{ }\edtext{\skt{शोणितार्शस्सु }\edtext{\skt{मुरुं\-ग्या\-दी\-ना\-ङ्क\-षायैः}}{\lemma{\skt{मुरुंग्यादीनाङ्}}\Afootnote{H; \skt{मञ्जि\appbrk{}ष्ठा\appbrk{}मु\appbrk{}रु\appbrk{}ङ्ग्या\appbrk{}दीनां}\ A}}}{\lemma{\skt{शोणितार्शस्सु}}\Afootnote{H; \skt{शोणितार्शःसु}\ A}}\skt{~। }\edtext{\skt{श्लेष्मार्शस्सु सुर\-सा\-दी\-ना\-ङ्क\-षायैः}}{\lemma{\skt{श्लेष्मार्शस्सु}}\Afootnote{H; \skt{श्लेष्मार्शःसु}\ A}}\skt{~। उपद्रवाश्च }\edtext{\edtext{\skt{यथा\-स्व\-मु\-प\-च\-रेत्}}{\lemma{\skt{यथास्वम्}}\Afootnote{H; \skt{यथास्व}\ A}}\skt{~।}}{\lemma{\skt{उपचरेत्}}\Afootnote{H; \skt{मुपाचरेत्}\ A}}\edtext{}{\lemma{\skt{सर्प्पींषि च}}\Afootnote{H\textsuperscript{ac}; \skt{सर्पीषि च}\ A H\textsuperscript{pc}}}\edtext{}{\lemma{\skt{कषायैः}}\Afootnote{H; \skt{कषाये पाचयेत्}\ A}}\edtext{}{\lemma{\skt{कषायैः}}\Afootnote{H; \skt{कषाये}\ A}}\pend

\pstart \ledsidenote{\passref{4.6.10}}\skt{परञ्च यत्नमास्थाय गुदे }\edtext{\skt{शस्त्रा\-न्य\-व\-चा\-र\-येत्~। त}}{\lemma{\skt{शस्त्रा\appbrk{}न्य\appbrk{}व\appbrk{}चा\appbrk{}र\appbrk{}येत्~। त}}\Afootnote{H; \skt{क्षारा\appbrk{}ग्नि\appbrk{}श\appbrk{}स्त्रा\appbrk{}ण्य\appbrk{}व\appbrk{}चा\appbrk{}र\appbrk{}येत्, त}\ A}}\skt{ }\edtext{\edtext{\skt{द्विभ्र\-मा\-च्छा\-ण्ढ्य\-शो\-फ\-दा\-ह\-म\-द\-मू\-र्च्छा\-टो\-पा\-ना\-हा\-ती\-सा\-र\-प्र\-वा\-ह\-णानि}}{\lemma{\skt{द्विभ्रमाच्}}\Afootnote{H; \skt{द्विभ्रमाद्धि}\ A}}}{\lemma{\skt{छाण्ढ्य\appbrk{}शो\appbrk{}फ\appbrk{}दा\appbrk{}ह\appbrk{}म\appbrk{}द\appbrk{}मू\appbrk{}र्च्छा\appbrk{}टो\appbrk{}पा\appbrk{}ना\appbrk{}हा\appbrk{}ती\appbrk{}सा\appbrk{}र\appbrk{}प्र\appbrk{}वा\appbrk{}ह\appbrk{}णानि}}\Afootnote{H; \skt{षाण्ढ्य\appbrk{}शो\appbrk{}फ\appbrk{}दा\appbrk{}ह\appbrk{}म\appbrk{}द\appbrk{}मू\appbrk{}र्छा\appbrk{}टो\appbrk{}पा\appbrk{}ना\appbrk{}हा\appbrk{}ती\appbrk{}सा\appbrk{}र\appbrk{}प्र\appbrk{}वा\appbrk{}ह\appbrk{}णानि}\ A}}\skt{ भवन्ति~। मरणम्वा~॥}\pend

\pstart \ledsidenote{\passref{4.6.11}}\skt{अत ऊर्द्ध्वं यन्त्र\-प्र\-मा\-ण\-मु\-प\-दे\-क्ष्यामः~। तत्र यन्त्रं लौहन्दान्तं }\edtext{\skt{शांर्ग्गम्वा}}{\lemma{\skt{शांर्ग्गम्वा}}\Afootnote{H\textsuperscript{pc}; \skt{शार्ङ्गं वा}\ A; \skt{शांर्ग्गम्वा म्}\ H\textsuperscript{ac}}}\skt{ }\edtext{\skt{र्क्षम्वा}}{\lemma{\skt{र्क्षम्वा}}\Afootnote{A H\textsuperscript{pc}; \skt{वा}\ H\textsuperscript{ac}}}\skt{~। }\edtext{\skt{गोस्त\-ना\-का\-र\-ञ्च\-तु\-र\-ङ्गु\-ला\-य\-त\-म्प\-ञ्चा\-ङ्गु\-ल\-प\-रि\-णाहं}}{\lemma{\skt{पञ्चा\appbrk{}ङ्गु\appbrk{}ल\appbrk{}प\appbrk{}रि\appbrk{}णाहं}}\Afootnote{H; \skt{पञ्चा\appbrk{}ङ्गु\appbrk{}ल\appbrk{}प\appbrk{}रि\appbrk{}णाहं पुंसां, षड\appbrk{}ङ्गु\appbrk{}ल\appbrk{}प\appbrk{}रि\appbrk{}णाहं}\ A}}\skt{~। }\edtext{\edtext{\skt{नारी\-णा\-न्त\-ला\-य\-त\-मि\-त्येके}}{\lemma{\skt{इत्य्}}\Afootnote{H; \emph{om.}\ A}}}{\lemma{\skt{एके}}\Afootnote{H; \emph{om.}\ A}}\skt{~। तद्द्विछिद्रं~। }\edtext{\skt{दर्श\-ना\-र्थ\-मे\-क\-श्छि\-द्रन्तु}}{\lemma{\skt{एकश्छिद्रन्तु}}\Afootnote{H; \skt{एकच्छिद्रं तु}\ A}}\skt{ कर्म्मणि~। }\edtext{\skt{एकछिद्रे हि}}{\lemma{\skt{एकछिद्रे हि}}\Afootnote{H; \skt{एकद्वारे हि}\ A}}\skt{ }\edtext{\skt{शस्त्रा}}{\lemma{\skt{शस्त्रा}}\Afootnote{H; \skt{शस्त्र}\ A}}\skt{ }\edtext{\edtext{\skt{ग्निक्षा\-रा\-णा\-म\-ति\-क्रमो}}{\lemma{\skt{ग्निक्षाराणाम्}}\Afootnote{H; \skt{क्षाराग्नीना}\ A}}\skt{ न}}{\lemma{\skt{अतिक्रमो न}}\Afootnote{H; \skt{मतिक्रमो न}\ A}}\skt{ भवति~। छिद्रप्रमाणं तु }\edtext{\skt{त्र्यङ्गु\-ला\-य\-त\-म\-ङ्गु\-ष्ठो\-द\-र\-वि\-स्तारं}}{\lemma{\skt{त्र्यङ्गु\appbrk{}ला\appbrk{}य\appbrk{}त\appbrk{}म\appbrk{}ङ्गु\appbrk{}ष्ठो\appbrk{}द\appbrk{}र\appbrk{}वि\appbrk{}स्तारं}}\Afootnote{H; \skt{त्र्यङ्गु\appbrk{}ला\appbrk{}य\appbrk{}त\appbrk{}म\appbrk{}ङ्गु\appbrk{}ष्ठो\appbrk{}द\appbrk{}र\appbrk{}प\appbrk{}रि\appbrk{}णाहं}\ A}}\skt{ यद\-ङ्गु\-ल\-म\-व\-शिष्टं }\edtext{\edtext{\edtext{\edtext{\skt{तत्स्या\-द्द्व्य\-ङ्गु\-ल\-म\-ध\-स्ता\-दु\-प\-र्य\-र्द्धा\-ङ्गु\-लो\-छ्रि\-त\-म\-र्द्धा\-ङ्गु\-ष्ठो\-प\-रि\-वृ\-त्त\-क\-र्ण्णि\-क\-मेष}}{\lemma{\skt{तत्}}\Afootnote{H; \skt{त}\ A}}}{\lemma{\skt{स्याद्द्व्यङ्गुलम्}}\Afootnote{H; \skt{स्यार्धाङ्गुला}\ A}}}{\lemma{\skt{अध\appbrk{}स्ता\appbrk{}दु\appbrk{}प\appbrk{}र्य\appbrk{}र्द्धा\appbrk{}ङ्गु\appbrk{}लो\appbrk{}छ्रि\appbrk{}तम्}}\Afootnote{H; \skt{दध\appbrk{}स्ता\appbrk{}द\appbrk{}र्धा\appbrk{}ङ्गु\appbrk{}लोच्छ्रि}\ A}}}{\lemma{\skt{अर्द्धा\appbrk{}ङ्गु\appbrk{}ष्ठो\appbrk{}प\appbrk{}रि\appbrk{}वृ\appbrk{}त्त\appbrk{}क\appbrk{}र्ण्णि\appbrk{}कम्}}\Afootnote{H; \skt{तोप\appbrk{}रि\appbrk{}वृ\appbrk{}त्त\appbrk{}क\appbrk{}र्णि\appbrk{}कम्}\ A}}\skt{ }\edtext{\skt{यन्त्रा\-कृ\-ति\-स\-मासः~॥छ्॥}}{\lemma{\skt{यन्त्रा\appbrk{}कृ\appbrk{}ति\appbrk{}स\appbrk{}मासः~॥छ्}}\Afootnote{H; \skt{यन्त्रा\appbrk{}कृ\appbrk{}ति\appbrk{}स\appbrk{}मासः}\ A}}\pend

\pstart \ledsidenote{\passref{4.6.12}}\skt{अत }\edtext{\skt{ऊर्द्ध्व\-म\-र्श\-सा\-मा\-लेपो}}{\lemma{\skt{ऊर्द्ध्व\appbrk{}म\appbrk{}र्श\appbrk{}सा\appbrk{}मा\appbrk{}लेपो}}\Afootnote{H; \skt{ऊर्ध्व\appbrk{}म\appbrk{}र्श\appbrk{}सा\appbrk{}मा\appbrk{}ले\appbrk{}पान्}\ A}}\skt{ वक्ष्यामः~॥ }\edtext{\skt{स्नुहाक्षीरयुक्तं}}{\lemma{\skt{स्नुहाक्षीरयुक्तं}}\Afootnote{H; \skt{स्नुहीक्षीरयुक्तं}\ A}}\skt{ हरि\-द्रा\-चू\-र्ण्ण\-मा\-लेपः प्रथमः~। कुक्कु\-ट\-पु\-री\-ष\-गु\-ञ्जा\-ह\-रि\-द्रा\-पि\-प्प\-ली\-चू\-र्ण्ण\-मिति }\edtext{\skt{गोमूत्रपिष्टो}}{\lemma{\skt{गोमूत्रपिष्टो}}\Afootnote{H; \skt{गोमू\appbrk{}त्र\appbrk{}पि\appbrk{}त्त\appbrk{}पिष्टो}\ A}}\skt{ द्वितीयः~। }\edtext{\skt{दन्ती\-चि\-त्र\-क\-सु\-व\-र्च्चि\-का\-कल्को}}{\lemma{\skt{दन्ती\appbrk{}चि\appbrk{}त्र\appbrk{}क\appbrk{}सु\appbrk{}व\appbrk{}र्च्चि\appbrk{}का\appbrk{}कल्को}}\Afootnote{H; \skt{दन्ती\appbrk{}चि\appbrk{}त्र\appbrk{}क\appbrk{}सु\appbrk{}व\appbrk{}र्चि\appbrk{}का\appbrk{}ला\appbrk{}ङ्ग\appbrk{}ली\appbrk{}कल्को}\ A}}\skt{ }\edtext{\skt{गोमू\-त्र\-यु\-क्त\-स्तृ\-तीयः}}{\lemma{\skt{गोमू\appbrk{}त्र\appbrk{}यु\appbrk{}क्त\appbrk{}स्तृ\appbrk{}तीयः}}\Afootnote{H; \skt{वा गोपि\appbrk{}त्त\appbrk{}पि\appbrk{}ष्ट\appbrk{}स्तृ\appbrk{}तीयः}\ A}}\skt{~। }\edtext{\edtext{\skt{पिप्प\-ली\-सै\-न्ध\-व\-कु\-ष्ठ\-शि\-री\-ष\-क\-ल्कोऽ\-र्क\-क्षी\-र\-यु\-क्तश्च}}{\lemma{\skt{पिप्प\appbrk{}ली\appbrk{}सै\appbrk{}न्ध\appbrk{}व\appbrk{}कु\appbrk{}ष्ठ\appbrk{}शि\appbrk{}री\appbrk{}ष\appbrk{}कल्को}}\Afootnote{H; \skt{पिप्प\appbrk{}ली\appbrk{}सै\appbrk{}न्ध\appbrk{}व\appbrk{}कु\appbrk{}ष्ठ\appbrk{}शि\appbrk{}री\appbrk{}ष\appbrk{}फ\appbrk{}ल\appbrk{}कल्कः}\ A}}}{\lemma{\skt{ऽर्क\appbrk{}क्षी\appbrk{}र\appbrk{}यु\appbrk{}क्तश्च}}\Afootnote{H; \skt{स्नुही\appbrk{}क्षी\appbrk{}र\appbrk{}पि\appbrk{}ष्टोऽ\appbrk{}र्क\appbrk{}क्षी\appbrk{}र\appbrk{}पिष्टो वा}\ A}}\skt{ चतुर्थः~॥ कासी }\edtext{\skt{सह\-रि\-ता\-ल\-सै\-न्ध\-वा\-श्व\-मा\-र\-क\-वि\-ड\-ङ्ग\-पू\-ती\-क\-कृ\-त\-वे\-ध\-ना\-ज\-म्बु\-का\-श्व\-को\-त्त\-म\-कि\-णि\-ही\-स्नु\-हा\-प\-यस्सु च}}{\lemma{\skt{सह\appbrk{}रि\appbrk{}ता\appbrk{}ल\appbrk{}सै\appbrk{}न्ध\appbrk{}वा\appbrk{}श्व\appbrk{}मा\appbrk{}र\appbrk{}क\appbrk{}वि\appbrk{}ड\appbrk{}ङ्ग\appbrk{}पू\appbrk{}ती\appbrk{}क\appbrk{}कृ\appbrk{}त\appbrk{}वे\appbrk{}ध\appbrk{}ना\appbrk{}ज\appbrk{}म्बु\appbrk{}का\appbrk{}श्व\appbrk{}को\appbrk{}त्त\appbrk{}म\appbrk{}कि\appbrk{}णि\appbrk{}ही\appbrk{}स्नु\appbrk{}हा\appbrk{}प\appbrk{}यस्सु च}}\Afootnote{H; \skt{सह\appbrk{}रि\appbrk{}ता\appbrk{}ल\appbrk{}सै\appbrk{}न्ध\appbrk{}वा\appbrk{}श्व\appbrk{}मा\appbrk{}र\appbrk{}क\appbrk{}वि\appbrk{}ड\appbrk{}ङ्ग\appbrk{}पू\appbrk{}ती\appbrk{}क\appbrk{}कृ\appbrk{}त\appbrk{}वे\appbrk{}ध\appbrk{}न\appbrk{}ज\appbrk{}म्ब्व\appbrk{}र्को\appbrk{}त्त\appbrk{}मा\appbrk{}र\appbrk{}णी\appbrk{}द\appbrk{}न्ती\appbrk{}चि\appbrk{}त्र\appbrk{}का\appbrk{}ल\appbrk{}र्क\appbrk{}स्नु\appbrk{}ही\appbrk{}प\appbrk{}यःसु}\ A}}\skt{~। }\edtext{\skt{तैल\-म्वि\-प\-क्व\-म\-भ्य\-ञ्ज\-ने\-ना\-र्शसः}}{\lemma{\skt{विप\appbrk{}क्व\appbrk{}म\appbrk{}भ्य\appbrk{}ञ्ज\appbrk{}ने\appbrk{}ना\appbrk{}र्शसः}}\Afootnote{\skt{विप\appbrk{}क्व\appbrk{}म\appbrk{}भ्य\appbrk{}ञ्ज\appbrk{}ने\appbrk{}नार्शः}\ A; \skt{विप\appbrk{}क्व\appbrk{}म\appbrk{}भ्य\appbrk{}ञ्ज\appbrk{}ने\appbrk{}ना\appbrk{}र्शसः}\ H}}\skt{ }\edtext{\edtext{\skt{सातयेयुरिति}}{\lemma{\skt{सातयेयुर्}}\Afootnote{H; \skt{शातयति}\ A}}\skt{~॥~॥ ं}}{\lemma{\skt{इति~॥~॥ ं}}\Afootnote{H; \emph{om.}\ A}}\skt{~॥}\pend

\pstart \ledsidenote{\passref{4.6.13}}\skt{अत }\edtext{\edtext{\skt{ऊर्द्ध\-म\-दृ\-श्ये\-ष्व\-र्शस्सु}}{\lemma{\skt{ऊर्द्धम्}}\Afootnote{H; \skt{ऊर्ध्वं भेषजसाध्}\ A}}}{\lemma{\skt{अदृश्येष्व्}}\Afootnote{H; \skt{येष्वदृश्येष्व्}\ A}}\skt{ }\edtext{\skt{योगा\-न्य\-प\-न\-यार्थं}}{\lemma{\skt{योगा\appbrk{}न्य\appbrk{}प\appbrk{}न\appbrk{}यार्थं}}\Afootnote{H; \skt{योगान्यापनार्थं}\ A}}\skt{ वक्ष्यामः॥ }\edtext{\skt{तत्र}}{\lemma{\skt{तत्र}}\Afootnote{H; \skt{प्रातः प्रतर्}\ A}}\skt{ }\edtext{\skt{गुड\-ह\-री\-त\-की\-मा\-से\-वेत}}{\lemma{\skt{गुड\appbrk{}ह\appbrk{}री\appbrk{}त\appbrk{}की\appbrk{}मा\appbrk{}से\appbrk{}वेत}}\Afootnote{\skt{गुड\appbrk{}ह\appbrk{}री\appbrk{}त\appbrk{}की\appbrk{}मा\appbrk{}से\appbrk{}वेत}\ A H}}\skt{ ब्रह्मचारी~। गोमू\-त्र\-द्रो\-ण\-सि\-द्धम्वा }\edtext{\skt{हरी\-त\-की\-श\-त\-मु\-प\-से\-वेत य}}{\lemma{\skt{उपसेवेत य}}\Afootnote{H; \skt{प्रातः प्रातर्य}\ A}}\skt{ थावलं क्षौद्रेण~। अपा\-मा\-र्ग्ग\-मू\-लम्वा तण्डुलोदकेन~। }\edtext{\skt{शता\-व\-री\-मू\-ल\-क\-ल्कम्वा}}{\lemma{\skt{शता\appbrk{}व\appbrk{}री\appbrk{}मू\appbrk{}ल\appbrk{}क\appbrk{}ल्कम्वा}}\Afootnote{H; \skt{सक्षौ\appbrk{}द्र\appbrk{}म\appbrk{}ह\appbrk{}रहः, शता\appbrk{}व\appbrk{}री\appbrk{}मू\appbrk{}ल\appbrk{}कल्कं वा}\ A}}\skt{ क्षीरेण~। }\edtext{\skt{चित्त्रकचूर्ण्णम्वा}}{\lemma{\skt{चित्त्रकचूर्ण्णम्वा}}\Afootnote{H; \skt{चित्र\appbrk{}क\appbrk{}चू\appbrk{}र्ण\appbrk{}युक्तं वा}\ A}}\skt{ }\edtext{\skt{सीधुना}}{\lemma{\skt{सीधुना}}\Afootnote{H; \skt{सीधुं}\ A}}\skt{~। }\edtext{\skt{मधु\-रा\-र्द्ध\-भ\-ल्ला\-त\-क\-चू\-र्ण्ण\-यु\-क्तम्वा}}{\lemma{\skt{मधु\appbrk{}रा\appbrk{}र्द्ध\appbrk{}भ\appbrk{}ल्ला\appbrk{}त\appbrk{}क\appbrk{}चू\appbrk{}र्ण्ण\appbrk{}यु\appbrk{}क्तम्वा}}\Afootnote{H; \skt{परार्ध्यं भल्ला\appbrk{}त\appbrk{}चू\appbrk{}र्ण\appbrk{}युक्तं वा}\ A}}\skt{ }\edtext{\skt{शक्तून्}}{\lemma{\skt{शक्तून्}}\Afootnote{H; \skt{सक्तुमन्}\ A}}\skt{~। }\edtext{\skt{ल}}{\lemma{\skt{ल}}\Afootnote{H; \skt{थमल}\ A}}\skt{ }\edtext{\skt{वणो}}{\lemma{\skt{वणो}}\Afootnote{H; \skt{वणं}\ A}}\skt{ तक्रेण~। }\edtext{\skt{कलसे}}{\lemma{\skt{कलसे}}\Afootnote{H; \skt{कलशे}\ A}}\skt{ }\edtext{\edtext{\skt{वान्त\-श्चि\-त्त्र\-क\-मू\-ल\-क\-ल्क\-लिप्ते}}{\lemma{\skt{वान्तश्}}\Afootnote{H; \skt{वाऽन्तश्}\ A}}}{\lemma{\skt{चित्त्र\appbrk{}क\appbrk{}मू\appbrk{}ल\appbrk{}क\appbrk{}ल्क\appbrk{}लिप्ते}}\Afootnote{H; \skt{चित्र\appbrk{}क\appbrk{}मू\appbrk{}ल\appbrk{}क\appbrk{}ल्का\appbrk{}व\appbrk{}लिप्ते}\ A}}\skt{ निषिक्तं }\edtext{\edtext{\skt{तक्र\-म\-न\-म्ल\-म्पा\-न\-भो\-ज\-ने\-षू\-प\-यु\-ञ्जीत}}{\lemma{\skt{तक्रम्}}\Afootnote{H; \skt{तक्र}\ A}}}{\lemma{\skt{अनम्लम्}}\Afootnote{H; \skt{मम्लमनम्लं वा}\ A}}\skt{~। एष एव }\edtext{\skt{भार्ग्ग्या\-स्फा\-ता\-गु\-डू\-चीषु}}{\lemma{\skt{भार्ग्ग्या\appbrk{}स्फा\appbrk{}ता\appbrk{}गु\appbrk{}डू\appbrk{}चीषु}}\Afootnote{H; \skt{भार्ग्या\appbrk{}स्फो\appbrk{}ता\appbrk{}य\appbrk{}वा\appbrk{}न्या\appbrk{}म\appbrk{}ल\appbrk{}क\appbrk{}गु\appbrk{}डू\appbrk{}चीषु}\ A}}\skt{ }\edtext{\skt{तक्रकल्पः।}}{\lemma{\skt{तक्रकल्पः}}\Afootnote{A; \skt{तक्रकल्पः।𑑛}\ H}}\skt{ }\edtext{\skt{पिप्प\-ली\-पि\-प्प\-ली\-मूलं}}{\lemma{\skt{पिप्प\appbrk{}ली\appbrk{}पि\appbrk{}प्प\appbrk{}ली\appbrk{}मूलं}}\Afootnote{H; \skt{पिप्प\appbrk{}ली\appbrk{}पि\appbrk{}प्प\appbrk{}ली\appbrk{}मूल}\ A}}\skt{ }\edtext{\skt{विड\-ङ्ग\-शु\-ण्ठी\-ह\-री\-त\-कीषु वा}}{\lemma{\skt{विड\appbrk{}ङ्ग\appbrk{}शु\appbrk{}ण्ठी\appbrk{}ह\appbrk{}री\appbrk{}त\appbrk{}कीषु वा}}\Afootnote{H; \skt{चव्य\appbrk{}चि\appbrk{}त्र\appbrk{}क\appbrk{}वि\appbrk{}ड\appbrk{}ङ्ग\appbrk{}शु\appbrk{}ण्ठी\appbrk{}ह\appbrk{}री\appbrk{}त\appbrk{}कीषु च}\ A}}\skt{~। }\edtext{\skt{पूर्व्ववदेव}}{\lemma{\skt{पूर्व्ववदेव}}\Afootnote{H; \skt{पूर्ववदेव}\ A}}\skt{ निरन्नो वा }\edtext{\skt{तक्र\-म\-ह\-र\-ह\-र्म्मा\-स\-मा\-से\-वेत}}{\lemma{\skt{अह\appbrk{}र\appbrk{}ह\appbrk{}र्म्मा\appbrk{}स\appbrk{}मा\appbrk{}से\appbrk{}वेत}}\Afootnote{H; \skt{अह\appbrk{}र\appbrk{}ह\appbrk{}र्मा\appbrk{}स\appbrk{}मु\appbrk{}प\appbrk{}से\appbrk{}वेत}\ A}}\skt{~। }\edtext{\skt{शृङ्ग\-वे\-र\-पु\-न\-र्ण्ण\-सिद्धं वा}}{\lemma{\skt{शृङ्ग\appbrk{}वे\appbrk{}र\appbrk{}पु\appbrk{}न\appbrk{}र्ण्ण\appbrk{}सिद्धं वा}}\Afootnote{H; \skt{शृङ्ग\appbrk{}वे\appbrk{}र\appbrk{}पु\appbrk{}न\appbrk{}र्न\appbrk{}वा\appbrk{}चि\appbrk{}त्र\appbrk{}क\appbrk{}क\appbrk{}षा\appbrk{}य\appbrk{}सिद्धं वा}\ A}}\skt{ पयः~। }\edtext{\skt{क्रटजमूलत्व}}{\lemma{\skt{क्रटजमूलत्व}}\Afootnote{H; \skt{कृटजमूलत्व}\ A}}\skt{ क्फाणितम्वा पिप्प\-ल्या\-दि\-प्र\-ती\-वापं क्षौद्रेण }\edtext{\skt{वातव्याध्योक्तं}}{\lemma{\skt{वातव्याध्योक्तं}}\Afootnote{H; \skt{महा\appbrk{}वा\appbrk{}त\appbrk{}व्या\appbrk{}ध्युक्तं}\ A}}\skt{ }\edtext{\skt{हिंग्वादिवटके}}{\lemma{\skt{हिंग्वादिवटके}}\Afootnote{H; \skt{हिङ्ग्वादिचूर्णम्}\ A}}\skt{ उपसेवेत~। }\edtext{\skt{तक्रा\-हा\-र\-क्षी\-रा\-हारो वा}}{\lemma{\skt{तक्रा\appbrk{}हा\appbrk{}र\appbrk{}क्षी\appbrk{}रा\appbrk{}हारो वा}}\Afootnote{H; \skt{तक्राहारः क्षीराहारो वा}\ A}}\skt{~। }\edtext{\skt{चित्र\-क\-मू\-ल\-क\-षा\-य\-सिद्धं वा}}{\lemma{\skt{चित्र\appbrk{}क\appbrk{}मू\appbrk{}ल\appbrk{}क\appbrk{}षा\appbrk{}य\appbrk{}सिद्धं वा}}\Afootnote{H; \skt{क्षार\appbrk{}ल\appbrk{}व\appbrk{}णां\appbrk{}श्चि\appbrk{}त्र\appbrk{}क\appbrk{}मू\appbrk{}ल\appbrk{}क्षा\appbrk{}रो\appbrk{}द\appbrk{}क\appbrk{}सिद्धं वा पयः पला\appbrk{}श\appbrk{}त\appbrk{}रु\appbrk{}क्षा\appbrk{}र\appbrk{}सिद्धं वा}\ A}}\skt{ }\edtext{\edtext{\skt{कुल्माषाम्भक्ष}}{\lemma{\skt{कुल्माषाम्}}\Afootnote{H; \emph{om.}\ A}}}{\lemma{\skt{भक्ष}}\Afootnote{H; \emph{om.}\ A}}\skt{ }\edtext{\skt{येत्}}{\lemma{\skt{येत्}}\Afootnote{H; \emph{om.}\ A}}\skt{~। }\edtext{\skt{पला\-श\-त\-रु\-ण\-सि\-द्धम्वा}}{\lemma{\skt{पला\appbrk{}श\appbrk{}त\appbrk{}रु\appbrk{}ण\appbrk{}सि\appbrk{}द्धम्वा}}\Afootnote{H; \skt{पला\appbrk{}श\appbrk{}त\appbrk{}रु\appbrk{}क्षा\appbrk{}र\appbrk{}सि\appbrk{}द्धान्वा कुल्माषान्}\ A}}\skt{ }\edtext{\skt{पाट\-ल्य\-पा\-मा\-र्ग्ग\-वृ\-ह\-ती\-प\-ला\-श\-क्षा\-रम्वा}}{\lemma{\skt{पाट\appbrk{}ल्य\appbrk{}पा\appbrk{}मा\appbrk{}र्ग्ग\appbrk{}वृ\appbrk{}ह\appbrk{}ती\appbrk{}प\appbrk{}ला\appbrk{}श\appbrk{}क्षा\appbrk{}रम्वा}}\Afootnote{H; \skt{पाट\appbrk{}ला\appbrk{}पा\appbrk{}मा\appbrk{}र्ग\appbrk{}बृं\appbrk{}ह\appbrk{}ती\appbrk{}प\appbrk{}ला\appbrk{}श\appbrk{}क्षारं वा}\ A}}\skt{ }\edtext{\skt{परि\-स्रु\-त\-म\-ह\-र\-ह\-र्घृ\-त\-युक्तं}}{\lemma{\skt{अह\appbrk{}र\appbrk{}ह\appbrk{}र्घृ\appbrk{}त\appbrk{}युक्तं}}\Afootnote{H; \skt{अह\appbrk{}र\appbrk{}ह\appbrk{}र्घृ\appbrk{}त\appbrk{}सं\appbrk{}सृष्टं}\ A}}\skt{~। }\edtext{\skt{कुट\-ज\-च\-न्दा\-क\-मू\-ल\-क\-ल्कम्वा}}{\lemma{\skt{कुट\appbrk{}ज\appbrk{}च\appbrk{}न्दा\appbrk{}क\appbrk{}मू\appbrk{}ल\appbrk{}क\appbrk{}ल्कम्वा}}\Afootnote{H; \skt{कुट\appbrk{}ज\appbrk{}ब\appbrk{}न्दा\appbrk{}क\appbrk{}मू\appbrk{}ल\appbrk{}कल्कं वा}\ A}}\skt{~। }\edtext{\skt{तक्रेण}}{\lemma{\skt{तक्रेण}}\Afootnote{H; \skt{तक्रेण, चित्र\appbrk{}क\appbrk{}पू\appbrk{}ती\appbrk{}क\appbrk{}ना\appbrk{}ग\appbrk{}र\appbrk{}कल्कं वा पूतीकक्षारेण}\ A}}\skt{~। क्षारो\-द\-क\-सि\-द्धम्वा स र्प्पिः }\edtext{\skt{पिप्प\-ल्या\-दि\-प्र\-ति\-वापं}}{\lemma{\skt{पिप्प\appbrk{}ल्या\appbrk{}दि\appbrk{}प्र\appbrk{}ति\appbrk{}वापं}}\Afootnote{H; \skt{पिप्प\appbrk{}ल्या\appbrk{}दि\appbrk{}प्र\appbrk{}ती\appbrk{}वापं}\ A}}\skt{~। }\edtext{\skt{कृष्ण\-ति\-ल\-प्र\-कु\-ञ्चम्वा}}{\lemma{\skt{कृष्ण\appbrk{}ति\appbrk{}ल\appbrk{}प्र\appbrk{}कु\appbrk{}ञ्चम्वा}}\Afootnote{H; \skt{कृष्ण\appbrk{}ति\appbrk{}ल\appbrk{}प्र\appbrk{}सृतं प्रकुञ्चं वा}\ A}}\skt{ }\edtext{\skt{शीत\-तो\-या\-नु\-पा\-न\-म्प्रातः}}{\lemma{\skt{शीत\appbrk{}तो\appbrk{}या\appbrk{}नु\appbrk{}पा\appbrk{}नम्}}\Afootnote{H; \emph{om.}\ A}}\skt{ प्रात\-रु\-प\-से\-वेत~। }\edtext{\edtext{\edtext{\skt{एव\-मे\-भि\-र्न्न\-व\-म्व\-र्द्ध\-तेऽ\-ग्निश्च}}{\lemma{\skt{एवम्}}\Afootnote{H; \skt{शीतो\appbrk{}द\appbrk{}का\appbrk{}नु\appbrk{}पा\appbrk{}नम्}\ A}}}{\lemma{\skt{न्नवम्}}\Afootnote{H; \skt{अभि}\ A}}}{\lemma{\skt{ऽग्निश्च}}\Afootnote{H\textsuperscript{pc}; \skt{ऽग्नि}\ A; \skt{ग्निश्च}\ H\textsuperscript{ac}}}\skt{ }\edtext{\skt{भवति}}{\lemma{\skt{भवति}}\Afootnote{H; \skt{र}\ A}}\skt{ }\edtext{\skt{अर्शांसि}}{\lemma{\skt{अर्शांसि}}\Afootnote{H; \skt{र्शांसि}\ A}}\skt{ }\edtext{\skt{चोपशाम्यति}}{\lemma{\skt{चोपशाम्यति}}\Afootnote{H; \skt{चोपशाम्यन्ति}\ A}}\skt{~।}\pend

\pstart \ledsidenote{\passref{4.6.14}}\skt{द्वि }\edtext{\skt{पञ्च\-मू\-ल\-द\-न्ती\-चि\-त्त्र\-क\-प\-थ्या\-ना\-न्तु\-ला\-मा\-हृत्य}}{\lemma{\skt{पञ्च\appbrk{}मू\appbrk{}ल\appbrk{}द\appbrk{}न्ती\appbrk{}चि\appbrk{}त्त्र\appbrk{}क\appbrk{}प\appbrk{}थ्या\appbrk{}नान्}}\Afootnote{H; \skt{पञ्च\appbrk{}मू\appbrk{}ली\appbrk{}द\appbrk{}न्ती\appbrk{}चि\appbrk{}त्र\appbrk{}क\appbrk{}प\appbrk{}थ्यानां}\ A}}\skt{ }\edtext{\skt{जलद्रोणे}}{\lemma{\skt{जलद्रोणे}}\Afootnote{H; \skt{जलचतुर्द्रोणे}\ A}}\skt{ }\edtext{\skt{चि पाचयेत्}}{\lemma{\skt{चि पाचयेत्}}\Afootnote{H; \skt{विपाचयेत्}\ A}}\skt{~। }\edtext{\skt{पादा\-व\-शि\-ष्ट\-ङ्क\-षा\-य\-मा\-दाय}}{\lemma{\skt{पादावशिष्टङ्}}\Afootnote{H; \skt{ततः पादावशिष्टं}\ A}}\skt{ सुशी\-त\-ङ्गु\-ड\-तु\-लया सहोन्मिश्र्य घृतभाजने }\edtext{\skt{निषिच्य}}{\lemma{\skt{निषिच्य}}\Afootnote{H; \skt{निःक्षि}\ A}}\skt{ }\edtext{\skt{यवपल्लेप्लवे}}{\lemma{\skt{यवपल्लेप्लवे}}\Afootnote{H; \skt{प्य मा}\ A}}\skt{ समुपेक्षेत }\edtext{\skt{ततः}}{\lemma{\skt{ततः}}\Afootnote{H; \skt{यवपल्ले, ततः}\ A}}\skt{ प्रातः }\edtext{\skt{प्रात\-र्म्मा\-त्रा\-म्पा\-य\-येत्}}{\lemma{\skt{पाययेत्}}\Afootnote{H; \skt{पाययेत}\ A}}\skt{~। }\edtext{\skt{तेना\-र्श\-सो\-ग्र\-ह\-णी\-दो\-ष\-पा\-ण्डु\-रो\-गौ\-दा\-व\-र्त्त\-श्वा\-सका सा}}{\lemma{\skt{तेना\appbrk{}र्श\appbrk{}सो\appbrk{}ग्र\appbrk{}ह\appbrk{}णी\appbrk{}दो\appbrk{}ष\appbrk{}पा\appbrk{}ण्डु\appbrk{}रो\appbrk{}गौ\appbrk{}दा\appbrk{}व\appbrk{}र्त्त\appbrk{}श्वा\appbrk{}सका सा}}\Afootnote{H; \skt{तेना\appbrk{}र्शो\appbrk{}ग्र\appbrk{}ह\appbrk{}णी\appbrk{}दो\appbrk{}ष\appbrk{}पा\appbrk{}ण्डु\appbrk{}रो\appbrk{}गो\appbrk{}दा\appbrk{}वर्ता}\ A}}\skt{ }\edtext{\skt{रोचकाग्नि}}{\lemma{\skt{रोचकाग्नि}}\Afootnote{H; \skt{रोचका}\ A}}\skt{ }\edtext{\skt{दौर्व्वल्यानि न}}{\lemma{\skt{दौर्व्वल्यानि न}}\Afootnote{H; \skt{न}\ A}}\skt{ }\edtext{\skt{भवन्ति}}{\lemma{\skt{भवन्ति}}\Afootnote{\skt{भवन्ति दीप्त\appbrk{}श्चा\appbrk{}ग्नि\appbrk{}र्भ\appbrk{}वति}\ A; \skt{भवन्ति~॥ ं}\ H}}\skt{~॥~॥}\pend

\pstart \ledsidenote{\passref{4.6.15}}\skt{पिप्प\-ली\-म\-रि\-च\-वि\-डं\-गै\-ल\-वा\-लु\-क\-लो\-ध्राणां }\edtext{\edtext{\skt{द्विप\-लि\-का\-न्भा\-गा\-नि\-न्द्र\-वा\-रुण्याः}}{\lemma{\skt{द्विपलिकान्}}\Afootnote{H; \skt{द्वे द्वे पले}\ A}}}{\lemma{\skt{भागा\appbrk{}नि\appbrk{}न्द्र\appbrk{}वा\appbrk{}रुण्याः}}\Afootnote{H; \skt{इन्द्रवारुण्याः}\ A}}\skt{ पञ्चपलानि }\edtext{\skt{हरी\-त\-क्या\-म\-ल\-क\-पि\-त्थ\-फ\-लानां}}{\lemma{\skt{हरी\appbrk{}त\appbrk{}क्या\appbrk{}म\appbrk{}ल\appbrk{}क\appbrk{}पि\appbrk{}त्थ\appbrk{}फ\appbrk{}लानां}}\Afootnote{H; \skt{कपित्थमध्यस्य दश, पथ्या\appbrk{}फ\appbrk{}ला\appbrk{}ना\appbrk{}म\appbrk{}र्ध\appbrk{}प्रस्थः, प्रस्थो धात्रीफलानां}\ A}}\skt{ }\edtext{\skt{दश}}{\lemma{\skt{दश}}\Afootnote{H; \emph{om.}\ A}}\skt{ }\edtext{\skt{दश}}{\lemma{\skt{दश}}\Afootnote{H; \emph{om.}\ A}}\skt{ एतदैकध्यं जलचतुर्द्रोणे वि पाच्य पादावशेषं }\edtext{\skt{विश्राव्य}}{\lemma{\skt{विश्राव्य}}\Afootnote{H; \skt{परिस्राव्य}\ A}}\skt{ सुशीतं }\edtext{\skt{गुड\-तु\-ला\-द्व\-यो\-न्न्मिश्र्य}}{\lemma{\skt{गुड\appbrk{}तु\appbrk{}ला\appbrk{}द्व\appbrk{}यो\appbrk{}न्न्मिश्र्य}}\Afootnote{H; \skt{गुड\appbrk{}तु\appbrk{}ला\appbrk{}द्व\appbrk{}ये\appbrk{}नो\appbrk{}न्मिश्र्य}\ A}}\skt{ घृतभाजने }\edtext{\skt{निक्षिच्य}}{\lemma{\skt{निक्षिच्य}}\Afootnote{H; \skt{निःक्षिप्य}\ A}}\skt{ }\edtext{\edtext{\skt{पक्ष\-मा\-त्र\-मु\-प\-क्षि\-येत}}{\lemma{\skt{पक्षमात्रम्}}\Afootnote{H; \skt{पक्षमु}\ A}}}{\lemma{\skt{उपक्षियेत}}\Afootnote{H; \skt{पेक्षेत}\ A}}\skt{ यवपल्ले ततः प्रातः प्रात\-र्य\-था\-व\-ल\-मु\-प\-युञ्जी त~। एष }\edtext{\skt{खल्व\-रि\-ष्टोऽ\-र्शो\-ग्र\-ह\-णी\-प्ली\-ह\-हृ\-त्पा\-ण्डु\-रो\-ग\-शो\-फ\-कु\-ष्ठ\-गु\-ल्मो\-द\-र\-क्रि\-मि\-हरो}}{\lemma{\skt{ऽर्शो\appbrk{}ग्र\appbrk{}ह\appbrk{}णी\appbrk{}प्ली\appbrk{}ह\appbrk{}हृ\appbrk{}त्पा\appbrk{}ण्डु\appbrk{}रो\appbrk{}ग\appbrk{}शो\appbrk{}फ\appbrk{}कु\appbrk{}ष्ठ\appbrk{}गु\appbrk{}ल्मो\appbrk{}द\appbrk{}र\appbrk{}क्रि\appbrk{}मि\appbrk{}हरो}}\Afootnote{H; \skt{प्लीहा\appbrk{}ग्नि\appbrk{}ष\appbrk{}ङ्गा\appbrk{}र्शो\appbrk{}ग्र\appbrk{}ह\appbrk{}णी\appbrk{}हृ\appbrk{}त्पा\appbrk{}ण्डु\appbrk{}रो\appbrk{}ग\appbrk{}शो\appbrk{}फ\appbrk{}कु\appbrk{}ष्ठ\appbrk{}गु\appbrk{}ल्मो\appbrk{}द\appbrk{}र\appbrk{}कृ\appbrk{}मि\appbrk{}हरो}\ A}}\skt{ वल\-व\-र्ण्ण\-क\-र\-श्चेति~॥}\pend

\pstart \ledsidenote{\passref{4.6.16}}\skt{तत्र वातप्रायेषु }\edtext{\skt{स्नेह\-स्वे\-द\-स्ने\-ह\-वि\-रे\-च\-नास्था}}{\lemma{\skt{स्नेह\appbrk{}स्वे\appbrk{}द\appbrk{}स्ने\appbrk{}ह\appbrk{}वि\appbrk{}रे\appbrk{}च\appbrk{}नास्था}}\Afootnote{H; \skt{स्नेह\appbrk{}स्वे\appbrk{}द\appbrk{}व\appbrk{}म\appbrk{}न\appbrk{}वि\appbrk{}रे\appbrk{}च\appbrk{}नास्था}\ A}}\skt{ पना\-नु\-वा\-स\-न\-म\-प्र\-ति\-सिद्धं~। पित्तजेषु }\edtext{\skt{विरेचनमेव}}{\lemma{\skt{एव}}\Afootnote{H; \skt{ए}\ A}}\skt{~। }\edtext{\skt{संशमनं}}{\lemma{\skt{संशमनं}}\Afootnote{H; \skt{वं}\ A}}\skt{ रक्तजेषु~। }\edtext{\skt{कफजेषु}}{\lemma{\skt{कफजेषु}}\Afootnote{H; \skt{संशमनं, कफजेषु}\ A}}\skt{ शृङ्ग\-वे\-र\-कु\-ल\-त्थो\-प\-योगो }\edtext{\skt{यथा\-स्वौ\-ष\-ध\-सिद्धं}}{\lemma{\skt{यथा\appbrk{}स्वौ\appbrk{}ष\appbrk{}ध\appbrk{}सिद्धं}}\Afootnote{H; \skt{सर्वदोषहरं यथोक्तं सर्वजेषु, यथा\appbrk{}स्वौ\appbrk{}ष\appbrk{}धि\appbrk{}सिद्धं च}\ A}}\skt{ }\edtext{\skt{सर्व्वजेषु~॥ ं}}{\lemma{\skt{सर्व्वजेषु~॥ ं}}\Afootnote{H; \skt{पयः सर्वेष्विति}\ A}}\skt{~॥}\pend

\pstart \ledsidenote{\passref{4.6.17}}\skt{अत ऊर्द्ध्व\-म्भ\-ल्ला\-त\-क\-वि\-धा\-न\-मु\-प\-दे\-क्ष्यामः~॥ }\edtext{\edtext{\skt{भल्ला\-त\-का\-न्प\-रि\-प\-क्वा\-न\-नु\-प\-ह\-ता\-ना\-हृत्य}}{\lemma{\skt{भल्लातकान्}}\Afootnote{H; \skt{भल्लातकानि}\ A}}}{\lemma{\skt{परि\appbrk{}प\appbrk{}क्वा\appbrk{}न\appbrk{}नु\appbrk{}प\appbrk{}ह\appbrk{}ता\appbrk{}ना\appbrk{}हृत्य}}\Afootnote{H; \skt{परि\appbrk{}प\appbrk{}क्वा\appbrk{}न्य\appbrk{}नु\appbrk{}प\appbrk{}ह\appbrk{}ता\appbrk{}न्या\appbrk{}हृत्य}\ A}}\skt{ तत एकमादाय }\edtext{\skt{त्रिधा}}{\lemma{\skt{त्रिधा}}\Afootnote{H; \skt{द्विधा त्रिधा}\ A}}\skt{ }\edtext{\skt{चतुर्द्धा}}{\lemma{\skt{चतुर्द्धा}}\Afootnote{H; \skt{चतुर्धा वा}\ A}}\skt{ छेदयित्वा कषायकल्पेन विपाच्य तस्य कषायस्य शुक्तिमनुष्णां }\edtext{\skt{घृता\-भ्य\-क्त\-ता\-ल्वो\-ष्ठ\-जिह्वः}}{\lemma{\skt{घृता\appbrk{}भ्य\appbrk{}क्त\appbrk{}ता\appbrk{}ल्वो\appbrk{}ष्ठ\appbrk{}जिह्वः}}\Afootnote{H; \skt{घृता\appbrk{}भ्य\appbrk{}क्त\appbrk{}ता\appbrk{}लु\appbrk{}जि\appbrk{}ह्वौष्ठः}\ A}}\skt{ प्रातः }\edtext{\skt{प्रात\-रु\-प\-से\-वेत}}{\lemma{\skt{उपसेवेत}}\Afootnote{H; \skt{उपसेवेत, ततो}\ A}}\skt{~। अपराह्णे }\edtext{\skt{क्षीर\-स\-र्प्पि\-रो\-दन}}{\lemma{\skt{क्षीर\appbrk{}स\appbrk{}र्प्पि\appbrk{}रो\appbrk{}दन}}\Afootnote{H; \skt{क्षीरं सर्पिरोदन}\ A}}\skt{ इत्याहारः~। }\edtext{\skt{एव\-मे\-कै\-क\-म्व\-र्द्ध\-ये\-द्या\-व\-त्पञ्च}}{\lemma{\skt{पञ्च}}\Afootnote{H; \skt{पञ्चे}\ A}}\skt{~। }\edtext{\edtext{\skt{ततश्चोर्द्ध}}{\lemma{\skt{ततश्}}\Afootnote{H; \skt{ति, त}\ A}}}{\lemma{\skt{चोर्द्ध}}\Afootnote{H; \skt{तः}\ A}}\skt{ पञ्च }\edtext{\skt{पञ्चा\-भि\-व\-र्द्ध\-ये\-दा\-स\-ता\-दिति}}{\lemma{\skt{पञ्चा\appbrk{}भि\appbrk{}व\appbrk{}र्द्ध\appbrk{}ये\appbrk{}दा\appbrk{}स\appbrk{}ता\appbrk{}दिति}}\Afootnote{H; \skt{पञ्चा\appbrk{}भि\appbrk{}व\appbrk{}र्ध\appbrk{}ये\appbrk{}द्या\appbrk{}व\appbrk{}त्स\appbrk{}प्त\appbrk{}ति\appbrk{}रिति}\ A}}\skt{~। प्राप्य च }\edtext{\skt{शत}}{\lemma{\skt{शत}}\Afootnote{H; \skt{सप्}\ A}}\skt{ }\edtext{\edtext{\skt{शत\-म\-प\-क\-र्ष\-ये\-द्भूयः}}{\lemma{\skt{शतम्}}\Afootnote{H; \skt{तति}\ A}}}{\lemma{\skt{अपकर्षयेद्}}\Afootnote{H; \skt{मपकर्षयेद्}\ A}}\skt{ पञ्च पञ्च यावत्पञ्चेति~। }\edtext{\edtext{\skt{पञ्चेभ्यस्च}}{\lemma{\skt{पञ्चेभ्यस्}}\Afootnote{H; \skt{पञ्चभ्य}\ A}}}{\lemma{\skt{च}}\Afootnote{H; \skt{स्त्वेकैकं}\ A}}\skt{ }\edtext{\skt{यावदेक}}{\lemma{\skt{यावदेक}}\Afootnote{H; \skt{यावदेकम्}\ A}}\skt{ इति~। एवं }\edtext{\skt{भल्ला\-त\-क\-स\-ह\-स्र\-द्व\-य\-मु\-प\-युज्य}}{\lemma{\skt{भल्ला\appbrk{}त\appbrk{}क\appbrk{}स\appbrk{}ह\appbrk{}स्र\appbrk{}द्व\appbrk{}यम्}}\Afootnote{H; \skt{भल्ला\appbrk{}त\appbrk{}क\appbrk{}स\appbrk{}ह\appbrk{}स्राम्}\ A}}\skt{ }\edtext{\skt{सर्व्वाभिर्व्विमुक्तो}}{\lemma{\skt{सर्व्वाभिर्}}\Afootnote{H; \skt{सर्व\appbrk{}कु\appbrk{}ष्ठा\appbrk{}र्शो\appbrk{}भिर्}\ A}}\skt{ वलवान्नरोगः }\edtext{\skt{शतायुर्ब्भवति}}{\lemma{\skt{ब्भवति}}\Afootnote{H; \skt{भवति}\ A}}\skt{~।}\pend

\pstart \ledsidenote{\passref{4.6.18}}\edtext{\skt{भल्ला\-त\-क\-म\-ज्जेभ्यः}}{\lemma{\skt{भल्ला\appbrk{}त\appbrk{}क\appbrk{}म\appbrk{}ज्जेभ्यः}}\Afootnote{H; \skt{द्विव्रणीयोक्तेन विधानेन भल्ला\appbrk{}त\appbrk{}क\appbrk{}नि\appbrk{}श्च्यु\appbrk{}तितं स्नेहमादाय प्रातः प्रातः शुक्ति\appbrk{}मा\appbrk{}त्र\appbrk{}मु\appbrk{}प\appbrk{}यु\appbrk{}ञ्जीत, जीर्णे पूर्ववदाहारः फलप्रकर्षश्च~। भल्ला\appbrk{}त\appbrk{}क\appbrk{}म\appbrk{}ज्जभ्यो वा}\ A}}\skt{ स्नेह\-मा\-दा\-या\-प\-कृ\-ष्ट\-दोषः प्रति\-सं\-सृ\-ष्ट\-भक्तः~। }\edtext{\edtext{\skt{निवातमगारं}}{\lemma{\skt{निवातम्}}\Afootnote{H; \skt{निवात}\ A}}}{\lemma{\skt{अगारं}}\Afootnote{H; \skt{मागारं}\ A}}\skt{ प्रविश्य }\edtext{\skt{यथा\-व\-ल\-मु\-प\-यु\-ञ्जीत}}{\lemma{\skt{उपयुञ्जीत}}\Afootnote{H; \skt{प्रसृतिं प्रकुञ्चं वोपयुञ्जीत}\ A}}\skt{~। }\edtext{\skt{तस्मि न्जीर्ण्णे}}{\lemma{\skt{तस्मि न्}}\Afootnote{H; \skt{तस्मिञ्}\ A}}\skt{ }\edtext{\skt{क्षीर\-स\-र्प्पि\-रो\-दन}}{\lemma{\skt{क्षीर\appbrk{}स\appbrk{}र्प्पि\appbrk{}रो\appbrk{}दन}}\Afootnote{H; \skt{क्षीरं सर्पिरोदन}\ A}}\skt{ इत्याहारः~। एवं }\edtext{\edtext{\skt{मासमुपयुज्य}}{\lemma{\skt{मासम्}}\Afootnote{H; \skt{मास}\ A}}}{\lemma{\skt{उपयुज्य}}\Afootnote{H; \skt{मुपयुज्य}\ A}}\skt{ }\edtext{\edtext{\skt{मास\-त्र\-य\-मु\-प\-दि\-ष्टा\-हारो}}{\lemma{\skt{मासत्रयम्}}\Afootnote{H; \skt{मासत्रय}\ A}}}{\lemma{\skt{उपदिष्टाहारो}}\Afootnote{H; \skt{मादिष्टाहारो}\ A}}\skt{ }\edtext{\skt{रक्षे\-दा\-त्मा\-न\-न्ततः}}{\lemma{\skt{रक्षेदात्मानन्}}\Afootnote{H; \skt{रक्षेदात्मानं}\ A}}\skt{ सर्व्वो\-प\-ता\-पा\-न\-प\-हृत्य }\edtext{\edtext{\skt{वर्ण्ण\-वा\-न्ब\-ल\-वा\-ञ्छ्र\-वण}}{\lemma{\skt{बलवाञ्}}\Afootnote{H; \skt{बलवाञ्}\ A}}}{\lemma{\skt{छ्रवण}}\Afootnote{H; \skt{श्रवण}\ A}}\skt{ ग्रह\-ण\-धा\-र\-ण\-श\-क्ति\-स\-म्पन्नो }\edtext{\skt{वर्ष\-श\-ता\-यु\-र्ब्भ\-वति}}{\lemma{\skt{वर्ष\appbrk{}श\appbrk{}ता\appbrk{}यु\appbrk{}र्ब्भ\appbrk{}वति}}\Afootnote{H; \skt{वर्ष\appbrk{}श\appbrk{}ता\appbrk{}यु\appbrk{}र्भ\appbrk{}वति}\ A}}\skt{~। }\edtext{\skt{मासि}}{\lemma{\skt{मासि}}\Afootnote{H; \skt{मासे}\ A}}\skt{ }\edtext{\skt{मासि च}}{\lemma{\skt{मासि च}}\Afootnote{H; \skt{मासे च}\ A}}\skt{ प्रयोगे वर्षशतं }\edtext{\edtext{\skt{वर्ष\-श\-त\-मा\-यु\-षोऽ\-भि\-वृ\-द्धि\-र्ब्भ\-वति}}{\lemma{\skt{ऽभिवृद्धिर्}}\Afootnote{H; \skt{ऽभिवृद्धि}\ A}}}{\lemma{\skt{ब्भवति}}\Afootnote{H; \skt{र्भवति}\ A}}\skt{~। एव\-न्द\-श\-मा\-सा\-नु\-प\-युज्य वर्ष }\edtext{\skt{सह\-स्रा\-यु\-र्ब्भ\-वति}}{\lemma{\skt{ब्भवति}}\Afootnote{H; \skt{भवति}\ A}}\skt{~॥ }\edtext{\skt{भवन्ति}}{\lemma{\skt{भवन्ति}}\Afootnote{H; \emph{om.}\ A}}\skt{ }\edtext{\skt{चात्र}}{\lemma{\skt{चात्र}}\Afootnote{H; \emph{om.}\ A}}\skt{ }\edtext{\skt{श्लोकाः}}{\lemma{\skt{श्लोकाः}}\Afootnote{H; \emph{om.}\ A}}\skt{~॥}\pend

\pstart \ledsidenote{\passref{4.6.19}}\edtext{\skt{अर्शांसि}}{\lemma{\skt{अर्शांसि}}\Afootnote{H; \skt{भवन्ति चात्र यथा सर्वाणि कुष्ठानि हतः स्वदिरबीजकौ~। तथैवार्शांसि}\ A}}\skt{ सर्व्वाणि }\edtext{\skt{सदा}}{\lemma{\skt{सदा}}\Afootnote{H; \emph{om.}\ A}}\skt{ , }\edtext{\skt{व्र्क्षकारुष्करौ}}{\lemma{\skt{व्र्क्षकारुष्करौ}}\Afootnote{H; \skt{वृक्षकारुष्करौ}\ A}}\skt{ }\edtext{\skt{घ्नतः}}{\lemma{\skt{घ्नतः}}\Afootnote{H; \skt{हतः}\ A}}\skt{~। }\edtext{\skt{यथा}}{\lemma{\skt{यथा}}\Afootnote{H; \emph{om.}\ A}}\skt{ }\edtext{\skt{सर्व्वाणि}}{\lemma{\skt{सर्व्वाणि}}\Afootnote{H; \emph{om.}\ A}}\skt{ }\edtext{\skt{कुष्ठानि}}{\lemma{\skt{कुष्ठानि}}\Afootnote{H; \emph{om.}\ A}}\skt{ }\edtext{\skt{गाय\-त्र्य\-श\-न\-पा\-दयोः}}{\lemma{\skt{गाय\appbrk{}त्र्य\appbrk{}श\appbrk{}न\appbrk{}पा\appbrk{}दयोः}}\Afootnote{H; \emph{om.}\ A}}\skt{~॥}\pend

\pstart \ledsidenote{\passref{4.6.20}}\edtext{\skt{क्षाराग्निना}}{\lemma{\skt{क्षाराग्निना}}\Afootnote{H; \skt{हरिद्रायाः प्रयोगेण प्रमेहा इव षोडश~। क्षाराग्नी ना}\ A}}\skt{ }\edtext{\skt{निवर्त्तन्ते}}{\lemma{\skt{निवर्त्तन्ते}}\Afootnote{H; \skt{तिवर्तन्ते}\ A}}\skt{ }\edtext{\skt{दृश्यान्यर्शांसि}}{\lemma{\skt{दृश्यान्यर्शांसि}}\Afootnote{H; \skt{तथा दृश्या}\ A}}\skt{ }\edtext{\skt{सर्व्वदा}}{\lemma{\skt{सर्व्वदा}}\Afootnote{H; \emph{om.}\ A}}\skt{~। }\edtext{\skt{रजन्यो}}{\lemma{\skt{रजन्यो}}\Afootnote{H; \emph{om.}\ A}}\skt{ }\edtext{\skt{नाभिवर्त्तन्ते}}{\lemma{\skt{नाभिवर्त्तन्ते}}\Afootnote{H; \emph{om.}\ A}}\skt{ }\edtext{\skt{प्रमेहा}}{\lemma{\skt{प्रमेहा}}\Afootnote{H; \emph{om.}\ A}}\skt{ }\edtext{\skt{इव}}{\lemma{\skt{इव}}\Afootnote{H; \emph{om.}\ A}}\skt{ }\edtext{\skt{षोडश}}{\lemma{\skt{षोडश}}\Afootnote{H; \skt{गुदोद्भवाः}\ A}}\skt{~॥}\pend

\pstart \ledsidenote{\passref{4.6.21}}\skt{घृता\-नि\-दी\-प\-नी\-यानि }\edtext{\skt{लेहायकृतयः}}{\lemma{\skt{लेहायकृतयः}}\Afootnote{H; \skt{लेहायस्कृतयः}\ A}}\skt{ सुराः~। }\edtext{\skt{असवाश्च}}{\lemma{\skt{असवाश्च}}\Afootnote{H; \skt{आसवाश्च}\ A}}\skt{ प्रयोक्तव्या }\edtext{\skt{वीक्ष्य\-दो\-ष\-स\-मु\-छ्रेयं}}{\lemma{\skt{वीक्ष्य\appbrk{}दो\appbrk{}ष\appbrk{}स\appbrk{}मु\appbrk{}छ्रेयं}}\Afootnote{H; \skt{वीक्ष्य दोष\appbrk{}स\appbrk{}मु\appbrk{}च्छ्रि\appbrk{}तम्}\ A}}\skt{~॥}\pend

\pstart \ledsidenote{\passref{4.6.22}}\edtext{\skt{वेगा\-व\-रो\-ध\-स्त्री\-पृ\-ष्ठ\-या\-नान्य्}}{\lemma{\skt{वेगा\appbrk{}व\appbrk{}रो\appbrk{}ध\appbrk{}स्त्री\appbrk{}पृ\appbrk{}ष्ठ\appbrk{}या\appbrk{}नान्य्}}\Afootnote{H; \skt{वेगा\appbrk{}व\appbrk{}रो\appbrk{}ध\appbrk{}स्त्री\appbrk{}पृ\appbrk{}ष्ठ\appbrk{}या\appbrk{}नान्}\ A}}\skt{ ७ }\edtext{\skt{उत्कु}}{\lemma{\skt{उत्कु}}\Afootnote{H; \skt{युत्कु}\ A}}\skt{ }\edtext{\skt{दुकासनं}}{\lemma{\skt{दुकासनं}}\Afootnote{H\textsuperscript{pc}; \skt{टुकासनम्}\ A H\textsuperscript{ac}}}\skt{~। }\edtext{\edtext{\skt{यथा\-स्व\-न्दो\-ष\-ल\-ञ्च\-न्न\-म\-र्शसी}}{\lemma{\skt{यथा\appbrk{}स्व\appbrk{}न्दो\appbrk{}ष\appbrk{}ल\appbrk{}ञ्च\appbrk{}न्नम्}}\Afootnote{H; \skt{यथास्वं दोषलं चान्न}\ A}}}{\lemma{\skt{अर्शसी}}\Afootnote{H; \skt{मर्शःसु}\ A}}\skt{ }\edtext{\skt{परि\-व\-र्ज्ज\-ये\-दिति}}{\lemma{\skt{परि\appbrk{}व\appbrk{}र्ज्ज\appbrk{}ये\appbrk{}दिति}}\Afootnote{H; \skt{परिवर्जयेत्}\ A}}\skt{~॥}\pend

\pstart \ledsidenote{\passref{4.6.trailer}}\edtext{\skt{चिकित्साषट्}}{\lemma{\skt{चिकित्साषट्}}\Afootnote{H; \skt{इति सुश्रु\appbrk{}त\appbrk{}सं\appbrk{}हि\appbrk{}तायां चिकि\appbrk{}त्सा\appbrk{}स्था\appbrk{}नेऽ\appbrk{}र्श\appbrk{}श्चि\appbrk{}कि\appbrk{}त्सितं नाम षष्ठोऽध्यायः}\ A}}\skt{~॥}\pend

\endnumbering
\beginnumbering
\eledsection*{Adhyāya 7}
\pstart \ledsidenote{\passref{4.7.1}}\skt{॥ }\edtext{\skt{अथा\-तोऽ\-श्म\-री\-णा\-ञ्चि\-कि\-त्सि\-त\-म्व्या\-ख्या\-स्यामः}}{\lemma{\skt{ऽश्मरीणाञ्}}\Afootnote{H; \skt{ऽश्मरी}\ A}}\skt{~॥}\pend

\pstart\versepar \ledsidenote{\passref{4.7.3}}\skt{अश्मरी दारुणो }\edtext{\skt{व्याधिरन्तकः}}{\lemma{\skt{व्याधिरन्तकः}}\Afootnote{H; \skt{व्याधिरन्तक}\ A}}\skt{ प्रतिमो मतः~। }\pend
\pstart\versepar \skt{}\edtext{\skt{भेषजैस्तरुणः}}{\lemma{\skt{भेषजैस्}}\Afootnote{H; \skt{औषधैस्}\ A}}\skt{ साध्यः प्रवृ\-द्ध\-श्च्छे\-द\-म\-र्हति~॥}\pend

\pstart \ledsidenote{\passref{4.7.4}}\skt{॥ तस्य }\edtext{\skt{सर्व्वेषु}}{\lemma{\skt{सर्व्वेषु}}\Afootnote{H\textsuperscript{pc}; \skt{पूर्वे}\ A; \skt{सर्व्वे}\ H\textsuperscript{ac}}}\skt{ }\edtext{\edtext{\skt{रोगेष्विति}}{\lemma{\skt{रोगेष्व्}}\Afootnote{H\textsuperscript{pc}; \skt{षु}\ A H\textsuperscript{ac}}}}{\lemma{\skt{इति}}\Afootnote{H\textsuperscript{pc}; \emph{om.}\ A H\textsuperscript{ac}}}\skt{ }\edtext{\skt{पाठः~। ३रूपेषु}}{\lemma{\skt{पाठः~। ३रूपेषु}}\Afootnote{H\textsuperscript{pc}; \skt{रूपेषु}\ A H\textsuperscript{ac}}}\skt{ स्नेहादि क्रम }\edtext{\skt{इष्यते}}{\lemma{\skt{इष्यते}}\Afootnote{H; \skt{इष्यते~। तेनास्यापचयं यान्ति व्याधे\appbrk{}र्मू\appbrk{}ला\appbrk{}न्य\appbrk{}शे\appbrk{}षतः}\ A}}\skt{~।}\pend

\pstart\versepar \ledsidenote{\passref{4.7.5}}\edtext{\skt{पाषाणभेदो}}{\lemma{\skt{पाषाणभेदो}}\Afootnote{H; \skt{पाषाणमेदो}\ A}}\skt{ वसुको }\edtext{\edtext{\skt{वसि\-रोऽ\-श्म\-न्त\-क\-स्तथा}}{\lemma{\skt{वसिरो}}\Afootnote{H; \skt{वशिरा}\ A}}}{\lemma{\skt{ऽश्मन्तकस्}}\Afootnote{H; \skt{श्मन्तकौ}\ A}}\skt{~॥ }\pend
\pstart\versepar \skt{शतावरी श्वदंष्ट्रा च वृहती कण्टकारिका~।}\pend

\pstart\versepar \ledsidenote{\passref{4.7.6}}\edtext{\skt{कपो\-त\-व\-ङ्का\-र्त्त\-गलः}}{\lemma{\skt{कपो\appbrk{}त\appbrk{}व\appbrk{}ङ्का\appbrk{}र्त्त\appbrk{}गलः}}\Afootnote{H; \skt{कपो\appbrk{}त\appbrk{}व\appbrk{}ङ्काऽऽ\appbrk{}आ\appbrk{}र्त\appbrk{}गलः}\ A}}\skt{ }\edtext{\edtext{\skt{काञ्च\-नो\-शी\-र\-क\-स्तथा}}{\lemma{\skt{काञ्च\appbrk{}नो\appbrk{}शी\appbrk{}र\appbrk{}कस्}}\Afootnote{H; \skt{कच्चकोशीरकु}\ A}}}{\lemma{\skt{तथा}}\Afootnote{H; \skt{ब्जकाः}\ A}}\skt{~॥ }\pend
\pstart\versepar \skt{वृक्षादनी }\edtext{\skt{सभल्लूको}}{\lemma{\skt{सभल्लूको}}\Afootnote{H; \skt{भल्लुकश्च}\ A}}\skt{ वरुणः शाकजं फलं~।}\pend

\pstart\versepar \ledsidenote{\passref{4.7.7}}\edtext{\skt{यव\-को\-ल\-कु\-ल\-त्थानि}}{\lemma{\skt{यव\appbrk{}को\appbrk{}ल\appbrk{}कु\appbrk{}ल\appbrk{}त्थानि}}\Afootnote{H; \skt{यवाः कुलत्थाः कोलानि}\ A}}\skt{ कतकस्य फलानि च~॥ }\pend
\pstart\versepar \skt{}\edtext{\skt{ऊष\-का\-दि\-प्र\-ती\-वाय}}{\lemma{\skt{ऊष\appbrk{}का\appbrk{}दि\appbrk{}प्र\appbrk{}ती\appbrk{}वाय}}\Afootnote{H; \skt{ऊष\appbrk{}का\appbrk{}दि\appbrk{}प्र\appbrk{}ती\appbrk{}वाप}\ A}}\skt{ }\edtext{\skt{एषां}}{\lemma{\skt{एषां}}\Afootnote{H; \skt{मेषां}\ A}}\skt{ }\edtext{\skt{क्वाथे}}{\lemma{\skt{क्वाथे}}\Afootnote{H; \skt{क्वाथैर्}\ A}}\skt{ घृतङ्कृतं~।}\pend

\pstart\versepar \ledsidenote{\passref{4.7.8}}\skt{भिनत्ति }\edtext{\edtext{\skt{वात\-स\-म्भू\-ता\-म\-श्मरीं}}{\lemma{\skt{वातसम्भूताम्}}\Afootnote{H; \skt{वातसंभूता}\ A}}}{\lemma{\skt{अश्मरीं}}\Afootnote{H; \skt{मश्मरीं}\ A}}\skt{ क्षिप्रमेव तु~॥ }\pend
\pstart\versepar \skt{}\edtext{\skt{क्षारान्यवागूः}}{\lemma{\skt{यवागूः}}\Afootnote{H; \skt{यवागू}\ A}}\skt{ }\edtext{\skt{पेयानि}}{\lemma{\skt{पेयानि}}\Afootnote{H; \skt{र्यूषांश्च}\ A}}\skt{ }\edtext{\skt{कषायानि}}{\lemma{\skt{कषायानि}}\Afootnote{H; \skt{कषायाणि}\ A}}\skt{ पयांसि च~॥}\pend

\pstart\versepar \ledsidenote{\passref{4.7.9}}\skt{भोजनानि च कुर्व्वीत वर्ग्गेऽ\-स्मि\-न्वा\-त\-ना\-शने~॥ }\pend
\pstart\versepar \skt{कुशः काशः }\edtext{\skt{शरः}}{\lemma{\skt{शरः}}\Afootnote{H; \skt{सरो}\ A}}\skt{ }\edtext{\skt{कुच्छूरित्कटो}}{\lemma{\skt{कुच्छूरित्कटो}}\Afootnote{H; \skt{गुन्द्रा इत्कटो}\ A}}\skt{ }\edtext{\skt{मोरटोऽश्मभित्}}{\lemma{\skt{ऽश्मभित्}}\Afootnote{H; \skt{ऽश्मभित्}\ A}}\skt{~।}\pend

\pstart\versepar \ledsidenote{\passref{4.7.10}}\edtext{\skt{दर्ब्भा}}{\lemma{\skt{दर्ब्भा}}\Afootnote{H; \skt{वरी}\ A}}\skt{ विदारी वाराही }\edtext{\skt{शालिमूलं}}{\lemma{\skt{शालिमूलं}}\Afootnote{H; \skt{शालिमूल}\ A}}\skt{ त्रिकण्टकं~॥ }\pend
\pstart\versepar \skt{भल्लूकः }\edtext{\skt{पाटली}}{\lemma{\skt{पाटली}}\Afootnote{H; \skt{पाटला}\ A}}\skt{ }\edtext{\skt{पाठ}}{\lemma{\skt{पाठ}}\Afootnote{H; \skt{पाठा}\ A}}\skt{ पत्तूरोऽथ }\edtext{\skt{कुरूटिकाः}}{\lemma{\skt{कुरूटिकाः}}\Afootnote{H; \skt{कुरुष्टिका}\ A}}\skt{~।}\pend

\pstart\versepar \ledsidenote{\passref{4.7.11}}\edtext{\skt{पुनर्न्नवो}}{\lemma{\skt{पुनर्न्नवो}}\Afootnote{H; \skt{पुनर्नवा}\ A}}\skt{ शिरीषश्च }\edtext{\edtext{\skt{क्वथितास्तेषु}}{\lemma{\skt{क्वथितास्}}\Afootnote{H; \skt{क्वथिता}\ A}}}{\lemma{\skt{तेषु}}\Afootnote{H; \skt{स्तेषु}\ A}}\skt{ साधितं~॥ }\pend
\pstart\versepar \skt{घृतं }\edtext{\skt{शिलाजै}}{\lemma{\skt{शिलाजै}}\Afootnote{H; \skt{शिलाज}\ A}}\skt{ }\edtext{\edtext{\skt{मधु\-कै\-र्व्व\-द\-रे\-न्दी\-व\-रस्य}}{\lemma{\skt{मधुकैर्}}\Afootnote{H; \skt{मधुक}\ A}}}{\lemma{\skt{व्वद\appbrk{}रे\appbrk{}न्दी\appbrk{}व\appbrk{}रस्य}}\Afootnote{H; \skt{बीजै\appbrk{}रि\appbrk{}न्दी\appbrk{}व\appbrk{}रस्य}\ A}}\skt{ च~।}\pend

\pstart\versepar \ledsidenote{\passref{4.7.12}}\edtext{\skt{त्रपु\-षे\-र्व्वा\-रु\-का\-दीनां}}{\lemma{\skt{त्रपुषेर्}}\Afootnote{H; \skt{त्रपुसैर्}\ A}}\skt{ }\edtext{\skt{वीजैश्चावाधितं}}{\lemma{\skt{चावाधितं}}\Afootnote{H; \skt{चावापितं}\ A}}\skt{ शुभं~॥ }\pend
\pstart\versepar \skt{भिनत्ति }\edtext{\edtext{\skt{पित्त\-सं\-भू\-ता\-म\-श्मरीं}}{\lemma{\skt{पित्तसंभूताम्}}\Afootnote{H; \skt{पित्तसंभूता}\ A}}}{\lemma{\skt{अश्मरीं}}\Afootnote{H; \skt{मश्मरीं}\ A}}\skt{ क्षिप्रमेव तु~।}\pend

\pstart\versepar \ledsidenote{\passref{4.7.13}}\edtext{\skt{क्षारान्यवागूः}}{\lemma{\skt{यवागूः}}\Afootnote{H; \skt{यवागू}\ A}}\skt{ }\edtext{\skt{पेयानि}}{\lemma{\skt{पेयानि}}\Afootnote{H; \skt{र्यूषांश्च}\ A}}\skt{ }\edtext{\skt{कषायानि}}{\lemma{\skt{कषायानि}}\Afootnote{H; \skt{कषायाणि}\ A}}\skt{ पयांसि च~॥ }\pend
\pstart\versepar \skt{भोजनानि च कुर्व्वीत }\edtext{\skt{वर्ग्गे\-स्मि\-न्पि\-त्त\-ना\-शने}}{\lemma{\skt{वर्ग्गेस्मिन्}}\Afootnote{H; \skt{वर्गेऽस्मिन्}\ A}}\skt{~।}\pend

\pstart\versepar \ledsidenote{\passref{4.7.14}}\edtext{\skt{गणे}}{\lemma{\skt{गणे}}\Afootnote{H; \skt{गणो}\ A}}\skt{ }\edtext{\skt{वरुणकादौ तु}}{\lemma{\skt{वरुणकादौ तु}}\Afootnote{H; \skt{वरुणकादिस्तु}\ A}}\skt{ }\edtext{\skt{गुग्गु\-ल्वे\-ला\-ह\-रे\-णुभिः}}{\lemma{\skt{गुग्गु\appbrk{}ल्वे\appbrk{}ला\appbrk{}ह\appbrk{}रे\appbrk{}णुभिः}}\Afootnote{H; \skt{गुग्गु\appbrk{}ल्वे\appbrk{}ला\appbrk{}ह\appbrk{}रे\appbrk{}णवः}\ A}}\skt{~॥ }\pend
\pstart\versepar \skt{}\edtext{\skt{कुष्ठं}}{\lemma{\skt{कुष्ठं}}\Afootnote{H; \skt{कुष्ठ}\ A}}\skt{ }\edtext{\skt{भद्रादिमरिचं}}{\lemma{\skt{भद्रादिमरिचं}}\Afootnote{H; \skt{भद्रादिमरिच}\ A}}\skt{ }\edtext{\skt{चित्त्रकैः}}{\lemma{\skt{चित्त्रकैः}}\Afootnote{H; \skt{चिक्त्रकैः}\ A}}\skt{ ससुराह्वयैः~।}\pend

\pstart\versepar \ledsidenote{\passref{4.7.15}}\skt{एतैः सिद्ध\-म\-जा\-स\-र्प्पि\-रू\-ष\-कादि गणेन च~॥ }\pend
\pstart\versepar \skt{भिनत्ति }\edtext{\edtext{\skt{कफ\-सं\-भू\-ता\-म\-श्मरीं}}{\lemma{\skt{कफसंभूताम्}}\Afootnote{H; \skt{कफसंभूता}\ A}}}{\lemma{\skt{अश्मरीं}}\Afootnote{H; \skt{मश्मरीं}\ A}}\skt{ क्षिप्रमेव तु~।}\pend

\pstart\versepar \ledsidenote{\passref{4.7.16}}\edtext{\skt{क्षारान्यवागूः}}{\lemma{\skt{यवागूः}}\Afootnote{H; \skt{यवागू}\ A}}\skt{ }\edtext{\skt{पेयानि}}{\lemma{\skt{पेयानि}}\Afootnote{H; \skt{र्यूषांश्च}\ A}}\skt{ }\edtext{\skt{कषायानि}}{\lemma{\skt{कषायानि}}\Afootnote{H; \skt{कषायाणि}\ A}}\skt{ पयान्सि च~॥ }\pend
\pstart\versepar \skt{भोजनानि च कुर्व्वीत वर्ग्गेऽ\-स्मि\-न्क\-फ\-ना\-शने~।}\pend

\pstart\versepar \ledsidenote{\passref{4.7.17}}\edtext{\skt{पैचु\-का\-को\-ल्ल\-क\-तकं}}{\lemma{\skt{पैचु\appbrk{}का\appbrk{}को\appbrk{}ल्ल\appbrk{}क\appbrk{}तकं}}\Afootnote{H; \skt{पिचु\appbrk{}का\appbrk{}ङ्को\appbrk{}ल\appbrk{}क\appbrk{}तक}\ A}}\skt{ शाकेन्दीवरजैः फलैः~॥ }\pend
\pstart\versepar \skt{}\edtext{\edtext{\skt{पीतमुष्णाम्वु}}{\lemma{\skt{पीतम्}}\Afootnote{H; \skt{चू}\ A}}}{\lemma{\skt{उष्णाम्वु}}\Afootnote{H; \skt{र्णितैः}\ A}}\skt{ सगुडं }\edtext{\edtext{\skt{शर्क्क\-रा\-म्या\-त\-य\-त्यधः}}{\lemma{\skt{शर्क्क\appbrk{}रा\appbrk{}म्या\appbrk{}त\appbrk{}यत्य्}}\Afootnote{H; \skt{तोयं शर्करा}\ A}}}{\lemma{\skt{अधः}}\Afootnote{H; \skt{शमनं पिबेत्}\ A}}\skt{~।}\pend

\pstart\versepar \ledsidenote{\passref{4.7.18}}\edtext{\skt{क्रोञ्चो\-ष्ट्र\-रा\-स\-भा\-स्थीनि}}{\lemma{\skt{क्रोञ्चो\appbrk{}ष्ट्र\appbrk{}रा\appbrk{}स\appbrk{}भा\appbrk{}स्थीनि}}\Afootnote{H; \skt{क्रौञ्चो\appbrk{}ष्ट्र\appbrk{}रा\appbrk{}स\appbrk{}भा\appbrk{}स्थीनि}\ A}}\skt{ श्वदंष्ट्रा }\edtext{\skt{तालपत्त्रिका}}{\lemma{\skt{तालपत्त्रिका}}\Afootnote{H; \skt{तालमूलिका}\ A}}\skt{~॥ }\pend
\pstart\versepar \skt{अजमोदा }\edtext{\skt{कदम्वस्य}}{\lemma{\skt{कदम्वस्य}}\Afootnote{H; \skt{कदम्बरस्य}\ A}}\skt{ मूल\-न्ना\-ग\-र\-मेव च~।}\pend

\pstart\versepar \ledsidenote{\passref{4.7.19}}\skt{पीतानि शर्क्करां भिन्द्युः सुर\-यो\-ष्णो\-द\-केन वा~॥ }\pend
\pstart\versepar \skt{}\edtext{\skt{वृत्तकण्टक}}{\lemma{\skt{वृत्तकण्टक}}\Afootnote{H; \skt{त्रिकण्टक}\ A}}\skt{ }\edtext{\skt{वीजानि}}{\lemma{\skt{वीजानि}}\Afootnote{H; \skt{स्य बीजानां}\ A}}\skt{ चूर्ण्ण\-म्मा\-क्षि\-क\-सं\-युतं~।}\pend

\pstart\versepar \ledsidenote{\passref{4.7.20}}\skt{अविक्षीरेण सप्ताहं }\edtext{\edtext{\skt{पेय\-म\-श्म\-रि\-भे\-दनं}}{\lemma{\skt{पेयम्}}\Afootnote{H; \emph{om.}\ A}}}{\lemma{\skt{अश्मरिभेदनं}}\Afootnote{H; \skt{अश्मरीभेदनं}\ A}}\skt{~॥ }\pend
\pstart\versepar \skt{}\edtext{\skt{द्रव्याणाञ्च}}{\lemma{\skt{द्रव्याणाञ्च}}\Afootnote{H; \skt{पिबेत्~। द्रव्याणां तु}\ A}}\skt{ घृतोक्तानां }\edtext{\skt{क्षारा\-वि\-ण्मू\-त्र\-गा\-लितः}}{\lemma{\skt{क्षारा\appbrk{}वि\appbrk{}ण्मू\appbrk{}त्र\appbrk{}गा\appbrk{}लितः}}\Afootnote{H; \skt{क्षारोऽ\appbrk{}वि\appbrk{}मू\appbrk{}त्र\appbrk{}गा\appbrk{}लितः}\ A}}\skt{~।}\pend

\pstart\versepar \ledsidenote{\passref{4.7.21}}\skt{ग्राम्य\-स\-त्त्व\-श\-कृ\-त्क्षारैः }\edtext{\skt{संयुक्तं}}{\lemma{\skt{संयुक्तं}}\Afootnote{H; \skt{संयुक्तः}\ A}}\skt{ }\edtext{\skt{क्षारसाधितं}}{\lemma{\skt{क्षारसाधितं}}\Afootnote{H; \skt{साधितः}\ A}}\skt{~॥ }\pend
\pstart\versepar \skt{}\edtext{\skt{तत्रो\-ष\-का\-दि\-रा\-वापः}}{\lemma{\skt{तत्रो\appbrk{}ष\appbrk{}का\appbrk{}दि\appbrk{}रा\appbrk{}वापः}}\Afootnote{H; \skt{शनैः~। तत्रो\appbrk{}ष\appbrk{}का\appbrk{}दि\appbrk{}रा\appbrk{}वापः}\ A}}\skt{ }\edtext{\edtext{\skt{कार्य\-स्त्रि\-क\-टु\-का\-युतः}}{\lemma{\skt{कार्यस्}}\Afootnote{H; \skt{कार्य}\ A}}}{\lemma{\skt{त्रिकटुकायुतः}}\Afootnote{H; \skt{स्त्रिकटुकान्वितः}\ A}}\skt{~।}\pend

\pstart\versepar \ledsidenote{\passref{4.7.22}}\edtext{\skt{क्षार}}{\lemma{\skt{क्षार}}\Afootnote{H; \skt{एष क्षारो}\ A}}\skt{ }\edtext{\skt{एषोश्मरीं}}{\lemma{\skt{एषोश्मरीं}}\Afootnote{H; \skt{ऽश्मरीं}\ A}}\skt{ गुल्मं }\edtext{\skt{शर्क्करांश्च}}{\lemma{\skt{शर्क्करांश्च}}\Afootnote{H; \skt{शर्करां च}\ A}}\skt{ }\edtext{\skt{भिनत्त्यपि}}{\lemma{\skt{भिनत्त्यपि}}\Afootnote{H; \skt{भिनत्त्यपि}\ A}}\skt{~॥ }\pend
\pstart\versepar \skt{}\edtext{\skt{तिला\-पा\-मा\-र्ग्ग\-क\-द\-ली\-प\-ला\-स\-य\-व\-व\-ल्वजाः}}{\lemma{\skt{तिला\appbrk{}पा\appbrk{}मा\appbrk{}र्ग्ग\appbrk{}क\appbrk{}द\appbrk{}ली\appbrk{}प\appbrk{}ला\appbrk{}स\appbrk{}य\appbrk{}व\appbrk{}व\appbrk{}ल्वजाः}}\Afootnote{H; \skt{तिला\appbrk{}पा\appbrk{}मा\appbrk{}र्ग\appbrk{}क\appbrk{}द\appbrk{}ली\appbrk{}प\appbrk{}ला\appbrk{}श\appbrk{}य\appbrk{}व\appbrk{}क\appbrk{}ल्कजः}\ A}}\skt{~।}\pend

\pstart\versepar \ledsidenote{\passref{4.7.23}}\skt{क्षारः पेयोऽविमूत्रेण }\edtext{\skt{शर्क्क\-रा\-श्म\-रि\-भे\-दनः}}{\lemma{\skt{शर्क्क\appbrk{}रा\appbrk{}श्म\appbrk{}रि\appbrk{}भे\appbrk{}दनः}}\Afootnote{H; \skt{शर्करानाशनः परः}\ A}}\skt{~॥ }\pend
\pstart\versepar \skt{}\edtext{\skt{पाट\-ली\-क\-र\-वी\-राणां}}{\lemma{\skt{पाट\appbrk{}ली\appbrk{}क\appbrk{}र\appbrk{}वी\appbrk{}राणां}}\Afootnote{H; \skt{पाट\appbrk{}ला\appbrk{}क\appbrk{}र\appbrk{}वी\appbrk{}राणां}\ A}}\skt{ }\edtext{\edtext{\skt{क्षारमेवं}}{\lemma{\skt{क्षारम्}}\Afootnote{H; \skt{क्षार}\ A}}}{\lemma{\skt{एवं}}\Afootnote{H; \skt{मेवं}\ A}}\skt{ समाचरेत्~।}\pend

\pstart\versepar \ledsidenote{\passref{4.7.24}}\skt{श्वदं\-ष्ट्रा\-य\-ष्टि\-का\-व्रह्मी }\edtext{\skt{कल्क\-म्वा\-क्ष\-स\-म\-म्पि\-वेत्}}{\lemma{\skt{वाक्षसमम्}}\Afootnote{H; \skt{वाऽक्षसमं}\ A}}\skt{~॥ }\pend
\pstart\versepar \skt{}\edtext{\skt{सहैलकाक्षै}}{\lemma{\skt{सहैलकाक्षै}}\Afootnote{H; \skt{सहैडकाख्यौ}\ A}}\skt{ पेयौ वा }\edtext{\skt{तथासौ}}{\lemma{\skt{तथासौ}}\Afootnote{H; \skt{शोभाञ्ज}\ A}}\skt{ }\edtext{\skt{भञ्जनार्क्कजौ}}{\lemma{\skt{भञ्जनार्क्कजौ}}\Afootnote{H; \skt{नकमार्कवौ}\ A}}\skt{~।}\pend

\pstart\versepar \ledsidenote{\passref{4.7.25}}\edtext{\skt{कपो\-त\-व\-ङ्क\-मू\-लम्वा}}{\lemma{\skt{कपो\appbrk{}त\appbrk{}व\appbrk{}ङ्क\appbrk{}मू\appbrk{}लम्वा}}\Afootnote{H; \skt{कपो\appbrk{}त\appbrk{}व\appbrk{}ङ्का\appbrk{}मूलं वा}\ A}}\skt{ }\edtext{\skt{पिवेदम्लं}}{\lemma{\skt{पिवेदम्लं}}\Afootnote{H; \skt{पिबेदम्लैः}\ A}}\skt{ सुरादिभिः~॥ }\pend
\pstart\versepar \skt{तत्सिद्धम्वा पिवेत्क्षीरं वेद\-ना\-भि\-रु\-प\-द्रुतः~।}\pend

\pstart\versepar \ledsidenote{\passref{4.7.26}}\skt{हरीतक्यादि }\edtext{\edtext{\skt{सिद्ध\-म्वा\-सि\-द्ध\-म्वापि}}{\lemma{\skt{वासिद्धम्}}\Afootnote{H; \skt{वा वर्षाभूसिद्धम्}\ A}}}{\lemma{\skt{वापि}}\Afootnote{H; \emph{om.}\ A}}\skt{ }\edtext{\skt{पुनर्न्नवे}}{\lemma{\skt{पुनर्न्नवे}}\Afootnote{H; \skt{एव वा}\ A}}\skt{~॥ }\pend
\pstart\versepar \skt{}\edtext{\skt{सर्व्व\-थै\-वो\-प\-युज्याः}}{\lemma{\skt{सर्व्व\appbrk{}थै\appbrk{}वो\appbrk{}प\appbrk{}युज्याः}}\Afootnote{H; \skt{सर्व\appbrk{}थै\appbrk{}वो\appbrk{}प\appbrk{}योज्यः}\ A}}\skt{ }\edtext{\edtext{\skt{स्युर्ग्ग्णो}}{\lemma{\skt{स्युर्}}\Afootnote{H; \skt{स्याद्}\ A}}}{\lemma{\skt{ग्ग्णो}}\Afootnote{H; \skt{गणो}\ A}}\skt{ वीरतरादिकः~।}\pend

\pstart\versepar \ledsidenote{\passref{4.7.27}}\skt{घृतैः क्षारैः कषायैश्च क्षीरैः सोत्तरबस्तिभिः~॥ }\pend
\pstart\versepar \skt{यदि }\edtext{\skt{नैवोपशाम्येत}}{\lemma{\skt{नैवोपशाम्येत}}\Afootnote{H; \skt{नोपशमं गच्छेच्}\ A}}\skt{ च्छेदस्तत्रोत्तरो विधिः~।}\pend

\pstart\versepar \ledsidenote{\passref{4.7.28}}\edtext{\skt{निश्चितस्यास्य}}{\lemma{\skt{निश्चितस्यास्य}}\Afootnote{H; \skt{कुशलस्यापि}\ A}}\skt{ वैद्यस्य यतः }\edtext{\edtext{\skt{सिद्धिरिह}}{\lemma{\skt{सिद्धिर्}}\Afootnote{H; \skt{सिद्धि}\ A}}}{\lemma{\skt{इह}}\Afootnote{H; \skt{रिह}\ A}}\skt{ ध्रुवा~॥ }\pend
\pstart\versepar \skt{उपक्रमो जघन्योयमतः सम्परिकीर्त्तितः~।}\pend

\pstart\versepar \ledsidenote{\passref{4.7.29}}\edtext{\skt{अक्रियायां}}{\lemma{\skt{अक्रियायां}}\Afootnote{H; \skt{अक्रियायं}\ A}}\skt{ ध्रुवो मृत्युः क्रियायां संशयो भवेत्~॥ }\pend
\pstart\versepar \skt{तस्मादापृच्छ्य कर्त्तव्यमीश्वरं साधुकारिणा~॥}\pend

\pstart \ledsidenote{\passref{4.7.30}}\edtext{\edtext{\edtext{\edtext{\skt{अथा\-रो\-ग्या\-गा\-र्थि\-न\-मु\-प\-स्नि\-ग्ध\-म\-प\-कृ\-ष्ट\-दो\-ष\-मी\-ष\-त्क\-र्षि\-त\-म\-भ्य\-क्त\-म\-भु\-क्त\-वन्तं}}{\lemma{\skt{अथा\appbrk{}रो\appbrk{}ग्या\appbrk{}गा\appbrk{}र्थि\appbrk{}नम्}}\Afootnote{H\textsuperscript{pc}; \skt{अथ रोगान्वितम्}\ A; \skt{अथा\appbrk{}रो\appbrk{}ग्या\appbrk{}र्थि\appbrk{}नम्}\ H\textsuperscript{ac}}}}{\lemma{\skt{कर्षितम्}}\Afootnote{H; \skt{कर्शितम्}\ A}}}{\lemma{\skt{अभ्यक्तम्}}\Afootnote{H; \skt{अभ्य}\ A}}}{\lemma{\skt{अभुक्तवन्तं}}\Afootnote{H; \skt{क्त}\ A}}\skt{ }\edtext{\skt{स्निन्नशरीरं}}{\lemma{\skt{स्निन्नशरीरं}}\Afootnote{H; \skt{स्विन्नशरीरं}\ A}}\skt{ }\edtext{\skt{कृत\-म\-ङ्ग\-ल\-स्व\-स्ति\-वा\-च\-न\-म\-ग्रो\-प\-ह\-र\-णी\-यो\-क्तेन}}{\lemma{\skt{कृत\appbrk{}म\appbrk{}ङ्ग\appbrk{}ल\appbrk{}स्व\appbrk{}स्ति\appbrk{}वा\appbrk{}च\appbrk{}नम्}}\Afootnote{H; \skt{भुक्तवन्तं कृत\appbrk{}ब\appbrk{}लि\appbrk{}म\appbrk{}ङ्ग\appbrk{}ल\appbrk{}स्व\appbrk{}स्ति\appbrk{}वा\appbrk{}च\appbrk{}नम्}\ A}}\skt{ }\edtext{\skt{विधिनानेन}}{\lemma{\skt{विधिनानेन}}\Afootnote{H; \skt{विधानेनोप}\ A}}\skt{ }\edtext{\skt{सस\-म्भृ\-त\-स\-म्भा\-र\-मा\-श्वास्य}}{\lemma{\skt{सस\appbrk{}म्भृ\appbrk{}त\appbrk{}स\appbrk{}म्भा\appbrk{}रम्}}\Afootnote{H; \skt{कल्पि\appbrk{}त\appbrk{}स\appbrk{}म्भा\appbrk{}रम्}\ A}}\skt{ ततो वल\-व\-न्त\-म\-वि\-क्ल\-व\-मा\-जा\-नु\-समे फलके }\edtext{\skt{प्रागुपवेश्य}}{\lemma{\skt{प्रागुपवेश्य}}\Afootnote{H; \skt{प्रागुपविष्टान्य}\ A}}\skt{ }\edtext{\skt{पुरु }\edtext{\skt{षन्तस्योत्सङ्गे}}{\lemma{\skt{तस्योत्सङ्गे}}\Afootnote{H; \skt{षस्योत्सङ्गे}\ A}}}{\lemma{\skt{पुरु षन्}}\Afootnote{H; \skt{पुरु}\ A}}\skt{ }\edtext{\skt{निष\-ण्ण\-पू\-र्व\-का\-य\-मु\-त्तानं}}{\lemma{\skt{निष\appbrk{}ण्ण\appbrk{}पू\appbrk{}र्व\appbrk{}का\appbrk{}य\appbrk{}मु\appbrk{}त्तानं}}\Afootnote{H; \skt{निष\appbrk{}ण्ण\appbrk{}पू\appbrk{}र्व\appbrk{}का\appbrk{}य\appbrk{}मु\appbrk{}त्ता\appbrk{}नम्}\ A}}\skt{ उन्नतकटीकं }\edtext{\skt{वस्त्र\-चु\-म्भ\-ल\-को\-प\-विष्टं}}{\lemma{\skt{वस्त्र\appbrk{}चु\appbrk{}म्भ\appbrk{}ल\appbrk{}को\appbrk{}प\appbrk{}विष्टं}}\Afootnote{H; \skt{वस्त्रा\appbrk{}धा\appbrk{}र\appbrk{}को\appbrk{}प\appbrk{}विष्टं}\ A}}\skt{ }\edtext{\skt{सङ्कु\-च\-ति\-जा\-नु\-कु\-र्प्प\-र\-मि\-त\-रेण}}{\lemma{\skt{सङ्कु\appbrk{}च\appbrk{}ति\appbrk{}जा\appbrk{}नु\appbrk{}कु\appbrk{}र्प्प\appbrk{}र\appbrk{}मि\appbrk{}त\appbrk{}रेण}}\Afootnote{H; \skt{सङ्कु\appbrk{}चि\appbrk{}त\appbrk{}जा\appbrk{}नु\appbrk{}कू\appbrk{}र्प\appbrk{}र\appbrk{}मि\appbrk{}त\appbrk{}रेण}\ A}}\skt{ सहाववद्धं~। }\edtext{\skt{सूत्रकेन}}{\lemma{\skt{सूत्रकेन}}\Afootnote{H; \skt{सूत्रेण}\ A}}\skt{ }\edtext{\skt{शाटकेन}}{\lemma{\skt{शाटकेन}}\Afootnote{H; \skt{शाटकैर्}\ A}}\skt{ वाततः }\edtext{\skt{स्वभ्य\-क्त\-ना\-भि\-प्र\-दे\-श\-पा\-र्श्व\-स्या\-व\-मृद्य}}{\lemma{\skt{स्वभ्य\appbrk{}क्त\appbrk{}ना\appbrk{}भि\appbrk{}प्र\appbrk{}दे\appbrk{}श\appbrk{}पा\appbrk{}र्श्व\appbrk{}स्या\appbrk{}व\appbrk{}मृद्य}}\Afootnote{H; \skt{स्वभ्य\appbrk{}क्त\appbrk{}ना\appbrk{}भि\appbrk{}प्र\appbrk{}दे\appbrk{}शस्य वामपार्श्वं विमृद्य}\ A}}\skt{ मुष्टिना }\edtext{\skt{पीडयेदधो}}{\lemma{\skt{पीडयेदधो}}\Afootnote{H; \skt{वपीडयेदधो}\ A}}\skt{ नाभेर्यावदश्म र्यधः }\edtext{\skt{प्रतिपन्नेति}}{\lemma{\skt{प्रतिपन्नेति}}\Afootnote{H; \skt{प्रपन्नेति}\ A}}\skt{~। ततः }\edtext{\skt{स्नेहाभ्यक्त}}{\lemma{\skt{स्नेहाभ्यक्त}}\Afootnote{H; \skt{स्नेहाभ्यक्ते}\ A}}\skt{ }\edtext{\skt{क्लिप्तनखो}}{\lemma{\skt{क्लिप्तनखो}}\Afootnote{H; \skt{क्लृप्तनखे}\ A}}\skt{ वाम\-ह\-स्त\-प्र\-दे\-शिनी }\edtext{\skt{मध्यमाङ्गुल्यौ}}{\lemma{\skt{मध्यमाङ्गुल्यौ}}\Afootnote{H; \skt{मध्यमे अङ्गुल्यौ पायौ}\ A}}\skt{ प्रणि\-धा\-या\-नु\-से\-व\-नी\-मा\-साद्य प्रयत्नवलाभ्यां }\edtext{\edtext{\skt{मेढ्र\-गु\-द\-यो\-र\-न्त\-र\-मा\-न\-यित्वा}}{\lemma{\skt{मेढ्रगुदयोर्}}\Afootnote{H; \skt{पायुमेढ्}\ A}}}{\lemma{\skt{अन्त\appbrk{}र\appbrk{}मा\appbrk{}न\appbrk{}यित्वा}}\Afootnote{H; \skt{रान्तरमानीय}\ A}}\skt{ }\edtext{\skt{निर्व्व\-ली\-क\-म\-ना\-य\-त\-न\-म\-वि\-ष\-मञ्च}}{\lemma{\skt{निर्व्व\appbrk{}ली\appbrk{}क\appbrk{}म\appbrk{}ना\appbrk{}य\appbrk{}त\appbrk{}न\appbrk{}म\appbrk{}वि\appbrk{}ष\appbrk{}मञ्च}}\Afootnote{H; \skt{निर्व्य\appbrk{}ली\appbrk{}क\appbrk{}म\appbrk{}ना\appbrk{}य\appbrk{}त\appbrk{}म\appbrk{}वि\appbrk{}षमं च}\ A}}\skt{ }\edtext{\skt{बस्ति}}{\lemma{\skt{बस्ति}}\Afootnote{H; \skt{वस्तिं}\ A}}\skt{ सन्निवेश्य }\edtext{\skt{भृश\-मु\-त्पी\-ड\-ये\-द\-ङ्गु\-लीभ्यां}}{\lemma{\skt{भृश\appbrk{}मु\appbrk{}त्पी\appbrk{}ड\appbrk{}ये\appbrk{}द\appbrk{}ङ्गु\appbrk{}लीभ्यां}}\Afootnote{H; \skt{भृश\appbrk{}मु\appbrk{}त्पी\appbrk{}ड\appbrk{}ये\appbrk{}द\appbrk{}ङ्गु\appbrk{}लिभ्यां}\ A}}\skt{ }\edtext{\skt{यदा}}{\lemma{\skt{यदा}}\Afootnote{H; \skt{यथा}\ A}}\skt{ ग्रन्थिरिवोन्नतं शल्यम्भवति~॥}\pend

\pstart\versepar \ledsidenote{\passref{4.7.31}}\skt{स चेद्गृहीतशल्ये तु विवृत्ताक्षो विचेतनः~। }\pend
\pstart\versepar \skt{}\edtext{\skt{हतवल्लम्वते}}{\lemma{\skt{हतवल्लम्वते}}\Afootnote{H; \skt{हतवल्लम्ब}\ A}}\skt{ }\edtext{\skt{शीर्षा}}{\lemma{\skt{शीर्षा}}\Afootnote{H; \skt{शीर्षश्च}\ A}}\skt{ निर्व्विकारो मृतोपमः~॥}\pend

\pstart\versepar \ledsidenote{\passref{4.7.32}}\skt{न तस्य }\edtext{\skt{निर्ह\-रे\-च्छ\-ल्य\-न्नि\-र्ह\-तेपि}}{\lemma{\skt{निर्हतेपि}}\Afootnote{H; \skt{निर्हरेत्तु}\ A}}\skt{ म्रियेत सः~। }\pend
\pstart\versepar \skt{विना }\edtext{\edtext{\skt{त्वेतैर्व्विकारैस्तु}}{\lemma{\skt{त्वेतैर्}}\Afootnote{H; \skt{त्वेते}\ A}}}{\lemma{\skt{व्विकारैस्तु}}\Afootnote{H; \skt{षु रूपेषु}\ A}}\skt{ }\edtext{\edtext{\skt{समा\-ह\-र्त्तु\-मु\-पा\-च\-रेत्}}{\lemma{\skt{समाहर्त्तुम्}}\Afootnote{H; \skt{निर्हर्तुं}\ A}}}{\lemma{\skt{उपाचरेत्}}\Afootnote{H; \skt{प्रयतेत वै}\ A}}\skt{~॥}\pend

\pstart \ledsidenote{\passref{4.7.33}}\skt{ततः }\edtext{\skt{सेवन्याः}}{\lemma{\skt{सेवन्याः}}\Afootnote{H; \emph{om.}\ A}}\skt{ सव्ये पार्श्वे }\edtext{\skt{सेवनी}}{\lemma{\skt{सेवनी}}\Afootnote{H; \skt{सेवनीं}\ A}}\skt{ }\edtext{\skt{यव\-मा\-त्रे\-णा\-व\-मु\-च्या\-व\-चा\-र\-ये\-च्छस्त्रं}}{\lemma{\skt{यव\appbrk{}मा\appbrk{}त्रे\appbrk{}णा\appbrk{}व\appbrk{}मु\appbrk{}च्या\appbrk{}व\appbrk{}चा\appbrk{}र\appbrk{}येच्}}\Afootnote{H; \skt{यवमात्रेण मुक्त्वाऽ\appbrk{}व\appbrk{}चा\appbrk{}र\appbrk{}येच्}\ A}}\skt{~। }\edtext{\skt{अश्म\-री\-प्र\-मा\-णा\-द्द\-क्षि\-णतो}}{\lemma{\skt{अश्म\appbrk{}री\appbrk{}प्र\appbrk{}मा\appbrk{}णाद्}}\Afootnote{H; \skt{अश्मरीप्रमाणं}\ A}}\skt{ वा क्रिया सौक\-र्य\-हे\-तो\-रि\-त्येके~। }\edtext{\skt{यथा न}}{\lemma{\skt{यथा न}}\Afootnote{H; \skt{यथा सा न}\ A}}\skt{ }\edtext{\skt{भिद्यते न}}{\lemma{\skt{भिद्यते न}}\Afootnote{H; \skt{भिद्यते}\ A}}\skt{ }\edtext{\skt{वचूर्ण्ण्यते वा}}{\lemma{\skt{वचूर्ण्ण्यते वा}}\Afootnote{H; \skt{चूर्ण्यते वा}\ A}}\skt{ तथा }\edtext{\skt{प्रयतेत}}{\lemma{\skt{प्रयतेत}}\Afootnote{H; \skt{प्रयते}\ A}}\skt{~॥ }\edtext{\edtext{\skt{तच्चू\-र्ण्ण\-म\-ल्प\-म\-प्ये\-व\-स्थितं}}{\lemma{\skt{तच्चूर्ण्णम्}}\Afootnote{\skt{त, चूर्णम्}\ A; \skt{तच्चूर्ण्णम्}\ H}}\skt{ हि}}{\lemma{\skt{एवस्थितं हि}}\Afootnote{H; \skt{अवस्थितं हि}\ A}}\skt{ पुनः }\edtext{\skt{परिवृद्धमेति}}{\lemma{\skt{परिवृद्धम्}}\Afootnote{H; \skt{परिवृद्धिम्}\ A}}\skt{~। }\edtext{\edtext{\skt{तस्मा\-त्म\-म\-ग्रा\-म\-ग्र\-व\-क्रे\-ना\-द\-दीत}}{\lemma{\skt{ममग्राम्}}\Afootnote{H; \skt{समस्ताम्}\ A}}}{\lemma{\skt{अग्र\appbrk{}व\appbrk{}क्रे\appbrk{}ना\appbrk{}द\appbrk{}दीत}}\Afootnote{H; \skt{अग्र\appbrk{}व\appbrk{}क्रे\appbrk{}णा\appbrk{}द\appbrk{}दीत}\ A}}\skt{~। स्त्रीणान्तु }\edtext{\skt{बस्तिः}}{\lemma{\skt{बस्तिः}}\Afootnote{H; \skt{बस्ति}\ A}}\skt{ पार्श्वगतो }\edtext{\skt{गर्ब्भा\-स\-य\-स\-न्नि\-च\-य\-स्तस्मा}}{\lemma{\skt{गर्ब्भा\appbrk{}स\appbrk{}य\appbrk{}स\appbrk{}न्नि\appbrk{}च\appbrk{}यस्}}\Afootnote{H; \skt{गर्भाशयः सन्निकृष्टः}\ A}}\skt{ }\edtext{\skt{त्तासा\-म्मु\-त्स\-ङ्ग\-व\-च्छस्त्रं}}{\lemma{\skt{मुत्सङ्गवच्}}\Afootnote{H; \skt{उत्सङ्गवच्छ्}\ A}}\skt{ पातयेत्~॥ अतोन्यथा खल्वा\-सा\-म्मू\-त्र\-स्रावी व्रणो }\edtext{\skt{भवति}}{\lemma{\skt{भवति}}\Afootnote{H; \skt{भवेत्}\ A}}\skt{~॥ पुरुषस्य वा }\edtext{\skt{मूत्र\-प्र\-से\-क\-क्ष\-ण\-ना\-द\-श्मरी}}{\lemma{\skt{मूत्र\appbrk{}प्र\appbrk{}से\appbrk{}क\appbrk{}क्ष\appbrk{}ण\appbrk{}ना\appbrk{}द\appbrk{}श्मरी}}\Afootnote{H; \skt{मूत्र\appbrk{}प्र\appbrk{}से\appbrk{}क\appbrk{}क्ष\appbrk{}ण\appbrk{}ना\appbrk{}न्मू\appbrk{}त्र\appbrk{}क्ष\appbrk{}र\appbrk{}णम्, अश्मरी}\ A}}\skt{ व्रणादृते भिन्न\-ब\-स्ति\-रे\-कधा }\edtext{\skt{द्विधा वा न}}{\lemma{\skt{द्विधा वा न}}\Afootnote{H; \skt{ऽपि न}\ A}}\skt{ भवति~। द्विधा }\edtext{\skt{भिन्न\-ब\-स्ति\-का\-श्म\-रीको पि न}}{\lemma{\skt{भिन्न\appbrk{}ब\appbrk{}स्ति\appbrk{}का\appbrk{}श्म\appbrk{}रीको पि न}}\Afootnote{H; \skt{भिन्न\appbrk{}ब\appbrk{}स्ति\appbrk{}रा\appbrk{}श्म\appbrk{}रिको न}\ A}}\skt{ सिध्यति~। }\edtext{\skt{अश्म\-री\-नि\-मि\-त्तश्च}}{\lemma{\skt{अश्म\appbrk{}री\appbrk{}नि\appbrk{}मि\appbrk{}त्तश्च}}\Afootnote{H; \skt{अश्मरी}\ A}}\skt{ }\edtext{\skt{व्रणमेकधा}}{\lemma{\skt{व्रणम्}}\Afootnote{H; \skt{व्रणनिमित्तम्}\ A}}\skt{ }\edtext{\skt{भिन्नबस्ति}}{\lemma{\skt{भिन्नबस्ति}}\Afootnote{H; \skt{भिन्नबस्तिर्}\ A}}\skt{ }\edtext{\skt{विधिना}}{\lemma{\skt{विधिना}}\Afootnote{H; \emph{om.}\ A}}\skt{ जीवति }\edtext{\skt{शस्त्रक्रियाभ्यां}}{\lemma{\skt{शस्त्रक्रियाभ्यां}}\Afootnote{H; \skt{क्रियाभ्या}\ A}}\skt{ }\edtext{\edtext{\skt{शस्त्र\-वि\-हि\-ता\-च्छे\-दा\-भि\-ष्य\-न्द\-प\-रि\-वृ\-द्ध\-त्वाच्च}}{\lemma{\skt{शस्त्रविहिताच्}}\Afootnote{H; \skt{साच्छा\appbrk{}स्त्र\appbrk{}वि\appbrk{}हि\appbrk{}तच्छ्}\ A}}}{\lemma{\skt{छेदा\appbrk{}भि\appbrk{}ष्य\appbrk{}न्द\appbrk{}प\appbrk{}रि\appbrk{}वृ\appbrk{}द्ध\appbrk{}त्वाच्च}}\Afootnote{H; \skt{एदा\appbrk{}न्निः\appbrk{}स्य\appbrk{}न्द\appbrk{}प\appbrk{}रि\appbrk{}वृ\appbrk{}द्ध\appbrk{}त्वाच्च}\ A}}\skt{ शल्यस्य }\edtext{\skt{जीवेदिति}}{\lemma{\skt{जीवेदिति}}\Afootnote{H; \skt{इति}\ A}}\skt{~। }\edtext{\skt{उद्धृ\-त\-श\-ल्य\-मु\-ष्णो\-द\-क\-द्रोण्यां}}{\lemma{\skt{उष्णोदकद्रोण्यां}}\Afootnote{H; \skt{तूष्णो\appbrk{}द\appbrk{}क\appbrk{}द्रो\appbrk{}ण्याम्}\ A}}\skt{ }\edtext{\skt{चावगाह्य}}{\lemma{\skt{चावगाह्य}}\Afootnote{H; \skt{अवगाह्य}\ A}}\skt{ स्वेदयेत्तथाहि }\edtext{\edtext{\skt{बस्तिरसृजा}}{\lemma{\skt{बस्तिर्}}\Afootnote{\skt{हस्तिर्}\ A; \skt{वस्ति}\ H}}\skt{ न}}{\lemma{\skt{असृजा न}}\Afootnote{A; \skt{रसृजा न}\ H}}\skt{ }\edtext{\skt{पूर्यते}}{\lemma{\skt{पूर्यते}}\Afootnote{H; \skt{पूर्यते, पूर्णं वा क्षीर\appbrk{}वृ\appbrk{}क्ष\appbrk{}क\appbrk{}षायं पुष्पनेत्रेण विदध्यात्}\ A}}\skt{~।}\pend

\pstart\versepar \ledsidenote{\passref{4.7.34}}\edtext{\skt{पूर्ण्णे वा}}{\lemma{\skt{पूर्ण्णे वा}}\Afootnote{H; \skt{भवति चात्र}\ A}}\skt{ }\edtext{\skt{क्षीर\-वृ\-क्ष\-क\-षायं}}{\lemma{\skt{क्षीर\appbrk{}वृ\appbrk{}क्ष\appbrk{}क\appbrk{}षायं}}\Afootnote{H; \skt{क्षीर\appbrk{}वृ\appbrk{}क्ष\appbrk{}क\appbrk{}षाय}\ A}}\skt{ }\edtext{\skt{तु}}{\lemma{\skt{तु}}\Afootnote{H; \skt{स्तु}\ A}}\skt{ पुष्पनेत्रेण }\edtext{\skt{योजितं}}{\lemma{\skt{योजितं}}\Afootnote{H; \skt{योजितः}\ A}}\skt{~। }\pend
\pstart\versepar \skt{}\edtext{\skt{निर्ह\-रे\-द\-श्म\-री\-चूर्ण्णं}}{\lemma{\skt{निर्ह\appbrk{}रे\appbrk{}द\appbrk{}श्म\appbrk{}री\appbrk{}चूर्ण्णं}}\Afootnote{H; \skt{निर्हरेदश्मरीं तूर्णं}\ A}}\skt{ रक्त\-म्ब\-स्ति\-ग\-तञ्च यत्~॥}\pend

\pstart \ledsidenote{\passref{4.7.35}}\edtext{\skt{मूत्र\-वि\-शु\-द्ध्य\-र्थञ्च}}{\lemma{\skt{मूत्र\appbrk{}वि\appbrk{}शु\appbrk{}द्ध्य\appbrk{}र्थञ्च}}\Afootnote{H; \skt{मूत्र\appbrk{}मा\appbrk{}र्ग\appbrk{}वि\appbrk{}शो\appbrk{}ध\appbrk{}नार्थं चा}\ A}}\skt{ }\edtext{\skt{अस्मै}}{\lemma{\skt{अस्मै}}\Afootnote{H; \skt{स्मै}\ A}}\skt{ }\edtext{\edtext{\skt{गुड\-सौ\-हि\-त्य\-म्वि\-त\-च\-रेत्}}{\lemma{\skt{गुडसौहित्यम्}}\Afootnote{H; \skt{गुडसौहित्यं}\ A}}\skt{~। उद्धृत्य}}{\lemma{\skt{वितचरेत्~। उद्धृत्य}}\Afootnote{H\textsuperscript{pc}; \skt{वितरेतुद्धृत्य}\ A H\textsuperscript{ac}}}\skt{ चैनं }\edtext{\skt{मधुघृताभ्यक्तं}}{\lemma{\skt{मधुघृताभ्यक्तं}}\Afootnote{H; \skt{मधुघृताभ्यक्त}\ A}}\skt{ व्रणं }\edtext{\skt{मूत्र\-वि\-शो\-ध\-नेन}}{\lemma{\skt{मूत्र\appbrk{}वि\appbrk{}शो\appbrk{}ध\appbrk{}नेन}}\Afootnote{H; \skt{मूत्रविशोधन}\ A}}\skt{ }\edtext{\skt{द्रव्येण}}{\lemma{\skt{द्रव्येण}}\Afootnote{H; \skt{द्रव्य}\ A}}\skt{ सिद्धामुष्णां }\edtext{\skt{यवागूं}}{\lemma{\skt{यवागूं}}\Afootnote{H; \emph{om.}\ A}}\skt{ सघृतां }\edtext{\skt{पाय\-ये\-दु\-भ\-यतः}}{\lemma{\skt{पाय\appbrk{}ये\appbrk{}दु\appbrk{}भ\appbrk{}यतः}}\Afootnote{H; \skt{यवागूं पाययेतोभय}\ A}}\skt{ कालं त्रिरात्रं~। त्रिरा\-त्रा\-दू\-र्द्ध्व\-ङ्गु\-ड\-प्र\-गा\-ढेन पयसा }\edtext{\skt{मृद्वो\-द\-न\-म\-ल्प\-म्भो\-ज\-येत्}}{\lemma{\skt{भोजयेत्}}\Afootnote{H; \skt{भोज\appbrk{}ये\appbrk{}द्द\appbrk{}श\appbrk{}रात्रं (मूत्}\ A}}\skt{~। }\edtext{\skt{दशरात्रादूर्द्ध्वं}}{\lemma{\skt{दशरात्रादूर्द्ध्वं}}\Afootnote{H; \skt{रासृग्विशुद्ध्यर्थं व्रणक्लेदनार्थं च), दशरात्रादूर्ध्वं}\ A}}\skt{ }\edtext{\skt{फला\-म्लै\-र्जा\-ङ्ग\-ल\-र\-सै\-र्द्द\-श\-रा\-त्र\-ञ्चै\-न\-म\-प्र\-मत्तः}}{\lemma{\skt{द्दशरात्रञ्}}\Afootnote{H; \skt{उपाचरेत्, ततो दशरात्रं}\ A}}\skt{ स्वेदयेत्~॥ }\edtext{\skt{क्षीर\-वृ\-क्ष\-क\-षा\-येण}}{\lemma{\skt{क्षीर\appbrk{}वृ\appbrk{}क्ष\appbrk{}क\appbrk{}षा\appbrk{}येण}}\Afootnote{H; \skt{स्नेहेन द्रवस्वेदेन वा, क्षीर\appbrk{}वृ\appbrk{}क्ष\appbrk{}क\appbrk{}षा\appbrk{}येण}\ A}}\skt{ चास्य व्रणं प्रक्षालयेत्~। }\edtext{\skt{लोध्र\-म\-धु\-क\-म\-ञ्जि\-ष्ठा\-प्र\-पु\-ण्ड\-री\-क\-क\-ल्कै\-र्व्व्र\-ण\-म्प्र\-ति\-ग्रा\-ह\-येत्}}{\lemma{\skt{लोध्र\appbrk{}म\appbrk{}धु\appbrk{}क\appbrk{}म\appbrk{}ञ्जि\appbrk{}ष्ठा\appbrk{}प्र\appbrk{}पु\appbrk{}ण्ड\appbrk{}री\appbrk{}क\appbrk{}क\appbrk{}ल्कैर्}}\Afootnote{H; \skt{रोध्र\appbrk{}म\appbrk{}धु\appbrk{}क\appbrk{}म\appbrk{}ञ्जि\appbrk{}ष्ठा\appbrk{}प्र\appbrk{}पौ\appbrk{}ण्ड\appbrk{}री\appbrk{}क\appbrk{}क\appbrk{}ल्कैर्}\ A}}\skt{~। }\edtext{\skt{एतेष्वेव च}}{\lemma{\skt{एव च}}\Afootnote{H; \skt{एव}\ A}}\skt{ हरिद्रायुतेषु }\edtext{\skt{घृत\-म्वि\-प\-क्व\-म्व्र\-णा\-भ्य\-ञ्ज\-न\-मिति}}{\lemma{\skt{घृतम्}}\Afootnote{H; \skt{तैलं घृतं वा}\ A}}\skt{~॥ स्त्यान\-शो\-णि\-त\-ञ्चो\-त्त\-र\-ब\-स्ति\-भि\-रु\-पा\-च\-रेत्~। }\edtext{\skt{सप्तरात्राच च}}{\lemma{\skt{सप्तरात्राच च}}\Afootnote{H; \skt{सप्तरात्राच्च स्}\ A}}\skt{ }\edtext{\skt{मार्ग्ग\-म\-प्र\-ति\-प\-द्य\-माने}}{\lemma{\skt{मार्ग्गम्}}\Afootnote{H; \skt{वमार्गम्}\ A}}\skt{ मूत्रे }\edtext{\skt{यथो\-क्त\-वि\-धा\-नेन}}{\lemma{\skt{यथो\appbrk{}क्त\appbrk{}वि\appbrk{}धा\appbrk{}नेन}}\Afootnote{H; \skt{व्रणं यथोक्तेन विधिना द}\ A}}\skt{ }\edtext{\edtext{\skt{व्रणमग्निना}}{\lemma{\skt{व्रणम्}}\Afootnote{H; \skt{हेद}\ A}}}{\lemma{\skt{अग्निना}}\Afootnote{H; \skt{ग्निना}\ A}}\skt{ }\edtext{\skt{दहेत्}}{\lemma{\skt{दहेत्}}\Afootnote{H; \emph{om.}\ A}}\skt{~। }\edtext{\skt{स्वमा\-र्ग्ग\-म्प्र\-ति\-पन्ने}}{\lemma{\skt{स्वमार्ग्गम्}}\Afootnote{H; \skt{स्वमार्ग}\ A}}\skt{ चोत्तरवस्त्या }\edtext{\edtext{\skt{स्थाप\-ना\-नु\-वा\-स\-नै\-रु\-पा\-च\-रे\-त्म\-धु\-क\-क\-षा\-यै\-रिति}}{\lemma{\skt{उपाचरेत्}}\Afootnote{H; \skt{उपाचरेन्}\ A}}}{\lemma{\skt{मधु\appbrk{}क\appbrk{}क\appbrk{}षा\appbrk{}यैर्}}\Afootnote{H; \skt{मधु\appbrk{}र\appbrk{}क\appbrk{}षा\appbrk{}यैर्}\ A}}\skt{~। }\edtext{\skt{यदृच्छया}}{\lemma{\skt{यदृच्छया}}\Afootnote{H; \skt{यदृच्छया वा}\ A}}\skt{ }\edtext{\edtext{\skt{मूत्र\-मा\-र्ग्ग\-म\-प्र\-ति\-पन्नां}}{\lemma{\skt{मूत्रमार्ग्गम्}}\Afootnote{H; \skt{मूत्रमार्ग}\ A}}}{\lemma{\skt{अप्रतिपन्नां}}\Afootnote{H; \skt{प्रतिपन्नाम्}\ A}}\skt{ }\edtext{\skt{सक्ताश्मरीं}}{\lemma{\skt{सक्ताश्मरीं}}\Afootnote{H; \skt{अन्तरासक्तां शुक्राश्मरीं}\ A}}\skt{ }\edtext{\skt{वा}}{\lemma{\skt{वा}}\Afootnote{H; \skt{शर्करां वा}\ A}}\skt{ }\edtext{\skt{स्रोत\-सा\-प\-ह\-रेत्}}{\lemma{\skt{स्रोत\appbrk{}सा\appbrk{}प\appbrk{}ह\appbrk{}रेत्}}\Afootnote{H; \skt{स्रोत\appbrk{}साऽ\appbrk{}प\appbrk{}ह\appbrk{}रेत्, एवं}\ A}}\skt{~॥ }\edtext{\skt{विदार्यवानालं}}{\lemma{\skt{विदार्यवानालं}}\Afootnote{H; \skt{चाशक्ये विदार्य नाडीं}\ A}}\skt{ }\edtext{\skt{शस्त्रे\-णा\-ग्र\-च\-क्रेण}}{\lemma{\skt{शस्त्रे\appbrk{}णा\appbrk{}ग्र\appbrk{}च\appbrk{}क्रेण}}\Afootnote{H; \skt{शस्त्रेण बडिशेनोद्ध}\ A}}\skt{ }\edtext{\edtext{\skt{रूढ\-व्र\-ण\-श्चा\-ङ्ग\-ना\-श्व\-न\-ग\-ना\-ग\-र\-थ\-द्रु\-मा\-न्ना\-रो\-हे\-द्वर्षं}}{\lemma{\skt{रूढव्रणश्}}\Afootnote{H; \skt{रेत्~। रूढव्रणश्}\ A}}}{\lemma{\skt{नारोहेद्}}\Afootnote{H; \skt{नारोहेत}\ A}}\skt{~। नाप्सु }\edtext{\skt{प्लवेद्भुञ्जीत}}{\lemma{\skt{प्लवेद्}}\Afootnote{H; \skt{प्लवेत}\ A}}\skt{ वागुरु~।}\pend

\pstart \ledsidenote{\passref{4.7.36}}\edtext{\skt{तत्र}}{\lemma{\skt{तत्र}}\Afootnote{H; \skt{मूत्र\appbrk{}व\appbrk{}ह\appbrk{}शु\appbrk{}क्र\appbrk{}व\appbrk{}ह\appbrk{}मु\appbrk{}ष्क\appbrk{}स्रोतो मूत्र\appbrk{}प्र\appbrk{}से\appbrk{}क\appbrk{}से\appbrk{}व\appbrk{}नी\appbrk{}यो\appbrk{}नि\appbrk{}गु\appbrk{}द\appbrk{}ब\appbrk{}स्ती\appbrk{}नष्टौ परिहरेत्~। तत्र}\ A}}\skt{ }\edtext{\skt{मूत्र\-व\-ह\-च्छे\-दा\-त्म\-रणं}}{\lemma{\skt{मूत्र\appbrk{}व\appbrk{}ह\appbrk{}च्छे\appbrk{}दात्}}\Afootnote{H; \skt{मूत्र\appbrk{}व\appbrk{}ह\appbrk{}च्छे\appbrk{}दान्}\ A}}\skt{~। मूत्रपूर्ण्णवस्तेः }\edtext{\skt{शुक्र\-व\-ह\-च्छे\-दा\-त्म\-रणं}}{\lemma{\skt{शुक्र\appbrk{}व\appbrk{}ह\appbrk{}च्छे\appbrk{}दात्}}\Afootnote{H; \skt{शुक्र\appbrk{}व\appbrk{}ह\appbrk{}च्छे\appbrk{}दान्}\ A}}\skt{~। क्लैव्यम्वा }\edtext{\skt{फल\-स्रो\-तो\-प\-घा\-ता\-द्ध्व\-ज\-भङ्गः}}{\lemma{\skt{फल\appbrk{}स्रो\appbrk{}तो\appbrk{}प\appbrk{}घा\appbrk{}ताद्}}\Afootnote{H; \skt{मुष्क\appbrk{}स्रो\appbrk{}तü\appbrk{}प\appbrk{}घा\appbrk{}ताद्}\ A}}\skt{~। }\edtext{\edtext{\skt{मूत्र\-प्र\-से\-व\-क\-च्छे\-दा\-न्मू\-त्र\-श्र\-वणं}}{\lemma{\skt{मूत्र\appbrk{}प्र\appbrk{}से\appbrk{}व\appbrk{}क\appbrk{}च्छे\appbrk{}दान्}}\Afootnote{H; \skt{मूत्र\appbrk{}प्र\appbrk{}से\appbrk{}क\appbrk{}क्ष\appbrk{}ण\appbrk{}नान्}\ A}}}{\lemma{\skt{मूत्रश्रवणं}}\Afootnote{H; \skt{मूत्रप्रक्षरणं}\ A}}\skt{~। }\edtext{\skt{सेव\-नी\-च्छे\-दा\-द्रुजा}}{\lemma{\skt{सेव\appbrk{}नी\appbrk{}च्छे\appbrk{}दा\appbrk{}द्रुजा}}\Afootnote{H; \skt{सेव\appbrk{}नी\appbrk{}यो\appbrk{}नि\appbrk{}च्छे\appbrk{}दा\appbrk{}द्रुजः}\ A}}\skt{ }\edtext{\skt{प्रादुर्ब्भावः}}{\lemma{\skt{प्रादुर्ब्भावः}}\Afootnote{H; \skt{प्रादुर्भावः}\ A}}\skt{~। बस्ति\-गु\-द\-वि\-द्ध\-ल\-क्षणं प्रागुक्तमिति~॥}\pend

\pstart\versepar \ledsidenote{\passref{4.7.37}}\skt{भवति चात्र~॥ }\pend
\pstart\versepar \skt{}\edtext{\skt{मर्म्मा\-ण्ण्ये\-ता\-न्य\-सं\-वुध्य}}{\lemma{\skt{मर्म्मा\appbrk{}ण्ण्ये\appbrk{}ता\appbrk{}न्य\appbrk{}सं\appbrk{}वुध्य}}\Afootnote{H; \skt{मर्मा\appbrk{}ण्य\appbrk{}ष्टा\appbrk{}व\appbrk{}सं\appbrk{}बुध्य}\ A}}\skt{ }\edtext{\skt{स्रोतोजानि}}{\lemma{\skt{स्रोतोजानि}}\Afootnote{H; \skt{स्रोतजानि}\ A}}\skt{ शरीरिणां~। }\pend
\pstart\versepar \skt{}\edtext{\edtext{\skt{व्यापा\-द\-ये\-द्व\-हू\-न्मर्त्या}}{\lemma{\skt{व्यापादयेद्}}\Afootnote{H; \skt{व्यापादयेद्}\ A}}}{\lemma{\skt{मर्त्या}}\Afootnote{H; \skt{मर्त्यान्}\ A}}\skt{ }\edtext{\edtext{\skt{च्छस्त्र\-क\-र्म्मा\-प\-टु\-र्ब्भि\-षक्}}{\lemma{\skt{च्छस्त्र\appbrk{}क\appbrk{}र्म्मा\appbrk{}प\appbrk{}टुर्}}\Afootnote{H; \skt{शस्त्रकर्मापटु}\ A}}}{\lemma{\skt{ब्भिषक्}}\Afootnote{H; \skt{र्भिषक्}\ A}}\skt{~॥}\pend

\pstart\versepar \ledsidenote{\passref{4.7.38}}\skt{सेवनी शुक्रहरणी }\edtext{\edtext{\skt{स्रोत\-सी\-फ\-ल\-यो\-र्ग्गुदः}}{\lemma{\skt{स्रोत\appbrk{}सी\appbrk{}फ\appbrk{}ल\appbrk{}योर्}}\Afootnote{H; \skt{स्रोतसी फलयो}\ A}}}{\lemma{\skt{ग्गुदः}}\Afootnote{H; \skt{र्गुदम्}\ A}}\skt{~। }\pend
\pstart\versepar \skt{मूत्रसेकं मूत्रवहं }\edtext{\skt{तथा}}{\lemma{\skt{तथा}}\Afootnote{H; \skt{योनिर्बस्तिस्तथाऽ}\ A}}\skt{ }\edtext{\edtext{\edtext{\skt{बस्ति\-मि\-हा\-ष्ट\-म\-मिति}}{\lemma{\skt{बस्तिम्}}\Afootnote{H; \emph{om.}\ A}}}{\lemma{\skt{इहाष्टमम्}}\Afootnote{H; \skt{ष्टमः}\ A}}}{\lemma{\skt{इति}}\Afootnote{H; \emph{om.}\ A}}\skt{~॥}\pend

\pstart \ledsidenote{\passref{4.7.trailer}}\edtext{\skt{चिकित्सायां}}{\lemma{\skt{चिकित्सायां}}\Afootnote{H; \skt{इति सुश्रु\appbrk{}त\appbrk{}सं\appbrk{}हि\appbrk{}तायां चिकि\appbrk{}त्सा\appbrk{}स्था\appbrk{}नेऽ\appbrk{}श्म\appbrk{}री\appbrk{}चि\appbrk{}कि\appbrk{}त्सितं}\ A}}\skt{ }\edtext{\edtext{\skt{सप्तमोऽध्यायः}}{\lemma{\skt{सप्तमो}}\Afootnote{H; \skt{नाम सप्तमोऽ}\ A}}}{\lemma{\skt{ऽध्यायः}}\Afootnote{H; \skt{ध्यायः}\ A}}\skt{~॥}\pend

\endnumbering
\beginnumbering
\eledsection*{Adhyāya 8}
\pstart \ledsidenote{\passref{4.8.1}}\skt{अथातो भग\-न्द\-रा\-णा\-ञ्चि\-कि\-त्सि\-त\-म्व्या\-ख्या\-स्यामः~॥}\pend

\pstart \ledsidenote{\passref{4.8.2}}\edtext{\skt{पञ्च}}{\lemma{\skt{पञ्च}}\Afootnote{H; \skt{यथोवाच भ}\ A}}\skt{ }\edtext{\skt{भगन्दरा}}{\lemma{\skt{भगन्दरा}}\Afootnote{H; \skt{गवान्धन्वन्तरिः}\ A}}\skt{ }\edtext{\skt{व्याख्याताः}}{\lemma{\skt{व्याख्याताः}}\Afootnote{H; \emph{om.}\ A}}\skt{~।}\pend

\pstart \ledsidenote{\passref{4.8.3}}\edtext{\skt{तेष्वसाध्यः}}{\lemma{\skt{तेष्व्}}\Afootnote{H; \skt{पञ्च भगन्दरा व्याख्याताः, तेष्व्}\ A}}\skt{ शम्बूकावर्त्तः शल्यनिमित्तश्च~। शेषाः कृच्छ्रसाध्याः~।}\pend

\pstart \ledsidenote{\passref{4.8.4}}\skt{तत्र }\edtext{\edtext{\edtext{\edtext{\skt{भग\-न्द\-र\-पि\-ट\-को\-प\-द्रु\-त\-मा\-तु\-र\-म\-प\-त\-र्प्प\-णा\-द्येन}}{\lemma{\skt{भग\appbrk{}न्द\appbrk{}र\appbrk{}पि\appbrk{}ट\appbrk{}को\appbrk{}प\appbrk{}द्रु\appbrk{}तम्}}\Afootnote{H; \skt{भग\appbrk{}न्द\appbrk{}र\appbrk{}पि\appbrk{}ड\appbrk{}को\appbrk{}प\appbrk{}द्रुत}\ A}}}{\lemma{\skt{आतुरम्}}\Afootnote{H; \skt{मातुर}\ A}}}{\lemma{\skt{अपतर्प्पणाद्}}\Afootnote{H; \skt{मपतर्पणा}\ A}}}{\lemma{\skt{येन}}\Afootnote{H; \skt{दि}\ A}}\skt{ }\edtext{\skt{विरे\-च\-ना\-न्ते\-नै\-का\-द\-श\-वि\-धे\-नो\-प\-क्र\-मे\-णो\-प\-च\-रेत्~। अपक्क्वपिटकं}}{\lemma{\skt{विरे\appbrk{}च\appbrk{}ना\appbrk{}न्ते\appbrk{}नै\appbrk{}का\appbrk{}द\appbrk{}श\appbrk{}वि\appbrk{}धे\appbrk{}नो\appbrk{}प\appbrk{}क्र\appbrk{}मे\appbrk{}णो\appbrk{}प\appbrk{}च\appbrk{}रेत्~। अपक्क्वपिटकं}}\Afootnote{H; \skt{विरे\appbrk{}च\appbrk{}ना\appbrk{}न्ते\appbrk{}नै\appbrk{}का\appbrk{}द\appbrk{}श\appbrk{}वि\appbrk{}धे\appbrk{}नो\appbrk{}प\appbrk{}क्र\appbrk{}मे\appbrk{}णो\appbrk{}प\appbrk{}क्र\appbrk{}मे\appbrk{}ता\appbrk{}प\appbrk{}क्व\appbrk{}पि\appbrk{}डकं}\ A}}\skt{~। }\edtext{\skt{पक्क्वेषु}}{\lemma{\skt{पक्क्वेषु}}\Afootnote{H; \skt{पक्वेषु}\ A}}\skt{ }\edtext{\skt{चोप\-स्नि\-ग्ध\-म\-व\-गा\-ढ\-स्विन्नं}}{\lemma{\skt{चोप\appbrk{}स्नि\appbrk{}ग्ध\appbrk{}म\appbrk{}व\appbrk{}गा\appbrk{}ढ\appbrk{}स्विन्नं}}\Afootnote{H; \skt{चोप\appbrk{}स्नि\appbrk{}ग्ध\appbrk{}म\appbrk{}व\appbrk{}गा\appbrk{}ह\appbrk{}स्विन्नं}\ A}}\skt{ शय्यायां }\edtext{\skt{संवे\-श्या\-र्श\-स\-मिव}}{\lemma{\skt{संवे\appbrk{}श्या\appbrk{}र्श\appbrk{}स\appbrk{}मिव}}\Afootnote{H; \skt{सन्नि\appbrk{}वे\appbrk{}श्या\appbrk{}र्श\appbrk{}स\appbrk{}मिव}\ A}}\skt{ यन्त्रयित्वा }\edtext{\edtext{\skt{भग\-न्द\-र\-म\-भि\-स\-मीक्ष्य}}{\lemma{\skt{अभिस्}}\Afootnote{H; \emph{om.}\ A}}}{\lemma{\skt{अमीक्ष्य}}\Afootnote{H; \skt{समीक्ष्य}\ A}}\skt{ }\edtext{\skt{परा\-ची\-न\-म\-र्वा\-क्चीनं}}{\lemma{\skt{परा\appbrk{}ची\appbrk{}न\appbrk{}म\appbrk{}र्वा\appbrk{}क्चीनं}}\Afootnote{H; \skt{परा\appbrk{}ची\appbrk{}न\appbrk{}म\appbrk{}वा\appbrk{}चीनं}\ A}}\skt{ }\edtext{\skt{बहि\-र्म्मु\-ख\-म\-न्त\-र्मु\-खम्वा}}{\lemma{\skt{बहि\appbrk{}र्म्मु\appbrk{}ख\appbrk{}म\appbrk{}न्त\appbrk{}र्मु\appbrk{}खम्वा}}\Afootnote{H; \skt{वा}\ A}}\skt{ ततः }\edtext{\skt{प्रादा\-ये\-ष\-णी\-न्न\-म\-यित्वा}}{\lemma{\skt{प्रादा\appbrk{}ये\appbrk{}ष\appbrk{}णी\appbrk{}न्न\appbrk{}म\appbrk{}यित्वा}}\Afootnote{H; \skt{प्रणि\appbrk{}धा\appbrk{}यै\appbrk{}ष\appbrk{}णी\appbrk{}मु\appbrk{}न्नम्य}\ A}}\skt{ }\edtext{\skt{शाय\-स\-मु\-द्ध\-रे\-च्छ\-स्त्रेण}}{\lemma{\skt{शाय\appbrk{}स\appbrk{}मु\appbrk{}द्ध\appbrk{}रे\appbrk{}च्छ\appbrk{}स्त्रेण}}\Afootnote{H; \skt{साश\appbrk{}य\appbrk{}मु\appbrk{}द्ध\appbrk{}रे\appbrk{}च्छ\appbrk{}स्त्रेण}\ A}}\skt{~। }\edtext{\edtext{\skt{अन्त\-र्मु\-ख\-ञ्चै\-वम}}{\lemma{\skt{अन्तर्मुखञ्}}\Afootnote{H; \skt{अन्तर्मुखे}\ A}}}{\lemma{\skt{चैवम}}\Afootnote{H; \skt{चैवं स}\ A}}\skt{ }\edtext{\skt{शक्यं यन्त्रम्प्रणिधाय}}{\lemma{\skt{शक्यं यन्त्रम्}}\Afootnote{H; \skt{म्यग्यन्त्रं}\ A}}\skt{ }\edtext{\skt{प्रवाहतो}}{\lemma{\skt{प्रवाहतो}}\Afootnote{H; \skt{प्रवाहमाणस्य}\ A}}\skt{ }\edtext{\skt{भग\-न्द\-र\-मौ\-ख\-मा\-सा\-द्ये\-ष\-णी\-न्दत्वा}}{\lemma{\skt{भग\appbrk{}न्द\appbrk{}र\appbrk{}मौ\appbrk{}ख\appbrk{}मा\appbrk{}सा\appbrk{}द्ये\appbrk{}ष\appbrk{}णीन्}}\Afootnote{H; \skt{भग\appbrk{}न्द\appbrk{}र\appbrk{}मु\appbrk{}ख\appbrk{}मा\appbrk{}सा\appbrk{}द्यै\appbrk{}षणीं}\ A}}\skt{ शस्त्रम्पातयेत्~। आसाद्य }\edtext{\skt{चाग्निं}}{\lemma{\skt{चाग्निं}}\Afootnote{H; \skt{वाऽग्निं}\ A}}\skt{ क्षारं चेति }\edtext{\skt{सामान्यं}}{\lemma{\skt{सामान्यं}}\Afootnote{H; \skt{एतत्सामान्यं}\ A}}\skt{ सर्व्वेषु}\pend

\pstart\versepar \ledsidenote{\passref{4.8.5}}\skt{विशेषतस्तु }\edtext{\skt{नान्यतरे}}{\lemma{\skt{नान्यतरे}}\Afootnote{H; \skt{नाड्यन्तरे}\ A}}\skt{ }\edtext{\skt{व्रणं}}{\lemma{\skt{व्रणं}}\Afootnote{H; \skt{व्रणान्}\ A}}\skt{ कुर्याद्भिषक्तु शतपोनके~। }\pend
\pstart\versepar \skt{}\edtext{\edtext{\skt{तत\-स्ता\-सू\-प\-रू\-ढासु}}{\lemma{\skt{ततस्}}\Afootnote{H; \skt{तत}\ A}}}{\lemma{\skt{तासूपरूढासु}}\Afootnote{H; \skt{स्तेषूपरूढेषु}\ A}}\skt{ शेषा नाडी\-रु\-पा\-च\-रेत्~।}\pend

\pstart\versepar \ledsidenote{\passref{4.8.6}}\skt{गतयो }\edtext{\skt{न्योन्यसंबन्धा}}{\lemma{\skt{न्योन्यसंबन्धा}}\Afootnote{H; \skt{ऽन्योन्यसंबद्धा}\ A}}\skt{ }\edtext{\edtext{\skt{बह्व्या\-श्छि\-न्द्या\-न्य\-ने\-कधा}}{\lemma{\skt{बह्व्याश्}}\Afootnote{H; \skt{बाह्याश्}\ A}}}{\lemma{\skt{छिन्द्यान्यनेकधा}}\Afootnote{H; \skt{छेद्यास्त्वनेकधा}\ A}}\skt{~॥ }\pend
\pstart\versepar \skt{}\edtext{\edtext{\skt{नाडी\-र\-न\-भि\-स\-म्बद्धा}}{\lemma{\skt{नाडीरनभिस्}}\Afootnote{H; \skt{नाडीरनभि}\ A}}\skt{ }\edtext{\skt{यश्छिनत्येकधा}}{\lemma{\skt{छिनत्य्}}\Afootnote{H; \skt{छिनत्त्य्}\ A}}}{\lemma{\skt{अम्बद्धा यश्}}\Afootnote{H; \skt{संबद्धा यश्}\ A}}\skt{ भिषक्~।}\pend

\pstart\versepar \ledsidenote{\passref{4.8.7}}\skt{स कुर्याद्विवृतं }\edtext{\skt{जन्तो\-र्म्मा\-र्ग्ग\-ङ्गु\-द\-वि\-दा\-रणं}}{\lemma{\skt{म्मार्ग्गङ्}}\Afootnote{H; \skt{व्रणं}\ A}}\skt{~॥ }\pend
\pstart\versepar \skt{तस्य }\edtext{\skt{त्त्रं}}{\lemma{\skt{त्त्रं}}\Afootnote{H; \skt{तद्}\ A}}\skt{ विवृतं }\edtext{\skt{मार्ग्गम्विण्मूत्रे}}{\lemma{\skt{विण्मूत्रे}}\Afootnote{H; \skt{विण्मूत्रम्}\ A}}\skt{ }\edtext{\skt{तु}}{\lemma{\skt{तु}}\Afootnote{H; \skt{अनुगच्छति}\ A}}\skt{ }\edtext{\skt{प्रपद्यतः}}{\lemma{\skt{प्रपद्यतः}}\Afootnote{H; \emph{om.}\ A}}\skt{~।}\pend

\pstart\versepar \ledsidenote{\passref{4.8.8}}\edtext{\skt{साटोपं}}{\lemma{\skt{साटोपं}}\Afootnote{H; \skt{आटोपं}\ A}}\skt{ गुदशूलञ्च करोति पवनो भृशं~॥ }\pend
\pstart\versepar \skt{}\edtext{\skt{तत्रापि}}{\lemma{\skt{तत्रापि}}\Afootnote{H; \skt{तत्राधि}\ A}}\skt{ }\edtext{\skt{शततम्त्रोऽपि}}{\lemma{\skt{शततम्त्रो}}\Afootnote{H; \skt{गततन्त्रो}\ A}}\skt{ }\edtext{\skt{भिष\-क्गु\-ह्ये\-द\-सं\-शयं}}{\lemma{\skt{भिषक्}}\Afootnote{H; \skt{भिषङ्}\ A}}\skt{~।}\pend

\pstart\versepar \ledsidenote{\passref{4.8.9}}\skt{तस्मान्न विवृतः कार्यो व्रणस्तु शतपोनके~॥ }\pend
\pstart\versepar \skt{व्याधौ }\edtext{\skt{तनु}}{\lemma{\skt{तनु}}\Afootnote{H; \skt{तत्र}\ A}}\skt{ बहुच्छिद्रे भिषजा }\edtext{\skt{साधु}}{\lemma{\skt{साधु}}\Afootnote{H; \skt{वै वि}\ A}}\skt{ जानता~।}\pend

\pstart\versepar \ledsidenote{\passref{4.8.10}}\edtext{\skt{अर्द्ध\-ला\-ङ्ग\-ल\-ल\-क\-श्छेदः}}{\lemma{\skt{अर्द्ध\appbrk{}ला\appbrk{}ङ्ग\appbrk{}ल\appbrk{}ल\appbrk{}कश्}}\Afootnote{H; \skt{अर्ध\appbrk{}ला\appbrk{}ङ्ग\appbrk{}ल\appbrk{}कश्}\ A}}\skt{ कार्यो लाङ्गलकोऽपि वा~॥ }\pend
\pstart\versepar \skt{सर्व्वतोभद्रको }\edtext{\skt{वापि}}{\lemma{\skt{वापि}}\Afootnote{H; \skt{वाऽपि}\ A}}\skt{ कार्यो गोतीर्थकोऽपि वा~।}\pend

\pstart\versepar \ledsidenote{\passref{4.8.11}}\edtext{\skt{सर्व्वाना}}{\lemma{\skt{सर्व्वाना}}\Afootnote{H; \skt{सर्वतः}\ A}}\skt{ }\edtext{\skt{स्रावमार्ग्गस्तु}}{\lemma{\skt{स्रावमार्ग्गस्तु}}\Afootnote{H; \skt{स्नावमार्गास्तु}\ A}}\skt{ }\edtext{\edtext{\skt{प्रद\-हे\-द्भि\-ष\-ग\-ग्निना}}{\lemma{\skt{प्रदहेद्}}\Afootnote{H; \skt{दहेद्}\ A}}}{\lemma{\skt{भिषगग्निना}}\Afootnote{H; \skt{वैद्यस्तथाऽग्निना}\ A}}\skt{~॥ }\pend
\pstart\versepar \skt{सुकुमारस्य }\edtext{\edtext{\skt{भीरोश्च}}{\lemma{\skt{भीरोश्}}\Afootnote{H; \skt{भीरो}\ A}}}{\lemma{\skt{च}}\Afootnote{H; \skt{र्हि}\ A}}\skt{ दुःकरः शतपोनकः~।}\pend

\pstart\versepar \ledsidenote{\passref{4.8.12}}\edtext{\skt{तत्र}}{\lemma{\skt{तत्र}}\Afootnote{H; \skt{रुजास्नावापहं तत्}\ A}}\skt{ }\edtext{\skt{धीरः}}{\lemma{\skt{धीरः}}\Afootnote{H; \emph{om.}\ A}}\skt{ }\edtext{\skt{प्रमाणज्ञः}}{\lemma{\skt{प्रमाणज्ञः}}\Afootnote{H; \skt{र}\ A}}\skt{ }\edtext{\edtext{\skt{स्वेदमाशु}}{\lemma{\skt{स्वेदम्}}\Afootnote{A; \skt{स्वेद}\ H}}}{\lemma{\skt{आशु}}\Afootnote{A; \skt{माशु}\ H}}\skt{ प्रयोजयेत्~॥ }\pend
\pstart\versepar \skt{}\edtext{\skt{स्वेद\-द्र\-व्यै\-र्य\-थो\-क्तैस्तु}}{\lemma{\skt{यथोक्तैस्तु}}\Afootnote{H; \skt{यथोद्दिष्टैः}\ A}}\skt{ }\edtext{\edtext{\skt{कृस\-रा\-पा\-य\-सै\-स्तथा}}{\lemma{\skt{कृस\appbrk{}रा\appbrk{}पा\appbrk{}य\appbrk{}सैस्}}\Afootnote{H; \skt{कृशरापाय}\ A}}}{\lemma{\skt{तथा}}\Afootnote{H; \skt{सादिभिः}\ A}}\skt{~।}\pend

\pstart\versepar \ledsidenote{\passref{4.8.13}}\edtext{\edtext{\skt{ग्राम्या\-नू\-पो\-द\-कै\-र्म्मां\-सै\-र्ल्ला\-वा\-द्यै\-र्व्वापि}}{\lemma{\skt{ग्राम्या\appbrk{}नू\appbrk{}पो\appbrk{}द\appbrk{}कैर्}}\Afootnote{H; \skt{ग्राम्या\appbrk{}नू\appbrk{}पौ\appbrk{}द\appbrk{}कैर्}\ A}}}{\lemma{\skt{व्वापि}}\Afootnote{H; \skt{वाऽपि}\ A}}\skt{ विष्किरैः~॥ }\pend
\pstart\versepar \skt{}\edtext{\skt{वृक्षा\-द\-नी\-म\-थै\-र\-ण्ड\-म्बि\-त्वा\-दिञ्च}}{\lemma{\skt{बित्वादिञ्च}}\Afootnote{H; \skt{बिल्वादिं च}\ A}}\skt{ गणन्तथा~।}\pend

\pstart\versepar \ledsidenote{\passref{4.8.14}}\skt{कषायं सुकृतं कृत्वा }\edtext{\skt{ततः}}{\lemma{\skt{ततः}}\Afootnote{H; \skt{स्नेह}\ A}}\skt{ कुम्भे }\edtext{\skt{निधापयेत्}}{\lemma{\skt{निधापयेत्}}\Afootnote{H; \skt{निषेचयेत्}\ A}}\skt{~॥ }\pend
\pstart\versepar \skt{नाडीस्वेदेन तेनास्य }\edtext{\skt{तद्व्रणं}}{\lemma{\skt{तद्}}\Afootnote{H; \skt{तं}\ A}}\skt{ }\edtext{\skt{स्वेदयेत्ततः}}{\lemma{\skt{ततः}}\Afootnote{H; \skt{भिषक्}\ A}}\skt{~।}\pend

\pstart\versepar \ledsidenote{\passref{4.8.15}}\edtext{\skt{तिला\-न\-त\-सि\-बी\-जानि}}{\lemma{\skt{तिला\appbrk{}न\appbrk{}त\appbrk{}सि\appbrk{}बी\appbrk{}जानि}}\Afootnote{H; \skt{तिलै\appbrk{}र\appbrk{}ण्डा\appbrk{}त\appbrk{}सी\appbrk{}माष}\ A}}\skt{ }\edtext{\skt{यव\-गो\-धू\-म\-को\-द्र\-वान्}}{\lemma{\skt{यव\appbrk{}गो\appbrk{}धू\appbrk{}म\appbrk{}को\appbrk{}द्र\appbrk{}वान्}}\Afootnote{H; \skt{यव\appbrk{}गो\appbrk{}धू\appbrk{}म\appbrk{}स\appbrk{}र्ष\appbrk{}पान्}\ A}}\skt{~॥ }\pend
\pstart\versepar \skt{}\edtext{\skt{लव\-णा\-म्ल\-क\-व\-र्ग्गश्च}}{\lemma{\skt{लव\appbrk{}णा\appbrk{}म्ल\appbrk{}क\appbrk{}व\appbrk{}र्ग्गश्च}}\Afootnote{H; \skt{लव\appbrk{}णा\appbrk{}न्य\appbrk{}म्ल\appbrk{}वर्गं च}\ A}}\skt{ }\edtext{\edtext{\skt{स्थाल्यामाधाय}}{\lemma{\skt{स्थाल्याम्}}\Afootnote{H; \skt{स्थाल्या}\ A}}}{\lemma{\skt{आधाय}}\Afootnote{H; \skt{मेवोप}\ A}}\skt{ साधयेत्~।}\pend

\pstart\versepar \ledsidenote{\passref{4.8.16}}\skt{आतुरं }\edtext{\skt{स्वेदयेत्तत्र}}{\lemma{\skt{तत्र}}\Afootnote{H; \skt{तेन}\ A}}\skt{ तथा सिध्यति कुर्व्वतः~॥ }\pend
\pstart\versepar \skt{स्विन्नञ्च पाययेदेनं कुष्ठञ्च लवणानि च~।}\pend

\pstart\versepar \ledsidenote{\passref{4.8.17}}\edtext{\skt{वचा\-हि\-ङ्ग्व\-ज\-मो\-दश्च}}{\lemma{\skt{वचा\appbrk{}हि\appbrk{}ङ्ग्व\appbrk{}ज\appbrk{}मो\appbrk{}दश्च}}\Afootnote{H; \skt{वचा\appbrk{}हि\appbrk{}ङ्ग्व\appbrk{}ज\appbrk{}मोदं च}\ A}}\skt{ समभागानि सर्प्पिषा~॥ }\pend
\pstart\versepar \skt{}\edtext{\skt{मार्द्वी\-के\-ना\-थ\-वा\-म्लेन}}{\lemma{\skt{मार्द्वी\appbrk{}के\appbrk{}ना\appbrk{}थ\appbrk{}वा\appbrk{}म्लेन}}\Afootnote{H; \skt{मार्द्वी\appbrk{}के\appbrk{}ना\appbrk{}थ\appbrk{}वाऽ\appbrk{}म्लेन}\ A}}\skt{ सुरा\-सौ\-वी\-र\-केण वा~।}\pend

\pstart\versepar \ledsidenote{\passref{4.8.18}}\skt{ततो मधुकतैलेन तस्य }\edtext{\edtext{\skt{सिञ्चे\-द्व्र\-ण\-म्भि\-षक्}}{\lemma{\skt{व्रणम्}}\Afootnote{H; \emph{om.}\ A}}}{\lemma{\skt{भिषक्}}\Afootnote{H; \skt{भिषग्}\ A}}\skt{~॥ }\pend
\pstart\versepar \skt{}\edtext{\skt{प्रति\-सा\-र\-ये\-द्गु\-द\-ञ्चास्य}}{\lemma{\skt{प्रतिसारयेद्}}\Afootnote{H; \skt{ब्रणम्~। परिषिञ्चेद्}\ A}}\skt{ तैलै\-र्व्वा\-त\-रु\-जा\-पहैः~।}\pend

\pstart\versepar \ledsidenote{\passref{4.8.19}}\edtext{\skt{विधिनानेन}}{\lemma{\skt{विधिनानेन}}\Afootnote{H; \skt{विधिनाऽनेन}\ A}}\skt{ विण्मूत्रं }\edtext{\skt{स्वमा\-र्ग्ग\-म\-नु\-प\-द्यते}}{\lemma{\skt{अनुपद्यते}}\Afootnote{H; \skt{अधिगच्छति}\ A}}\skt{~॥ }\pend
\pstart\versepar \skt{अन्ये }\edtext{\edtext{\skt{चोपद्रवास्तीव्राः}}{\lemma{\skt{चोपद्रवास्}}\Afootnote{H; \skt{चोपद्रवा}\ A}}}{\lemma{\skt{तीव्राः}}\Afootnote{H; \skt{स्तीव्राः}\ A}}\skt{ }\edtext{\edtext{\skt{सिध्यन्त्यत्र}}{\lemma{\skt{सिध्यन्त्य्}}\Afootnote{H; \skt{सिध्यन्त्}\ A}}\skt{ न}}{\lemma{\skt{अत्र न}}\Afootnote{H; \skt{यत्र न}\ A}}\skt{ संशयः~।}\pend

\pstart\versepar \ledsidenote{\passref{4.8.20}}\edtext{\skt{तपोनके}}{\lemma{\skt{तपोनके}}\Afootnote{H; \skt{शतपोन}\ A}}\skt{ }\edtext{\skt{समा\-ख्या\-त\-मु\-ष्ट्र\-ग्रीवे}}{\lemma{\skt{समाख्यातम्}}\Afootnote{H; \skt{क आख्यात}\ A}}\skt{ क्रियां शृणु~॥ }\pend
\pstart\versepar \skt{}\edtext{\edtext{\skt{एष्यो\-ष्ट्र\-ग्री\-व\-मा\-साद्य}}{\lemma{\skt{एष्योष्ट्रग्रीवम्}}\Afootnote{H; \skt{अथोष्ट्रग्रीव}\ A}}}{\lemma{\skt{आसाद्य}}\Afootnote{H; \skt{मेषित्वा}\ A}}\skt{ च्छित्वा क्षार\-न्नि\-पा\-त\-येत्~।}\pend

\pstart\versepar \ledsidenote{\passref{4.8.21}}\edtext{\skt{पूतिमांसं}}{\lemma{\skt{पूतिमांसं}}\Afootnote{H; \skt{पूतिमांस}\ A}}\skt{ }\edtext{\edtext{\skt{व्यपो\-हा\-र्थ\-म\-ग्नि\-रत्र}}{\lemma{\skt{व्यपोहार्थम्}}\Afootnote{H; \skt{व्यपोहार्थ}\ A}}\skt{ न}}{\lemma{\skt{अग्निरत्र न}}\Afootnote{H; \skt{मग्निरत्र न}\ A}}\skt{ पूजितः~। }\pend
\pstart\versepar \skt{अथैनं }\edtext{\skt{घृत\-सं\-वृ\-ष्टै\-स्तिलैः}}{\lemma{\skt{घृतसंवृष्टैस्}}\Afootnote{H; \skt{घृतसंसृष्टैस्}\ A}}\skt{ पिष्टैः }\edtext{\skt{प्रलेपयेत्}}{\lemma{\skt{प्रलेपयेत्}}\Afootnote{H; \skt{प्रपेपयेत्}\ A}}\skt{~।}\pend

\pstart\versepar \ledsidenote{\passref{4.8.22}}\skt{बन्धं }\edtext{\skt{यथावत्कुर्व्वीत}}{\lemma{\skt{यथावत्}}\Afootnote{H; \skt{ततोऽनु}\ A}}\skt{ }\edtext{\skt{परिषेकश्च}}{\lemma{\skt{परिषेकश्च}}\Afootnote{H; \skt{परिषेकं तु}\ A}}\skt{ सर्प्पिषा~॥ }\pend
\pstart\versepar \skt{तृतीये दिवसे मुक्त्वा यथास्वं शोधयेद्भिषक्~।}\pend

\pstart \ledsidenote{\passref{4.8.23ab}}\skt{ततः }\edtext{\edtext{\skt{शुद्धिम्विदित्वा}}{\lemma{\skt{शुद्धिम्}}\Afootnote{H; \skt{शुद्धं}\ A}}\skt{ तु}}{\lemma{\skt{विदित्वा तु}}\Afootnote{H; \skt{विदित्वा च}\ A}}\skt{ }\edtext{\skt{रोपयीत}}{\lemma{\skt{रोपयीत}}\Afootnote{H; \skt{रोपयेत्तु}\ A}}\skt{ यथाक्रमं~॥}\pend

\pstart \ledsidenote{\passref{4.8.23.add}}\skt{एतत्कर्म्म समाख्यातं सर्व्वे\-षा\-म\-नु\-पू\-र्व्वशः~।}\pend

\pstart \ledsidenote{\passref{4.8.23cd}}\edtext{\skt{उत्कृ\-त्या\-श्रा\-व\-मा\-र्ग्गन्तु}}{\lemma{\skt{उत्कृ\appbrk{}त्या\appbrk{}श्रा\appbrk{}व\appbrk{}मा\appbrk{}र्ग्गन्तु}}\Afootnote{H; \skt{उत्कृ\appbrk{}त्या\appbrk{}स्रा\appbrk{}व\appbrk{}मा\appbrk{}र्गंस्तु}\ A}}\skt{ परिस्राविणि बुद्धिमान्~॥}\pend

\pstart\versepar \ledsidenote{\passref{4.8.24}}\skt{क्षारेण वा स्राव\-ग\-ति\-न्द\-हे\-द्धु\-त\-व\-हेन वा~। }\pend
\pstart\versepar \skt{}\edtext{\skt{सुखो\-ष्णे\-ना\-नु\-तै\-लेन}}{\lemma{\skt{सुखो\appbrk{}ष्णे\appbrk{}ना\appbrk{}नु\appbrk{}तै\appbrk{}लेन}}\Afootnote{H; \skt{सुखो\appbrk{}ष्णे\appbrk{}ना\appbrk{}णु\appbrk{}तै\appbrk{}लेन}\ A}}\skt{ सेच\-ये\-द्गु\-द\-म\-ण्डलं~॥}\pend

\pstart\versepar \ledsidenote{\passref{4.8.25}}\edtext{\edtext{\skt{उप\-ना\-ह\-म्प्र\-दे\-हांश्च}}{\lemma{\skt{उपनाहम्}}\Afootnote{H; \skt{उपनाहाः}\ A}}}{\lemma{\skt{प्रदेहांश्च}}\Afootnote{H; \skt{प्रदेहाश्च}\ A}}\skt{ }\edtext{\skt{मूत्र\-क्षा\-र\-स\-म\-न्वितं}}{\lemma{\skt{मूत्र\appbrk{}क्षा\appbrk{}र\appbrk{}स\appbrk{}म\appbrk{}न्वितं}}\Afootnote{H; \skt{मूत्र\appbrk{}क्षा\appbrk{}र\appbrk{}स\appbrk{}म\appbrk{}न्विताः}\ A}}\skt{~। }\pend
\pstart\versepar \skt{}\edtext{\edtext{\skt{वाम\-नी\-यौ\-ष\-धै\-र्यु\-क्ता\-म्प\-रि\-षे\-कांश्च}}{\lemma{\skt{युक्ताम्}}\Afootnote{H; \skt{कार्याः}\ A}}}{\lemma{\skt{परिषेकांश्च}}\Afootnote{H; \skt{परिषेकाश्च}\ A}}\skt{ मात्रया~॥}\pend

\pstart\versepar \ledsidenote{\passref{4.8.26}}\skt{मृदुभूतं }\edtext{\edtext{\skt{विदि\-त्वै\-न\-म\-ल्प\-स्रावं}}{\lemma{\skt{विदित्वैनम्}}\Afootnote{H; \skt{विदित्वैन}\ A}}}{\lemma{\skt{अल्पस्रावं}}\Afootnote{H; \skt{मल्पस्राव}\ A}}\skt{ }\edtext{\skt{रुजान्वितं}}{\lemma{\skt{रुजान्वितं}}\Afootnote{H; \skt{रुगन्वितम्}\ A}}\skt{~। }\pend
\pstart\versepar \skt{गतिमन्विष्य शस्त्रेण }\edtext{\skt{च्छिन्द्या\-त्ख\-र्ज्जू\-र\-प\-त्र\-कम्}}{\lemma{\skt{खर्ज्जू\appbrk{}र\appbrk{}प\appbrk{}त्र\appbrk{}कम्}}\Afootnote{H; \skt{खरु\appbrk{}जू\appbrk{}र\appbrk{}प\appbrk{}त्र\appbrk{}कम्}\ A}}\skt{~॥}\pend

\pstart\versepar \ledsidenote{\passref{4.8.27}}\edtext{\skt{चन्द्रार्ध}}{\lemma{\skt{चन्द्रार्ध}}\Afootnote{H; \skt{चन्द्रार्धं}\ A}}\skt{ }\edtext{\skt{चन्द्रं}}{\lemma{\skt{चन्द्रं}}\Afootnote{H; \skt{चन्द्र}\ A}}\skt{ }\edtext{\skt{वज्रं च}}{\lemma{\skt{वज्रं च}}\Afootnote{H; \skt{चक्रं च}\ A}}\skt{ }\edtext{\skt{सूचीमुखं}}{\lemma{\skt{सूचीमुखं}}\Afootnote{H; \skt{सूचीमुस्व}\ A}}\skt{ }\edtext{\skt{पराङ्मुखं}}{\lemma{\skt{पराङ्मुखं}}\Afootnote{H; \skt{मवाङ्मुखम्}\ A}}\skt{~। }\pend
\pstart\versepar \skt{}\edtext{\skt{छित्त्वाग्निना}}{\lemma{\skt{छित्त्वाग्निना}}\Afootnote{H; \skt{छित्त्वाऽग्निना}\ A}}\skt{ }\edtext{\skt{दहेच्चापि}}{\lemma{\skt{चापि}}\Afootnote{H; \skt{सम्}\ A}}\skt{ }\edtext{\skt{पाद्यक्षारेण वा}}{\lemma{\skt{पाद्यक्षारेण वा}}\Afootnote{H; \skt{यगेवं क्षारेण वा}\ A}}\skt{ }\edtext{\skt{ततः}}{\lemma{\skt{ततः}}\Afootnote{H; \skt{पुनः}\ A}}\skt{~॥}\pend

\pstart\versepar \ledsidenote{\passref{4.8.28}}\skt{ततः }\edtext{\edtext{\skt{संशो\-ध\-नै\-स्ती\-क्ष्णै\-र्म्मृ\-दु\-पू\-र्व्व\-म्वि\-शो\-ध\-येत्}}{\lemma{\skt{तीक्ष्णैर्}}\Afootnote{H; \skt{एव}\ A}}}{\lemma{\skt{म्मृदुपूर्व्वम्}}\Afootnote{H; \skt{मृदुपूर्वैर्}\ A}}\skt{~। }\pend
\pstart\versepar \skt{}\edtext{\skt{बहिरन्तर्म्मुखं}}{\lemma{\skt{बहिरन्तर्म्मुखं}}\Afootnote{H; \skt{बहि\appbrk{}र\appbrk{}न्त\appbrk{}र्मु\appbrk{}खश्}\ A}}\skt{ चापि }\edtext{\edtext{\skt{शिशोर्यस्य}}{\lemma{\skt{शिशोर्}}\Afootnote{H; \skt{शिशो}\ A}}}{\lemma{\skt{यस्य}}\Afootnote{H; \skt{र्यस्य}\ A}}\skt{ }\edtext{\skt{भगन्दरं}}{\lemma{\skt{भगन्दरं}}\Afootnote{H; \skt{भगन्दरः}\ A}}\skt{~॥}\pend

\pstart\versepar \ledsidenote{\passref{4.8.29}}\edtext{\skt{तस्या\-हि\-त\-म्वि\-रे\-काग्निः}}{\lemma{\skt{विरेकाग्निः}}\Afootnote{H; \skt{विरेकाग्नि}\ A}}\skt{ }\edtext{\skt{शस्त्र\-क्षा\-रा\-व\-ता\-रणं}}{\lemma{\skt{शस्त्र\appbrk{}क्षा\appbrk{}रा\appbrk{}व\appbrk{}ता\appbrk{}रणं}}\Afootnote{H; \skt{शस्त्र\appbrk{}क्षा\appbrk{}रा\appbrk{}व\appbrk{}चा\appbrk{}र\appbrk{}णम्}\ A}}\skt{~। }\pend
\pstart\versepar \skt{यद्यत्मृदु च तीक्ष्णञ्च }\edtext{\edtext{\edtext{\skt{तत्त\-स्या\-म\-व\-ता\-र\-येत्}}{\lemma{\skt{तत्}}\Afootnote{H; \skt{तत्तत्}\ A}}}{\lemma{\skt{तस्याम्}}\Afootnote{H; \skt{तस्या}\ A}}}{\lemma{\skt{अवतारयेत्}}\Afootnote{H; \skt{वचारयेत्}\ A}}\skt{~॥}\pend

\pstart\versepar \ledsidenote{\passref{4.8.30}}\edtext{\skt{रज\-न्या\-र\-ग्व\-धौ\-का\-ला\-चू\-र्ण्ण\-म्म\-धु\-घृ\-ता\-प्लुतं}}{\lemma{\skt{रज\appbrk{}न्या\appbrk{}र\appbrk{}ग्व\appbrk{}धौ\appbrk{}का\appbrk{}ला\appbrk{}चू\appbrk{}र्ण्णम्}}\Afootnote{H; \skt{आर\appbrk{}ग्व\appbrk{}ध\appbrk{}नि\appbrk{}शा\appbrk{}का\appbrk{}ला\appbrk{}चूर्णं}\ A}}\skt{~। }\pend
\pstart\versepar \skt{}\edtext{\skt{सूत्र\-व\-र्त्ति\-प्र\-णि\-हितं}}{\lemma{\skt{सूत्र\appbrk{}व\appbrk{}र्त्ति\appbrk{}प्र\appbrk{}णि\appbrk{}हितं}}\Afootnote{H; \skt{अग्र\appbrk{}व\appbrk{}र्ति\appbrk{}प्र\appbrk{}णि\appbrk{}हितं}\ A}}\skt{ व्रणानां }\edtext{\skt{शोधनम्परं}}{\lemma{\skt{परं}}\Afootnote{H; \skt{हितम्}\ A}}\skt{~॥}\pend

\pstart\versepar \ledsidenote{\passref{4.8.31}}\skt{योगोऽ\-य\-न्ना\-श\-य\-त्याशु }\edtext{\edtext{\skt{गति\-म्मे\-घ\-मि\-वा\-निलः}}{\lemma{\skt{मेघम्}}\Afootnote{H; \skt{मेघ}\ A}}}{\lemma{\skt{इवानिलः}}\Afootnote{H; \skt{मिवानिलः}\ A}}\skt{~। }\pend
\pstart\versepar \skt{आगन्तुजे }\edtext{\skt{भिषडीं}}{\lemma{\skt{भिषडीं}}\Afootnote{H; \skt{भिषङ्नाडीं}\ A}}\skt{ शस्त्रेणोत्कृत्य यत्नतः~॥}\pend

\pstart\versepar \ledsidenote{\passref{4.8.32}}\edtext{\skt{जम्वो\-ष्टे\-ना\-ग्नि\-व\-र्ण्णेन}}{\lemma{\skt{जम्वो\appbrk{}ष्टे\appbrk{}ना\appbrk{}ग्नि\appbrk{}व\appbrk{}र्ण्णेन}}\Afootnote{H; \skt{जम्ब्वो\appbrk{}ष्टे\appbrk{}ना\appbrk{}ग्नि\appbrk{}व\appbrk{}र्णेन}\ A}}\skt{ तप्तया वा शलाकया~। }\pend
\pstart\versepar \skt{}\edtext{\skt{दहे\-द्य\-थो\-क्त\-म्म\-ति\-मान्तं}}{\lemma{\skt{तं}}\Afootnote{H; \skt{स्तं}\ A}}\skt{ व्रणं सुसमाहितः~॥}\pend

\pstart\versepar \ledsidenote{\passref{4.8.33}}\edtext{\skt{क्रिमिघ्नञ्च}}{\lemma{\skt{क्रिमिघ्नञ्च}}\Afootnote{H; \skt{कृमिघ्नं च}\ A}}\skt{ }\edtext{\skt{विधि\-ङ्कु\-र्या\-च्छ\-ल्या\-प\-न\-य\-मेव}}{\lemma{\skt{छल्यापनयम्}}\Afootnote{H; \skt{छल्यानयनम्}\ A}}\skt{ च~। }\pend
\pstart\versepar \skt{प्रत्याख्यायैष चारभ्यो }\edtext{\skt{वर्ज्याश्चापि}}{\lemma{\skt{वर्ज्याश्}}\Afootnote{H; \skt{वर्ज्यश्}\ A}}\skt{ }\edtext{\skt{त्रिदोषजाः}}{\lemma{\skt{त्रिदोषजाः}}\Afootnote{H; \skt{त्रिदोषजः}\ A}}\skt{~॥}\pend

\pstart\versepar \ledsidenote{\passref{4.8.34}}\skt{एतत्कर्म्म समाख्यातं }\edtext{\edtext{\skt{सर्व्वे\-षा\-म\-नु\-पू\-र्व्वशः}}{\lemma{\skt{सर्व्वेषाम्}}\Afootnote{H; \skt{सर्वेषा}\ A}}}{\lemma{\skt{अनुपूर्व्वशः}}\Afootnote{H; \skt{मनुपूर्वशः}\ A}}\skt{~। }\pend
\pstart\versepar \skt{एषां तु }\edtext{\skt{शस्त्रपातेन}}{\lemma{\skt{शस्त्रपातेन}}\Afootnote{H; \skt{शस्त्रपतना}\ A}}\skt{ }\edtext{\skt{वेदना}}{\lemma{\skt{वेदना}}\Afootnote{H; \skt{द्वेदना}\ A}}\skt{ }\edtext{\skt{यदि}}{\lemma{\skt{यदि}}\Afootnote{H; \skt{यत्र}\ A}}\skt{ जायते~॥}\pend

\pstart\versepar \ledsidenote{\passref{4.8.35}}\edtext{\skt{तत्रा\-नु\-तै\-ले\-नो\-ष्णेन}}{\lemma{\skt{तत्रा\appbrk{}नु\appbrk{}तै\appbrk{}ले\appbrk{}नो\appbrk{}ष्णेन}}\Afootnote{H; \skt{तत्रा\appbrk{}णु\appbrk{}तै\appbrk{}ले\appbrk{}नो\appbrk{}ष्णेन}\ A}}\skt{ }\edtext{\skt{पायोःसेकः}}{\lemma{\skt{पायोःसेकः}}\Afootnote{H; \skt{परिषेकः}\ A}}\skt{ प्रशस्यते~। }\pend
\pstart\versepar \skt{}\edtext{\skt{वात\-घ्न\-क्वा\-थ\-सं\-पूर्ण्णां}}{\lemma{\skt{वात\appbrk{}घ्न\appbrk{}क्वा\appbrk{}थ\appbrk{}सं\appbrk{}पूर्ण्णां}}\Afootnote{H; \skt{वात\appbrk{}घ्नौ\appbrk{}ष\appbrk{}ध\appbrk{}सं\appbrk{}पूर्णां}\ A}}\skt{ स्थालीं }\edtext{\skt{च्छिद्रसराविकां}}{\lemma{\skt{च्छिद्रसराविकां}}\Afootnote{H; \skt{छिद्र\appbrk{}श\appbrk{}रा\appbrk{}वि\appbrk{}काम्}\ A}}\skt{~॥}\pend

\pstart\versepar \ledsidenote{\passref{4.8.36}}\skt{स्नेहा\-भ्य\-क्त\-गु\-द\-स्त\-प्ता\-म\-ध्या\-सीत सबाष्पिकान्~। }\pend
\pstart\versepar \skt{}\edtext{\skt{नान्या}}{\lemma{\skt{नान्या}}\Afootnote{H; \skt{नाड्या}\ A}}\skt{ }\edtext{\skt{वास्पाहरेत्स्वेदं}}{\lemma{\skt{वास्पाहरेत्}}\Afootnote{H; \skt{वाऽस्याहरेत्}\ A}}\skt{ शयानस्य रुजाहरन्~॥ }\pend
\pstart\versepar \skt{उदकोष्ठे }\edtext{\skt{वगाह्ये वा}}{\lemma{\skt{वगाह्ये वा}}\Afootnote{H; \skt{वगाह्यो वा}\ A}}\skt{ तथा शाम्यति वेदना~।}\edtext{}{\lemma{\skt{रुजाहरन्~॥ उदकोष्ठे}}\Afootnote{H; \skt{रुजापहम्~। उष्णोदकेऽ}\ A}}\pend

\pstart\versepar \ledsidenote{\passref{4.8.37}}\edtext{\skt{कद\-ली\-मृ\-ग\-लो\-पा\-क\-पृ\-ष\-ता\-जि\-न\-सं\-वृतां}}{\lemma{\skt{कद\appbrk{}ली\appbrk{}मृ\appbrk{}ग\appbrk{}लो\appbrk{}पा\appbrk{}क\appbrk{}पृ\appbrk{}ष\appbrk{}ता\appbrk{}जि\appbrk{}न\appbrk{}सं\appbrk{}वृतां}}\Afootnote{H; \skt{कद\appbrk{}ली\appbrk{}मृ\appbrk{}ग\appbrk{}लो\appbrk{}पा\appbrk{}क\appbrk{}प्रि\appbrk{}य\appbrk{}का\appbrk{}जि\appbrk{}न\appbrk{}सं\appbrk{}भृ\appbrk{}तान्}\ A}}\skt{~॥ }\pend
\pstart\versepar \skt{}\edtext{\skt{कार\-ये\-दु\-प\-ना\-हाम्वा}}{\lemma{\skt{कार\appbrk{}ये\appbrk{}दु\appbrk{}प\appbrk{}ना\appbrk{}हाम्वा}}\Afootnote{H; \skt{कार\appbrk{}ये\appbrk{}दु\appbrk{}प\appbrk{}ना\appbrk{}हांश्च}\ A}}\skt{ }\edtext{\skt{साल्व\-ला\-दी\-म्वि\-च\-क्षणः}}{\lemma{\skt{साल्वलादीम्}}\Afootnote{H; \skt{साल्वणादीन}\ A}}\skt{~।}\pend

\pstart\versepar \ledsidenote{\passref{4.8.38}}\skt{कटुत्रिकं }\edtext{\edtext{\skt{वचा\-हि\-ङ्गु\-ल\-व\-णा\-न्यथ}}{\lemma{\skt{वचा\appbrk{}हि\appbrk{}ङ्गु\appbrk{}ल\appbrk{}व\appbrk{}णान्य्}}\Afootnote{H; \skt{वचा\appbrk{}हि\appbrk{}ङ्गु\appbrk{}ल\appbrk{}व\appbrk{}णान्}\ A}}}{\lemma{\skt{अथ}}\Afootnote{H; \skt{यथ}\ A}}\skt{ दीप्यकं~॥ }\pend
\pstart\versepar \skt{}\edtext{\skt{पाय\-ये\-द\-म्ल\-कौ\-लन्थां}}{\lemma{\skt{पाय\appbrk{}ये\appbrk{}द\appbrk{}म्ल\appbrk{}कौ\appbrk{}लन्थां}}\Afootnote{H; \skt{पाय\appbrk{}ये\appbrk{}च्चा\appbrk{}म्ल\appbrk{}कौ\appbrk{}लत्थ}\ A}}\skt{ सुरा\-सौ\-वी\-र\-का\-दिभिः~।}\pend

\pstart\versepar \ledsidenote{\passref{4.8.39}}\edtext{\skt{ज्योति\-ष्म\-ती\-ला\-ङ्ग\-ल\-की\-श्या\-मा\-द\-न्ती\-तृ\-वृ\-त्तिला}}{\lemma{\skt{ज्योति\appbrk{}ष्म\appbrk{}ती\appbrk{}ला\appbrk{}ङ्ग\appbrk{}ल\appbrk{}की\appbrk{}श्या\appbrk{}मा\appbrk{}द\appbrk{}न्ती\appbrk{}तृ\appbrk{}वृ\appbrk{}त्तिला}}\Afootnote{H; \skt{ज्योति\appbrk{}ष्म\appbrk{}ती\appbrk{}ला\appbrk{}ङ्ग\appbrk{}ल\appbrk{}की\appbrk{}श्या\appbrk{}मा\appbrk{}द\appbrk{}न्ती\appbrk{}त्रि\appbrk{}वृ\appbrk{}त्तिलाः}\ A}}\skt{~। }\pend
\pstart\versepar \skt{कुष्ठं शताह्वा गोलोमी तिल्वको गिरिकर्ण्णिका~।}\pend

\pstart\versepar \ledsidenote{\passref{4.8.40}}\edtext{\skt{काञ्चनक्सीर्का}}{\lemma{\skt{काञ्चनक्सीर्का}}\Afootnote{H; \skt{कासीसं काञ्चनक्सीर्}\ A}}\skt{ }\edtext{\skt{चैव}}{\lemma{\skt{चैव}}\Afootnote{H; \skt{यौ}\ A}}\skt{ वर्ग्गः }\edtext{\skt{शोधनमिष्यते}}{\lemma{\skt{शोधनम्}}\Afootnote{H; \skt{शोधन}\ A}}\skt{~॥ }\pend
\pstart\versepar \skt{}\edtext{\skt{तृवृत्तिला}}{\lemma{\skt{तृवृत्तिला}}\Afootnote{H; \skt{त्रिवृत्तिला}\ A}}\skt{ नागदन्ती मञ्जिष्ठा }\edtext{\skt{सहसर्प्पिषा}}{\lemma{\skt{सहसर्प्पिषा}}\Afootnote{H; \skt{पयाअ सह}\ A}}\skt{~।}\pend

\pstart\versepar \ledsidenote{\passref{4.8.41}}\edtext{\skt{उत्सा\-द\-न\-म्भ\-वे\-दे\-त\-त्सै\-न्ध\-व\-क्षौ\-द्र\-सं\-युतं}}{\lemma{\skt{भवेदेतत्}}\Afootnote{H; \skt{भवेदेतत्}\ A}}\skt{~॥ }\pend
\pstart\versepar \skt{रसाञ्जनं हरिद्रे द्वे }\edtext{\skt{मञ्जि\-ष्ठा\-नि\-म्ब\-प\-ल्लवा}}{\lemma{\skt{मञ्जि\appbrk{}ष्ठा\appbrk{}नि\appbrk{}म्ब\appbrk{}प\appbrk{}ल्लवा}}\Afootnote{H; \skt{मञ्जि\appbrk{}ष्ठा\appbrk{}नि\appbrk{}म्ब\appbrk{}प\appbrk{}ल्लवाः}\ A}}\skt{~।}\pend

\pstart\versepar \ledsidenote{\passref{4.8.42}}\edtext{\skt{तृवृ\-त्ते\-तो\-व\-ती\-द\-न्ती\-कल्को}}{\lemma{\skt{तृवृ\appbrk{}त्ते\appbrk{}तो\appbrk{}व\appbrk{}ती\appbrk{}द\appbrk{}न्ती\appbrk{}कल्को}}\Afootnote{H; \skt{त्रिवृ\appbrk{}त्ते\appbrk{}जो\appbrk{}व\appbrk{}ती\appbrk{}द\appbrk{}न्ती\appbrk{}कल्को}\ A}}\skt{ नाडीव्रणापहः~॥ }\pend
\pstart\versepar \skt{}\edtext{\skt{कुष्ठ\-न्ति\-ला\-त्रि\-वृ\-द्दन्ती}}{\lemma{\skt{तिलात्रिवृद्}}\Afootnote{H; \skt{त्रिवृत्तिला}\ A}}\skt{ मागध्यः }\edtext{\skt{सैन्धवम्मधुः}}{\lemma{\skt{मधुः}}\Afootnote{H; \skt{मधु}\ A}}\skt{~।}\pend

\pstart\versepar \ledsidenote{\passref{4.8.43}}\skt{रजनी त्रिफला तुत्थं हितं स्याद्व्र\-ण\-शो\-ध\-नम्~॥ }\pend
\pstart\versepar \skt{मागध्यो मधुकं }\edtext{\skt{लोध्रं}}{\lemma{\skt{लोध्रं}}\Afootnote{H; \skt{रोध्रं}\ A}}\skt{ कुष्ठमेला हरेणवः~।}\pend

\pstart\versepar \ledsidenote{\passref{4.8.44}}\skt{समङ्गा धातकी चैव सारिवा रजनीद्वयं~॥ }\pend
\pstart\versepar \skt{प्रियङ्गवः }\edtext{\skt{सर्ज्जरसं}}{\lemma{\skt{सर्ज्जरसं}}\Afootnote{H; \skt{सर्जरसः}\ A}}\skt{ पद्म\-क\-म्प\-द्म\-के\-सरं~।}\pend

\pstart\versepar \ledsidenote{\passref{4.8.45}}\edtext{\skt{मातुलुङ्गस्य}}{\lemma{\skt{मातुलुङ्गस्य}}\Afootnote{H; \skt{सुधा वचा लाङ्ग}\ A}}\skt{ }\edtext{\skt{पत्राणि}}{\lemma{\skt{पत्राणि}}\Afootnote{H; \skt{लकी}\ A}}\skt{ मधूच्छिष्टं ससैन्धवं~॥ }\pend
\pstart\versepar \skt{एतत्सम्भृत्य सम्भारं तैलन्धीरो विपाचयेत्~।}\pend

\pstart\versepar \ledsidenote{\passref{4.8.46}}\edtext{\skt{एतद्धि}}{\lemma{\skt{एतद्धि}}\Afootnote{H; \skt{एतद्वै}\ A}}\skt{ गण्डमालासु }\edtext{\skt{मण्डले}}{\lemma{\skt{मण्डले}}\Afootnote{H; \skt{मण्डलेष्व्}\ A}}\skt{ }\edtext{\skt{मधुमेहिषु}}{\lemma{\skt{मधुमेहिषु}}\Afootnote{H; \skt{अथ मेहिषु}\ A}}\skt{~॥ }\pend
\pstart\versepar \skt{रोपणार्थं हितन्तैलं भग\-न्द\-र\-वि\-ना\-शनं~।}\pend

\pstart \ledsidenote{\passref{4.8.47ab}}\edtext{\edtext{\skt{न्यग्रोधादिर्ग्गणो}}{\lemma{\skt{न्यग्रोधादिर्}}\Afootnote{H; \skt{न्यग्रोधादि}\ A}}\skt{ यश्च}}{\lemma{\skt{ग्गणो यश्च}}\Afootnote{H; \skt{गणश्चैव}\ A}}\skt{ हितः }\edtext{\skt{शोधनरोपणः}}{\lemma{\skt{शोधनरोपणः}}\Afootnote{H; \skt{शोधनरोपणे}\ A}}\skt{~॥}\pend

\pstart\versepar \ledsidenote{\passref{4.8.48}}\edtext{\skt{त्रृवृ\-द्द\-न्ती\-ह\-रि\-द्रा\-र्क्क\-मूलं}}{\lemma{\skt{त्रृवृ\appbrk{}द्द\appbrk{}न्ती\appbrk{}ह\appbrk{}रि\appbrk{}द्रा\appbrk{}र्क्क\appbrk{}मूलं}}\Afootnote{H; \skt{त्रिवृ\appbrk{}द्द\appbrk{}न्ती\appbrk{}ह\appbrk{}रि\appbrk{}द्रा\appbrk{}र्क\appbrk{}मूलं}\ A}}\skt{ }\edtext{\skt{लोहाश्वमारकौः}}{\lemma{\skt{लोहाश्वमारकौः}}\Afootnote{H; \skt{लोहाश्वमारकौ}\ A}}\skt{~। }\pend
\pstart\versepar \skt{विड\-ङ्ग\-सा\-र\-न्त्रि\-फला }\edtext{\skt{अर्क्कस्नुहीपयो}}{\lemma{\skt{अर्क्कस्नुहीपयो}}\Afootnote{H; \skt{स्नुह्यर्कपयसी}\ A}}\skt{ मधु~॥}\pend

\pstart\versepar \ledsidenote{\passref{4.8.49}}\edtext{\skt{मधू\-च्छि\-ष्ट\-स\-मा\-यु\-क्तै\-स्तै\-ल\-मे\-तै\-र्व्वि\-पा\-चितं}}{\lemma{\skt{व्विपाचितं}}\Afootnote{H; \skt{विपाचयेत्}\ A}}\skt{~। }\pend
\pstart\versepar \skt{}\edtext{\edtext{\skt{भग\-न्द\-र\-वि\-ना\-शा\-र्थ\-म्व्र\-ण्रो\-गा\-पहं}}{\lemma{\skt{भग\appbrk{}न्द\appbrk{}र\appbrk{}वि\appbrk{}ना\appbrk{}शा\appbrk{}र्थम्}}\Afootnote{H; \skt{भग\appbrk{}न्द\appbrk{}र\appbrk{}वि\appbrk{}ना\appbrk{}शार्}\ A}}}{\lemma{\skt{व्रण्रोगापहं}}\Afootnote{H; \skt{थ}\ A}}\skt{ }\edtext{\skt{शुभं}}{\lemma{\skt{शुभं}}\Afootnote{H; \skt{मेतद्योज्यं विशेषतः}\ A}}\skt{~॥}\pend

\pstart\versepar \ledsidenote{\passref{4.8.50}}\skt{चित्रकार्क्कौ }\edtext{\skt{त्रृवृत्पाठा}}{\lemma{\skt{त्रृवृत्पाठा}}\Afootnote{H; \skt{त्रिवृत्पाठे}\ A}}\skt{ }\edtext{\skt{मलयं}}{\lemma{\skt{मलयं}}\Afootnote{H; \skt{मलपूं}\ A}}\skt{ हयमारकं~। }\pend
\pstart\versepar \skt{सुधां वचां लाङ्गलकीं }\edtext{\skt{हरितालं}}{\lemma{\skt{हरितालं}}\Afootnote{H; \skt{सप्तपर्णं}\ A}}\skt{ सुवर्च्चिकां~॥}\pend

\pstart \ledsidenote{\passref{4.8.51ab}}\skt{ज्योतिष्मतीञ्च सम्भृत्य }\edtext{\edtext{\skt{तैल\-न्नि\-ष्प\-न्द\-न\-म्प\-चेत्}}{\lemma{\skt{निष्पन्दनम्}}\Afootnote{H; \skt{धीरो विपा}\ A}}}{\lemma{\skt{पचेत्}}\Afootnote{H; \skt{चयेत्}\ A}}\skt{~।}\pend

\pstart \ledsidenote{\passref{4.8.51.add}}\skt{लाक्षा\-दि\-क\-प्र\-ती\-वा\-प\-मे\-त\-द्वि\-द्या\-द्भ\-ग\-न्दरे~॥}\pend

\pstart\versepar \ledsidenote{\passref{4.8.52}}\skt{शोधनं रोपणञ्चैव }\edtext{\skt{सर्व\-ण्णी\-क\-र\-ण\-न्तथा}}{\lemma{\skt{सर्व\appbrk{}ण्णी\appbrk{}क\appbrk{}र\appbrk{}णन्}}\Afootnote{H; \skt{सर्वणकरणं}\ A}}\skt{~। }\pend
\pstart\versepar \skt{}\edtext{\edtext{\skt{द्विव्र\-णी\-य\-म\-व\-क्ष्येत}}{\lemma{\skt{द्विव्रणीयम्}}\Afootnote{H; \skt{द्विव्रणीय}\ A}}}{\lemma{\skt{अवक्ष्येत}}\Afootnote{H; \skt{मवेक्षेत}\ A}}\skt{ व्रणावस्थासु बुद्धिमान्~॥}\pend

\pstart\versepar \ledsidenote{\passref{4.8.53}}\skt{छिद्रादूर्द्ध्वं }\edtext{\edtext{\skt{हरे\-दो\-ष्ठ\-म\-र्शो\-य\-न्त्रस्य}}{\lemma{\skt{हरेदोष्ठम्}}\Afootnote{H; \skt{हरेदोष्ठ}\ A}}}{\lemma{\skt{अर्शोयन्त्रस्य}}\Afootnote{H; \skt{मर्शोयन्त्रस्य}\ A}}\skt{ }\edtext{\skt{तत्ववित्}}{\lemma{\skt{तत्ववित्}}\Afootnote{H; \skt{यन्त्रवित्}\ A}}\skt{~। }\pend
\pstart\versepar \skt{ततो भगन्दरे }\edtext{\skt{दद्या\-दे\-त\-द\-र्द्धे\-न्दु\-स\-न्निभं}}{\lemma{\skt{दद्या\appbrk{}दे\appbrk{}त\appbrk{}द\appbrk{}र्द्धे\appbrk{}न्दु\appbrk{}स\appbrk{}न्निभं}}\Afootnote{H; \skt{दद्या\appbrk{}दे\appbrk{}त\appbrk{}द\appbrk{}र्धे\appbrk{}न्दु\appbrk{}स\appbrk{}न्नि\appbrk{}भम्}\ A}}\skt{~॥}\pend

\pstart\versepar \ledsidenote{\passref{4.8.54}}\skt{व्यायामम्मैथुनं }\edtext{\skt{युद्धं}}{\lemma{\skt{युद्धं}}\Afootnote{H; \skt{कोपं}\ A}}\skt{ पृष्ठ\-या\-न\-ङ्गु\-रूणि च~। }\pend
\pstart\versepar \skt{संव\-त्स\-र\-म्प\-रि\-ह\-रे\-दु\-प\-रू\-ढ\-व्रणो नरः~॥}\pend

\pstart \ledsidenote{\passref{4.8.trailer}}\edtext{\skt{इति~॥ चिकि}}{\lemma{\skt{इति~॥ चिकि}}\Afootnote{H; \skt{इति सुश्रु\appbrk{}त\appbrk{}सं\appbrk{}हि\appbrk{}तायां चिकित्सास्थाने भग\appbrk{}न्द\appbrk{}र\appbrk{}चि\appbrk{}कि\appbrk{}त्सितं नामाष्टमोऽध्यायः}\ A}}\skt{~॥}\pend

\endnumbering
\beginnumbering
\eledsection*{Adhyāya 9}
\pstart \ledsidenote{\passref{4.9.1}}\skt{अथातः कुष्ठ\-चि\-कि\-त्सि\-त\-म्व्या\-ख्या\-स्यामः~॥}\pend

\pstart \ledsidenote{\passref{4.9.4}}\edtext{\skt{त्वग्दोषी तु}}{\lemma{\skt{त्वग्दोषी तु}}\Afootnote{H; \skt{तत्र त्वग्दोषी}\ A}}\skt{ }\edtext{\skt{मान्स\-दु\-ग्ध\-द\-धि\-ति\-ल\-कु\-ल\-त्थ\-मा\-ष\-नि\-ष्पा\-वे\-क्षु\-वि\-का\-रा\-म्ल\-पि\-ष्ट\-म\-या\-नि\-वि\-दा\-ह्या\-भि\-ष्य\-न्दीनि}}{\lemma{\skt{मान्स\appbrk{}दु\appbrk{}ग्ध\appbrk{}द\appbrk{}धि\appbrk{}ति\appbrk{}ल\appbrk{}कु\appbrk{}ल\appbrk{}त्थ\appbrk{}मा\appbrk{}ष\appbrk{}नि\appbrk{}ष्पा\appbrk{}वे\appbrk{}क्षु\appbrk{}वि\appbrk{}का\appbrk{}रा\appbrk{}म्ल\appbrk{}पि\appbrk{}ष्ट\appbrk{}म\appbrk{}या\appbrk{}नि\appbrk{}वि\appbrk{}दा\appbrk{}ह्या\appbrk{}भि\appbrk{}ष्य\appbrk{}न्दीनि}}\Afootnote{H; \skt{मांस\appbrk{}व\appbrk{}सा\appbrk{}दु\appbrk{}ग्ध\appbrk{}द\appbrk{}धि\appbrk{}तै\appbrk{}ल\appbrk{}कु\appbrk{}ल\appbrk{}त्थ\appbrk{}मा\appbrk{}ष\appbrk{}नि\appbrk{}ष्पा\appbrk{}वे\appbrk{}क्षु\appbrk{}पि\appbrk{}ष्ट\appbrk{}वि\appbrk{}का\appbrk{}रा\appbrk{}म्ल\appbrk{}वि\appbrk{}रु\appbrk{}द्धा\appbrk{}ध्य\appbrk{}श\appbrk{}ना\appbrk{}जी\appbrk{}र्ण\appbrk{}वि\appbrk{}दा\appbrk{}ह्य\appbrk{}भि\appbrk{}ष्य\appbrk{}न्दीनि}\ A}}\skt{ }\edtext{\skt{दिवा\-स्व\-प्न\-व्य\-वा\-यञ्च}}{\lemma{\skt{दिवा\appbrk{}स्व\appbrk{}प्न\appbrk{}व्य\appbrk{}वा\appbrk{}यञ्च}}\Afootnote{\skt{दिवास्वप्नं व्यवायं च}\ A; \skt{दिवा\appbrk{}स्व\appbrk{}प्न\appbrk{}व्य\appbrk{}वा\appbrk{}अ\appbrk{}यञ्च}\ H}}\skt{ }\edtext{\skt{दूरत}}{\lemma{\skt{दूरत}}\Afootnote{H; \emph{om.}\ A}}\skt{ }\edtext{\skt{एव}}{\lemma{\skt{एव}}\Afootnote{H; \emph{om.}\ A}}\skt{ परिहरेत्~।}\pend

\pstart \ledsidenote{\passref{4.9.5}}\skt{ततः }\edtext{\skt{शालि\-ष\-ष्टि\-क\-य\-व\-गो\-धू\-म\-श्या\-मा\-क\-को\-र\-दू\-षो\-द्दा\-ल\-का\-न\-न\-वा\-म्भु\-ञ्जीत}}{\lemma{\skt{शालि\appbrk{}ष\appbrk{}ष्टि\appbrk{}क\appbrk{}य\appbrk{}व\appbrk{}गो\appbrk{}धू\appbrk{}म\appbrk{}श्या\appbrk{}मा\appbrk{}क\appbrk{}को\appbrk{}र\appbrk{}दू\appbrk{}षो\appbrk{}द्दा\appbrk{}ल\appbrk{}का\appbrk{}न\appbrk{}न\appbrk{}वा\appbrk{}म्भु\appbrk{}ञ्जीत}}\Afootnote{H; \skt{शालि\appbrk{}ष\appbrk{}ष्टि\appbrk{}क\appbrk{}य\appbrk{}व\appbrk{}गो\appbrk{}धू\appbrk{}म\appbrk{}को\appbrk{}र\appbrk{}दू\appbrk{}ष\appbrk{}श्या\appbrk{}मा\appbrk{}को\appbrk{}द्दा\appbrk{}ल\appbrk{}का\appbrk{}दी\appbrk{}न\appbrk{}न\appbrk{}वा\appbrk{}न्भु\appbrk{}ञ्जीत}\ A}}\skt{~। मुद्गा\-ढ\-क्यो\-र\-न्य\-त\-रस्य }\edtext{\skt{सूपेन}}{\lemma{\skt{सूपेन}}\Afootnote{\skt{यूषेण सूपेन}\ A; \skt{सूपेन}\ H}}\skt{ }\edtext{\skt{यूषेण वा}}{\lemma{\skt{यूषेण वा}}\Afootnote{H; \skt{वा}\ A}}\skt{~॥ }\edtext{\skt{नम्ब\-प\-त्त्रा\-रु\-ष्क\-र\-वि\-या\-मि\-श्रेण}}{\lemma{\skt{नम्ब\appbrk{}प\appbrk{}त्त्रा\appbrk{}रु\appbrk{}ष्क\appbrk{}र\appbrk{}वि\appbrk{}या\appbrk{}मि\appbrk{}श्रेण}}\Afootnote{H; \skt{निम्ब\appbrk{}प\appbrk{}त्रा\appbrk{}रु\appbrk{}ष्क\appbrk{}र\appbrk{}व्या\appbrk{}मि\appbrk{}श्रेण}\ A}}\skt{ }\edtext{\skt{मण्डू\-क\-प\-ण्त्त्य\-व\-ल्गु\-जा\-ट\-रू\-ष\-क\-पुष्पैः}}{\lemma{\skt{मण्डू\appbrk{}क\appbrk{}प\appbrk{}ण्त्त्य\appbrk{}व\appbrk{}ल्गु\appbrk{}जा\appbrk{}ट\appbrk{}रू\appbrk{}ष\appbrk{}क\appbrk{}पुष्पैः}}\Afootnote{\skt{मण्डू\appbrk{}क\appbrk{}प\appbrk{}र्ण्य\appbrk{}व\appbrk{}ल्गु\appbrk{}जा\appbrk{}ट\appbrk{}रू\appbrk{}पि\appbrk{}का\appbrk{}पुष्पैः}\ A; \skt{मण्डू\appbrk{}क\appbrk{}प\appbrk{}ण्त्त्य\appbrk{}व\appbrk{}ल्गु\appbrk{}जा\appbrk{}ट\appbrk{}रू\appbrk{}ष\appbrk{}क\appbrk{}पुष्पैः ः}\ H}}\skt{ सर्प्पिःसिद्धैः सर्ष\-प\-तै\-ल\-सि\-द्धैर्व्वा~। तिक्तवर्ग्गेण }\edtext{\skt{वाभिहितेन}}{\lemma{\skt{वाभिहितेन}}\Afootnote{H; \skt{वाऽभिहितेन}\ A}}\skt{ मांससात्म्याय वा }\edtext{\skt{जाङ्ग\-ल\-मा\-न्स\-स\-मे\-द\-ष्क\-म्वि\-त\-रेत्}}{\lemma{\skt{जाङ्ग\appbrk{}ल\appbrk{}मा\appbrk{}न्स\appbrk{}स\appbrk{}मे\appbrk{}द\appbrk{}ष्कम्}}\Afootnote{H; \skt{जाङ्ग\appbrk{}ल\appbrk{}मां\appbrk{}स\appbrk{}म\appbrk{}मे\appbrk{}दस्कं}\ A}}\skt{~। तैल\-म्व\-ज्र\-क\-म\-भ्य\-ङ्गार्थे आर\-ग्व\-धा\-दि\-क\-षा\-य\-मु\-त्सा\-द\-नार्थे~। }\edtext{\skt{पान\-प\-रि\-षे\-का\-व\-गा\-हेषु}}{\lemma{\skt{पान\appbrk{}प\appbrk{}रि\appbrk{}षे\appbrk{}का\appbrk{}व\appbrk{}गा\appbrk{}हेषु}}\Afootnote{H; \skt{पान\appbrk{}प\appbrk{}रि\appbrk{}षे\appbrk{}का\appbrk{}व\appbrk{}गा\appbrk{}हा\appbrk{}दिषु च}\ A}}\skt{ खदि\-र\-क\-षा\-य\-मि\-त्येष }\edtext{\skt{चाहारविभागः}}{\lemma{\skt{चाहारविभागः}}\Afootnote{\skt{आहा\appbrk{}रा\appbrk{}चा\appbrk{}र\appbrk{}वि\appbrk{}भागः}\ A; \skt{विभागः}\ H}}\skt{~।}\pend

\pstart \ledsidenote{\passref{4.9.6}}\skt{तत्र पूर्व्व\-रू\-पे\-षू\-भ\-यतः~। }\edtext{\skt{संशो\-ध\-न\-मा\-से\-वेत}}{\lemma{\skt{संशो\appbrk{}ध\appbrk{}न\appbrk{}मा\appbrk{}से\appbrk{}वेत}}\Afootnote{\skt{संशो\appbrk{}ध\appbrk{}न\appbrk{}मा\appbrk{}से\appbrk{}वेत~। त}\ A; \skt{संशो\appbrk{}ध\appbrk{}न\appbrk{}मा\appbrk{}से\appbrk{}वेत}\ H}}\skt{~। }\edtext{\skt{त्वक्प्राप्ते}}{\lemma{\skt{त्वक्प्राप्ते}}\Afootnote{H; \skt{त्र त्वक्संप्राप्ते}\ A}}\skt{ }\edtext{\skt{संशो\-ध\-ना\-ले\-प\-नानि}}{\lemma{\skt{संशो\appbrk{}ध\appbrk{}ना\appbrk{}ले\appbrk{}प\appbrk{}नानि}}\Afootnote{H; \skt{शोध\appbrk{}ना\appbrk{}ले\appbrk{}प\appbrk{}नानि}\ A}}\skt{~। शोणितप्राप्ते संशो\-ध\-ना\-ले\-प\-न\-क\-षा\-य\-पा।\-न\-शो\-णि\-ता\-व\-से\-च\-नानि~। मांसप्राप्ते }\edtext{\skt{संशो\-ध\-ना\-ले\-प\-न\-क\-षा\-य\-पा\-न\-शो\-णि\-ता\-व\-से\-का\-रि\-ष्ट\-म\-न्थ\-प्रा\-सान्}}{\lemma{\skt{संशो\appbrk{}ध\appbrk{}ना\appbrk{}ले\appbrk{}प\appbrk{}न\appbrk{}क\appbrk{}षा\appbrk{}य\appbrk{}पा\appbrk{}न\appbrk{}शो\appbrk{}णि\appbrk{}ता\appbrk{}व\appbrk{}से\appbrk{}का\appbrk{}रि\appbrk{}ष्ट\appbrk{}म\appbrk{}न्थ\appbrk{}प्रा\appbrk{}सान्}}\Afootnote{H; \skt{शोध\appbrk{}ना\appbrk{}ले\appbrk{}प\appbrk{}न\appbrk{}क\appbrk{}षा\appbrk{}य\appbrk{}पा\appbrk{}न\appbrk{}शो\appbrk{}णि\appbrk{}ता\appbrk{}व\appbrk{}से\appbrk{}च\appbrk{}ना\appbrk{}रि\appbrk{}ष्ट\appbrk{}म\appbrk{}न्थ\appbrk{}प्राशाः}\ A}}\skt{~। }\edtext{\skt{चतुर्थे}}{\lemma{\skt{चतुर्थे}}\Afootnote{H; \skt{चतुर्थ}\ A}}\skt{ }\edtext{\skt{धातुप्राप्ते}}{\lemma{\skt{धातुप्राप्ते}}\Afootnote{H; \skt{कर्मगुणप्राप्तं}\ A}}\skt{ }\edtext{\skt{याप्येनात्मवतः}}{\lemma{\skt{याप्येनात्मवतः}}\Afootnote{H; \skt{याप्यमात्मवतः}\ A}}\skt{ सम्वि\-धा\-न\-व\-तश्च~। तत्र }\edtext{\edtext{\edtext{\skt{संशो\-ध\-ना\-च्छो\-णि\-ता\-व\-से\-का\-च्चो\-र्द्ध्व\-म्भ\-ल्ला\-त\-क\-शि\-ला\-ज\-तु\-कु\-रु\-व\-क\-ख\-दि\-रा\-स\-ना\-य\-स्कृ\-ति\-वि\-धा\-न\-मु\-प\-से\-वेत}}{\lemma{\skt{छोणि\appbrk{}ता\appbrk{}व\appbrk{}से\appbrk{}काच्}}\Afootnote{H; \skt{छोणि\appbrk{}ता\appbrk{}व\appbrk{}से\appbrk{}च\appbrk{}नाच्}\ A}}}{\lemma{\skt{भल्ला\appbrk{}त\appbrk{}क\appbrk{}शि\appbrk{}ला\appbrk{}ज\appbrk{}तु\appbrk{}कु\appbrk{}रु\appbrk{}व\appbrk{}क\appbrk{}ख\appbrk{}दि\appbrk{}रा\appbrk{}स\appbrk{}ना\appbrk{}य\appbrk{}स्कृ\appbrk{}ति\appbrk{}वि\appbrk{}धा\appbrk{}नम्}}\Afootnote{H; \skt{भल्ला\appbrk{}त\appbrk{}शि\appbrk{}ला\appbrk{}ज\appbrk{}तु\appbrk{}धा\appbrk{}तु\appbrk{}मा\appbrk{}क्षी\appbrk{}क\appbrk{}गु\appbrk{}ग्गु\appbrk{}ल्व\appbrk{}गु\appbrk{}रु\appbrk{}तु\appbrk{}व\appbrk{}र\appbrk{}क\appbrk{}ख\appbrk{}दि\appbrk{}रा\appbrk{}स\appbrk{}ना\appbrk{}य\appbrk{}स्कृ\appbrk{}ति\appbrk{}वि\appbrk{}धा\appbrk{}नम्}\ A}}}{\lemma{\skt{उपसेवेत}}\Afootnote{H; \skt{आसेवेत}\ A}}\skt{~। }\edtext{\skt{पञ्चमे}}{\lemma{\skt{पञ्चमे}}\Afootnote{H; \skt{पञ्चमं}\ A}}\skt{ नैवोपक्रमेत्}\pend

\pstart \ledsidenote{\passref{4.9.7}}\skt{तत्र प्रथममेव कुष्ठिनं }\edtext{\skt{स्नेह\-पा\-ने\-नो\-प\-क्र\-मेत्}}{\lemma{\skt{स्नेह\appbrk{}पा\appbrk{}ने\appbrk{}नो\appbrk{}प\appbrk{}क्र\appbrk{}मेत्}}\Afootnote{H; \skt{स्नेह\appbrk{}पा\appbrk{}न\appbrk{}वि\appbrk{}धा\appbrk{}ने\appbrk{}नो\appbrk{}प\appbrk{}पा\appbrk{}द\appbrk{}येत्}\ A}}\skt{~। }\edtext{\skt{मेष\-शृ\-ङ्गी\-श्व\-दं\-ष्ट्रा\-सा\-ङ्ग\-ष्टा\-गु\-डू\-ची\-द्वि\-प\-ञ्च\-मू\-ल\-सि\-द्ध\-न्तैलं}}{\lemma{\skt{मेष\appbrk{}शृ\appbrk{}ङ्गी\appbrk{}श्व\appbrk{}दं\appbrk{}ष्ट्रा\appbrk{}सा\appbrk{}ङ्ग\appbrk{}ष्टा\appbrk{}गु\appbrk{}डू\appbrk{}ची\appbrk{}द्वि\appbrk{}प\appbrk{}ञ्च\appbrk{}मू\appbrk{}ल\appbrk{}सि\appbrk{}द्ध\appbrk{}न्तैलं}}\Afootnote{H; \skt{मेष\appbrk{}शृ\appbrk{}ङ्गी\appbrk{}श्व\appbrk{}दं\appbrk{}ष्ट्रा\appbrk{}शा\appbrk{}र्ङ्गे\appbrk{}ष्टा\appbrk{}गु\appbrk{}दू\appbrk{}ची\appbrk{}द्वि\appbrk{}प\appbrk{}ञ्च\appbrk{}मू\appbrk{}ली\appbrk{}सिद्धं तैलं}\ A}}\skt{ घृतम्वा }\edtext{\skt{वात\-कु\-ष्ठि\-ना\-म्पा\-ना\-भ्य\-ङ्ग\-यो\-र्व्वि\-द\-ध्यात्}}{\lemma{\skt{वातकुष्ठिनाम्}}\Afootnote{H; \skt{वातकुष्टिनां}\ A}}\skt{~। }\edtext{\skt{धवा\-श्व\-क\-र्ण्ण\-क\-कु\-भ\-प\-ला\-श\-पि\-चु\-म\-र्द्द\-म\-धु\-क\-लो\-ध्र\-स\-म\-ङ्गा\-सिद्धं}}{\lemma{\skt{धवा\appbrk{}श्व\appbrk{}क\appbrk{}र्ण्ण\appbrk{}क\appbrk{}कु\appbrk{}भ\appbrk{}प\appbrk{}ला\appbrk{}श\appbrk{}पि\appbrk{}चु\appbrk{}म\appbrk{}र्द्द\appbrk{}म\appbrk{}धु\appbrk{}क\appbrk{}लो\appbrk{}ध्र\appbrk{}स\appbrk{}म\appbrk{}ङ्गा\appbrk{}सिद्धं}}\Afootnote{H; \skt{धवा\appbrk{}श्व\appbrk{}क\appbrk{}र्ण\appbrk{}क\appbrk{}कु\appbrk{}भ\appbrk{}प\appbrk{}ला\appbrk{}श\appbrk{}पि\appbrk{}चु\appbrk{}म\appbrk{}र्द\appbrk{}प\appbrk{}र्प\appbrk{}ट\appbrk{}क\appbrk{}म\appbrk{}धु\appbrk{}क\appbrk{}रो\appbrk{}ध्र\appbrk{}स\appbrk{}म\appbrk{}ङ्गा\appbrk{}सिद्धं}\ A}}\skt{ }\edtext{\skt{पित्तकुष्ठिनां}}{\lemma{\skt{पित्तकुष्ठिनां}}\Afootnote{H; \skt{सर्पिः पित्तकुष्ठिनां}\ A}}\skt{~। }\edtext{\skt{प्रिया\-ल\-स\-ना\-र\-ग्व\-ध\-नि\-म्ब\-स\-प्त\-प\-र्ण्ण\-चि\-त्र\-क\-म\-रि\-च\-व\-चा\-कु\-ष्ठ\-सिद्धं}}{\lemma{\skt{प्रिया\appbrk{}ल\appbrk{}स\appbrk{}ना\appbrk{}र\appbrk{}ग्व\appbrk{}ध\appbrk{}नि\appbrk{}म्ब\appbrk{}स\appbrk{}प्त\appbrk{}प\appbrk{}र्ण्ण\appbrk{}चि\appbrk{}त्र\appbrk{}क\appbrk{}म\appbrk{}रि\appbrk{}च\appbrk{}व\appbrk{}चा\appbrk{}कु\appbrk{}ष्ठ\appbrk{}सिद्धं}}\Afootnote{H; \skt{प्रिया\appbrk{}ल\appbrk{}शा\appbrk{}ला\appbrk{}र\appbrk{}ग्व\appbrk{}ध\appbrk{}नि\appbrk{}म्ब\appbrk{}स\appbrk{}प्त\appbrk{}प\appbrk{}र्ण\appbrk{}चि\appbrk{}त्र\appbrk{}क\appbrk{}म\appbrk{}रि\appbrk{}च\appbrk{}व\appbrk{}चा\appbrk{}कु\appbrk{}ष्ठ\appbrk{}सिद्धं}\ A}}\skt{ श्लेष्मकुष्ठिनां~। }\edtext{\skt{भल्ला\-त\-क\-वि\-ड\-ङ्ग\-सि\-द्धम्वा}}{\lemma{\skt{भल्ला\appbrk{}त\appbrk{}क\appbrk{}वि\appbrk{}ड\appbrk{}ङ्ग\appbrk{}सि\appbrk{}द्धम्वा}}\Afootnote{H; \skt{भल्ला\appbrk{}त\appbrk{}का\appbrk{}भ\appbrk{}या\appbrk{}वि\appbrk{}ड\appbrk{}ङ्ग\appbrk{}सिद्धं वा}\ A}}\skt{ सर्व्वेषां }\edtext{\edtext{\skt{हरी\-त\-व्या\-म\-ल\-क\-वि\-भी\-त\-क\-प\-वो\-ल\-पि\-चु\-म\-र्द्दा\-ट\-रू\-ष\-क\-क\-टु\-रो\-हि\-णी\-भा\-र\-द्रा\-दुं\-रा\-ल\-भा\-त्रा\-य\-मा\-णा\-प\-र्प्प\-ट\-क\-ञ्चे\-त्ये\-त\-द्द्वि\-प\-लि\-क्या\-नि\-ज\-ला\-ढके}}{\lemma{\skt{हरी\appbrk{}त\appbrk{}व्या\appbrk{}म\appbrk{}ल\appbrk{}क\appbrk{}वि\appbrk{}भी\appbrk{}त\appbrk{}क\appbrk{}प\appbrk{}वो\appbrk{}ल\appbrk{}पि\appbrk{}चु\appbrk{}म\appbrk{}र्द्दा\appbrk{}ट\appbrk{}रू\appbrk{}ष\appbrk{}क\appbrk{}क\appbrk{}टु\appbrk{}रो\appbrk{}हि\appbrk{}णी\appbrk{}भा\appbrk{}र\appbrk{}द्रा\appbrk{}दुं\appbrk{}रा\appbrk{}ल\appbrk{}भा\appbrk{}त्रा\appbrk{}य\appbrk{}मा\appbrk{}णा\appbrk{}प\appbrk{}र्प्प\appbrk{}ट\appbrk{}क\appbrk{}ञ्चे\appbrk{}त्ये\appbrk{}तद्}}\Afootnote{H; \emph{om.}\ A}}}{\lemma{\skt{द्विप\appbrk{}लि\appbrk{}क्या\appbrk{}नि\appbrk{}ज\appbrk{}ला\appbrk{}ढके}}\Afootnote{H; \emph{om.}\ A}}\skt{ }\edtext{\skt{प्रक्षिप्य}}{\lemma{\skt{प्रक्षिप्य}}\Afootnote{H; \emph{om.}\ A}}\skt{ }\edtext{\edtext{\edtext{\skt{तत्क\-षा\-य\-म्पा\-दा\-व\-शि\-ष्ट\-म\-व\-तार्य}}{\lemma{\skt{तत्कषायम्}}\Afootnote{H; \emph{om.}\ A}}}{\lemma{\skt{पादावशिष्टम्}}\Afootnote{H; \emph{om.}\ A}}}{\lemma{\skt{अवतार्य}}\Afootnote{H; \emph{om.}\ A}}\skt{ }\edtext{\edtext{\skt{घृत\-प्र\-स्थ\-मा\-व\-पे\-त्क\-ल्क\-पे\-ष्याणि}}{\lemma{\skt{घृत\appbrk{}प्र\appbrk{}स्थ\appbrk{}मा\appbrk{}व\appbrk{}पेत्}}\Afootnote{\emph{om.}\ A; \skt{घृत\appbrk{}प्र\appbrk{}स्थ\appbrk{}मा\appbrk{}व\appbrk{}पेत्}\ H}}}{\lemma{\skt{कल्कपेष्याणि}}\Afootnote{H; \emph{om.}\ A}}\skt{ }\edtext{\skt{चेमानि}}{\lemma{\skt{चेमानि}}\Afootnote{H; \emph{om.}\ A}}\skt{ }\edtext{\edtext{\skt{भेष\-जा\-न्य\-र्द्ध\-प\-लि\-क्यानि}}{\lemma{\skt{भेषजान्य्}}\Afootnote{H; \emph{om.}\ A}}}{\lemma{\skt{अर्द्धपलिक्यानि}}\Afootnote{H; \emph{om.}\ A}}\skt{~। }\edtext{\edtext{\edtext{\edtext{\edtext{\skt{त्राय\-मा\-णा\-मु\-स्ते\-न्द्र\-य\-वा\-च\-न्द\-न\-कि\-रा\-त\-ति\-क्त\-क\-पि\-प्प\-ल्य\-श्चे\-त्ये\-त\-त्मृ\-द्व\-ग्नि\-सि\-द्ध\-म\-व\-तार्य}}{\lemma{\skt{त्राय\appbrk{}मा\appbrk{}णा\appbrk{}मु\appbrk{}स्ते\appbrk{}न्द्र\appbrk{}य\appbrk{}वा\appbrk{}च\appbrk{}न्द\appbrk{}न\appbrk{}कि\appbrk{}रा\appbrk{}त\appbrk{}ति\appbrk{}क्त\appbrk{}क\appbrk{}पि\appbrk{}प्प\appbrk{}ल्यश्}}\Afootnote{H; \emph{om.}\ A}}}{\lemma{\skt{चेत्य्}}\Afootnote{H; \emph{om.}\ A}}}{\lemma{\skt{एतत्}}\Afootnote{H; \emph{om.}\ A}}}{\lemma{\skt{मृद्वग्निसिद्धम्}}\Afootnote{H; \emph{om.}\ A}}}{\lemma{\skt{अवतार्य}}\Afootnote{H; \emph{om.}\ A}}\skt{ }\edtext{\edtext{\edtext{\skt{पाय\-ये\-दे\-त\-त्ति\-क्त\-क\-न्नाम}}{\lemma{\skt{पाययेदेतत्}}\Afootnote{H; \emph{om.}\ A}}}{\lemma{\skt{तिक्तकन्}}\Afootnote{H; \emph{om.}\ A}}}{\lemma{\skt{नाम}}\Afootnote{H; \emph{om.}\ A}}\skt{ }\edtext{\skt{सर्प्पिः}}{\lemma{\skt{सर्प्पिः}}\Afootnote{H; \emph{om.}\ A}}\skt{ }\edtext{\edtext{\skt{कुष्ठ\-ज्व\-र\-गु\-ल्मा\-र्शो\-ग्र\-ह\-णी\-दो\-ष\-पा\-ण्डु\-रो\-ग\-वि\-स\-र्प्प\-षा\-ण्ढ्य\-श\-म\-न\-श्चेति}}{\lemma{\skt{कुष्ठ\appbrk{}ज्व\appbrk{}र\appbrk{}गु\appbrk{}ल्मा\appbrk{}र्शो\appbrk{}ग्र\appbrk{}ह\appbrk{}णी\appbrk{}दो\appbrk{}ष\appbrk{}पा\appbrk{}ण्डु\appbrk{}रो\appbrk{}ग\appbrk{}वि\appbrk{}स\appbrk{}र्प्प\appbrk{}षा\appbrk{}ण्ढ्य\appbrk{}श\appbrk{}म\appbrk{}नश्}}\Afootnote{H; \skt{तुवरकतैलं भल्लातकतैलं}\ A}}}{\lemma{\skt{चेति}}\Afootnote{H; \skt{वेति}\ A}}\skt{~॥}\pend

\pstart \ledsidenote{\passref{4.9.8}}\edtext{\skt{सप्त\-प\-र्ण्णा\-ति\-वि\-षा\-श\-म्या\-क\-क\-टु\-रो\-हि\-ण्य\-मृ\-णा\-ल\-ह\-री\-त\-क्या\-म\-ल\-क\-वि\-भी\-त\-क\-प\-टो\-ल\-पि\-चु\-म\-र्द्द\-प\-र्प्प\-ट\-कि\-रा\-त\-ति\-क्त\-क\-ह\-रि\-द्रा\-दु\-रा\-ल\-भा\-त्रा\-य\-मा\-णा\-मु\-स्ता\-च\-न्द\-न\-प\-द्म\-क\-ह\-रि\-द्रो\-प\-कु\-ल्या\-वि\-श्वा\-मू\-र्व्वा\-श\-ता\-व\-री\-शा\-रि\-वा\-ट\-रू\-ष\-क\-कु\-ट\-जे\-न्द्र\-वा\-रु\-णी\-पा\-त्ः\-अ\-वृ\-क्षा\-द\-नी\-ग्र\-न्थिका}}{\lemma{\skt{सप्त\appbrk{}प\appbrk{}र्ण्णा\appbrk{}ति\appbrk{}वि\appbrk{}षा\appbrk{}श\appbrk{}म्या\appbrk{}क\appbrk{}क\appbrk{}टु\appbrk{}रो\appbrk{}हि\appbrk{}ण्य\appbrk{}मृ\appbrk{}णा\appbrk{}ल\appbrk{}ह\appbrk{}री\appbrk{}त\appbrk{}क्या\appbrk{}म\appbrk{}ल\appbrk{}क\appbrk{}वि\appbrk{}भी\appbrk{}त\appbrk{}क\appbrk{}प\appbrk{}टो\appbrk{}ल\appbrk{}पि\appbrk{}चु\appbrk{}म\appbrk{}र्द्द\appbrk{}प\appbrk{}र्प्प\appbrk{}ट\appbrk{}कि\appbrk{}रा\appbrk{}त\appbrk{}ति\appbrk{}क्त\appbrk{}क\appbrk{}ह\appbrk{}रि\appbrk{}द्रा\appbrk{}दु\appbrk{}रा\appbrk{}ल\appbrk{}भा\appbrk{}त्रा\appbrk{}य\appbrk{}मा\appbrk{}णा\appbrk{}मु\appbrk{}स्ता\appbrk{}च\appbrk{}न्द\appbrk{}न\appbrk{}प\appbrk{}द्म\appbrk{}क\appbrk{}ह\appbrk{}रि\appbrk{}द्रो\appbrk{}प\appbrk{}कु\appbrk{}ल्या\appbrk{}वि\appbrk{}श्वा\appbrk{}मू\appbrk{}र्व्वा\appbrk{}श\appbrk{}ता\appbrk{}व\appbrk{}री\appbrk{}शा\appbrk{}रि\appbrk{}वा\appbrk{}ट\appbrk{}रू\appbrk{}ष\appbrk{}क\appbrk{}कु\appbrk{}ट\appbrk{}जे\appbrk{}न्द्र\appbrk{}वा\appbrk{}रु\appbrk{}णी\appbrk{}पा\appbrk{}त्ः\appbrk{}अ\appbrk{}वृ\appbrk{}क्षा\appbrk{}द\appbrk{}नी\appbrk{}ग्र\appbrk{}न्थिका}}\Afootnote{H; \skt{सप्त\appbrk{}प\appbrk{}र्णा\appbrk{}र\appbrk{}ग्व\appbrk{}धा\appbrk{}ति\appbrk{}वि\appbrk{}षे\appbrk{}क्षु\appbrk{}र\appbrk{}पा\appbrk{}ठा\appbrk{}क\appbrk{}टु\appbrk{}रो\appbrk{}हि\appbrk{}ण्य\appbrk{}मृ\appbrk{}ता\appbrk{}त्रि\appbrk{}फ\appbrk{}ला\appbrk{}प\appbrk{}टो\appbrk{}ल\appbrk{}पि\appbrk{}चु\appbrk{}म\appbrk{}र्द\appbrk{}प\appbrk{}र्प\appbrk{}ट\appbrk{}क\appbrk{}दु\appbrk{}रा\appbrk{}ल\appbrk{}भा\appbrk{}त्रा\appbrk{}य\appbrk{}मा\appbrk{}णा\appbrk{}मु\appbrk{}स्ता\appbrk{}च\appbrk{}न्द\appbrk{}न\appbrk{}प\appbrk{}द्म\appbrk{}क\appbrk{}ह\appbrk{}रि\appbrk{}द्रो\appbrk{}प\appbrk{}कु\appbrk{}ल्या\appbrk{}वि\appbrk{}शा\appbrk{}ला\appbrk{}मू\appbrk{}र्वा\appbrk{}श\appbrk{}ता\appbrk{}व\appbrk{}री\appbrk{}सा\appbrk{}रि\appbrk{}वे\appbrk{}न्द्र\appbrk{}य\appbrk{}वा\appbrk{}ट\appbrk{}रू\appbrk{}ष\appbrk{}क\appbrk{}ष\appbrk{}ट्ग्र\appbrk{}न्था\appbrk{}म\appbrk{}धु\appbrk{}क\appbrk{}भू\appbrk{}नि\appbrk{}म्ब\appbrk{}गृ\appbrk{}ष्टिका}\ A}}\skt{ }\edtext{\skt{चेति}}{\lemma{\skt{चेति}}\Afootnote{H; \skt{इति}\ A}}\skt{~। }\edtext{\skt{समभागानि}}{\lemma{\skt{समभागानि}}\Afootnote{H; \skt{समभागाः}\ A}}\skt{ कल्कः }\edtext{\skt{स्याच्चतुर्ग्गुणं}}{\lemma{\skt{स्याच्}}\Afootnote{H; \skt{स्यात्, कल्काच्}\ A}}\skt{ सर्प्पिः प्रक्षिप्य }\edtext{\edtext{\edtext{\edtext{\edtext{\edtext{\skt{तद्द्वि\-गु\-ण\-मा\-म\-कं\-ल\-र\-स\-स्त\-च्च\-तु\-र्ग्गु\-ण\-मा\-प\-स्त\-दै\-कध्यं}}{\lemma{\skt{द्विगुणम्}}\Afootnote{H; \skt{द्विगुणो}\ A}}}{\lemma{\skt{आम\appbrk{}कं\appbrk{}ल\appbrk{}र\appbrk{}सस्}}\Afootnote{H; \skt{धात्री\appbrk{}फ\appbrk{}ल\appbrk{}र\appbrk{}सस्}\ A}}}{\lemma{\skt{तच्}}\Afootnote{H; \skt{तच्}\ A}}}{\lemma{\skt{चतुर्ग्गुणम्}}\Afootnote{H; \skt{चतुर्गुणा}\ A}}}{\lemma{\skt{आपस्}}\Afootnote{H; \skt{आपा}\ A}}}{\lemma{\skt{तदैकध्यं}}\Afootnote{H; \skt{स्तदैकध्यं}\ A}}\skt{ }\edtext{\edtext{\skt{मृद्व\-ग्नि\-सि\-द्ध\-म\-व\-तार्य}}{\lemma{\skt{मृद्वग्निसिद्धम्}}\Afootnote{H; \emph{om.}\ A}}}{\lemma{\skt{अवतार्य}}\Afootnote{H; \skt{समालोड्य}\ A}}\skt{ }\edtext{\skt{पाय\-ये\-दे\-त\-न्म\-हा\-ति\-क्तकं}}{\lemma{\skt{पाययेदेतन्}}\Afootnote{H; \skt{विपचेत्, एतन्}\ A}}\skt{ नाम सर्प्पिः }\edtext{\skt{कुष्ठा\-र्शो\-वा\-त\-पि\-त्त\-वि\-ष\-म\-ज्व\-र\-हृ\-द्रो\-गो\-न्मा\-द\-ग\-ल\-ग\-ण्ड\-ग\-ण्ड\-मा\-ला\-स\-श्ली\-प\-द\-का\-म\-ला\-पा\-ण्डु\-रो\-ग\-षा\-ण्ढ्य\-क\-ण्डू\-पा\-मा\-दीश्च}}{\lemma{\skt{कुष्ठा\appbrk{}र्शो\appbrk{}वा\appbrk{}त\appbrk{}पि\appbrk{}त्त\appbrk{}वि\appbrk{}ष\appbrk{}म\appbrk{}ज्व\appbrk{}र\appbrk{}हृ\appbrk{}द्रो\appbrk{}गो\appbrk{}न्मा\appbrk{}द\appbrk{}ग\appbrk{}ल\appbrk{}ग\appbrk{}ण्ड\appbrk{}ग\appbrk{}ण्ड\appbrk{}मा\appbrk{}ला\appbrk{}स\appbrk{}श्ली\appbrk{}प\appbrk{}द\appbrk{}का\appbrk{}म\appbrk{}ला\appbrk{}पा\appbrk{}ण्डु\appbrk{}रो\appbrk{}ग\appbrk{}षा\appbrk{}ण्ढ्य\appbrk{}क\appbrk{}ण्डू\appbrk{}पा\appbrk{}मा\appbrk{}दीश्च}}\Afootnote{H; \skt{कुष्ठ\appbrk{}वि\appbrk{}ष\appbrk{}म\appbrk{}ज्व\appbrk{}र\appbrk{}र\appbrk{}क्त\appbrk{}पि\appbrk{}त्त\appbrk{}हृ\appbrk{}द्रो\appbrk{}गो\appbrk{}न्मा\appbrk{}दा\appbrk{}प\appbrk{}स्मा\appbrk{}र\appbrk{}गु\appbrk{}ल्म\appbrk{}पि\appbrk{}डा\appbrk{}का\appbrk{}सृ\appbrk{}ग्द\appbrk{}र\appbrk{}ग\appbrk{}ल\appbrk{}ग\appbrk{}ण्ड\appbrk{}ग\appbrk{}ण्ड\appbrk{}मा\appbrk{}ला\appbrk{}श्ली\appbrk{}प\appbrk{}द\appbrk{}पा\appbrk{}ण्डु\appbrk{}रो\appbrk{}ग\appbrk{}वि\appbrk{}स\appbrk{}र्पा\appbrk{}र्शः\appbrk{}षा\appbrk{}ण्ढ्य\appbrk{}क\appbrk{}ण्डू\appbrk{}पा\appbrk{}मा\appbrk{}दीञ्}\ A}}\skt{ }\edtext{\skt{शमयेदित्य्}}{\lemma{\skt{शमयेदित्य्}}\Afootnote{H; \skt{छमयेदिति}\ A}}\pend

\pstart \ledsidenote{\passref{4.9.10}}\edtext{\skt{अन्यतमेन}}{\lemma{\skt{अन्यतमेन}}\Afootnote{H; \skt{अतोऽन्यतमेन}\ A}}\skt{ }\edtext{\skt{घृते\-नो\-प\-स्नि\-ग्ध\-स्यैकां}}{\lemma{\skt{घृते\appbrk{}नो\appbrk{}प\appbrk{}स्नि\appbrk{}ग्ध\appbrk{}स्यैकां}}\Afootnote{H; \skt{घृतेन स्निग्धस्विन्नस्यैकां}\ A}}\skt{ द्वे तिस्रश्चतस्रः पञ्च वा }\edtext{\skt{सिराम्विध्येत्}}{\lemma{\skt{सिराम्}}\Afootnote{H; \skt{सिरा}\ A}}\skt{~। मण्डलानि }\edtext{\skt{चोत्सन्नान्येव}}{\lemma{\skt{एव}}\Afootnote{H; \skt{अव}\ A}}\skt{ }\edtext{\skt{लिखेत्~। अभीक्ष्णं}}{\lemma{\skt{लिखेत्~। अभीक्ष्णं}}\Afootnote{H; \skt{लिखेदभीष्णं}\ A}}\skt{ }\edtext{\skt{प्रच्छेद्वा}}{\lemma{\skt{प्रच्छेद्वा}}\Afootnote{H; \skt{प्रच्छयेद्वा}\ A}}\skt{~। }\edtext{\skt{समुद्रफेनेन}}{\lemma{\skt{समुद्रफेनेन}}\Afootnote{H; \skt{समु\appbrk{}द्र\appbrk{}फे\appbrk{}न\appbrk{}शा\appbrk{}क\appbrk{}गो\appbrk{}जी\appbrk{}का\appbrk{}को\appbrk{}दु\appbrk{}म्ब\appbrk{}रि\appbrk{}काप}\ A}}\skt{ }\edtext{\skt{चाव\-घृ\-ष्या\-व\-ले\-प\-येत्}}{\lemma{\skt{चाव\appbrk{}घृ\appbrk{}ष्या\appbrk{}व\appbrk{}ले\appbrk{}प\appbrk{}येत्}}\Afootnote{H; \skt{त्रैर्वा+\appbrk{}अ\appbrk{}व\appbrk{}घृ\appbrk{}ष्या\appbrk{}ले\appbrk{}प\appbrk{}येल्}\ A}}\skt{~। }\edtext{\skt{लाक्षा\-स\-र्ज्ज\-र\-सा\-ञ्ज\-न\-प्र\-पु\-ना\-ट\-ते\-जो\-व\-त्या\-श्व\-मा\-र\-का\-र्क्क\-कु\-ट\-ज\-रे\-व\-त\-मू\-ल\-क\-ल्कै\-र्म्मू\-त्र\-पिष्टैः}}{\lemma{\skt{लाक्षा\appbrk{}स\appbrk{}र्ज्ज\appbrk{}र\appbrk{}सा\appbrk{}ञ्ज\appbrk{}न\appbrk{}प्र\appbrk{}पु\appbrk{}ना\appbrk{}ट\appbrk{}ते\appbrk{}जो\appbrk{}व\appbrk{}त्या\appbrk{}श्व\appbrk{}मा\appbrk{}र\appbrk{}का\appbrk{}र्क्क\appbrk{}कु\appbrk{}ट\appbrk{}ज\appbrk{}रे\appbrk{}व\appbrk{}त\appbrk{}मू\appbrk{}ल\appbrk{}क\appbrk{}ल्कै\appbrk{}र्म्मू\appbrk{}त्र\appbrk{}पिष्टैः}}\Afootnote{H; \skt{लाक्षा\appbrk{}स\appbrk{}र्ज\appbrk{}र\appbrk{}स\appbrk{}र\appbrk{}सा\appbrk{}ञ्ज\appbrk{}न\appbrk{}प्र\appbrk{}पु\appbrk{}न्ना\appbrk{}डा\appbrk{}व\appbrk{}ल्गु\appbrk{}ज\appbrk{}ते\appbrk{}जो\appbrk{}व\appbrk{}त्य\appbrk{}श्व\appbrk{}मा\appbrk{}र\appbrk{}का\appbrk{}र्क\appbrk{}कु\appbrk{}ट\appbrk{}जा\appbrk{}रे\appbrk{}व\appbrk{}त\appbrk{}मू\appbrk{}ल\appbrk{}क\appbrk{}ल्कै\appbrk{}र्मू\appbrk{}त्र\appbrk{}पिष्टैः}\ A}}\skt{ पित्तपिष्टैर्व्वा~॥ }\edtext{\skt{स्वर्ज्जि\-का\-तु\-त्थ\-का\-सी\-स\-वि\-ड\-ङ्गा\-गा\-र\-धू\-म\-चि\-त्र\-क\-त्रि\-क\-टु\-क\-सु\-धा\-ह\-रि\-द्रा\-सै\-न्ध\-वैर्व्वा~॥ एतैरेव}}{\lemma{\skt{स्वर्ज्जि\appbrk{}का\appbrk{}तु\appbrk{}त्थ\appbrk{}का\appbrk{}सी\appbrk{}स\appbrk{}वि\appbrk{}ड\appbrk{}ङ्गा\appbrk{}गा\appbrk{}र\appbrk{}धू\appbrk{}म\appbrk{}चि\appbrk{}त्र\appbrk{}क\appbrk{}त्रि\appbrk{}क\appbrk{}टु\appbrk{}क\appbrk{}सु\appbrk{}धा\appbrk{}ह\appbrk{}रि\appbrk{}द्रा\appbrk{}सै\appbrk{}न्ध\appbrk{}वैर्व्वा~॥ एतैर्}}\Afootnote{H; \skt{स्वर्जि\appbrk{}का\appbrk{}तु\appbrk{}त्थ\appbrk{}का\appbrk{}सी\appbrk{}स\appbrk{}वि\appbrk{}ड\appbrk{}ङ्गा\appbrk{}गा\appbrk{}र\appbrk{}धू\appbrk{}म\appbrk{}चि\appbrk{}त्र\appbrk{}क\appbrk{}क\appbrk{}टु\appbrk{}क\appbrk{}सु\appbrk{}धा\appbrk{}ह\appbrk{}रि\appbrk{}द्रा\appbrk{}सै\appbrk{}न्ध\appbrk{}व\appbrk{}क\appbrk{}ल्कैर्वा, एतान्येवावाप्य क्षारकल्पेन निः}\ A}}\skt{ च }\edtext{\skt{क्षारविपक्वैः}}{\lemma{\skt{क्षारविपक्वैः}}\Afootnote{H; \skt{क्षारे ततो विपाच्य}\ A}}\skt{ }\edtext{\skt{फाणितमेव}}{\lemma{\skt{फाणितमेव}}\Afootnote{\skt{फाणीतमिव}\ A; \skt{फाणितमेव}\ H}}\skt{ }\edtext{\skt{सञ्जातरसैः}}{\lemma{\skt{सञ्जातरसैः}}\Afootnote{H; \skt{संजा\appbrk{}त\appbrk{}म\appbrk{}व\appbrk{}तार्य लेपयेत्, ज्योति\appbrk{}ष्क\appbrk{}फ\appbrk{}ल\appbrk{}ला\appbrk{}क्षा\appbrk{}म\appbrk{}रि\appbrk{}च\appbrk{}प्पि\appbrk{}ली\appbrk{}सु\appbrk{}म\appbrk{}नः\appbrk{}प\appbrk{}त्रैर्वा}\ A}}\skt{~।}\edtext{}{\lemma{\skt{एव च}}\Afootnote{H; \skt{स्रुते पालाशे}\ A}}\pend

\pstart \ledsidenote{\passref{4.9.10.2}}\edtext{\skt{स्वर्ज्जि\-का\-तु\-त्थ\-कु\-ट\-ज\-वि\-ड\-ङ्ग\-म\-रि\-च\-लो\-ध्र\-म\-नः\-शि\-ला\-चू\-र्ण्ण\-यु\-क्तैर्व्वा}}{\lemma{\skt{स्वर्ज्जि\appbrk{}का\appbrk{}तु\appbrk{}त्थ\appbrk{}कु\appbrk{}ट\appbrk{}ज\appbrk{}वि\appbrk{}ड\appbrk{}ङ्ग\appbrk{}म\appbrk{}रि\appbrk{}च\appbrk{}लो\appbrk{}ध्र\appbrk{}म\appbrk{}नः\appbrk{}शि\appbrk{}ला\appbrk{}चू\appbrk{}र्ण्ण\appbrk{}यु\appbrk{}क्तैर्व्वा}}\Afootnote{H; \skt{स्वर्जि\appbrk{}का\appbrk{}कु\appbrk{}ष्ठ\appbrk{}तु\appbrk{}त्थ\appbrk{}कु\appbrk{}ट\appbrk{}ज\appbrk{}चि\appbrk{}त्र\appbrk{}क\appbrk{}वि\appbrk{}ड\appbrk{}ङ्ग\appbrk{}म\appbrk{}रि\appbrk{}च\appbrk{}रो\appbrk{}ध्र\appbrk{}म\appbrk{}नः\appbrk{}शि\appbrk{}ला\appbrk{}क\appbrk{}लैर्वा}\ A}}\skt{~। }\edtext{\skt{ज्योति\-ष्क\-फ\-ल\-ला\-क्षा\-म\-रि\-च\-पि\-प्प\-ली\-सु\-म\-ना\-प\-त्त्रैर्व्वा}}{\lemma{\skt{ज्योति\appbrk{}ष्क\appbrk{}फ\appbrk{}ल\appbrk{}ला\appbrk{}क्षा\appbrk{}म\appbrk{}रि\appbrk{}च\appbrk{}पि\appbrk{}प्प\appbrk{}ली\appbrk{}सु\appbrk{}म\appbrk{}ना\appbrk{}प\appbrk{}त्त्रैर्}}\Afootnote{H; \skt{हरी\appbrk{}त\appbrk{}की\appbrk{}क\appbrk{}र\appbrk{}ञ्जि\appbrk{}का\appbrk{}वि\appbrk{}ड\appbrk{}ङ्ग\appbrk{}सि\appbrk{}द्धा\appbrk{}र्थ\appbrk{}क\appbrk{}ल\appbrk{}व\appbrk{}ण\appbrk{}रो\appbrk{}च\appbrk{}ना\appbrk{}व\appbrk{}ल्गु\appbrk{}ज\appbrk{}ह\appbrk{}रि\appbrk{}द्रा\appbrk{}क\appbrk{}ल्कैर्}\ A}}\skt{~।}\pend

\pstart \ledsidenote{\passref{4.9.10.add}}\skt{हरी\-त\-की\-क\-र\-ञ्जि\-का\-वि\-ड\-ङ्ग\-सि\-द्धा\-र्थ\-क\-ल\-व\-ण\-रो\-च\-ना\-व\-ल्गु\-ज\-ह\-रि\-द्रा\-क\-ल्कैर्व्वा~।}\pend

\pstart \ledsidenote{\passref{4.9.10.1}}\edtext{\skt{मनः\-शि\-ला\-र्क्क\-क्षी\-र\-म\-रि\-च\-क\-ल्कैर्व्वा}}{\lemma{\skt{मनः\appbrk{}शि\appbrk{}ला\appbrk{}र्क्क\appbrk{}क्षी\appbrk{}र\appbrk{}म\appbrk{}रि\appbrk{}च\appbrk{}क\appbrk{}ल्कैर्}}\Afootnote{H; \skt{हरि\appbrk{}ता\appbrk{}ल\appbrk{}म\appbrk{}नः\appbrk{}शि\appbrk{}ला\appbrk{}र्क\appbrk{}क्षी\appbrk{}र\appbrk{}ति\appbrk{}ल\appbrk{}शि\appbrk{}ग्रु\appbrk{}म\appbrk{}रि\appbrk{}च\appbrk{}क\appbrk{}ल्कैर्}\ A}}\skt{~।}\pend

\pstart\versepar \ledsidenote{\passref{4.9.11}}\edtext{\skt{सर्व्वकुष्ठापहाः}}{\lemma{\skt{सर्व्वकुष्ठापहाः}}\Afootnote{H; \skt{सर्वे कुष्ठापहाः}\ A}}\skt{ सिद्धा लेपाः }\edtext{\skt{षट्प्रकीर्त्तिताः}}{\lemma{\skt{षट्प्रकीर्त्तिताः}}\Afootnote{H; \skt{सप्त प्रकीर्तिताः}\ A}}\skt{~। }\pend
\pstart\versepar \skt{}\edtext{\skt{वैशे\-षि\-का\-न\-त\-स्तू\-र्द्ध्व\-न्द\-र्द्त्रू\-श्वि\-त्रेषु}}{\lemma{\skt{दर्द्त्रूश्वित्रेषु}}\Afootnote{H; \skt{दद्रूश्वित्रेषु}\ A}}\skt{ }\edtext{\skt{वक्ष्यते}}{\lemma{\skt{वक्ष्यते}}\Afootnote{H; \skt{मे शृणु}\ A}}\skt{~।}\pend

\pstart\versepar \ledsidenote{\passref{4.9.12}}\edtext{\skt{श्रीवेष्टकुष्ठं}}{\lemma{\skt{श्रीवेष्टकुष्ठं}}\Afootnote{H\textsuperscript{pc}; \skt{लाक्षा कुष्ठं}\ A; \skt{श्रीवेष्टककुष्ठं}\ H\textsuperscript{ac}}}\skt{ }\edtext{\skt{कुटजं}}{\lemma{\skt{कुटजं}}\Afootnote{H; \emph{om.}\ A}}\skt{ }\edtext{\skt{सर्षपो}}{\lemma{\skt{सर्षपो}}\Afootnote{H; \skt{सर्षपाः}\ A}}\skt{ }\edtext{\skt{लाक्षा}}{\lemma{\skt{लाक्षा}}\Afootnote{H; \skt{श्रीनिकेतं रात्रिर्}\ A}}\skt{ व्योषं }\edtext{\skt{प्रपुनाटस्य}}{\lemma{\skt{प्रपुनाटस्य}}\Afootnote{H; \skt{चक्रमर्दस्य}\ A}}\skt{ बीजं~। }\pend
\pstart\versepar \skt{}\edtext{\skt{तक्रेण}}{\lemma{\skt{तक्रेण}}\Afootnote{H; \skt{कृत्वैकस्थं तक्र}\ A}}\skt{ पिष्टः }\edtext{\skt{सहरिद्रलेपो}}{\lemma{\skt{सहरिद्रलेपो}}\Afootnote{\skt{प्रलेपो}\ A; \skt{सहरिद्रलेपो}\ H}}\skt{ }\edtext{\skt{दद्द्रूषूक्तो}}{\lemma{\skt{दद्द्रूषूक्तो}}\Afootnote{\skt{दद्रूषूक्तो}\ A; \skt{दद्द्रूषूक्तो}\ H}}\skt{ }\edtext{\skt{मूल\-का\-बी\-ज\-युक्तः}}{\lemma{\skt{मूल\appbrk{}का\appbrk{}बी\appbrk{}ज\appbrk{}युक्तः}}\Afootnote{H; \skt{मूल\appbrk{}का\appbrk{}द्बी\appbrk{}ज\appbrk{}युक्तः}\ A}}\skt{~।}\footnoteA{Aṣṭāṅgahṛdaya Ci.19.85: \emph{lākṣā vyoṣaṃ prāpunāṭa ca bījaṃ | takronmiṣraḥ syād haridrā ca lepo dadrūṣūkto mūlakotthaṃ ca bījam ||} (Kunte 1939: 717).}\pend

\pstart\versepar \ledsidenote{\passref{4.9.13}}\edtext{\skt{ससैन्धवं प्रपुनाटस्य}}{\lemma{\skt{ससैन्धवं प्रपुनाटस्य}}\Afootnote{H; \skt{सिन्धूद्भूतं चक्रमर्दस्य}\ A}}\skt{ }\edtext{\skt{बीजं}}{\lemma{\skt{बीजं}}\Afootnote{H; \skt{बी}\ A}}\skt{ }\edtext{\skt{रसाञ्जनं}}{\lemma{\skt{रसाञ्जनं}}\Afootnote{H; \skt{जमि}\ A}}\skt{ }\edtext{\skt{सगुडं}}{\lemma{\skt{सगुडं}}\Afootnote{H; \skt{क्षूद्भूतं}\ A}}\skt{ }\edtext{\skt{केसरञ्च}}{\lemma{\skt{केसरञ्च}}\Afootnote{H; \skt{केशरं}\ A}}\skt{~। }\edtext{\skt{रसेन}}{\lemma{\skt{रसेन}}\Afootnote{H; \emph{om.}\ A}}\skt{ }\edtext{\skt{पिष्टानि}}{\lemma{\skt{पिष्टानि}}\Afootnote{H; \skt{तार्}\ A}}\skt{ }\edtext{\skt{कपित्थजेन}}{\lemma{\skt{कपित्थजेन}}\Afootnote{\skt{क्ष्यशैलम्}\ A; \skt{कपित्थजे}\ H}}\skt{~। }\pend
\pstart\versepar \skt{}\edtext{\skt{शीघ्रं}}{\lemma{\skt{शीघ्रं}}\Afootnote{\skt{पिष्टो}\ A; \emph{om.}\ H}}\skt{ }\edtext{\skt{तीव्रादर्द्द्रवो नाश इति पाठा शीघ्रं}}{\lemma{\skt{तीव्रादर्द्द्रवो नाश इति पाठा शीघ्रं}}\Afootnote{\skt{लेपोऽयं कपित्थाद्रसेन दद्रूस्}\ A; \skt{न~। शीघ्रं}\ H}}\skt{ }\edtext{\skt{तीव्रं\-ना\-श\-ये\-युश्च}}{\lemma{\skt{तीव्रं\appbrk{}ना\appbrk{}श\appbrk{}ये\appbrk{}युश्च}}\Afootnote{H; \skt{तूर्णं नाशयत्येष योगः}\ A}}\skt{ }\edtext{\skt{दर्द्द्रूं}}{\lemma{\skt{दर्द्द्रूं}}\Afootnote{H; \emph{om.}\ A}}\skt{~।}\pend

\pstart \ledsidenote{\passref{4.9.14}}\edtext{\skt{ज्योति\-ष्म\-ती\-दृ\-मालः}}{\lemma{\skt{ज्योति\appbrk{}ष्म\appbrk{}ती\appbrk{}दृ\appbrk{}मालः}}\Afootnote{H; \skt{हेमक्षीरी व्याधिघातः}\ A}}\skt{ शिरीषो निम्बः सर्ज्जः }\edtext{\skt{कुटजः}}{\lemma{\skt{कुटजः}}\Afootnote{H; \skt{वत्सकः}\ A}}\skt{ साजकर्ण्णः~। }\edtext{\edtext{\skt{शीघ्र\-न्ती\-व्रा\-ना\-श\-य\-न्त्येव}}{\lemma{\skt{तीव्रानाशयन्त्य्}}\Afootnote{H; \skt{तीव्रा नाशयन्ती}\ A}}}{\lemma{\skt{एव}}\Afootnote{H; \skt{ह}\ A}}\skt{ }\edtext{\skt{दर्द्द्रीं}}{\lemma{\skt{दर्द्द्रीं}}\Afootnote{H; \skt{दद्रूः}\ A}}\skt{ }\edtext{\skt{स्नाना\-ले\-पो\-द्घ\-र्ष\-णे\-षू\-प\-युक्तः}}{\lemma{\skt{स्नाना\appbrk{}ले\appbrk{}पो\appbrk{}द्घ\appbrk{}र्ष\appbrk{}णे\appbrk{}षू\appbrk{}प\appbrk{}युक्तः}}\Afootnote{H; \skt{स्नाना\appbrk{}ले\appbrk{}पो\appbrk{}द्घ\appbrk{}र्ष\appbrk{}णे\appbrk{}षू\appbrk{}प\appbrk{}युक्ताः}\ A}}\skt{~॥}\pend

\pstart\versepar \ledsidenote{\passref{4.9.15}}\edtext{\skt{उदु\-म्ब\-री\-भ\-द्र\-ना\-मा\-नु\-यावै}}{\lemma{\skt{उदु\appbrk{}म्ब\appbrk{}री\appbrk{}भ\appbrk{}द्र\appbrk{}ना\appbrk{}मा\appbrk{}नु\appbrk{}यावै}}\Afootnote{H; \skt{भद्}\ A}}\skt{ }\edtext{\skt{तस्या}}{\lemma{\skt{तस्या}}\Afootnote{H; \skt{रा}\ A}}\skt{ }\edtext{\edtext{\skt{मूल\-म्म\-ल\-यू\-मू\-ल\-मिस्रं}}{\lemma{\skt{मूलम्}}\Afootnote{H; \skt{सं}\ A}}}{\lemma{\skt{मल\appbrk{}यू\appbrk{}मू\appbrk{}ल\appbrk{}मिस्रं}}\Afootnote{H; \skt{ज्ञोदु\appbrk{}म्ब\appbrk{}री\appbrk{}मू\appbrk{}ल\appbrk{}तुल्यं दत्त्वा मूलं}\ A}}\skt{~। }\pend
\pstart\versepar \skt{}\edtext{\edtext{\skt{सिद्धन्तोये}}{\lemma{\skt{सिद्धन्}}\Afootnote{H; \skt{क्षोदयित्वा मलप्वाः। सिद्धं}\ A}}}{\lemma{\skt{तोये}}\Afootnote{H; \skt{तोयं}\ A}}\skt{ }\edtext{\edtext{\skt{पीतमुष्णे}}{\lemma{\skt{पीतम्}}\Afootnote{A; \skt{पीत}\ H}}}{\lemma{\skt{उष्णे}}\Afootnote{A; \skt{मुष्णे}\ H}}\skt{ सुखोष्णं स्फोटाञ्छ्वित्रे पुण्डरीके च कुर्यात्~।}\pend

\pstart\versepar \ledsidenote{\passref{4.9.16}}\edtext{\skt{द्विप्यं}}{\lemma{\skt{द्विप्यं}}\Afootnote{H; \skt{द्वैपं}\ A}}\skt{ }\edtext{\edtext{\skt{दग्ध\-ञ्च\-र्म्म\-मा\-त\-ङ्ग\-जञ्च}}{\lemma{\skt{चर्म्मम्}}\Afootnote{\skt{चर्म}\ A H}}}{\lemma{\skt{आतङ्गजञ्च}}\Afootnote{\skt{मातङ्गजं वा}\ A; \skt{मातङ्गजञ्च}\ H}}\skt{ भिन्ने स्फोटे }\edtext{\skt{तैलयुक्तः}}{\lemma{\skt{तैलयुक्तः}}\Afootnote{H; \skt{तैलयुक्तं}\ A}}\skt{ प्रलेपः~। }\pend
\pstart\versepar \skt{पूतिः कीटो राज\-वृ\-क्षो\-द्भ\-वेन क्षारेणाक्तः }\edtext{\edtext{\skt{श्वित्रमेको}}{\lemma{\skt{श्वित्रम्}}\Afootnote{A; \skt{श्वित्र}\ H}}}{\lemma{\skt{एको}}\Afootnote{A; \skt{मेको}\ H}}\skt{ निहन्ति~॥}\pend

\pstart\versepar \ledsidenote{\passref{4.9.17}}\edtext{\skt{अव\-ल्गु\-जा\-फ\-ल\-मग्र्यं}}{\lemma{\skt{अव\appbrk{}ल्गु\appbrk{}जा\appbrk{}फ\appbrk{}ल\appbrk{}मग्र्यं}}\Afootnote{\skt{कृष्णस्य सर्पस्य मसी सुदग्धा बैभीतकं तैलमथ द्वितीयम्। एतत्समस्तं मृदितं प्रलेपाच्छ्वित्राणि सर्वाण्यपहन्ति शीघ्रम्}\ A; \skt{अव\appbrk{}ल्गु\appbrk{}जा\appbrk{}फ\appbrk{}ल\appbrk{}मग्र्यं}\ H}}\skt{ }\edtext{\skt{नदी\-ज\-म\-यो\-र\-जो\-मा\-ग\-धि\-का\-स्ति\-लाश्च}}{\lemma{\skt{नदी\appbrk{}ज\appbrk{}म\appbrk{}यो\appbrk{}र\appbrk{}जो\appbrk{}मा\appbrk{}ग\appbrk{}धि\appbrk{}का\appbrk{}स्ति\appbrk{}लाश्च}}\Afootnote{\emph{om.}\ A; \skt{नदी\appbrk{}ज\appbrk{}म\appbrk{}यो\appbrk{}र\appbrk{}जो\appbrk{}मा\appbrk{}ग\appbrk{}धि\appbrk{}का\appbrk{}स्ति\appbrk{}लाश्च}\ H}}\skt{~। }\pend
\pstart\versepar \skt{}\edtext{\skt{उदुम्बरी}}{\lemma{\skt{उदुम्बरी}}\Afootnote{H; \emph{om.}\ A}}\skt{ }\edtext{\skt{काकनामा}}{\lemma{\skt{काकनामा}}\Afootnote{H; \emph{om.}\ A}}\skt{ }\edtext{\skt{सलाक्षा}}{\lemma{\skt{सलाक्षा}}\Afootnote{H; \emph{om.}\ A}}\skt{ }\edtext{\skt{रसाञ्जनञ्चैव}}{\lemma{\skt{रसाञ्जनञ्चैव}}\Afootnote{H; \emph{om.}\ A}}\skt{ }\edtext{\skt{वर्ग्गः}}{\lemma{\skt{वर्ग्गः}}\Afootnote{H; \emph{om.}\ A}}\skt{ }\edtext{\skt{समांशः}}{\lemma{\skt{समांशः}}\Afootnote{H; \emph{om.}\ A}}\skt{~॥}\pend

\pstart\versepar \ledsidenote{\passref{4.9.17.add}}\skt{गोमूत्रपिष्टो गुटिकाकृतस्तु लेपः कार्याश्वित्रिणां श्वित्रहन्ता~। }\pend
\pstart\versepar \skt{लेपोच्छ्वित्रं शिखिनो हन्ति पित्तं क्रीवे\-च\-म्वा\-द\-ग्ध\-मे\-तेन युक्तं~॥}\pend

\pstart\versepar \ledsidenote{\passref{4.9.39}}\edtext{\skt{मूत्र\-ङ्ग\-व्य\-ञ्चि\-त्रकं}}{\lemma{\skt{चित्रकं}}\Afootnote{H; \skt{चित्रक}\ A}}\skt{ व्योषयुक्तं सर्प्पिः कुम्भे क्षौद्रयुक्तं स्थितं हि~। }\pend
\pstart\versepar \skt{}\edtext{\skt{पक्षादूर्द्ध्वं}}{\lemma{\skt{पक्षादूर्द्ध्वं}}\Afootnote{A; \skt{पक्षादूर्द्ध्वं}\ H}}\skt{ श्वित्रिभिः }\edtext{\skt{पेय\-मे\-त\-त्कु\-र्या\-च्चास्मै}}{\lemma{\skt{चास्मै}}\Afootnote{H; \skt{चास्मिन्}\ A}}\skt{ }\edtext{\skt{कुष्ठ\-दृ\-ष्ट\-म्वि\-धानं}}{\lemma{\skt{कुष्ठदृष्टम्}}\Afootnote{H; \skt{कुष्ठदिष्टं}\ A}}\skt{~॥}\pend

\pstart \ledsidenote{\passref{4.9.40}}\edtext{\skt{पूति\-का\-र्क्क\-स्नु\-हा\-रा\-ज\-वृ\-क्ष\-द्रु\-माणां}}{\lemma{\skt{पूति\appbrk{}का\appbrk{}र्क्क\appbrk{}स्नु\appbrk{}हा\appbrk{}रा\appbrk{}ज\appbrk{}वृ\appbrk{}क्ष\appbrk{}द्रु\appbrk{}माणां}}\Afootnote{H; \skt{पूती\appbrk{}का\appbrk{}र्क\appbrk{}स्रु\appbrk{}ङ्न\appbrk{}रे\appbrk{}न्द्र\appbrk{}द्रु\appbrk{}माणां}\ A}}\skt{ मूत्रैः पिष्टाः }\edtext{\skt{पल्ल\-वा\-व\-ल्क\-लाश्च}}{\lemma{\skt{पल्ल\appbrk{}वा\appbrk{}व\appbrk{}ल्क\appbrk{}लाश्च}}\Afootnote{H; \skt{पल्लवाः सौमनाश्च}\ A}}\skt{~। }\edtext{\skt{लेपैः}}{\lemma{\skt{लेपैः}}\Afootnote{H; \skt{लेपः}\ A}}\skt{ श्वित्रं }\edtext{\skt{घ्नन्ति}}{\lemma{\skt{घ्नन्ति}}\Afootnote{H; \skt{हन्ति}\ A}}\skt{ }\edtext{\skt{दर्द्द्रूम्व्रणाश्च}}{\lemma{\skt{दर्द्द्रूम्}}\Afootnote{H; \skt{दद्रूर्}\ A}}\skt{ }\edtext{\skt{दुष्ट\-न्य\-र्शां\-स्यु\-ग्र\-ना\-डीं\-व्र\-णाश्च}}{\lemma{\skt{दुष्ट\appbrk{}न्य\appbrk{}र्शां\appbrk{}स्यु\appbrk{}ग्र\appbrk{}ना\appbrk{}डीं\appbrk{}व्र\appbrk{}णाश्च}}\Afootnote{H; \skt{दुष्टान्यर्शांस्येष पाठाऽ\appbrk{}गु\appbrk{}रु\appbrk{}ह\appbrk{}रि\appbrk{}द्रा\appbrk{}क\appbrk{}षायं, पिष्टमेहिनं हरि\appbrk{}द्रा\appbrk{}दा\appbrk{}रु\appbrk{}ह\appbrk{}रि\appbrk{}द्रा\appbrk{}क\appbrk{}षायं, सान्द्रमेहिनं सप्त\appbrk{}प\appbrk{}र्ण\appbrk{}क\appbrk{}षायं, शुक्रमेहिनं दूर्वा\appbrk{}शै\appbrk{}व\appbrk{}ल\appbrk{}प्ल\appbrk{}व\appbrk{}ह\appbrk{}ठ\appbrk{}क\appbrk{}र\appbrk{}ञ्ज\appbrk{}क\appbrk{}से\appbrk{}रु\appbrk{}क\appbrk{}क\appbrk{}षायं ककु\appbrk{}भ\appbrk{}च\appbrk{}न्द\appbrk{}न\appbrk{}क\appbrk{}षायं वा, फेनमेहिनं त्रिफ\appbrk{}ला\appbrk{}र\appbrk{}ग्व\appbrk{}घ\appbrk{}मृ\appbrk{}द्वी\appbrk{}का\appbrk{}क\appbrk{}षायं, मधु\appbrk{}म\appbrk{}धु\appbrk{}र\appbrk{}मिति; पैत्तिकेषु नीलमेहिनं शाल\appbrk{}सा\appbrk{}रा\appbrk{}दि\appbrk{}क\appbrk{}षा\appbrk{}य\appbrk{}म\appbrk{}श्व\appbrk{}त्थ\appbrk{}क\appbrk{}षायं वा पाययेत्, हरिद्रामेहिनं राज\appbrk{}वृ\appbrk{}क्ष\appbrk{}क\appbrk{}षायं, अम्लमेहिनं न्यग्रो\appbrk{}धा\appbrk{}दि\appbrk{}क\appbrk{}षायं, क्षारमेहिनं त्रिफलाकषायं, मञ्जिष्ठामेहिनं मञ्जि\appbrk{}ष्ठा\appbrk{}च\appbrk{}न्द\appbrk{}न\appbrk{}क\appbrk{}षायं, शोणितमेहिनं गुडू\appbrk{}ची\appbrk{}ति\appbrk{}न्दु\appbrk{}का\appbrk{}स्थि\appbrk{}का\appbrk{}श्म\appbrk{}र्य\appbrk{}ख\appbrk{}र्जू\appbrk{}र\appbrk{}क\appbrk{}षायं मधुमिश्रं; अत ऊर्ध्व\appbrk{}म\appbrk{}सा\appbrk{}ध्ये\appbrk{}ष्वपि योगान्यापनार्थं वक्ष्यामः, तद्यथा सर्पिर्मेहिनं कुष्ठ\appbrk{}कु\appbrk{}ट\appbrk{}ज\appbrk{}पा\appbrk{}ठा\appbrk{}हि\appbrk{}ङ्गु\appbrk{}क\appbrk{}टु\appbrk{}रो\appbrk{}हि\appbrk{}णी\appbrk{}कल्कं गुडू\appbrk{}ची\appbrk{}चि\appbrk{}त्र\appbrk{}क\appbrk{}क\appbrk{}षा\appbrk{}येण पाययेत्, वसा\appbrk{}मे\appbrk{}हि\appbrk{}न\appbrk{}म\appbrk{}ग्नि\appbrk{}म\appbrk{}न्थ\appbrk{}क\appbrk{}षायं शिंशपाकषायं वा, क्षौद्रमेहिनं कद\appbrk{}र\appbrk{}क्र\appbrk{}मु\appbrk{}क\appbrk{}क\appbrk{}षायं, हस्तिमेहिनं तिन्दु\appbrk{}क\appbrk{}क\appbrk{}पि\appbrk{}त्थ\appbrk{}शि\appbrk{}री\appbrk{}ष\appbrk{}प\appbrk{}ला\appbrk{}श\appbrk{}पा\appbrk{}ठा\appbrk{}मू\appbrk{}र्वा\appbrk{}दुः\appbrk{}स्थि\appbrk{}क्षारं चेति; दह्य\appbrk{}मा\appbrk{}न\appbrk{}मौ\appbrk{}द\appbrk{}क\appbrk{}क\appbrk{}न्द\appbrk{}क्वा\appbrk{}थ\appbrk{}सिद्धां यवागूं क्षीरे\appbrk{}क्षु\appbrk{}र\appbrk{}स\appbrk{}म\appbrk{}धुरां पाययेत्}\ A}}\skt{~॥}\pend

\pstart\versepar \ledsidenote{\passref{4.9.41}}\edtext{\skt{ततो}}{\lemma{\skt{ततो}}\Afootnote{H; \skt{अस्}\ A}}\skt{ }\edtext{\edtext{\skt{जन्तुभिर्न्निःस्रुते}}{\lemma{\skt{जन्तुभिर्}}\Afootnote{H; \skt{मादूर्}\ A}}}{\lemma{\skt{न्निःस्रुते}}\Afootnote{H; \skt{ध्वं निःस्रुते}\ A}}\skt{ दुष्टरक्ते जातप्राणं सर्प्पिषा स्नेहयित्वा~। }\pend
\pstart\versepar \skt{}\edtext{\skt{तीक्ष्णै\-र्यो\-गै\-र्व्वा\-म\-यित्वा}}{\lemma{\skt{व्वामयित्वा}}\Afootnote{H; \skt{छर्दयित्वा}\ A}}\skt{ }\edtext{\skt{प्रगाढं}}{\lemma{\skt{प्रगाढं}}\Afootnote{H; \skt{प्रागाढं}\ A}}\skt{ }\edtext{\edtext{\edtext{\skt{सर्व्वा\-न्दो\-षा\-न्नि\-र्ह\-रे\-च्चा\-प्र\-मत्ताः}}{\lemma{\skt{सर्व्वान्}}\Afootnote{H; \skt{पश्चाद्}\ A}}}{\lemma{\skt{दोषान्}}\Afootnote{H; \skt{दोषं}\ A}}}{\lemma{\skt{चाप्रमत्ताः}}\Afootnote{H; \skt{चाप्रमत्तः}\ A}}\skt{~॥}\pend

\pstart \ledsidenote{\passref{4.9.42}}\skt{दुर्व्वान्तो वा }\edtext{\skt{दुर्व्विरिक्तोऽथवा}}{\lemma{\skt{ऽथवा}}\Afootnote{H; \skt{ऽपि वा}\ A}}\skt{ स्यात्कु\-ष्ठी\-दो\-षै\-रु\-द्ध\-तै\-र्व्व्या\-प्त\-देहः~। }\edtext{\skt{निःसन्देहं}}{\lemma{\skt{निःसन्देहं}}\Afootnote{H; \skt{निःसन्दिग्धं}\ A}}\skt{ }\edtext{\edtext{\edtext{\skt{यात्य\-सा\-ध्य\-त्व\-माशु}}{\lemma{\skt{यात्य्}}\Afootnote{H; \skt{यात्}\ A}}}{\lemma{\skt{असाध्यत्वम्}}\Afootnote{H; \skt{यसाध्यत्व}\ A}}}{\lemma{\skt{आशु}}\Afootnote{H; \skt{माशु}\ A}}\skt{ }\edtext{\skt{तस्मा\-द्दो\-षा\-न्नि\-र्ह\-रे\-त्तस्य}}{\lemma{\skt{दोषान्}}\Afootnote{H; \skt{कृत्स्नान्}\ A}}\skt{ }\edtext{\skt{कृत्स्नान्}}{\lemma{\skt{कृत्स्नान्}}\Afootnote{H; \skt{दोषान्}\ A}}\skt{~॥}\pend

\pstart\versepar \ledsidenote{\passref{4.9.43}}\edtext{\skt{पक्षा\-त्प\-क्षा\-च्छ\-र्द्द\-ना\-न्य\-भ्यु\-पेया }\edtext{\skt{त्मासा\-त्मा\-सा\-त्स्रं\-स\-न\-ञ्चा\-प्र\-मत्तः}}{\lemma{\skt{चाप्रमत्तः}}\Afootnote{H; \skt{चापि देयम्~। स्राव्यं रक्तं वत्सरे हि द्विरल्पं नस्यं दद्याच्च त्रिरात्रा}\ A}}}{\lemma{\skt{अभ्युपेया त्मासात्}}\Afootnote{H; \skt{अभ्यु\appbrk{}पे\appbrk{}या\appbrk{}न्मा\appbrk{}सान्}\ A}}\skt{~। }\pend
\pstart\versepar \skt{}\edtext{\edtext{\edtext{\edtext{\skt{त्र्यहा\-त्त्र्य\-हा\-न्न\-स्त\-त\-श्चा\-व\-पीडां}}{\lemma{\skt{त्र्यहात्}}\Afootnote{H; \skt{त्त्रिरात्}\ A}}}{\lemma{\skt{त्र्यहान्}}\Afootnote{H; \skt{रात्}\ A}}}{\lemma{\skt{नस्ततश्}}\Afootnote{H; \emph{om.}\ A}}}{\lemma{\skt{चावपीडां}}\Afootnote{H; \emph{om.}\ A}}\skt{ }\edtext{\skt{मासे}}{\lemma{\skt{मासे}}\Afootnote{H; \emph{om.}\ A}}\skt{ }\edtext{\edtext{\edtext{\skt{मासे\-ष्व\-सृ\-क्मो\-क्ष\-ये\-त्षट्सु}}{\lemma{\skt{मासेष्व्}}\Afootnote{H; \emph{om.}\ A}}}{\lemma{\skt{असृक्मोक्षयेत्}}\Afootnote{H; \emph{om.}\ A}}}{\lemma{\skt{षट्सु}}\Afootnote{H; \emph{om.}\ A}}\skt{ }\edtext{\skt{षट्सु}}{\lemma{\skt{षट्सु}}\Afootnote{H; \emph{om.}\ A}}\skt{~॥}\pend

\pstart \ledsidenote{\passref{4.9.44}}\edtext{\skt{हरीतकीं}}{\lemma{\skt{हरीतकीं}}\Afootnote{H; \emph{om.}\ A}}\skt{ }\edtext{\skt{सगुडां}}{\lemma{\skt{सगुडां}}\Afootnote{H; \skt{पथ्या}\ A}}\skt{ }\edtext{\skt{व्योषयुक्तां}}{\lemma{\skt{व्योषयुक्तां}}\Afootnote{H; \skt{व्योषं सेक्षुजातं सतैलं लीढ्वा शीघ्रं}\ A}}\skt{ }\edtext{\skt{मुच्येत}}{\lemma{\skt{मुच्येत}}\Afootnote{H; \skt{मुच्यते}\ A}}\skt{ }\edtext{\edtext{\skt{कुष्ठाच्च}}{\lemma{\skt{कुष्ठाच्}}\Afootnote{H; \skt{कुष्ठरोगात्~। धात्री\appbrk{}प\appbrk{}थ्या\appbrk{}क्षो\appbrk{}प\appbrk{}कु\appbrk{}ल्या\appbrk{}वि\appbrk{}ड\appbrk{}ङ्गा\appbrk{}न्क्षौ\appbrk{}द्रा\appbrk{}ज्या\appbrk{}भ्या\appbrk{}मे\appbrk{}कतो वावलिह्यात्}\ A}}\skt{ लिहं स तैलान्}}{\lemma{\skt{च लिहं स तैलान्}}\Afootnote{H; \emph{om.}\ A}}\skt{~। }\edtext{\skt{फलत्रिकं}}{\lemma{\skt{फलत्रिकं}}\Afootnote{H; \emph{om.}\ A}}\skt{ }\edtext{\skt{सविडङ्गं}}{\lemma{\skt{सविडङ्गं}}\Afootnote{H; \emph{om.}\ A}}\skt{ }\edtext{\skt{सकृष्णं}}{\lemma{\skt{सकृष्णं}}\Afootnote{H; \emph{om.}\ A}}\skt{ }\edtext{\skt{सर्प्पिस्तैलं}}{\lemma{\skt{सर्प्पिस्तैलं}}\Afootnote{\emph{om.}\ A; \skt{सर्प्पिस्तैलं}\ H}}\skt{ }\edtext{\skt{क्षौद्रयुक्तं}}{\lemma{\skt{क्षौद्रयुक्तं}}\Afootnote{H; \emph{om.}\ A}}\skt{ }\edtext{\skt{लिहेद्वा}}{\lemma{\skt{लिहेद्वा}}\Afootnote{H; \emph{om.}\ A}}\skt{~॥}\pend

\pstart\versepar \ledsidenote{\passref{4.9.45}}\edtext{\skt{पिवं}}{\lemma{\skt{पिवं}}\Afootnote{H; \skt{पीत्वा मासं वा पलांशां}\ A}}\skt{ }\edtext{\skt{हरिद्राम्पलशः}}{\lemma{\skt{पलशः}}\Afootnote{H; \emph{om.}\ A}}\skt{ }\edtext{\skt{समूत्रां}}{\lemma{\skt{समूत्रां}}\Afootnote{H; \skt{मूत्रे}\ A}}\skt{ }\edtext{\edtext{\edtext{\skt{मासा\-न्मु\-च्ये\-त्पा\-प\-रो\-गा\-त्म\-नुष्यः}}{\lemma{\skt{मासान्मुच्येत्}}\Afootnote{H; \skt{णान्तं}\ A}}}{\lemma{\skt{पापरोगात्}}\Afootnote{H; \skt{पापरोगस्}\ A}}}{\lemma{\skt{मनुष्यः}}\Afootnote{H; \skt{य गच्छे}\ A}}\skt{~॥ }\pend
\pstart\versepar \skt{}\edtext{\skt{एवं}}{\lemma{\skt{एवं}}\Afootnote{H; \skt{त्~। एवं}\ A}}\skt{ }\edtext{\edtext{\skt{पेयश्चित्त्रकः}}{\lemma{\skt{पेयश्}}\Afootnote{A; \skt{पेय}\ H}}}{\lemma{\skt{चित्त्रकः}}\Afootnote{A; \skt{श्चित्त्रकः}\ H}}\skt{ श्लक्ष्णपिष्टः पिप्पल्यो वा पूर्व्व\-व\-त्मू\-त्र\-युक्ताः~।}\pend

\pstart\versepar \ledsidenote{\passref{4.9.46}}\edtext{\edtext{\edtext{\skt{रसा\-ञ्ज\-न\-म्पे\-य\-मे\-वञ्च}}{\lemma{\skt{रसाञ्जनम्}}\Afootnote{H; \skt{तद्वत्तार्क्ष्यं मासमा}\ A}}}{\lemma{\skt{पेयम्}}\Afootnote{H; \emph{om.}\ A}}}{\lemma{\skt{एवञ्च}}\Afootnote{H; \skt{त्रं च}\ A}}\skt{ }\edtext{\skt{मास\-न्ते\-ना\-ज\-स्र\-न्दे\-ह\-मा\-ले\-प\-येच्च}}{\lemma{\skt{मास\appbrk{}न्ते\appbrk{}ना\appbrk{}ज\appbrk{}स्रन्}}\Afootnote{H; \skt{पेयं तेनाजस्रं}\ A}}\skt{~॥ }\pend
\pstart\versepar \skt{}\edtext{\edtext{\skt{सप्त\-च्छ\-द\-म्पि\-चु\-म\-र्द्दा\-ट\-रूषं}}{\lemma{\skt{सप्तच्छदम्}}\Afootnote{H; \skt{आरिष्टीत्वक्साप्त}\ A}}\skt{ म्मुस्तं\-ला\-क्षा\-प\-ञ्च\-मूलं}}{\lemma{\skt{पिचु\appbrk{}म\appbrk{}र्द्दा\appbrk{}ट\appbrk{}रूषं म्मुस्तं\appbrk{}ला\appbrk{}क्षा\appbrk{}प\appbrk{}ञ्च\appbrk{}मूलं}}\Afootnote{H; \skt{पर्णी च तुल्या लाक्षा मुस्तं पञ्चमूल्यौ}\ A}}\skt{ हरिद्रे~।}\pend

\pstart\versepar \ledsidenote{\passref{4.9.47}}\edtext{\skt{मञ्जिष्ठा च}}{\lemma{\skt{मञ्जिष्ठा च}}\Afootnote{H; \skt{मञ्जिष्ठा}\ A}}\skt{ }\edtext{\skt{चित्त्रकः}}{\lemma{\skt{चित्त्रकः}}\Afootnote{H; \emph{om.}\ A}}\skt{ }\edtext{\skt{क्षौद्रनामा}}{\lemma{\skt{क्षौद्रनामा}}\Afootnote{H; \skt{क्षौ}\ A}}\skt{ }\edtext{\skt{कटु\-त्रि\-क\-न्त्रि\-फला}}{\lemma{\skt{कटु\appbrk{}त्रि\appbrk{}क\appbrk{}न्त्रि\appbrk{}फला}}\Afootnote{H; \skt{वासको}\ A}}\skt{ देवदारु~॥ }\pend
\pstart\versepar \skt{}\edtext{\skt{चूर्ण्णीकृतं}}{\lemma{\skt{चूर्ण्णीकृतं}}\Afootnote{H; \skt{पथ्यावह्नी व्योष\appbrk{}धा\appbrk{}त्री\appbrk{}वि\appbrk{}डङ्गाः~। सामान्यांशं योजयित्वा विडङ्गैश्चूर्णं कृत्}\ A}}\skt{ }\edtext{\skt{सविडङ्गं}}{\lemma{\skt{सविडङ्गं}}\Afootnote{H; \emph{om.}\ A}}\skt{ }\edtext{\skt{समा\-न्स।\-मे\-त\-दु\-र्ग्ग\-म्प\-लशः}}{\lemma{\skt{समा\appbrk{}न्स।\appbrk{}मे\appbrk{}त\appbrk{}दु\appbrk{}र्ग्ग\appbrk{}म्प\appbrk{}लशः}}\Afootnote{H; \skt{वा}\ A}}\skt{ }\edtext{\skt{सेव्यमानः}}{\lemma{\skt{सेव्यमानः}}\Afootnote{H; \skt{तत्प\appbrk{}लो\appbrk{}न्मा\appbrk{}न\appbrk{}म\appbrk{}श्नन्}\ A}}\skt{~॥}\pend

\pstart\versepar \ledsidenote{\passref{4.9.48}}\edtext{\edtext{\skt{मुच्ये\-त्कु\-ष्ठा\-त्त्रि\-फ\-ला\-व्यो\-ष\-युक्तं}}{\lemma{\skt{मुच्येत्}}\Afootnote{H; \emph{om.}\ A}}}{\lemma{\skt{त्रिफ\appbrk{}ला\appbrk{}व्यो\appbrk{}ष\appbrk{}युक्तं}}\Afootnote{H; \skt{जन्तुर्मुच्यते त्रैफलं वा सर्पिर्द्रोणं व्योषयुक्तं}\ A}}\skt{ }\pend
\pstart\versepar \skt{}\edtext{\skt{चतुर्ग्गुणं}}{\lemma{\skt{चतुर्ग्गुणं}}\Afootnote{H; \skt{च युञ्}\ A}}\skt{ }\edtext{\skt{सर्प्पिषः}}{\lemma{\skt{सर्प्पिषः}}\Afootnote{H; \emph{om.}\ A}}\skt{ }\edtext{\edtext{\skt{प्रस्थमाशु}}{\lemma{\skt{प्रस्थम्}}\Afootnote{H; \skt{ज}\ A}}}{\lemma{\skt{आशु}}\Afootnote{H; \skt{न्}\ A}}\skt{~। }\pend
\pstart\versepar \skt{}\edtext{\skt{गोमूत्रापां}}{\lemma{\skt{गोमूत्रापां}}\Afootnote{H; \skt{गोमूत्राम्बु}\ A}}\skt{ }\edtext{\skt{द्रोण\-सि\-द्धे\-क्षि\-पीले}}{\lemma{\skt{द्रोण\appbrk{}सि\appbrk{}द्धे\appbrk{}क्षि\appbrk{}पीले}}\Afootnote{H; \skt{द्रोण\appbrk{}सि\appbrk{}द्धेऽ\appbrk{}क्ष\appbrk{}पीडे}\ A}}\skt{ }\pend
\pstart\versepar \skt{सिद्धं }\edtext{\skt{सजि\-र्न्ना\-श\-ये\-च्चापि}}{\lemma{\skt{सजिर्न्नाशयेच्}}\Afootnote{H; \skt{सर्पिर्नाशयेच्}\ A}}\skt{ कुष्ठं~॥}\pend

\pstart\versepar \ledsidenote{\passref{4.9.49}}\edtext{\skt{व्याधिघाते}}{\lemma{\skt{व्याधिघाते}}\Afootnote{H; \skt{आरग्वधे}\ A}}\skt{ }\edtext{\skt{सप्त\-प\-र्ण्णे\-व\-ढोले}}{\lemma{\skt{सप्त\appbrk{}प\appbrk{}र्ण्णे\appbrk{}व\appbrk{}ढोले}}\Afootnote{H; \skt{सप्तपर्णे पटोले}\ A}}\skt{ सवृक्षके }\edtext{\skt{सकरञ्जे}}{\lemma{\skt{सकरञ्जे}}\Afootnote{H; \skt{नक्तमाले}\ A}}\skt{ }\edtext{\skt{सनिम्बे}}{\lemma{\skt{सनिम्बे}}\Afootnote{H; \skt{सनिम्}\ A}}\skt{~। }\pend
\pstart\versepar \skt{}\edtext{\skt{समुष्ककद्वे}}{\lemma{\skt{समुष्ककद्वे}}\Afootnote{H; \skt{बे}\ A}}\skt{ }\edtext{\skt{हरिद्रे}}{\lemma{\skt{हरिद्रे}}\Afootnote{H; \skt{जीर्णं}\ A}}\skt{ }\edtext{\skt{विपक्वं}}{\lemma{\skt{विपक्वं}}\Afootnote{H; \skt{पक्वं}\ A}}\skt{ }\edtext{\edtext{\skt{हन्याद्द्वृतं}}{\lemma{\skt{हन्याद्}}\Afootnote{\skt{तद्ध\appbrk{}रि\appbrk{}द्रा\appbrk{}द्व\appbrk{}येन हन्यात्}\ A; \skt{हन्याद्}\ H}}}{\lemma{\skt{द्वृतं}}\Afootnote{H; \emph{om.}\ A}}\skt{ }\edtext{\skt{सर्व्वकुष्ठानि}}{\lemma{\skt{सर्व्वकुष्ठानि}}\Afootnote{H; \skt{कुष्}\ A}}\skt{ }\edtext{\skt{शीघ्रं}}{\lemma{\skt{शीघ्रं}}\Afootnote{H; \skt{ठं मुष्कके चापि सर्पिः}\ A}}\skt{~॥}\pend

\pstart\versepar \ledsidenote{\passref{4.9.50}}\edtext{\skt{निम्बा\-स\-न\-म्प\-द्मकं}}{\lemma{\skt{निम्बासनम्}}\Afootnote{H; \skt{रोध्रारिष्टं}\ A}}\skt{ }\edtext{\skt{समृ\-णा\-ल\-स\-प्त\-च्छ\-द\-ञ्च\-न्दनं}}{\lemma{\skt{समृ\appbrk{}णा\appbrk{}ल\appbrk{}स\appbrk{}प्त\appbrk{}च्छ\appbrk{}द\appbrk{}ञ्च\appbrk{}न्दनं}}\Afootnote{H; \skt{रक्तसारः सप्ताह्वाक्षौ वृक्षको}\ A}}\skt{ }\edtext{\edtext{\skt{लोध्रमेव}}{\lemma{\skt{लोध्रम्}}\Afootnote{H; \skt{बीजकश्}\ A}}\skt{ च}}{\lemma{\skt{एव च}}\Afootnote{H; \skt{च}\ A}}\skt{~। }\pend
\pstart\versepar \skt{}\edtext{\skt{योज्योत्स्नाने}}{\lemma{\skt{योज्योत्स्नाने}}\Afootnote{H; \skt{योज्याः स्नाने}\ A}}\skt{ दह्यमानस्य जन्तोः }\edtext{\skt{पेया\-वा\-स्या\-न्क्षौ\-द्र\-युक्ता}}{\lemma{\skt{पेयावास्यान्}}\Afootnote{H; \skt{पेया वा स्यात्}\ A}}\skt{ त्रिभण्डी~॥}\pend

\pstart\versepar \ledsidenote{\passref{4.9.51}}\edtext{\edtext{\skt{कुष्ठी\-मां\-से\-ष्व\-व\-दी\-र्यत्सु}}{\lemma{\skt{कुष्ठीमांसेष्व्}}\Afootnote{H; \skt{खादेत्कुष्ठी मां}\ A}}}{\lemma{\skt{अवदीर्यत्सु}}\Afootnote{H; \skt{सशा(पा)ते पु}\ A}}\skt{ }\edtext{\edtext{\skt{खादेत्मुद्गं}}{\lemma{\skt{खादेत्}}\Afootnote{H; \skt{राणान्}\ A}}}{\lemma{\skt{मुद्गं}}\Afootnote{H; \skt{मुद्गान्}\ A}}\skt{ }\edtext{\skt{सिद्धन्निम्बतोये}}{\lemma{\skt{सिद्धन्निम्बतोये}}\Afootnote{H; \skt{सिद्धान्निम्बतोये}\ A}}\skt{ }\edtext{\skt{सतैलं}}{\lemma{\skt{सतैलं}}\Afootnote{H; \skt{सतैलान्}\ A}}\skt{~। }\pend
\pstart\versepar \skt{निम्बक्वाथं जातसत्त्वः }\edtext{\skt{पिबेद्वा}}{\lemma{\skt{पिबेद्वा}}\Afootnote{A; \skt{पिबेद्वा}\ H}}\skt{ क्वाथम्वा }\edtext{\skt{स्याद\-र्क्क\-स\-प्त\-च्छ\-दाभ्यां}}{\lemma{\skt{स्याद\appbrk{}र्क्क\appbrk{}स\appbrk{}प्त\appbrk{}च्छ\appbrk{}दाभ्यां}}\Afootnote{\skt{अर्का\appbrk{}ल\appbrk{}र्क\appbrk{}स\appbrk{}प्त\appbrk{}च्छ\appbrk{}दा\appbrk{}नाम्}\ A; \skt{स्याद\appbrk{}र्क्क\appbrk{}स\appbrk{}प्त\appbrk{}च्छ\appbrk{}दाभ्यां}\ H}}\skt{~॥}\pend

\pstart\versepar \ledsidenote{\passref{4.9.52}}\edtext{\skt{दग्धे\-ष्व\-ङ्गे\-ष्व\-श्व\-मा\-र\-स्य\-मूलं}}{\lemma{\skt{दग्धेष्व्}}\Afootnote{H; \skt{जग्धेष्व्}\ A}}\skt{ }\edtext{\skt{लेहोयोज्यः}}{\lemma{\skt{लेहोयोज्यः}}\Afootnote{H; \skt{लेपो युक्तः}\ A}}\skt{ }\edtext{\skt{सविडङ्गः}}{\lemma{\skt{सविडङ्गः}}\Afootnote{H; \skt{स्याद्विडङ्गैः}\ A}}\skt{ }\edtext{\skt{समूत्रः}}{\lemma{\skt{समूत्रः}}\Afootnote{H; \skt{समूत्रैः}\ A}}\skt{~। }\pend
\pstart\versepar \skt{मूत्रैश्चैनं }\edtext{\skt{सेच\-ये\-दा\-दि\-सेच्च}}{\lemma{\skt{सेच\appbrk{}ये\appbrk{}दा\appbrk{}दि\appbrk{}सेच्च}}\Afootnote{H; \skt{सेच\appbrk{}ये\appbrk{}द्भो\appbrk{}ज\appbrk{}येच्च}\ A}}\skt{ }\edtext{\skt{सर्व्वाहाराः}}{\lemma{\skt{सर्व्वाहाराः}}\Afootnote{H; \skt{सर्वाहारान्}\ A}}\skt{ }\edtext{\skt{सविडंगा}}{\lemma{\skt{सविडंगा}}\Afootnote{H; \skt{संप्र\appbrk{}यु\appbrk{}क्ता\appbrk{}न्वि\appbrk{}डङ्गैः}\ A}}\skt{ }\edtext{\skt{यथा}}{\lemma{\skt{यथा}}\Afootnote{H; \emph{om.}\ A}}\skt{ }\edtext{\skt{स्युः}}{\lemma{\skt{स्युः}}\Afootnote{H; \emph{om.}\ A}}\skt{~॥}\pend

\pstart\versepar \ledsidenote{\passref{4.9.53}}\skt{कारञ्जम्वा सार्षपं }\edtext{\skt{वास्यतैलं}}{\lemma{\skt{वास्यतैलं}}\Afootnote{H; \skt{वा क्षतेषु क्षेप्यं तैलं}\ A}}\skt{ }\edtext{\skt{व्रणे}}{\lemma{\skt{व्रणे}}\Afootnote{H; \emph{om.}\ A}}\skt{ }\edtext{\skt{देयं}}{\lemma{\skt{देयं}}\Afootnote{H; \emph{om.}\ A}}\skt{ }\edtext{\edtext{\skt{शिग्रु\-को\-सा\-म्व्र\-योर्व्वा}}{\lemma{\skt{शिग्रु\appbrk{}को\appbrk{}सा\appbrk{}म्व्र\appbrk{}योर्}}\Afootnote{H; \skt{शिग्रुकोशाम्रयो}\ A}}}{\lemma{\skt{व्वा}}\Afootnote{H; \skt{र्वा}\ A}}\skt{~। }\pend
\pstart\versepar \skt{पक्वं }\edtext{\skt{सर्व्व\-ङ्क\-टु\-कै\-स्ति\-क्त\-कैश्च}}{\lemma{\skt{सर्व्व\appbrk{}ङ्क\appbrk{}टु\appbrk{}कै\appbrk{}स्ति\appbrk{}क्त\appbrk{}कैश्च}}\Afootnote{H; \skt{सर्वैर्वा कटूष्णैः सतिकैः}\ A}}\skt{ शेषञ्च }\edtext{\edtext{\skt{स्यात्कु\-ष्ठ\-व\-त्स\-म्वि\-धानं}}{\lemma{\skt{स्यात्}}\Afootnote{H; \skt{स्याद्}\ A}}}{\lemma{\skt{कुष्ठ\appbrk{}व\appbrk{}त्स\appbrk{}म्वि\appbrk{}धानं}}\Afootnote{H; \skt{दुष्ट\appbrk{}व\appbrk{}त्सं\appbrk{}वि\appbrk{}धा\appbrk{}नम्}\ A}}\skt{~॥}\pend

\pstart\versepar \ledsidenote{\passref{4.9.54}}\skt{सप्त\-प\-र्ण्ण\-क\-र\-ञ्जा\-र्क्क\-मा\-ल\-ती\-क\-र\-वी\-रजं~। }\pend
\pstart\versepar \skt{}\edtext{\skt{मूलं}}{\lemma{\skt{मूलं}}\Afootnote{H; \emph{om.}\ A}}\skt{ }\edtext{\skt{स्नुहा\-शि\-री\-षा\-भ्या\-ञ्चि\-त्र\-का\-स्फो\-त\-यो\-रपि}}{\lemma{\skt{स्नुहा\appbrk{}शि\appbrk{}री\appbrk{}षा\appbrk{}भ्या\appbrk{}ञ्चि\appbrk{}त्र\appbrk{}का\appbrk{}स्फो\appbrk{}त\appbrk{}योर्}}\Afootnote{H; \skt{स्नुही\appbrk{}शि\appbrk{}री\appbrk{}ष\appbrk{}यो\appbrk{}र्मूलं चित्र\appbrk{}का\appbrk{}स्फो\appbrk{}त\appbrk{}योर्}\ A}}\skt{~॥}\pend

\pstart \ledsidenote{\passref{4.9.55cd}}\edtext{\skt{कर\-ञ्ज\-बी\-ज\-न्त्रि\-फ\-ला\-त्रि\-कटुं}}{\lemma{\skt{कर\appbrk{}ञ्ज\appbrk{}बी\appbrk{}ज\appbrk{}न्त्रि\appbrk{}फ\appbrk{}ला\appbrk{}त्रि\appbrk{}कटुं}}\Afootnote{H; \skt{करञ्जबीजं त्रिकटु त्रिफलां}\ A}}\skt{ रजनीद्वयं~।}\pend

\pstart\versepar \ledsidenote{\passref{4.9.56}}\edtext{\skt{सिद्धार्थकं}}{\lemma{\skt{सिद्धार्थकं}}\Afootnote{H; \skt{सिद्धार्थकान्}\ A}}\skt{ विडङ्गानि }\edtext{\skt{प्रपुनाटञ्च}}{\lemma{\skt{प्रपुनाटञ्च}}\Afootnote{H; \skt{प्रपुन्नाडं च}\ A}}\skt{ संहरेत्~॥ }\pend
\pstart\versepar \skt{मूत्रपिष्टैः }\edtext{\edtext{\skt{पचेत्तैलमेभिः}}{\lemma{\skt{पचेत्}}\Afootnote{H; \skt{पचेदेतैस्}\ A}}}{\lemma{\skt{एभिः}}\Afootnote{H; \emph{om.}\ A}}\skt{ कुष्ठविनाशनं~।}\pend

\pstart\versepar \ledsidenote{\passref{4.9.64}}\edtext{\skt{अत्यङ्गाद्वजकं}}{\lemma{\skt{अत्यङ्गाद्वजकं}}\Afootnote{H; \emph{om.}\ A}}\skt{ }\edtext{\skt{नाम}}{\lemma{\skt{नाम}}\Afootnote{H; \skt{म}\ A}}\skt{ }\edtext{\skt{नाडी\-दु\-ष्ट\-व्र\-णा\-पहं}}{\lemma{\skt{नाडी\appbrk{}दु\appbrk{}ष्ट\appbrk{}व्र\appbrk{}णा\appbrk{}पहं}}\Afootnote{H; \skt{हाव\appbrk{}ज्र\appbrk{}क\appbrk{}मि\appbrk{}त्ये\appbrk{}त\appbrk{}न्नाम्ना तैलं महागुणम्}\ A}}\skt{~॥ }\pend
\pstart\versepar \skt{}\edtext{\skt{पित्तावाप्ये}}{\lemma{\skt{पित्तावाप्ये}}\Afootnote{H; \skt{पित्तावापैर्}\ A}}\skt{ }\edtext{\skt{मूत्रपिष्टे}}{\lemma{\skt{मूत्रपिष्टे}}\Afootnote{H; \skt{मूत्रपिष्ट}\ A}}\skt{ }\edtext{\skt{तैलं}}{\lemma{\skt{तैलं}}\Afootnote{H; \skt{स्तैलं}\ A}}\skt{ }\edtext{\skt{लाक्षादिके}}{\lemma{\skt{लाक्षादिके}}\Afootnote{H; \skt{लाक्षादिकैः}\ A}}\skt{ कृतं~।}\pend

\pstart\versepar \ledsidenote{\passref{4.9.65}}\edtext{\edtext{\skt{कटु\-का\-ला\-वु\-नि\-भि\-ष\-क्स\-प्ताहं}}{\lemma{\skt{कटु\appbrk{}का\appbrk{}ला\appbrk{}वु\appbrk{}नि\appbrk{}भि\appbrk{}षक्}}\Afootnote{H; \skt{सप्ताहं कटुकालाब्}\ A}}}{\lemma{\skt{सप्ताहं}}\Afootnote{H; \skt{वां}\ A}}\skt{ }\edtext{\skt{विदधीत}}{\lemma{\skt{विदधीत}}\Afootnote{H; \skt{निदधीत}\ A}}\skt{ }\edtext{\skt{च}}{\lemma{\skt{च}}\Afootnote{H; \skt{चिकित्सकः/ पीत}\ A}}\skt{~॥ }\pend
\pstart\versepar \skt{}\edtext{\skt{ततो}}{\lemma{\skt{ततो}}\Afootnote{H; \skt{वन्तं ततो}\ A}}\skt{ }\edtext{\skt{मात्रा\-पी\-त\-व\-न्त\-न्ते\-ना\-त्य\-क्तञ्च}}{\lemma{\skt{मात्रा\appbrk{}पी\appbrk{}त\appbrk{}व\appbrk{}न्त\appbrk{}न्ते\appbrk{}ना\appbrk{}त्य\appbrk{}क्तञ्च}}\Afootnote{H; \skt{मात्रां तेनाभ्यक्तं च}\ A}}\skt{ मानवं~।}\pend

\pstart\versepar \ledsidenote{\passref{4.9.66}}\edtext{\skt{शाय\-यी\-ता\-त\-पे\-तस्य}}{\lemma{\skt{शाय\appbrk{}यी\appbrk{}ता\appbrk{}त\appbrk{}पे\appbrk{}तस्य}}\Afootnote{H; \skt{शाययेदातपे तस्य}\ A}}\skt{ दोषा गच्छन्ति सर्व्वशः~॥ }\pend
\pstart\versepar \skt{}\edtext{\skt{स्थितदोषं}}{\lemma{\skt{स्थितदोषं}}\Afootnote{H; \skt{स्रुतदोषं}\ A}}\skt{ समुत्थाप्य }\edtext{\skt{सुस्नातं}}{\lemma{\skt{सुस्नातं}}\Afootnote{H; \skt{स्नातं}\ A}}\skt{ }\edtext{\skt{खदिराम्बुना}}{\lemma{\skt{खदिराम्बुना}}\Afootnote{H; \skt{खदिरवारिणा}\ A}}\skt{~।}\pend

\pstart\versepar \ledsidenote{\passref{4.9.67}}\edtext{\skt{तेनैव}}{\lemma{\skt{तेनैव}}\Afootnote{H; \emph{om.}\ A}}\skt{ }\edtext{\skt{भोजयेदेनं}}{\lemma{\skt{भोजयेदेनं}}\Afootnote{H; \emph{om.}\ A}}\skt{ यवागूं }\edtext{\skt{वापि}}{\lemma{\skt{वापि}}\Afootnote{H; \emph{om.}\ A}}\skt{ पाययेत्~॥ }\pend
\pstart\versepar \skt{एवं संशोधने वर्ग्गे }\edtext{\skt{कुष्ठ\-द्वे\-ष्यौ\-ष\-धेषु}}{\lemma{\skt{कुष्ठ\appbrk{}द्वे\appbrk{}ष्यौ\appbrk{}ष\appbrk{}धेषु}}\Afootnote{H; \skt{कुष्ठ\appbrk{}घ्ने\appbrk{}ष्वौ\appbrk{}ष\appbrk{}धेषु}\ A}}\skt{ च~।}\edtext{}{\lemma{\skt{पाययेत्~॥ एवं}}\Afootnote{H; \skt{पाययेदेनं साधितां खदिराम्बुना/ एवं}\ A}}\pend

\pstart\versepar \ledsidenote{\passref{4.9.68}}\skt{कुर्यात्तैलानि }\edtext{\skt{सर्प्पींषि}}{\lemma{\skt{सर्प्पींषि}}\Afootnote{H; \skt{सर्पीषि}\ A}}\skt{ }\edtext{\skt{प्रद\-हो\-द्य\-र्ष\-णानि च}}{\lemma{\skt{प्रद\appbrk{}हो\appbrk{}द्य\appbrk{}र्ष\appbrk{}णानि च}}\Afootnote{H; \skt{प्रदे\appbrk{}हो\appbrk{}द्घ\appbrk{}र्ष\appbrk{}णानि च}\ A}}\skt{~॥ }\pend
\pstart\versepar \skt{प्रातः प्रातश्च सेवेत }\edtext{\skt{योगा\-वै\-रे\-च\-नां\-शुभां}}{\lemma{\skt{योगा\appbrk{}वै\appbrk{}रे\appbrk{}च\appbrk{}नां\appbrk{}शुभां}}\Afootnote{H; \skt{योगा\appbrk{}न्वै\appbrk{}रे\appbrk{}च\appbrk{}नाञ्शु}\ A}}\skt{~। }\pend
\pstart\versepar \skt{}\edtext{\skt{यथो\-त्था\-ना\-नि\-ग\-च्छन्ति}}{\lemma{\skt{यथो\appbrk{}त्था\appbrk{}ना\appbrk{}नि\appbrk{}ग\appbrk{}च्छन्ति}}\Afootnote{H; \skt{भान्}\ A}}\skt{ पञ्च षट्सप्त }\edtext{\skt{वाष्टवा}}{\lemma{\skt{वाष्टवा}}\Afootnote{H; \skt{चाष्टौ वा यैरुत्थानं न गच्छति}\ A}}\skt{~॥}\pend

\pstart\versepar \ledsidenote{\passref{4.9.69}}\skt{कारभम्वा }\edtext{\skt{पिबेत्मूत्रं}}{\lemma{\skt{पिबेत्}}\Afootnote{H; \skt{पिबेन्}\ A}}\skt{ जीर्ण्णे }\edtext{\skt{तत्क्षीरभोजनः}}{\lemma{\skt{तत्क्षीरभोजनः}}\Afootnote{H; \skt{तत्क्षी\appbrk{}र\appbrk{}भो\appbrk{}ज\appbrk{}नम्}\ A}}\skt{~। }\pend
\pstart\versepar \skt{जातसत्वानि कुष्ठानि मासैः }\edtext{\edtext{\skt{षड्भिर्व्व्यपोहति}}{\lemma{\skt{षड्भिर्}}\Afootnote{H; \skt{षड्भि}\ A}}}{\lemma{\skt{व्व्यपोहति}}\Afootnote{H; \skt{रपोहति}\ A}}\skt{~॥}\pend

\pstart\versepar \ledsidenote{\passref{4.9.70}}\edtext{\skt{दिदृक्षुरन्तं}}{\lemma{\skt{दिदृक्षुरन्तं}}\Afootnote{\skt{दिदृक्षुरन्तं}\ A H}}\skt{ कुष्ठस्य खदिरं कुष्ठपीडितः~। }\pend
\pstart\versepar \skt{}\edtext{\skt{सर्व्व\-थै\-वो\-प\-यु\-ञ्जीत}}{\lemma{\skt{सर्व्व\appbrk{}थै\appbrk{}वो\appbrk{}प\appbrk{}यु\appbrk{}ञ्जीत}}\Afootnote{H; \skt{सर्वथैव प्रयुञ्जीत}\ A}}\skt{ स्नान\-पा\-ना\-श\-ना\-दिषु~॥}\pend

\pstart\versepar \ledsidenote{\passref{4.9.71}}\skt{यथा हन्ति }\edtext{\edtext{\edtext{\skt{प्रवृ\-द्ध\-त्वा\-न्कु\-ष्ठ\-मा\-तु\-र\-मो\-जसा}}{\lemma{\skt{प्रवृद्धत्वान्}}\Afootnote{\skt{प्रवृद्धत्वात्}\ A; \skt{प्रवृद्धत्वान्}\ H}}}{\lemma{\skt{कुष्ठम्}}\Afootnote{A; \skt{कुष्ठ}\ H}}}{\lemma{\skt{आतुरम्}}\Afootnote{A; \skt{मातुरम्}\ H}}\skt{~। }\pend
\pstart\versepar \skt{}\edtext{\skt{सेव्यमानस्तथा}}{\lemma{\skt{सेव्यमानस्}}\Afootnote{H; \emph{om.}\ A}}\skt{ हन्ति }\edtext{\skt{खदिरः}}{\lemma{\skt{खदिरः}}\Afootnote{H; \skt{उपयुक्तस्तु खदिरः}\ A}}\skt{ }\edtext{\edtext{\skt{कुष्ठमोजसा}}{\lemma{\skt{कुष्ठम्}}\Afootnote{A; \skt{कुष्ठ}\ H}}}{\lemma{\skt{ओजसा}}\Afootnote{A; \skt{मोजसा}\ H}}\skt{~॥}\pend

\pstart\versepar \ledsidenote{\passref{4.9.72}}\edtext{\edtext{\edtext{\skt{नीच\-रो\-म\-न\-खो\-नि\-त्य\-न्नि\-त्य\-मौ\-ष\-ध\-त\-त्परः}}{\lemma{\skt{नीच\appbrk{}रो\appbrk{}म\appbrk{}न\appbrk{}खो\appbrk{}नि\appbrk{}त्यन्}}\Afootnote{H; \skt{नीचरोमनखः श्रान्तो}\ A}}}{\lemma{\skt{नित्यम्}}\Afootnote{H; \skt{हिताश्}\ A}}}{\lemma{\skt{औषधतत्परः}}\Afootnote{H; \skt{यौषधतत्परः}\ A}}\skt{~। }\pend
\pstart\versepar \skt{}\edtext{\skt{योषि\-त्मां\-स\-सु\-रा\-वर्ज्जी}}{\lemma{\skt{योषि\appbrk{}त्मां\appbrk{}स\appbrk{}सु\appbrk{}रा\appbrk{}वर्ज्जी}}\Afootnote{H; \skt{योषि\appbrk{}न्मां\appbrk{}स\appbrk{}सु\appbrk{}रा\appbrk{}वर्जी}\ A}}\skt{ कुष्ठी कुष्ठमपोहति~॥}\pend

\pstart \ledsidenote{\passref{4.9.trailer}}\edtext{\skt{चिकित्सा}}{\lemma{\skt{चिकित्सा}}\Afootnote{H; \skt{इति सुश्रु\appbrk{}त\appbrk{}सं\appbrk{}हि\appbrk{}तायां चिकित्सास्थाने कुष्ठचिकित्सितं}\ A}}\skt{ }\edtext{\skt{नवमः}}{\lemma{\skt{नवमः}}\Afootnote{H; \skt{नाम नवमोऽध्यायः}\ A}}\skt{~॥}\pend

\endnumbering
\beginnumbering
\eledsection*{Adhyāya 10}
\pstart \ledsidenote{\passref{4.10.1}}\skt{अथातो महा\-कु\-ष्ठ\-चि\-कि\-त्सि\-तम्व्या ख्यास्यामः~॥}\pend

\pstart\versepar \ledsidenote{\passref{4.10.3}}\skt{कुष्ठेषु }\edtext{\skt{महेषु}}{\lemma{\skt{महेषु}}\Afootnote{H; \skt{मेहेषु}\ A}}\skt{ कफामयेषु सर्व्वाङ्गशोफेषु च दारुणेषु~। }\pend
\pstart\versepar \skt{}\edtext{\edtext{\skt{कृषत्वमिच्सत्सु}}{\lemma{\skt{कृषत्वम्}}\Afootnote{H; \skt{कृशत्व}\ A}}\skt{ च}}{\lemma{\skt{इच्सत्सु च}}\Afootnote{H; \skt{मिच्छत्सु च}\ A}}\skt{ }\edtext{\skt{मेदूरेषु}}{\lemma{\skt{मेदूरेषु}}\Afootnote{H; \skt{मेदुरेषु}\ A}}\skt{ }\edtext{\skt{योगा\-नि\-मा\-न\-ग्र\-मतिः}}{\lemma{\skt{योगा\appbrk{}नि\appbrk{}मा\appbrk{}न\appbrk{}ग्र\appbrk{}मतिः}}\Afootnote{H; \skt{योगानिमानग्र्}\ A}}\skt{ }\edtext{\skt{प्रयुञ्जते}}{\lemma{\skt{प्रयुञ्जते}}\Afootnote{H; \skt{यम\appbrk{}ति\appbrk{}र्वि\appbrk{}द\appbrk{}ध्यात्}\ A}}\skt{~॥}\pend

\pstart \ledsidenote{\passref{4.10.4}}\edtext{\skt{क्षुण्णां}}{\lemma{\skt{क्षुण्णां}}\Afootnote{H; \skt{क्षुण्णान्}\ A}}\skt{ }\edtext{\skt{यवान्निः}}{\lemma{\skt{यवान्निः}}\Afootnote{H; \skt{यवान्नि}\ A}}\skt{ }\edtext{\skt{पूतां}}{\lemma{\skt{पूतां}}\Afootnote{H; \skt{ष्पूतान्}\ A}}\skt{ रात्रौ }\edtext{\skt{गोमू\-त्र\-प\-र्यु\-षि\-ता\-म्म\-हति}}{\lemma{\skt{गोमू\appbrk{}त्र\appbrk{}प\appbrk{}र्यु\appbrk{}षि\appbrk{}ताम्}}\Afootnote{H; \skt{गोमू\appbrk{}त्र\appbrk{}प\appbrk{}र्यु\appbrk{}षि\appbrk{}तान्}\ A}}\skt{ किलिञ्जे शोषयेदेवं सप्तरात्रं भाव\-ये\-च्छो\-ष\-येच्च~। }\edtext{\skt{ततस्तां}}{\lemma{\skt{ततस्तां}}\Afootnote{H; \skt{ततस्तान्}\ A}}\skt{ }\edtext{\skt{कपालभृष्टां}}{\lemma{\skt{कपालभृष्टां}}\Afootnote{H; \skt{कपालभृष्टान्}\ A}}\skt{ }\edtext{\skt{शक्तुं}}{\lemma{\skt{शक्तुं}}\Afootnote{H; \skt{शक्तून्}\ A}}\skt{ कारयित्वा }\edtext{\skt{प्रातरेव}}{\lemma{\skt{एव}}\Afootnote{H; \skt{प्रातरेव}\ A}}\skt{ }\edtext{\edtext{\skt{कुष्ठिम्मेहिनम्वा}}{\lemma{\skt{कुष्ठिम्}}\Afootnote{H; \skt{कुष्ठिनं}\ A}}}{\lemma{\skt{मेहिनम्वा}}\Afootnote{H; \skt{प्रमेहिणं वा}\ A}}\skt{ }\edtext{\skt{साल\-सा\-रा\-दि\-क\-षा\-येण}}{\lemma{\skt{साल\appbrk{}सा\appbrk{}रा\appbrk{}दि\appbrk{}क\appbrk{}षा\appbrk{}येण}}\Afootnote{H; \skt{साल\appbrk{}सा\appbrk{}रा\appbrk{}दि\appbrk{}क\appbrk{}षा\appbrk{}येण कण्ट\appbrk{}कि\appbrk{}वृ\appbrk{}क्ष\appbrk{}क\appbrk{}षा\appbrk{}येण वा}\ A}}\skt{ पाययेत्~॥ }\edtext{\skt{भल्ला\-त\-क\-प्र\-पु\-ना\-ढा\-व\-ल्गु\-ज\-वि\-ड\-ङ्ग\-मु\-स्ता\-चू\-र्ण्ण\-च\-तु\-र्ब्भा\-ग\-युक्तं}}{\lemma{\skt{भल्ला\appbrk{}त\appbrk{}क\appbrk{}प्र\appbrk{}पु\appbrk{}ना\appbrk{}ढा\appbrk{}व\appbrk{}ल्गु\appbrk{}ज\appbrk{}वि\appbrk{}ड\appbrk{}ङ्ग\appbrk{}मु\appbrk{}स्ता\appbrk{}चू\appbrk{}र्ण्ण\appbrk{}च\appbrk{}तु\appbrk{}र्ब्भा\appbrk{}ग\appbrk{}युक्तं}}\Afootnote{H; \skt{भल्ला\appbrk{}त\appbrk{}क\appbrk{}प्र\appbrk{}पु\appbrk{}न्ना\appbrk{}डा\appbrk{}व\appbrk{}ल्गु\appbrk{}जा\appbrk{}र्क\appbrk{}चि\appbrk{}त्र\appbrk{}क\appbrk{}वि\appbrk{}ड\appbrk{}ङ्ग\appbrk{}मु\appbrk{}स्त\appbrk{}चू\appbrk{}र्ण\appbrk{}च\appbrk{}तु\appbrk{}र्भा\appbrk{}ग\appbrk{}युक्ता}\ A}}\skt{~। }\edtext{\edtext{\skt{एवमेव}}{\lemma{\skt{एवम्}}\Afootnote{H; \skt{न्; एव}\ A}}}{\lemma{\skt{एव}}\Afootnote{H; \skt{मेव}\ A}}\skt{ }\edtext{\skt{शाल\-सा\-रा\-दि\-क\-षा\-य\-प\-रि\-पी\-ता\-न\-ग्व\-धा\-दि\-क\-षा\-य\-क\-प\-रि\-पी\-तम्वा}}{\lemma{\skt{शाल\appbrk{}सा\appbrk{}रा\appbrk{}दि\appbrk{}क\appbrk{}षा\appbrk{}य\appbrk{}प\appbrk{}रि\appbrk{}पी\appbrk{}ता\appbrk{}न\appbrk{}ग्व\appbrk{}धा\appbrk{}दि\appbrk{}क\appbrk{}षा\appbrk{}य\appbrk{}क\appbrk{}प\appbrk{}रि\appbrk{}पी\appbrk{}तम्वा}}\Afootnote{H; \skt{साल\appbrk{}सा\appbrk{}रा\appbrk{}दि\appbrk{}क\appbrk{}षा\appbrk{}य\appbrk{}प\appbrk{}रि\appbrk{}पी\appbrk{}ता\appbrk{}ना\appbrk{}मा\appbrk{}र\appbrk{}ग्व\appbrk{}धा\appbrk{}दि\appbrk{}क\appbrk{}षा\appbrk{}य\appbrk{}प\appbrk{}रि\appbrk{}पी\appbrk{}तानां वा गवा\appbrk{}श्वा\appbrk{}श\appbrk{}कृ\appbrk{}द्भू\appbrk{}तानां वा}\ A}}\skt{~। }\edtext{\edtext{\skt{एवमेव}}{\lemma{\skt{एवम्}}\Afootnote{H; \emph{om.}\ A}}}{\lemma{\skt{एव}}\Afootnote{H; \emph{om.}\ A}}\skt{ }\edtext{\skt{गोहन्नकां}}{\lemma{\skt{गोहन्नकां}}\Afootnote{H; \emph{om.}\ A}}\skt{ }\edtext{\skt{हयन्नकाम्वा}}{\lemma{\skt{हयन्नकाम्वा}}\Afootnote{H; \skt{यवानां}\ A}}\skt{ }\edtext{\skt{शक्तुं}}{\lemma{\skt{शक्तुं}}\Afootnote{H; \skt{शक्तून्}\ A}}\skt{ कारयित्वा भल्लातकादीनां चूर्ण्णान्यावाप्य खदिरा }\edtext{\edtext{\skt{नरा\-ज\-वृ\-क्ष\-रो\-ही\-त\-का\-दी\-ना\-म\-न्य\-त\-मस्य}}{\lemma{\skt{नरा\appbrk{}ज\appbrk{}वृ\appbrk{}क्ष\appbrk{}रो\appbrk{}ही\appbrk{}त\appbrk{}का\appbrk{}दी\appbrk{}नाम्}}\Afootnote{H; \skt{शन\appbrk{}नि\appbrk{}म्ब\appbrk{}रा\appbrk{}ज\appbrk{}वृ\appbrk{}क्ष\appbrk{}रो\appbrk{}ही\appbrk{}त\appbrk{}क\appbrk{}गु\appbrk{}डू\appbrk{}चीना}\ A}}}{\lemma{\skt{अन्यतमस्य}}\Afootnote{H; \skt{मन्यतमस्य}\ A}}\skt{ }\edtext{\skt{कस्̤आये}}{\lemma{\skt{कस्̤आये}}\Afootnote{H; \skt{कषायेण}\ A}}\skt{ }\edtext{\skt{शोलामधुरेण}}{\lemma{\skt{शोलामधुरेण}}\Afootnote{H; \skt{शर्क\appbrk{}रा\appbrk{}म\appbrk{}धु\appbrk{}म\appbrk{}धु\appbrk{}रेण}\ A}}\skt{ }\edtext{\skt{द्राक्षायुक्तेन वा}}{\lemma{\skt{द्राक्षायुक्तेन वा}}\Afootnote{H; \skt{द्राक्षायुक्तेन दा}\ A}}\skt{ }\edtext{\skt{दाडि\-म\-वे\-त\-सा\-म्लेन वा}}{\lemma{\skt{दाडि\appbrk{}म\appbrk{}वे\appbrk{}त\appbrk{}सा\appbrk{}म्लेन वा}}\Afootnote{H; \skt{डिमा\appbrk{}म\appbrk{}ल\appbrk{}क\appbrk{}वे\appbrk{}त\appbrk{}सा\appbrk{}म्लेन}\ A}}\skt{ सैन्ध }\edtext{\skt{वान्वितेन वा}}{\lemma{\skt{वान्वितेन वा}}\Afootnote{H; \skt{वल\appbrk{}व\appbrk{}णा\appbrk{}न्वि\appbrk{}तेन}\ A}}\skt{ पाययेदेष }\edtext{\skt{मन्थकल्पः}}{\lemma{\skt{मन्थकल्पः}}\Afootnote{H; \skt{सर्वमन्थकल्पः}\ A}}\skt{~॥}\pend

\pstart \ledsidenote{\passref{4.10.5}}\edtext{\skt{यवाश्च}}{\lemma{\skt{यवाश्च}}\Afootnote{H; \skt{यावकांश्च}\ A}}\skt{ }\edtext{\skt{भक्ष्याम्वा}}{\lemma{\skt{भक्ष्याम्वा}}\Afootnote{H; \skt{भक्ष्यान्}\ A}}\skt{ }\edtext{\skt{धानोल्लविकां}}{\lemma{\skt{धानोल्लविकां}}\Afootnote{H; \skt{धानोलुम्बक}\ A}}\skt{ }\edtext{\edtext{\skt{कुल्मा\-स\-चू\-र्ण्ण\-को\-सो\-त्का\-रि\-का\-श\-ष्कु\-लि\-का\-ना\-म्वि\-ष\-क\-लो\-ढ\-क\-प्र\-भृ\-ती\-न्यु\-प\-से\-वेत}}{\lemma{\skt{कुल्मा\appbrk{}स\appbrk{}चू\appbrk{}र्ण्ण\appbrk{}को\appbrk{}सो\appbrk{}त्का\appbrk{}रि\appbrk{}का\appbrk{}श\appbrk{}ष्कु\appbrk{}लि\appbrk{}का\appbrk{}नाम्}}\Afootnote{H; \skt{कुल्मा\appbrk{}षा\appbrk{}पू\appbrk{}प\appbrk{}पू\appbrk{}र्ण\appbrk{}को\appbrk{}शो\appbrk{}त्का\appbrk{}रि\appbrk{}का\appbrk{}श\appbrk{}ष्कु\appbrk{}लिका}\ A}}}{\lemma{\skt{विष\appbrk{}क\appbrk{}लो\appbrk{}ढ\appbrk{}क\appbrk{}प्र\appbrk{}भृ\appbrk{}ती\appbrk{}न्यु\appbrk{}प\appbrk{}से\appbrk{}वेत}}\Afootnote{H; \skt{कुणा\appbrk{}वी\appbrk{}प्र\appbrk{}भृ\appbrk{}ती\appbrk{}न्से\appbrk{}वेत}\ A}}\skt{~। }\edtext{\skt{यवविधान}}{\lemma{\skt{यवविधान}}\Afootnote{H; \skt{यवविधाने}\ A}}\skt{ }\edtext{\edtext{\skt{वत्गो\-धू\-म\-वि\-धा\-न\-वे\-णु\-य\-व\-वि\-धानं}}{\lemma{\skt{वत्}}\Afootnote{H; \skt{न}\ A}}}{\lemma{\skt{गोधू\appbrk{}म\appbrk{}वि\appbrk{}धा\appbrk{}न\appbrk{}वे\appbrk{}णु\appbrk{}य\appbrk{}व\appbrk{}वि\appbrk{}धानं}}\Afootnote{H; \skt{गोधू\appbrk{}म\appbrk{}वे\appbrk{}णु\appbrk{}य\appbrk{}वा\appbrk{}नु\appbrk{}प\appbrk{}यु\appbrk{}ञ्जीत}\ A}}\skt{~॥}\pend

\pstart \ledsidenote{\passref{4.10.6}}\skt{अरिष्टानतो वक्ष्यामः~। }\edtext{\skt{पूती\-क\-चि\-त्र\-क\-च\-व्य\-सु\-र\-दा\-रु\-शा\-रि\-वा\-द\-न्ती\-तृ\-वृत्क }\edtext{\skt{त्रिक\-टु\-का\-ना\-मिति}}{\lemma{\skt{इति}}\Afootnote{H\textsuperscript{pc}; \emph{om.}\ A H\textsuperscript{ac}}}}{\lemma{\skt{पूती\appbrk{}क\appbrk{}चि\appbrk{}त्र\appbrk{}क\appbrk{}च\appbrk{}व्य\appbrk{}सु\appbrk{}र\appbrk{}दा\appbrk{}रु\appbrk{}शा\appbrk{}रि\appbrk{}वा\appbrk{}द\appbrk{}न्ती\appbrk{}तृ\appbrk{}वृत्क त्रिकटुकानाम्}}\Afootnote{H\textsuperscript{pc}; \skt{पूती\appbrk{}क\appbrk{}च\appbrk{}व्य\appbrk{}चि\appbrk{}त्र\appbrk{}क\appbrk{}सु\appbrk{}र\appbrk{}दा\appbrk{}रु\appbrk{}सा\appbrk{}रि\appbrk{}वा\appbrk{}द\appbrk{}न्ती\appbrk{}त्रि\appbrk{}वृ\appbrk{}त्त्रि\appbrk{}क\appbrk{}टु\appbrk{}कानां}\ A; \skt{पूती\appbrk{}क\appbrk{}चि\appbrk{}त्र\appbrk{}क\appbrk{}च\appbrk{}व्य\appbrk{}सु\appbrk{}र\appbrk{}दा\appbrk{}रु\appbrk{}शा\appbrk{}रि\appbrk{}वा\appbrk{}द\appbrk{}न्ती\appbrk{}तृ\appbrk{}वृत्क}\ H\textsuperscript{ac}}}\skt{ }\edtext{\skt{पाठः प्रत्येकांशेन}}{\lemma{\skt{पाठः प्रत्येकांशेन}}\Afootnote{H\textsuperscript{pc}; \skt{प्रत्येकं}\ A; \skt{प्रत्येकांशेन}\ H\textsuperscript{ac}}}\skt{ }\edtext{\skt{षद्यलिकाभागाः}}{\lemma{\skt{षद्यलिकाभागाः}}\Afootnote{H; \skt{षट्पलिका भागा}\ A}}\skt{~। }\edtext{\skt{वद\-र\-कु\-ड\-व\-स्त्रि\-फ\-ल\-कु\-डवः}}{\lemma{\skt{वद\appbrk{}र\appbrk{}कु\appbrk{}ड\appbrk{}व\appbrk{}स्त्रि\appbrk{}फ\appbrk{}ल\appbrk{}कु\appbrk{}डवः}}\Afootnote{H; \skt{बद\appbrk{}र\appbrk{}कु\appbrk{}ड\appbrk{}व\appbrk{}स्त्रि\appbrk{}फ\appbrk{}ला\appbrk{}कु\appbrk{}डव}\ A}}\skt{~। }\edtext{\edtext{\skt{इत्येषां}}{\lemma{\skt{इत्य्}}\Afootnote{H; \skt{इत्}\ A}}}{\lemma{\skt{एषां}}\Afootnote{H; \skt{येतेषां}\ A}}\skt{ चूर्ण्णानि ततः }\edtext{\skt{पिप्प\-ली\-म\-धु\-घृ\-तै\-र\-न्य\-त\-रतः}}{\lemma{\skt{अन्यतरतः}}\Afootnote{H; \skt{अन्तः}\ A}}\skt{ प्रलिप्ते घृतभाजने प्राक्कृतसंस्कारे }\edtext{\edtext{\edtext{\skt{सप्तो\-द\-क\-कु\-ड\-वा\-न\-यो\-ज\-जोऽ\-र्द्ध\-कु\-ड\-म\-म\-र्द्ध\-तुलां}}{\lemma{\skt{सप्तो\appbrk{}द\appbrk{}क\appbrk{}कु\appbrk{}ड\appbrk{}वा\appbrk{}न\appbrk{}यो\appbrk{}जजो}}\Afootnote{H; \skt{सप्तो\appbrk{}द\appbrk{}क\appbrk{}कु\appbrk{}ड\appbrk{}वा\appbrk{}न\appbrk{}यो\appbrk{}रजो}\ A}}}{\lemma{\skt{ऽर्द्धकुडमम्}}\Afootnote{H; \skt{ऽर्धकुडव}\ A}}}{\lemma{\skt{अर्द्धतुलां}}\Afootnote{H; \skt{मर्धतुलां च}\ A}}\skt{ गुड\-स्या\-भि\-हि\-तानि }\edtext{\skt{चूर्ण्णा\-न\-वा\-प्या\-नु\-गुप्तं}}{\lemma{\skt{चूर्ण्णा\appbrk{}न\appbrk{}वा\appbrk{}प्या\appbrk{}नु\appbrk{}गुप्तं}}\Afootnote{H; \skt{चूर्णान्यावाप्य स्वनुगुप्तं}\ A}}\skt{ }\edtext{\skt{यवपल्ले}}{\lemma{\skt{यवपल्ले}}\Afootnote{H; \skt{कृत्वा यवपल्ले}\ A}}\skt{ }\edtext{\skt{सप्तचात्रं}}{\lemma{\skt{सप्तचात्रं}}\Afootnote{H; \skt{सप्तरात्रं}\ A}}\skt{ वासयेत्~। ततो }\edtext{\edtext{\skt{यथा\-व\-ल\-मु\-प\-यु\-ञ्जीत}}{\lemma{\skt{यथावलम्}}\Afootnote{H; \skt{यथाबल}\ A}}}{\lemma{\skt{उपयुञ्जीत}}\Afootnote{H; \skt{मुपयुञ्जीत}\ A}}\skt{~। }\edtext{\skt{एव}}{\lemma{\skt{एव}}\Afootnote{H; \skt{ए}\ A}}\skt{ }\edtext{\skt{खल्वरिष्टः}}{\lemma{\skt{खल्वरिष्टः}}\Afootnote{H; \skt{षोऽरिष्टः}\ A}}\skt{ }\edtext{\edtext{\skt{कुष्ठ\-मे\-दो\-र\-पा\-ण्डु\-रो\-गश्च}}{\lemma{\skt{कुष्ठ\appbrk{}मे\appbrk{}दो\appbrk{}र\appbrk{}पा\appbrk{}ण्डु\appbrk{}रो\appbrk{}गश्}}\Afootnote{H; \skt{कुष्ठ\appbrk{}मे\appbrk{}ह\appbrk{}मे\appbrk{}दः\appbrk{}पा\appbrk{}ण्डु\appbrk{}रोग}\ A}}}{\lemma{\skt{च}}\Afootnote{H; \skt{श्व}\ A}}\skt{ यथूनपहन्ति~। एवं }\edtext{\skt{शालसारादौ}}{\lemma{\skt{शालसारादौ}}\Afootnote{H; \skt{शालसारादौ न्यग्रो\appbrk{}धा\appbrk{}दा\appbrk{}वा\appbrk{}र\appbrk{}ग्व\appbrk{}धादौ}\ A}}\skt{ }\edtext{\skt{चारिष्टं}}{\lemma{\skt{चारिष्टं}}\Afootnote{H; \skt{चारिष्टान्}\ A}}\skt{ कुर्व्वीत~॥}\pend

\pstart \ledsidenote{\passref{4.10.7}}\skt{आसवानतो वक्ष्यामः~। }\edtext{\skt{पला\-श\-भ\-स्म\-प\-रि\-स्नु\-त\-स्यो\-ष्णो\-द\-कस्य}}{\lemma{\skt{पला\appbrk{}श\appbrk{}भ\appbrk{}स्म\appbrk{}प\appbrk{}रि\appbrk{}स्नु\appbrk{}त\appbrk{}स्यो\appbrk{}ष्णो\appbrk{}द\appbrk{}कस्य}}\Afootnote{H; \skt{पला\appbrk{}श\appbrk{}भ\appbrk{}स्म\appbrk{}प\appbrk{}रि\appbrk{}स्रु\appbrk{}त\appbrk{}स्यो\appbrk{}ष्णो\appbrk{}द\appbrk{}कस्य शीतीभूतस्य}\ A}}\skt{ त्रयो भागा द्वौ }\edtext{\skt{फाणितस्य}}{\lemma{\skt{फाणितस्य}}\Afootnote{H; \skt{फाणितस्यै}\ A}}\skt{ }\edtext{\edtext{\skt{तदे\-क\-ध्य\-म\-रि\-ष्ट\-क\-ल्पेन}}{\lemma{\skt{तदेकध्यम्}}\Afootnote{H; \skt{कध्य}\ A}}}{\lemma{\skt{अरिष्टकल्पेन}}\Afootnote{H; \skt{मरिष्टकल्पेन}\ A}}\skt{ विदध्यात्~। एवन्तिलादीनां क्षारेषु~। }\edtext{\skt{शालसारादौ}}{\lemma{\skt{शालसारादौ}}\Afootnote{H; \skt{शालसारादौ व्यग्रो\appbrk{}धा\appbrk{}दा\appbrk{}वा\appbrk{}र\appbrk{}ग्व\appbrk{}धादौ}\ A}}\skt{ }\edtext{\skt{मूत्रे}}{\lemma{\skt{मूत्रे}}\Afootnote{H; \skt{मूत्रेषु}\ A}}\skt{ }\edtext{\edtext{\skt{चास\-वा\-म्वि\-द\-द्यात्}}{\lemma{\skt{चासवाम्}}\Afootnote{H; \skt{चासवान्}\ A}}}{\lemma{\skt{विदद्यात्}}\Afootnote{H; \skt{विदध्यात्}\ A}}\skt{~॥}\pend

\pstart \ledsidenote{\passref{4.10.8}}\edtext{\skt{अतः}}{\lemma{\skt{अतः}}\Afootnote{H; \skt{अथ}\ A}}\skt{ }\edtext{\skt{सुराम्वक्ष्यामः}}{\lemma{\skt{सुराम्}}\Afootnote{H; \skt{सुरा}\ A}}\skt{~॥ }\edtext{\skt{शिंस\-पा\-ख\-दि\-रयोः}}{\lemma{\skt{शिंस\appbrk{}पा\appbrk{}ख\appbrk{}दि\appbrk{}रयोः}}\Afootnote{H; \skt{शिंश\appbrk{}पा\appbrk{}ख\appbrk{}दि\appbrk{}रयोः}\ A}}\skt{ }\edtext{\edtext{\skt{सार\-मा\-दा\-यो\-त्पाट्य}}{\lemma{\skt{सारम्}}\Afootnote{H; \skt{सार}\ A}}}{\lemma{\skt{आदायोत्पाट्य}}\Afootnote{H; \skt{मादायोत्पाठ्य}\ A}}\skt{ }\edtext{\skt{चोत्तमक }\edtext{\edtext{\skt{चणी\-या\-व्र\-ह्मी\-को\-श\-व\-ती\-स्त\-त्स\-र्व्व\-मे\-कतः}}{\lemma{\skt{सर्व्वम्}}\Afootnote{H; \skt{सर्व}\ A}}}{\lemma{\skt{एकतः}}\Afootnote{H; \skt{मकतः}\ A}}}{\lemma{\skt{चोत्तमक चणी\appbrk{}या\appbrk{}व्र\appbrk{}ह्मी\appbrk{}को\appbrk{}श\appbrk{}व\appbrk{}ती\appbrk{}स्तत्}}\Afootnote{H\textsuperscript{pc}; \skt{चोत्त\appbrk{}मा\appbrk{}र\appbrk{}णी\appbrk{}ब्रा\appbrk{}ह्मी\appbrk{}को\appbrk{}श\appbrk{}व\appbrk{}ती\appbrk{}स्तत्}\ A; \skt{चोत्त\appbrk{}म\appbrk{}च\appbrk{}णी\appbrk{}या\appbrk{}व्र\appbrk{}ह्मी\appbrk{}को\appbrk{}श\appbrk{}व\appbrk{}ती\appbrk{}स्तत्}\ H\textsuperscript{ac}}}\skt{ }\edtext{\skt{कषायकल्के न}}{\lemma{\skt{कषायकल्के न}}\Afootnote{\skt{कषायकल्पेन}\ A; \skt{कषायक ल्पे ल्के न}\ H}}\skt{ }\edtext{\skt{विपा\-च्यो\-द\-क\-मा\-द\-दीत}}{\lemma{\skt{विपा\appbrk{}च्यो\appbrk{}द\appbrk{}क\appbrk{}मा\appbrk{}द\appbrk{}दीत}}\Afootnote{H; \skt{विपा\appbrk{}च्यो\appbrk{}द\appbrk{}क\appbrk{}मा\appbrk{}ददी}\ A}}\skt{~। }\edtext{\edtext{\skt{तन्म\-ण्डो\-द\-का\-र्थ\-म\-भि\-षु\-णु\-यात्}}{\lemma{\skt{तन्}}\Afootnote{H; \skt{त}\ A}}}{\lemma{\skt{अभिषुणुयात्}}\Afootnote{H; \skt{किण्व\appbrk{}पि\appbrk{}ष्ट\appbrk{}म\appbrk{}भि\appbrk{}षु\appbrk{}णु\appbrk{}याच्च}\ A}}\skt{~। यथोक्तमेवं }\edtext{\skt{सुरां}}{\lemma{\skt{सुरां}}\Afootnote{H; \skt{सुराः}\ A}}\skt{ शालसारादौ न्यग्रो\-धा\-दा\-वा\-र\-ग्व\-धादौ च विदध्यात्~॥}\pend

\pstart \ledsidenote{\passref{4.10.9}}\edtext{\skt{अतो}}{\lemma{\skt{अतो}}\Afootnote{H; \emph{om.}\ A}}\skt{ }\edtext{\edtext{\skt{लेहान्वक्ष्यामः}}{\lemma{\skt{लेहान्}}\Afootnote{H; \emph{om.}\ A}}}{\lemma{\skt{वक्ष्यामः}}\Afootnote{H; \emph{om.}\ A}}\skt{~॥ }\edtext{\skt{खदि\-रा\-स\-न\-नि\-म्व\-रा\-ज\-वृ\-क्ष\-सा\-र\-क्वा\-थेन}}{\lemma{\skt{खदि\appbrk{}रा\appbrk{}स\appbrk{}न\appbrk{}नि\appbrk{}म्व\appbrk{}रा\appbrk{}ज\appbrk{}वृ\appbrk{}क्ष\appbrk{}सा\appbrk{}र\appbrk{}क्वा\appbrk{}थेन}}\Afootnote{H; \emph{om.}\ A}}\skt{ }\edtext{\skt{तत्सा\-र\-पि\-ण्डा\-ने\-व\-श्ले\-क्ष्ण\-पि\-ष्टा\-म्प्र\-क्षिप्य}}{\lemma{\skt{तत्सा\appbrk{}र\appbrk{}पि\appbrk{}ण्डा\appbrk{}ने\appbrk{}व\appbrk{}श्ले\appbrk{}क्ष्ण\appbrk{}पि\appbrk{}ष्टाम्}}\Afootnote{H; \emph{om.}\ A}}\skt{ }\edtext{\edtext{\edtext{\skt{विप\-चे\-त्त\-त्प्ना\-ति\-द्र\-व\-न्ना\-ति\-सा\-न्द्र\-म\-व\-तार्य}}{\lemma{\skt{तत्प्}}\Afootnote{H; \skt{ततो}\ A}}}{\lemma{\skt{नातिसान्द्रम्}}\Afootnote{H; \skt{नातिसान्द्र}\ A}}}{\lemma{\skt{अवतार्य}}\Afootnote{H; \skt{मवतार्य तस्य}\ A}}\skt{ }\edtext{\skt{पाणि\-त\-ल\-म्पू\-र्ण्ण\-म\-प्रा\-त\-राशी}}{\lemma{\skt{पूर्ण्ण\appbrk{}म\appbrk{}प्रा\appbrk{}त\appbrk{}राशी}}\Afootnote{H; \skt{पूर्ण\appbrk{}म\appbrk{}प्रा\appbrk{}त\appbrk{}राशो}\ A}}\skt{ }\edtext{\skt{मधुनाविमिश्र}}{\lemma{\skt{मधुनाविमिश्र}}\Afootnote{H; \skt{मधुमिश्रं}\ A}}\skt{ लिह्यादेवं शालसारादौ }\edtext{\skt{न्यग्रोधादौ च}}{\lemma{\skt{न्यग्रोधादौ च}}\Afootnote{H; \skt{न्यग्रो\appbrk{}धा\appbrk{}दा\appbrk{}वा\appbrk{}र\appbrk{}ग्व\appbrk{}धादौ च}\ A}}\skt{ लेहान्कारयेत्~॥}\pend

\pstart \ledsidenote{\passref{4.10.10}}\edtext{\edtext{\skt{अत\-श्चू\-र्ण्ण\-क्रि\-या\-म्व\-क्ष्यामः}}{\lemma{\skt{अतश्}}\Afootnote{H; \skt{अत}\ A}}}{\lemma{\skt{चूर्ण्णक्रियाम्}}\Afootnote{H; \skt{श्चूर्णक्रियां}\ A}}\skt{~॥ }\edtext{\edtext{\skt{शाल\-सा\-रा\-दि\-चू\-र्ण्ण\-प्र\-स्थ\-मा\-हृत्य}}{\lemma{\skt{शाल\appbrk{}सा\appbrk{}रा\appbrk{}दि\appbrk{}चू\appbrk{}र्ण्ण\appbrk{}प्र\appbrk{}स्थम्}}\Afootnote{H; \skt{शालसारादीनां सारचूर्णप्रस्थ}\ A}}}{\lemma{\skt{आहृत्य}}\Afootnote{H; \skt{माहृत्या}\ A}}\skt{ }\edtext{\edtext{\skt{आर\-ग्व\-धा\-दि\-क\-षा\-य\-प\-रि\-पी\-त\-म\-ने\-कशः}}{\lemma{\skt{आर\appbrk{}ग्व\appbrk{}धा\appbrk{}दि\appbrk{}क\appbrk{}षा\appbrk{}य\appbrk{}प\appbrk{}रि\appbrk{}पी\appbrk{}तम्}}\Afootnote{H; \skt{रग्व\appbrk{}धा\appbrk{}दि\appbrk{}क\appbrk{}षा\appbrk{}य\appbrk{}प\appbrk{}रि\appbrk{}पीत}\ A}}}{\lemma{\skt{अनेकशः}}\Afootnote{H; \skt{मनेकशः}\ A}}\skt{ शाल\-सा\-रा\-दि\-क\-षा\-ये\-णैव पाययेत्~। एवं न्यग्रोधादीनां~। फलेषु पुष्पे\-ष्वा\-र\-ग्व\-धा\-दीनां }\edtext{\edtext{\skt{चूर्ण्णङ्कारयेत्}}{\lemma{\skt{चूर्ण्णङ्}}\Afootnote{H; \skt{चूर्ण}\ A}}}{\lemma{\skt{कारयेत्}}\Afootnote{H; \skt{क्रियां कारयेत्}\ A}}\skt{~॥}\pend

\pstart \ledsidenote{\passref{4.10.11}}\skt{अत }\edtext{\skt{ऊर्द्ध\-म\-य\-स्कृ\-ती\-न्व\-क्ष्यामः}}{\lemma{\skt{ऊर्द्ध\appbrk{}म\appbrk{}य\appbrk{}स्कृ\appbrk{}तीन्}}\Afootnote{H; \skt{ऊर्ध्व\appbrk{}म\appbrk{}य\appbrk{}स्कृ\appbrk{}तीर्}\ A}}\skt{~। }\edtext{\skt{तीक्ष्ण\-लो\-ह\-प\-त्रा\-ण्या\-हृ\-त्य\-त\-न्त\-नूनि}}{\lemma{\skt{तीक्ष्ण\appbrk{}लो\appbrk{}ह\appbrk{}प\appbrk{}त्रा\appbrk{}ण्या\appbrk{}हृ\appbrk{}त्य\appbrk{}तन्}}\Afootnote{H; \skt{तीक्ष्ण\appbrk{}लो\appbrk{}ह\appbrk{}प\appbrk{}त्राणि}\ A}}\skt{ लव\-ण\-व\-र्ग्ग\-प्र\-दि\-ग्धानि }\edtext{\skt{गोम\-या\-ग्नि\-त\-प्तानि}}{\lemma{\skt{गोम\appbrk{}या\appbrk{}ग्नि\appbrk{}त\appbrk{}प्तानि}}\Afootnote{H; \skt{गोम\appbrk{}या\appbrk{}ग्नि\appbrk{}प्र\appbrk{}त\appbrk{}प्तानि}\ A}}\skt{ }\edtext{\skt{कृत्वा}}{\lemma{\skt{कृत्वा}}\Afootnote{H; \emph{om.}\ A}}\skt{ }\edtext{\skt{त्रिफ\-ला\-शा\-ल\-सा\-रा\-दि\-क\-षाये}}{\lemma{\skt{त्रिफ\appbrk{}ला\appbrk{}शा\appbrk{}ल\appbrk{}सा\appbrk{}रा\appbrk{}दि\appbrk{}क\appbrk{}षाये}}\Afootnote{H; \skt{त्रिफ\appbrk{}ला\appbrk{}शा\appbrk{}ल\appbrk{}सा\appbrk{}रा\appbrk{}दि\appbrk{}क\appbrk{}षा\appbrk{}येण}\ A}}\skt{ }\edtext{\edtext{\skt{निर्व्वयेद्देवं}}{\lemma{\skt{निर्व्वयेद्}}\Afootnote{H; \skt{निर्वापयेत्}\ A}}}{\lemma{\skt{देवं}}\Afootnote{H; \emph{om.}\ A}}\skt{ }\edtext{\skt{द्विर\-ष्टौ\-वा\-रा\-न्ततः}}{\lemma{\skt{द्विरष्टौवारान्}}\Afootnote{H; \skt{षोडशवारान्}\ A}}\skt{ }\edtext{\skt{खदि\-रा\-ङ्गा\-र\-त\-प्ता\-न्यू\-प\-शा\-न्ता\-न्य\-शू\-नि\-चू\-र्ण्णानि}}{\lemma{\skt{खदि\appbrk{}रा\appbrk{}ङ्गा\appbrk{}र\appbrk{}त\appbrk{}प्ता\appbrk{}न्यू\appbrk{}प\appbrk{}शा\appbrk{}न्ता\appbrk{}न्य\appbrk{}शू\appbrk{}नि\appbrk{}चू\appbrk{}र्ण्णानि}}\Afootnote{H; \skt{खदि\appbrk{}रा\appbrk{}ङ्गा\appbrk{}र\appbrk{}त\appbrk{}प्ता\appbrk{}न्यु\appbrk{}प\appbrk{}शा\appbrk{}न्त\appbrk{}ता\appbrk{}पानि सूक्ष्मचूर्णानि}\ A}}\skt{ कारयेत्~। घनतान्तव }\edtext{\edtext{\skt{परि\-स्रा\-वि\-ता\-न्नि\-र्गतो}}{\lemma{\skt{परिस्रावितान्}}\Afootnote{H; \skt{परिस्राविता}\ A}}}{\lemma{\skt{निर्गतो}}\Afootnote{H; \skt{नि, ततो}\ A}}\skt{ यथा वलम्मात्रां सर्प्पिर्म्मधुभ्यां }\edtext{\skt{संसृ\-ज्यो\-प\-यु\-ञित}}{\lemma{\skt{संसृ\appbrk{}ज्यो\appbrk{}प\appbrk{}यु\appbrk{}ञित}}\Afootnote{H; \skt{संसृ\appbrk{}ज्यो\appbrk{}प\appbrk{}यु\appbrk{}ञ्ञीत}\ A}}\skt{~॥ जीर्ण्णे यथा }\edtext{\edtext{\skt{व्याधि\-म\-न\-म्ल\-ल\-व\-णा\-हा\-र\-ञ्कु\-र्व्वीत}}{\lemma{\skt{व्याधिम्}}\Afootnote{H; \skt{व्याध्}\ A}}}{\lemma{\skt{अन\appbrk{}म्ल\appbrk{}ल\appbrk{}व\appbrk{}णा\appbrk{}हा\appbrk{}रञ्}}\Afootnote{H; \skt{यन\appbrk{}म्ल\appbrk{}म\appbrk{}ल\appbrk{}व\appbrk{}ण\appbrk{}मा\appbrk{}हारं}\ A}}\skt{~। एवं }\edtext{\edtext{\skt{तुलामुपयुज्य}}{\lemma{\skt{तुलाम्}}\Afootnote{H; \skt{तुल}\ A}}}{\lemma{\skt{उपयुज्य}}\Afootnote{H; \skt{मुपयुज्य}\ A}}\skt{ }\edtext{\edtext{\skt{कुष्ठ\-मे\-ह\-मे\-हश्च}}{\lemma{\skt{कुष्ठ\appbrk{}मे\appbrk{}ह\appbrk{}मे\appbrk{}हश्}}\Afootnote{H; \skt{कुष्ठमेदमेदः}\ A}}}{\lemma{\skt{च}}\Afootnote{H; \skt{श्व}\ A}}\skt{ }\edtext{\skt{यथु\-पा\-ण्डु\-रो\-गो\-न्मा\-दा\-प\-स्मा\-रा\-न\-प\-हृत्य}}{\lemma{\skt{यथु\appbrk{}पा\appbrk{}ण्डु\appbrk{}रो\appbrk{}गो\appbrk{}न्मा\appbrk{}दा\appbrk{}प\appbrk{}स्मा\appbrk{}रा\appbrk{}न\appbrk{}प\appbrk{}हृत्य}}\Afootnote{H; \skt{यथु\appbrk{}पा\appbrk{}ण्डु\appbrk{}रो\appbrk{}गो\appbrk{}न्मा\appbrk{}दा\appbrk{}प\appbrk{}स्मा\appbrk{}रा\appbrk{}न\appbrk{}प\appbrk{}हत्य}\ A}}\skt{ वर्ष\-श\-त\-ञ्जी\-वति~। }\edtext{\edtext{\skt{तुला\-या\-न्तु\-ला\-या\-म्व\-र्ष\-श\-त\-मु\-त्कर्षः}}{\lemma{\skt{वर्षशतम्}}\Afootnote{H; \skt{वर्षशत}\ A}}}{\lemma{\skt{उत्कर्षः}}\Afootnote{H; \skt{मुत्कर्षः}\ A}}\skt{~। एतेन }\edtext{\edtext{\skt{सर्व्व\-लो\-हे\-ष\-य\-स्कृ\-तयो}}{\lemma{\skt{सर्व्वलोहेष्}}\Afootnote{H; \skt{सर्वलौहे}\ A}}}{\lemma{\skt{अयस्कृतयो}}\Afootnote{H; \skt{ष्वयस्कृतयो}\ A}}\skt{ व्याख्याताः~।}\pend

\pstart \ledsidenote{\passref{4.10.12}}\edtext{\skt{तृवृ\-च्छ्या\-मा\-ग्नि\-म\-न्थ\-स\-प्त\-ला\-के\-वु\-क\-शं\-खि\-नी\-शं\-ख\-पु\-ष्पी\-ति\-ल्व\-क\-त्रि\-फ\-ला\-प\-ला\-स\-शिं\-श\-पानां}}{\lemma{\skt{तृवृ\appbrk{}च्छ्या\appbrk{}मा\appbrk{}ग्नि\appbrk{}म\appbrk{}न्थ\appbrk{}स\appbrk{}प्त\appbrk{}ला\appbrk{}के\appbrk{}वु\appbrk{}क\appbrk{}शं\appbrk{}खि\appbrk{}नी\appbrk{}शं\appbrk{}ख\appbrk{}पु\appbrk{}ष्पी\appbrk{}ति\appbrk{}ल्व\appbrk{}क\appbrk{}त्रि\appbrk{}फ\appbrk{}ला\appbrk{}प\appbrk{}ला\appbrk{}स\appbrk{}शिं\appbrk{}श\appbrk{}पानां}}\Afootnote{H; \skt{त्रिवृ\appbrk{}च्छ्या\appbrk{}मा\appbrk{}ग्नि\appbrk{}म\appbrk{}न्थ\appbrk{}स\appbrk{}प्त\appbrk{}ला\appbrk{}के\appbrk{}वु\appbrk{}क\appbrk{}श\appbrk{}ङ्खि\appbrk{}नी\appbrk{}ति\appbrk{}ल्व\appbrk{}क\appbrk{}त्रि\appbrk{}फ\appbrk{}ला\appbrk{}प\appbrk{}ला\appbrk{}श\appbrk{}शिं\appbrk{}श\appbrk{}पानां}\ A}}\skt{ स्वरसमादाय पालाश्यां }\edtext{\skt{द्रोण्या}}{\lemma{\skt{द्रोण्या}}\Afootnote{H; \skt{द्रोण्यामभ्या}\ A}}\skt{ }\edtext{\skt{सिंच्य}}{\lemma{\skt{सिंच्य}}\Afootnote{H; \skt{सिच्य}\ A}}\skt{ }\edtext{\skt{तस्यां}}{\lemma{\skt{तस्यां}}\Afootnote{H; \emph{om.}\ A}}\skt{ }\edtext{\skt{खदि\-रा\-ङ्गा\-र\-त\-प्त\-म\-य\-स्पि\-ण्ड\-मे\-क\-विं\-शति}}{\lemma{\skt{एकविंशति}}\Afootnote{H; \skt{त्रिसप्त}\ A}}\skt{ }\edtext{\skt{वारान्निर्व्वाप्य}}{\lemma{\skt{वारान्}}\Afootnote{H; \skt{कृत्वो}\ A}}\skt{ }\edtext{\edtext{\skt{तत\-स्त\-मा\-दा\-दाय}}{\lemma{\skt{ततस्}}\Afootnote{H; \emph{om.}\ A}}}{\lemma{\skt{तमादादाय}}\Afootnote{H; \skt{तमादाय}\ A}}\skt{ }\edtext{\edtext{\skt{रसमासिच्य}}{\lemma{\skt{रसम्}}\Afootnote{H; \skt{पुन}\ A}}}{\lemma{\skt{आसिच्य}}\Afootnote{H; \skt{रासिच्य}\ A}}\skt{ स्थाल्यां गोमयाग्निना }\edtext{\skt{विपचेत्}}{\lemma{\skt{विपचेत्}}\Afootnote{H; \skt{विपचे}\ A}}\skt{~। }\edtext{\skt{सिद्ध्यति}}{\lemma{\skt{सिद्ध्यति}}\Afootnote{H; \emph{om.}\ A}}\skt{ }\edtext{\edtext{\skt{चास्मि\-न्पि\-प्प\-ल्या\-दि\-चू\-र्ण्णद्वौ}}{\lemma{\skt{चास्मिन्}}\Afootnote{H; \emph{om.}\ A}}}{\lemma{\skt{पिप्प\appbrk{}ल्या\appbrk{}दि\appbrk{}चू\appbrk{}र्ण्णद्वौ}}\Afootnote{H; \emph{om.}\ A}}\skt{ }\edtext{\skt{मधुनः}}{\lemma{\skt{मधुनः}}\Afootnote{H; \emph{om.}\ A}}\skt{ }\edtext{\edtext{\skt{प्रक्षिपेत्तावन्तं}}{\lemma{\skt{प्रक्षिपेत्}}\Afootnote{H; \skt{त्}\ A}}}{\lemma{\skt{तावन्तं}}\Afootnote{H; \emph{om.}\ A}}\skt{ }\edtext{\skt{घृतस्येति}}{\lemma{\skt{घृतस्येति}}\Afootnote{H; \skt{त}\ A}}\skt{~। }\edtext{\edtext{\edtext{\skt{तत\-श्च\-तु\-भा\-गा\-व\-शि\-ष्ट\-म\-व\-ता\-र्या\-ग्नि\-त\-प्ता\-न्य\-य\-स्प\-त्राणि}}{\lemma{\skt{ततश्}}\Afootnote{H; \skt{तश्}\ A}}}{\lemma{\skt{चतु\appbrk{}भा\appbrk{}गा\appbrk{}व\appbrk{}शि\appbrk{}ष्टम्}}\Afootnote{H; \skt{चतु\appbrk{}र्थ\appbrk{}भा\appbrk{}गा\appbrk{}व\appbrk{}शिष्ट}\ A}}}{\lemma{\skt{अव\appbrk{}ता\appbrk{}र्या\appbrk{}ग्नि\appbrk{}त\appbrk{}प्ता\appbrk{}न्य\appbrk{}यस्}}\Afootnote{H; \skt{मवतार्य परिस्राव्य भूयोऽ\appbrk{}ग्नि\appbrk{}त\appbrk{}प्ता\appbrk{}न्ययः}\ A}}\skt{ }\edtext{\skt{प्रक्षिपेत्ततः}}{\lemma{\skt{प्रक्षिपेत्}}\Afootnote{H; \skt{प्रक्षिपेत्, सिध्यति चास्मि\appbrk{}न्पि\appbrk{}प्प\appbrk{}ल्या\appbrk{}दि\appbrk{}चू\appbrk{}र्ण\appbrk{}भागं द्वौ मधु\appbrk{}न\appbrk{}स्ता\appbrk{}व\appbrk{}द्घृ\appbrk{}त\appbrk{}स्येति दद्यात्}\ A}}\skt{ }\edtext{\edtext{\skt{प्रशान्तमायसे}}{\lemma{\skt{प्रशान्तम्}}\Afootnote{H; \skt{प्रशान्त}\ A}}}{\lemma{\skt{आयसे}}\Afootnote{H; \skt{मायसे}\ A}}\skt{ }\edtext{\edtext{\edtext{\skt{पात्रे\-ष्व\-नु\-गु\-प्त\-न्नि\-द\-द्यात्}}{\lemma{\skt{पात्रेष्व्}}\Afootnote{H; \skt{पात्रे स्}\ A}}}{\lemma{\skt{अनुगुप्तन्}}\Afootnote{H; \skt{वनुगुप्तं}\ A}}}{\lemma{\skt{निदद्यात्}}\Afootnote{H; \skt{निदध्यात्}\ A}}\skt{~। ततो }\edtext{\skt{यथा यो}}{\lemma{\skt{यथा यो}}\Afootnote{H; \skt{यथायोगं}\ A}}\skt{ }\edtext{\edtext{\skt{शुक्ति\-प्र\-कु\-ञ्च\-मे\-वो\-प\-यु\-ञ्जीत}}{\lemma{\skt{शुक्तिप्रकुञ्चम्}}\Afootnote{H; \skt{शुक्तिं प्रकुञ्चं}\ A}}}{\lemma{\skt{एवोपयुञ्जीत}}\Afootnote{H; \skt{वोपयुञ्जीत}\ A}}\skt{~। जीर्णे यथा }\edtext{\edtext{\edtext{\skt{व्याधि\-मा\-हा\-र\-मु\-प\-से\-वेत}}{\lemma{\skt{व्याधिम्}}\Afootnote{H; \skt{व्याध्}\ A}}}{\lemma{\skt{आहारम्}}\Afootnote{H; \skt{याहार}\ A}}}{\lemma{\skt{उपसेवेत}}\Afootnote{H; \skt{मुपसेवेत}\ A}}\skt{~। }\edtext{\skt{एषौ\-ष\-धा\-य\-स्कृ\-ति\-र\-सा\-ध्य\-मपि}}{\lemma{\skt{अपि}}\Afootnote{H; \emph{om.}\ A}}\skt{ }\edtext{\skt{कुष्ठम्मेहम्वा}}{\lemma{\skt{मेहम्वा}}\Afootnote{H; \skt{प्रमेहं वा}\ A}}\skt{ साधयति~। }\edtext{\edtext{\skt{स्थूल\-मा\-शु\-लि\-ख\-येति}}{\lemma{\skt{स्थूलम्}}\Afootnote{H; \skt{स्थूल}\ A}}}{\lemma{\skt{आशुलिखयेति}}\Afootnote{H; \skt{मपकर्षति}\ A}}\skt{~। }\edtext{\edtext{\skt{शोफमुप}}{\lemma{\skt{शोफम्}}\Afootnote{H; \skt{शोफ}\ A}}}{\lemma{\skt{उप}}\Afootnote{H; \skt{मुप}\ A}}\skt{ }\edtext{\edtext{\edtext{\skt{सन्न\-म\-ग्नि\-मु\-द्ध\-रति}}{\lemma{\skt{सन्नम्}}\Afootnote{H; \skt{हन्ति, सन्न्}\ A}}}{\lemma{\skt{अग्निम्}}\Afootnote{H; \skt{मग्नि}\ A}}}{\lemma{\skt{उद्धरति}}\Afootnote{H; \skt{मुद्धरति}\ A}}\skt{~। विशेषेण चोपदिश्यते }\edtext{\skt{राज\-य\-क्ष्मि\-णा\-मेष}}{\lemma{\skt{एष}}\Afootnote{H; \emph{om.}\ A}}\skt{ }\edtext{\skt{प्रयोगो}}{\lemma{\skt{प्रयोगो}}\Afootnote{H; \emph{om.}\ A}}\skt{ }\edtext{\edtext{\skt{वर्ष\-श\-ता\-यु\-र\-नेन}}{\lemma{\skt{वर्षशतायुर्}}\Afootnote{H; \skt{वर्षशतायु}\ A}}}{\lemma{\skt{अनेन}}\Afootnote{H; \skt{श्चानया}\ A}}\skt{ पुरुषो भवति~। }\edtext{\edtext{\skt{शाल\-सा\-रा\-दि\-क्वा\-थ\-मा\-सिंच्य}}{\lemma{\skt{शाल\appbrk{}सा\appbrk{}रा\appbrk{}दि\appbrk{}क्वा\appbrk{}थम्}}\Afootnote{H; \skt{शाल\appbrk{}सा\appbrk{}रा\appbrk{}दि\appbrk{}क्वाथ}\ A}}}{\lemma{\skt{आसिंच्य}}\Afootnote{H; \skt{मासिच्य}\ A}}\skt{ }\edtext{\edtext{\skt{पाला\-स्या\-न्द्रो\-ण्या\-मयो}}{\lemma{\skt{पालास्यान्}}\Afootnote{H; \skt{पालाश्यां}\ A}}}{\lemma{\skt{द्रोण्यामयो}}\Afootnote{H; \skt{द्रोण्यामयो}\ A}}\skt{ }\edtext{\edtext{\skt{घनैस्तप्तैः}}{\lemma{\skt{घनैस्}}\Afootnote{H; \skt{घनां}\ A}}}{\lemma{\skt{तप्तैः}}\Afootnote{H; \skt{स्तप्}\ A}}\skt{ }\edtext{\skt{संयोजा}}{\lemma{\skt{संयोजा}}\Afootnote{H; \skt{तान्}\ A}}\skt{ निर्व्वा\-प्य\-कृ\-त\-सं\-स्कारे }\edtext{\skt{कलसेऽभ्यासिच्य}}{\lemma{\skt{कलसे}}\Afootnote{H; \skt{कलशे}\ A}}\skt{ }\edtext{\skt{पिप्पल्यादिचूर्ण्णं}}{\lemma{\skt{पिप्पल्यादिचूर्ण्णं}}\Afootnote{H; \skt{पिप्प\appbrk{}ल्या\appbrk{}दि\appbrk{}चू\appbrk{}र्ण\appbrk{}भागं}\ A}}\skt{ क्षौद्रं }\edtext{\edtext{\skt{गुडमिति}}{\lemma{\skt{गुडम्}}\Afootnote{H; \skt{गुड}\ A}}\skt{ च}}{\lemma{\skt{इति च}}\Afootnote{H; \skt{मिति च}\ A}}\skt{ }\edtext{\edtext{\edtext{\skt{दत्वा\-नु\-गु\-प्त\-न्नि\-द\-ध्या\-त्त\-मौ\-ष\-धा\-य\-स्कृतिं}}{\lemma{\skt{दत्वानुगुप्तन्}}\Afootnote{H; \skt{दत्त्वा स्वनुगुप्तं}\ A}}}{\lemma{\skt{निदध्यात्}}\Afootnote{H; \skt{निदध्यात्, एतां}\ A}}}{\lemma{\skt{तमौ\appbrk{}ष\appbrk{}धा\appbrk{}य\appbrk{}स्कृतिं}}\Afootnote{H; \skt{महौ\appbrk{}ष\appbrk{}धा\appbrk{}य\appbrk{}स्कृतिं}\ A}}\skt{ }\edtext{\edtext{\skt{मास\-म\-र्द्ध\-मा\-सम्वा}}{\lemma{\skt{मासम्}}\Afootnote{H; \skt{मास}\ A}}}{\lemma{\skt{अर्द्धमासम्वा}}\Afootnote{H; \skt{मर्धमासं वा}\ A}}\skt{ स्थितां }\edtext{\edtext{\skt{यथा\-व\-ल\-मु\-प\-यु\-ञ्जा\-देवं}}{\lemma{\skt{यथावलम्}}\Afootnote{H; \skt{यथाबल}\ A}}}{\lemma{\skt{उपयुञ्जादेवं}}\Afootnote{H; \skt{मुपयुञ्जीत/ एवं}\ A}}\skt{ }\edtext{\skt{न्यग्रोधादिषु च}}{\lemma{\skt{न्यग्रोधादिषु च}}\Afootnote{H; \skt{न्यग्रो\appbrk{}धा\appbrk{}दा\appbrk{}वा\appbrk{}रे\appbrk{}व\appbrk{}ता\appbrk{}दिषु च}\ A}}\skt{ }\edtext{\skt{निदध्यात्}}{\lemma{\skt{निदध्यात्}}\Afootnote{H; \skt{विदध्यात्}\ A}}\skt{~॥ ०~॥}\pend

\pstart \ledsidenote{\passref{4.10.13}}\skt{अतः }\edtext{\skt{खदिरविधानं}}{\lemma{\skt{खदिरविधानं}}\Afootnote{H; \skt{खदिरविधान}\ A}}\skt{ }\edtext{\skt{वक्ष्यामः}}{\lemma{\skt{वक्ष्यामः}}\Afootnote{H; \skt{मुपदेक्ष्यामः}\ A}}\skt{~॥ }\edtext{\edtext{\skt{प्रश\-स्त\-दे\-श\-जा\-त\-म\-नु\-प\-ह\-त\-म्म\-ध्य\-म\-व\-यसं}}{\lemma{\skt{प्रश\appbrk{}स्त\appbrk{}दे\appbrk{}श\appbrk{}जा\appbrk{}तम्}}\Afootnote{H; \skt{प्रश\appbrk{}स्त\appbrk{}दे\appbrk{}श\appbrk{}जात}\ A}}}{\lemma{\skt{अनुपहतम्}}\Afootnote{H; \skt{मनुपहतं}\ A}}\skt{ }\edtext{\skt{खदि\-र\-म्प\-रि\-स\-म\-न्तात्}}{\lemma{\skt{परिसमन्तात्}}\Afootnote{H; \skt{परितः}\ A}}\skt{~॥ }\edtext{\skt{खनयित्वा}}{\lemma{\skt{खनयित्वा}}\Afootnote{H; \skt{खानयित्वा}\ A}}\skt{ तस्य मध्यमं मूलं च्छित्वा }\edtext{\skt{यस्मयं}}{\lemma{\skt{यस्मयं}}\Afootnote{H; \skt{अयोमयं}\ A}}\skt{ }\edtext{\edtext{\skt{कुम्भ\-न्त\-स्मि\-न्न\-न्त\-रेऽ\-धस्ता}}{\lemma{\skt{तस्मिन्नन्तरे}}\Afootnote{H; \skt{तस्मिन्नन्तरे}\ A}}}{\lemma{\skt{ऽधस्ता}}\Afootnote{H; \emph{om.}\ A}}\skt{ }\edtext{\skt{निदध्याद्यथा}}{\lemma{\skt{निदध्याद्}}\Afootnote{H; \skt{निदध्याद्}\ A}}\skt{ रस\-ग्र\-ह\-ण\-स\-मर्थो भवति~। }\edtext{\edtext{\skt{तम\-व\-की\-र्यै\-न्ध\-नै\-र्गो\-म\-य\-मि\-श्रै\-रा\-दी\-प\-येत्}}{\lemma{\skt{तम्}}\Afootnote{H; \skt{ततस्तं गोम\appbrk{}य\appbrk{}मृ\appbrk{}दा+\appbrk{}अ\appbrk{}व\appbrk{}लिप्त}\ A}}}{\lemma{\skt{अव\appbrk{}की\appbrk{}र्यै\appbrk{}न्ध\appbrk{}नैर्}}\Afootnote{H; \skt{मव\appbrk{}की\appbrk{}र्ये\appbrk{}न्ध\appbrk{}नैर्}\ A}}\skt{~। }\edtext{\skt{यथास्य}}{\lemma{\skt{यथास्य}}\Afootnote{H; \skt{यथा+अस्य}\ A}}\skt{ }\edtext{\skt{दह्यमानस्याधो}}{\lemma{\skt{दह्यमानस्याधो}}\Afootnote{H; \skt{दह्यमानस्य}\ A}}\skt{ रसः }\edtext{\skt{प्रसवति}}{\lemma{\skt{प्रसवति}}\Afootnote{H; \skt{स्रवत्य्}\ A}}\skt{~। }\edtext{\edtext{\skt{तद्यथा}}{\lemma{\skt{तद्}}\Afootnote{H; \skt{अधस्तात्, तद्}\ A}}}{\lemma{\skt{यथा}}\Afootnote{H; \skt{यदा}\ A}}\skt{ }\edtext{\edtext{\skt{जानी\-या\-त्पू\-र्ण्ण\-भा\-ज\-न\-म\-थै\-न\-मु\-द्धृत्य}}{\lemma{\skt{पूर्ण्णभाजनम्}}\Afootnote{H; \skt{पूर्णं भाजनम्}\ A}}}{\lemma{\skt{अथैनम्}}\Afootnote{H; \skt{इति, अथैनम्}\ A}}\skt{ परिस्राव्य रस\-म\-न्य\-स्मि\-न्पात्रे निधा\-या\-नु\-गु\-प्त\-न्नि\-द\-ध्यात्~। ततो }\edtext{\edtext{\skt{यथा\-यो\-ग\-म्मा\-त्रा\-मा\-म\-ल\-क\-र\-स\-म\-धु\-स\-र्प्पिभिः}}{\lemma{\skt{मात्राम्}}\Afootnote{H; \skt{मात्रा}\ A}}}{\lemma{\skt{आम\appbrk{}ल\appbrk{}क\appbrk{}र\appbrk{}स\appbrk{}म\appbrk{}धु\appbrk{}स\appbrk{}र्प्पिभिः}}\Afootnote{H; \skt{माम\appbrk{}ल\appbrk{}क\appbrk{}र\appbrk{}स\appbrk{}म\appbrk{}धु\appbrk{}स\appbrk{}र्पिर्भिः}\ A}}\skt{ }\edtext{\skt{संसृ\-द्यो\-प\-यु\-ञ्जीत}}{\lemma{\skt{संसृ\appbrk{}द्यो\appbrk{}प\appbrk{}यु\appbrk{}ञ्जीत}}\Afootnote{H; \skt{संसृ\appbrk{}ज्यो\appbrk{}प\appbrk{}यु\appbrk{}ञ्जीत}\ A}}\skt{~। जीर्ण्णे भल्ला\-त\-क\-वि\-धा\-न\-व\-दा\-हारः परिहारश्च }\edtext{\skt{कुष्ठिनां}}{\lemma{\skt{कुष्ठिनां}}\Afootnote{H\textsuperscript{pc}; \emph{om.}\ A H\textsuperscript{ac}}}\skt{ }\edtext{\skt{क्षीरवर्प्रस्थे}}{\lemma{\skt{क्षीरवर्}}\Afootnote{H\textsuperscript{pc}; \emph{om.}\ A H\textsuperscript{ac}}}\skt{ }\edtext{\skt{प्रयुक्ते}}{\lemma{\skt{प्रयुक्ते}}\Afootnote{H; \skt{चोपयुक्ते}\ A}}\skt{ शतं }\edtext{\edtext{\edtext{\edtext{\skt{वर्षा\-णा\-मा\-यु\-षोऽ\-भि\-वृ\-द्धि\-र्भ\-वति}}{\lemma{\skt{वर्षाणाम्}}\Afootnote{H; \skt{वर्षाणा}\ A}}}{\lemma{\skt{आयुषो}}\Afootnote{H; \skt{मायुषो}\ A}}}{\lemma{\skt{ऽभिवृद्धिर्}}\Afootnote{H; \skt{ऽभिवृद्धि}\ A}}}{\lemma{\skt{भवति}}\Afootnote{H; \skt{र्भवति}\ A}}\skt{~। खदि\-र\-सा\-र\-तु\-ला\-मु\-द\-क\-द्रोणे विपाच्य }\edtext{\edtext{\edtext{\skt{षोड\-शां\-शा\-व\-शि\-ष्ट\-म\-व\-ता\-र्या\-नु\-गु\-प्त\-न्नि\-द\-ध्य\-ध्यात्}}{\lemma{\skt{षोड\appbrk{}शां\appbrk{}शा\appbrk{}व\appbrk{}शि\appbrk{}ष्टम्}}\Afootnote{H; \skt{षोड\appbrk{}शां\appbrk{}शा\appbrk{}व\appbrk{}शिष्ट}\ A}}}{\lemma{\skt{अव\appbrk{}ता\appbrk{}र्या\appbrk{}नु\appbrk{}गु\appbrk{}प्तन्}}\Afootnote{H; \skt{मव\appbrk{}ता\appbrk{}र्या\appbrk{}नु\appbrk{}गुप्तं}\ A}}}{\lemma{\skt{निदध्यध्यात्}}\Afootnote{H; \skt{निदध्यात्}\ A}}\skt{~। }\edtext{\skt{ततो}}{\lemma{\skt{ततो}}\Afootnote{H; \skt{त}\ A}}\skt{ }\edtext{\edtext{\skt{यथा\-यो\-ग\-मा\-म\-ल\-क\-र\-स\-म\-धु\-स\-र्प्पिर्भिः}}{\lemma{\skt{यथायोगम्}}\Afootnote{H; \emph{om.}\ A}}}{\lemma{\skt{आम\appbrk{}ल\appbrk{}क\appbrk{}र\appbrk{}स\appbrk{}म\appbrk{}धु\appbrk{}स\appbrk{}र्प्पिर्भिः}}\Afootnote{H; \skt{माम\appbrk{}ल\appbrk{}क\appbrk{}र\appbrk{}स\appbrk{}म\appbrk{}धु\appbrk{}स\appbrk{}र्पिर्भिः}\ A}}\skt{ संसृ\-ज्यो\-प\-यु\-ञ्जीत~। एष }\edtext{\skt{सर्व्व\-वृ\-क्ष\-सा\-र\-कल्पः}}{\lemma{\skt{सर्व्व\appbrk{}वृ\appbrk{}क्ष\appbrk{}सा\appbrk{}र\appbrk{}कल्पः}}\Afootnote{H; \skt{एव सर्व\appbrk{}वृ\appbrk{}क्षा\appbrk{}सा\appbrk{}रेषु कल्पः}\ A}}\skt{~। }\edtext{\skt{खदि\-र\-सा\-र\-तु\-लाम्वा}}{\lemma{\skt{खदि\appbrk{}र\appbrk{}सा\appbrk{}र\appbrk{}तु\appbrk{}लाम्वा}}\Afootnote{H; \skt{खदि\appbrk{}र\appbrk{}सा\appbrk{}र\appbrk{}चू\appbrk{}र्ण\appbrk{}तुलां खदि\appbrk{}र\appbrk{}सा\appbrk{}र\appbrk{}क्वा\appbrk{}थ\appbrk{}मात्रां वा}\ A}}\skt{ प्रातः }\edtext{\edtext{\skt{प्रात\-रु\-प\-से\-वेत}}{\lemma{\skt{प्रातर्}}\Afootnote{H; \skt{प्रात}\ A}}}{\lemma{\skt{उपसेवेत}}\Afootnote{H; \skt{रुपसेवेत}\ A}}\skt{~। }\edtext{\edtext{\skt{खदि\-र\-सा\-र\-क्वा\-थ\-सि\-द्ध\-मा\-वि\-कम्वा}}{\lemma{\skt{खदि\appbrk{}र\appbrk{}सा\appbrk{}र\appbrk{}क्वा\appbrk{}थ\appbrk{}सि\appbrk{}द्धम्}}\Afootnote{H; \skt{खदि\appbrk{}र\appbrk{}सा\appbrk{}र\appbrk{}क्वा\appbrk{}थ\appbrk{}सिद्ध}\ A}}}{\lemma{\skt{आविकम्वा}}\Afootnote{H; \skt{माविकं वा}\ A}}\skt{ सर्प्पिर्}\pend

\pstart \ledsidenote{\passref{4.10.14}}\edtext{\skt{अमृ\-त\-व\-ल्ली\-स्व\-र\-स\-सि\-द्धम्वा}}{\lemma{\skt{अमृ\appbrk{}त\appbrk{}व\appbrk{}ल्ली\appbrk{}स्व\appbrk{}र\appbrk{}स\appbrk{}सि\appbrk{}द्धम्वा}}\Afootnote{H; \skt{अमृ\appbrk{}त\appbrk{}व\appbrk{}ल्ली\appbrk{}स्व\appbrk{}रसं क्वाथं वा}\ A}}\skt{ प्रातः }\edtext{\edtext{\skt{प्रात\-रु\-प\-से\-वेत}}{\lemma{\skt{प्रातर्}}\Afootnote{H; \skt{प्रात}\ A}}}{\lemma{\skt{उपसेवेत}}\Afootnote{H; \skt{रुपसेवेत, तत्सिद्धं वा सर्पिः, अपराह्णे सस\appbrk{}र्पि\appbrk{}ष्क\appbrk{}मो\appbrk{}द\appbrk{}न\appbrk{}मा\appbrk{}म\appbrk{}ल\appbrk{}क\appbrk{}यू\appbrk{}षेण भुञ्जीत; एवं मासमुपयुज्य सर्व\appbrk{}कु\appbrk{}ष्ठै\appbrk{}र्वि\appbrk{}मु\appbrk{}च्यत इति}\ A}}\skt{~।}\pend

\pstart \ledsidenote{\passref{4.10.15}}\edtext{\skt{अपराह्णे}}{\lemma{\skt{अपराह्णे}}\Afootnote{H; \skt{कृष्ण\appbrk{}ति\appbrk{}ल\appbrk{}भ\appbrk{}ल्ला\appbrk{}त\appbrk{}क\appbrk{}तै\appbrk{}ला\appbrk{}म\appbrk{}ल\appbrk{}क\appbrk{}रस}\ A}}\skt{ }\edtext{\skt{सर्प्पि\-ष्म\-न्त\-सा\-रा\-दि\-क\-षा\-यस्य च}}{\lemma{\skt{सर्प्पि\appbrk{}ष्म\appbrk{}न्त\appbrk{}सा\appbrk{}रा\appbrk{}दि\appbrk{}क\appbrk{}षा\appbrk{}यस्य च}}\Afootnote{H; \skt{सर्पिषां द्रोणं शाल\appbrk{}सा\appbrk{}रा\appbrk{}दि\appbrk{}क\appbrk{}षा\appbrk{}य\appbrk{}सय च}\ A}}\skt{ }\edtext{\skt{त्रिफ\-ला\-त्रि\-क\-टु\-त्रि\-व\-र्णा\-का\-व\-ल्गु\-ज\-फ\-रू\-ष\-क\-फ\-ल\-वि\-डं\-ग\-फ\-ल\-म\-ज्ज\-चि\-त्र\-क\-ह\-रि\-द्रा\-द्व\-य\-तृ\-वृ\-द्द\-न्ती\-न्द्र\-य\-व\-य\-ष्ठी\-म\-धु\-का\-ति\-वि\-षा\-र\-सा\-ञ्ज\-न\-प्रि\-यं\-गूनां}}{\lemma{\skt{त्रिफ\appbrk{}ला\appbrk{}त्रि\appbrk{}क\appbrk{}टु\appbrk{}त्रि\appbrk{}व\appbrk{}र्णा\appbrk{}का\appbrk{}व\appbrk{}ल्गु\appbrk{}ज\appbrk{}फ\appbrk{}रू\appbrk{}ष\appbrk{}क\appbrk{}फ\appbrk{}ल\appbrk{}वि\appbrk{}डं\appbrk{}ग\appbrk{}फ\appbrk{}ल\appbrk{}म\appbrk{}ज्ज\appbrk{}चि\appbrk{}त्र\appbrk{}क\appbrk{}ह\appbrk{}रि\appbrk{}द्रा\appbrk{}द्व\appbrk{}य\appbrk{}तृ\appbrk{}वृ\appbrk{}द्द\appbrk{}न्ती\appbrk{}न्द्र\appbrk{}य\appbrk{}व\appbrk{}य\appbrk{}ष्ठी\appbrk{}म\appbrk{}धु\appbrk{}का\appbrk{}ति\appbrk{}वि\appbrk{}षा\appbrk{}र\appbrk{}सा\appbrk{}ञ्ज\appbrk{}न\appbrk{}प्रि\appbrk{}यं\appbrk{}गूनां}}\Afootnote{H; \skt{त्रिफ\appbrk{}ला\appbrk{}त्रि\appbrk{}क\appbrk{}टु\appbrk{}क\appbrk{}प\appbrk{}रू\appbrk{}ष\appbrk{}फ\appbrk{}ल\appbrk{}म\appbrk{}ज्ज\appbrk{}वि\appbrk{}ड\appbrk{}ङ्ग\appbrk{}फ\appbrk{}ल\appbrk{}सा\appbrk{}र\appbrk{}चि\appbrk{}त्रा\appbrk{}र्का\appbrk{}व\appbrk{}ल्गु\appbrk{}ज\appbrk{}ह\appbrk{}रि\appbrk{}द्रा\appbrk{}द्व\appbrk{}य\appbrk{}त्रि\appbrk{}वृ\appbrk{}द्द\appbrk{}न्ती\appbrk{}द्र\appbrk{}व\appbrk{}न्ती\appbrk{}न्द्र\appbrk{}य\appbrk{}व\appbrk{}य\appbrk{}ष्टी\appbrk{}म\appbrk{}धु\appbrk{}का\appbrk{}ति\appbrk{}वि\appbrk{}षा\appbrk{}र\appbrk{}सा\appbrk{}ञ्ज\appbrk{}न\appbrk{}प्रि\appbrk{}य\appbrk{}ङ्गूणां}\ A}}\skt{ }\edtext{\skt{पलिका}}{\lemma{\skt{पलिका}}\Afootnote{H; \skt{पालिका}\ A}}\skt{ }\edtext{\skt{भागाः}}{\lemma{\skt{भागाः}}\Afootnote{H; \skt{भागा}\ A}}\skt{~। }\edtext{\skt{तदैकध्यं}}{\lemma{\skt{तदैकध्यं}}\Afootnote{H; \skt{स्तानैकध्यं}\ A}}\skt{ स्नेह\-पा\-क\-वि\-धा\-नेन }\edtext{\edtext{\edtext{\skt{विप\-चे\-त्त\-त्सा\-धु\-सि\-द्ध\-म\-व\-ता\-र्या\-नु\-गु\-प्त\-न्नि\-द\-ध्या\-त्तत}}{\lemma{\skt{विपचेत्}}\Afootnote{H; \skt{पचेत्}\ A}}}{\lemma{\skt{तत्साधुसिद्धम्}}\Afootnote{H; \skt{तत्साधुसिद्ध}\ A}}}{\lemma{\skt{अव\appbrk{}ता\appbrk{}र्या\appbrk{}नु\appbrk{}गु\appbrk{}प्तन्}}\Afootnote{H; \skt{मवतार्य परि\appbrk{}स्रा\appbrk{}व्या\appbrk{}नु\appbrk{}गुप्तं}\ A}}\skt{ }\edtext{\skt{उप\-सं\-स्कृ\-त\-श\-रीरं}}{\lemma{\skt{उप\appbrk{}सं\appbrk{}स्कृ\appbrk{}त\appbrk{}श\appbrk{}रीरं}}\Afootnote{H; \skt{उप\appbrk{}सं\appbrk{}स्कृ\appbrk{}त\appbrk{}श\appbrk{}रीरः}\ A}}\skt{~। }\edtext{\skt{प्रातरुत्थाय}}{\lemma{\skt{उत्थाय}}\Afootnote{H; \skt{प्रातरुत्थाय}\ A}}\skt{ पाणिशुक्तिमात्रं }\edtext{\skt{क्षौद्रे\-णो\-प\-सं\-सृ\-ज्यो\-प\-यु\-ञ्जीत}}{\lemma{\skt{क्षौद्रे\appbrk{}णो\appbrk{}प\appbrk{}सं\appbrk{}सृ\appbrk{}ज्यो\appbrk{}प\appbrk{}यु\appbrk{}ञ्जीत}}\Afootnote{H; \skt{क्षौद्रेण प्रति\appbrk{}सं\appbrk{}सृ\appbrk{}ज्यो\appbrk{}प\appbrk{}यु\appbrk{}ञ्जि\appbrk{}आत}\ A}}\skt{~। जीर्ण्णे }\edtext{\skt{मुद्ग\-यू\-षे\-णा\-ल\-व\-णेन}}{\lemma{\skt{मुद्ग\appbrk{}यू\appbrk{}षे\appbrk{}णा\appbrk{}ल\appbrk{}व\appbrk{}णेन}}\Afootnote{H; \skt{मुद्गा\appbrk{}म\appbrk{}ल\appbrk{}क\appbrk{}यू\appbrk{}षे\appbrk{}णा\appbrk{}ल\appbrk{}व\appbrk{}णेन}\ A}}\skt{ सर्प्पिष्मन्तं }\edtext{\edtext{\skt{खदि\-रो\-द\-क\-सि\-द्ध\-म्मृ\-दु\-मो\-द\-न\-म\-श्नी\-यात्}}{\lemma{\skt{मृदुमोदनम्}}\Afootnote{H; \skt{मृद्वोदन}\ A}}}{\lemma{\skt{अश्नीयात्}}\Afootnote{H; \skt{मश्नीयात्}\ A}}\skt{~। }\edtext{\edtext{\skt{खदि\-रो\-द\-क\-से\-वी\-त्येवं}}{\lemma{\skt{खदि\appbrk{}रो\appbrk{}द\appbrk{}क\appbrk{}से\appbrk{}वीत्य्}}\Afootnote{H; \skt{खदि\appbrk{}रो\appbrk{}द\appbrk{}क\appbrk{}सेवी, इत्}\ A}}}{\lemma{\skt{एवं}}\Afootnote{H; \skt{येवं}\ A}}\skt{ }\edtext{\edtext{\skt{द्रोणमुपयुञ्ज्य}}{\lemma{\skt{द्रोणम्}}\Afootnote{H; \skt{द्रोण}\ A}}}{\lemma{\skt{उपयुञ्ज्य}}\Afootnote{H; \skt{मुपयुज्य}\ A}}\skt{ सर्व्व\-कु\-ष्ठै\-र्व्वि\-मुक्तः }\edtext{\skt{शुद्धमनाः}}{\lemma{\skt{शुद्धमनाः}}\Afootnote{H; \skt{शुद्धतनुः}\ A}}\skt{ }\edtext{\edtext{\skt{स्मृति\-मा\-न्व\-र्ष\-श\-ता\-यु\-र्ब्भ\-वति}}{\lemma{\skt{वर्षशतायुर्}}\Afootnote{H; \skt{वर्षशतायु}\ A}}}{\lemma{\skt{ब्भवति}}\Afootnote{H; \skt{ररोगो भवति}\ A}}\skt{~।}\pend

\pstart\versepar \ledsidenote{\passref{4.10.16}}\edtext{\edtext{\skt{मन्था\-रि\-ष्टा\-स\-व\-सु\-रा\-ले\-ह\-श्चू\-र्ण्ण\-न्य\-य\-स्कृतीः}}{\lemma{\skt{मन्था\appbrk{}रि\appbrk{}ष्टा\appbrk{}स\appbrk{}व\appbrk{}सु\appbrk{}रा\appbrk{}ले\appbrk{}हश्}}\Afootnote{H; \skt{भवति चात्र सुरा\appbrk{}म\appbrk{}न्था\appbrk{}स\appbrk{}वा\appbrk{}रि\appbrk{}ष्टां\appbrk{}ल्लेहां}\ A}}}{\lemma{\skt{चूर्ण्णन्ययस्कृतीः}}\Afootnote{H; \skt{श्चूर्णान्ययस्कृतीः}\ A}}\skt{~। }\pend
\pstart\versepar \skt{सहस्रशोऽपि कुर्व्वीत वीजेनानेन }\edtext{\skt{बुद्धिमानिति}}{\lemma{\skt{बुद्धिमानिति}}\Afootnote{H; \skt{बुद्धिमान्}\ A}}\skt{~॥ ०~॥}\pend

\pstart \ledsidenote{\passref{4.10.trailer}}\edtext{\skt{चिकित्सा}}{\lemma{\skt{चिकित्सा}}\Afootnote{H; \skt{इति सुश्रु\appbrk{}त\appbrk{}सं\appbrk{}हि\appbrk{}तायां चिकित्सा}\ A}}\skt{ }\edtext{\skt{दशमः}}{\lemma{\skt{दशमः}}\Afootnote{H; \emph{om.}\ A}}\skt{~॥ ०~॥ }\edtext{\skt{द्विव्र\-णी\-प\-ञ्च\-स\-द्यश्च}}{\lemma{\skt{द्विव्र\appbrk{}णी\appbrk{}प\appbrk{}ञ्च\appbrk{}स\appbrk{}द्यश्च}}\Afootnote{H; \emph{om.}\ A}}\skt{ }\edtext{\edtext{\edtext{\skt{भग\-न\-म्वा\-त\-म्म\-हा\-म\-रुत्}}{\lemma{\skt{भगनम्}}\Afootnote{H; \emph{om.}\ A}}}{\lemma{\skt{वातम्}}\Afootnote{H; \emph{om.}\ A}}}{\lemma{\skt{महामरुत्}}\Afootnote{H; \emph{om.}\ A}}\skt{~। }\edtext{\edtext{\skt{दुर्न्ना\-मा\-चा\-श्म\-री\-चैव}}{\lemma{\skt{दुर्न्नामाचाश्}}\Afootnote{H; \emph{om.}\ A}}}{\lemma{\skt{मरीचैव}}\Afootnote{H; \emph{om.}\ A}}\skt{ }\edtext{\skt{भगस्फोटं}}{\lemma{\skt{भगस्फोटं}}\Afootnote{H; \emph{om.}\ A}}\skt{ }\edtext{\skt{त्वगामयं}}{\lemma{\skt{त्वगामयं}}\Afootnote{H; \skt{स्थाने}\ A}}\skt{~॥ }\edtext{\skt{महा\-कु\-ष्ठ\-चि\-कित्सा च}}{\lemma{\skt{महा\appbrk{}कु\appbrk{}ष्ठ\appbrk{}चि\appbrk{}कित्सा च}}\Afootnote{H; \skt{महा\appbrk{}कु\appbrk{}ष्ठ\appbrk{}चि\appbrk{}कित्सि}\ A}}\skt{ }\edtext{\skt{प्रोक्तं वै}}{\lemma{\skt{प्रोक्तं वै}}\Afootnote{H; \emph{om.}\ A}}\skt{ }\edtext{\skt{प्रथमो}}{\lemma{\skt{प्रथमो}}\Afootnote{H; \skt{तं नाम}\ A}}\skt{ }\edtext{\skt{दश~॥ अ}}{\lemma{\skt{दश~॥ अ}}\Afootnote{H; \skt{दशमोऽध्यायः}\ A}}\pend

\endnumbering
\beginnumbering
\eledsection*{Adhyāya 11}
\pstart \ledsidenote{\passref{4.11.1}}\skt{अथातः }\edtext{\skt{प्रमे\-हा\-णा\-ञ्चि\-कि\-त्सि\-त\-म्व्या\-ख्या\-स्यामः}}{\lemma{\skt{प्रमेहाणाञ्}}\Afootnote{H; \skt{प्रमेह}\ A}}\skt{~॥}\pend

\pstart \ledsidenote{\passref{4.11.3}}\skt{द्वौ }\edtext{\skt{प्रमहौ}}{\lemma{\skt{प्रमहौ}}\Afootnote{H; \skt{प्रमेहौ}\ A}}\skt{ भवतः }\edtext{\edtext{\skt{सह\-जोऽ\-प\-थ्य\-नि\-मि\-त्तश्च}}{\lemma{\skt{ऽप\appbrk{}थ्य\appbrk{}नि\appbrk{}मि\appbrk{}त्तश्}}\Afootnote{H; \skt{ऽपथ्यनिमित्त}\ A}}}{\lemma{\skt{च}}\Afootnote{H; \skt{श्च}\ A}}\skt{~। तत्र सहजो }\edtext{\skt{नाम}}{\lemma{\skt{नाम}}\Afootnote{H; \emph{om.}\ A}}\skt{ }\edtext{\skt{माता\-पि\-तृ\-बी\-ज\-दो\-ष\-कृतः}}{\lemma{\skt{माता\appbrk{}पि\appbrk{}तृ\appbrk{}बी\appbrk{}ज\appbrk{}दो\appbrk{}ष\appbrk{}कृतः}}\Afootnote{H; \skt{मातृ\appbrk{}पि\appbrk{}तृ\appbrk{}बी\appbrk{}ज\appbrk{}दो\appbrk{}ष\appbrk{}कृतः}\ A}}\skt{~॥ }\edtext{\skt{अहि\-ता\-हा\-र\-नि\-मि\-त्तेऽ\-प\-थ्य\-नि\-मित्तः}}{\lemma{\skt{अहि\appbrk{}ता\appbrk{}हा\appbrk{}र\appbrk{}नि\appbrk{}मित्ते}}\Afootnote{H; \skt{अहिताहारजो}\ A}}\skt{~। तयोः पूर्व्वेणोपद्रुतः }\edtext{\skt{कृषो}}{\lemma{\skt{कृषो}}\Afootnote{H; \skt{कृशो}\ A}}\skt{ }\edtext{\edtext{\skt{रूक्षोऽ\-ल्पा\-शी\-भृ\-श\-म्प\-र्य\-ट\-न\-शी\-लश्च}}{\lemma{\skt{ऽल्पाशीभृशम्}}\Afootnote{H; \skt{ऽल्पाशी पिपासुर्भृशं}\ A}}}{\lemma{\skt{पर्य\appbrk{}ट\appbrk{}न\appbrk{}शी\appbrk{}लश्च}}\Afootnote{H; \skt{परि\appbrk{}स\appbrk{}र\appbrk{}ण\appbrk{}शी\appbrk{}लश्च}\ A}}\skt{ }\edtext{\skt{भवत्युत्तरेण}}{\lemma{\skt{भवत्युत्तरेण}}\Afootnote{H; \skt{भवति; उत्तरेण}\ A}}\skt{ स्थूलो }\edtext{\skt{वह्वाशी}}{\lemma{\skt{वह्वाशी}}\Afootnote{H; \skt{बह्वाई}\ A}}\skt{ स्निग्धः }\edtext{\skt{शय्यास्वप्नशीलः}}{\lemma{\skt{शय्यास्वप्नशीलः}}\Afootnote{H; \skt{शय्या\appbrk{}स\appbrk{}न\appbrk{}स्व\appbrk{}प्न\appbrk{}शीलः}\ A}}\skt{ प्रायेणेति~।}\pend

\pstart \ledsidenote{\passref{4.11.4}}\skt{तत्र }\edtext{\edtext{\skt{कृष\-श\-म\-न्न\-पा\-न\-सं\-स्कृ\-ताभिः}}{\lemma{\skt{कृषशम्}}\Afootnote{H\textsuperscript{pc}; \skt{कृशमन्}\ A; \skt{कृषम्}\ H\textsuperscript{ac}}}}{\lemma{\skt{अन्न\appbrk{}पा\appbrk{}न\appbrk{}सं\appbrk{}स्कृ\appbrk{}ताभिः}}\Afootnote{H; \skt{नपा\appbrk{}न\appbrk{}प्र\appbrk{}ति\appbrk{}सं\appbrk{}स्कृ\appbrk{}ताभिः}\ A}}\skt{ }\edtext{\skt{क्रियाभिश्च}}{\lemma{\skt{क्रियाभिश्च}}\Afootnote{H; \skt{क्रियाभिश्}\ A}}\skt{ }\edtext{\skt{चिकित्सेत्}}{\lemma{\skt{चिकित्सेत्}}\Afootnote{H; \skt{चिकित्सेत}\ A}}\skt{~। स्थूल\-म\-प\-त\-र्प्प\-ण\-यु\-क्ताभिः~।}\pend

\pstart \ledsidenote{\passref{4.11.5}}\skt{सर्व्व }\edtext{\skt{एव}}{\lemma{\skt{एव}}\Afootnote{H; \skt{एव च}\ A}}\skt{ परिहरेयुः }\edtext{\skt{सौवी\-र\-क\-तु\-षो\-द\-क\-शु\-क्त\-मै\-रे\-य\-सु\-रा\-स\-व\-प\-य\-स्तै\-ल\-गु\-ड\-घृ\-ते\-क्षु\-र\-स\-पि\-ष्टा\-म्ल\-पा\-न\-कानि}}{\lemma{\skt{सौवी\appbrk{}र\appbrk{}क\appbrk{}तु\appbrk{}षो\appbrk{}द\appbrk{}क\appbrk{}शु\appbrk{}क्त\appbrk{}मै\appbrk{}रे\appbrk{}य\appbrk{}सु\appbrk{}रा\appbrk{}स\appbrk{}व\appbrk{}प\appbrk{}य\appbrk{}स्तै\appbrk{}ल\appbrk{}गु\appbrk{}ड\appbrk{}घृ\appbrk{}ते\appbrk{}क्षु\appbrk{}र\appbrk{}स\appbrk{}पि\appbrk{}ष्टा\appbrk{}म्ल\appbrk{}पा\appbrk{}न\appbrk{}कानि}}\Afootnote{H; \skt{सौवी\appbrk{}र\appbrk{}क\appbrk{}तु\appbrk{}षो\appbrk{}द\appbrk{}क\appbrk{}शु\appbrk{}क्त\appbrk{}मै\appbrk{}रे\appbrk{}य\appbrk{}सु\appbrk{}रा\appbrk{}स\appbrk{}व\appbrk{}तो\appbrk{}य\appbrk{}प\appbrk{}य\appbrk{}स्तै\appbrk{}ल\appbrk{}घृ\appbrk{}ते\appbrk{}क्षु\appbrk{}वि\appbrk{}का\appbrk{}र\appbrk{}द\appbrk{}धि\appbrk{}पि\appbrk{}ष्टा\appbrk{}न्ना\appbrk{}म्ल\appbrk{}य\appbrk{}वा\appbrk{}गू\appbrk{}णा\appbrk{}न\appbrk{}कानि}\ A}}\skt{ }\edtext{\edtext{\skt{ग्राम्या\-नू\-पो\-द\-का\-ना\-म्मां\-सा\-नीति}}{\lemma{\skt{ग्राम्या\appbrk{}नू\appbrk{}पो\appbrk{}द\appbrk{}का\appbrk{}नाम्}}\Afootnote{H; \skt{ग्राम्या\appbrk{}नू\appbrk{}पौ\appbrk{}द\appbrk{}कमां}\ A}}}{\lemma{\skt{मांसानीति}}\Afootnote{H; \skt{सानि चेति}\ A}}\skt{~॥}\pend

\pstart \ledsidenote{\passref{4.11.6}}\skt{ततः }\edtext{\skt{शालि\-ष\-ष्टि\-क\-श्या\-मा\-क\-यव}}{\lemma{\skt{शालि\appbrk{}ष\appbrk{}ष्टि\appbrk{}क\appbrk{}श्या\appbrk{}मा\appbrk{}क\appbrk{}यव}}\Afootnote{H\textsuperscript{pc}; \skt{शालि\appbrk{}ष\appbrk{}ष्टि\appbrk{}क\appbrk{}यव}\ A; \skt{शालि\appbrk{}ष\appbrk{}ष्टि\appbrk{}क\appbrk{}श्या\appbrk{}माक}\ H\textsuperscript{ac}}}\skt{ ३ गोधू\-म\-को\-द्र\-वो\-द्दा\-ल\-का\-न\-न\-वा\-न्भु\-ञ्जीत }\edtext{\skt{चन\-क\-तु\-व\-रि\-क\-कु\-ल\-त्थ\-मु\-द्ग\-वि\-क\-ल्पेन}}{\lemma{\skt{चन\appbrk{}क\appbrk{}तु\appbrk{}व\appbrk{}रि\appbrk{}क\appbrk{}कु\appbrk{}ल\appbrk{}त्थ\appbrk{}मु\appbrk{}द्ग\appbrk{}वि\appbrk{}क\appbrk{}ल्पेन}}\Afootnote{H; \skt{चण\appbrk{}का\appbrk{}ढ\appbrk{}की\appbrk{}कु\appbrk{}ल\appbrk{}त्थ\appbrk{}मु\appbrk{}द्ग\appbrk{}वि\appbrk{}क\appbrk{}ल्पेन}\ A}}\skt{ तिक्तकषायाभ्यां च शाक\-ग\-णा\-भ्या\-न्नि\-कु\-म्भे\-ङ्गु\-दी\-स\-र्ष\-पा\-त\-सी\-तै\-ल\-सि\-द्धाभ्यां }\edtext{\skt{वस्तुमूत्रैर्व्वा}}{\lemma{\skt{वस्तुमूत्रैर्व्वा}}\Afootnote{H; \skt{बद्धमूत्रैर्वा}\ A}}\skt{ }\edtext{\edtext{\skt{जाङ्ग\-लै\-र्म्मां\-स\-र\-सै\-र\-न\-म्ल\-घृ\-त\-मे\-दोभिः}}{\lemma{\skt{म्मांसरसैर्}}\Afootnote{H; \skt{मांसैर्}\ A}}}{\lemma{\skt{अन\appbrk{}म्ल\appbrk{}घृ\appbrk{}त\appbrk{}मे\appbrk{}दोभिः}}\Afootnote{H; \skt{अप\appbrk{}हृ\appbrk{}त\appbrk{}मे\appbrk{}दो\appbrk{}भि\appbrk{}र\appbrk{}न\appbrk{}म्लै\appbrk{}र\appbrk{}घृ\appbrk{}तै\appbrk{}श्चेति}\ A}}\pend

\pstart \ledsidenote{\passref{4.11.7}}\skt{तत्रादित एव }\edtext{\skt{प्रमे\-हि\-ण\-मु\-प\-स्नि\-ग्ध\-म\-न्य\-तमे}}{\lemma{\skt{उपस्निग्धम्}}\Afootnote{H; \skt{स्निग्धम्}\ A}}\skt{ न तैलेन प्रियङ्ग्वादि सिद्धेन वा घृतेन वामयेत्~। }\edtext{\skt{प्रगा\-ढ\-म्वि\-रे\-च\-येत्}}{\lemma{\skt{विरेचयेत्}}\Afootnote{H; \skt{विरेचयेच्च}\ A}}\skt{~। }\edtext{\skt{विरे\-च\-ना\-न\-न्तरं}}{\lemma{\skt{विरे\appbrk{}च\appbrk{}ना\appbrk{}न\appbrk{}न्तरं}}\Afootnote{H; \skt{विरे\appbrk{}च\appbrk{}ना\appbrk{}द\appbrk{}न\appbrk{}न्तरं}\ A}}\skt{ }\edtext{\skt{सुर\-सा\-दि\-क\-षा\-ये\-णा\-स्था\-प\-येत्}}{\lemma{\skt{सुर\appbrk{}सा\appbrk{}दि\appbrk{}क\appbrk{}षा\appbrk{}ये\appbrk{}णा\appbrk{}स्था\appbrk{}प\appbrk{}येत्}}\Afootnote{H; \skt{सुर\appbrk{}सा\appbrk{}दि\appbrk{}क\appbrk{}षा\appbrk{}ये\appbrk{}णा\appbrk{}स्था\appbrk{}प\appbrk{}येन्}\ A}}\skt{~। }\edtext{\skt{महौ\-ष\-ध\-भ\-द्र\-दा\-रु\-मु\-स्ता\-वा\-येन}}{\lemma{\skt{महौ\appbrk{}ष\appbrk{}ध\appbrk{}भ\appbrk{}द्र\appbrk{}दा\appbrk{}रु\appbrk{}मु\appbrk{}स्ता\appbrk{}वा\appbrk{}येन}}\Afootnote{H; \skt{महौ\appbrk{}ष\appbrk{}ध\appbrk{}भ\appbrk{}द्र\appbrk{}दा\appbrk{}रु\appbrk{}मु\appbrk{}स्ता\appbrk{}वा\appbrk{}पेन}\ A}}\skt{ मधु\-सै\-न्ध\-व\-यु\-क्तेन }\edtext{\skt{तह्यमानम्वा}}{\lemma{\skt{तह्यमानम्वा}}\Afootnote{H; \skt{दह्यमानं च}\ A}}\skt{ न्यग्रो\-धा\-दि\-क\-षा\-येण निस्तैलेन~।}\pend

\pstart \ledsidenote{\passref{4.11.8}}\edtext{\skt{शुद्ध\-दे\-हा\-मा\-म\-ल\-क\-र\-सेन}}{\lemma{\skt{शुद्ध\appbrk{}दे\appbrk{}हा\appbrk{}मा\appbrk{}म\appbrk{}ल\appbrk{}क\appbrk{}र\appbrk{}सेन}}\Afootnote{H; \skt{ततः शुद्ध\appbrk{}दे\appbrk{}ह\appbrk{}मा\appbrk{}म\appbrk{}ल\appbrk{}क\appbrk{}र\appbrk{}सेन}\ A}}\skt{ हरि\-द्रा\-म्म\-धु\-सं\-यु\-क्ता\-म्पा\-य\-येत्~। }\edtext{\skt{त्रिफ\-ला\-दे\-व\-दा\-रु\-मु\-स्त\-क\-क\-षा\-यम्वा}}{\lemma{\skt{त्रिफ\appbrk{}ला\appbrk{}दे\appbrk{}व\appbrk{}दा\appbrk{}रु\appbrk{}मु\appbrk{}स्त\appbrk{}क\appbrk{}क\appbrk{}षा\appbrk{}यम्वा}}\Afootnote{H; \skt{त्रिफ\appbrk{}ला\appbrk{}वि\appbrk{}शा\appbrk{}ला\appbrk{}दे\appbrk{}व\appbrk{}दा\appbrk{}रु\appbrk{}मु\appbrk{}स्त\appbrk{}क\appbrk{}षायं वा}\ A}}\skt{~। शालकम्पि }\edtext{\skt{ल्यक\-मु\-ष्क\-क\-क\-ल्क\-म\-क्ष\-मात्र वा}}{\lemma{\skt{ल्यक\appbrk{}मु\appbrk{}ष्क\appbrk{}क\appbrk{}क\appbrk{}ल्क\appbrk{}म\appbrk{}क्ष\appbrk{}मात्र वा}}\Afootnote{H; \skt{ल्लक\appbrk{}मु\appbrk{}ष्क\appbrk{}क\appbrk{}क\appbrk{}ल्क\appbrk{}म\appbrk{}क्ष\appbrk{}मात्रं वा}\ A}}\skt{ }\edtext{\edtext{\skt{मधु\-म\-धु\-र\-मा\-म\-ल\-क\-र\-सेन}}{\lemma{\skt{मधुमधुरम्}}\Afootnote{H; \skt{मधुमधुर}\ A}}}{\lemma{\skt{आम\appbrk{}ल\appbrk{}क\appbrk{}र\appbrk{}सेन}}\Afootnote{H; \skt{माम\appbrk{}ल\appbrk{}क\appbrk{}र\appbrk{}सेन}\ A}}\skt{~। }\edtext{\skt{कुट\-ज\-क\-पि\-त्थ\-रो\-ही\-त\-क\-वि\-भी\-त\-क\-स\-प्त\-प\-र्ण्ण\-पु\-ष्य\-क\-ल्कम्वा}}{\lemma{\skt{कुट\appbrk{}ज\appbrk{}क\appbrk{}पि\appbrk{}त्थ\appbrk{}रो\appbrk{}ही\appbrk{}त\appbrk{}क\appbrk{}वि\appbrk{}भी\appbrk{}त\appbrk{}क\appbrk{}स\appbrk{}प्त\appbrk{}प\appbrk{}र्ण्ण\appbrk{}पु\appbrk{}ष्य\appbrk{}क\appbrk{}ल्कम्वा}}\Afootnote{H\textsuperscript{pc}; \skt{हरिद्रायुतं, कुट\appbrk{}ज\appbrk{}क\appbrk{}पि\appbrk{}त्थ\appbrk{}रो\appbrk{}ही\appbrk{}त\appbrk{}क\appbrk{}बि\appbrk{}भी\appbrk{}त\appbrk{}क\appbrk{}स\appbrk{}प्त\appbrk{}प\appbrk{}र्ण\appbrk{}पु\appbrk{}ष्प\appbrk{}कल्कं वा}\ A; \skt{कुट\appbrk{}ज\appbrk{}क\appbrk{}पि\appbrk{}त्थ\appbrk{}रो\appbrk{}ही\appbrk{}त\appbrk{}क\appbrk{}वि\appbrk{}भी\appbrk{}त\appbrk{}क\appbrk{}स\appbrk{}प्त\appbrk{}प\appbrk{}र्ण्ण\appbrk{}क\appbrk{}ल्कम्वा}\ H\textsuperscript{ac}}}\skt{~। }\edtext{\skt{निम्वा\-र\-ग्व\-ध\-स\-प्त\-प\-र्ण्ण\-पू\-र्व्वा\-कु\-ट\-ज\-सो\-म\-वृ\-क्ष\-प\-ला\-शा\-नाम्वा}}{\lemma{\skt{निम्वा\appbrk{}र\appbrk{}ग्व\appbrk{}ध\appbrk{}स\appbrk{}प्त\appbrk{}प\appbrk{}र्ण्ण\appbrk{}पू\appbrk{}र्व्वा\appbrk{}कु\appbrk{}ट\appbrk{}ज\appbrk{}सो\appbrk{}म\appbrk{}वृ\appbrk{}क्ष\appbrk{}प\appbrk{}ला\appbrk{}शा\appbrk{}नाम्वा}}\Afootnote{H; \skt{निम्बा\appbrk{}र\appbrk{}ग्व\appbrk{}ध\appbrk{}स\appbrk{}प्त\appbrk{}प\appbrk{}र्ण\appbrk{}मू\appbrk{}र्वा\appbrk{}कु\appbrk{}ट\appbrk{}ज\appbrk{}सो\appbrk{}म\appbrk{}वृ\appbrk{}क्ष\appbrk{}प\appbrk{}ला\appbrk{}शानां वा}\ A}}\skt{~। }\edtext{\skt{त्वक्प\-त्त्र\-पु\-ष्प\-फ\-ल\-मू\-ल\-क\-षा\-या\-ण्येते}}{\lemma{\skt{त्वक्प\appbrk{}त्त्र\appbrk{}पु\appbrk{}ष्प\appbrk{}फ\appbrk{}ल\appbrk{}मू\appbrk{}ल\appbrk{}क\appbrk{}षा\appbrk{}याण्य्}}\Afootnote{H; \skt{त्वक्प\appbrk{}त्र\appbrk{}मू\appbrk{}ल\appbrk{}फ\appbrk{}ल\appbrk{}पु\appbrk{}ष्प\appbrk{}क\appbrk{}षा\appbrk{}याणि}\ A}}\skt{ पञ्च }\edtext{\skt{प्रयोगाः}}{\lemma{\skt{प्रयोगाः}}\Afootnote{H; \skt{योगाः}\ A}}\skt{ }\edtext{\skt{सर्व्व\-मे\-हो\-प\-ह\-न्तारो}}{\lemma{\skt{सर्व्व\appbrk{}मे\appbrk{}हो\appbrk{}प\appbrk{}ह\appbrk{}न्तारो}}\Afootnote{H; \skt{सर्व\appbrk{}मे\appbrk{}हा\appbrk{}ना\appbrk{}म\appbrk{}प\appbrk{}ह\appbrk{}न्तारो}\ A}}\skt{ व्याख्याताः~।}\pend

\pstart \ledsidenote{\passref{4.11.9}}\edtext{\skt{विशेषत}}{\lemma{\skt{विशेषत}}\Afootnote{H; \skt{विशेषश्चात}\ A}}\skt{ ऊर्द्ध्वं~। }\edtext{\skt{तत्रो\-द\-क\-मे\-हि\-न\-म्पा\-रि\-जा\-त\-क\-षा\-य\-म्पा\-य\-येत्~। इक्षु\-वा\-लि\-का\-मे\-हिनां}}{\lemma{\skt{पाययेत्~। इक्षु\appbrk{}वा\appbrk{}लि\appbrk{}का\appbrk{}मे\appbrk{}हिनां}}\Afootnote{H; \skt{पाययेत्, इक्षु}\ A}}\skt{ }\edtext{\skt{वैज\-य\-न्ती\-क\-षायं}}{\lemma{\skt{वैज\appbrk{}य\appbrk{}न्ती\appbrk{}क\appbrk{}षायं}}\Afootnote{H; \emph{om.}\ A}}\skt{~। }\edtext{\edtext{\skt{सुरा\-मे\-हि\-ना\-न्नि\-म्व\-क\-षायं}}{\lemma{\skt{सुरामेहिनान्}}\Afootnote{H; \emph{om.}\ A}}}{\lemma{\skt{निम्वकषायं}}\Afootnote{H; \emph{om.}\ A}}\skt{~। }\edtext{\skt{सिकतामेहिनं}}{\lemma{\skt{सिकतामेहिनं}}\Afootnote{H; \skt{मेहिनं}\ A}}\skt{ }\edtext{\skt{चित्त्रकषायं}}{\lemma{\skt{चित्त्रकषायं}}\Afootnote{H; \skt{चित्रककषायं}\ A}}\skt{~। शनैर्म्मेहिनं खदिरकषायं~। }\edtext{\skt{लव\-ण\-मे\-हि\-न\-म्पा\-ठा\-ग\-रु\-क\-षायं}}{\lemma{\skt{पाठा\appbrk{}ग\appbrk{}रु\appbrk{}क\appbrk{}षायं}}\Afootnote{H; \skt{पाठा+\appbrk{}अ\appbrk{}गु\appbrk{}रु\appbrk{}ह\appbrk{}रि\appbrk{}द्रा\appbrk{}क\appbrk{}षायं}\ A}}\skt{~। पिष्टमेहिनं }\edtext{\skt{रजन्योः}}{\lemma{\skt{रजन्योः}}\Afootnote{H; \skt{हरि\appbrk{}द्रा\appbrk{}दा\appbrk{}रु\appbrk{}ह\appbrk{}रिद्रा}\ A}}\skt{ कषायं~। सान्द्रमेहिनं सप्त\-प\-र्ण्ण\-क\-षायं~। शुक्रमेहिनं }\edtext{\skt{दूर्व्वा\-शै\-वा\-ल\-प्ल\-व\-ह\-ठ\-क\-र\-ञ्च\-क\-से\-रु\-क\-क\-षायं}}{\lemma{\skt{दूर्व्वा\appbrk{}शै\appbrk{}वा\appbrk{}ल\appbrk{}प्ल\appbrk{}व\appbrk{}ह\appbrk{}ठ\appbrk{}क\appbrk{}र\appbrk{}ञ्च\appbrk{}क\appbrk{}से\appbrk{}रु\appbrk{}क\appbrk{}क\appbrk{}षायं}}\Afootnote{H; \skt{दूर्वा\appbrk{}शै\appbrk{}व\appbrk{}ल\appbrk{}प्ल\appbrk{}व\appbrk{}ह\appbrk{}ठ\appbrk{}क\appbrk{}र\appbrk{}ञ्ज\appbrk{}क\appbrk{}से\appbrk{}रु\appbrk{}क\appbrk{}क\appbrk{}षायं ककु\appbrk{}भ\appbrk{}च\appbrk{}न्द\appbrk{}न\appbrk{}क\appbrk{}षायं वा}\ A}}\skt{~। फेनमेहिनं }\edtext{\skt{त्रिफ\-ला\-र\-ग्व\-ध\-मृ\-द्वी\-का\-क\-षायं}}{\lemma{\skt{त्रिफ\appbrk{}ला\appbrk{}र\appbrk{}ग्व\appbrk{}ध\appbrk{}मृ\appbrk{}द्वी\appbrk{}का\appbrk{}क\appbrk{}षायं}}\Afootnote{H; \skt{त्रिफ\appbrk{}ला\appbrk{}र\appbrk{}ग्व\appbrk{}घ\appbrk{}मृ\appbrk{}द्वी\appbrk{}का\appbrk{}क\appbrk{}षायं}\ A}}\skt{ मधुमधुरं~। }\edtext{\skt{पैत्तिकेषु}}{\lemma{\skt{पैत्तिकेषु}}\Afootnote{H; \skt{इति; पैत्तिकेषु}\ A}}\skt{ }\edtext{\skt{नील\-मे\-हि\-न\-म\-श्व\-त्थ\-क\-षायं}}{\lemma{\skt{अश्वत्थकषायं}}\Afootnote{H; \skt{शाल\appbrk{}सा\appbrk{}रा\appbrk{}दि\appbrk{}क\appbrk{}षा\appbrk{}य\appbrk{}म\appbrk{}श्व\appbrk{}त्थ\appbrk{}क\appbrk{}षायं वा}\ A}}\skt{~। }\edtext{\skt{हरिद्रामेहिनं}}{\lemma{\skt{हरिद्रामेहिनं}}\Afootnote{H; \skt{पाययेत्, हरिद्रामेहिनं}\ A}}\skt{ राज\-वृ\-क्ष\-क\-षायं~। }\edtext{\skt{अम्व्लमेहिनं}}{\lemma{\skt{अम्व्लमेहिनं}}\Afootnote{H; \skt{अम्लमेहिनं}\ A}}\skt{ न्य ग्रोधादिकषायं~। क्षारमेहिनं त्रिफलाकषायं~। मञ्जिष्ठामेहिनं मञ्जि\-ष्ठा\-च\-न्द\-न\-क\-षायं~। शोणितमेहिनं }\edtext{\edtext{\skt{गुडू\-ची\-ति\-न्दु\-का\-स्थि\-का\-स्म\-र्य\-ख\-र्ज्जू\-र\-क\-षा\-य\-म्म\-धु\-मिश्र}}{\lemma{\skt{गुडू\appbrk{}ची\appbrk{}ति\appbrk{}न्दु\appbrk{}का\appbrk{}स्थि\appbrk{}का\appbrk{}स्म\appbrk{}र्य\appbrk{}ख\appbrk{}र्ज्जू\appbrk{}र\appbrk{}क\appbrk{}षा\appbrk{}यम्}}\Afootnote{H; \skt{गुडू\appbrk{}ची\appbrk{}ति\appbrk{}न्दु\appbrk{}का\appbrk{}स्थि\appbrk{}का\appbrk{}श्म\appbrk{}र्य\appbrk{}ख\appbrk{}र्जू\appbrk{}र\appbrk{}क\appbrk{}षायं}\ A}}}{\lemma{\skt{मधुमिश्र}}\Afootnote{H; \skt{मधुमिश्रं}\ A}}\skt{ }\edtext{\skt{पाययेत्~॥ अत}}{\lemma{\skt{पाययेत्~॥ अत}}\Afootnote{H; \skt{अत}\ A}}\skt{ }\edtext{\edtext{\edtext{\skt{ऊर्द्ध्व\-म\-सा\-ध्ये\-ष्वपि}}{\lemma{\skt{ऊर्द्ध्वम्}}\Afootnote{H; \skt{ऊर्ध्व}\ A}}}{\lemma{\skt{असाध्येष्व्}}\Afootnote{H; \skt{मसाध्येष्}\ A}}}{\lemma{\skt{अपि}}\Afootnote{H; \skt{वपि}\ A}}\skt{ }\edtext{\skt{योगा\-न्या\-प\-ना\-र्था\-न्व\-क्ष्यामः}}{\lemma{\skt{यापनार्थान्}}\Afootnote{H; \skt{यापनार्थं}\ A}}\skt{~। }\edtext{\skt{तद्यथा}}{\lemma{\skt{तद्}}\Afootnote{H; \skt{तद्}\ A}}\skt{~। }\edtext{\edtext{\skt{सर्प्पिर्म्मेहिनं}}{\lemma{\skt{सर्प्पिर्}}\Afootnote{H; \skt{सर्पि}\ A}}}{\lemma{\skt{म्मेहिनं}}\Afootnote{H; \skt{र्मेहिनं}\ A}}\skt{ }\edtext{\skt{कुष्ठ\-कु\-ट\-ज\-पा\-ठा\-हि\-ङ्गु\-त्रि\-क\-टु\-क\-रो\-हिणी}}{\lemma{\skt{कुष्ठ\appbrk{}कु\appbrk{}ट\appbrk{}ज\appbrk{}पा\appbrk{}ठा\appbrk{}हि\appbrk{}ङ्गु\appbrk{}त्रि\appbrk{}क\appbrk{}टु\appbrk{}क\appbrk{}रो\appbrk{}हिणी}}\Afootnote{H; \skt{कुष्ठ\appbrk{}कु\appbrk{}ट\appbrk{}ज\appbrk{}पा\appbrk{}ठा\appbrk{}हि\appbrk{}ङ्गु\appbrk{}क\appbrk{}टु\appbrk{}रो\appbrk{}हिणी}\ A}}\skt{ कल्कं गुडू\-ची\-चि\-त्त्र\-क\-क\-षा\-येण पाययेत्~। }\edtext{\edtext{\skt{वसा\-मे\-हि\-न\-म\-ग्नि\-म\-न्थ\-क\-षायं}}{\lemma{\skt{वसामेहिनम्}}\Afootnote{H; \skt{वसामेहिन}\ A}}}{\lemma{\skt{अग्नि\appbrk{}म\appbrk{}न्थ\appbrk{}क\appbrk{}षायं}}\Afootnote{H; \skt{मग्नि\appbrk{}म\appbrk{}न्थ\appbrk{}क\appbrk{}षायं शिंशपाकषायं वा}\ A}}\skt{~। }\edtext{\skt{क्षौद्र\-मे\-हि\-न\-ङ्क\-द\-र\-क\-क्र\-मु\-क\-क\-षायं}}{\lemma{\skt{कद\appbrk{}र\appbrk{}क\appbrk{}क्र\appbrk{}मु\appbrk{}क\appbrk{}क\appbrk{}षायं}}\Afootnote{H; \skt{कद\appbrk{}र\appbrk{}क्र\appbrk{}मु\appbrk{}क\appbrk{}क\appbrk{}षायं}\ A}}\skt{~। }\edtext{\skt{हस्ति\-मे\-हि\-न\-न्ति\-न्दु\-क\-क\-पि\-त्थ\-शि\-री\-ष\-प\-ला\-श\-पा\-ठा\-मू\-र्वा\-दु\-ष्य\-र्शा\-क\-षायं}}{\lemma{\skt{तिन्दु\appbrk{}क\appbrk{}क\appbrk{}पि\appbrk{}त्थ\appbrk{}शि\appbrk{}री\appbrk{}ष\appbrk{}प\appbrk{}ला\appbrk{}श\appbrk{}पा\appbrk{}ठा\appbrk{}मू\appbrk{}र्वा\appbrk{}दु\appbrk{}ष्य\appbrk{}र्शा\appbrk{}क\appbrk{}षायं}}\Afootnote{H; \skt{तिन्दु\appbrk{}क\appbrk{}क\appbrk{}पि\appbrk{}त्थ\appbrk{}शि\appbrk{}री\appbrk{}ष\appbrk{}प\appbrk{}ला\appbrk{}श\appbrk{}पा\appbrk{}ठा\appbrk{}मू\appbrk{}र्वादुः}\ A}}\skt{~। }\edtext{\skt{हस्त्य\-श्व\-शू\-क\-र\-ख\-रो\-ष्ट्रा\-स्थि\-क्षार}}{\lemma{\skt{हस्त्य\appbrk{}श्व\appbrk{}शू\appbrk{}क\appbrk{}र\appbrk{}ख\appbrk{}रो\appbrk{}ष्ट्रा\appbrk{}स्थि\appbrk{}क्षार}}\Afootnote{H; \skt{स्थिक्षारं}\ A}}\skt{ }\edtext{\skt{वा}}{\lemma{\skt{वा}}\Afootnote{H; \skt{चेति}\ A}}\skt{ }\edtext{\edtext{\skt{दह्य\-मा\-न\-मौ\-द\-क\-क\-न्द\-क्वा\-थ\-सिद्धं}}{\lemma{\skt{दह्यमानम्}}\Afootnote{H; \skt{दह्यमान}\ A}}}{\lemma{\skt{औद\appbrk{}क\appbrk{}क\appbrk{}न्द\appbrk{}क्वा\appbrk{}थ\appbrk{}सिद्धं}}\Afootnote{H; \skt{मौद\appbrk{}क\appbrk{}क\appbrk{}न्द\appbrk{}क्वा\appbrk{}थ\appbrk{}सिद्धां}\ A}}\skt{ यवागूं }\edtext{\skt{क्षीरे\-क्षु\-र\-स\-म\-धु\-म\-धु\-रा\-म्पा\-य\-येत्}}{\lemma{\skt{क्षीरे\appbrk{}क्षु\appbrk{}र\appbrk{}स\appbrk{}म\appbrk{}धु\appbrk{}म\appbrk{}धु\appbrk{}राम्}}\Afootnote{H; \skt{क्षीरे\appbrk{}क्षु\appbrk{}र\appbrk{}स\appbrk{}म\appbrk{}धुरां}\ A}}\skt{~।}\pend

\pstart \ledsidenote{\passref{4.11.10}}\skt{ततः }\edtext{\skt{प्रियं\-ग्व\-न\-न्ता\-घृ\-षि\-का\-प\-द्मा\-तु\-र्य\-न्ती\-लो\-हि\-ति\-का\-म्व\-ष्ठा\-दा\-डि\-म\-त्वक्}}{\lemma{\skt{प्रियं\appbrk{}ग्व\appbrk{}न\appbrk{}न्ता\appbrk{}घृ\appbrk{}षि\appbrk{}का\appbrk{}प\appbrk{}द्मा\appbrk{}तु\appbrk{}र्य\appbrk{}न्ती\appbrk{}लो\appbrk{}हि\appbrk{}ति\appbrk{}का\appbrk{}म्व\appbrk{}ष्ठा\appbrk{}दा\appbrk{}डि\appbrk{}म\appbrk{}त्वक्}}\Afootnote{H; \skt{प्रिय\appbrk{}ङ्ग्व\appbrk{}न\appbrk{}न्ता\appbrk{}यू\appbrk{}थि\appbrk{}का\appbrk{}प\appbrk{}द्मा\appbrk{}त्रा\appbrk{}य\appbrk{}न्ति\appbrk{}का\appbrk{}लो\appbrk{}हि\appbrk{}ति\appbrk{}का\appbrk{}म्ब\appbrk{}ष्ठा\appbrk{}दा\appbrk{}डि\appbrk{}म\appbrk{}त्वक्}\ A}}\skt{~। शालपर्ण्णी }\edtext{\skt{पद्म\-के\-स\-र\-धा\-त\-कीव}}{\lemma{\skt{पद्म\appbrk{}के\appbrk{}स\appbrk{}र\appbrk{}धा\appbrk{}त\appbrk{}कीव}}\Afootnote{H; \skt{पद्म\appbrk{}तु\appbrk{}ङ्ग\appbrk{}के\appbrk{}श\appbrk{}र\appbrk{}धा\appbrk{}त\appbrk{}कीव}\ A}}\skt{ }\edtext{\skt{कुल\-शा\-ल्म\-ली\-श्री\-वे\-ष्ट\-क\-मो\-च\-र\-से\-ष्व\-य\-स्कृ\-ती\-ले\-हा\-ना\-स\-वा\-न्कु\-र्व्वीत}}{\lemma{\skt{अय\appbrk{}स्कृ\appbrk{}ती\appbrk{}ले\appbrk{}हा\appbrk{}ना\appbrk{}स\appbrk{}वान्}}\Afootnote{H; \skt{अरि\appbrk{}ष्टा\appbrk{}न\appbrk{}य\appbrk{}स्कृ\appbrk{}ती\appbrk{}र्ले\appbrk{}हा\appbrk{}ना\appbrk{}स\appbrk{}वांश्च}\ A}}\skt{~। }\edtext{\skt{शृङ्गा\-ट\-क\-क\-लो\-टी\-वि\-स\-मृ\-णा\-ल\-क\-से\-रु\-क\-म\-धु\-का\-म्र\-ज\-म्व्वा\-स\-नति}}{\lemma{\skt{शृङ्गा\appbrk{}ट\appbrk{}क\appbrk{}क\appbrk{}लो\appbrk{}टी\appbrk{}वि\appbrk{}स\appbrk{}मृ\appbrk{}णा\appbrk{}ल\appbrk{}क\appbrk{}से\appbrk{}रु\appbrk{}क\appbrk{}म\appbrk{}धु\appbrk{}का\appbrk{}म्र\appbrk{}ज\appbrk{}म्व्वा\appbrk{}स\appbrk{}नति}}\Afootnote{H; \skt{शृङ्गा\appbrk{}ट\appbrk{}क\appbrk{}गि\appbrk{}लो\appbrk{}ड्य\appbrk{}वि\appbrk{}स\appbrk{}मृ\appbrk{}णा\appbrk{}ल\appbrk{}का\appbrk{}श\appbrk{}क\appbrk{}से\appbrk{}रु\appbrk{}क\appbrk{}म\appbrk{}धु\appbrk{}का\appbrk{}म्र\appbrk{}ज\appbrk{}म्ब्व\appbrk{}स\appbrk{}नति}\ A}}\skt{ }\edtext{\skt{निश\-क\-कु\-भ\-क\-ट्व\-ङ्ग\-लो\-ध्र\-भ\-ल्ला\-त\-क\-प\-ला\-श\-च\-र्म्म\-वृ\-क्ष\-गि\-रि\-क\-र्ण्णि\-का\-शी\-त\-शि\-व\-नि\-चु\-ल\-दा\-डि\-मा\-ज\-क\-र्ण्ण\-ह\-रि\-वृ\-क्ष\-रा\-जा\-द\-न\-घो\-ण्टा\-वि\-कं\-क\-तेषु}}{\lemma{\skt{निश\appbrk{}क\appbrk{}कु\appbrk{}भ\appbrk{}क\appbrk{}ट्व\appbrk{}ङ्ग\appbrk{}लो\appbrk{}ध्र\appbrk{}भ\appbrk{}ल्ला\appbrk{}त\appbrk{}क\appbrk{}प\appbrk{}ला\appbrk{}श\appbrk{}च\appbrk{}र्म्म\appbrk{}वृ\appbrk{}क्ष\appbrk{}गि\appbrk{}रि\appbrk{}क\appbrk{}र्ण्णि\appbrk{}का\appbrk{}शी\appbrk{}त\appbrk{}शि\appbrk{}व\appbrk{}नि\appbrk{}चु\appbrk{}ल\appbrk{}दा\appbrk{}डि\appbrk{}मा\appbrk{}ज\appbrk{}क\appbrk{}र्ण्ण\appbrk{}ह\appbrk{}रि\appbrk{}वृ\appbrk{}क्ष\appbrk{}रा\appbrk{}जा\appbrk{}द\appbrk{}न\appbrk{}घो\appbrk{}ण्टा\appbrk{}वि\appbrk{}कं\appbrk{}क\appbrk{}तेषु}}\Afootnote{H; \skt{निश\appbrk{}क\appbrk{}कु\appbrk{}भ\appbrk{}क\appbrk{}ट्व\appbrk{}ङ्ग\appbrk{}रो\appbrk{}ध्र\appbrk{}भ\appbrk{}ल्ला\appbrk{}त\appbrk{}क\appbrk{}प\appbrk{}ला\appbrk{}श\appbrk{}च\appbrk{}र्म\appbrk{}वृ\appbrk{}क्ष\appbrk{}गि\appbrk{}रि\appbrk{}क\appbrk{}र्णि\appbrk{}का\appbrk{}शी\appbrk{}त\appbrk{}शि\appbrk{}व\appbrk{}नि\appbrk{}चु\appbrk{}ल\appbrk{}दा\appbrk{}डि\appbrk{}मा\appbrk{}ज\appbrk{}क\appbrk{}र्ण\appbrk{}ह\appbrk{}रि\appbrk{}वृ\appbrk{}क्ष\appbrk{}रा\appbrk{}जा\appbrk{}द\appbrk{}न\appbrk{}गो\appbrk{}प\appbrk{}घो\appbrk{}ण्टा\appbrk{}वि\appbrk{}क\appbrk{}ङ्ग\appbrk{}तेषु वा}\ A}}\skt{ यवान्न }\edtext{\skt{विका\-रा\-श्चो\-प\-से\-वेत}}{\lemma{\skt{चोपसेवेत}}\Afootnote{H; \skt{च सेवेत}\ A}}\skt{~। }\edtext{\skt{यथो\-क्त\-क\-षा\-य\-सि\-द्धम्वा}}{\lemma{\skt{यथो\appbrk{}क्त\appbrk{}क\appbrk{}षा\appbrk{}य\appbrk{}सि\appbrk{}द्धम्वा}}\Afootnote{H; \skt{यथो\appbrk{}क्त\appbrk{}क\appbrk{}षा\appbrk{}य\appbrk{}सिद्धां}\ A}}\skt{ }\edtext{\skt{चास्मै}}{\lemma{\skt{चास्मै}}\Afootnote{H; \emph{om.}\ A}}\skt{ }\edtext{\skt{यवा\-गू\-म्प्र\-य\-च्छे\-त्क\-षा\-याणि}}{\lemma{\skt{प्रयच्छेत्}}\Afootnote{H; \skt{चास्मै प्रयच्छेत्}\ A}}\skt{ वा पातुं~।}\pend

\pstart \ledsidenote{\passref{4.11.11}}\edtext{\edtext{\edtext{\skt{महा\-ध\-न\-म\-हि\-ता\-चा\-र\-मौ\-ष\-ध\-द्वे\-षि\-ण\-मी\-श्व\-रम्वा}}{\lemma{\skt{महा\appbrk{}ध\appbrk{}न\appbrk{}म\appbrk{}हि\appbrk{}ता\appbrk{}चा\appbrk{}रम्}}\Afootnote{H; \skt{महा\appbrk{}ध\appbrk{}न\appbrk{}म\appbrk{}हि\appbrk{}ता\appbrk{}हार}\ A}}}{\lemma{\skt{औष\appbrk{}ध\appbrk{}द्वे\appbrk{}षि\appbrk{}णम्}}\Afootnote{H; \skt{मौषधद्वेषिण}\ A}}}{\lemma{\skt{ईश्वरम्वा}}\Afootnote{H; \skt{मीश्वरं वा}\ A}}\skt{ }\edtext{\edtext{\edtext{\edtext{\skt{पाठा\-भ\-या\-चि\-त्र\-क\-प्र\-गा\-ढ\-म\-न\-ल्प\-मा\-क्षि\-क\-म\-न्य\-त\-म\-मा\-सवं}}{\lemma{\skt{पाठा\appbrk{}भ\appbrk{}या\appbrk{}चि\appbrk{}त्र\appbrk{}क\appbrk{}प्र\appbrk{}गा\appbrk{}ढम्}}\Afootnote{H; \skt{पाठा\appbrk{}भ\appbrk{}या\appbrk{}चि\appbrk{}त्र\appbrk{}क\appbrk{}प्र\appbrk{}गाढ}\ A}}}{\lemma{\skt{अन\appbrk{}ल्प\appbrk{}मा\appbrk{}क्षि\appbrk{}कम्}}\Afootnote{H; \skt{मनल्पमाक्षिक}\ A}}}{\lemma{\skt{अन्यतमम्}}\Afootnote{H; \skt{मन्यतम}\ A}}}{\lemma{\skt{आसवं}}\Afootnote{H; \skt{मासवं}\ A}}\skt{ }\edtext{\skt{पाययेत्~। अंगा\-र\-क\-शू\-ल्यो\-प\-दं\-शम्वा}}{\lemma{\skt{पाययेत्~। अंगा\appbrk{}र\appbrk{}क\appbrk{}शू\appbrk{}ल्यो\appbrk{}प\appbrk{}दं\appbrk{}शम्वा}}\Afootnote{H; \skt{पाययेत्, अङ्गा\appbrk{}अ\appbrk{}र\appbrk{}शू\appbrk{}ल्यो\appbrk{}प\appbrk{}दंशं वा}\ A}}\skt{~। }\edtext{\edtext{\edtext{\skt{मार्द्वी\-न्त\-म\-भी\-क्ष्ण\-म्म\-धु\-क\-पि\-त्थ\-म\-रि\-चा\-नु\-सि\-द्धानि}}{\lemma{\skt{मार्द्वीन्तम्}}\Afootnote{H; \skt{माध्वीक}\ A}}}{\lemma{\skt{अभीक्ष्णम्}}\Afootnote{H; \skt{मभीक्ष्णं}\ A}}}{\lemma{\skt{मधु\appbrk{}क\appbrk{}पि\appbrk{}त्थ\appbrk{}म\appbrk{}रि\appbrk{}चा\appbrk{}नु\appbrk{}सि\appbrk{}द्धानि}}\Afootnote{H; \skt{क्षौद्र\appbrk{}क\appbrk{}पि\appbrk{}त्थ\appbrk{}म\appbrk{}रि\appbrk{}चा\appbrk{}नु\appbrk{}विद्धा}\ A}}\skt{ }\edtext{\skt{पानकानि}}{\lemma{\skt{पानकानि}}\Afootnote{H; \skt{नि}\ A}}\skt{ चास्मै }\edtext{\skt{उपहरेत्~। उष्ट्रा\-श्व\-त\-र\-पु\-री\-ष\-चू\-र्ण्णानि}}{\lemma{\skt{उपहरेत्~। उष्ट्रा\appbrk{}श्व\appbrk{}त\appbrk{}र\appbrk{}पु\appbrk{}री\appbrk{}ष\appbrk{}चू\appbrk{}र्ण्णानि}}\Afootnote{H; \skt{पान\appbrk{}भो\appbrk{}ज\appbrk{}ना\appbrk{}न्यु\appbrk{}प\appbrk{}ह\appbrk{}रेत्, उष्ट्रा\appbrk{}श्व\appbrk{}त\appbrk{}र\appbrk{}ख\appbrk{}र\appbrk{}पु\appbrk{}री\appbrk{}ष\appbrk{}चू\appbrk{}र्णानि}\ A}}\skt{ चास्मै }\edtext{\skt{दद्यादशनेषु}}{\lemma{\skt{दद्यादशनेषु}}\Afootnote{H; \skt{दद्यादशनेषु}\ A}}\skt{~॥ }\edtext{\skt{हिङ्गुसिद्धैश्च}}{\lemma{\skt{हिङ्गुसिद्धैश्च}}\Afootnote{H; \skt{हिङ्गुसैन्धव}\ A}}\skt{ }\edtext{\skt{यूषैः}}{\lemma{\skt{यूषैः}}\Afootnote{H; \skt{युक्तैर्यूषैः}\ A}}\skt{ सार्षपैश्च }\edtext{\skt{रागै}}{\lemma{\skt{रागै}}\Afootnote{H; \skt{रागैर्}\ A}}\skt{ भोजयेत्~। अविरुद्धानि चास्मै पान\-भो\-ज\-ना\-न्यु\-प\-ह\-रेत्~। रसगन्धवन्ति च~। }\edtext{\skt{प्रवृद्धमेदस्तु}}{\lemma{\skt{प्रवृद्धमेदस्तु}}\Afootnote{H; \skt{प्रवृद्धमेहास्तु}\ A}}\skt{ व्यायामनियुद्ध }\edtext{\skt{क्रीडापहाय}}{\lemma{\skt{क्रीडापहाय}}\Afootnote{H; \skt{क्रीडागजतु}\ A}}\skt{ }\edtext{\edtext{\skt{गज\-र\-थ\-पा\-द\-च\-र्या\-प\-रि\-क्र\-म\-णा\-न्य\-स्त्रा\-पा\-स्त्रम्वा}}{\lemma{\skt{गज\appbrk{}र\appbrk{}थ\appbrk{}पा\appbrk{}द\appbrk{}च\appbrk{}र्या\appbrk{}प\appbrk{}रि\appbrk{}क्र\appbrk{}म\appbrk{}णा\appbrk{}न्यस्}}\Afootnote{H; \skt{रग\appbrk{}र\appbrk{}थ\appbrk{}प\appbrk{}दा\appbrk{}ति\appbrk{}च\appbrk{}र्या\appbrk{}प\appbrk{}रि\appbrk{}क्र\appbrk{}म\appbrk{}णान्य}\ A}}}{\lemma{\skt{त्रापास्त्रम्वा}}\Afootnote{H; \skt{स्त्रोपास्त्रे वा}\ A}}\skt{ }\edtext{\skt{सेवेत}}{\lemma{\skt{सेवेत}}\Afootnote{H; \skt{सेवेरन्}\ A}}\skt{~।}\pend

\pstart \ledsidenote{\passref{4.11.12}}\edtext{\edtext{\skt{अध\-न\-स्त्व\-वा\-न्धवो}}{\lemma{\skt{अधनस्}}\Afootnote{H; \skt{अधन}\ A}}\skt{ वा}}{\lemma{\skt{त्ववान्धवो वा}}\Afootnote{H; \skt{स्त्वबान्धवो वा}\ A}}\skt{ }\edtext{\skt{पाद\-त्र\-च्छ\-त्त्र\-र\-हितो}}{\lemma{\skt{पाद\appbrk{}त्र\appbrk{}च्छ\appbrk{}त्त्र\appbrk{}र\appbrk{}हितो}}\Afootnote{H; \skt{पाद\appbrk{}त्रा\appbrk{}णा\appbrk{}त\appbrk{}प\appbrk{}त्र\appbrk{}वि\appbrk{}र\appbrk{}हितो}\ A}}\skt{ }\edtext{\edtext{\skt{भैक्षा\-शी\-ग्रा\-मै\-क\-रा\-त्र\-वासी}}{\lemma{\skt{भैक्षाशीग्राम्}}\Afootnote{H; \skt{भैक्ष्याशी ग्रा}\ A}}}{\lemma{\skt{ऐकरात्रवासी}}\Afootnote{H; \skt{मैकरात्रवासी}\ A}}\skt{ }\edtext{\edtext{\skt{मुनिरिव}}{\lemma{\skt{मुनिर्}}\Afootnote{H; \skt{मुनि}\ A}}}{\lemma{\skt{इव}}\Afootnote{H; \skt{रिव}\ A}}\skt{ }\edtext{\skt{संयतात्मा यो}}{\lemma{\skt{संयतात्मा यो}}\Afootnote{H; \skt{संयतात्मा यो}\ A}}\skt{ }\edtext{\edtext{\skt{जन\-श\-त\-म\-धिकं}}{\lemma{\skt{जनशतम्}}\Afootnote{H; \skt{जनशत}\ A}}}{\lemma{\skt{अधिकं}}\Afootnote{H; \skt{मधिकं वा}\ A}}\skt{ गच्छेत्~। महाधनो वा }\edtext{\skt{श्यामा\-क\-नी\-वा\-रा\-म\-ल\-क\-ति\-न्दु\-क\-क\-पि\-त्था\-स्म\-न्त\-क\-फ\-ल\-मू\-ला\-हारो}}{\lemma{\skt{श्यामा\appbrk{}क\appbrk{}नी\appbrk{}वा\appbrk{}रा\appbrk{}म\appbrk{}ल\appbrk{}क\appbrk{}ति\appbrk{}न्दु\appbrk{}क\appbrk{}क\appbrk{}पि\appbrk{}त्था\appbrk{}स्म\appbrk{}न्त\appbrk{}क\appbrk{}फ\appbrk{}ल\appbrk{}मू\appbrk{}ला\appbrk{}हारो}}\Afootnote{H; \skt{श्यामा\appbrk{}क\appbrk{}नी\appbrk{}वा\appbrk{}र\appbrk{}वृ\appbrk{}त्ति\appbrk{}रा\appbrk{}म\appbrk{}ल\appbrk{}क\appbrk{}क\appbrk{}पि\appbrk{}त्थ\appbrk{}ति\appbrk{}न्दु\appbrk{}का\appbrk{}श्म\appbrk{}न्त\appbrk{}क\appbrk{}फ\appbrk{}ला\appbrk{}हारो}\ A}}\skt{ मृगैस्सह }\edtext{\skt{वसे\-त्मू\-त्र\-श\-कृ\-द्भक्षः}}{\lemma{\skt{मूत्रशकृद्भक्षः}}\Afootnote{H; \skt{तन्मू\appbrk{}त्र\appbrk{}श\appbrk{}कृ\appbrk{}द्भक्षः}\ A}}\skt{~। }\edtext{\skt{सत\-त\-म\-नु\-व्र\-जे\-द्वागाः}}{\lemma{\skt{वागाः}}\Afootnote{H; \skt{गाः}\ A}}\skt{~। }\edtext{\skt{व्राह्मणो वा}}{\lemma{\skt{व्राह्मणो वा}}\Afootnote{H; \skt{ब्रह्मणो वा}\ A}}\skt{ }\edtext{\skt{शिलो\-ञ्छ\-वृ\-त्ति\-र्ब्भूत्वा}}{\lemma{\skt{ब्भूत्वा}}\Afootnote{H; \skt{भूत्वा}\ A}}\skt{ ब्रह्म\-र\-थ\-मु\-द्ध\-रेत्~। }\edtext{\skt{कृसेत्कृशोदत्स}}{\lemma{\skt{कृसेत्कृशोदत्स}}\Afootnote{H\textsuperscript{pc}; \skt{कृषेत्}\ A; \skt{कृसेत्}\ H\textsuperscript{ac}}}\skt{ २ }\edtext{\skt{सततसमितं}}{\lemma{\skt{सततसमितं}}\Afootnote{H; \skt{सततमितरः}\ A}}\skt{ खनेद्वा }\edtext{\skt{कूपं}}{\lemma{\skt{कूपं}}\Afootnote{H; \skt{कूपं, कृशं तु सततं रक्षेत्}\ A}}\skt{~॥}\pend

\pstart\versepar \ledsidenote{\passref{4.11.13}}\edtext{\skt{भवति}}{\lemma{\skt{भवति}}\Afootnote{A; \emph{om.}\ H}}\skt{ }\edtext{\skt{चात्र्~॥ अधनो}}{\lemma{\skt{चात्र्~॥ अधनो}}\Afootnote{\skt{चात्र अधनो}\ A; \skt{भ{\symfont ¦}{\symfont ¦}{\symfont ¦}{\symfont ¦}{\symfont ¦}{\symfont ¦}{\symfont ¦}{\symfont ¦}{\symfont ¦}{\symfont ¦}{\symfont ¦}{\symfont ¦}{\symfont ¦}{\symfont ¦}{\symfont ¦}{\symfont ¦}{\symfont ¦}{\symfont ¦}{\symfont ¦}~॥ अधनो}\ H}}\skt{ वैद्यसन्देशा देवं }\edtext{\skt{कुर्व्वणतन्द्रितः}}{\lemma{\skt{कुर्व्वणतन्द्रितः}}\Afootnote{H; \skt{कुर्वन्नतन्द्रितः}\ A}}\skt{~। }\pend
\pstart\versepar \skt{}\edtext{\edtext{\skt{सम्व\-त्स\-रा\-न्त\-रा\-त्पुंसः}}{\lemma{\skt{सम्व\appbrk{}त्स\appbrk{}रा\appbrk{}न्त\appbrk{}रात्}}\Afootnote{H; \skt{संव\appbrk{}त्स\appbrk{}रा\appbrk{}द\appbrk{}न्त\appbrk{}राद्}\ A}}}{\lemma{\skt{पुंसः}}\Afootnote{H; \skt{वा}\ A}}\skt{ }\edtext{\skt{प्रमे\-हा\-न्प्र\-ति\-मु\-च्यते}}{\lemma{\skt{प्रमेहान्}}\Afootnote{H; \skt{प्रमेहात्}\ A}}\skt{ }\edtext{\skt{इति}}{\lemma{\skt{इति}}\Afootnote{H; \emph{om.}\ A}}\skt{~॥ }\edtext{\skt{प्र}}{\lemma{\skt{प्र}}\Afootnote{H; \emph{om.}\ A}}\skt{~॥}\pend

\pstart \ledsidenote{\passref{4.11.trailer}}\edtext{\skt{प्रमेहचिकित्सा}}{\lemma{\skt{प्रमेहचिकित्सा}}\Afootnote{H; \skt{इति सुश्रु\appbrk{}त\appbrk{}सं\appbrk{}हि\appbrk{}तायां चिकित्सास्थाने प्रमेहचिकित्सि}\ A}}\skt{ }\edtext{\skt{एकादश}}{\lemma{\skt{एकादश}}\Afootnote{H; \skt{तं नामै\appbrk{}क\appbrk{}द\appbrk{}शोऽ\appbrk{}ध्यायः}\ A}}\skt{~॥}\pend

\endnumbering
\beginnumbering
\eledsection*{Adhyāya 12}
\pstart \ledsidenote{\passref{4.12.1}}\skt{अथातः }\edtext{\skt{प्रमे\-ह\-पि\-ट\-क\-चि\-कि\-त्सितं}}{\lemma{\skt{प्रमे\appbrk{}ह\appbrk{}पि\appbrk{}ट\appbrk{}क\appbrk{}चि\appbrk{}कि\appbrk{}त्सितं}}\Afootnote{H; \skt{प्रमे\appbrk{}ह\appbrk{}पि\appbrk{}ड\appbrk{}का\appbrk{}चि\appbrk{}कि\appbrk{}त्सितं}\ A}}\skt{ व्याख्यास्यामः~॥}\pend

\pstart \ledsidenote{\passref{4.12.3}}\edtext{\skt{शराविकाद्याः}}{\lemma{\skt{शराविकाद्याः}}\Afootnote{H; \skt{शराविकाद्या नव पिडकाः}\ A}}\skt{ }\edtext{\skt{प्रागुक्ता}}{\lemma{\skt{प्रागुक्ता}}\Afootnote{H; \skt{प्रागुक्}\ A}}\skt{ }\edtext{\skt{नव\-पि\-ट\-का\-स्तास्तु}}{\lemma{\skt{नव\appbrk{}पि\appbrk{}ट\appbrk{}का\appbrk{}स्तास्तु}}\Afootnote{H\textsuperscript{pc}; \skt{ताः; ताः}\ A; \skt{नवपिटका तास्तु}\ H\textsuperscript{ac}}}\skt{ }\edtext{\edtext{\edtext{\skt{व्रण\-व\-तोऽ\-ल्पा\-स्त्व\-ङ्मा\-न्स\-प्राप्ताः}}{\lemma{\skt{व्रणवतो}}\Afootnote{H; \skt{प्राणवतो}\ A}}}{\lemma{\skt{ऽल्पास्}}\Afootnote{H; \skt{ऽल्पा}\ A}}}{\lemma{\skt{त्वङ्मान्सप्राप्ताः}}\Afootnote{H; \skt{स्त्वङ्मांसप्राप्ता}\ A}}\skt{~। मृद्व्यो }\edtext{\skt{ल्परुजः}}{\lemma{\skt{ल्परुजः}}\Afootnote{H; \skt{ऽल्परुजः}\ A}}\skt{ }\edtext{\skt{क्षिप्रभेदिन्यश्च}}{\lemma{\skt{क्षिप्रभेदिन्यश्च}}\Afootnote{H; \skt{क्षिप्र\appbrk{}पा\appbrk{}क\appbrk{}भे\appbrk{}दि\appbrk{}न्यश्च}\ A}}\skt{ साध्याः~।}\pend

\pstart \ledsidenote{\passref{4.12.4}}\skt{ताभिरुपद्रुतं }\edtext{\skt{प्रमे\-हि\-ण\-मु\-प\-क्र\-मेत}}{\lemma{\skt{उपक्रमेत}}\Afootnote{H; \skt{उपचरेत्}\ A}}\skt{~। तत्र }\edtext{\skt{पूर्व्व\-रू\-पे\-ष्व\-प\-त\-र्प्प\-ण\-म्व\-न\-स्प\-ति\-क\-षा\-य\-पानं}}{\lemma{\skt{वन\appbrk{}स्प\appbrk{}ति\appbrk{}क\appbrk{}षा\appbrk{}य\appbrk{}पानं}}\Afootnote{H; \skt{वन\appbrk{}स्प\appbrk{}ति\appbrk{}क\appbrk{}षायं}\ A}}\skt{~। }\edtext{\skt{वस्तमूत्रस्य}}{\lemma{\skt{वस्तमूत्रस्य}}\Afootnote{H; \skt{बस्तमूत्रं}\ A}}\skt{ }\edtext{\edtext{\skt{चोप\-दि\-शे\-दे\-व\-म\-कु\-र्व्वतो}}{\lemma{\skt{चोप\appbrk{}दि\appbrk{}शे\appbrk{}दे\appbrk{}वम्}}\Afootnote{H; \skt{चोपदिशेत्; एव}\ A}}}{\lemma{\skt{अकुर्व्वतो}}\Afootnote{H; \skt{मकुर्वत}\ A}}\skt{ }\edtext{\skt{वातस्य}}{\lemma{\skt{वातस्य}}\Afootnote{H; \skt{स्तसय}\ A}}\skt{ मधुराहारस्य }\edtext{\skt{पूर्व्वस्वेदः}}{\lemma{\skt{पूर्व्वस्वेदः}}\Afootnote{H; \skt{मूत्रं स्वेदः}\ A}}\skt{ श्लेष्मा च मधुरी भवति~। }\edtext{\skt{प्रमेह}}{\lemma{\skt{प्रमेह}}\Afootnote{H; \skt{प्रमेहश्}\ A}}\skt{ चाभिव्यक्तो भवति~। }\edtext{\skt{तस्योभयतः}}{\lemma{\skt{तस्योभयतः}}\Afootnote{H; \skt{तत्रोभयतः}\ A}}\skt{ }\edtext{\edtext{\edtext{\skt{संशो\-ध\-न\-मे\-वे\-म\-कु\-र्व्वतः}}{\lemma{\skt{संशोधनम्}}\Afootnote{H; \skt{संशोधन}\ A}}}{\lemma{\skt{एवेम्}}\Afootnote{H; \skt{मासेवेत; एव}\ A}}}{\lemma{\skt{अकुर्व्वतः}}\Afootnote{H; \skt{मकुर्वतस्}\ A}}\skt{~। तस्य दोषाः प्रवृद्धा }\edtext{\skt{मांसशोणितं}}{\lemma{\skt{मांसशोणितं}}\Afootnote{H; \skt{मांसशोणिते}\ A}}\skt{ प्रदूष्य }\edtext{\skt{शोफञ्जनयन्ति}}{\lemma{\skt{जनयन्ति}}\Afootnote{H; \skt{जनयन्त्}\ A}}\skt{~। }\edtext{\edtext{\skt{उप\-द्र\-वा\-श्का\-श्चि\-त्त\-त्रोक्तः}}{\lemma{\skt{उपद्रवाश्}}\Afootnote{H; \skt{युपद्रवान्}\ A}}}{\lemma{\skt{काश्चित्}}\Afootnote{H; \skt{वा कांश्चित्}\ A}}\skt{ प्रतीकारः सिरामोक्षश्च~। }\edtext{\edtext{\skt{एवमकुर्व्वतो}}{\lemma{\skt{एवम्}}\Afootnote{H; \skt{एव}\ A}}\skt{ वा}}{\lemma{\skt{अकुर्व्वतो वा}}\Afootnote{H; \skt{मकुर्वतस्}\ A}}\skt{ तस्य शोफः }\edtext{\skt{प्रवृद्धोऽतिमात्रं}}{\lemma{\skt{प्रवृद्धो}}\Afootnote{H; \skt{वृद्धो}\ A}}\skt{ }\edtext{\edtext{\edtext{\skt{तद्वि\-दा\-ह\-मा\-प\-द्यते}}{\lemma{\skt{तद्}}\Afootnote{H; \skt{रुजो}\ A}}}{\lemma{\skt{विदाहम्}}\Afootnote{H; \skt{विदाह}\ A}}}{\lemma{\skt{आपद्यते}}\Afootnote{H; \skt{मापद्यते}\ A}}\skt{~। तत्र }\edtext{\skt{शस्त्र\-प्र\-णि\-धा\-न\-म्व्र\-ण\-क्रि\-यो\-प\-सेवा}}{\lemma{\skt{शस्त्र\appbrk{}प्र\appbrk{}णि\appbrk{}धा\appbrk{}नम्}}\Afootnote{H; \skt{शस्त्र\appbrk{}प्र\appbrk{}णि\appbrk{}धा\appbrk{}न\appbrk{}मुक्तं}\ A}}\skt{ च~। }\edtext{\edtext{\skt{एवमकुर्व्वतो}}{\lemma{\skt{एवम्}}\Afootnote{H; \skt{एव}\ A}}\skt{ वा}}{\lemma{\skt{अकुर्व्वतो वा}}\Afootnote{H; \skt{मकुर्वतस्}\ A}}\skt{ तस्य }\edtext{\edtext{\edtext{\skt{पूयोऽ\-भ्य\-न्त\-र\-म\-व\-दा\-र्यो\-त्स\-ङ्ग\-म्म\-हा\-व\-का\-श\-ङ्कृत्वा}}{\lemma{\skt{ऽभ्यन्तरम्}}\Afootnote{H; \skt{ऽभ्यन्तर}\ A}}}{\lemma{\skt{अव\appbrk{}दा\appbrk{}र्यो\appbrk{}त्स\appbrk{}ङ्गम्}}\Afootnote{H; \skt{मवदार्योत्सङ्गं}\ A}}}{\lemma{\skt{महावकाशङ्}}\Afootnote{H; \skt{महा\appbrk{}न्त\appbrk{}म\appbrk{}व\appbrk{}काशं}\ A}}\skt{ }\edtext{\skt{वृद्धो}}{\lemma{\skt{वृद्धो}}\Afootnote{H; \skt{प्रवृद्धो}\ A}}\skt{ भवत्यसाध्यः~। तस्मादादित एव }\edtext{\skt{प्रमे\-हि\-ण\-मु\-प\-क्र\-मेत}}{\lemma{\skt{उपक्रमेत}}\Afootnote{H; \skt{उपक्रमेत्}\ A}}\pend

\pstart \ledsidenote{\passref{4.12.6}}\edtext{\skt{दुर्व्वि\-रे\-च्य\-म\-धु\-मे\-हिनो}}{\lemma{\skt{दुर्व्वि\appbrk{}रे\appbrk{}च्य\appbrk{}म\appbrk{}धु\appbrk{}मे\appbrk{}हिनो}}\Afootnote{H; \skt{दुर्विरेच्या हि मधुमेहिनो}\ A}}\skt{ भवन्ति~। मेदोऽ\-भि\-व्या\-प्त\-श\-री\-र\-त्वात्~। }\edtext{\edtext{\skt{तस्मात्तीक्ष्णमेषां}}{\lemma{\skt{तीक्ष्णम्}}\Afootnote{H; \skt{तीक्ष्ण}\ A}}}{\lemma{\skt{एषां}}\Afootnote{H; \skt{मेतेषां}\ A}}\skt{ }\edtext{\skt{शोधनङ्कुर्वीत}}{\lemma{\skt{कुर्वीत}}\Afootnote{H; \skt{कुर्वीत~। पिडकापीडिताः सोपद्रवाः सर्व एव प्रमेहा मूत्रादिमाधुर्ये मधु\appbrk{}ग\appbrk{}न्ध\appbrk{}सा\appbrk{}मा\appbrk{}न्या\appbrk{}त्पा\appbrk{}रि\appbrk{}भा\appbrk{}षिकीं मधुमेहाख्यां लभन्ते}\ A}}\skt{~।}\pend

\pstart \ledsidenote{\passref{4.12.7}}\edtext{\skt{न }\edtext{\skt{चैत\-ङ्क\-थं\-चि\-दपि}}{\lemma{\skt{कथंचिदपि}}\Afootnote{H; \skt{कथंचिदपि}\ A}}}{\lemma{\skt{न चैतङ्}}\Afootnote{H; \skt{न चैतान्}\ A}}\skt{ स्वेदयेत्~। }\edtext{\edtext{\skt{मेदो\-ब\-हु\-त्वा\-त्ख\-ल्वेषां}}{\lemma{\skt{मेदोबहुत्वात्}}\Afootnote{H; \skt{मेदोबहुत्वादे}\ A}}}{\lemma{\skt{खल्वेषां}}\Afootnote{H; \skt{तेषां}\ A}}\skt{ }\edtext{\skt{सीर्यते}}{\lemma{\skt{सीर्यते}}\Afootnote{H; \skt{विशीर्यते}\ A}}\skt{ देहः~। स्वेदेन}\pend

\pstart \ledsidenote{\passref{4.12.8}}\edtext{\skt{रसहरणीनां}}{\lemma{\skt{रसहरणीनां}}\Afootnote{H; \skt{रसायनीनां च}\ A}}\skt{ }\edtext{\edtext{\skt{दौर्व्व\-ल्या\-न्नो\-र्द्ध\-मु\-त्ति\-ष्ठंति}}{\lemma{\skt{नोर्द्धम्}}\Afootnote{H; \skt{नोर्ध्}\ A}}}{\lemma{\skt{उत्तिष्ठंति}}\Afootnote{H; \skt{वमुत्तिष्ठन्ति}\ A}}\skt{~। प्रमेहिणां }\edtext{\skt{दोषाः}}{\lemma{\skt{दोषाः}}\Afootnote{H; \skt{दोषाः}\ A}}\skt{~। }\edtext{\skt{यतो}}{\lemma{\skt{यतो}}\Afootnote{H; \skt{ततो}\ A}}\skt{ }\edtext{\edtext{\skt{मधु\-मे\-हि\-ना\-म\-धः\-काये}}{\lemma{\skt{मधुमेहिनाम्}}\Afootnote{H; \skt{मधुमेहिना}\ A}}}{\lemma{\skt{अधःकाये}}\Afootnote{H; \skt{मधःकाये}\ A}}\skt{ }\edtext{\skt{पिटकाः}}{\lemma{\skt{पिटकाः}}\Afootnote{H; \skt{पिडकाः}\ A}}\skt{ }\edtext{\edtext{\skt{प्रादुर्ब्भवन्ति}}{\lemma{\skt{प्रादुर्}}\Afootnote{H; \skt{प्रादु}\ A}}}{\lemma{\skt{ब्भवन्ति}}\Afootnote{H; \skt{र्भवन्ति}\ A}}\skt{~।}\pend

\pstart \ledsidenote{\passref{4.12.9}}\skt{अपक्वानां }\edtext{\skt{खलु}}{\lemma{\skt{खलु}}\Afootnote{H; \skt{तु}\ A}}\skt{ }\edtext{\skt{पिटकानां}}{\lemma{\skt{पिटकानां}}\Afootnote{H; \skt{पिडकानां}\ A}}\skt{ शोफ\-व\-त्प्र\-ती\-कारः~। पक्वा\-ना\-म्व्र\-ण\-व\-दिति~। तैलन्तु }\edtext{\edtext{\skt{व्रण\-रो\-प\-ण\-मे\-वादौ}}{\lemma{\skt{व्रणरोपणम्}}\Afootnote{H; \skt{व्रणरोपण}\ A}}}{\lemma{\skt{एवादौ}}\Afootnote{H; \skt{मेवादौ}\ A}}\skt{ कुर्वीत~। }\edtext{\edtext{\skt{आर\-ग्व\-धा\-दि\-क\-षा\-य\-मु\-त्सा\-द\-नार्थे}}{\lemma{\skt{आर\appbrk{}ग्व\appbrk{}धा\appbrk{}दि\appbrk{}क\appbrk{}षा\appbrk{}यम्}}\Afootnote{H; \skt{आर\appbrk{}ग्व\appbrk{}धा\appbrk{}दि\appbrk{}क\appbrk{}षाय}\ A}}}{\lemma{\skt{उत्सादनार्थे}}\Afootnote{H; \skt{मुत्सादनार्थे}\ A}}\skt{~। }\edtext{\skt{शाल\-सा\-रा\-दि\-क\-षा\-य\-म्प\-रि\-से\-चने}}{\lemma{\skt{परिसेचने}}\Afootnote{H; \skt{परिषेचने}\ A}}\skt{~। पिप्प\-ल्या\-दि\-क\-षा\-य\-म्पा\-न\-भो\-ज\-नेषु~। }\edtext{\skt{पाठा\-चि\-त्र\-क\-शा\-र्ङ्ग्गे\-ष्टा\-क्षु\-द्र\-वृ\-हती}}{\lemma{\skt{पाठा\appbrk{}चि\appbrk{}त्र\appbrk{}क\appbrk{}शा\appbrk{}र्ङ्ग्गे\appbrk{}ष्टा\appbrk{}क्षु\appbrk{}द्र\appbrk{}वृ\appbrk{}हती}}\Afootnote{H; \skt{पाठा\appbrk{}चि\appbrk{}त्र\appbrk{}क\appbrk{}शा\appbrk{}र्ङ्गे\appbrk{}ष्टा\appbrk{}क्षु\appbrk{}द्रा\appbrk{}बृ\appbrk{}हती}\ A}}\skt{ }\edtext{\skt{शारि\-वा\-सो\-म\-क\-ल्क\-स\-प्त\-प\-र्ण्णा\-र\-ग्व\-ध\-कु\-ट\-ज\-मू\-ल\-चू\-र्ण्णानि}}{\lemma{\skt{शारि\appbrk{}वा\appbrk{}सो\appbrk{}म\appbrk{}क\appbrk{}ल्क\appbrk{}स\appbrk{}प्त\appbrk{}प\appbrk{}र्ण्णा\appbrk{}र\appbrk{}ग्व\appbrk{}ध\appbrk{}कु\appbrk{}ट\appbrk{}ज\appbrk{}मू\appbrk{}ल\appbrk{}चू\appbrk{}र्ण्णानि}}\Afootnote{H\textsuperscript{pc}; \skt{सारि\appbrk{}वा\appbrk{}सो\appbrk{}म\appbrk{}व\appbrk{}ल्क\appbrk{}स\appbrk{}प्त\appbrk{}प\appbrk{}र्णा\appbrk{}र\appbrk{}ग्व\appbrk{}ध\appbrk{}कु\appbrk{}ट\appbrk{}ज\appbrk{}मू\appbrk{}ल\appbrk{}चू\appbrk{}र्णानि}\ A; \skt{शारि\appbrk{}वा\appbrk{}सो\appbrk{}म\appbrk{}व\appbrk{}ल्क\appbrk{}स\appbrk{}प्त\appbrk{}प\appbrk{}र्ण्णा\appbrk{}र\appbrk{}ग्व\appbrk{}ध\appbrk{}कु\appbrk{}ट\appbrk{}ज\appbrk{}मू\appbrk{}ल\appbrk{}चू\appbrk{}र्ण्णानि}\ H\textsuperscript{ac}}}\skt{ }\edtext{\skt{मधुविमिश्राणि}}{\lemma{\skt{मधुविमिश्राणि}}\Afootnote{H; \skt{मधुमिथाणि}\ A}}\skt{ }\edtext{\skt{चाश्नीयात्}}{\lemma{\skt{चाश्नीयात्}}\Afootnote{H; \skt{प्रश्नीयात्}\ A}}\skt{~।}\pend

\pstart \ledsidenote{\passref{4.12.10}}\edtext{\edtext{\skt{शाल\-सा\-रा\-दि\-क\-षा\-य\-च\-तु\-र्था\-ञ्शा\-व\-शि\-ष्ट\-म\-व\-तार्य}}{\lemma{\skt{शाल\appbrk{}सा\appbrk{}रा\appbrk{}दि\appbrk{}क\appbrk{}षा\appbrk{}य\appbrk{}च\appbrk{}तु\appbrk{}र्था\appbrk{}ञ्शा\appbrk{}व\appbrk{}शि\appbrk{}ष्टम्}}\Afootnote{H; \skt{शाल\appbrk{}सा\appbrk{}रा\appbrk{}दि\appbrk{}व\appbrk{}र्ग\appbrk{}क\appbrk{}षायं चतु\appbrk{}र्भा\appbrk{}गा\appbrk{}व\appbrk{}शिष्ट}\ A}}}{\lemma{\skt{अवतार्य}}\Afootnote{H; \skt{मवतार्य}\ A}}\skt{ परिस्राव्य पुनरुपनीय साधयेत्~। सिध्यति }\edtext{\skt{चाम\-ल\-क\-लो\-ध्र\-प्रि\-य\-ङ्गु\-दन्ती}}{\lemma{\skt{चाम\appbrk{}ल\appbrk{}क\appbrk{}लो\appbrk{}ध्र\appbrk{}प्रि\appbrk{}य\appbrk{}ङ्गु\appbrk{}दन्ती}}\Afootnote{H; \skt{चाम\appbrk{}ल\appbrk{}क\appbrk{}रो\appbrk{}ध्र\appbrk{}प्रि\appbrk{}य\appbrk{}ङ्गु\appbrk{}दन्ती}\ A}}\skt{ }\edtext{\edtext{\skt{कृष्णा\-य\-स्ता\-म्र\-चू\-र्णा\-न्या\-व\-ये\-त्त\-दे\-त\-द\-नु\-प\-दग्धं}}{\lemma{\skt{कृष्णा\appbrk{}य\appbrk{}स्ता\appbrk{}म्र\appbrk{}चू\appbrk{}र्णा\appbrk{}न्या\appbrk{}व\appbrk{}येत्}}\Afootnote{H; \skt{कृष्णा\appbrk{}य\appbrk{}स्ता\appbrk{}म्र\appbrk{}चू\appbrk{}र्णा\appbrk{}न्या\appbrk{}वपे}\ A}}}{\lemma{\skt{तदे\appbrk{}त\appbrk{}द\appbrk{}नु\appbrk{}प\appbrk{}दग्धं}}\Afootnote{H; \skt{त्, एत\appbrk{}द\appbrk{}नु\appbrk{}प\appbrk{}दग्धं}\ A}}\skt{ }\edtext{\edtext{\skt{लेही\-भू\-त\-म\-व\-ता\-र्या\-नु\-गु\-प्त\-न्नि\-द\-ध्या\-त्ततो}}{\lemma{\skt{लेहीभूतम्}}\Afootnote{H; \skt{लेहीभूत}\ A}}}{\lemma{\skt{अव\appbrk{}ता\appbrk{}र्या\appbrk{}नु\appbrk{}गु\appbrk{}प्तन्}}\Afootnote{H; \skt{मव\appbrk{}ता\appbrk{}र्या\appbrk{}नु\appbrk{}गुप्तं}\ A}}\skt{ }\edtext{\skt{यथा\-यो\-ग\-मु\-प\-यु\-ञ्ज्येत}}{\lemma{\skt{उपयुञ्ज्येत}}\Afootnote{H; \skt{उपयुञ्जीत}\ A}}\skt{~। एष लेहः }\edtext{\edtext{\skt{सर्व्व\-म\-धु\-मे\-हा\-ना\-न्नि\-हन्ता}}{\lemma{\skt{सर्व्व\appbrk{}म\appbrk{}धु\appbrk{}मे\appbrk{}हा\appbrk{}नान्}}\Afootnote{H; \skt{सर्वमेहानां}\ A}}}{\lemma{\skt{निहन्ता}}\Afootnote{H; \skt{हन्ता}\ A}}\skt{~॥}\pend

\pstart \ledsidenote{\passref{4.12.11}}\skt{त्रिफला }\edtext{\skt{त्रिक\-टु\-क\-चि\-त्र\-क\-वि\-डं\-ग\-मु\-स्ता\-ना\-न्न\-व\-भा\-गा\-स्ता\-वन्त}}{\lemma{\skt{त्रिक\appbrk{}टु\appbrk{}क\appbrk{}चि\appbrk{}त्र\appbrk{}क\appbrk{}वि\appbrk{}डं\appbrk{}ग\appbrk{}मु\appbrk{}स्ता\appbrk{}नान्}}\Afootnote{H; \skt{नित्र\appbrk{}क\appbrk{}त्रि\appbrk{}क\appbrk{}टु\appbrk{}क\appbrk{}वि\appbrk{}ड\appbrk{}ङ्ग\appbrk{}मु\appbrk{}स्तानां}\ A}}\skt{ एव कृष्णायश्चूर्णस्य }\edtext{\edtext{\skt{तत्स\-र्व्व\-मे\-क\-ध्य\-ङ्कृत्वा}}{\lemma{\skt{सर्व्वम्}}\Afootnote{H; \skt{सर्व}\ A}}}{\lemma{\skt{एकध्यङ्}}\Afootnote{H; \skt{मैकध्यं}\ A}}\skt{ }\edtext{\edtext{\skt{यथा\-यो\-ग\-म्मा\-त्रा\-स\-र्प्पि\-र्म्म\-धुभ्यां}}{\lemma{\skt{मात्रासर्प्पिर्}}\Afootnote{H; \skt{मात्रां सर्पि}\ A}}}{\lemma{\skt{म्मधुभ्यां}}\Afootnote{H; \skt{र्मधुभ्यां}\ A}}\skt{ }\edtext{\skt{संसृ\-ज्यो\-प\-यु\-ञ्जीत}}{\lemma{\skt{संसृ\appbrk{}ज्यो\appbrk{}प\appbrk{}यु\appbrk{}ञ्जीत}}\Afootnote{H; \skt{संसृज्योपयुञ्जी}\ A}}\skt{~। }\edtext{\edtext{\skt{तदे\-त\-न्न\-वा\-य\-स\-न्नाम}}{\lemma{\skt{तदेतन्}}\Afootnote{H; \skt{त, एतन्}\ A}}}{\lemma{\skt{नाम}}\Afootnote{H; \emph{om.}\ A}}\skt{ एतेन }\edtext{\skt{खलु}}{\lemma{\skt{खलु}}\Afootnote{H; \emph{om.}\ A}}\skt{ जाठ }\edtext{\skt{र्यन्न}}{\lemma{\skt{र्यन्न}}\Afootnote{A H\textsuperscript{pc}; \skt{र्यणम्}\ H\textsuperscript{ac}}}\skt{ भवति~। }\edtext{\edtext{\skt{सन्ना\-ग्नि\-रा\-प्या\-य्यते}}{\lemma{\skt{सन्नाग्निर्}}\Afootnote{H; \skt{सन्नोऽग्नि}\ A}}}{\lemma{\skt{आप्याय्यते}}\Afootnote{H; \skt{राप्यायते}\ A}}\skt{~। }\edtext{\skt{दुर्न्नाम्नो न}}{\lemma{\skt{दुर्न्नाम्नो न}}\Afootnote{H; \skt{दुर्ना}\ A}}\skt{ }\edtext{\skt{भवन्ति}}{\lemma{\skt{भवन्ति}}\Afootnote{H; \skt{म}\ A}}\skt{~। }\edtext{\skt{शोफ\-पा\-ण्डु\-रो\-गा\-वि\-पा\-क\-का\-स\-श्वा\-सश्च}}{\lemma{\skt{शोफ\appbrk{}पा\appbrk{}ण्डु\appbrk{}रो\appbrk{}गा\appbrk{}वि\appbrk{}पा\appbrk{}क\appbrk{}का\appbrk{}स\appbrk{}श्वा\appbrk{}सश्च}}\Afootnote{H; \skt{शोफ\appbrk{}पा\appbrk{}ण्डु\appbrk{}कु\appbrk{}ष्ठ\appbrk{}रो\appbrk{}गा\appbrk{}वि\appbrk{}पा\appbrk{}क\appbrk{}का\appbrk{}स\appbrk{}श्वा\appbrk{}स\appbrk{}प्र\appbrk{}मे\appbrk{}हाश्च न भवन्ति}\ A}}\skt{~।}\pend

\pstart \ledsidenote{\passref{4.12.12}}\edtext{\edtext{\skt{शाल\-सा\-रा\-दि\-नि\-र्यू\-ह\-श्च\-तु\-र्थां\-शा\-व\-शे\-षितं}}{\lemma{\skt{शाल\appbrk{}सा\appbrk{}रा\appbrk{}दि\appbrk{}नि\appbrk{}र्यू\appbrk{}हश्}}\Afootnote{H; \skt{शाल\appbrk{}सा\appbrk{}रा\appbrk{}दि\appbrk{}नि\appbrk{}र्यूहे}\ A}}}{\lemma{\skt{चतु\appbrk{}र्थां\appbrk{}शा\appbrk{}व\appbrk{}शे\appbrk{}षितं}}\Afootnote{H; \skt{चतु\appbrk{}र्थां\appbrk{}शा\appbrk{}व\appbrk{}शे\appbrk{}षिते}\ A}}\skt{~। }\edtext{\skt{परिस्रुतं}}{\lemma{\skt{परिस्रुतं}}\Afootnote{H; \skt{परिस्रु}\ A}}\skt{ }\edtext{\skt{सुशीतञ्च}}{\lemma{\skt{सुशीतञ्च}}\Afootnote{H; \skt{ते त}\ A}}\skt{ }\edtext{\skt{मधु\-ना\-म\-धु\-री\-कृतं}}{\lemma{\skt{मधु\appbrk{}ना\appbrk{}म\appbrk{}धु\appbrk{}री\appbrk{}कृतं}}\Afootnote{H; \skt{तः शीते मधु माक्षि\appbrk{}क\appbrk{}मा\appbrk{}व\appbrk{}पेत्}\ A}}\skt{~।}\pend

\pstart \ledsidenote{\passref{4.12.13}}\edtext{\edtext{\skt{फाणि\-ती\-भा\-व\-मा\-पन्नं}}{\lemma{\skt{फाणितीभावम्}}\Afootnote{H; \skt{फाणितीभाव}\ A}}}{\lemma{\skt{आपन्नं}}\Afootnote{H; \skt{मापन्नं}\ A}}\skt{ गुडं शोधितमेव च। श्लक्ष्णपिष्टानि चूर्णानि }\edtext{\skt{पिप्प\-ल्या\-दे\-र्ग्ग\-णस्य}}{\lemma{\skt{पिप्पल्यादेर्}}\Afootnote{H; \skt{पिप्पल्यादि}\ A}}\skt{ च~॥}\pend

\pstart \ledsidenote{\passref{4.12.14}}\edtext{\edtext{\skt{ऐक\-ध्य\-मा\-व\-पे\-त्कुम्भे}}{\lemma{\skt{ऐकध्यम्}}\Afootnote{H; \skt{ऐकध्य}\ A}}}{\lemma{\skt{आवपेत्}}\Afootnote{H; \skt{मावपेत्}\ A}}\skt{ संस्कृते }\edtext{\skt{घृतभाविते}}{\lemma{\skt{घृतभाविते}}\Afootnote{H; \skt{घृतभावैते}\ A}}\skt{~। पिप्प\-ली\-चू\-र्ण्ण\-म\-धुभिः }\edtext{\skt{प्रलिप्तेऽभ्यन्तरं}}{\lemma{\skt{ऽभ्यन्तरं}}\Afootnote{H; \skt{ऽन्तः}\ A}}\skt{ }\edtext{\skt{शुचौ}}{\lemma{\skt{शुचौ}}\Afootnote{H; \skt{शुचौ दृढे}\ A}}\skt{~॥}\pend

\pstart \ledsidenote{\passref{4.12.15}}\skt{श्लक्ष्णानि तीक्ष्णलोहस्य तत्र पत्राणि बुद्धिमान्~। खदि\-रा\-ङ्गा\-र\-त\-प्तानि बहुशः }\edtext{\edtext{\skt{प्रक्षिप्येद्वुधः}}{\lemma{\skt{प्रक्षिप्येद्}}\Afootnote{H; \skt{सन्निपातयेत्}\ A}}}{\lemma{\skt{वुधः}}\Afootnote{H; \emph{om.}\ A}}\skt{~॥}\pend

\pstart \ledsidenote{\passref{4.12.16}}\edtext{\skt{सपिधानन्तु}}{\lemma{\skt{सपिधानन्तु}}\Afootnote{H; \skt{सुपिधानं तु}\ A}}\skt{ }\edtext{\skt{तत्कृत्वा}}{\lemma{\skt{तत्}}\Afootnote{H; \skt{तं}\ A}}\skt{ यवपल्ले निधापयेत्~। }\edtext{\skt{मासास्त्रीश्चतुरो}}{\lemma{\skt{त्रीश्}}\Afootnote{H; \skt{स्त्रींश्}\ A}}\skt{ }\edtext{\skt{वापि}}{\lemma{\skt{वापि}}\Afootnote{H; \skt{वाऽपि}\ A}}\skt{ }\edtext{\skt{यावद्वा}}{\lemma{\skt{यावद्वा}}\Afootnote{H; \skt{यावदा}\ A}}\skt{ लोहसंक्षयात्~॥}\pend

\pstart \ledsidenote{\passref{4.12.17}}\skt{ततो }\edtext{\skt{जातरसञ्जन्तुः}}{\lemma{\skt{जन्तुः}}\Afootnote{H; \skt{तं तु}\ A}}\skt{ प्रातः }\edtext{\skt{प्रातर्यथावलं}}{\lemma{\skt{प्रातर्यथावलं}}\Afootnote{H; \skt{प्रात\appbrk{}र्य\appbrk{}था\appbrk{}ब\appbrk{}लम्}\ A}}\skt{~। }\edtext{\edtext{\edtext{\edtext{\skt{उप\-यु\-ञ्ज्या\-द्य\-था\-यो\-ग\-मा\-हा\-र\-म्वास्य}}{\lemma{\skt{उपयुञ्ज्याद्}}\Afootnote{H; \skt{निषेवेत}\ A}}}{\lemma{\skt{यथायोगम्}}\Afootnote{H; \skt{यथायोग}\ A}}}{\lemma{\skt{आहारम्}}\Afootnote{H; \skt{माहारं}\ A}}}{\lemma{\skt{वास्य}}\Afootnote{H; \skt{चास्य}\ A}}\skt{ कल्पयेत्~॥}\pend

\pstart \ledsidenote{\passref{4.12.18}}\edtext{\skt{एष}}{\lemma{\skt{एष}}\Afootnote{H; \emph{om.}\ A}}\skt{ }\edtext{\skt{स्थूलस्य कृष\-क\-त्स्न\-न्न\-स्याग्नेः}}{\lemma{\skt{स्थूलस्य कृष\appbrk{}क\appbrk{}त्स्न\appbrk{}न्न\appbrk{}स्याग्नेः}}\Afootnote{H; \skt{कार्श्य\appbrk{}कृ\appbrk{}द्ब\appbrk{}लि\appbrk{}ना\appbrk{}मेष सन्नस्याग्नेः}\ A}}\skt{ प्रसाधकः~। }\edtext{\skt{शोफघ्नः}}{\lemma{\skt{शोफघ्नः}}\Afootnote{H; \skt{शोफ\appbrk{}नु\appbrk{}द्गु\appbrk{}ल्म\appbrk{}हृत्}\ A}}\skt{ }\edtext{\skt{कुष्ठमेहघ्नो}}{\lemma{\skt{कुष्ठमेहघ्नो}}\Afootnote{H; \skt{कुष्ठमेह}\ A}}\skt{ }\edtext{\skt{गुल्म\-पा\-ण्ड्वा\-म\-या\-पहः}}{\lemma{\skt{गुल्म\appbrk{}पा\appbrk{}ण्ड्वा\appbrk{}म\appbrk{}या\appbrk{}पहः}}\Afootnote{H; \skt{पाण्ड्वामयापहः}\ A}}\skt{~॥}\pend

\pstart \ledsidenote{\passref{4.12.19}}\skt{प्लीहोदरहरः }\edtext{\skt{शीघ्रं म्विष\-म\-ज्व\-र\-ना\-शनः}}{\lemma{\skt{शीघ्रं म्विष\appbrk{}म\appbrk{}ज्व\appbrk{}र\appbrk{}ना\appbrk{}शनः}}\Afootnote{H; \skt{शीघ्रं विष\appbrk{}म\appbrk{}ज्व\appbrk{}र\appbrk{}ना\appbrk{}शनः}\ A}}\skt{~। अभि\-ष्य\-न्दा\-प\-ह\-रणो लोहारिष्टो महागुणः~॥}\pend

\pstart \ledsidenote{\passref{4.12.20}}\skt{प्रमेहिणो यदा }\edtext{\edtext{\edtext{\skt{मूत्र\-म\-ना\-वि\-ल\-मपि}}{\lemma{\skt{मूत्रम्}}\Afootnote{H; \skt{मूत्र}\ A}}}{\lemma{\skt{अनाविलम्}}\Afootnote{H; \skt{मपि\appbrk{}च्छि\appbrk{}ल\appbrk{}म\appbrk{}नावि}\ A}}}{\lemma{\skt{अपि}}\Afootnote{H; \emph{om.}\ A}}\skt{ }\edtext{\skt{छिलं}}{\lemma{\skt{छिलं}}\Afootnote{H; \skt{लम्}\ A}}\skt{~। }\edtext{\edtext{\skt{विश\-द\-न्ति\-क्त\-क\-टु\-क\-न्त\-दा\-रो\-ग\-म्प्र\-च\-क्षत}}{\lemma{\skt{तदारोगम्}}\Afootnote{H; \skt{तदाऽऽआरोग्यं}\ A}}}{\lemma{\skt{प्रचक्षत}}\Afootnote{H; \skt{प्रचक्षते}\ A}}\skt{ }\edtext{\skt{इति}}{\lemma{\skt{इति}}\Afootnote{H; \emph{om.}\ A}}\skt{~॥}\pend

\pstart \ledsidenote{\passref{4.12.trailer}}\edtext{\skt{ं~॥ चिकित्साद्वादश}}{\lemma{\skt{ं~॥ चिकित्साद्वादश}}\Afootnote{H; \skt{इति श्रीसु\appbrk{}श्रु\appbrk{}त\appbrk{}सं\appbrk{}हि\appbrk{}तायां चिकित्सास्थाने प्रमे\appbrk{}ह\appbrk{}पि\appbrk{}ड\appbrk{}का\appbrk{}चि\appbrk{}कि\appbrk{}त्सितं नाम द्वादशोऽध्यायः}\ A}}\skt{~॥}\pend

\endnumbering
\beginnumbering
\eledsection*{Adhyāya 13}
\pstart \ledsidenote{\passref{4.13.1}}\skt{अथातो }\edtext{\skt{मधुमेहिं}}{\lemma{\skt{मधुमेहिं}}\Afootnote{H; \skt{मधुमेह}\ A}}\skt{ चिकित्सितं व्याख्यास्यामः~॥}\pend

\pstart\versepar \ledsidenote{\passref{4.13.3}}\edtext{\edtext{\skt{मधु\-मे\-हि\-त्व\-मा\-पन्नं}}{\lemma{\skt{मधुमेहित्वम्}}\Afootnote{H; \skt{मधुमेहित्व}\ A}}}{\lemma{\skt{आपन्नं}}\Afootnote{H; \skt{मापन्नं}\ A}}\skt{ भिषग्भिः परिवर्ज्जितं~। }\pend
\pstart\versepar \skt{योगेनानेन }\edtext{\edtext{\skt{मति\-मा\-न्म\-धु\-मे\-हि\-न\-मा\-च\-रेत्}}{\lemma{\skt{मधुमेहिनम्}}\Afootnote{H; \skt{प्रमेहिण}\ A}}}{\lemma{\skt{आचरेत्}}\Afootnote{H; \skt{मुपाचरेत्}\ A}}\skt{~॥}\pend

\pstart\versepar \ledsidenote{\passref{4.13.4}}\edtext{\skt{मासिशुक्ले}}{\lemma{\skt{मासिशुक्ले}}\Afootnote{H; \skt{मासे शुक्रे}\ A}}\skt{ शुचौ चैव शैलाः सूर्याङ्शुतापिताः~॥ }\pend
\pstart\versepar \skt{जतुप्रकाशं स्वरसं शिलाभ्यः }\edtext{\skt{प्रभवन्ति}}{\lemma{\skt{प्रभवन्ति}}\Afootnote{H; \skt{प्रस्रवन्ति}\ A}}\skt{ हि~॥}\pend

\pstart\versepar \ledsidenote{\passref{4.13.5}}\skt{शिलाजत्विति }\edtext{\skt{विख्या\-त\-म्न\-हा\-व्या\-धि\-वि\-ना\-शनं}}{\lemma{\skt{नहा\appbrk{}व्या\appbrk{}धि\appbrk{}वि\appbrk{}ना\appbrk{}शनं}}\Afootnote{H; \skt{सर्व\appbrk{}व्या\appbrk{}धि\appbrk{}वि\appbrk{}ना\appbrk{}श\appbrk{}नम्}\ A}}\skt{~। }\pend
\pstart\versepar \skt{}\edtext{\skt{त्रप्वा\-दी\-ना\-न्त्य्लो\-हानां}}{\lemma{\skt{त्य्}}\Afootnote{H; \skt{तु}\ A}}\skt{ }\edtext{\edtext{\skt{षण्णा\-म\-न्य\-त\-मा\-श्रितं}}{\lemma{\skt{षण्णाम्}}\Afootnote{H; \skt{षण्णा}\ A}}}{\lemma{\skt{अन्यतमाश्रितं}}\Afootnote{H; \skt{मन्य\appbrk{}त\appbrk{}मा\appbrk{}न्व\appbrk{}यम्}\ A}}\skt{~॥}\pend

\pstart\versepar \ledsidenote{\passref{4.13.6}}\edtext{\edtext{\skt{ज्ञेय\-न्त\-त्ख\-ल्त\-श्चापि}}{\lemma{\skt{तत्}}\Afootnote{H; \skt{स्वगन्}\ A}}}{\lemma{\skt{खल्तश्}}\Afootnote{H; \skt{धतश्}\ A}}\skt{ }\edtext{\skt{षड्योनिः}}{\lemma{\skt{षड्योनिः}}\Afootnote{H; \skt{षड्योनि}\ A}}\skt{ प्रथितं क्षितौ~। }\pend
\pstart\versepar \skt{}\edtext{\skt{लोहाद्यस्मात्तु}}{\lemma{\skt{लोहाद्यस्मात्तु}}\Afootnote{H; \skt{लोहाद्भवति तद्यस्मा}\ A}}\skt{ }\edtext{\skt{यज्जातं}}{\lemma{\skt{यज्जातं}}\Afootnote{H; \skt{च्छ्}\ A}}\skt{ }\edtext{\skt{शिलाजतु}}{\lemma{\skt{शिलाजतु}}\Afootnote{H; \skt{इलाजतु}\ A}}\skt{ }\edtext{\skt{हिताय वै}}{\lemma{\skt{हिताय वै}}\Afootnote{H; \skt{जतुप्रभम्}\ A}}\skt{~॥}\pend

\pstart\versepar \ledsidenote{\passref{4.13.7}}\skt{तस्य लोहस्य }\edtext{\skt{वीर्यञ्च}}{\lemma{\skt{वीर्यञ्च}}\Afootnote{H; \skt{तद्वीर्यं}\ A}}\skt{ रसञ्चापि विभर्त्ति तत्~। }\pend
\pstart\versepar \skt{त्रपु\-सी\-सा\-य\-सा\-दीनि }\edtext{\skt{प्रधा\-न्यु\-त्त\-रो\-त्तरं}}{\lemma{\skt{प्रधा\appbrk{}न्यु\appbrk{}त्त\appbrk{}रो\appbrk{}त्तरं}}\Afootnote{H; \skt{प्रधा\appbrk{}ना\appbrk{}न्यु\appbrk{}त्त\appbrk{}रो\appbrk{}त्त\appbrk{}रम्}\ A}}\skt{~॥}\pend

\pstart\versepar \ledsidenote{\passref{4.13.8}}\skt{यथा तथा }\edtext{\skt{प्रयोगेपि}}{\lemma{\skt{प्रयोगेपि}}\Afootnote{H; \skt{प्रयोगेऽपि}\ A}}\skt{ }\edtext{\skt{श्रेष्ठ}}{\lemma{\skt{श्रेष्ठ}}\Afootnote{H; \skt{श्रेष्ठे}\ A}}\skt{ }\edtext{\skt{श्रेष्था}}{\lemma{\skt{श्रेष्था}}\Afootnote{H; \skt{श्रेष्ठ}\ A}}\skt{ गुणाः स्मृताः~। }\pend
\pstart\versepar \skt{तत्स\-र्व्व\-न्ति\-क्त\-क\-टुकं कषायानुरसं }\edtext{\skt{रसं}}{\lemma{\skt{रसं}}\Afootnote{H; \skt{सरम्}\ A}}\skt{~॥}\pend

\pstart\versepar \ledsidenote{\passref{4.13.9}}\skt{कटु\-पा\-क्यु\-ष्ण\-वी\-र्यञ्च }\edtext{\skt{शोष\-ण\-ञ्छे\-द\-न\-म्परं}}{\lemma{\skt{परं}}\Afootnote{H; \skt{तथा}\ A}}\skt{~। }\pend
\pstart\versepar \skt{तेषु }\edtext{\skt{यत्कृष्णमलघुः}}{\lemma{\skt{कृष्णमलघुः}}\Afootnote{H; \skt{कृष्णमलघु}\ A}}\skt{ स्निग्धन्निः शर्क्करञ्च यत्~॥}\pend

\pstart\versepar \ledsidenote{\passref{4.13.10}}\edtext{\skt{गोमूत्रगन्धी}}{\lemma{\skt{गोमूत्रगन्धी}}\Afootnote{H; \skt{गोमूत्रगन्धि}\ A}}\skt{ }\edtext{\skt{नीलम्वा}}{\lemma{\skt{नीलम्वा}}\Afootnote{H; \skt{यच्चापि}\ A}}\skt{ }\edtext{\skt{तत्प्रचक्षते}}{\lemma{\skt{प्रचक्षते}}\Afootnote{H; \skt{प्रधानं प्रचक्षते}\ A}}\skt{~। }\pend
\pstart\versepar \skt{}\edtext{\skt{तद्भावितं}}{\lemma{\skt{तद्}}\Afootnote{H; \skt{तद्}\ A}}\skt{ }\edtext{\skt{सारगणे}}{\lemma{\skt{सारगणे}}\Afootnote{H; \skt{सारगणैर्}\ A}}\skt{ }\edtext{\skt{हृतदोषादिना}}{\lemma{\skt{हृतदोषादिना}}\Afootnote{H; \skt{हृतदोषो दिनो}\ A}}\skt{ }\edtext{\skt{दितः}}{\lemma{\skt{दितः}}\Afootnote{H; \skt{दये}\ A}}\skt{~॥}\pend

\pstart\versepar \ledsidenote{\passref{4.13.11}}\skt{पिवे\-त्सा\-रो\-द\-के\-नैव }\edtext{\skt{श्लक्ष्णम्पिष्टं}}{\lemma{\skt{श्लक्ष्णम्}}\Afootnote{H; \skt{श्लक्ष्ण}\ A}}\skt{ यथा वलं~। }\pend
\pstart\versepar \skt{जाङ्गलेन }\edtext{\skt{रसे\-ना\-न्न\-न्त\-स्मि\-न्जीर्ण्णे}}{\lemma{\skt{तस्मिन्}}\Afootnote{H; \skt{तस्मिञ्}\ A}}\skt{ तु भोजयेत्~॥}\pend

\pstart\versepar \ledsidenote{\passref{4.13.12}}\skt{उपयुज्य }\edtext{\skt{तुलामेव}}{\lemma{\skt{एव}}\Afootnote{H; \skt{एवं}\ A}}\skt{ }\edtext{\skt{अमृ\-त\-स्या\-स्म\-ज\-न्मनः}}{\lemma{\skt{अमृ\appbrk{}त\appbrk{}स्या\appbrk{}स्म\appbrk{}ज\appbrk{}न्मनः}}\Afootnote{H; \skt{गिरिजादमृ}\ A}}\skt{~। }\pend
\pstart\versepar \skt{}\edtext{\edtext{\skt{व्ययो\-म\-धु\-मे\-हा\-ख्य\-म\-न्तकं}}{\lemma{\skt{व्ययो\appbrk{}म\appbrk{}धु\appbrk{}मे\appbrk{}हा\appbrk{}ख्यम्}}\Afootnote{H; \skt{तोपमा}\ A}}}{\lemma{\skt{अन्तकं}}\Afootnote{H; \skt{त्}\ A}}\skt{ }\edtext{\skt{रोगसंकरं}}{\lemma{\skt{रोगसंकरं}}\Afootnote{H; \emph{om.}\ A}}\skt{~॥ }\pend
\pstart\versepar \skt{वपु\-र्व्व\-र्ण्ण\-व\-लो\-पेतः }\edtext{\skt{शतं}}{\lemma{\skt{शतं}}\Afootnote{H; \skt{मधु\appbrk{}मे\appbrk{}ह\appbrk{}वि\appbrk{}व\appbrk{}र्जितः}\ A}}\skt{ }\edtext{\skt{जीवत्यनामयः}}{\lemma{\skt{जीवत्यनामयः}}\Afootnote{\emph{om.}\ A; \skt{जीवत्यनामयः}\ H}}\skt{~।}\pend

\pstart \ledsidenote{\passref{4.13.13cd}}\skt{शतं शतं तुलायान्तु सहस्रं }\edtext{\skt{दशतौलिकं}}{\lemma{\skt{दशतौलिकं}}\Afootnote{H; \skt{दशतौलिके}\ A}}\skt{~॥}\pend

\pstart\versepar \ledsidenote{\passref{4.13.14}}\skt{भल्ला\-त\-क\-वि\-धा\-नेन परिहारविधिः स्मृतः~। }\pend
\pstart\versepar \skt{मेहं }\edtext{\skt{कुष्थ\-म\-प\-स्मा\-र\-मु\-न्मादं}}{\lemma{\skt{कुष्थ\appbrk{}म\appbrk{}प\appbrk{}स्मा\appbrk{}रम्}}\Afootnote{H; \skt{कुष्ठ\appbrk{}म\appbrk{}प\appbrk{}स्मा\appbrk{}रम्}\ A}}\skt{ }\edtext{\skt{शोफश्लीपदं}}{\lemma{\skt{शोफश्लीपदं}}\Afootnote{H; \skt{श्लीपदं गरम्}\ A}}\skt{~॥}\pend

\pstart \ledsidenote{\passref{4.13.15cd}}\edtext{\skt{व्ययो}}{\lemma{\skt{व्ययो}}\Afootnote{H; \skt{अपो}\ A}}\skt{ }\edtext{\skt{हत्य\-चि\-रा\-त्का\-ला\-च्छि\-ला\-जतु}}{\lemma{\skt{हत्यचिरात्}}\Afootnote{H; \skt{हत्यचिरात्}\ A}}\skt{ निषेवितं~।}\pend

\pstart\versepar \ledsidenote{\passref{4.13.16}}\edtext{\skt{न सोस्ति}}{\lemma{\skt{न सोस्ति}}\Afootnote{H; \skt{न सोऽस्ति}\ A}}\skt{ रोगो }\edtext{\skt{यञ्चापि न}}{\lemma{\skt{चापि न}}\Afootnote{H; \skt{चापि}\ A}}\skt{ }\edtext{\edtext{\skt{निह\-न्या\-च्छि\-ला\-जतुः}}{\lemma{\skt{निहन्याच्}}\Afootnote{H; \skt{निहन्यान्न}\ A}}}{\lemma{\skt{छिलाजतुः}}\Afootnote{H; \skt{शिलाजतु}\ A}}\skt{~॥ }\pend
\pstart\versepar \skt{}\edtext{\edtext{\skt{शर्क्क\-रा\-ञ्चि\-र\-सं\-वृ\-द्वा\-म\-श्म\-रीञ्च}}{\lemma{\skt{चिरसंवृद्वाम्}}\Afootnote{H; \skt{चिरसंभूतां भिनत्ति च तथा}\ A}}}{\lemma{\skt{अश्मरीञ्च}}\Afootnote{H; \skt{ऽश्मरीम्}\ A}}\skt{ }\edtext{\edtext{\skt{भिनत्यपि}}{\lemma{\skt{भिनत्य्}}\Afootnote{H; \emph{om.}\ A}}}{\lemma{\skt{अपि}}\Afootnote{H; \emph{om.}\ A}}\skt{~।}\pend

\pstart\versepar \ledsidenote{\passref{4.13.17}}\skt{भावनालोडने }\edtext{\skt{चैव}}{\lemma{\skt{चैव}}\Afootnote{H; \skt{चास्}\ A}}\skt{ }\edtext{\skt{पथ्याः}}{\lemma{\skt{पथ्याः}}\Afootnote{H; \skt{य}\ A}}\skt{ }\edtext{\skt{कार्यास्तु ये}}{\lemma{\skt{कार्यास्तु ये}}\Afootnote{H; \skt{कर्}\ A}}\skt{ }\edtext{\skt{गणाः}}{\lemma{\skt{गणाः}}\Afootnote{H; \emph{om.}\ A}}\skt{~॥ }\pend
\pstart\versepar \skt{}\edtext{\edtext{\edtext{\skt{एव\-ञ्चा\-द्भू\-द्भ\-व\-न्धा\-तु\-म्मा\-क्षिकं}}{\lemma{\skt{एवञ्}}\Afootnote{H; \emph{om.}\ A}}}{\lemma{\skt{चाद्भूद्भवन्}}\Afootnote{H; \skt{तव्ये भेषजैर्हितैः~। एवं च}\ A}}}{\lemma{\skt{धातुम्}}\Afootnote{H; \emph{om.}\ A}}\skt{ }\edtext{\skt{मृतोपमं}}{\lemma{\skt{मृतोपमं}}\Afootnote{H; \skt{धातुं तापी\appbrk{}ज\appbrk{}म\appbrk{}मृ\appbrk{}तो\appbrk{}प\appbrk{}मम्}\ A}}\skt{~।}\pend

\pstart\versepar \ledsidenote{\passref{4.13.18}}\edtext{\edtext{\skt{मधु\-र\-ञ्का\-ञ्च\-ना\-भा\-स\-म्लम्वा}}{\lemma{\skt{मधुरञ्}}\Afootnote{H; \skt{मधुरं}\ A}}}{\lemma{\skt{काञ्च\appbrk{}ना\appbrk{}भा\appbrk{}स\appbrk{}म्लम्वा}}\Afootnote{H; \skt{काञ्च\appbrk{}ना\appbrk{}भा\appbrk{}स\appbrk{}मम्लं वा}\ A}}\skt{ रजतप्रभं~॥ }\pend
\pstart\versepar \skt{पिवन्हन्ति }\edtext{\skt{जरां}}{\lemma{\skt{जरां}}\Afootnote{H; \skt{जरा}\ A}}\skt{ }\edtext{\skt{कुष्ठं}}{\lemma{\skt{कुष्ठं}}\Afootnote{H; \skt{कुष्ठ}\ A}}\skt{ }\edtext{\edtext{\skt{शोफ\-पा\-ण्ड्वा\-म\-य\-ञ्क्षरं}}{\lemma{\skt{शोफ\appbrk{}पा\appbrk{}ण्ड्वा\appbrk{}म\appbrk{}यञ्}}\Afootnote{H; \skt{मेहपाण्ड्वामय}\ A}}}{\lemma{\skt{क्षरं}}\Afootnote{H; \skt{क्षयान्}\ A}}\skt{~।}\pend

\pstart\versepar \ledsidenote{\passref{4.13.19}}\edtext{\edtext{\skt{तद्भा\-वि\-त\-कु\-ल\-काश्च}}{\lemma{\skt{तद्}}\Afootnote{H; \skt{तद्}\ A}}}{\lemma{\skt{भावि\appbrk{}त\appbrk{}कु\appbrk{}ल\appbrk{}काश्च}}\Afootnote{H; \skt{भावितः कपोतांश्च कुलत्थांश्च}\ A}}\skt{ विवर्ज्जयेत्~॥ }\pend
\pstart\versepar \skt{पञ्च\-क\-र्म्म\-गु\-णा\-तीतं श्रद्धावन्तं }\edtext{\skt{जिजीविषं}}{\lemma{\skt{जिजीविषं}}\Afootnote{H; \skt{जिजीविषुम्}\ A}}\skt{~॥}\pend

\pstart\versepar \ledsidenote{\passref{4.13.20}}\skt{योगेनानेन }\edtext{\edtext{\edtext{\skt{मति\-मा\-न्सा\-ध\-ये\-त्कु\-ष्ठि\-न\-न्नरं}}{\lemma{\skt{साधयेत्}}\Afootnote{H; \skt{साधयेदपि}\ A}}}{\lemma{\skt{कुष्ठिनन्}}\Afootnote{H; \skt{कुष्ठनम्}\ A}}}{\lemma{\skt{नरं}}\Afootnote{H; \emph{om.}\ A}}\skt{~। }\pend
\pstart\versepar \skt{}\edtext{\skt{वृक्षास्तौवरिका}}{\lemma{\skt{तौवरिका}}\Afootnote{H; \skt{तुवरका}\ A}}\skt{ }\edtext{\skt{नाम}}{\lemma{\skt{नाम}}\Afootnote{H; \skt{ये स्युः}\ A}}\skt{ }\edtext{\skt{पश्चि\-मा\-र्ण्ण\-व\-ती\-रजाः}}{\lemma{\skt{पश्चि\appbrk{}मा\appbrk{}र्ण्ण\appbrk{}व\appbrk{}ती\appbrk{}रजाः}}\Afootnote{H; \skt{पश्चि\appbrk{}मा\appbrk{}र्ण\appbrk{}व\appbrk{}भू\appbrk{}मिषु}\ A}}\skt{~॥}\pend

\pstart \ledsidenote{\passref{4.13.21cd}}\skt{तेषां }\edtext{\skt{फलान्याददीत}}{\lemma{\skt{आददीत}}\Afootnote{H; \skt{गृह्णीयात्}\ A}}\skt{ }\edtext{\edtext{\skt{सुप\-क्वा\-न्य\-म्वु\-दा\-गमे}}{\lemma{\skt{सुपक्वान्य्}}\Afootnote{H; \skt{सुपक्वान्}\ A}}}{\lemma{\skt{अम्वुदागमे}}\Afootnote{H; \skt{यम्बुदागमे}\ A}}\skt{~।}\pend

\pstart\versepar \ledsidenote{\passref{4.13.22}}\edtext{\skt{मज्जां}}{\lemma{\skt{मज्जां}}\Afootnote{A; \skt{ः}\ H}}\skt{ }\edtext{\skt{फलेभ्यः स}}{\lemma{\skt{फलेभ्यः स}}\Afootnote{H; \skt{तेभ्योऽ}\ A}}\skt{ }\edtext{\skt{विशोष्य}}{\lemma{\skt{विशोष्य}}\Afootnote{H; \skt{पि संहृत्य शोषयि}\ A}}\skt{ }\edtext{\skt{परिचूर्ण्णयेत्}}{\lemma{\skt{परिचूर्ण्णयेत्}}\Afootnote{H; \skt{त्वा विचूर्ण्य च}\ A}}\skt{~॥ }\pend
\pstart\versepar \skt{}\edtext{\skt{तिल\-व\-त्पी\-ड\-ये\-द्द्रोण्यां}}{\lemma{\skt{द्रोण्यां}}\Afootnote{A; \skt{रोण्यां}\ H}}\skt{ }\edtext{\skt{गालयेद्वा}}{\lemma{\skt{गालयेद्वा}}\Afootnote{H; \skt{स्रावयेद्वा}\ A}}\skt{ कुसुम्भवत्~।}\pend

\pstart\versepar \ledsidenote{\passref{4.13.23}}\skt{तत्तैलं }\edtext{\skt{संभृतं}}{\lemma{\skt{संभृतं}}\Afootnote{H; \skt{संहृतं}\ A}}\skt{ भूयः पचे\-दा\-दु\-द\-क\-क्ष\-यात्~॥ }\pend
\pstart\versepar \skt{अवतार्य करीषे च पक्षमात्रं निधापयेत्~।}\edtext{}{\lemma{\skt{पचे\appbrk{}दा\appbrk{}दु\appbrk{}द\appbrk{}क\appbrk{}क्ष\appbrk{}यात्~॥ अवतार्य}}\Afootnote{H; \skt{पचे\appbrk{}दा\appbrk{}तो\appbrk{}य\appbrk{}सं\appbrk{}क्ष\appbrk{}यात्~। अवतार्य}\ A}}\pend

\pstart\versepar \ledsidenote{\passref{4.13.24}}\edtext{\skt{स्निग्धं}}{\lemma{\skt{स्निग्धं}}\Afootnote{H; \skt{स्निग्धः}\ A}}\skt{ स्विन्नो हृतमलः }\edtext{\skt{पक्षादुद्धृत्य}}{\lemma{\skt{पक्षादुद्धृत्य}}\Afootnote{H; \skt{पक्षादूर्ध्वं प्र}\ A}}\skt{ }\edtext{\skt{यत्नतः}}{\lemma{\skt{यत्नतः}}\Afootnote{H; \skt{यत्नवान्}\ A}}\skt{~॥ }\pend
\pstart\versepar \skt{चतु\-र्थ\-भ\-क्ता\-न्त\-रितः शुक्लादौ दिवसे शुभे~।}\pend

\pstart\versepar \ledsidenote{\passref{4.13.25}}\skt{मन्त्रपूतस्य तैलस्य }\edtext{\edtext{\skt{प्राश्नियात्पाणि}}{\lemma{\skt{प्राश्नियात्}}\Afootnote{H; \emph{om.}\ A}}}{\lemma{\skt{पाणि}}\Afootnote{H; \skt{पिबेन्}\ A}}\skt{ }\edtext{\skt{मात्रिकां}}{\lemma{\skt{मात्रिकां}}\Afootnote{H; \skt{मात्रां यथाबलम्}\ A}}\skt{~॥ }\pend
\pstart\versepar \skt{तत्र मन्त्रम्प्रवक्ष्यामि येने\-द\-म\-भि\-म\-न्त्र्यते~।}\pend

\pstart\versepar \ledsidenote{\passref{4.13.26}}\edtext{\skt{मज्जासार}}{\lemma{\skt{मज्जासार}}\Afootnote{H; \skt{मज्जसार}\ A}}\skt{ महावीर्य सर्वा\-न्धा\-तू\-न्वि\-शो\-धय~॥ }\pend
\pstart\versepar \skt{शङ्ख\-च\-क्र\-ग\-दा\-पा\-णि\-स्त्वा\-मा\-ज्ञा\-प\-य\-तेऽ\-च्युतः~।}\pend

\pstart\versepar \ledsidenote{\passref{4.13.27}}\edtext{\edtext{\skt{तेना\-स्यो\-र्द्ध\-म\-ध\-श्चापि}}{\lemma{\skt{तेनास्योर्द्धम्}}\Afootnote{H; \skt{तेनास्योर्ध्व}\ A}}}{\lemma{\skt{अधश्}}\Afootnote{H; \skt{मधश्}\ A}}\skt{ दोषा }\edtext{\edtext{\skt{यान्त्यसकृत्ततः}}{\lemma{\skt{यान्त्य्}}\Afootnote{H; \skt{यान्त्}\ A}}}{\lemma{\skt{असकृत्ततः}}\Afootnote{H; \skt{यसकृत्ततः}\ A}}\skt{~॥ }\pend
\pstart\versepar \skt{}\edtext{\skt{साय\-म\-स्ने\-ह\-ल\-वणां}}{\lemma{\skt{सायम्}}\Afootnote{H; \emph{om.}\ A}}\skt{ }\edtext{\skt{यवागूं}}{\lemma{\skt{यवागूं}}\Afootnote{H; \skt{सायं यवागूं}\ A}}\skt{ शीतलाम्पिवेत्~।}\pend

\pstart\versepar \ledsidenote{\passref{4.13.28}}\edtext{\skt{पञ्चहानि}}{\lemma{\skt{पञ्चहानि}}\Afootnote{H; \skt{पञ्चाहं प्र}\ A}}\skt{ }\edtext{\edtext{\skt{पिवे\-त्तै\-ल\-म\-नेन}}{\lemma{\skt{तैलम्}}\Afootnote{H; \skt{तैल}\ A}}}{\lemma{\skt{अनेन}}\Afootnote{H; \skt{मनेन}\ A}}\skt{ विधिना नरः~॥ }\pend
\pstart\versepar \skt{}\edtext{\skt{पक्वं}}{\lemma{\skt{पक्वं}}\Afootnote{H; \skt{पक्षं}\ A}}\skt{ परिहरेच्चापि मुद्ग\-यू\-षौ\-द\-ना\-शनः~।}\pend

\pstart\versepar \ledsidenote{\passref{4.13.29}}\skt{पञ्च\-भि\-र्द्दि\-व\-सै\-रेवं सर्व्वकुष्ठैः }\edtext{\skt{प्रमुच्यते}}{\lemma{\skt{प्रमुच्यते}}\Afootnote{H; \skt{विमुच्यते}\ A}}\skt{~॥ }\pend
\pstart\versepar \skt{तदेव खदिरक्वाथे त्रिगुणे साधु साधितं~॥}\pend

\pstart\versepar \ledsidenote{\passref{4.13.30}}\edtext{\skt{निहि\-त\-म्पू\-र्व्व\-व\-त्पक्षं}}{\lemma{\skt{पक्षं}}\Afootnote{H; \skt{पक्षात्}\ A}}\skt{ }\edtext{\edtext{\edtext{\skt{पिवे\-त्सा\-म\-म\-त\-न्द्रितः}}{\lemma{\skt{पिवेत्}}\Afootnote{H; \skt{पिबेन्}\ A}}}{\lemma{\skt{सामम्}}\Afootnote{H; \skt{मास}\ A}}}{\lemma{\skt{अतन्द्रितः}}\Afootnote{H; \skt{मतन्द्रितः}\ A}}\skt{~॥ }\pend
\pstart\versepar \skt{}\edtext{\skt{तेना\-भ्य\-क्त\-श\-री\-रस्तु}}{\lemma{\skt{तेना\appbrk{}भ्य\appbrk{}क्त\appbrk{}श\appbrk{}री\appbrk{}रस्तु}}\Afootnote{H; \skt{तेना\appbrk{}भ्य\appbrk{}क्त\appbrk{}श\appbrk{}री\appbrk{}रश्च}\ A}}\skt{ }\edtext{\edtext{\skt{कुर्व्वी\-ता\-हा\-र\-मी\-रितं}}{\lemma{\skt{कुर्व्वीताहारम्}}\Afootnote{H; \skt{कुर्वीताहार}\ A}}}{\lemma{\skt{ईरितं}}\Afootnote{H; \skt{मीरितम्}\ A}}\skt{~।}\pend

\pstart\versepar \ledsidenote{\passref{4.13.31}}\skt{भिन्नस्वरं }\edtext{\skt{रक्त\-ने\-त्र\-म्प्र\-शीर्णं}}{\lemma{\skt{प्रशीर्णं}}\Afootnote{H; \skt{विशीर्णं}\ A}}\skt{ }\edtext{\skt{क्रिमि}}{\lemma{\skt{क्रिमि}}\Afootnote{H; \skt{कृमि}\ A}}\skt{ भक्षितं~॥ }\pend
\pstart\versepar \skt{अनेनाशु }\edtext{\skt{प्रयोगेन}}{\lemma{\skt{प्रयोगेन}}\Afootnote{H; \skt{प्रयोगेण}\ A}}\skt{ साध\-ये\-त्कु\-ष्ठि\-न\-न्नरं~।}\pend

\pstart\versepar \ledsidenote{\passref{4.13.32}}\edtext{\edtext{\skt{सर्प्पि\-र्म्म\-धु\-यु\-त\-म्पी\-त\-न्त\-देव}}{\lemma{\skt{सर्प्पिर्}}\Afootnote{H; \skt{सर्पि}\ A}}}{\lemma{\skt{म्मधुयुतम्}}\Afootnote{H; \skt{र्मधुयुतं}\ A}}\skt{ खदिराम्वुना~॥ }\pend
\pstart\versepar \skt{}\edtext{\skt{पक्ष\-म्मां\-स\-र\-सा\-हारं}}{\lemma{\skt{पक्षम्}}\Afootnote{H; \skt{पक्षि}\ A}}\skt{ करोति द्विशतायुषं~।}\pend

\pstart\versepar \ledsidenote{\passref{4.13.33}}\skt{तदेव नस्ये }\edtext{\edtext{\skt{पञ्चा\-श\-द्दि\-व\-सा\-न्यु\-प\-यो\-जितं}}{\lemma{\skt{पञ्चाशद्}}\Afootnote{H; \skt{पञ्चाशद्}\ A}}}{\lemma{\skt{दिव\appbrk{}सा\appbrk{}न्यु\appbrk{}प\appbrk{}यो\appbrk{}जितं}}\Afootnote{H; \skt{दिव\appbrk{}सा\appbrk{}नु\appbrk{}प\appbrk{}यो\appbrk{}जि\appbrk{}तम्}\ A}}\skt{~॥ }\pend
\pstart\versepar \skt{वपुष्मन्तं श्रुति\-ध\-र\-ङ्क\-रोति त्रिशतायुषं~।}\pend

\pstart \ledsidenote{\passref{4.13.34ab}}\edtext{\skt{शोधयेन्तु}}{\lemma{\skt{शोधयेन्तु}}\Afootnote{H; \skt{शोधयन्ति}\ A}}\skt{ }\edtext{\skt{नरम्मज्जा}}{\lemma{\skt{मज्जा}}\Afootnote{H; \skt{पीता मज्जा}\ A}}\skt{ }\edtext{\skt{तस्य}}{\lemma{\skt{तस्य}}\Afootnote{H; \skt{नस्}\ A}}\skt{ }\edtext{\skt{पीतस्य}}{\lemma{\skt{पीतस्य}}\Afootnote{H; \skt{तस्य}\ A}}\skt{ मात्रया~॥}\pend

\pstart\versepar \ledsidenote{\passref{4.13.35}}\edtext{\skt{शान्तर्द्धूमा}}{\lemma{\skt{शान्तर्द्धूमा}}\Afootnote{H; \skt{सान्तर्धूमस्}\ A}}\skt{ तस्य मज्जा तु }\edtext{\skt{दग्धा}}{\lemma{\skt{दग्धा}}\Afootnote{H; \skt{दग्धः}\ A}}\skt{ }\edtext{\skt{तैले}}{\lemma{\skt{तैले}}\Afootnote{H; \emph{om.}\ A}}\skt{ }\edtext{\skt{क्षिप्त्वा}}{\lemma{\skt{क्षिप्त्वा}}\Afootnote{H; \skt{क्षिप्तस्तैले}\ A}}\skt{ सैन्धवं }\edtext{\skt{चाग्निवर्ण्णं}}{\lemma{\skt{चाग्निवर्ण्णं}}\Afootnote{H; \skt{चाञ्जनं च}\ A}}\skt{~। }\pend
\pstart\versepar \skt{}\edtext{\skt{पिल्लं}}{\lemma{\skt{पिल्लं}}\Afootnote{H; \skt{पैल्ल्यं}\ A}}\skt{ }\edtext{\edtext{\skt{हन्या\-द\-र्म्म\-न\-का\-न्ध्य\-का\-चा\-न्ने\-त्र\-दक्तं}}{\lemma{\skt{हन्या\appbrk{}द\appbrk{}र्म्म\appbrk{}न\appbrk{}का\appbrk{}न्ध्य\appbrk{}का\appbrk{}चान्}}\Afootnote{H; \skt{हन्या\appbrk{}द\appbrk{}र्म\appbrk{}न\appbrk{}क्ता\appbrk{}न्ध्य\appbrk{}का\appbrk{}चान्}\ A}}}{\lemma{\skt{नेत्रदक्तं}}\Afootnote{H; \skt{नीलीरोगं}\ A}}\skt{ }\edtext{\edtext{\skt{तिमि\-र\-न्नी\-लि\-का\-ञ्जति}}{\lemma{\skt{तिमिरन्}}\Afootnote{H; \skt{तैमिरं}\ A}}}{\lemma{\skt{नीलिकाञ्जति}}\Afootnote{H; \skt{चाञ्जनेन}\ A}}\skt{~॥ ०~॥}\pend

\pstart \ledsidenote{\passref{4.13.trailer}}\edtext{\skt{चिकित्सा}}{\lemma{\skt{चिकित्सा}}\Afootnote{H; \skt{इति सुश्रु\appbrk{}त\appbrk{}सं\appbrk{}हि\appbrk{}तायां चिकित्सा}\ A}}\skt{ }\edtext{\skt{त्रि॥}}{\lemma{\skt{त्रि}}\Afootnote{H; \skt{स्थाने मधु\appbrk{}मे\appbrk{}ह\appbrk{}चि\appbrk{}कि\appbrk{}त्सितं नाम त्रयोदशोऽध्यायः}\ A}}\skt{ ०~॥}\pend

\endnumbering
\beginnumbering
\eledsection*{Adhyāya 14}
\pstart \ledsidenote{\passref{4.14.1}}\skt{अथात उद\-रा\-णा\-ञ्चि\-कि\-त्सि\-त\-म्व्या\-ख्या\-स्यामः~॥}\pend

\pstart\versepar \ledsidenote{\passref{4.14.3}}\skt{अष्टावुदराणि }\edtext{\skt{पूर्व्वमुक्ताणि}}{\lemma{\skt{उक्ताणि}}\Afootnote{H; \skt{उद्दिष्टानि}\ A}}\skt{ }\pend
\pstart\versepar \skt{तेष्व\-सा\-ध्य\-म्ब\-द्ध\-गुदं }\edtext{\skt{परिश्रावि च}}{\lemma{\skt{परिश्रावि च}}\Afootnote{H; \skt{परिस्रावि च}\ A}}\skt{~। अवशिष्टानि कृच्छ्रसाध्यानि सर्व्वाण्येव प्रत्या\-ख्या\-यो\-प\-क्र\-मेत~। तेष्वा\-द्य\-श्च\-तु\-र्व्वर्ग्गो भेषजसाध्यः~। उत्तरः शस्त्रसाध्यः~। }\edtext{\edtext{\skt{काल\-प्र\-क\-र्षा\-त्स\-र्व्वा\-ण्येव}}{\lemma{\skt{सर्व्वाण्य्}}\Afootnote{H; \skt{सर्वाण्}\ A}}}{\lemma{\skt{एव}}\Afootnote{H; \skt{येव}\ A}}\skt{ शस्त्रसाध्यानि }\edtext{\skt{भवन्ति}}{\lemma{\skt{भवन्ति}}\Afootnote{H; \emph{om.}\ A}}\skt{~। वर्ज्जयितव्यानि वा~।}\pend

\pstart \ledsidenote{\passref{4.14.4}}\skt{उदरी तु }\edtext{\skt{गुर्व\-भि\-ष्य\-न्दि\-रू\-क्ष\-वि\-दा\-ह्य\-ति\-स्नि\-ग्ध\-पि\-शि\-त\-प\-रि\-षे\-का\-व\-गा\-हा\-न्प\-रि\-ह\-रेत्}}{\lemma{\skt{गुर्व\appbrk{}भि\appbrk{}ष्य\appbrk{}न्दि\appbrk{}रू\appbrk{}क्ष\appbrk{}वि\appbrk{}दा\appbrk{}ह्य\appbrk{}ति\appbrk{}स्नि\appbrk{}ग्ध\appbrk{}पि\appbrk{}शि\appbrk{}त\appbrk{}प\appbrk{}रि\appbrk{}षे\appbrk{}का\appbrk{}व\appbrk{}गा\appbrk{}हान्}}\Afootnote{H; \skt{गुर्व\appbrk{}भि\appbrk{}ष्य\appbrk{}न्दि\appbrk{}रू\appbrk{}क्ष\appbrk{}वि\appbrk{}दा\appbrk{}हि\appbrk{}स्नि\appbrk{}ग्ध\appbrk{}पि\appbrk{}शि\appbrk{}त\appbrk{}प\appbrk{}रि\appbrk{}षे\appbrk{}का\appbrk{}व\appbrk{}गा\appbrk{}हान्}\ A}}\skt{~। शालि\-ष\-ष्टि\-क\-य\-व\-गो\-धू\-म\-नी\-वा\-रा\-न्नि\-त्य\-म\-श्नी\-यात्~।}\pend

\pstart \ledsidenote{\passref{4.14.5}}\skt{तत्र वातो\-द\-रि\-ण\-म्वि\-दा\-रि\-ग\-न्धा\-दि\-सि\-द्धेन सर्प्पिषा स्नेहयित्वा, तिल्व\-क\-वि\-प\-क्वे\-ना\-नु\-लोम्य चित्रा\-फ\-ल\-तै\-ल\-प्र\-गा\-ढेन }\edtext{\skt{विदा\-रि\-ग\-न्धा\-दि\-क\-षा\-ये\-णा\-स्था\-प\-ये\-द\-नु\-वा\-स\-येच्च}}{\lemma{\skt{विदा\appbrk{}रि\appbrk{}ग\appbrk{}न्धा\appbrk{}दि\appbrk{}क\appbrk{}षा\appbrk{}ये\appbrk{}णा\appbrk{}स्था\appbrk{}प\appbrk{}ये\appbrk{}द\appbrk{}नु\appbrk{}वा\appbrk{}स\appbrk{}येच्च}}\Afootnote{H; \skt{विदा\appbrk{}रि\appbrk{}ग\appbrk{}न्दा\appbrk{}दि\appbrk{}क\appbrk{}षा\appbrk{}ये\appbrk{}णा\appbrk{}स्था\appbrk{}प\appbrk{}ये\appbrk{}द\appbrk{}नु\appbrk{}वा\appbrk{}स\appbrk{}येच्च}\ A}}\skt{~। }\edtext{\skt{साल्व\-णे\-नो\-प\-ना\-ह\-ये\-दु\-दरं}}{\lemma{\skt{साल्व\appbrk{}णे\appbrk{}नो\appbrk{}प\appbrk{}ना\appbrk{}ह\appbrk{}ये\appbrk{}दु\appbrk{}दरं}}\Afootnote{H; \skt{साल्वणेन चोप\appbrk{}ना\appbrk{}ह\appbrk{}ये\appbrk{}दु\appbrk{}दरं}\ A}}\skt{~। भोज\-ये\-च्चै\-न\-म्वि\-दा\-रि\-ग\-न्धा\-दि\-सि\-द्धेन }\edtext{\skt{पयसा}}{\lemma{\skt{पयसा}}\Afootnote{H; \skt{क्षीरेण}\ A}}\skt{ }\edtext{\skt{जाङ्गलरसेन वा}}{\lemma{\skt{जाङ्गलरसेन वा}}\Afootnote{H; \skt{जाङ्गलरसेन च}\ A}}\skt{~। स्वेद\-ये\-च्चा\-भी\-क्ष्णम्~।}\pend

\pstart \ledsidenote{\passref{4.14.6}}\skt{पित्तोदरिणं तु मधु\-र\-ग\-ण\-वि\-प\-क्वेन सर्प्पिषा स्नेहयित्वा }\edtext{\skt{श्यामा\-त्रि\-वृ\-त्त्रि\-फ\-ला\-वि\-प\-क्वे\-ना\-नु\-लोम्य}}{\lemma{\skt{श्यामा\appbrk{}त्रि\appbrk{}वृ\appbrk{}त्त्रि\appbrk{}फ\appbrk{}ला\appbrk{}वि\appbrk{}प\appbrk{}क्वे\appbrk{}ना\appbrk{}नु\appbrk{}लोम्य}}\Afootnote{H; \skt{श्यामा\appbrk{}त्रि\appbrk{}फ\appbrk{}ला\appbrk{}त्रि\appbrk{}वृ\appbrk{}द्वि\appbrk{}प\appbrk{}क्वे\appbrk{}ना\appbrk{}नु\appbrk{}लोम्य}\ A}}\skt{ शर्क्क\-रा\-म\-धु\-घृ\-त\-प्र\-गा\-ढेन न्यग्रो\-धा\-दि\-क\-षा\-ये\-णा\-स्था\-प\-ये\-द\-नु\-वा\-स\-येच्च~। पाय\-से\-नो\-प\-ना\-ह\-ये\-दु\-दरं~। भोजयेच्चैनं विदा\-रि\-ग\-न्धा\-दि\-सि\-द्धेन }\edtext{\skt{क्षीरेण}}{\lemma{\skt{क्षीरेण}}\Afootnote{H; \skt{पयसा}\ A}}\skt{~॥}\pend

\pstart \ledsidenote{\passref{4.14.7}}\skt{श्लेष्मोदरिणन्तु }\edtext{\skt{पिप्प\-ल्या\-दि\-सि\-द्धेन}}{\lemma{\skt{पिप्प\appbrk{}ल्या\appbrk{}दि\appbrk{}सि\appbrk{}द्धेन}}\Afootnote{H; \skt{पिप्प\appbrk{}ल्या\appbrk{}दि\appbrk{}क\appbrk{}षा\appbrk{}य\appbrk{}सि\appbrk{}द्धेन}\ A}}\skt{ }\edtext{\skt{सर्प्पिषा}}{\lemma{\skt{सर्प्पिषा}}\Afootnote{H; \skt{सर्पिषो}\ A}}\skt{ }\edtext{\skt{स्नेहयित्वा}}{\lemma{\skt{स्नेहयित्वा}}\Afootnote{H; \skt{पस्नेह्य}\ A}}\skt{ }\edtext{\skt{स्नुहा\-क्षी\-र\-वि\-प\-क्वेन}}{\lemma{\skt{स्नुहा\appbrk{}क्षी\appbrk{}र\appbrk{}वि\appbrk{}प\appbrk{}क्वेन}}\Afootnote{H; \skt{स्नुही\appbrk{}क्षी\appbrk{}र\appbrk{}वि\appbrk{}पक्वे}\ A}}\skt{ }\edtext{\skt{चानुलोम्य}}{\lemma{\skt{चानुलोम्य}}\Afootnote{H; \skt{नानुलोम्य}\ A}}\skt{ }\edtext{\skt{त्रिक\-टु\-क\-मू\-त्र\-तै\-ल\-प्र\-गा\-ढेन}}{\lemma{\skt{त्रिक\appbrk{}टु\appbrk{}क\appbrk{}मू\appbrk{}त्र\appbrk{}तै\appbrk{}ल\appbrk{}प्र\appbrk{}गा\appbrk{}ढेन}}\Afootnote{H; \skt{त्रिक\appbrk{}टु\appbrk{}क\appbrk{}मू\appbrk{}त्र\appbrk{}क्षा\appbrk{}र\appbrk{}तै\appbrk{}ल\appbrk{}प्र\appbrk{}गा\appbrk{}ढेन}\ A}}\skt{ }\edtext{\skt{मुष्क\-का\-दि\-क्वा\-थे\-ना\-स्था\-प\-ये\-द\-नु\-वा\-स\-येच्च}}{\lemma{\skt{मुष्क\appbrk{}का\appbrk{}दि\appbrk{}क्वा\appbrk{}थे\appbrk{}ना\appbrk{}स्था\appbrk{}प\appbrk{}ये\appbrk{}द\appbrk{}नु\appbrk{}वा\appbrk{}स\appbrk{}येच्च}}\Afootnote{H; \skt{मुष्क\appbrk{}का\appbrk{}दि\appbrk{}क\appbrk{}षा\appbrk{}ये\appbrk{}णा\appbrk{}स्था\appbrk{}प\appbrk{}ये\appbrk{}द\appbrk{}नु\appbrk{}वा\appbrk{}स\appbrk{}येच्च}\ A}}\skt{~। }\edtext{\skt{कीकि\-ण्व\-स\-र्ष\-प\-मू\-ल\-क\-बी\-जै\-श्चो\-प\-ना\-ह\-ये\-दु\-दरं}}{\lemma{\skt{कीकि\appbrk{}ण्व\appbrk{}स\appbrk{}र्ष\appbrk{}प\appbrk{}मू\appbrk{}ल\appbrk{}क\appbrk{}बी\appbrk{}जैश्}}\Afootnote{H; \skt{शणा\appbrk{}त\appbrk{}सी\appbrk{}धा\appbrk{}त\appbrk{}की\appbrk{}कि\appbrk{}ण्व\appbrk{}स\appbrk{}र्ष\appbrk{}प\appbrk{}मू\appbrk{}ल\appbrk{}क\appbrk{}बी\appbrk{}ज\appbrk{}क\appbrk{}ल्कैश्}\ A}}\skt{~। भोज\-ये\-च्चै\-न\-न्त्रि\-क\-टु\-क\-प्र\-गा\-ढेन कुलत्थयूषेण~। पायसेन वा~। स्वेदयेच्चाभीक्ष्णं~।}\pend

\pstart \ledsidenote{\passref{4.14.8}}\skt{दूष्योदरिणं तु }\edtext{\skt{सप्त\-ला\-शं\-खि\-नी\-स्व\-र\-स\-सि\-द्धेन}}{\lemma{\skt{सप्त\appbrk{}ला\appbrk{}शं\appbrk{}खि\appbrk{}नी\appbrk{}स्व\appbrk{}र\appbrk{}स\appbrk{}सि\appbrk{}द्धेन}}\Afootnote{H; \skt{प्रत्याख्याय सप्त\appbrk{}ला\appbrk{}श\appbrk{}ङ्खि\appbrk{}नी\appbrk{}स्व\appbrk{}र\appbrk{}स\appbrk{}सि\appbrk{}द्धेन}\ A}}\skt{ सर्प्पिषा }\edtext{\skt{विरेचयेत्}}{\lemma{\skt{विरेचयेत्}}\Afootnote{H; \skt{विरचयेन्}\ A}}\skt{~। मास\-म\-र्द्ध\-मा\-सम्वा }\edtext{\skt{महा\-स्नु\-ही\-वृ\-क्ष\-क्षी\-र\-गो\-मू\-त्र\-सि\-द्धेन वा}}{\lemma{\skt{महा\appbrk{}स्नु\appbrk{}ही\appbrk{}वृ\appbrk{}क्ष\appbrk{}क्षी\appbrk{}र\appbrk{}गो\appbrk{}मू\appbrk{}त्र\appbrk{}सि\appbrk{}द्धेन वा}}\Afootnote{\skt{महा\appbrk{}वृ\appbrk{}क्ष\appbrk{}क्षी\appbrk{}र\appbrk{}सु\appbrk{}रा\appbrk{}गो\appbrk{}मू\appbrk{}त्र\appbrk{}सि\appbrk{}द्धेन वा; शुद्धकोष्ठं तु मद्ये\appbrk{}ना\appbrk{}श्व\appbrk{}मा\appbrk{}र\appbrk{}क\appbrk{}गु\appbrk{}ञ्जा\appbrk{}का\appbrk{}का\appbrk{}द\appbrk{}नी\appbrk{}मू\appbrk{}ल\appbrk{}कल्कं पाययेत्; इक्षुकाण्डानि वा कृष्णसर्पेण दंशयित्वा भक्ष\appbrk{}ये\appbrk{}द्व\appbrk{}ल्ली\appbrk{}फ\appbrk{}लानि वा, मूलजं कन्दजं वा विषमा सेवयेत्, तेनागदो भवत्यन्यं वा भावमापद्यते}\ A; \skt{महा\appbrk{}वृ\appbrk{}क्ष\appbrk{}क्षी\appbrk{}र\appbrk{}गो\appbrk{}मू\appbrk{}त्र\appbrk{}सि\appbrk{}द्धेन वा}\ H\textsuperscript{ac}; \skt{स्नुही\appbrk{}वृ\appbrk{}क्ष\appbrk{}क्षी\appbrk{}र\appbrk{}गो\appbrk{}मू\appbrk{}त्र\appbrk{}सि\appbrk{}द्धेन वा}\ H\textsuperscript{pc}}}\skt{~॥}\pend

\pstart\versepar \ledsidenote{\passref{4.14.9}}\edtext{\edtext{\skt{गोपि\-ता\-नि\-ल\-मू\-ल\-त्वा\-त्स\-ञ्चि\-त\-त्वा\-त्म\-लस्य}}{\lemma{\skt{गोपि\appbrk{}ता\appbrk{}नि\appbrk{}ल\appbrk{}मू\appbrk{}ल\appbrk{}त्वात्}}\Afootnote{H; \skt{भवति चात्र कुपि\appbrk{}ता\appbrk{}नि\appbrk{}ल\appbrk{}मू\appbrk{}ल\appbrk{}त्वात्}\ A}}}{\lemma{\skt{सञ्चितत्वात्}}\Afootnote{H; \skt{संचितत्वान्}\ A}}\skt{ च~। }\pend
\pstart\versepar \skt{सर्व्वोदरेषु }\edtext{\skt{संसन्ति}}{\lemma{\skt{संसन्ति}}\Afootnote{H; \skt{शंसन्ति}\ A}}\skt{ बहु\-श\-स्त्व\-नु\-लो\-मनं~॥०॥}\pend

\pstart \ledsidenote{\passref{4.14.10}}\skt{अत }\edtext{\edtext{\edtext{\skt{ऊर्द्ध्व\-न्सा\-मा\-न्या\-न्यो\-गा\-न्व\-क्ष्यामः}}{\lemma{\skt{ऊर्द्ध्वन्}}\Afootnote{H; \skt{ऊर्ध्वं}\ A}}}{\lemma{\skt{सामान्यान्योगान्}}\Afootnote{H; \skt{सामान्ययोगान्}\ A}}}{\lemma{\skt{वक्ष्यामः}}\Afootnote{H; \skt{वक्ष्यामः~। तद्यथा}\ A}}\skt{~॥ }\edtext{\edtext{\skt{एर\-ण्ड\-तै\-ल\-मे\-वा\-ह\-र\-ह\-र्म्मा\-स\-मेकं}}{\lemma{\skt{एवाहरहर्}}\Afootnote{H; \skt{अहरहर्}\ A}}}{\lemma{\skt{एकं}}\Afootnote{H; \emph{om.}\ A}}\skt{ द्वौ वा केवलं मूत्रयुक्तं क्षीरयुक्तं वा सेवे\-तो\-द\-क\-वर्ज्जी~। }\edtext{\skt{महिषीमूत्रं}}{\lemma{\skt{महिषीमूत्रं}}\Afootnote{H; \skt{माहिषं वा मूत्रं}\ A}}\skt{ क्षीरेण निराहारः सप्तरात्रं~। उष्ट्रीक्षीराहारो }\edtext{\skt{वान्न\-पा\-न\-प\-रि\-वा\-रि\-वर्ज्जी}}{\lemma{\skt{वान्न\appbrk{}पा\appbrk{}न\appbrk{}प\appbrk{}रि\appbrk{}वा\appbrk{}रि\appbrk{}वर्ज्जी}}\Afootnote{H; \skt{वाऽन्नवारिवर्जी}\ A}}\skt{ पक्षं~। }\edtext{\skt{पिप्पलीर्व्वा}}{\lemma{\skt{पिप्पलीर्व्वा}}\Afootnote{H; \skt{पिप्पलीं वा}\ A}}\skt{ मासं पूर्व्वोक्तेन }\edtext{\skt{विधानेन}}{\lemma{\skt{विधानेन}}\Afootnote{H; \skt{विधाने}\ A}}\skt{ }\edtext{\skt{वासेवेत}}{\lemma{\skt{वासेवेत}}\Afootnote{H; \skt{नासेवेत}\ A}}\skt{~। }\edtext{\skt{सैन्ध\-वा\-ज\-मो\-द\-नि\-म्ब\-सि\-द्धम्वा}}{\lemma{\skt{सैन्ध\appbrk{}वा\appbrk{}ज\appbrk{}मो\appbrk{}द\appbrk{}नि\appbrk{}म्ब\appbrk{}सि\appbrk{}द्धम्वा}}\Afootnote{H; \skt{सैन्ध\appbrk{}वा\appbrk{}ज\appbrk{}मो\appbrk{}दा\appbrk{}युक्तं वा}\ A}}\skt{ निकुम्भतैलं~। }\edtext{\skt{आर्द्द्र\-क\-ष्ट\-ङ्ग\-ण\-क्षा\-र\-पा\-त्रा\-र्द्ध\-सि\-द्धम्वा}}{\lemma{\skt{आर्द्द्र\appbrk{}क\appbrk{}ष्ट\appbrk{}ङ्ग\appbrk{}ण\appbrk{}क्षा\appbrk{}र\appbrk{}पा\appbrk{}त्रा\appbrk{}र्द्ध\appbrk{}सि\appbrk{}द्धम्वा}}\Afootnote{H; \skt{आर्द्र\appbrk{}शृ\appbrk{}ङ्ग\appbrk{}वे\appbrk{}र\appbrk{}र\appbrk{}स\appbrk{}पात्र}\ A}}\skt{~। }\edtext{\skt{शत\-पु\-ष्प\-सि\-द्धम्वा}}{\lemma{\skt{शत\appbrk{}पु\appbrk{}ष्प\appbrk{}सि\appbrk{}द्धम्वा}}\Afootnote{H; \skt{शतसिद्धं वा}\ A}}\skt{ वातशूले }\edtext{\skt{वचार्यं}}{\lemma{\skt{वचार्यं}}\Afootnote{H; \skt{ऽवचार्यं}\ A}}\skt{~। }\edtext{\skt{शृङ्ग\-वे\-र\-र\-स\-वि\-प\-क्वम्वा}}{\lemma{\skt{शृङ्ग\appbrk{}वे\appbrk{}र\appbrk{}र\appbrk{}स\appbrk{}वि\appbrk{}प\appbrk{}क्वम्वा}}\Afootnote{H; \skt{शृङ्ग\appbrk{}वे\appbrk{}र\appbrk{}र\appbrk{}स\appbrk{}वि\appbrk{}पक्वं}\ A}}\skt{ क्षीरमासेवेत~। }\edtext{\skt{चव्या\-शृ\-ङ्ग\-वे\-र\-क\-ल्कम्वा}}{\lemma{\skt{चव्या\appbrk{}शृ\appbrk{}ङ्ग\appbrk{}वे\appbrk{}र\appbrk{}क\appbrk{}ल्कम्वा}}\Afootnote{H; \skt{चव्य\appbrk{}शृ\appbrk{}ङ्ग\appbrk{}वे\appbrk{}र\appbrk{}कल्कं वा}\ A}}\skt{ पयसा~। }\edtext{\skt{सर\-ला\-म\-र\-दा\-रु\-चि\-त्त्र\-क\-मेव}}{\lemma{\skt{सर\appbrk{}ला\appbrk{}म\appbrk{}र\appbrk{}दा\appbrk{}रु\appbrk{}चि\appbrk{}त्त्र\appbrk{}कम्}}\Afootnote{H; \skt{सर\appbrk{}ल\appbrk{}दे\appbrk{}व\appbrk{}दा\appbrk{}रु\appbrk{}चि\appbrk{}त्र\appbrk{}कम्}\ A}}\skt{ वा~। }\edtext{\skt{मुरु\-ङ्गी\-शा\-ल\-प\-र्ण्णी\-श्या\-मा\-पु\-न\-र्ण्ण\-वा\-ज्यो\-ति\-ष्क\-फ\-ला\-तै\-लम्वा}}{\lemma{\skt{मुरु\appbrk{}ङ्गी\appbrk{}शा\appbrk{}ल\appbrk{}प\appbrk{}र्ण्णी\appbrk{}श्या\appbrk{}मा\appbrk{}पु\appbrk{}न\appbrk{}र्ण्ण\appbrk{}वा\appbrk{}ज्यो\appbrk{}ति\appbrk{}ष्क\appbrk{}फ\appbrk{}ला\appbrk{}तै\appbrk{}लम्वा}}\Afootnote{H; \skt{मुर\appbrk{}ङ्गी\appbrk{}शा\appbrk{}ल\appbrk{}प\appbrk{}र्णी\appbrk{}श्या\appbrk{}मा\appbrk{}पु\appbrk{}न\appbrk{}र्न\appbrk{}वा\appbrk{}कल्कं वा, ज्योति\appbrk{}ष्क\appbrk{}फ\appbrk{}ल\appbrk{}तैलं वा}\ A}}\skt{ क्षीरेण~। स्वर्ज्जि\-का\-हिं\-गु\-मि\-श्र\-म्पि\-बेत्~। }\edtext{\skt{गुडद्वितीयो वा}}{\lemma{\skt{गुडद्वितीयो वा}}\Afootnote{H; \skt{*गुडद्वितीयां वा}\ A}}\skt{ हरी\-त\-की\-म्भ\-क्ष\-येत्~। }\edtext{\skt{स्नुहा\-क्षी\-र\-प\-रि\-श्रा\-वि\-ता\-नाम्वा}}{\lemma{\skt{स्नुहा\appbrk{}क्षी\appbrk{}र\appbrk{}प\appbrk{}रि\appbrk{}श्रा\appbrk{}वि\appbrk{}ता\appbrk{}नाम्वा}}\Afootnote{H; \skt{स्रुती\appbrk{}क्षी\appbrk{}र\appbrk{}भा\appbrk{}वि\appbrk{}तानां वा}\ A}}\skt{ पिप्पलीनां सहस्रङ्कालेन~। }\edtext{\skt{पथ्याचूर्ण्णम्वा}}{\lemma{\skt{पथ्याचूर्ण्णम्वा}}\Afootnote{H; \skt{पथ्याकृष्णाचूर्णं वा}\ A}}\skt{ }\edtext{\skt{स्नुहा\-क्षी\-र\-प\-रि\-भा\-वि\-त\-मु\-त्का\-रि\-काम्वा}}{\lemma{\skt{स्नुहा\appbrk{}क्षी\appbrk{}र\appbrk{}प\appbrk{}रि\appbrk{}भा\appbrk{}वि\appbrk{}त\appbrk{}मु\appbrk{}त्का\appbrk{}रि\appbrk{}काम्वा}}\Afootnote{H; \skt{स्नुही\appbrk{}क्षी\appbrk{}र\appbrk{}भा\appbrk{}वि\appbrk{}त\appbrk{}मु\appbrk{}त्का\appbrk{}रिकां}\ A}}\skt{ }\edtext{\skt{पक्त्वा}}{\lemma{\skt{पक्त्वा}}\Afootnote{H; \skt{पक्वां}\ A}}\skt{ दापयेत्~। }\edtext{\edtext{\skt{हरी\-त\-की\-चू\-र्ण्ण\-प्र\-स्थ\-मा\-ढके}}{\lemma{\skt{हरी\appbrk{}त\appbrk{}की\appbrk{}चू\appbrk{}र्ण्ण\appbrk{}प्र\appbrk{}स्थम्}}\Afootnote{H; \skt{हरीतकीचूर्णं प्रस्थ}\ A}}}{\lemma{\skt{आढके}}\Afootnote{H; \skt{माढके}\ A}}\skt{ }\edtext{\skt{घृतस्यावाप्य}}{\lemma{\skt{घृतस्यावाप्य}}\Afootnote{H; \skt{घृत\appbrk{}स्या\appbrk{}वा\appbrk{}प्या\appbrk{}ङ्गा\appbrk{}रे\appbrk{}ष्व\appbrk{}भि\appbrk{}वि\appbrk{}लाप्य}\ A}}\skt{ }\edtext{\skt{खजे\-ना\-भि\-म\-थ्या\-नु\-गु\-प्त\-न्नि\-धा\-या\-र्द्ध\-मासं}}{\lemma{\skt{निधायार्द्धमासं}}\Afootnote{H; \skt{कृत्वाऽर्धमासं}\ A}}\skt{ यवपल्ले वासयेत्~। ततश्चोद्धृत्य }\edtext{\skt{परिश्राव्य}}{\lemma{\skt{परिश्राव्य}}\Afootnote{H; \skt{परिस्राव्य}\ A}}\skt{ }\edtext{\skt{मह\-ती\-ह\-री\-त\-की\-क्वा\-था\-म्ल\-द\-धी\-न्या\-वाप्य}}{\lemma{\skt{मह\appbrk{}ती\appbrk{}ह\appbrk{}री\appbrk{}त\appbrk{}की\appbrk{}क्वा\appbrk{}था\appbrk{}म्ल\appbrk{}द\appbrk{}धी\appbrk{}न्या\appbrk{}वाप्य}}\Afootnote{H; \skt{हरी\appbrk{}त\appbrk{}की\appbrk{}क्वा\appbrk{}था\appbrk{}म्ल\appbrk{}द\appbrk{}धी\appbrk{}न्या\appbrk{}वाप्य}\ A}}\skt{ }\edtext{\skt{विप\-चे\-द्य\-था\-यो\-ग\-म्मा\-स\-म\-र्द्ध\-मा\-सम्वा}}{\lemma{\skt{विपचेद्}}\Afootnote{H; \skt{विपचेत्, तद्}\ A}}\skt{ पाययेत्~। गव्ये पयसि }\edtext{\edtext{\skt{महा\-वृ\-क्ष\-क्षी\-र\-मा\-वाप्य}}{\lemma{\skt{महा\appbrk{}वृ\appbrk{}क्ष\appbrk{}क्षी\appbrk{}रम्}}\Afootnote{H; \skt{महावृक्षक्षीर}\ A}}}{\lemma{\skt{आवाप्य}}\Afootnote{H; \skt{मावाप्य}\ A}}\skt{ विपचेत्~॥ विप\-क्व\-ञ्चा\-व\-तार्य शीतीभूतं मन्था\-ने\-ना\-भि\-मथ्य }\edtext{\edtext{\skt{नव\-नी\-त\-मा\-दाय}}{\lemma{\skt{नवनीतम्}}\Afootnote{H; \skt{नवनीत}\ A}}}{\lemma{\skt{आदाय}}\Afootnote{H; \skt{मादाय}\ A}}\skt{ भूयो }\edtext{\skt{महा\-वृ\-क्ष\-क्षी\-र\-क\-षा\-ये\-णैव}}{\lemma{\skt{महा\appbrk{}वृ\appbrk{}क्ष\appbrk{}क्षी\appbrk{}र\appbrk{}क\appbrk{}षा\appbrk{}ये\appbrk{}णैव}}\Afootnote{H; \skt{महा\appbrk{}वृ\appbrk{}क्ष\appbrk{}क्षी\appbrk{}रे\appbrk{}णैव}\ A}}\skt{ }\edtext{\skt{सह}}{\lemma{\skt{सह}}\Afootnote{H; \emph{om.}\ A}}\skt{ विपचेत्~। }\edtext{\skt{तद्य\-था\-यो\-ग\-म्मा\-स\-म\-र्द्ध\-मा\-सम्वा}}{\lemma{\skt{अर्द्धमासम्वा}}\Afootnote{H; \skt{मासार्धं वा}\ A}}\skt{ पाययेत्~। }\edtext{\skt{चव्य\-चि\-त्र\-क\-द\-न्त्य\-ति\-वि\-षा\-कु\-ष्ठ\-शा\-रि\-वा\-वि\-डं\-ग\-त्रि\-फ\-ला\-ज\-मो\-द\-ह\-रि\-द्रा\-दा\-रु\-ह\-रि\-द्रा\-शं\-खि\-नी\-तृ\-वृ\-त्क\-टु\-का\-ना\-म\-र्द्ध\-का\-र्षिका}}{\lemma{\skt{चव्य\appbrk{}चि\appbrk{}त्र\appbrk{}क\appbrk{}द\appbrk{}न्त्य\appbrk{}ति\appbrk{}वि\appbrk{}षा\appbrk{}कु\appbrk{}ष्ठ\appbrk{}शा\appbrk{}रि\appbrk{}वा\appbrk{}वि\appbrk{}डं\appbrk{}ग\appbrk{}त्रि\appbrk{}फ\appbrk{}ला\appbrk{}ज\appbrk{}मो\appbrk{}द\appbrk{}ह\appbrk{}रि\appbrk{}द्रा\appbrk{}दा\appbrk{}रु\appbrk{}ह\appbrk{}रि\appbrk{}द्रा\appbrk{}शं\appbrk{}खि\appbrk{}नी\appbrk{}तृ\appbrk{}वृ\appbrk{}त्क\appbrk{}टु\appbrk{}का\appbrk{}नाम्}}\Afootnote{H; \skt{चव्य\appbrk{}चि\appbrk{}त्र\appbrk{}क\appbrk{}द\appbrk{}न्त्य\appbrk{}ति\appbrk{}वि\appbrk{}षा\appbrk{}कु\appbrk{}ष्ठ\appbrk{}सा\appbrk{}रि\appbrk{}वा\appbrk{}त्रि\appbrk{}फ\appbrk{}ला\appbrk{}ज\appbrk{}मो\appbrk{}द\appbrk{}ह\appbrk{}रि\appbrk{}द्रा\appbrk{}श\appbrk{}ङ्खि\appbrk{}नी\appbrk{}त्रि\appbrk{}वृ\appbrk{}त्त्रि\appbrk{}क\appbrk{}टु\appbrk{}का\appbrk{}नाम्}\ A}}\skt{ भागाः~। }\edtext{\edtext{\skt{राज\-वृ\-क्ष\-फ\-ल\-म\-ज्जोऽ\-अष्टौ}}{\lemma{\skt{राज\appbrk{}वृ\appbrk{}क्ष\appbrk{}फ\appbrk{}ल\appbrk{}मज्जो}}\Afootnote{H; \skt{राज\appbrk{}वृ\appbrk{}क्ष\appbrk{}फ\appbrk{}ल\appbrk{}मज्ज्}\ A}}}{\lemma{\skt{ऽअष्टौ}}\Afootnote{H; \skt{ञामष्टौ}\ A}}\skt{ कर्षाः~। महा\-वृ\-क्ष\-क्षी\-र\-पले द्वे~। गवां क्षीर\-मू\-त्र\-यो\-र\-ष्टा\-वष्टौ पलानि~। }\edtext{\skt{तत्सर्व्वं}}{\lemma{\skt{तत्}}\Afootnote{H; \skt{एतत्}\ A}}\skt{ घृतप्रस्थे समावाप्य विपचेत्~। }\edtext{\edtext{\skt{तथा\-यो\-ग\-म्मा\-स\-म\-र्द्ध\-मा\-सम्वा}}{\lemma{\skt{तथायोगम्}}\Afootnote{H; \skt{तद्यथा योगं}\ A}}}{\lemma{\skt{अर्द्धमासम्वा}}\Afootnote{H; \skt{अर्धमास्ं वा}\ A}}\skt{ पाययेत्~। एतानि }\edtext{\skt{खलु}}{\lemma{\skt{खलु}}\Afootnote{H; \emph{om.}\ A}}\skt{ तिल्व\-क\-घृ\-त\-च\-तु\-र्थानि }\edtext{\skt{सर्प्पीं\-ष्यु\-द\-र\-गु\-ल्म\-वि\-द्र\-ध्य\-ष्ठी\-ला\-ना\-ह\-कु\-ष्ठो\-न्मा\-दा\-प\-स्मा\-रे\-षू\-प\-यो\-ज्यानि}}{\lemma{\skt{सर्प्पींष्य्}}\Afootnote{H; \skt{सर्पीष्य्}\ A}}\skt{~। विरे\-च\-ना\-र्थ\-म्मू\-त्रा\-स\-वा\-रि\-ष्ट\-सु\-रा\-श्चा\-भी\-क्ष्ण\-म्म\-हा\-वृ\-क्ष\-क्षी\-र\-स\-म्भृताः सेवेत~। विरे\-च\-न\-द्र\-व्य\-क\-षा\-यम्वा शृङ्ग\-वे\-र\-दे\-व\-दा\-रु\-प्र\-गाढं~।}\pend

\pstart \ledsidenote{\passref{4.14.11}}\edtext{\edtext{\skt{वम\-न\-वि\-रे\-च\-न\-द्र\-व्या\-णा\-म्प\-लि\-का\-भागाः}}{\lemma{\skt{वम\appbrk{}न\appbrk{}वि\appbrk{}रे\appbrk{}च\appbrk{}न\appbrk{}द्र\appbrk{}व्या\appbrk{}णाम्}}\Afootnote{H; \skt{वम\appbrk{}न\appbrk{}वि\appbrk{}रे\appbrk{}च\appbrk{}न\appbrk{}शि\appbrk{}रो\appbrk{}वि\appbrk{}रे\appbrk{}च\appbrk{}न\appbrk{}द्र\appbrk{}व्याणां}\ A}}}{\lemma{\skt{पलिकाभागाः}}\Afootnote{H; \skt{पालिकाभागाः}\ A}}\skt{~। पिप्प\-ल्या\-दि\-व\-चा\-दि\-ह\-रि\-द्रा\-दि\-प\-रि\-प\-ठि\-ता\-नाञ्च द्रव्याणां श्लक्ष्णपिष्टानां यथोक्तानाञ्च लव\-णा\-ना\-न्त\-त्स\-र्व्व\-म्मू\-त्र\-गणे प्रक्षिप्य महा\-वृ\-क्ष\-क्षी\-र\-प्र\-स्थञ्च }\edtext{\edtext{\skt{मृद्व\-ग्नि\-ना\-व\-घ\-ट्ट\-य\-न्वि\-प\-चेत्}}{\lemma{\skt{मृद्व\appbrk{}ग्नि\appbrk{}ना\appbrk{}व\appbrk{}घ\appbrk{}ट्ट\appbrk{}यन्}}\Afootnote{H; \skt{मृद्व\appbrk{}ग्नि\appbrk{}नाऽ\appbrk{}व\appbrk{}घ\appbrk{}ट्ट\appbrk{}यन्}\ A}}\skt{~। अवि\-द\-ग्ध\-क\-ल्क\-न्त\-त्सा\-धु\-सि\-द्ध\-म\-व\-तार्य}}{\lemma{\skt{विपचेत्~। अवि\appbrk{}द\appbrk{}ग्ध\appbrk{}क\appbrk{}ल्कन्}}\Afootnote{H; \skt{विप\appbrk{}चे\appbrk{}द\appbrk{}प्र\appbrk{}द\appbrk{}ग्ध\appbrk{}कल्कं}\ A}}\skt{ }\edtext{\edtext{\edtext{\skt{शीती\-भू\-त\-म\-ङ्गु\-ष्ठ\-मा\-त्रा\-ङ्गु\-टि\-का\-म्व\-र्त्त\-येत्}}{\lemma{\skt{शीतीभूतम्}}\Afootnote{H; \skt{शीतीभूत}\ A}}}{\lemma{\skt{अङ्गुष्ठमात्राङ्}}\Afootnote{H; \skt{मक्षमात्रा}\ A}}}{\lemma{\skt{गुटिकाम्}}\Afootnote{H; \skt{गुटिका}\ A}}\skt{~। तासामेकां द्वे तिस्रो वा }\edtext{\skt{यथा\-यो\-ग\-म्मा\-सा\-स्त्रीं\-श्च\-तुरो}}{\lemma{\skt{यथायोगम्}}\Afootnote{H; \skt{गुटिका बलापेक्षया}\ A}}\skt{ वा सेवेत~। एषा }\edtext{\skt{खलु}}{\lemma{\skt{खलु}}\Afootnote{H; \skt{*आ}\ A}}\skt{ }\edtext{\skt{मूत्र\-व\-र्त्ति\-क्रि\-या\-वि\-शे\-षेण}}{\lemma{\skt{मूत्र\appbrk{}व\appbrk{}र्त्ति\appbrk{}क्रि\appbrk{}या\appbrk{}वि\appbrk{}शे\appbrk{}षेण}}\Afootnote{H; \skt{नाहवर्तिक्रिया विशेषेण}\ A}}\skt{ महा\-व्या\-धि\-षू\-प\-यु\-ज्यते~। }\edtext{\skt{कोष्ठजाश्च}}{\lemma{\skt{कोष्ठजाश्च}}\Afootnote{H; \skt{(विशेषेण) कोष्ठजांश्च}\ A}}\skt{ }\edtext{\skt{क्रिमीनुपहन्ति}}{\lemma{\skt{क्रिमीनुपहन्ति}}\Afootnote{H; \skt{कृमीनपहन्ति}\ A}}\skt{ }\edtext{\skt{कास\-श्वा\-स\-प्र\-ति\-श्या\-या\-रो\-च\-का\-वि\-पा\-को\-दा\-व\-र्त्ताश्च}}{\lemma{\skt{कास\appbrk{}श्वा\appbrk{}स\appbrk{}प्र\appbrk{}ति\appbrk{}श्या\appbrk{}या\appbrk{}रो\appbrk{}च\appbrk{}का\appbrk{}वि\appbrk{}पा\appbrk{}को\appbrk{}दा\appbrk{}व\appbrk{}र्त्ताश्च}}\Afootnote{H; \skt{कास\appbrk{}श्वा\appbrk{}स\appbrk{}कृ\appbrk{}मि\appbrk{}कु\appbrk{}ष्ठ\appbrk{}प्र\appbrk{}ति\appbrk{}श्या\appbrk{}या\appbrk{}रो\appbrk{}च\appbrk{}का\appbrk{}वि\appbrk{}पा\appbrk{}को\appbrk{}दा\appbrk{}व\appbrk{}र्तांश्च नाशयति}\ A}}\skt{~।}\pend

\pstart \ledsidenote{\passref{4.14.12}}\edtext{\skt{मद\-न\-फ\-ल\-म\-ज्जा\-जी\-मू\-त\-के\-क्ष्वा\-कु\-त्रि\-क\-टु\-क\-स\-र्ष\-प\-ल\-व\-णानि}}{\lemma{\skt{मद\appbrk{}न\appbrk{}फ\appbrk{}ल\appbrk{}म\appbrk{}ज्जा\appbrk{}जी\appbrk{}मू\appbrk{}त\appbrk{}के\appbrk{}क्ष्वा\appbrk{}कु\appbrk{}त्रि\appbrk{}क\appbrk{}टु\appbrk{}क\appbrk{}स\appbrk{}र्ष\appbrk{}प\appbrk{}ल\appbrk{}व\appbrk{}णानि}}\Afootnote{H; \skt{मद\appbrk{}न\appbrk{}फ\appbrk{}ल\appbrk{}म\appbrk{}ज्ज\appbrk{}कु\appbrk{}ट\appbrk{}ज\appbrk{}जी\appbrk{}मू\appbrk{}त\appbrk{}के\appbrk{}क्ष्वा\appbrk{}कु\appbrk{}धा\appbrk{}मा\appbrk{}र्ग\appbrk{}व\appbrk{}त्रि\appbrk{}वृ\appbrk{}त्त्रि\appbrk{}क\appbrk{}टु\appbrk{}क\appbrk{}स\appbrk{}र्ष\appbrk{}प\appbrk{}ल\appbrk{}व\appbrk{}णानि}\ A}}\skt{ }\edtext{\skt{मूत्र\-क्षा\-र\-यो\-र\-न्य\-त\-रेण}}{\lemma{\skt{मूत्रक्षारयोर्}}\Afootnote{H; \skt{महा\appbrk{}वृ\appbrk{}क्ष\appbrk{}क्षी\appbrk{}र\appbrk{}मू\appbrk{}त्र\appbrk{}योर्}\ A}}\skt{ }\edtext{\skt{पिष्ट्वा\-ङ्गु\-ष्ठ\-मा\-त्रा\-वर्त्ती}}{\lemma{\skt{पिष्ट्वा\appbrk{}ङ्गु\appbrk{}ष्ठ\appbrk{}मा\appbrk{}त्रा\appbrk{}वर्त्ती}}\Afootnote{H; \skt{पिष्ट्वाऽङ्गुष्ठमात्रां वर्तिं}\ A}}\skt{ कृत्वोदरिणः~। }\edtext{\skt{आनाहे तु}}{\lemma{\skt{आनाहे तु}}\Afootnote{H; \skt{आनाहे}\ A}}\skt{ }\edtext{\skt{तैल\-ल\-व\-ण\-गु\-दा\-भ्य\-क्त\-स्यैकां}}{\lemma{\skt{तैल\appbrk{}ल\appbrk{}व\appbrk{}ण\appbrk{}गु\appbrk{}दा\appbrk{}भ्य\appbrk{}क्त\appbrk{}स्यैकां}}\Afootnote{H; \skt{तैल\appbrk{}ल\appbrk{}व\appbrk{}णा\appbrk{}भ्य\appbrk{}क्त\appbrk{}गु\appbrk{}द\appbrk{}स्यैकां}\ A}}\skt{ द्वे }\edtext{\skt{च}}{\lemma{\skt{च}}\Afootnote{H; \skt{तिस्रो वा}\ A}}\skt{ पायौ }\edtext{\skt{निदध्यात्~। एषा\-ना\-ह\-व\-र्त्ति\-क्रिया}}{\lemma{\skt{निदध्यात्~। एषा\appbrk{}ना\appbrk{}ह\appbrk{}व\appbrk{}र्त्ति\appbrk{}क्रिया}}\Afootnote{H; \skt{निदध्यात्, एषाऽऽ\appbrk{}आ\appbrk{}ना\appbrk{}ह\appbrk{}व\appbrk{}र्ति\appbrk{}क्रिया}\ A}}\skt{ वात\-मू\-त्र\-पु\-री\-षो\-दा\-व\-र्त्ता\-ध्मा\-ना\-ना\-हेषु विधेया~। }\edtext{\skt{दूष्योदरिणन्तु}}{\lemma{\skt{दूष्योदरिणन्तु}}\Afootnote{H; \emph{om.}\ A}}\skt{ }\edtext{\edtext{\skt{व्योषा\-श्व\-मा\-र\-क\-गु\-ञ्जा\-का\-का\-द\-नी\-मू\-ल\-क\-ल्क\-म्पा\-य\-येत्}}{\lemma{\skt{व्योषा\appbrk{}श्व\appbrk{}मा\appbrk{}र\appbrk{}क\appbrk{}गु\appbrk{}ञ्जा\appbrk{}का\appbrk{}का\appbrk{}द\appbrk{}नी\appbrk{}मू\appbrk{}ल\appbrk{}क\appbrk{}ल्कम्}}\Afootnote{H; \emph{om.}\ A}}\skt{~। इक्षुकाण्डानि वा}}{\lemma{\skt{पाययेत्~। इक्षुकाण्डानि वा}}\Afootnote{H; \emph{om.}\ A}}\skt{ }\edtext{\skt{कृष्णसर्पेण}}{\lemma{\skt{कृष्णसर्पेण}}\Afootnote{H; \emph{om.}\ A}}\skt{ }\edtext{\skt{दंशयित्वा}}{\lemma{\skt{दंशयित्वा}}\Afootnote{H; \emph{om.}\ A}}\skt{ }\edtext{\skt{भक्षयेत्}}{\lemma{\skt{भक्षयेत्}}\Afootnote{H; \emph{om.}\ A}}\skt{~। }\edtext{\skt{वल्लीफलानि वा}}{\lemma{\skt{वल्लीफलानि वा}}\Afootnote{H; \emph{om.}\ A}}\skt{ }\edtext{\skt{कन्दजम्वा}}{\lemma{\skt{कन्दजम्वा}}\Afootnote{H; \emph{om.}\ A}}\skt{ }\edtext{\skt{मूलजम्वा}}{\lemma{\skt{मूलजम्वा}}\Afootnote{H; \emph{om.}\ A}}\skt{ }\edtext{\edtext{\skt{विषमासेवेत}}{\lemma{\skt{विषम्}}\Afootnote{H; \emph{om.}\ A}}}{\lemma{\skt{आसेवेत}}\Afootnote{H; \emph{om.}\ A}}\skt{~। }\edtext{\edtext{\skt{तेना\-ग\-दी\-भ\-व\-त्य\-न्य\-त\-रम्वा}}{\lemma{\skt{तेना\appbrk{}ग\appbrk{}दी\appbrk{}भ\appbrk{}वत्य्}}\Afootnote{H; \emph{om.}\ A}}}{\lemma{\skt{अन्यतरम्वा}}\Afootnote{H; \emph{om.}\ A}}\skt{ }\edtext{\edtext{\skt{भावमापद्यते}}{\lemma{\skt{भावम्}}\Afootnote{H; \emph{om.}\ A}}}{\lemma{\skt{आपद्यते}}\Afootnote{H; \emph{om.}\ A}}\skt{~।}\pend

\pstart \ledsidenote{\passref{4.14.13}}\edtext{\skt{प्लीहो\-द\-रि\-ण\-स्तू\-प\-स्नि\-ग्ध\-स्वि\-न्नस्य}}{\lemma{\skt{तूप\appbrk{}स्नि\appbrk{}ग्ध\appbrk{}स्वि\appbrk{}न्नस्य}}\Afootnote{H; \skt{स्निग्धस्विन्नस्य}\ A}}\skt{ दध्ना भुक्तवतो }\edtext{\edtext{\skt{वाम\-बा\-हु\-कु\-र्प्प\-र\-स\-न्धा\-व\-भ्य\-न्त\-र\-सि\-रा\-म्वि\-ध्येत्}}{\lemma{\skt{वाम\appbrk{}बा\appbrk{}हु\appbrk{}कु\appbrk{}र्प्प\appbrk{}र\appbrk{}स\appbrk{}न्धाव्}}\Afootnote{H; \skt{वामबाहौ कूर्परा}\ A}}}{\lemma{\skt{अभ्य\appbrk{}न्त\appbrk{}र\appbrk{}सि\appbrk{}राम्}}\Afootnote{H; \skt{भ्यन्तरतः सिरां}\ A}}\skt{~। विमर्द्दयेच्च पाणिना प्लीहानं~। }\edtext{\skt{रुधि\-र\-प्र\-स्य\-न्द\-ना\-र्थ\-न्ततः}}{\lemma{\skt{रुधि\appbrk{}र\appbrk{}प्र\appbrk{}स्य\appbrk{}न्द\appbrk{}ना\appbrk{}र्थन्}}\Afootnote{H; \skt{रुधि\appbrk{}र\appbrk{}स्य\appbrk{}न्द\appbrk{}नार्थं}\ A}}\skt{ }\edtext{\skt{शुद्धदेहं}}{\lemma{\skt{शुद्धदेहं}}\Afootnote{H; \skt{संशुद्धदेहं}\ A}}\skt{ }\edtext{\skt{समु\-द्र\-शु\-क्ति\-क्षा\-र\-म्प\-यसा}}{\lemma{\skt{समु\appbrk{}द्र\appbrk{}शु\appbrk{}क्ति\appbrk{}क्षा\appbrk{}रम्}}\Afootnote{H; \skt{समु\appbrk{}द्र\appbrk{}शु\appbrk{}क्ति\appbrk{}का\appbrk{}क्षारं}\ A}}\skt{ पाययेत~। }\edtext{\skt{हिङ्गु\-सौ\-व\-र्च्चि\-काम्वा}}{\lemma{\skt{हिङ्गु\appbrk{}सौ\appbrk{}व\appbrk{}र्च्चि\appbrk{}काम्वा}}\Afootnote{H; \skt{हिङ्गुसौवर्चिके वा}\ A}}\skt{ क्षीरेण~। }\edtext{\skt{श्रुतेन वा}}{\lemma{\skt{श्रुतेन वा}}\Afootnote{H; \skt{स्रुतेन}\ A}}\skt{ }\edtext{\skt{पलाशक्षारेण}}{\lemma{\skt{पलाशक्षारेण}}\Afootnote{H; \skt{पलाशक्षारेण वा}\ A}}\skt{~। }\edtext{\skt{यव\-क्षा\-र\-पा\-रि\-जा\-त\-के\-क्षु\-र\-का\-पा\-मा\-र्ग्ग\-क्षा\-रम्वा}}{\lemma{\skt{यव\appbrk{}क्षा\appbrk{}र\appbrk{}पा\appbrk{}रि\appbrk{}जा\appbrk{}त\appbrk{}के\appbrk{}क्षु\appbrk{}र\appbrk{}का\appbrk{}पा\appbrk{}मा\appbrk{}र्ग्ग\appbrk{}क्षा\appbrk{}रम्वा}}\Afootnote{H; \skt{यवक्षारं, किंशु\appbrk{}क\appbrk{}क्षा\appbrk{}रो\appbrk{}द\appbrk{}केन वा बहुशः स्रुतेन यवक्षारं, पारि\appbrk{}जा\appbrk{}त\appbrk{}के\appbrk{}क्षु\appbrk{}र\appbrk{}का\appbrk{}पा\appbrk{}मा\appbrk{}र्ग\appbrk{}क्षारं वा}\ A}}\skt{ तैलसंसृष्टं~। }\edtext{\skt{सौभा\-ञ्ज\-न\-क\-यूषं वा}}{\lemma{\skt{सौभा\appbrk{}ञ्ज\appbrk{}न\appbrk{}क\appbrk{}यूषं वा}}\Afootnote{H; \skt{शोभा\appbrk{}ञ्ज\appbrk{}न\appbrk{}क\appbrk{}यूषं}\ A}}\skt{ }\edtext{\skt{पिप्प\-ली\-सै\-न्ध\-व\-चि\-त्र\-क\-यु\-क्त\-म\-म्ल\-पू\-ति\-क\-र\-ञ्ज\-क्षा\-रम्वा}}{\lemma{\skt{अम्ल\appbrk{}पू\appbrk{}ति\appbrk{}क\appbrk{}र\appbrk{}ञ्ज\appbrk{}क्षा\appbrk{}रम्वा}}\Afootnote{H; \skt{पूति\appbrk{}क\appbrk{}र\appbrk{}ञ्ज\appbrk{}क्षारं वा}\ A}}\skt{~। }\edtext{\skt{विड\-ङ्ग\-ल\-व\-ण\-पि\-प्प\-ली\-प्रा\-गाढं}}{\lemma{\skt{विड\appbrk{}ङ्ग\appbrk{}ल\appbrk{}व\appbrk{}ण\appbrk{}पि\appbrk{}प्प\appbrk{}ली\appbrk{}प्रा\appbrk{}गाढं}}\Afootnote{H; \skt{ऽम्लस्रुतं विड्ल\appbrk{}व\appbrk{}ण\appbrk{}पि\appbrk{}प्प\appbrk{}ली\appbrk{}प्रा\appbrk{}गा\appbrk{}ढम्}\ A}}\skt{~।}\pend

\pstart \ledsidenote{\passref{4.14.14}}\edtext{\skt{पिप्प\-ली\-पि\-प्प\-ली\-मूलं}}{\lemma{\skt{पिप्प\appbrk{}ली\appbrk{}पि\appbrk{}प्प\appbrk{}ली\appbrk{}मूलं}}\Afootnote{H; \skt{पिप्प\appbrk{}ली\appbrk{}पि\appbrk{}प्प\appbrk{}ली\appbrk{}मूल}\ A}}\skt{ }\edtext{\edtext{\skt{शृङ्ग\-वे\-र\-चि\-त्र\-क\-य\-व\-क्षा\-र\-सै\-न्ध\-वा\-ना\-म्प\-लिका}}{\lemma{\skt{शृङ्ग\appbrk{}वे\appbrk{}र\appbrk{}चि\appbrk{}त्र\appbrk{}क\appbrk{}य\appbrk{}व\appbrk{}क्षा\appbrk{}र\appbrk{}सै\appbrk{}न्ध\appbrk{}वा\appbrk{}नाम्}}\Afootnote{H; \skt{चव्य\appbrk{}चि\appbrk{}त्र\appbrk{}क\appbrk{}शृ\appbrk{}ङ्ग\appbrk{}वे\appbrk{}र\appbrk{}य\appbrk{}व\appbrk{}क्षा\appbrk{}र\appbrk{}सै\appbrk{}न्ध\appbrk{}वानां}\ A}}}{\lemma{\skt{पलिका}}\Afootnote{H; \skt{पालिका}\ A}}\skt{ भागाः~। घृतप्रस्थं }\edtext{\skt{तत्तुल्यं}}{\lemma{\skt{तत्तुल्यं}}\Afootnote{H; \skt{तत्तुल्यं च}\ A}}\skt{ }\edtext{\edtext{\skt{क्षीर\-न्त\-दे\-क\-ध्य\-म्वि\-पा\-च\-ये\-दे\-त\-त्खलु}}{\lemma{\skt{तदेकध्यम्}}\Afootnote{H; \skt{तदैकध्यं}\ A}}}{\lemma{\skt{खलु}}\Afootnote{H; \emph{om.}\ A}}\skt{ षट्पलकं }\edtext{\edtext{\skt{प्लीहा\-ग्नि\-सा\-द\-गु\-ल्मो\-दा\-व\-र्त्त\-श्व\-य\-थु\-पा\-ण्डु\-रो\-ग\-श्वा\-स\-का\-स\-प्र\-ति\-श्या\-यो\-र्द्ध्व\-वा\-त\-वि\-ष\-म\-ज्व\-रा\-नु\-प\-हन्ति}}{\lemma{\skt{प्लीहा\appbrk{}ग्नि\appbrk{}सा\appbrk{}द\appbrk{}गु\appbrk{}ल्मो\appbrk{}दा\appbrk{}व\appbrk{}र्त्त\appbrk{}श्व\appbrk{}य\appbrk{}थु\appbrk{}पा\appbrk{}ण्डु\appbrk{}रो\appbrk{}ग\appbrk{}श्वा\appbrk{}स\appbrk{}का\appbrk{}स\appbrk{}प्र\appbrk{}ति\appbrk{}श्या\appbrk{}योर्}}\Afootnote{H; \skt{नाम सर्पिः प्लीहा\appbrk{}ग्नि\appbrk{}स\appbrk{}ङ्ग\appbrk{}गु\appbrk{}ल्मो\appbrk{}द\appbrk{}रो\appbrk{}दा\appbrk{}व\appbrk{}र्त\appbrk{}श्व\appbrk{}य\appbrk{}थु\appbrk{}पा\appbrk{}ण्डु\appbrk{}रो\appbrk{}ग\appbrk{}का\appbrk{}स\appbrk{}श्वा\appbrk{}स\appbrk{}प्र\appbrk{}ति\appbrk{}श्या\appbrk{}योर्}\ A}}}{\lemma{\skt{द्ध्ववा\appbrk{}त\appbrk{}वि\appbrk{}ष\appbrk{}म\appbrk{}ज्व\appbrk{}रा\appbrk{}नु\appbrk{}प\appbrk{}हन्ति}}\Afootnote{H; \skt{ध्ववा\appbrk{}त\appbrk{}वि\appbrk{}ष\appbrk{}म\appbrk{}ज्व\appbrk{}रा\appbrk{}न\appbrk{}प\appbrk{}हन्ति}\ A}}\skt{~॥ }\edtext{\skt{पिप्प्ल\-ली\-म\-रि\-च\-शृ\-ङ्ग\-वे\-र\-हि\-ङ्गु\-सै\-न्ध\-व\-सौ\-व\-र्च्च\-ल\-वि\-ड\-दा\-डि\-म\-य\-वा\-नि\-का\-पु\-ष्क\-र\-मू\-ला\-जा\-जी\-धा\-न्या\-म्ल\-वे\-त\-स\-क्षा\-र\-द्व\-य\-व\-चा\-चि\-त्त्र\-कै\-ला\-ज\-ग\-न्धा\-तु\-ल\-सी\-त्येतैः}}{\lemma{\skt{पिप्प्ल\appbrk{}ली\appbrk{}म\appbrk{}रि\appbrk{}च\appbrk{}शृ\appbrk{}ङ्ग\appbrk{}वे\appbrk{}र\appbrk{}हि\appbrk{}ङ्गु\appbrk{}सै\appbrk{}न्ध\appbrk{}व\appbrk{}सौ\appbrk{}व\appbrk{}र्च्च\appbrk{}ल\appbrk{}वि\appbrk{}ड\appbrk{}दा\appbrk{}डि\appbrk{}म\appbrk{}य\appbrk{}वा\appbrk{}नि\appbrk{}का\appbrk{}पु\appbrk{}ष्क\appbrk{}र\appbrk{}मू\appbrk{}ला\appbrk{}जा\appbrk{}जी\appbrk{}धा\appbrk{}न्या\appbrk{}म्ल\appbrk{}वे\appbrk{}त\appbrk{}स\appbrk{}क्षा\appbrk{}र\appbrk{}द्व\appbrk{}य\appbrk{}व\appbrk{}चा\appbrk{}चि\appbrk{}त्त्र\appbrk{}कै\appbrk{}ला\appbrk{}ज\appbrk{}ग\appbrk{}न्धा\appbrk{}तु\appbrk{}ल\appbrk{}सी\appbrk{}त्येतैः}}\Afootnote{H; \emph{om.}\ A}}\skt{ }\edtext{\edtext{\edtext{\skt{समै\-र्घृ\-त\-प्र\-स्थ\-न्दधि}}{\lemma{\skt{समैर्}}\Afootnote{H; \emph{om.}\ A}}}{\lemma{\skt{घृतप्रस्थन्}}\Afootnote{H; \emph{om.}\ A}}}{\lemma{\skt{दधि}}\Afootnote{H; \emph{om.}\ A}}\skt{ }\edtext{\skt{चतुर्ग्गुणं}}{\lemma{\skt{चतुर्ग्गुणं}}\Afootnote{H; \emph{om.}\ A}}\skt{ }\edtext{\edtext{\edtext{\skt{बीज\-पू\-र\-क\-र\-स\-यु\-क्त\-म्पा\-च\-ये\-दे\-त\-द्दा\-धिकं}}{\lemma{\skt{बीज\appbrk{}पू\appbrk{}र\appbrk{}क\appbrk{}र\appbrk{}स\appbrk{}यु\appbrk{}क्तम्}}\Afootnote{H; \emph{om.}\ A}}}{\lemma{\skt{पाचयेदेतद्}}\Afootnote{H; \skt{मन्}\ A}}}{\lemma{\skt{दाधिकं}}\Afootnote{H; \skt{दाग्निर्वा हिङ्ग्वादिकं}\ A}}\skt{ }\edtext{\edtext{\edtext{\skt{सर्प्पि\-र्द्दी\-प\-नी\-य\-मु\-द\-र\-गु\-ल्म\-का\-स\-श्वा\-स\-ज्व\-र\-वि\-द्र\-द्य\-ष्ठी\-ला\-ना\-ह\-स\-र्व्व\-वा\-त\-रु\-क्छ\-मनं}}{\lemma{\skt{सर्प्पिर्}}\Afootnote{H; \skt{चूर्}\ A}}}{\lemma{\skt{द्दीपनीयम्}}\Afootnote{H; \skt{णम्}\ A}}}{\lemma{\skt{उद\appbrk{}र\appbrk{}गु\appbrk{}ल्म\appbrk{}का\appbrk{}स\appbrk{}श्वा\appbrk{}स\appbrk{}ज्व\appbrk{}र\appbrk{}वि\appbrk{}द्र\appbrk{}द्य\appbrk{}ष्ठी\appbrk{}ला\appbrk{}ना\appbrk{}ह\appbrk{}स\appbrk{}र्व्व\appbrk{}वा\appbrk{}त\appbrk{}रु\appbrk{}क्छ\appbrk{}मनं}}\Afootnote{H; \skt{उ}\ A}}\skt{ }\edtext{\edtext{\skt{शूलघ्नञ्चेति}}{\lemma{\skt{शूलघ्नञ्}}\Afootnote{H; \emph{om.}\ A}}}{\lemma{\skt{चेति}}\Afootnote{H; \emph{om.}\ A}}\skt{~॥ }\edtext{\skt{शृङ्ग\-वे\-र\-मा\-तु\-लु\-ङ्ग\-र\-स\-दा\-डि\-म\-द\-धि\-शु\-ष्क\-मू\-ल\-को\-ला\-म्ल\-क\-षा\-य\-प्र\-स्थेषु}}{\lemma{\skt{शृङ्ग\appbrk{}वे\appbrk{}र\appbrk{}मा\appbrk{}तु\appbrk{}लु\appbrk{}ङ्ग\appbrk{}र\appbrk{}स\appbrk{}दा\appbrk{}डि\appbrk{}म\appbrk{}द\appbrk{}धि\appbrk{}शु\appbrk{}ष्क\appbrk{}मू\appbrk{}ल\appbrk{}को\appbrk{}ला\appbrk{}म्ल\appbrk{}क\appbrk{}षा\appbrk{}य\appbrk{}प्र\appbrk{}स्थेषु}}\Afootnote{H; \skt{पयुञ्जीत}\ A}}\skt{ }\edtext{\edtext{\skt{घृत\-प्र\-स्थ\-ञ्चा\-वाप्य}}{\lemma{\skt{घृतप्रस्थञ्}}\Afootnote{H; \emph{om.}\ A}}}{\lemma{\skt{चावाप्य}}\Afootnote{H; \emph{om.}\ A}}\skt{~। }\edtext{\skt{पञ्च}}{\lemma{\skt{पञ्च}}\Afootnote{H; \emph{om.}\ A}}\skt{ }\edtext{\edtext{\edtext{\edtext{\skt{कोल\-वि\-ड\-ङ्ग\-य\-वा\-नी\-शु\-ष्क\-मू\-ल\-पा\-ठ\-ग\-र्ब्भ\-न्त\-न्सि\-द्ध\-म्पा\-य\-येत्}}{\lemma{\skt{कोल\appbrk{}वि\appbrk{}ड\appbrk{}ङ्ग\appbrk{}य\appbrk{}वा\appbrk{}नी\appbrk{}शु\appbrk{}ष्क\appbrk{}मू\appbrk{}ल\appbrk{}पा\appbrk{}ठ\appbrk{}ग\appbrk{}र्ब्भन्}}\Afootnote{H\textsuperscript{pc}; \emph{om.}\ A; \skt{कोल\appbrk{}वि\appbrk{}ड\appbrk{}ङ्ग\appbrk{}य\appbrk{}वा\appbrk{}नी\appbrk{}शु\appbrk{}ष्क\appbrk{}मू\appbrk{}ल\appbrk{}पा\appbrk{}ठ\appbrk{}भ\appbrk{}र्ग्गन्}\ H\textsuperscript{ac}}}}{\lemma{\skt{तन्}}\Afootnote{H; \emph{om.}\ A}}}{\lemma{\skt{सिद्धम्}}\Afootnote{H; \emph{om.}\ A}}}{\lemma{\skt{पाययेत्}}\Afootnote{H; \emph{om.}\ A}}\skt{~॥}\pend

\pstart\versepar \ledsidenote{\passref{4.14.15}}\skt{यकृद्दाल्येऽप्येष एव क्रियाविभागः~। }\pend
\pstart\versepar \skt{विशेषतस्तु दक्षिणबाहौ }\edtext{\edtext{\skt{सिराम्विध्येत्}}{\lemma{\skt{सिराम्}}\Afootnote{H; \skt{सिरा}\ A}}}{\lemma{\skt{विध्येत्}}\Afootnote{H; \skt{व्यधः}\ A}}\skt{~॥}\pend

\pstart\versepar \ledsidenote{\passref{4.14.16}}\edtext{\skt{मणिबन्धे}}{\lemma{\skt{मणिबन्धे}}\Afootnote{H; \skt{मणिबन्धं}\ A}}\skt{ }\edtext{\skt{समुन्नाम्य}}{\lemma{\skt{समुन्नाम्य}}\Afootnote{H; \skt{सकृन्नाम्य}\ A}}\skt{ }\edtext{\edtext{\edtext{\skt{वाम\-म\-ङ्गु\-ष्ठ\-मी\-रितं}}{\lemma{\skt{वामम्}}\Afootnote{H; \skt{वामाङ्}\ A}}}{\lemma{\skt{अङ्गुष्ठम्}}\Afootnote{H; \skt{गुष्ठ}\ A}}}{\lemma{\skt{ईरितं}}\Afootnote{H; \skt{समीरिताम्}\ A}}\skt{~। }\pend
\pstart\versepar \skt{दहेत्सिरां शरेणाशु }\edtext{\skt{प्लीहो}}{\lemma{\skt{प्लीहो}}\Afootnote{H; \skt{प्लीह्नो}\ A}}\skt{ वैद्यः प्रशान्तये~॥}\pend

\pstart \ledsidenote{\passref{4.14.17}}\skt{बद्धगुदे }\edtext{\skt{परिस्राविनि च}}{\lemma{\skt{परिस्राविनि च}}\Afootnote{H; \skt{परिस्राविणि च}\ A}}\skt{ }\edtext{\skt{स्निग्धस्विन्नस्याधो}}{\lemma{\skt{स्निग्धस्विन्नस्याधो}}\Afootnote{H; \skt{स्निग्ध\appbrk{}स्वि\appbrk{}न्न\appbrk{}स्या\appbrk{}भ्य\appbrk{}क्त\appbrk{}स्याधो}\ A}}\skt{ नाभे\-र्व्वा\-म\-त\-श्च\-तु\-र\-ङ्गु\-ल\-म\-प\-हाय }\edtext{\skt{रोम\-रा\-ज्या\-दु\-दरं}}{\lemma{\skt{रोम\appbrk{}रा\appbrk{}ज्या\appbrk{}दु\appbrk{}दरं}}\Afootnote{H; \skt{रोमराज्या उदरं}\ A}}\skt{ पाटयित्वा चतु\-र\-ङ्गु\-ल\-प्र\-मा\-ण\-म\-न्त्राणि निष्कृष्य निरीक्ष्य }\edtext{\skt{बद्ध\-गु\-द\-स्या\-न्त्र\-प्र\-ति\-रो\-ध\-क\-र\-म\-स्मानं}}{\lemma{\skt{बद्ध\appbrk{}गु\appbrk{}द\appbrk{}स्या\appbrk{}न्त्र\appbrk{}प्र\appbrk{}ति\appbrk{}रो\appbrk{}ध\appbrk{}क\appbrk{}र\appbrk{}म\appbrk{}स्मानं}}\Afootnote{H; \skt{बद्ध\appbrk{}गु\appbrk{}द\appbrk{}स्या\appbrk{}न्त्र\appbrk{}प्र\appbrk{}ति\appbrk{}रो\appbrk{}ध\appbrk{}क\appbrk{}र\appbrk{}म\appbrk{}श्मानं}\ A}}\skt{ }\edtext{\skt{बालम्वापोह्य}}{\lemma{\skt{वापोह्य}}\Afootnote{H; \skt{वाऽपोह्य}\ A}}\skt{ }\edtext{\skt{मलजातन्ततो}}{\lemma{\skt{मलजातन्}}\Afootnote{H; \skt{मलजातं वा}\ A}}\skt{ मधु }\edtext{\skt{सर्प्पि\-षा\-भ्य\-ज्या\-न्त्राणि}}{\lemma{\skt{सर्प्पि\appbrk{}षा\appbrk{}भ्य\appbrk{}ज्या\appbrk{}न्त्राणि}}\Afootnote{H; \skt{सर्पि\appbrk{}र्भ्या\appbrk{}म\appbrk{}भ्य\appbrk{}ज्या\appbrk{}न्त्राणि}\ A}}\skt{ यथास्थानं स्थापयित्वा बाह्यं व्रणमुदरस्य सीव्येत्~। }\edtext{\skt{परिस्राविणो}}{\lemma{\skt{परिस्राविणो}}\Afootnote{H; \skt{परिस्राविण्य्}\ A}}\skt{ }\edtext{\skt{प्येवमन्त्रं}}{\lemma{\skt{अन्त्रं}}\Afootnote{H; \skt{एव शल्य\appbrk{}मु\appbrk{}द्धृ\appbrk{}त्या\appbrk{}न्त्र\appbrk{}स्रा\appbrk{}वा\appbrk{}न्सं\appbrk{}शोध्य*, तच्छिद्रमान्त्रं}\ A}}\skt{ }\edtext{\skt{समानयित्वा}}{\lemma{\skt{समानयित्वा}}\Afootnote{H; \skt{समाधा}\ A}}\skt{ }\edtext{\skt{च्छिद्रं}}{\lemma{\skt{च्छिद्रं}}\Afootnote{H; \skt{य}\ A}}\skt{ काल\-पि\-पी\-लि\-का\-भि\-र्द्द\-ङ्श\-येत्~। दष्टे च तासां }\edtext{\edtext{\skt{काया\-नु\-प\-ह\-रे\-च्छि\-न्न\-शि\-रांसि}}{\lemma{\skt{काया\appbrk{}नु\appbrk{}प\appbrk{}ह\appbrk{}रेच्}}\Afootnote{H; \skt{काया\appbrk{}न\appbrk{}प\appbrk{}ह\appbrk{}रेन्}\ A}}}{\lemma{\skt{छिन्नशिरांसि}}\Afootnote{H; \skt{न शिरांसि}\ A}}\skt{ ततः }\edtext{\skt{पूर्व्ववत्सीव्येत्}}{\lemma{\skt{सीव्येत्}}\Afootnote{H; \skt{सीव्येत्, संधानं च यथोक्तं कारयेत्, यष्टी\appbrk{}म\appbrk{}धु\appbrk{}क\appbrk{}मि\appbrk{}श्रया च कृष्ण\appbrk{}मृ\appbrk{}दाऽ\appbrk{}व\appbrk{}लिप्य बन्धे\appbrk{}नो\appbrk{}प\appbrk{}च\appbrk{}रेत्, ततो निवातमागारं प्रवे\appbrk{}श्या\appbrk{}चा\appbrk{}रि\appbrk{}क\appbrk{}मु\appbrk{}प\appbrk{}दि\appbrk{}शेत्}\ A}}\skt{~। }\edtext{\skt{वास\-ये\-च्चै\-न\-न्तै\-ल\-द्रो\-ण्या\-म\-ल्पम्वा}}{\lemma{\skt{अल्पम्वा}}\Afootnote{H; \skt{सर्पिर्द्रोण्यां वा}\ A}}\skt{ }\edtext{\skt{पयोवृत्तिरिति}}{\lemma{\skt{पयोवृत्तिर्}}\Afootnote{H; \skt{पयोवृत्तिम्}\ A}}\skt{~॥}\pend

\pstart \ledsidenote{\passref{4.14.18}}\edtext{\skt{उद\-को\-द\-रि\-णस्तु}}{\lemma{\skt{उद\appbrk{}को\appbrk{}द\appbrk{}रि\appbrk{}णस्तु}}\Afootnote{H; \skt{दकोदरिणस्तु}\ A}}\skt{ }\edtext{\skt{तैला\-भ्य\-क्त\-स्यो\-ष्णो\-द\-क\-प\-रि\-स्वि\-न्न\-स्याप्तैः}}{\lemma{\skt{तैला\appbrk{}भ्य\appbrk{}क्त\appbrk{}स्यो\appbrk{}ष्णो\appbrk{}द\appbrk{}क\appbrk{}प\appbrk{}रि\appbrk{}स्वि\appbrk{}न्न\appbrk{}स्याप्तैः}}\Afootnote{H; \skt{वात\appbrk{}ह\appbrk{}र\appbrk{}तै\appbrk{}ला\appbrk{}भ्य\appbrk{}क्त\appbrk{}स्यो\appbrk{}ष्णो\appbrk{}द\appbrk{}क\appbrk{}स्वि\appbrk{}न्नस्य स्थितस्याप्तैः}\ A}}\skt{ }\edtext{\skt{सुप\-रि\-गृ\-ही\-त\-स्या\-क\-क्षा\-त्प\-टू\-वे\-ष्टि\-त\-स्याधो}}{\lemma{\skt{पटू\appbrk{}वे\appbrk{}ष्टि\appbrk{}त\appbrk{}स्याधो}}\Afootnote{\skt{परि\appbrk{}वे\appbrk{}ष्टि\appbrk{}त\appbrk{}स्याधो}\ A; \skt{पटू\appbrk{}वे\appbrk{}ष्टि\appbrk{}त\appbrk{}स्यार्द्धो}\ H}}\skt{ }\edtext{\skt{नाभे\-श्च\-तु\-र\-ङ्गु\-ल\-म\-प\-हाय}}{\lemma{\skt{चतुरङ्गुलम्}}\Afootnote{H; \skt{वाम\appbrk{}त\appbrk{}श्च\appbrk{}तु\appbrk{}र\appbrk{}ङ्गु\appbrk{}लम्}\ A}}\skt{ }\edtext{\edtext{\edtext{\skt{रोम\-रा\-ज्द्व्य\-ङ्गु\-ल\-म्व्री\-हि\-मु\-खे\-ना\-ङ्गु\-ष्ठो\-द\-र\-मा\-त्र\-म्वि\-दार्य}}{\lemma{\skt{रोम\appbrk{}रा\appbrk{}ज्द्व्य\appbrk{}ङ्गु\appbrk{}लम्}}\Afootnote{\skt{रोमराज्या}\ A; \skt{रोमराज्या द्व्यङ्गुलम्}\ H\textsuperscript{ac}; \skt{रोम\appbrk{}रा\appbrk{}मा\appbrk{}ण्य\appbrk{}र्द्व्यं\appbrk{}ङ्गु\appbrk{}लम्}\ H\textsuperscript{pc}}}}{\lemma{\skt{व्रीहि\appbrk{}मु\appbrk{}खे\appbrk{}ना\appbrk{}ङ्गु\appbrk{}ष्ठो\appbrk{}द\appbrk{}र\appbrk{}मा\appbrk{}त्रम्}}\Afootnote{\skt{व्रीहि\appbrk{}मु\appbrk{}खे\appbrk{}ना\appbrk{}ङ्गु\appbrk{}ष्ठो\appbrk{}द\appbrk{}र\appbrk{}प्र\appbrk{}मा\appbrk{}ण\appbrk{}म\appbrk{}व\appbrk{}गाढं}\ A; \skt{व्रीहि\appbrk{}मु\appbrk{}खे\appbrk{}ना\appbrk{}ङ्गु\appbrk{}ष्ठो\appbrk{}द\appbrk{}रमा}\ H}}}{\lemma{\skt{विदार्य}}\Afootnote{\skt{विध्येत्, तत्र}\ A; \skt{त्रम्विदार्य}\ H}}\skt{ }\edtext{\skt{त्रप्वा\-दी\-ना\-म\-न्य\-त\-मस्य}}{\lemma{\skt{अन्यतमस्य}}\Afootnote{A H\textsuperscript{pc}; \skt{अन्यतस्य}\ H\textsuperscript{ac}}}\skt{ }\edtext{\skt{नाडीमुभयतो}}{\lemma{\skt{उभयतो}}\Afootnote{H; \skt{द्वि}\ A}}\skt{ }\edtext{\skt{द्वारं}}{\lemma{\skt{द्वारं}}\Afootnote{H; \skt{द्वारां}\ A}}\skt{ }\edtext{\skt{संयोज्य}}{\lemma{\skt{संयोज्य}}\Afootnote{H; \skt{पक्षनाडीं वा संयोज्य}\ A}}\skt{ }\edtext{\skt{दोषो\-द\-क\-म\-व\-सि\-ञ्चेन्}}{\lemma{\skt{अवसिञ्चेन्}}\Afootnote{H; \skt{अवसिञ्चेत्}\ A}}\skt{~। }\edtext{\skt{न}}{\lemma{\skt{न}}\Afootnote{H; \skt{तोतो नाडीमपहृत्य तैल\appbrk{}ल\appbrk{}व\appbrk{}णे\appbrk{}ना\appbrk{}भ्यज्य व्रणं बन्धे\appbrk{}नो\appbrk{}प\appbrk{}च\appbrk{}रेत्, न}\ A}}\skt{ }\edtext{\edtext{\skt{चैकस्मिन्दिवसे}}{\lemma{\skt{चैकस्मिन्}}\Afootnote{H; \skt{चैकस्मि}\ A}}}{\lemma{\skt{दिवसे}}\Afootnote{H; \skt{न्नेव दिवसे}\ A}}\skt{ }\edtext{\skt{सर्व्व\-दो\-षो\-द\-क\-म\-प\-ह\-रेत्}}{\lemma{\skt{सर्व्व\appbrk{}दो\appbrk{}षो\appbrk{}द\appbrk{}कम्}}\Afootnote{H; \skt{सर्वं दोषोदकम्}\ A}}\skt{~॥ }\edtext{\skt{सहसापहृते}}{\lemma{\skt{सहसापहृते}}\Afootnote{H; \skt{सहसा ह्यपहृते}\ A}}\skt{ }\edtext{\skt{दोषे}}{\lemma{\skt{दोषे}}\Afootnote{H; \emph{om.}\ A}}\skt{ तृष्णा\-ज्व\-रा\-ङ्ग\-म\-र्द्दा\-ती\-सा\-र\-श्वा\-स\-का\-स\-पा\-द\-दाहा उत्पद्येरन्~। आपूर्यते वा }\edtext{\skt{दृढ\-त\-र\-म\-स\-ञ्जा\-त\-प्रा\-णस्य}}{\lemma{\skt{दृढतरम्}}\Afootnote{H; \skt{भृश\appbrk{}त\appbrk{}र\appbrk{}मु\appbrk{}द\appbrk{}रम्}\ A}}\skt{ }\edtext{\edtext{\edtext{\skt{तस्मा\-त्तृ\-ती\-य\-ञ्च\-तु\-र्थ\-म्प\-ञ्चमं}}{\lemma{\skt{तृतीयञ्}}\Afootnote{H; \skt{तृतीय}\ A}}}{\lemma{\skt{चतुर्थम्}}\Afootnote{H; \skt{चतुर्थ}\ A}}}{\lemma{\skt{पञ्चमं}}\Afootnote{H; \skt{पञ्चम}\ A}}\skt{ }\edtext{\edtext{\edtext{\skt{षष्ठमष्टमम्वा}}{\lemma{\skt{षष्ठम्}}\Afootnote{H; \skt{षष्ठा}\ A}}}{\lemma{\skt{अष्टमम्}}\Afootnote{H; \skt{ष्टम}\ A}}}{\lemma{\skt{वा}}\Afootnote{H; \skt{दशम}\ A}}\skt{ }\edtext{\edtext{\skt{द्वादशम्वा}}{\lemma{\skt{द्वादशम्}}\Afootnote{H; \skt{द्वादशषोडश}\ A}}}{\lemma{\skt{वा}}\Afootnote{H; \skt{रात्राणाम्}\ A}}\skt{ }\edtext{\edtext{\skt{दिवसमन्तरं}}{\lemma{\skt{दिवसम्}}\Afootnote{H; \skt{अन्यतमम्}\ A}}}{\lemma{\skt{अन्तरं}}\Afootnote{H; \skt{अन्तरीकृत्य}\ A}}\skt{ }\edtext{\edtext{\skt{दोषो\-द\-क\-म\-ल्प\-म\-ल्प\-म\-व\-सि\-ञ्चेत्}}{\lemma{\skt{अल्पम्}}\Afootnote{H; \skt{अल्पा}\ A}}}{\lemma{\skt{अल्पम्}}\Afootnote{H; \skt{ल्पम्}\ A}}\skt{~। }\edtext{\skt{निस्रुते}}{\lemma{\skt{निस्रुते}}\Afootnote{H; \emph{om.}\ A}}\skt{ }\edtext{\skt{निस्रुते}}{\lemma{\skt{निस्रुते}}\Afootnote{H; \skt{निःसृते च}\ A}}\skt{ दोषे }\edtext{\skt{गाढ\-त\-र\-मा\-वि\-क\-कौ\-षे\-य\-च\-र्म्म\-णा\-म\-न्य\-त\-मेन}}{\lemma{\skt{आवि\appbrk{}क\appbrk{}कौ\appbrk{}षे\appbrk{}य\appbrk{}च\appbrk{}र्म्म\appbrk{}णाम्}}\Afootnote{H; \skt{आवि\appbrk{}क\appbrk{}कौ\appbrk{}शे\appbrk{}य\appbrk{}च\appbrk{}र्म\appbrk{}णाम्}\ A}}\skt{ परि\-वे\-ष्ट\-ये\-दु\-द\-र\-न्तथा नाध्मापयति वायुः षण्मासाश्च }\edtext{\edtext{\skt{भोजयेत्पयसा}}{\lemma{\skt{भोजयेत्}}\Afootnote{H; \skt{पयसा भोजयेज्}\ A}}}{\lemma{\skt{पयसा}}\Afootnote{H; \skt{जाङ्गल}\ A}}\skt{ रसेन वा~। }\edtext{\skt{तत\-स्त्री\-न्मा\-सा\-न\-वो\-द\-केन}}{\lemma{\skt{मासा\appbrk{}न\appbrk{}वो\appbrk{}द\appbrk{}केन}}\Afootnote{H\textsuperscript{pc}; \skt{मासा\appbrk{}न\appbrk{}र्धो\appbrk{}द\appbrk{}केन}\ A H\textsuperscript{ac}}}\skt{ }\edtext{\skt{फलाम्लेन}}{\lemma{\skt{फलाम्लेन}}\Afootnote{H; \skt{पयसा फलाम्लेन जाङ्गलरसेन}\ A}}\skt{ }\edtext{\skt{वाव\-शि\-ष्ट\-म्मा\-स\-त्र\-य\-मन्नं}}{\lemma{\skt{वावशिष्टम्}}\Afootnote{H; \skt{वा, अवशिष्टं}\ A}}\skt{ लघु }\edtext{\skt{हितं}}{\lemma{\skt{हितं}}\Afootnote{H; \skt{हितं वा}\ A}}\skt{ सेवेत एवं }\edtext{\skt{सम्व\-त्स\-रे\-णा\-ग\-दी\-भ\-वति}}{\lemma{\skt{सम्व\appbrk{}त्स\appbrk{}रे\appbrk{}णा\appbrk{}ग\appbrk{}दी\appbrk{}भ\appbrk{}वति}}\Afootnote{H; \skt{संव\appbrk{}त्स\appbrk{}रे\appbrk{}णा\appbrk{}गदो भवति}\ A}}\skt{~॥}\pend

\pstart\versepar \ledsidenote{\passref{4.14.19}}\edtext{\skt{भवति}}{\lemma{\skt{भवति}}\Afootnote{A; \skt{भभवति}\ H}}\skt{ चात्र॥ }\pend
\pstart\versepar \skt{आस्थापने चैव विरेचने च पाने }\edtext{\skt{तथा\-हा\-र\-वि\-धि\-क्रि\-यासु}}{\lemma{\skt{तथा\appbrk{}हा\appbrk{}र\appbrk{}वि\appbrk{}धि\appbrk{}क्रि\appbrk{}यासु}}\Afootnote{H; \skt{तथाऽऽ\appbrk{}आ\appbrk{}हा\appbrk{}र\appbrk{}वि\appbrk{}धि\appbrk{}क्रि\appbrk{}यासु}\ A}}\skt{~। }\pend
\pstart\versepar \skt{}\edtext{\skt{सर्व्वोदराणां}}{\lemma{\skt{सर्व्वोदराणां}}\Afootnote{H; \skt{सर्वोदरिभ्यः}\ A}}\skt{ }\edtext{\skt{कुशलेन}}{\lemma{\skt{कुशलेन}}\Afootnote{H; \skt{कुशलैः प्र}\ A}}\skt{ योज्यं क्षीरं }\edtext{\skt{रसो}}{\lemma{\skt{रसो}}\Afootnote{H; \skt{शृतं}\ A}}\skt{ जाङ्गलजो }\edtext{\skt{थ वा}}{\lemma{\skt{थ वा}}\Afootnote{H; \skt{रसो वा}\ A}}\skt{ }\edtext{\skt{स्यादिति}}{\lemma{\skt{स्यादिति}}\Afootnote{H; \emph{om.}\ A}}\skt{~॥}\pend

\pstart \ledsidenote{\passref{4.14.trailer}}\edtext{\skt{चिकित्सा}}{\lemma{\skt{चिकित्सा}}\Afootnote{H; \skt{इति सुश्रु\appbrk{}त\appbrk{}सं\appbrk{}हि\appbrk{}तायां चिकित्सास्थाने उद\appbrk{}र\appbrk{}चि\appbrk{}कि\appbrk{}त्सितं नाम चतुर्दशोऽध्यायः}\ A}}\skt{ १४~॥ ०~॥}\pend

\endnumbering
\beginnumbering
\eledsection*{Adhyāya 15}
\pstart \ledsidenote{\passref{4.15.1}}\skt{अथातो }\edtext{\skt{मूढ\-ग\-र्भ्भ\-चि\-कि\-त्सि\-त\-म्व्या\-ख्या\-स्यामः}}{\lemma{\skt{मूढ\appbrk{}ग\appbrk{}र्भ्भ\appbrk{}चि\appbrk{}कि\appbrk{}त्सि\appbrk{}तम्}}\Afootnote{H; \skt{मूढ\appbrk{}ग\appbrk{}र्भ\appbrk{}चि\appbrk{}कि\appbrk{}त्सितं}\ A}}\skt{~॥}\pend

\pstart \ledsidenote{\passref{4.15.3}}\edtext{\skt{नातः}}{\lemma{\skt{नातः}}\Afootnote{H; \skt{नातोऽन्यत्}\ A}}\skt{ }\edtext{\skt{कष्टतरमस्ति}}{\lemma{\skt{कष्टतरमस्ति}}\Afootnote{H; \skt{कष्टतममस्ति}\ A}}\skt{ यथा }\edtext{\skt{मूढ\-ग\-र्भ्भो\-द्ध\-र\-ण\-मत्र हि}}{\lemma{\skt{मूढ\appbrk{}ग\appbrk{}र्भ्भो\appbrk{}द्ध\appbrk{}र\appbrk{}ण\appbrk{}मत्र हि}}\Afootnote{H; \skt{मूढ\appbrk{}ग\appbrk{}र्भ\appbrk{}श\appbrk{}ल्यो\appbrk{}द्ध\appbrk{}र\appbrk{}णम्; अत्र हि}\ A}}\skt{ }\edtext{\skt{योनि\-य\-कृ\-त्प्ली\-हा\-न्त्र\-ग\-र्भ्भा\-श\-या\-ना\-म्मध्ये}}{\lemma{\skt{योनि\appbrk{}य\appbrk{}कृ\appbrk{}त्प्ली\appbrk{}हा\appbrk{}न्त्र\appbrk{}ग\appbrk{}र्भ्भा\appbrk{}श\appbrk{}या\appbrk{}नाम्}}\Afootnote{H; \skt{योनि\appbrk{}य\appbrk{}कृ\appbrk{}त्प्ली\appbrk{}हा\appbrk{}न्त्र\appbrk{}वि\appbrk{}व\appbrk{}र\appbrk{}ग\appbrk{}र्भा\appbrk{}श\appbrk{}यानां}\ A}}\skt{ }\edtext{\skt{कर्तव्यं}}{\lemma{\skt{कर्तव्यं}}\Afootnote{H; \skt{कर्म कर्तव्यं}\ A}}\skt{~। }\edtext{\skt{स्पर्शेनो}}{\lemma{\skt{स्पर्शेनो}}\Afootnote{H; \skt{स्पर्शेन}\ A}}\skt{ }\edtext{\skt{त्कर्षा\-य\-क\-र्ष\-णा\-प\-क\-र्त्त\-न\-भे\-द\-न\-व्या\-व\-र्त्त\-न\-पी\-ड\-ना\-र्ज्जू\-क\-र\-णानि}}{\lemma{\skt{त्कर्षा\appbrk{}य\appbrk{}क\appbrk{}र्ष\appbrk{}णा\appbrk{}प\appbrk{}क\appbrk{}र्त्त\appbrk{}न\appbrk{}भे\appbrk{}द\appbrk{}न\appbrk{}व्या\appbrk{}व\appbrk{}र्त्त\appbrk{}न\appbrk{}पी\appbrk{}ड\appbrk{}ना\appbrk{}र्ज्जू\appbrk{}क\appbrk{}र\appbrk{}णानि}}\Afootnote{H; \skt{उत्क\appbrk{}र्ष\appbrk{}णा\appbrk{}प\appbrk{}क\appbrk{}र्ष\appbrk{}ण\appbrk{}स्था\appbrk{}ना\appbrk{}प\appbrk{}व\appbrk{}र्त\appbrk{}नो\appbrk{}त्क\appbrk{}र्त\appbrk{}न\appbrk{}भे\appbrk{}द\appbrk{}न\appbrk{}च्छे\appbrk{}द\appbrk{}न\appbrk{}पी\appbrk{}ड\appbrk{}न\appbrk{}र्जू\appbrk{}क\appbrk{}र\appbrk{}ण\appbrk{}दा\appbrk{}र\appbrk{}णानि}\ A}}\skt{ चैकहस्तेन }\edtext{\skt{गर्भ्भमहिंसता}}{\lemma{\skt{गर्भ्भमहिंसता}}\Afootnote{H; \skt{गर्भं}\ A}}\skt{ }\edtext{\skt{गर्भ्भि\-णी\-ञ्चा\-हिं\-सता}}{\lemma{\skt{गर्भ्भिणीञ्}}\Afootnote{H; \skt{गर्भिणीं}\ A}}\skt{ }\edtext{\skt{तस्मादधिकृत्य तमापृच्छ्य}}{\lemma{\skt{तस्मादधिकृत्य तमापृच्छ्य}}\Afootnote{\skt{तस्मा\appbrk{}द\appbrk{}धि\appbrk{}प\appbrk{}ति\appbrk{}मा\appbrk{}पृच्छ्य}\ A; \skt{तस्मा\appbrk{}द\appbrk{}धि\appbrk{}कृ\appbrk{}त्य\appbrk{}मा\appbrk{}पृच्छ्य}\ H\textsuperscript{ac}; \skt{तस्मा\appbrk{}द\appbrk{}धि\appbrk{}कृ\appbrk{}त्य\appbrk{}त\appbrk{}मा\appbrk{}पृच्छ्य}\ H\textsuperscript{pc}}}\skt{ }\edtext{\skt{परञ्च }\edtext{\skt{यन्न\-मा\-स्था\-यो\-प\-क्र\-मेत}}{\lemma{\skt{नमा\appbrk{}स्था\appbrk{}यो\appbrk{}प\appbrk{}क्र\appbrk{}मेत}}\Afootnote{H; \skt{त्नमा\appbrk{}स्था\appbrk{}यो\appbrk{}प\appbrk{}क्र\appbrk{}मेत}\ A}}}{\lemma{\skt{परञ्च यन्}}\Afootnote{H; \skt{परं च य}\ A}}\skt{~॥}\pend

\pstart \ledsidenote{\passref{4.15.4}}\skt{तत्र }\edtext{\skt{समा\-से\-ना\-ष्ट\-विध}}{\lemma{\skt{समा\appbrk{}से\appbrk{}ना\appbrk{}ष्ट\appbrk{}विध}}\Afootnote{H; \skt{समा\appbrk{}से\appbrk{}ना\appbrk{}ष्ट\appbrk{}विधा}\ A}}\skt{ }\edtext{\skt{मूढ\-ग\-र्भ्भ\-ग\-ति\-रु\-द्दिष्टा}}{\lemma{\skt{मूढ\appbrk{}ग\appbrk{}र्भ्भ\appbrk{}ग\appbrk{}ति\appbrk{}रु\appbrk{}द्दिष्टा}}\Afootnote{H; \skt{मूढ\appbrk{}ग\appbrk{}र्भ\appbrk{}ग\appbrk{}ति\appbrk{}रु\appbrk{}द्दिष्टा}\ A}}\skt{ स्वभावगता }\edtext{\skt{वपि तु}}{\lemma{\skt{वपि तु}}\Afootnote{H; \skt{अपि}\ A}}\skt{ त्रयः सङ्गा भवन्ति शिरसो }\edtext{\skt{वैपु\-ल्या\-दं\-श\-यो\-र्ज्ज\-घ\-नस्य}}{\lemma{\skt{वैपु\appbrk{}ल्या\appbrk{}दं\appbrk{}श\appbrk{}योर्}}\Afootnote{H; \skt{वैगु\appbrk{}ण्या\appbrk{}दं\appbrk{}स\appbrk{}योर्}\ A}}\skt{ वा}\pend

\pstart\versepar \ledsidenote{\passref{4.15.5}}\edtext{\skt{जीवति}}{\lemma{\skt{जीवति}}\Afootnote{H; \skt{जीवति तु}\ A}}\skt{ }\edtext{\skt{गर्भ्भे}}{\lemma{\skt{गर्भ्भे}}\Afootnote{H; \skt{गर्भे}\ A}}\skt{ }\edtext{\skt{सूत्यः}}{\lemma{\skt{सूत्यः}}\Afootnote{H; \skt{सूति\appbrk{}का\appbrk{}ग\appbrk{}र्भ\appbrk{}नि\appbrk{}र्ह\appbrk{}रणे}\ A}}\skt{ }\edtext{\skt{प्रयतेरन्}}{\lemma{\skt{प्रयतेरन्}}\Afootnote{H; \skt{प्रयतेत}\ A}}\skt{~। }\pend
\pstart\versepar \skt{}\edtext{\skt{निर्हरणे}}{\lemma{\skt{निर्हरणे}}\Afootnote{H; \skt{निर्हर्तुमशक्ये}\ A}}\skt{ }\edtext{\skt{च्यावनांश्च}}{\lemma{\skt{च्यावनांश्च}}\Afootnote{H; \skt{च्यावनान्}\ A}}\skt{ }\edtext{\skt{मन्त्रा\-न\-नु\-श्रा\-व\-येत्}}{\lemma{\skt{मन्त्रा\appbrk{}न\appbrk{}नु\appbrk{}श्रा\appbrk{}व\appbrk{}येत्}}\Afootnote{H; \skt{मन्त्रा\appbrk{}नु\appbrk{}प\appbrk{}शृ\appbrk{}णु\appbrk{}यात्; तान्वक्ष्यामः}\ A}}\skt{~॥}\pend

\pstart\versepar \ledsidenote{\passref{4.15.6}}\edtext{\skt{इहामृतश्च}}{\lemma{\skt{इहामृतश्च}}\Afootnote{H; \skt{इहामृतं च}\ A}}\skt{ सोमश्च चित्रभानुश्च }\edtext{\skt{भाविनि}}{\lemma{\skt{भाविनि}}\Afootnote{H; \skt{भामिनि}\ A}}\skt{~। }\pend
\pstart\versepar \skt{उच्चैःश्रवाश्च तुरगो मन्दिरे निवसन्तु ते~॥}\pend

\pstart\versepar \ledsidenote{\passref{4.15.7}}\skt{इद\-म\-मृ\-त\-मपां }\edtext{\skt{समुद्धृतां वै}}{\lemma{\skt{समुद्धृतां वै}}\Afootnote{H; \skt{समुद्धृतं वै}\ A}}\skt{ तव लघु }\edtext{\edtext{\skt{गर्भ्भ\-मि\-म\-म्प्र\-मु\-ञ्चहे}}{\lemma{\skt{गर्भ्भमिमम्}}\Afootnote{H; \skt{गर्भमिमं}\ A}}}{\lemma{\skt{प्रमुञ्चहे}}\Afootnote{H; \skt{प्रमुञ्चतु}\ A}}\skt{ स्त्रि~। }\pend
\pstart\versepar \skt{तद\-न\-ल\-प\-व\-ना\-र्क्क\-वा\-स\-वास्ते सह लव\-णा\-म्बु\-ध\-रै\-र्द्दि\-शन्तु शान्तिं~॥}\pend

\pstart \ledsidenote{\passref{4.15.9}}\edtext{\skt{ओषधानि च}}{\lemma{\skt{ओषधानि च}}\Afootnote{H; \skt{औषधानि च}\ A}}\skt{ }\edtext{\skt{विद\-ध्या\-द्य\-थो\-क्तानि च}}{\lemma{\skt{विद\appbrk{}ध्या\appbrk{}द्य\appbrk{}थो\appbrk{}क्तानि च}}\Afootnote{\skt{विद\appbrk{}ध्या\appbrk{}द्य\appbrk{}थो\appbrk{}क्तानि}\ A; \skt{विद\appbrk{}ध्या\appbrk{}द्य\appbrk{}थो\appbrk{}क्तानि च}\ H}}\skt{~। मृते चोत्तानाया }\edtext{\skt{भुग्नसक्थिन्या}}{\lemma{\skt{भुग्नसक्थिन्या}}\Afootnote{H; \skt{आभुग्नसक्थ्या}\ A}}\skt{ }\edtext{\skt{वस्त्र\-चु\-म्भ\-ल\-को\-न्ना\-मि\-त\-कट्या}}{\lemma{\skt{वस्त्र\appbrk{}चु\appbrk{}म्भ\appbrk{}ल\appbrk{}को\appbrk{}न्ना\appbrk{}मि\appbrk{}त\appbrk{}कट्या}}\Afootnote{H; \skt{वस्त्रा\appbrk{}धा\appbrk{}र\appbrk{}को\appbrk{}न्न\appbrk{}मि\appbrk{}त\appbrk{}कठ्या}\ A}}\skt{ }\edtext{\skt{धन्व\-क\-मृ\-त्ति\-का\-शा\-ल्म\-ली\-मृ\-त्स्न\-घृ\-ता\-भ्या\-म्म्र\-क्ष\-यित्वा}}{\lemma{\skt{धन्व\appbrk{}क\appbrk{}मृ\appbrk{}त्ति\appbrk{}का\appbrk{}शा\appbrk{}ल्म\appbrk{}ली\appbrk{}मृ\appbrk{}त्स्न\appbrk{}घृ\appbrk{}ता\appbrk{}भ्याम्}}\Afootnote{H; \skt{धन्व\appbrk{}न\appbrk{}न\appbrk{}ग\appbrk{}वृ\appbrk{}त्ति\appbrk{}का\appbrk{}शा\appbrk{}ल्म\appbrk{}ली\appbrk{}मृ\appbrk{}त्स्न\appbrk{}घृ\appbrk{}ताभ्यां}\ A}}\skt{ हस्तं }\edtext{\skt{योन्याम्प्रवेश्य}}{\lemma{\skt{योन्याम्}}\Afootnote{H; \skt{योनौ}\ A}}\skt{ }\edtext{\edtext{\skt{गर्भ्भ\-म\-प\-ह\-रे\-त्तत्र}}{\lemma{\skt{गर्भ्भम्}}\Afootnote{H; \skt{गर्भम्}\ A}}}{\lemma{\skt{अपहरेत्}}\Afootnote{H; \skt{उपहरेत्}\ A}}\skt{ }\edtext{\edtext{\skt{शक्थि\-भ्या\-मा\-ग\-त\-म\-नु\-लो\-म\-मे\-वा\-ञ्छेत्}}{\lemma{\skt{शक्थिभ्याम्}}\Afootnote{H; \skt{सक्थिभ्याम्}\ A}}\skt{~। एकसक्थि}}{\lemma{\skt{आग\appbrk{}त\appbrk{}म\appbrk{}नु\appbrk{}लो\appbrk{}म\appbrk{}मे\appbrk{}वा\appbrk{}ञ्छेत्~। एकसक्थि}}\Afootnote{\skt{आग\appbrk{}त\appbrk{}म\appbrk{}नु\appbrk{}लो\appbrk{}म\appbrk{}मे\appbrk{}वा\appbrk{}ञ्छेत्, एकसक्थ्ना}\ A; \skt{आग\appbrk{}त\appbrk{}म\appbrk{}नु\appbrk{}लो\appbrk{}म\appbrk{}मे\appbrk{}वा\appbrk{}ञ्छेत्~। एकसक्थि}\ H}}\skt{ }\edtext{\skt{प्रति\-प\-न्न\-स्ये\-तरं}}{\lemma{\skt{प्रति\appbrk{}प\appbrk{}न्न\appbrk{}स्ये\appbrk{}तरं}}\Afootnote{H; \skt{प्रति\appbrk{}प\appbrk{}न्न\appbrk{}स्ये\appbrk{}तर}\ A}}\skt{ सक्थि }\edtext{\skt{प्रसार्य}}{\lemma{\skt{प्रसार्य}}\Afootnote{H; \skt{प्रसार्यापह}\ A}}\skt{ }\edtext{\skt{स्फिग्देशागतं}}{\lemma{\skt{स्फिग्देशागतं}}\Afootnote{H; \skt{रेत्, स्फिग्दे\appbrk{}शे\appbrk{}ना\appbrk{}ग\appbrk{}तस्य}\ A}}\skt{ स्फिग्देशं }\edtext{\edtext{\skt{प्रपी\-ड्यो\-र्द्ध्व\-मु\-त्क्षिप्य}}{\lemma{\skt{प्रपीड्योर्द्ध्वम्}}\Afootnote{A; \skt{प्रपीड्योर्द्ध}\ H}}}{\lemma{\skt{उत्क्षिप्य}}\Afootnote{A; \skt{मुत्क्षिप्य}\ H}}\skt{ सक्थिनी }\edtext{\skt{प्रसा\-र्य\-स्या\-प\-ह\-रेत्}}{\lemma{\skt{प्रसा\appbrk{}र्य\appbrk{}स्या\appbrk{}प\appbrk{}ह\appbrk{}रेत्}}\Afootnote{H; \skt{प्रसा\appbrk{}र्या\appbrk{}प\appbrk{}ह\appbrk{}रेत्}\ A}}\skt{~। तिर्यगागतस्य परिघस्येव }\edtext{\skt{तिर्यक्चीनस्य}}{\lemma{\skt{तिर्यक्चीनस्य}}\Afootnote{H; \skt{तिरश्चीनस्य}\ A}}\skt{ }\edtext{\skt{पश्चा\-द\-र्द्ध\-मू\-र्द्ध\-मु\-त्क्षिप्य}}{\lemma{\skt{पश्चा\appbrk{}द\appbrk{}र्द्ध\appbrk{}मू\appbrk{}र्द्ध\appbrk{}मु\appbrk{}त्क्षिप्य}}\Afootnote{\skt{पश्चा\appbrk{}द\appbrk{}र्ध\appbrk{}मू\appbrk{}र्ध्व\appbrk{}मु\appbrk{}त्क्षिप्य}\ A; \skt{पश्चा\appbrk{}द\appbrk{}र्द्ध\appbrk{}मू\appbrk{}र्द्ध\appbrk{}मु\appbrk{}त्क्षिप्य}\ H}}\skt{ }\edtext{\skt{पूर्व्वा\-र्द्ध\-म\-प\-त्य\-पथं}}{\lemma{\skt{पूर्व्वा\appbrk{}र्द्ध\appbrk{}म\appbrk{}प\appbrk{}त्य\appbrk{}पथं}}\Afootnote{\skt{पूर्वा\appbrk{}र्ध\appbrk{}म\appbrk{}प\appbrk{}त्य\appbrk{}पथं}\ A H}}\skt{ }\edtext{\skt{प्रत्या\-र्ज्ज\-व\-मा\-न\-यि\-त्वा\-प\-ह\-रेत्}}{\lemma{\skt{प्रत्या\appbrk{}र्ज्ज\appbrk{}व\appbrk{}मा\appbrk{}न\appbrk{}यि\appbrk{}त्वा\appbrk{}प\appbrk{}ह\appbrk{}रेत्}}\Afootnote{H; \skt{प्रत्या\appbrk{}र्ज\appbrk{}व\appbrk{}मा\appbrk{}नी\appbrk{}या\appbrk{}प\appbrk{}ह\appbrk{}रेत्, पार्श्वा\appbrk{}प\appbrk{}वृ\appbrk{}त्त\appbrk{}शि\appbrk{}र\appbrk{}स\appbrk{}मंसं प्रपी\appbrk{}ड्यो\appbrk{}र्ध्व\appbrk{}मु\appbrk{}त्क्षिप्य शिरोऽ\appbrk{}प\appbrk{}त्य\appbrk{}प\appbrk{}थ\appbrk{}मा\appbrk{}नी\appbrk{}या\appbrk{}प\appbrk{}ह\appbrk{}रेत्, बाहु\appbrk{}द्व\appbrk{}य\appbrk{}प्र\appbrk{}ति\appbrk{}प\appbrk{}न्न\appbrk{}स्यो\appbrk{}र्ध्व\appbrk{}मु\appbrk{}त्पी\appbrk{}ड्यांसौ शिरोऽ\appbrk{}नु\appbrk{}लो\appbrk{}म\appbrk{}मा\appbrk{}नी\appbrk{}या\appbrk{}प\appbrk{}ह\appbrk{}रेत्}\ A}}\skt{~। }\edtext{\skt{द्वावन्त्यावसाध्यौ}}{\lemma{\skt{द्वावन्त्यावसाध्यौ}}\Afootnote{\skt{द्वावन्त्यावसाध्यौ}\ A H}}\skt{ }\edtext{\skt{मूढगर्भ्भौ}}{\lemma{\skt{मूढगर्भ्भौ}}\Afootnote{H; \skt{मूढगर्}\ A}}\skt{ }\edtext{\skt{भवतः}}{\lemma{\skt{भवतः}}\Afootnote{H; \skt{भौ}\ A}}\skt{~। }\edtext{\edtext{\skt{एवमशक्ये}}{\lemma{\skt{एवम्}}\Afootnote{H; \skt{एव}\ A}}}{\lemma{\skt{अशक्ये}}\Afootnote{H; \skt{मशक्ये}\ A}}\skt{ }\edtext{\edtext{\skt{शस्त्र\-म\-व\-ता\-र\-येत्}}{\lemma{\skt{शस्त्रम्}}\Afootnote{H; \skt{शस्त्र}\ A}}}{\lemma{\skt{अवतारयेत्}}\Afootnote{H; \skt{मवचारयेत्}\ A}}\skt{~॥}\pend

\pstart\versepar \ledsidenote{\passref{4.15.10}}\edtext{\skt{सचेतनन्तु}}{\lemma{\skt{सचेतनन्तु}}\Afootnote{H; \skt{सचेतनं च}\ A}}\skt{ शस्त्रेण न कथञ्चन दारयेत्~। }\pend
\pstart\versepar \skt{}\edtext{\skt{आत्मा\-न\-ञ्ज\-न\-नी\-ञ्चैव}}{\lemma{\skt{जननीञ्}}\Afootnote{H; \emph{om.}\ A}}\skt{ }\edtext{\skt{हन्यादाशु}}{\lemma{\skt{हन्यादाशु}}\Afootnote{H; \skt{घातयेत्}\ A}}\skt{ }\edtext{\skt{सचेतनः}}{\lemma{\skt{सचेतनः}}\Afootnote{H; \emph{om.}\ A}}\skt{~॥}\edtext{}{\lemma{\skt{दारयेत्~। आत्मानञ्}}\Afootnote{H; \skt{दारयेत्~। दार्यमाणो हि जननीमात्मानं}\ A}}\pend

\pstart \ledsidenote{\passref{4.15.12}}\skt{ततः }\edtext{\skt{स्त्रियमाश्वास्य}}{\lemma{\skt{स्त्रियमाश्वास्य}}\Afootnote{\skt{स्त्रियमाश्वास्य}\ A H}}\skt{ }\edtext{\skt{मण्ड\-ला\-ग्रे\-णा\-ङ्गु\-लि\-श\-स्त्रेण वा}}{\lemma{\skt{मण्ड\appbrk{}ला\appbrk{}ग्रे\appbrk{}णा\appbrk{}ङ्गु\appbrk{}लि\appbrk{}श\appbrk{}स्त्रेण वा}}\Afootnote{H; \skt{मण्ड\appbrk{}ला\appbrk{}ग्रे\appbrk{}णा\appbrk{}ङ्गु\appbrk{}ली\appbrk{}श\appbrk{}स्त्रेण वा}\ A}}\skt{ शिरो विदार्य }\edtext{\skt{शिरः\-क\-पा\-ला\-न्या\-हृत्य}}{\lemma{\skt{शिरः\appbrk{}क\appbrk{}पा\appbrk{}ला\appbrk{}न्या\appbrk{}हृत्य}}\Afootnote{\skt{शिरः\appbrk{}क\appbrk{}पा\appbrk{}ला\appbrk{}न्या\appbrk{}हृत्य}\ A H}}\skt{ शंकुना }\edtext{\edtext{\skt{संगृह्योरसि}}{\lemma{\skt{संगृह्योर्}}\Afootnote{H; \skt{गृहीत्वो}\ A}}}{\lemma{\skt{असि}}\Afootnote{H; \skt{रसि}\ A}}\skt{ }\edtext{\skt{कक्षे}}{\lemma{\skt{कक्षे}}\Afootnote{H; \skt{कक्षायां}\ A}}\skt{ }\edtext{\skt{चापहरेत्~। }\edtext{\skt{अभि\-न्न\-शि\-र\-म्वा\-क्षि\-कूटो}}{\lemma{\skt{वाक्षिकूटो}}\Afootnote{H; \skt{समक्षिकूटे}\ A}}}{\lemma{\skt{चापहरेत्~। अभिन्नशिरम्}}\Afootnote{H; \skt{वाऽपहरेत्; अभिन्नशिर}\ A}}\skt{ गण्डे वा }\edtext{\skt{असं\-स\-क्त\-स्यां\-श\-देशे}}{\lemma{\skt{असं\appbrk{}स\appbrk{}क्त\appbrk{}स्यां\appbrk{}श\appbrk{}देशे}}\Afootnote{\skt{अंस\appbrk{}सं\appbrk{}स\appbrk{}क्त\appbrk{}स्यां\appbrk{}स\appbrk{}देशे}\ A; \skt{असं\appbrk{}से\appbrk{}क्त\appbrk{}स्यां\appbrk{}श\appbrk{}देशे}\ H}}\skt{ }\edtext{\skt{बाहूं}}{\lemma{\skt{बाहूं}}\Afootnote{H\textsuperscript{pc}; \skt{बाहू}\ A; \skt{बाहूं>।}\ H\textsuperscript{ac}}}\skt{ }\edtext{\skt{च्छ्हित्वा}}{\lemma{\skt{च्छ्हित्वा}}\Afootnote{H; \skt{छित्त्वा}\ A}}\skt{ }\edtext{\edtext{\skt{दृति\-मि\-वा\-त\-त\-म्वा\-त\-पू\-र्ण्णो\-द\-र\-म्वि\-दार्य}}{\lemma{\skt{दृति\appbrk{}मि\appbrk{}वा\appbrk{}त\appbrk{}तम्}}\Afootnote{\skt{दृतिमिवाततं}\ A H}}}{\lemma{\skt{वात\appbrk{}पू\appbrk{}र्ण्णो\appbrk{}द\appbrk{}रम्}}\Afootnote{H; \skt{वातपूर्णोदरं वा}\ A}}\skt{ }\edtext{\skt{निरस्यात्राणि}}{\lemma{\skt{निरस्यात्राणि}}\Afootnote{H; \skt{निरस्यान्त्राणि}\ A}}\skt{ }\edtext{\skt{शिथिलीभूतां}}{\lemma{\skt{शिथिलीभूतां}}\Afootnote{H; \skt{शिथि\appbrk{}ली\appbrk{}भू\appbrk{}त\appbrk{}मा\appbrk{}ह\appbrk{}रेत्}\ A}}\skt{~। }\edtext{\skt{जघनशक्तस्य वा}}{\lemma{\skt{जघनशक्तस्य वा}}\Afootnote{H; \skt{जघनसक्तस्य वा}\ A}}\skt{ }\edtext{\skt{जघ\-न\-क\-पा\-ला\-नीति}}{\lemma{\skt{जघ\appbrk{}न\appbrk{}क\appbrk{}पा\appbrk{}ला\appbrk{}नीति}}\Afootnote{A H\textsuperscript{pc}; \skt{जघ\appbrk{}न\appbrk{}क\appbrk{}पा\appbrk{}नीति}\ H\textsuperscript{ac}}}\skt{~।}\pend

\pstart\versepar \ledsidenote{\passref{4.15.13}}\edtext{\skt{यस्मिन्नङ्गे हि यो}}{\lemma{\skt{यस्मिन्नङ्गे हि यो}}\Afootnote{\skt{किंबहुना यद्यदङ्गं हि}\ A; \skt{यस्मिन्नङ्गे हि यो}\ H}}\skt{ }\edtext{\skt{गर्भ्भः}}{\lemma{\skt{गर्भ्भः}}\Afootnote{H; \skt{गर्भस्य तस्य}\ A}}\skt{ }\edtext{\skt{सज्जते}}{\lemma{\skt{सज्जते}}\Afootnote{H; \skt{सज्जति}\ A}}\skt{ }\edtext{\skt{तस्य}}{\lemma{\skt{तस्य}}\Afootnote{H; \skt{तद्}\ A}}\skt{ }\edtext{\skt{तम्भिषक्}}{\lemma{\skt{तम्भिषक्}}\Afootnote{H; \skt{भिषक्}\ A}}\skt{~। }\pend
\pstart\versepar \skt{}\edtext{\edtext{\skt{सम्यक्छित्त्वा}}{\lemma{\skt{सम्यक्}}\Afootnote{H; \skt{सम्य}\ A}}}{\lemma{\skt{छित्त्वा}}\Afootnote{H; \skt{ग्वि}\ A}}\skt{ }\edtext{\skt{निर्ह\-रे\-द्र\-क्षे\-न्ना\-रीञ्च}}{\lemma{\skt{निर्हरेद्}}\Afootnote{\skt{निर्हरेच्छित्त्वा}\ A; \skt{निर्हरेद्}\ H}}\skt{ यत्नतः~॥}\pend

\pstart\versepar \ledsidenote{\passref{4.15.14}}\edtext{\skt{चित्रां हि}}{\lemma{\skt{चित्रां हि}}\Afootnote{H; \emph{om.}\ A}}\skt{ }\edtext{\skt{गतिं}}{\lemma{\skt{गतिं}}\Afootnote{H; \emph{om.}\ A}}\skt{ }\edtext{\skt{गर्भ्भस्य}}{\lemma{\skt{गर्भ्भस्य}}\Afootnote{H; \skt{गर्भस्य}\ A}}\skt{ }\edtext{\skt{करोति}}{\lemma{\skt{करोति}}\Afootnote{H; \skt{गतयश्चित्}\ A}}\skt{ }\edtext{\edtext{\skt{कुपितोऽनिलः}}{\lemma{\skt{कुपितो}}\Afootnote{H; \skt{रा जायन्ते}\ A}}}{\lemma{\skt{ऽनिलः}}\Afootnote{H; \skt{ऽनिलकोपतः}\ A}}\skt{~। }\pend
\pstart\versepar \skt{}\edtext{\skt{तत्रा\-न\-ल्प\-म\-ति\-र्वैद्यो}}{\lemma{\skt{तत्रा\appbrk{}न\appbrk{}ल्प\appbrk{}म\appbrk{}ति\appbrk{}र्वैद्यो}}\Afootnote{\skt{तत्रा\appbrk{}न\appbrk{}ल्प\appbrk{}म\appbrk{}ति\appbrk{}र्वैद्यो}\ A H}}\skt{ }\edtext{\skt{वर्त्येत}}{\lemma{\skt{वर्त्येत}}\Afootnote{H; \skt{वर्तेत}\ A}}\skt{ }\edtext{\skt{मतिपूर्व्वकं}}{\lemma{\skt{मतिपूर्व्वकं}}\Afootnote{H; \skt{विधिपूर्वकम्}\ A}}\skt{~॥}\pend

\pstart\versepar \ledsidenote{\passref{4.15.15}}\edtext{\skt{न चोपेक्षेत्मृतं}}{\lemma{\skt{न चोपेक्षेत्}}\Afootnote{H; \skt{नोपेक्षेत}\ A}}\skt{ }\edtext{\skt{गर्भ्भ\-म्मु\-हू\-र्त्त\-मपि}}{\lemma{\skt{गर्भ्भम्}}\Afootnote{H; \skt{गर्भं}\ A}}\skt{ पण्डितः~। }\pend
\pstart\versepar \skt{स ह्याशु जननीं हन्ति निरुच्छ्वासम्पशुं यथा~॥}\pend

\pstart\versepar \ledsidenote{\passref{4.15.16}}\skt{मण्डलाग्रेण कर्त्तव्यं }\edtext{\edtext{\skt{च्छेद्य\-म\-न्त\-र्व्वि\-जा\-नता}}{\lemma{\skt{च्छेद्यम्}}\Afootnote{A; \skt{च्छेद्य}\ H}}}{\lemma{\skt{अन्तर्व्विजानता}}\Afootnote{\skt{अन्तर्विजानता}\ A; \skt{मन्तर्व्विजानता}\ H}}\skt{~। }\pend
\pstart\versepar \skt{वृद्धिपत्रं हि तीक्ष्णाग्रं नारीं }\edtext{\edtext{\skt{हन्यात्कथञ्चन}}{\lemma{\skt{हन्यात्}}\Afootnote{H; \skt{हिंस्यात्}\ A}}}{\lemma{\skt{कथञ्}}\Afootnote{H; \skt{कदा}\ A}}\skt{~॥}\pend

\pstart\versepar \ledsidenote{\passref{4.15.17}}\edtext{\skt{अथा\-म\-प\-त\-न्ती\-म\-प\-रा\-म्पा\-त\-ये\-त्पू\-र्व्व\-व\-द्भि\-षक्}}{\lemma{\skt{अथा\appbrk{}म\appbrk{}प\appbrk{}त\appbrk{}न्तीम्}}\Afootnote{H; \skt{अथापतन्तीम्}\ A}}\skt{~। }\pend
\pstart\versepar \skt{हस्ते\-ना\-प\-ह\-रेद्वा }\edtext{\skt{पि}}{\lemma{\skt{पि}}\Afootnote{H; \skt{ऽपि}\ A}}\skt{ }\edtext{\skt{पार्श्वा\-भ्या\-म्प\-रि\-पीड्य}}{\lemma{\skt{पार्श्वाभ्याम्}}\Afootnote{H; \skt{पार्श्वभ्यां}\ A}}\skt{ वा~॥}\pend

\pstart\versepar \ledsidenote{\passref{4.15.18}}\edtext{\edtext{\edtext{\skt{मुहु\-र्न्ना\-री\-म\-व\-धु\-ने\-त्पी\-ड\-ये\-द्वां\-श\-पि\-ण्डि\-कान्}}{\lemma{\skt{मुहुर्न्नारीम्}}\Afootnote{H; \skt{धुनुयाच्च मुहुर्नारीं}\ A}}}{\lemma{\skt{अवधुनेत्}}\Afootnote{H; \emph{om.}\ A}}}{\lemma{\skt{वांशपिण्डिकान्}}\Afootnote{H; \skt{वाऽंसपिण्डिकाम्}\ A}}\skt{~। }\pend
\pstart\versepar \skt{}\edtext{\edtext{\skt{तैला\-भ्य\-क्त\-यो\-नि\-न्तै\-लेन}}{\lemma{\skt{तैला\appbrk{}भ्य\appbrk{}क्त\appbrk{}यो\appbrk{}निन्}}\Afootnote{H; \skt{तैलाक्तयोने}\ A}}}{\lemma{\skt{तैलेन}}\Afootnote{H; \skt{रेवं तां}\ A}}\skt{ }\edtext{\skt{पाय\-ये\-त्म\-ति\-मा\-न्भि\-षक्}}{\lemma{\skt{पाययेत्}}\Afootnote{H; \skt{पातयेन्}\ A}}\skt{~॥}\pend

\pstart \ledsidenote{\passref{4.15.18a}}\skt{हिङ्गु\-पि\-प्प\-लि\-पा\-ट\-ल्यौ\-भा\-र्ग्गी\-मे\-ला\-म\-हौ\-षधं~।}\pend

\pstart \ledsidenote{\passref{4.15.21}}\edtext{\skt{रास्नामतिविषा}}{\lemma{\skt{रास्नामतिविषा}}\Afootnote{H; \skt{वचामतिविषां रा}\ A}}\skt{ }\edtext{\edtext{\edtext{\skt{चव्यामेभिर्द्दोषः}}{\lemma{\skt{चव्याम्}}\Afootnote{H; \skt{स्नां चव्यं}\ A}}}{\lemma{\skt{एभिर्}}\Afootnote{H; \emph{om.}\ A}}}{\lemma{\skt{द्दोषः}}\Afootnote{H; \emph{om.}\ A}}\skt{ }\edtext{\skt{प्रसिचय्ते}}{\lemma{\skt{प्रसिचय्ते}}\Afootnote{H; \skt{संचूर्ण्य पाययेत्~। स्नेहेन दोषस्यन्दार्थं वेद\appbrk{}नो\appbrk{}प\appbrk{}श\appbrk{}माय च}\ A}}\skt{~॥}\pend

\pstart \ledsidenote{\passref{4.15.20}}\edtext{\skt{योनिश्च}}{\lemma{\skt{योनिश्च}}\Afootnote{H; \skt{एवं}\ A}}\skt{ मृद्वी }\edtext{\skt{भवति}}{\lemma{\skt{भवति}}\Afootnote{H; \skt{भवेद्}\ A}}\skt{ }\edtext{\edtext{\skt{योनिशूलश्च}}{\lemma{\skt{योनिशूलश्}}\Afootnote{H; \skt{योनिस्तच्छूलं}\ A}}}{\lemma{\skt{च}}\Afootnote{H; \skt{चोप}\ A}}\skt{ }\edtext{\skt{शाम्यति}}{\lemma{\skt{शाम्यति}}\Afootnote{H; \skt{शाम्यति~। कृष्णा\appbrk{}त\appbrk{}न्मू\appbrk{}ल\appbrk{}शु\appbrk{}ण्ठ्ये\appbrk{}ला\appbrk{}हि\appbrk{}ङ्गु\appbrk{}भार्गीः सदीप्यकाः}\ A}}\skt{~।}\pend

\pstart \ledsidenote{\passref{4.15.23}}\edtext{\edtext{\skt{त्रिरात्रमेवं}}{\lemma{\skt{त्रिरात्रम्}}\Afootnote{H; \skt{तथा तेजोवतीं चापि पाय\appbrk{}ये\appbrk{}त्पू\appbrk{}र्व\appbrk{}व\appbrk{}द्भि\appbrk{}षक्~। त्रिरात्रं}\ A}}}{\lemma{\skt{एवं}}\Afootnote{H; \skt{पञ्च}\ A}}\skt{ सप्ताहं ततः स्नेहम्पुनः }\edtext{\skt{पुनः}}{\lemma{\skt{पुनः}}\Afootnote{H; \skt{पिबेत्}\ A}}\skt{~॥}\pend

\pstart\versepar \ledsidenote{\passref{4.15.24}}\edtext{\skt{सायञ्च}}{\lemma{\skt{सायञ्च}}\Afootnote{H; \emph{om.}\ A}}\skt{ }\edtext{\edtext{\skt{पाययेद्दद्या}}{\lemma{\skt{पाययेद्}}\Afootnote{H; \skt{पाययेता}\ A}}}{\lemma{\skt{दद्या}}\Afootnote{H; \skt{सवं न}\ A}}\skt{ }\edtext{\edtext{\skt{मासम्वापि}}{\lemma{\skt{मासम्}}\Afootnote{H; \skt{क्तमरिष्टं}\ A}}}{\lemma{\skt{वापि}}\Afootnote{H; \skt{वा}\ A}}\skt{ }\edtext{\skt{सुसंस्कृतां}}{\lemma{\skt{सुसंस्कृतां}}\Afootnote{H; \skt{सुसंस्कृतम्}\ A}}\skt{~। }\pend
\pstart\versepar \skt{शिरी\-ष\-क\-कु\-भा\-भ्याञ्च तोयमाचमने हितं~॥}\pend

\pstart \ledsidenote{\passref{4.15.25}}\edtext{\skt{सर्व्वतः}}{\lemma{\skt{सर्व्वतः}}\Afootnote{H; \skt{उपद्रवाश्च येऽन्ये स्युस्ता\appbrk{}न्य\appbrk{}था\appbrk{}स्व\appbrk{}मु\appbrk{}पा\appbrk{}च\appbrk{}रेत्~। सर्वतः}\ A}}\skt{ }\edtext{\skt{परिशुद्द्धा च}}{\lemma{\skt{परिशुद्द्धा च}}\Afootnote{H; \skt{परिशुद्धा च}\ A}}\skt{ स्निग्ध\-प\-थ्या\-ल्प\-भो\-जना~॥}\pend

\pstart\versepar \ledsidenote{\passref{4.15.26}}\skt{स्वेदाभ्यङ्गपरा }\edtext{\skt{नित्य\-म्भ\-वे\-त्मा\-स\-म\-त\-न्द्रिता}}{\lemma{\skt{मासमतन्द्रिता}}\Afootnote{\skt{क्रोधविवर्जिता}\ A; \skt{मासमतन्द्रिता}\ H}}\skt{~॥ }\pend
\pstart\versepar \skt{पयो वातहरैः सिद्धन्दशाहं भोजने हितं~।}\pend

\pstart\versepar \ledsidenote{\passref{4.15.27}}\skt{रसन्दशाहं }\edtext{\skt{शेषन्तु}}{\lemma{\skt{शेषन्तु}}\Afootnote{H; \skt{शेषे तु}\ A}}\skt{ }\edtext{\skt{यथा\-यो\-ग\-मु\-पा\-च\-रेत्}}{\lemma{\skt{यथा\appbrk{}यो\appbrk{}ग\appbrk{}मु\appbrk{}पा\appbrk{}च\appbrk{}रेत्}}\Afootnote{\skt{यथा\appbrk{}यो\appbrk{}ग\appbrk{}मु\appbrk{}पा\appbrk{}च\appbrk{}रेत्}\ A H}}\skt{~॥ }\pend
\pstart\versepar \skt{}\edtext{\skt{निरु\-प\-द्र\-वा\-म्वि\-शु\-द्धाञ्च}}{\lemma{\skt{निरुपद्रवाम्}}\Afootnote{H; \skt{व्युपद्रवां}\ A}}\skt{ ज्ञात्वा च }\edtext{\skt{वलवर्ण्णिनीं}}{\lemma{\skt{वलवर्ण्णिनीं}}\Afootnote{H; \skt{वरवर्णिनीम्}\ A}}\skt{~॥}\pend

\pstart\versepar \ledsidenote{\passref{4.15.28}}\edtext{\skt{ऊर्द्ध्वञ्चतुर्षु}}{\lemma{\skt{चतुर्षु}}\Afootnote{H; \skt{चतुर्भ्यो}\ A}}\skt{ }\edtext{\skt{मासेषु}}{\lemma{\skt{मासेषु}}\Afootnote{H; \skt{मासेभ्यो}\ A}}\skt{ }\edtext{\skt{विसृजेत्तु}}{\lemma{\skt{विसृजेत्तु}}\Afootnote{H; \skt{विसृजेत्प}\ A}}\skt{ }\edtext{\skt{विचक्षणः}}{\lemma{\skt{विचक्षणः}}\Afootnote{H; \skt{रिहारतः}\ A}}\skt{~॥ }\pend
\pstart\versepar \skt{योनि\-स\-न्त\-र्प्प\-णेऽ\-भ्यङ्गे पाने बस्तिषु भोजने~।}\pend

\pstart\versepar \ledsidenote{\passref{4.15.29}}\skt{बला\-तै\-ल\-मि\-द\-ञ्चास्यै }\edtext{\skt{दद्या\-द\-नि\-ल\-वा\-रणं}}{\lemma{\skt{दद्या\appbrk{}द\appbrk{}नि\appbrk{}ल\appbrk{}वा\appbrk{}रणं}}\Afootnote{H; \skt{दद्या\appbrk{}द\appbrk{}नि\appbrk{}ल\appbrk{}वा\appbrk{}र\appbrk{}णम्}\ A}}\skt{~॥ }\pend
\pstart\versepar \skt{बला\-मू\-ल\-क\-षा\-यस्य }\edtext{\skt{दश\-मू\-ली\-कृ\-तस्य}}{\lemma{\skt{दश\appbrk{}मू\appbrk{}ली\appbrk{}कृ\appbrk{}तस्य}}\Afootnote{H; \skt{दश\appbrk{}मू\appbrk{}ली\appbrk{}शृ\appbrk{}तस्य}\ A}}\skt{ च~।}\pend

\pstart\versepar \ledsidenote{\passref{4.15.30}}\skt{यव\-को\-ल\-कु\-ल\-त्थानां क्वाथस्य पयसस्तथा~॥ }\pend
\pstart\versepar \skt{अष्टावष्टौ शुभा }\edtext{\edtext{\skt{भागा\-तै\-ल\-स्ये\-क\-स्त\-दे\-कतः}}{\lemma{\skt{भागा\appbrk{}तै\appbrk{}ल\appbrk{}स्ये\appbrk{}कस्}}\Afootnote{H; \skt{भागा\appbrk{}स्तै\appbrk{}ला\appbrk{}दे\appbrk{}कस्}\ A}}}{\lemma{\skt{तदेकतः}}\Afootnote{H; \skt{तदेकतः}\ A}}\skt{~।}\pend

\pstart\versepar \ledsidenote{\passref{4.15.31}}\edtext{\skt{पचेदावाप्य}}{\lemma{\skt{पचेदावाप्य}}\Afootnote{H; \skt{पचेदावाप्य}\ A}}\skt{ मधुरं गणं सैन्धवसंयुतं~॥ }\pend
\pstart\versepar \skt{}\edtext{\skt{तत्रागरुं}}{\lemma{\skt{तत्रागरुं}}\Afootnote{H; \skt{तथाऽगुरुं}\ A}}\skt{ सर्ज्जरसं सर\-ल\-न्दे\-व\-दारु च~।}\pend

\pstart\versepar \ledsidenote{\passref{4.15.32}}\skt{मञ्जि\-ष्ठा\-ञ्च\-न्द\-न\-ङ्कु\-ष्ठ\-मेलां कालानुसारिवां~॥ }\pend
\pstart\versepar \skt{}\edtext{\skt{मान्सी}}{\lemma{\skt{मान्सी}}\Afootnote{H; \skt{मांसीं}\ A}}\skt{ शैले\-य\-क\-म्प\-त्त्र\-न्त\-गरं }\edtext{\skt{शारिवां}}{\lemma{\skt{शारिवां}}\Afootnote{H; \skt{सारिवां}\ A}}\skt{ वचां~।}\pend

\pstart\versepar \ledsidenote{\passref{4.15.33}}\edtext{\skt{शता\-व\-री\-म\-श्व\-गन्धां}}{\lemma{\skt{शता\appbrk{}व\appbrk{}री\appbrk{}म\appbrk{}श्व\appbrk{}गन्धां}}\Afootnote{\skt{शता\appbrk{}व\appbrk{}री\appbrk{}म\appbrk{}श्व\appbrk{}गन्धां}\ A H}}\skt{ }\edtext{\skt{शत\-पु\-ष्पा\-म्पु\-न\-र्ण्णवां}}{\lemma{\skt{पुनर्ण्णवां}}\Afootnote{H; \skt{पुनर्नवाम्}\ A}}\skt{~॥ }\pend
\pstart\versepar \skt{तत्साधुसिद्धं सौवर्ण्णे राजते मृन्मयेऽपि वा~।}\pend

\pstart\versepar \ledsidenote{\passref{4.15.34}}\skt{प्रक्षिप्य }\edtext{\edtext{\skt{कलसेऽभ्यर्च्च्य}}{\lemma{\skt{कलसे}}\Afootnote{H; \skt{कलशे}\ A}}}{\lemma{\skt{ऽभ्यर्च्च्य}}\Afootnote{H; \skt{सम्यक्}\ A}}\skt{ स्वनु\-गु\-प्त\-न्नि\-धा\-प\-येत्~॥ }\pend
\pstart\versepar \skt{}\edtext{\skt{बला\-तै\-ल\-मि\-द\-न्नाम्ना}}{\lemma{\skt{नाम्ना}}\Afootnote{H; \skt{ख्यातं}\ A}}\skt{ सर्व्व\-वा\-त\-वि\-का\-र\-नुत्~।}\pend

\pstart\versepar \ledsidenote{\passref{4.15.35}}\edtext{\skt{यथाबलमतो}}{\lemma{\skt{यथाबलमतो}}\Afootnote{\skt{यथाबलमतो}\ A H}}\skt{ मात्रां सूतिकायै प्रदापयेत्~॥ }\pend
\pstart\versepar \skt{या च }\edtext{\skt{गर्भ्भार्थिनी}}{\lemma{\skt{गर्भ्भार्थिनी}}\Afootnote{H; \skt{गर्भार्थिनी}\ A}}\skt{ नारी क्षीणशुक्रश्च यः पुमान्~।}\pend

\pstart\versepar \ledsidenote{\passref{4.15.36}}\skt{वातक्षीणे मर्म्महते मथितेऽभिहते तथा~॥ }\pend
\pstart\versepar \skt{भग्ने }\edtext{\skt{स्रमाभिपन्ने च}}{\lemma{\skt{स्रमाभिपन्ने च}}\Afootnote{H; \skt{श्रमाभिपन्ने च}\ A}}\skt{ }\edtext{\skt{सर्व्वथैव प्रयुज्यते}}{\lemma{\skt{सर्व्वथैव प्रयुज्यते}}\Afootnote{H; \skt{सर्व\appbrk{}थै\appbrk{}वो\appbrk{}प\appbrk{}यु\appbrk{}ज्यते}\ A}}\skt{~।}\pend

\pstart\versepar \ledsidenote{\passref{4.15.37}}\edtext{\skt{सर्व्वा\-ना\-पे\-क्ष\-का\-दींश्च}}{\lemma{\skt{सर्व्वा\appbrk{}ना\appbrk{}पे\appbrk{}क्ष\appbrk{}का\appbrk{}दींश्च}}\Afootnote{H; \skt{एत\appbrk{}दा\appbrk{}क्षे\appbrk{}प\appbrk{}का\appbrk{}दीन्वै}\ A}}\skt{ वात\-व्या\-धी\-न\-पो\-हति~॥ }\pend
\pstart\versepar \skt{}\edtext{\edtext{\skt{हिक्का\-श्वा\-स\-म\-धी\-म\-न्थ\-ङ्गु\-ल्म\-ङ्का\-सञ्च}}{\lemma{\skt{हिक्का\appbrk{}श्वा\appbrk{}स\appbrk{}म\appbrk{}धी\appbrk{}म\appbrk{}न्थङ्}}\Afootnote{\skt{हिक्कां कासमधीमन्थं}\ A; \skt{हिक्का श्वास\appbrk{}म\appbrk{}धी\appbrk{}म\appbrk{}न्थङ्}\ H}}}{\lemma{\skt{कासञ्च}}\Afootnote{H; \skt{श्वासं च}\ A}}\skt{ दुस्तरं~।}\pend

\pstart\versepar \ledsidenote{\passref{4.15.38}}\edtext{\skt{षण्मा\-सा\-नु\-प\-यु\-ज्येन}}{\lemma{\skt{षण्मा\appbrk{}सा\appbrk{}नु\appbrk{}प\appbrk{}यु\appbrk{}ज्येन}}\Afootnote{H; \skt{षण्मा\appbrk{}सा\appbrk{}नु\appbrk{}प\appbrk{}यु\appbrk{}ज्यैत}\ A}}\skt{ }\edtext{\edtext{\skt{मन्त्र\-वृ\-द्धि\-र्व्व्य\-पो\-हति}}{\lemma{\skt{मन्त्रवृद्धिर्}}\Afootnote{\skt{दन्त्रवृद्धिम्}\ A; \skt{मन्त्रवृद्धि}\ H}}}{\lemma{\skt{व्व्यपोहति}}\Afootnote{\skt{अपोहति}\ A; \skt{र्व्व्यपोह-ति}\ H}}\skt{~॥ }\pend
\pstart\versepar \skt{प्रत्यग्रधातुः पुरुषो भवेच्च स्थिरयौवनः~।}\pend

\pstart\versepar \ledsidenote{\passref{4.15.39}}\edtext{\edtext{\skt{एत\-द्धि\-रा\-ज्ञा\-ङ्क\-र्त्तव्यं}}{\lemma{\skt{एतद्}}\Afootnote{H; \skt{राज्ञामेतद्}\ A}}}{\lemma{\skt{धिराज्ञाङ्}}\Afootnote{H; \skt{धि}\ A}}\skt{ }\edtext{\skt{राजामात्याश्च ये}}{\lemma{\skt{राजामात्याश्च ये}}\Afootnote{H; \skt{राजमात्राश्च ये}\ A}}\skt{ नराः~॥ }\pend
\pstart\versepar \skt{सुखिनः सुकुमाराश्च धनिनश्चापि ये नराः~।}\pend

\pstart\versepar \ledsidenote{\passref{4.15.40}}\skt{बला\-क\-षा\-य\-पी\-ते\-भ्य\-स्ति\-लेभ्यो वा }\edtext{\skt{प्यनेकशः}}{\lemma{\skt{प्य्}}\Afootnote{H; \skt{ऽप्य्}\ A}}\skt{~॥ }\pend
\pstart\versepar \skt{तैलमुत्पाद्य }\edtext{\skt{तत्तैलं}}{\lemma{\skt{तैलं}}\Afootnote{H; \skt{क्वाथ}\ A}}\skt{ }\edtext{\skt{शतपाककृतं}}{\lemma{\skt{शतपाककृतं}}\Afootnote{H; \skt{शतपाकं कृतं}\ A}}\skt{ शुभं~।}\pend

\pstart\versepar \ledsidenote{\passref{4.15.41}}\edtext{\edtext{\skt{शस्तेऽह्नि}}{\lemma{\skt{शस्ते}}\Afootnote{H; \emph{om.}\ A}}}{\lemma{\skt{ऽह्नि}}\Afootnote{H; \skt{नि}\ A}}\skt{ }\edtext{\edtext{\skt{निभृ\-त\-वे\-श्म\-न्यु\-प\-यु\-ञ्ज्या\-द्य\-था\-बलं}}{\lemma{\skt{निभृतवेश्मन्य्}}\Afootnote{\skt{वाते निभृतागारे}\ A; \skt{भृवृते वेश्मन्य्}\ H\textsuperscript{ac}; \skt{निभृवृते वेश्मन्य्}\ H\textsuperscript{pc}}}}{\lemma{\skt{उपयुञ्ज्याद्}}\Afootnote{H; \skt{प्रयुञ्जीत}\ A}}\skt{~॥ }\pend
\pstart\versepar \skt{}\edtext{\skt{जीर्ण्णेस्मिन्पयसा}}{\lemma{\skt{जीर्ण्णेस्मिन्}}\Afootnote{H; \skt{जीर्णेऽस्मिन्}\ A}}\skt{ }\edtext{\skt{सिद्धं}}{\lemma{\skt{सिद्धं}}\Afootnote{\skt{स्निग्धम्}\ A; \skt{सि}\ H}}\skt{ }\edtext{\edtext{\skt{षष्टि\-का\-न्न\-म\-वा\-प्नु\-यात्}}{\lemma{\skt{षष्टिकान्नम्}}\Afootnote{\skt{अश्नी\appbrk{}या\appbrk{}त्ष\appbrk{}ष्टि\appbrk{}कौ\appbrk{}द\appbrk{}नम्}\ A; \skt{द्धं षष्टिकान्नम्}\ H}}}{\lemma{\skt{अवाप्नुयात्}}\Afootnote{H; \emph{om.}\ A}}\skt{~।}\pend

\pstart\versepar \ledsidenote{\passref{4.15.42}}\skt{अनेन विधिना }\edtext{\skt{द्रोण\-मु\-प\-यु\-ञ्ज्या\-न्न\-मी\-रितं}}{\lemma{\skt{उप\appbrk{}यु\appbrk{}ञ्ज्या\appbrk{}न्न\appbrk{}मी\appbrk{}रितं}}\Afootnote{H; \skt{उप\appbrk{}यु\appbrk{}ज्या\appbrk{}न्न\appbrk{}मी\appbrk{}रि\appbrk{}तम्}\ A}}\skt{~॥ }\pend
\pstart\versepar \skt{भुञ्जीत द्विगुणं कालं बल\-व\-र्ण्णा\-न्वि\-त\-स्ततः~।}\pend

\pstart\versepar \ledsidenote{\passref{4.15.43}}\skt{सर्व्व\-पा\-पै\-र्व्वि\-नि\-र्म्मुक्तः शतायुः पुरुषो भवेत्~॥ }\pend
\pstart\versepar \skt{शतं शतं }\edtext{\skt{चोत्कर्षाद्द्रोणे}}{\lemma{\skt{चोत्कर्षाद्}}\Afootnote{H; \skt{तथोत्कर्षो}\ A}}\skt{ द्रोणे }\edtext{\skt{प्रकीर्त्त्यते}}{\lemma{\skt{प्रकीर्त्त्यते}}\Afootnote{H; \skt{प्रकीर्तितः}\ A}}\skt{~॥}\pend

\pstart\versepar \ledsidenote{\passref{4.15.44}}\edtext{\edtext{\skt{बला\-क\-ल्पे\-ना\-ति\-ब\-ला\-गु\-डू\-च्या\-दि\-त्य\-व\-ल्लिषु}}{\lemma{\skt{बला\appbrk{}क\appbrk{}ल्पे\appbrk{}ना\appbrk{}ति\appbrk{}ब\appbrk{}ला\appbrk{}गु\appbrk{}डू\appbrk{}च्या\appbrk{}दित्य्}}\Afootnote{H; \skt{बला\appbrk{}क\appbrk{}ल्पे\appbrk{}ना\appbrk{}ति\appbrk{}ब\appbrk{}ला\appbrk{}गु\appbrk{}डू\appbrk{}च्या\appbrk{}दित्}\ A}}}{\lemma{\skt{अवल्लिषु}}\Afootnote{H; \skt{यपर्णिषु}\ A}}\skt{~। }\pend
\pstart\versepar \skt{}\edtext{\skt{शैरीयके}}{\lemma{\skt{शैरीयके}}\Afootnote{H; \skt{सैरेयके}\ A}}\skt{ }\edtext{\skt{वीरतरे}}{\lemma{\skt{वीरतरे}}\Afootnote{H; \skt{वीरतरौ}\ A}}\skt{ शता\-व\-र्या\-न्त्रि\-क\-ण्टके~॥}\pend

\pstart \ledsidenote{\passref{4.15.45}}\skt{तैलानि }\edtext{\edtext{\skt{कुर्या\-त्म\-धु\-क\-प्र\-सा\-र\-ण्योश्च}}{\lemma{\skt{कुर्यात्}}\Afootnote{H; \skt{मधुके कुर्यात्}\ A}}}{\lemma{\skt{मधु\appbrk{}क\appbrk{}प्र\appbrk{}सा\appbrk{}र\appbrk{}ण्योश्च}}\Afootnote{H; \skt{प्रसारिण्यां च}\ A}}\skt{ }\edtext{\skt{बुद्धिमान्}}{\lemma{\skt{बुद्धिमान्}}\Afootnote{H; \skt{बुद्धिमान्~। नीलोत्पलं वरीमूलं गव्ये क्षीरे विपाचयेत्}\ A}}\skt{~।}\pend

\pstart\versepar \ledsidenote{\passref{4.15.47}}\skt{सर्व्वेषामेव }\edtext{\edtext{\skt{चैते\-षा\-मु\-प\-यो\-ग\-ञ्चि\-कि\-त्सकः}}{\lemma{\skt{चैतेषाम्}}\Afootnote{H; \skt{जानीया}\ A}}}{\lemma{\skt{उपयोगञ्}}\Afootnote{H; \skt{दुपयोगं}\ A}}\skt{~। }\pend
\pstart\versepar \skt{}\edtext{\skt{बला\-तै\-ल\-व\-दा\-चष्टे}}{\lemma{\skt{बला\appbrk{}तै\appbrk{}ल\appbrk{}व\appbrk{}दा\appbrk{}चष्टे}}\Afootnote{H; \skt{बला\appbrk{}तै\appbrk{}ल\appbrk{}व\appbrk{}दे\appbrk{}तेषां}\ A}}\skt{ }\edtext{\skt{गुणांश्चैवा}}{\lemma{\skt{चैवा}}\Afootnote{H; \skt{चैव}\ A}}\skt{ }\edtext{\skt{विशेषत}}{\lemma{\skt{विशेषत}}\Afootnote{H; \skt{विशेषतः}\ A}}\skt{ }\edtext{\skt{इति}}{\lemma{\skt{इति}}\Afootnote{H; \emph{om.}\ A}}\skt{~॥}\pend

\pstart \ledsidenote{\passref{4.15.trailer}}\edtext{\skt{॥ चिकित्सअ}}{\lemma{\skt{॥ चिकित्सअ}}\Afootnote{\skt{इति सुश्रु\appbrk{}त\appbrk{}सं\appbrk{}हि\appbrk{}तायां चिकित्सास्थाने मूढ\appbrk{}ग\appbrk{}र्भ\appbrk{}चि\appbrk{}कि\appbrk{}त्सितं नाम पञ्चदशोऽध्यायः}\ A; \skt{चिकित्सा १५}\ H}}\skt{ १५~॥}\pend

\endnumbering
\beginnumbering
\eledsection*{Adhyāya 16}
\pstart \ledsidenote{\passref{4.16.1}}\skt{अथातो }\edtext{\skt{विद्र\-धी\-ना\-ञ्चि\-कि\-त्सि\-त\-म्व्या\-ख्या\-स्यामः}}{\lemma{\skt{विद्रधीनाञ्}}\Afootnote{H; \skt{विद्रधीतां}\ A}}\skt{~॥}\pend

\pstart\versepar \ledsidenote{\passref{4.16.3}}\skt{उक्ता }\edtext{\skt{विद्रधयो ये}}{\lemma{\skt{विद्रधयो ये}}\Afootnote{H; \skt{विद्रधयः}\ A}}\skt{ }\edtext{\skt{षट्तेष्वसाध्यस्तु}}{\lemma{\skt{षट्तेष्व्}}\Afootnote{H; \skt{षड्ये तेष्व्}\ A}}\skt{ सर्व्वजः~। }\pend
\pstart\versepar \skt{}\edtext{\edtext{\skt{शेषेष्वामेषु}}{\lemma{\skt{शेषेष्व्}}\Afootnote{H; \skt{शेषेष्}\ A}}}{\lemma{\skt{आमेषु}}\Afootnote{H; \skt{वामेषु}\ A}}\skt{ कर्त्तव्या }\edtext{\skt{त्वरता}}{\lemma{\skt{त्वरता}}\Afootnote{H; \skt{त्वरितं}\ A}}\skt{ शोफवत्क्रिया~॥}\pend

\pstart\versepar \ledsidenote{\passref{4.16.4}}\skt{वात\-घ्न\-मू\-ल\-क\-ल्कैस्तु }\edtext{\skt{वसा\-तै\-ल\-घृ\-ता\-युतैः}}{\lemma{\skt{वसा\appbrk{}तै\appbrk{}ल\appbrk{}घृ\appbrk{}ता\appbrk{}युतैः}}\Afootnote{H; \skt{घृत\appbrk{}तै\appbrk{}ल\appbrk{}व\appbrk{}सा\appbrk{}युतैः}\ A}}\skt{~। }\pend
\pstart\versepar \skt{सुखोष्णो वहलो लेपः प्रयोज्यो वातविद्रधौ~॥}\pend

\pstart\versepar \ledsidenote{\passref{4.16.5}}\skt{सानू\-पौ\-द\-क\-मा\-न्सस्तु }\edtext{\skt{काकोल्यादिश्च}}{\lemma{\skt{काकोल्यादिश्च}}\Afootnote{H; \skt{काकोल्यादिः स}\ A}}\skt{ }\edtext{\skt{योगणः}}{\lemma{\skt{योगणः}}\Afootnote{H; \skt{तर्पणः}\ A}}\skt{~। }\pend
\pstart\versepar \skt{स्नेहाम्लसिद्धो लवणः प्रयो\-ज्य\-श्चो\-प\-ना\-हने~॥}\pend

\pstart\versepar \ledsidenote{\passref{4.16.6}}\skt{वेशवारैः }\edtext{\skt{सकृसरैः}}{\lemma{\skt{सकृसरैः}}\Afootnote{H; \skt{सकृशरैः}\ A}}\skt{ पयोभिः पायसैस्तथा~। }\pend
\pstart\versepar \skt{}\edtext{\skt{स्वेदयेच्चापि}}{\lemma{\skt{चापि}}\Afootnote{H; \emph{om.}\ A}}\skt{ }\edtext{\skt{सत\-त\-मा\-ह\-रे\-च्चापि}}{\lemma{\skt{आहरेच्}}\Afootnote{H; \skt{चापि निर्हरेच्}\ A}}\skt{ शोणितं~॥}\pend

\pstart\versepar \ledsidenote{\passref{4.16.7}}\skt{स चेदे\-व\-मु\-प\-क्रान्तः }\edtext{\skt{पाकाय}}{\lemma{\skt{पाकाय}}\Afootnote{H; \skt{पाकाया}\ A}}\skt{ }\edtext{\skt{विमुखो}}{\lemma{\skt{विमुखो}}\Afootnote{H; \skt{भिमुखो}\ A}}\skt{ यदि~। }\pend
\pstart\versepar \skt{तम्पाचयित्वा शस्त्रेण }\edtext{\edtext{\skt{भिद्याद्भित्वा}}{\lemma{\skt{भिद्याद्}}\Afootnote{H; \skt{भिन्द्याद्}\ A}}\skt{ च}}{\lemma{\skt{भित्वा च}}\Afootnote{H; \skt{भिन्नं च शोधयेत्}\ A}}\skt{ }\edtext{\skt{बुद्धिमान्}}{\lemma{\skt{बुद्धिमान्}}\Afootnote{H; \emph{om.}\ A}}\skt{~॥}\pend

\pstart\versepar \ledsidenote{\passref{4.16.8}}\edtext{\skt{पञ्च\-मू\-ली\-क\-षा\-येण}}{\lemma{\skt{पञ्च\appbrk{}मू\appbrk{}ली\appbrk{}क\appbrk{}षा\appbrk{}येण}}\Afootnote{H; \skt{पञ्च\appbrk{}मू\appbrk{}ल\appbrk{}क\appbrk{}षा\appbrk{}येण}\ A}}\skt{ प्रक्षाल्य लवणोत्तरैः~। }\pend
\pstart\versepar \skt{}\edtext{\edtext{\skt{तिलै\-र्ब्भ\-द्रा\-दि\-म\-धु\-क\-सं\-युतैः}}{\lemma{\skt{तिलैर्}}\Afootnote{H; \skt{तैलैर्}\ A}}}{\lemma{\skt{ब्भद्रा\appbrk{}दि\appbrk{}म\appbrk{}धु\appbrk{}क\appbrk{}सं\appbrk{}युतैः}}\Afootnote{H; \skt{भद्रा\appbrk{}दि\appbrk{}म\appbrk{}धु\appbrk{}क\appbrk{}सं\appbrk{}युक्तैः}\ A}}\skt{ }\edtext{\skt{पूर्य}}{\lemma{\skt{पूर्य}}\Afootnote{H; \skt{प्रति}\ A}}\skt{ }\edtext{\skt{वेष्टयेत्}}{\lemma{\skt{वेष्टयेत्}}\Afootnote{H; \skt{पूरयेत्}\ A}}\skt{~॥}\pend

\pstart\versepar \ledsidenote{\passref{4.16.9}}\edtext{\skt{वैरे\-च\-नि\-क\-यु\-क्तेन}}{\lemma{\skt{वैरे\appbrk{}च\appbrk{}नि\appbrk{}क\appbrk{}यु\appbrk{}क्तेन}}\Afootnote{H; \skt{वैरे\appbrk{}च\appbrk{}नि\appbrk{}क\appbrk{}यु\appbrk{}क्तेन}\ A}}\skt{ }\edtext{\skt{त्रिवृतेन}}{\lemma{\skt{त्रिवृतेन}}\Afootnote{H; \skt{त्रैवृतेन}\ A}}\skt{ विशोध्य च~। }\pend
\pstart\versepar \skt{पृथ\-क्प\-र्ण्ण्या\-दि\-सि\-द्धेन }\edtext{\skt{त्रिवृते}}{\lemma{\skt{त्रिवृते}}\Afootnote{H; \skt{त्रैवृते}\ A}}\skt{ }\edtext{\skt{नैव}}{\lemma{\skt{नैव}}\Afootnote{H; \skt{न च}\ A}}\skt{ रोपयेत्~॥}\pend

\pstart\versepar \ledsidenote{\passref{4.16.10}}\skt{पैत्तिकं }\edtext{\skt{शर्क्क\-रा\-ला\-ज\-म\-धुकैः}}{\lemma{\skt{शर्क्क\appbrk{}रा\appbrk{}ला\appbrk{}ज\appbrk{}म\appbrk{}धुकैः}}\Afootnote{H; \skt{शर्क\appbrk{}रा\appbrk{}ला\appbrk{}जा\appbrk{}म\appbrk{}धुकैः}\ A}}\skt{ }\edtext{\skt{शारिवायुतैः}}{\lemma{\skt{शारिवायुतैः}}\Afootnote{H; \skt{सारिवायुतैः}\ A}}\skt{~। }\pend
\pstart\versepar \skt{प्रदि\-ह्या\-त्क्षी\-र\-पि\-ष्टैर्व्वा पय\-स्यो\-शी\-र\-च\-न्दनैः~॥}\pend

\pstart\versepar \ledsidenote{\passref{4.16.11}}\skt{पाक्यैः शीतकषायैर्व्वा क्षीरै\-रि\-क्षु\-र\-सै\-स्तथा~। }\pend
\pstart\versepar \skt{}\edtext{\skt{जीव\-नी\-य\-घृ\-तै\-र्व्वापि}}{\lemma{\skt{व्वापि}}\Afootnote{H; \skt{वाऽपि}\ A}}\skt{ सेच\-ये\-च्छ\-र्क्क\-रा\-युतैः~॥}\pend

\pstart\versepar \ledsidenote{\passref{4.16.12}}\edtext{\skt{तृवृ\-द्ध\-री\-त\-की\-नाञ्च}}{\lemma{\skt{तृवृ\appbrk{}द्ध\appbrk{}री\appbrk{}त\appbrk{}की\appbrk{}नाञ्च}}\Afootnote{H; \skt{त्रिवृ\appbrk{}द्ध\appbrk{}री\appbrk{}त\appbrk{}कीनां च}\ A}}\skt{ }\edtext{\skt{चूर्ण्णम्मधुयुतं}}{\lemma{\skt{मधुयुतं}}\Afootnote{H; \skt{लिह्यान्मधुद्}\ A}}\skt{ }\edtext{\skt{लिहेत्}}{\lemma{\skt{लिहेत्}}\Afootnote{H; \skt{रवम्}\ A}}\skt{~। }\pend
\pstart\versepar \skt{}\edtext{\edtext{\skt{जलौ\-का\-भि\-र्ह\-रे\-च्चा\-सृ\-क्प\-क्व\-म्वा\-पाट्य}}{\lemma{\skt{जलौकाभिर्}}\Afootnote{H; \skt{जलौकोभिर्}\ A}}}{\lemma{\skt{वापाट्य}}\Afootnote{H; \skt{चापठ्य}\ A}}\skt{ बुद्धिमान्~॥}\pend

\pstart\versepar \ledsidenote{\passref{4.16.13}}\skt{क्षीर\-वृ\-क्ष\-क\-षा\-येण }\edtext{\skt{प्रक्षा\-ल्यो\-द\-क\-जेन वा}}{\lemma{\skt{प्रक्षा\appbrk{}ल्यो\appbrk{}द\appbrk{}क\appbrk{}जेन वा}}\Afootnote{H; \skt{प्रक्षा\appbrk{}ल्यौ\appbrk{}द\appbrk{}क\appbrk{}जेन वा}\ A}}\skt{~। }\pend
\pstart\versepar \skt{तिलैः सयष्टीमधुकैः }\edtext{\skt{सामृतैः}}{\lemma{\skt{सामृतैः}}\Afootnote{H; \skt{सक्षौद्रैः}\ A}}\skt{ सर्प्पिषायुतैः~॥}\pend

\pstart\versepar \ledsidenote{\passref{4.16.14}}\edtext{\skt{प्रतिग्राह्य}}{\lemma{\skt{प्रतिग्राह्य}}\Afootnote{H; \skt{उपदिह्य}\ A}}\skt{ प्रतनुना }\edtext{\skt{सर्प्पि\-षा\-वे\-ष्ट\-ये\-द्व्रणं}}{\lemma{\skt{सर्प्पि\appbrk{}षा\appbrk{}वे\appbrk{}ष्ट\appbrk{}येद्}}\Afootnote{H; \skt{वाससा वेष्टयेद्}\ A}}\skt{~। }\pend
\pstart\versepar \skt{प्रपौ\-ण्ड\-री\-क\-म\-ञ्जिष्ठा मधु\-को\-शी\-र\-प\-द्मकैः~॥}\pend

\pstart\versepar \ledsidenote{\passref{4.16.15}}\skt{सहरिद्रैः कृतं सर्प्पिः सक्षीरं व्रणरोपणं~। }\pend
\pstart\versepar \skt{}\edtext{\skt{क्षीर\-शु\-क्ला\-पृ\-थ\-क्प\-र्ण्णी\-स\-म\-ङ्गा\-लो\-ध्र\-च\-न्दनैः}}{\lemma{\skt{क्षीर\appbrk{}शु\appbrk{}क्ला\appbrk{}पृ\appbrk{}थ\appbrk{}क्प\appbrk{}र्ण्णी\appbrk{}स\appbrk{}म\appbrk{}ङ्गा\appbrk{}लो\appbrk{}ध्र\appbrk{}च\appbrk{}न्दनैः}}\Afootnote{H; \skt{क्षीर\appbrk{}शु\appbrk{}क्ला\appbrk{}पृ\appbrk{}थ\appbrk{}क्प\appbrk{}र्णी\appbrk{}स\appbrk{}म\appbrk{}ङ्गा\appbrk{}रो\appbrk{}ध्र\appbrk{}च\appbrk{}न्दनैः}\ A}}\skt{~॥}\pend

\pstart\versepar \ledsidenote{\passref{4.16.16}}\skt{न्यग्रो\-धा\-दि\-प्र\-वा\-लेषु तेषां त्वक्ष्वथवाकृतं~। }\pend
\pstart\versepar \skt{नक्तमालस्य पत्राणि }\edtext{\edtext{\skt{पटो\-ला\-रि\-ष्ट\-यो\-स्तथा}}{\lemma{\skt{पटो\appbrk{}ला\appbrk{}रि\appbrk{}ष्ट\appbrk{}योस्}}\Afootnote{H; \skt{तरुणानि}\ A}}}{\lemma{\skt{तथा}}\Afootnote{H; \skt{फलानि च}\ A}}\skt{~॥}\pend

\pstart \ledsidenote{\passref{4.16.17}}\edtext{\skt{द्वे}}{\lemma{\skt{द्वे}}\Afootnote{H; \skt{सुमनायाश्च पत्राणि पटो\appbrk{}ला\appbrk{}रि\appbrk{}ष्ट\appbrk{}यो\appbrk{}स्तथा~। द्वे}\ A}}\skt{ हरिद्रे मधूच्छिष्टं मधु\-क\-न्ति\-क्त\-रो\-हिणी~।}\pend

\pstart \ledsidenote{\passref{4.16.18}}\edtext{\skt{प्रियं\-गु\-कु\-श\-मू\-लञ्च}}{\lemma{\skt{प्रियं\appbrk{}गु\appbrk{}कु\appbrk{}श\appbrk{}मू\appbrk{}लञ्च}}\Afootnote{H; \skt{प्रियङ्गुः कुशमूलं च}\ A}}\skt{ निचुलस्य }\edtext{\skt{त्वगेव च}}{\lemma{\skt{त्वगेव च}}\Afootnote{H; \skt{त्वगेव च~। मञ्जि\appbrk{}ष्ठा\appbrk{}च\appbrk{}न्द\appbrk{}नो\appbrk{}शी\appbrk{}र\appbrk{}मु\appbrk{}त्पलं सारिवे त्रिवृत्}\ A}}\skt{~॥}\pend

\pstart\versepar \ledsidenote{\passref{4.16.19}}\edtext{\skt{एषां}}{\lemma{\skt{एषां}}\Afootnote{H; \skt{एतेषां}\ A}}\skt{ }\edtext{\edtext{\skt{कर्ष\-स\-मै\-र्ब्भा\-गै\-र्घृ\-त\-प्रस्थं}}{\lemma{\skt{कर्षसमैर्}}\Afootnote{H; \skt{कार्षि}\ A}}}{\lemma{\skt{ब्भागैर्}}\Afootnote{H; \skt{कैर्भागैर्}\ A}}\skt{ विपाचयेत्~। }\pend
\pstart\versepar \skt{दुष्ट\-व्र\-ण\-प्र\-श\-मनं }\edtext{\skt{तथा}}{\lemma{\skt{तथा}}\Afootnote{H; \emph{om.}\ A}}\skt{ }\edtext{\skt{नाडीविनाशनं}}{\lemma{\skt{नाडीविनाशनं}}\Afootnote{H; \skt{नाडी\appbrk{}व्र\appbrk{}ण\appbrk{}वि\appbrk{}शो\appbrk{}ध\appbrk{}नम्}\ A}}\skt{~॥}\pend

\pstart \ledsidenote{\passref{4.16.20}}\skt{सद्य\-श्च्छि\-न्न\-व्र\-णा\-नाञ्च करञ्जाद्यमिदं }\edtext{\skt{शुभं}}{\lemma{\skt{शुभं}}\Afootnote{H; \skt{शुभम्~। दुष्टव्रणाश्च ये केचिद्ये चोत्सृष्टक्रिया व्रणाः}\ A}}\skt{~।}\pend

\pstart \ledsidenote{\passref{4.16.22}}\edtext{\skt{इष्ट\-का\-सि\-क\-ता\-लो\-ह\-गो\-श\-कृ\-त्तु\-ष\-पा\-ङ्शुभिः}}{\lemma{\skt{इष्ट\appbrk{}का\appbrk{}सि\appbrk{}क\appbrk{}ता\appbrk{}लो\appbrk{}ह\appbrk{}गो\appbrk{}श\appbrk{}कृ\appbrk{}त्तु\appbrk{}ष\appbrk{}पा\appbrk{}ङ्शुभिः}}\Afootnote{H; \skt{करञ्जाद्येन हविषा प्रशाम्यन्ति न संशयः~। इष्ट\appbrk{}का\appbrk{}सि\appbrk{}क\appbrk{}ता\appbrk{}लो\appbrk{}ष्ट\appbrk{}गो\appbrk{}श\appbrk{}कृ\appbrk{}त्तु\appbrk{}ष\appbrk{}पां\appbrk{}शुभिः}\ A}}\skt{~॥}\pend

\pstart\versepar \ledsidenote{\passref{4.16.23}}\skt{मूत्रैरुष्णैश्च सततं }\edtext{\skt{स्वेद\-ये\-च्छ्ले\-ष्म\-वि\-द्रधिं}}{\lemma{\skt{छ्लेष्मविद्रधिं}}\Afootnote{H; \skt{छ्लष्मविद्रधिम्}\ A}}\skt{~। }\pend
\pstart\versepar \skt{कषा\-य\-पा\-नै\-र्व्व\-म\-नै\-रा\-ले\-पै\-रु\-प\-ना\-हनैः~॥}\pend

\pstart\versepar \ledsidenote{\passref{4.16.24}}\edtext{\skt{हरे\-द्दो\-षा\-न\-भी\-क्ष्णञ्च}}{\lemma{\skt{दोषानभीक्ष्णञ्च}}\Afootnote{H; \skt{दोषानभीक्ष्णं चाप्य्}\ A}}\skt{ }\edtext{\skt{तथै\-वा\-सृ\-ग\-ला\-वुभिः}}{\lemma{\skt{तथै\appbrk{}वा\appbrk{}सृ\appbrk{}ग\appbrk{}ला\appbrk{}वुभिः}}\Afootnote{H; \skt{अला\appbrk{}ब्वाऽ\appbrk{}सृ\appbrk{}क्त\appbrk{}थैव च}\ A}}\skt{~। }\pend
\pstart\versepar \skt{आर\-ग्व\-ध\-क\-षा\-येण }\edtext{\skt{पक्वम्वा}}{\lemma{\skt{पक्वम्वा}}\Afootnote{H; \skt{पक्वं चा}\ A}}\skt{ }\edtext{\skt{पाट्य}}{\lemma{\skt{पाट्य}}\Afootnote{H; \skt{पाठ्य}\ A}}\skt{ धावयेत्~॥}\pend

\pstart\versepar \ledsidenote{\passref{4.16.25}}\edtext{\skt{हरि\-द्रा\-त्रि\-फ\-ला\-श\-क्तु\-ति\-लै\-र्म्म\-धु\-स\-मा\-युतैः}}{\lemma{\skt{हरि\appbrk{}द्रा\appbrk{}त्रि\appbrk{}फ\appbrk{}ला\appbrk{}श\appbrk{}क्तु\appbrk{}ति\appbrk{}लैर्}}\Afootnote{H; \skt{हरि\appbrk{}द्रा\appbrk{}त्रि\appbrk{}वृ\appbrk{}ता\appbrk{}श\appbrk{}क्तु\appbrk{}ति\appbrk{}लैर्}\ A}}\skt{~। }\pend
\pstart\versepar \skt{पूरयित्वा व्रणं }\edtext{\skt{सम्य\-क्व\-ध्नी\-या\-त्की\-र्त्तितं}}{\lemma{\skt{सम्यक्}}\Afootnote{H; \skt{सम्यग्}\ A}}\skt{ यथा~॥}\pend

\pstart\versepar \ledsidenote{\passref{4.16.26}}\skt{ततः कुलत्थिकादन्ती }\edtext{\skt{तृवृ\-च्छ्या\-मा\-र्क्क\-ति\-ल्वकैः}}{\lemma{\skt{तृवृ\appbrk{}च्छ्या\appbrk{}मा\appbrk{}र्क्क\appbrk{}ति\appbrk{}ल्वकैः}}\Afootnote{H; \skt{त्रिवृ\appbrk{}च्छ्या\appbrk{}मा\appbrk{}र्क\appbrk{}ति\appbrk{}ल्वकैः}\ A}}\skt{~। }\pend
\pstart\versepar \skt{कुर्यात्तैलं सगोमूत्रं हितन्तत्र ससैन्धवं~॥}\pend

\pstart\versepar \ledsidenote{\passref{4.16.27}}\edtext{\skt{पित्त\-वि\-द्र\-धि\-व\-त्सर्व्वा}}{\lemma{\skt{सर्व्वा}}\Afootnote{H; \skt{सर्वाः}\ A}}\skt{ क्रिया निरवशेषतः~। }\pend
\pstart\versepar \skt{विद्रध्योः कुशलः कुर्या\-द्र\-क्ता\-ग\-न्तु\-नि\-मि\-त्तयोः~॥}\pend

\pstart\versepar \ledsidenote{\passref{4.16.28}}\edtext{\edtext{\edtext{\skt{वरु\-णा\-दि\-र्ग\-ण\-क्वा\-थ\-म\-प\-क्वेऽ\-भ्य\-न्त\-रो\-त्थिते}}{\lemma{\skt{वरुणादिर्}}\Afootnote{H; \skt{वरुणादि}\ A}}}{\lemma{\skt{गणक्वाथम्}}\Afootnote{H; \skt{गणक्वाथ}\ A}}}{\lemma{\skt{अपक्वे}}\Afootnote{H; \skt{मपक्वे}\ A}}\skt{~। }\pend
\pstart\versepar \skt{}\edtext{\skt{उष\-का\-दि\-प्र\-ती\-वापं}}{\lemma{\skt{उष\appbrk{}का\appbrk{}दि\appbrk{}प्र\appbrk{}ती\appbrk{}वापं}}\Afootnote{H; \skt{ऊष\appbrk{}का\appbrk{}दि\appbrk{}प्र\appbrk{}ती\appbrk{}वापं}\ A}}\skt{ }\edtext{\skt{पिवेच्छंकरणा}}{\lemma{\skt{पिवेच्छंकरणा}}\Afootnote{H; \skt{पिबेत्सुखकरं नरः}\ A}}\skt{ }\edtext{\skt{यदि}}{\lemma{\skt{यदि}}\Afootnote{H; \emph{om.}\ A}}\skt{~॥}\pend

\pstart\versepar \ledsidenote{\passref{4.16.29}}\skt{अनयोर्व्वर्ग्गयोः }\edtext{\skt{सर्प्पिः}}{\lemma{\skt{सर्प्पिः}}\Afootnote{H; \skt{सिद्धं सर्}\ A}}\skt{ }\edtext{\skt{सिद्धं}}{\lemma{\skt{सिद्धं}}\Afootnote{H; \skt{पि}\ A}}\skt{ }\edtext{\skt{वैरेचनेषु च}}{\lemma{\skt{वैरेचनेषु च}}\Afootnote{H; \skt{र्वैरेचनेन च}\ A}}\skt{~॥ }\pend
\pstart\versepar \skt{}\edtext{\skt{विद्रधिं}}{\lemma{\skt{विद्रधिं}}\Afootnote{H; \skt{अचिराद्विद्रधिं}\ A}}\skt{ }\edtext{\skt{हन्ति न}}{\lemma{\skt{हन्ति न}}\Afootnote{H; \skt{हन्ति}\ A}}\skt{ }\edtext{\skt{चिरात्प्रातः}}{\lemma{\skt{चिरात्}}\Afootnote{H; \emph{om.}\ A}}\skt{ प्रातर्न्निषेवितं~।}\pend

\pstart\versepar \ledsidenote{\passref{4.16.30}}\skt{एभिरेव गणैश्चापि संसिद्धं }\edtext{\skt{स्नेहयोजितं}}{\lemma{\skt{स्नेहयोजितं}}\Afootnote{H; \skt{स्नेहसंयुतम्}\ A}}\skt{~॥ }\pend
\pstart\versepar \skt{}\edtext{\edtext{\skt{कार्यमास्थापनं}}{\lemma{\skt{कार्यम्}}\Afootnote{H; \skt{कार्य}\ A}}}{\lemma{\skt{आस्थापनं}}\Afootnote{H; \skt{मास्थापनं}\ A}}\skt{ }\edtext{\skt{क्षिप्रं स च}}{\lemma{\skt{क्षिप्रं स च}}\Afootnote{H; \skt{क्षिप्रं तथै}\ A}}\skt{ }\edtext{\skt{स्नेहो}}{\lemma{\skt{स्नेहो}}\Afootnote{H; \skt{वाप्य}\ A}}\skt{ नुवासनं~।}\pend

\pstart\versepar \ledsidenote{\passref{4.16.31}}\skt{पानालेपन भोज्येषु }\edtext{\skt{मधु\-सि\-ग्रु\-द्रु\-मोऽपि}}{\lemma{\skt{मधुसिग्रुद्रुमो}}\Afootnote{H; \skt{मधुशिग्रुद्रुमो}\ A}}\skt{ वा~॥ }\pend
\pstart\versepar \skt{}\edtext{\skt{दत्तावायो}}{\lemma{\skt{दत्तावायो}}\Afootnote{H; \skt{दत्तावापो}\ A}}\skt{ }\edtext{\edtext{\skt{यथा\-दो\-ष\-म\-पक्वं}}{\lemma{\skt{यथादोषम्}}\Afootnote{H; \skt{यथादोष}\ A}}}{\lemma{\skt{अपक्वं}}\Afootnote{H; \skt{मपक्वं}\ A}}\skt{ हन्ति विद्रधिं~।}\pend

\pstart\versepar \ledsidenote{\passref{4.16.32}}\edtext{\skt{तोय\-धा\-न्या\-न्या\-म्ल\-मू\-त्रैस्तु}}{\lemma{\skt{तोय\appbrk{}धा\appbrk{}न्या\appbrk{}न्या\appbrk{}म्ल\appbrk{}मू\appbrk{}त्रैस्तु}}\Afootnote{H; \skt{तोय\appbrk{}धा\appbrk{}न्या\appbrk{}म्ल\appbrk{}मू\appbrk{}त्रैस्तु}\ A}}\skt{ पेयो }\edtext{\skt{वापि}}{\lemma{\skt{वापि}}\Afootnote{H; \skt{वाऽपि}\ A}}\skt{ सुरादिभिः~॥ }\pend
\pstart\versepar \skt{यथा\-दो\-ष\-ग\-ण\-क्वाथैः }\edtext{\skt{पिवेद्वापि}}{\lemma{\skt{वापि}}\Afootnote{H; \skt{वाऽपि}\ A}}\skt{ }\edtext{\skt{शिलाजतुं}}{\lemma{\skt{शिलाजतुं}}\Afootnote{H; \skt{शिलाजतु}\ A}}\skt{~।}\pend

\pstart\versepar \ledsidenote{\passref{4.16.33}}\edtext{\edtext{\skt{प्रधा\-न\-ङ्गु\-ग्गु\-लु\-म्वापि}}{\lemma{\skt{गुग्गुलुम्}}\Afootnote{H; \skt{गुग्गुलं}\ A}}}{\lemma{\skt{वापि}}\Afootnote{H; \skt{चापि}\ A}}\skt{ }\edtext{\edtext{\skt{दारु\-शु\-ण्ठी\-म\-थापि}}{\lemma{\skt{दारुशुण्ठीम्}}\Afootnote{H; \skt{शुण्ठीं च}\ A}}\skt{ वा}}{\lemma{\skt{अथापि वा}}\Afootnote{H; \skt{सुरदारु च}\ A}}\skt{~॥ }\pend
\pstart\versepar \skt{}\edtext{\skt{स्वेदोपनाहः}}{\lemma{\skt{स्वेदोपनाहः}}\Afootnote{H; \skt{स्नेहोपनाहौ}\ A}}\skt{ कुर्याच्च सदा }\edtext{\skt{चाप्य\-नु\-लो\-म\-येत्}}{\lemma{\skt{अनुलोमयेत्}}\Afootnote{H; \skt{अनुलोमनम्}\ A}}\skt{~।}\pend

\pstart\versepar \ledsidenote{\passref{4.16.34}}\edtext{\skt{शिरां}}{\lemma{\skt{शिरां}}\Afootnote{H; \emph{om.}\ A}}\skt{ }\edtext{\edtext{\skt{यथोक्ताम्विध्येत}}{\lemma{\skt{यथोक्ताम्}}\Afootnote{H; \skt{यथोद्दिष्टां सिरां}\ A}}}{\lemma{\skt{विध्येत}}\Afootnote{H; \skt{विध्येत्}\ A}}\skt{ कफजे विद्रधौ भिषक्~॥ }\pend
\pstart\versepar \skt{रक्त\-पि\-त्ता\-नि\-लो\-त्थेषु }\edtext{\skt{केचिद्वाहौ}}{\lemma{\skt{केचिद्}}\Afootnote{H; \skt{केचिद्}\ A}}\skt{ वदन्ति तु~।}\pend

\pstart\versepar \ledsidenote{\passref{4.16.35}}\edtext{\skt{पर्यागतञ्च}}{\lemma{\skt{पर्यागतञ्च}}\Afootnote{H; \skt{पक्वं वा}\ A}}\skt{ }\edtext{\skt{शस्त्रेण}}{\lemma{\skt{शस्त्रेण}}\Afootnote{H; \skt{बहिरुन्नद्धं}\ A}}\skt{ भित्त्वा }\edtext{\skt{व्रण\-व\-दा\-च\-रेत्}}{\lemma{\skt{व्रण\appbrk{}व\appbrk{}दा\appbrk{}च\appbrk{}रेत्}}\Afootnote{H; \skt{व्रण\appbrk{}व\appbrk{}दा\appbrk{}च\appbrk{}रेत्}\ A}}\skt{~॥ }\pend
\pstart\versepar \skt{}\edtext{\edtext{\skt{श्रुतेषूर्द्धमधो}}{\lemma{\skt{श्रुतेषूर्द्धम्}}\Afootnote{H; \skt{स्रुतेषूर्ध्}\ A}}}{\lemma{\skt{अधो}}\Afootnote{H; \skt{वमधो}\ A}}\skt{ }\edtext{\skt{वापि}}{\lemma{\skt{वापि}}\Afootnote{H; \skt{वाऽपि}\ A}}\skt{ }\edtext{\skt{मैरे\-या\-म्ल\-सु\-रा\-सवैः}}{\lemma{\skt{मैरे\appbrk{}या\appbrk{}म्ल\appbrk{}सु\appbrk{}रा\appbrk{}सवैः}}\Afootnote{H; \skt{मैरे\appbrk{}या\appbrk{}म्ल\appbrk{}सु\appbrk{}रा\appbrk{}सवैः}\ A}}\skt{~।}\pend

\pstart\versepar \ledsidenote{\passref{4.16.36}}\skt{पेयो वरुणकादिस्तु }\edtext{\skt{मधु\-सि\-ग्रु\-द्रु\-मोऽपि}}{\lemma{\skt{मधुसिग्रुद्रुमो}}\Afootnote{H; \skt{मधुशिग्रुद्रुमो}\ A}}\skt{ वा~॥ }\pend
\pstart\versepar \skt{}\edtext{\skt{मधुसिग्रोस्तु}}{\lemma{\skt{मधुसिग्रोस्तु}}\Afootnote{H; \skt{शिग्रुमू}\ A}}\skt{ }\edtext{\skt{मूलेषु}}{\lemma{\skt{मूलेषु}}\Afootnote{H; \skt{लजले सिद्धं स}\ A}}\skt{ }\edtext{\skt{सिद्धानां}}{\lemma{\skt{सिद्धानां}}\Afootnote{H; \skt{सिद्धा\appbrk{}र्थ\appbrk{}क\appbrk{}मो\appbrk{}द\appbrk{}नम्}\ A}}\skt{ }\edtext{\skt{सहसावुधैः}}{\lemma{\skt{सहसावुधैः}}\Afootnote{H; \emph{om.}\ A}}\skt{~॥}\pend

\pstart\versepar \ledsidenote{\passref{4.16.37}}\skt{यव\-को\-ल\-कु\-ल\-त्थानां }\edtext{\skt{यूषैर्ब्भुञ्जीत}}{\lemma{\skt{ब्भुञ्जीत}}\Afootnote{H; \skt{भुञ्जीत}\ A}}\skt{ मानवः~॥ }\pend
\pstart\versepar \skt{प्रातः प्रातश्च सेवेत मात्रया तैल्वकं घृतं~।}\pend

\pstart \ledsidenote{\passref{4.16.39}}\skt{पर्यागते }\edtext{\skt{विद्रधो तु}}{\lemma{\skt{विद्रधो तु}}\Afootnote{H; \skt{विद्रधौ तु}\ A}}\skt{ }\edtext{\edtext{\skt{सिद्धिर्न्नैकान्तिकी}}{\lemma{\skt{सिद्धिर्}}\Afootnote{H; \skt{सिद्धि}\ A}}}{\lemma{\skt{न्नैकान्तिकी}}\Afootnote{H; \skt{र्नैकान्तिकी}\ A}}\skt{ }\edtext{\skt{भवेदिति}}{\lemma{\skt{भवेदिति}}\Afootnote{H; \skt{स्मृता~। प्रत्याख्याय तु कुर्वीत मज्जजाते तु विद्रधौ}\ A}}\skt{~॥}\pend

\pstart \ledsidenote{\passref{4.16.trailer}}\edtext{\skt{चिकित्सा}}{\lemma{\skt{चिकित्सा}}\Afootnote{H; \skt{इति सुश्रु\appbrk{}त\appbrk{}सं\appbrk{}हि\appbrk{}तायां चिकित्सास्थाने विद्र\appbrk{}धि\appbrk{}चि\appbrk{}कि\appbrk{}त्सितं नाम षोडशोऽध्यायः}\ A}}\pend

\endnumbering
\beginnumbering
\eledsection*{Adhyāya 17}
\pstart \ledsidenote{\passref{4.17.1}}\skt{अथातो }\edtext{\skt{विस\-र्प्प\-ना\-डी\-स्त\-न\-रो\-गा\-णा\-ञ्चि\-कि\-त्सितं}}{\lemma{\skt{विस\appbrk{}र्प्प\appbrk{}ना\appbrk{}डी\appbrk{}स्त\appbrk{}न\appbrk{}रो\appbrk{}गा\appbrk{}णाञ्}}\Afootnote{H; \skt{विस\appbrk{}र्प\appbrk{}ना\appbrk{}डी\appbrk{}स्त\appbrk{}न\appbrk{}रोग}\ A}}\skt{ व्याख्यास्यामः~॥}\pend

\pstart\versepar \ledsidenote{\passref{4.17.3}}\skt{साध्या }\edtext{\edtext{\skt{विसर्प्पास्त्रय}}{\lemma{\skt{विसर्प्पास्}}\Afootnote{H; \skt{विसर्पा}\ A}}}{\lemma{\skt{त्रय}}\Afootnote{H; \skt{स्त्रय}\ A}}\skt{ आदितो ये }\edtext{\edtext{\skt{तत्स\-न्नि\-पा\-त\-क्ष\-त\-जा\-व\-साध्यौ}}{\lemma{\skt{तत्}}\Afootnote{H; \skt{न}\ A}}}{\lemma{\skt{सन्नि\appbrk{}पा\appbrk{}त\appbrk{}क्ष\appbrk{}त\appbrk{}जा\appbrk{}व\appbrk{}साध्यौ}}\Afootnote{H; \skt{सन्नि\appbrk{}पा\appbrk{}त\appbrk{}क्ष\appbrk{}तजौ हि साध्यौ}\ A}}\skt{~। }\pend
\pstart\versepar \skt{साध्येषु }\edtext{\skt{पथ्यैर्हि}}{\lemma{\skt{पथ्यैर्हि}}\Afootnote{H; \skt{तत्पथ्य}\ A}}\skt{ }\edtext{\skt{गणै\-र्व्वि\-द\-ध्या\-द्घृ\-ता\-व\-से\-कोश्च}}{\lemma{\skt{घृतावसेकोश्च}}\Afootnote{H; \skt{घृतानि सेकांश्च}\ A}}\skt{ }\edtext{\skt{तथा}}{\lemma{\skt{तथा}}\Afootnote{H; \skt{तथो}\ A}}\skt{ }\edtext{\skt{प्रदेहान्}}{\lemma{\skt{प्रदेहान्}}\Afootnote{H; \skt{पदेहान्}\ A}}\skt{~॥}\pend

\pstart\versepar \ledsidenote{\passref{4.17.4}}\edtext{\skt{कुष्ठं}}{\lemma{\skt{कुष्ठं}}\Afootnote{H; \skt{मुस्ता}\ A}}\skt{ }\edtext{\skt{शता\-ह्वा\-सु\-र\-दा\-रु\-मु\-स्ता\-वा\-रा\-हि\-कु\-स्तुं\-वु\-रु\-कृ\-ष्ण\-गन्धा}}{\lemma{\skt{शता\appbrk{}ह्वा\appbrk{}सु\appbrk{}र\appbrk{}दा\appbrk{}रु\appbrk{}मु\appbrk{}स्ता\appbrk{}वा\appbrk{}रा\appbrk{}हि\appbrk{}कु\appbrk{}स्तुं\appbrk{}वु\appbrk{}रु\appbrk{}कृ\appbrk{}ष्ण\appbrk{}गन्धा}}\Afootnote{H; \skt{शता\appbrk{}ह्वा\appbrk{}सु\appbrk{}र\appbrk{}दा\appbrk{}रु\appbrk{}कु\appbrk{}ष्ठ\appbrk{}वा\appbrk{}रा\appbrk{}हि\appbrk{}कु\appbrk{}स्तु\appbrk{}म्बु\appbrk{}रु\appbrk{}कृ\appbrk{}ष्ण\appbrk{}गन्धाः}\ A}}\skt{~। }\pend
\pstart\versepar \skt{}\edtext{\edtext{\skt{वाते\-र्क्क\-वं\-शा\-र्त्त\-ग\-णास्तु}}{\lemma{\skt{वाते\appbrk{}र्क्क\appbrk{}वं\appbrk{}शा\appbrk{}र्त्त\appbrk{}ग\appbrk{}णास्}}\Afootnote{H; \skt{वातात्मको चोष्णगणाः}\ A}}}{\lemma{\skt{तु}}\Afootnote{H; \skt{प्र}\ A}}\skt{ }\edtext{\skt{योज्या}}{\lemma{\skt{योज्या}}\Afootnote{H; \skt{योज्याः}\ A}}\skt{ सेकेषु लेपेषु तथा }\edtext{\skt{घृतेषु}}{\lemma{\skt{घृतेषु}}\Afootnote{H; \skt{सद\appbrk{}न्ती\appbrk{}न्सु\appbrk{}रा\appbrk{}ष्ट्र\appbrk{}जा\appbrk{}सै\appbrk{}न्ध\appbrk{}व\appbrk{}सं\appbrk{}प्र\appbrk{}यु\appbrk{}क्तान्~। प्रक्षालने चापि कर\appbrk{}ञ्ज\appbrk{}नि\appbrk{}म्ब\appbrk{}जा\appbrk{}त्य\appbrk{}क्ष\appbrk{}पी\appbrk{}लु\appbrk{}स्व\appbrk{}रसाः प्रयोज्याः}\ A}}\skt{~॥}\pend

\pstart\versepar \ledsidenote{\passref{4.17.5}}\skt{यत्पञ्चमूलं खलु }\edtext{\edtext{\edtext{\skt{कण्ट\-का\-ख्य\-म\-ल्प\-म्म\-ह\-च्चा\-प्य\-थ\-व\-ल्लि\-जाञ्च}}{\lemma{\skt{कण्टकाख्यम्}}\Afootnote{H; \skt{कण्टकाख्य}\ A}}}{\lemma{\skt{अल्पम्}}\Afootnote{H; \skt{मल्पं}\ A}}}{\lemma{\skt{अथवल्लिजाञ्च}}\Afootnote{H; \skt{अथ वल्लिजं च}\ A}}\skt{~। }\pend
\pstart\versepar \skt{तच्चो\-प\-यो\-ज्य\-म्भि\-षजा }\edtext{\skt{प्रदेहैः}}{\lemma{\skt{प्रदेहैः}}\Afootnote{H; \skt{पर्देहे}\ A}}\skt{ सेके घृते चापि तथैव तैले~॥}\pend

\pstart\versepar \ledsidenote{\passref{4.17.6}}\skt{कसे\-रु\-शृ\-ङ्गा\-ट\-क\-प\-द्म\-गुन्द्राः सशैवलाः }\edtext{\edtext{\skt{सोत्प\-ल\-क\-र्द्द\-माश्च}}{\lemma{\skt{सोत्प\appbrk{}ल\appbrk{}क\appbrk{}र्द्द\appbrk{}माश्}}\Afootnote{H; \skt{सोत्पलकर्दमा}\ A}}}{\lemma{\skt{च}}\Afootnote{H; \skt{श्च}\ A}}\skt{~। }\pend
\pstart\versepar \skt{}\edtext{\skt{वस्त्रान्तरा}}{\lemma{\skt{वस्त्रान्तरा}}\Afootnote{H; \skt{वस्त्रान्तराः}\ A}}\skt{ पित्तकृते }\edtext{\skt{विस\-र्प्प\-ले\-पा\-वि\-धेयाः}}{\lemma{\skt{विस\appbrk{}र्प्प\appbrk{}ले\appbrk{}पा\appbrk{}वि\appbrk{}धेयाः}}\Afootnote{H; \skt{विसर्पे लेपा विधेयाः}\ A}}\skt{ सघृताः सुशीताः~॥}\pend

\pstart\versepar \ledsidenote{\passref{4.17.7}}\skt{ह्रीवे\-र\-ला\-म\-ज्ज\-क\-च\-न्द\-नानि }\edtext{\skt{श्रोतो\-ज\-मु\-क्ता\-म\-णि\-गै\-रि\-कैश्च}}{\lemma{\skt{श्रोतो\appbrk{}ज\appbrk{}मु\appbrk{}क्ता\appbrk{}म\appbrk{}णि\appbrk{}गै\appbrk{}रि\appbrk{}कैश्च}}\Afootnote{H; \skt{स्रोतो\appbrk{}ज\appbrk{}मु\appbrk{}क्ता\appbrk{}म\appbrk{}णि\appbrk{}गै\appbrk{}रि\appbrk{}काश्च}\ A}}\skt{~। }\pend
\pstart\versepar \skt{}\edtext{\skt{क्षीराभिपिष्टाः}}{\lemma{\skt{क्षीराभिपिष्टाः}}\Afootnote{H; \skt{क्षीरेण पिष्टाः}\ A}}\skt{ सघृताः सुशीता लेपाः }\edtext{\edtext{\skt{प्रयोज्यास्तनवः}}{\lemma{\skt{प्रयोज्यास्}}\Afootnote{H; \skt{प्रयोज्या}\ A}}}{\lemma{\skt{तनवः}}\Afootnote{H; \skt{स्तनवः}\ A}}\skt{ सुखाय~॥}\pend

\pstart\versepar \ledsidenote{\passref{4.17.8}}\edtext{\edtext{\skt{मञ्जी\-ष्ठ\-य\-ष्टी\-म\-धु\-क\-म्प\-य\-स्या\-प्र\-पौ\-ण्ड\-रीकं}}{\lemma{\skt{मञ्जी\appbrk{}ष्ठ\appbrk{}य\appbrk{}ष्टी\appbrk{}म\appbrk{}धु\appbrk{}कम्}}\Afootnote{H; \skt{प्रपौण्डरीकं मधुकं}\ A}}}{\lemma{\skt{पय\appbrk{}स्या\appbrk{}प्र\appbrk{}पौ\appbrk{}ण्ड\appbrk{}रीकं}}\Afootnote{H; \skt{पयस्या मञ्}\ A}}\skt{ }\edtext{\skt{सह}}{\lemma{\skt{सह}}\Afootnote{H; \skt{जिष्ठिका}\ A}}\skt{ }\edtext{\skt{पद्मकेन}}{\lemma{\skt{पद्मकेन}}\Afootnote{H; \skt{पद्मकचन्दने च}\ A}}\skt{~। }\pend
\pstart\versepar \skt{}\edtext{\skt{सौगन्धिकाञ्चेति}}{\lemma{\skt{सौगन्धिकाञ्}}\Afootnote{H; \skt{सुगन्धिका}\ A}}\skt{ }\edtext{\skt{सुखप्रलेपः}}{\lemma{\skt{सुखप्रलेपः}}\Afootnote{H; \skt{सुखाय लेपः}\ A}}\skt{ पैत्ते विसर्प्पे भिषजा }\edtext{\skt{प्रयोज्याः}}{\lemma{\skt{प्रयोज्याः}}\Afootnote{H; \skt{प्रयोज्यः}\ A}}\skt{~॥}\pend

\pstart\versepar \ledsidenote{\passref{4.17.9}}\edtext{\skt{न्यग्रो\-ध\-व\-र्ग्ग\-म्प\-रि\-षे\-च\-नञ्च}}{\lemma{\skt{न्यग्रोधवर्ग्गम्}}\Afootnote{H; \skt{न्यग्रोधवर्गैः}\ A}}\skt{ घृतञ्च कुर्यात्स्वरसेन तस्य~। }\pend
\pstart\versepar \skt{शीतैः पयो भिश्च मधूदकैश्च सश\-र्क्क\-रै\-रि\-क्षु\-र\-सैश्च }\edtext{\skt{सेकां}}{\lemma{\skt{सेकां}}\Afootnote{H; \skt{सेकान्}\ A}}\skt{~॥}\pend

\pstart\versepar \ledsidenote{\passref{4.17.10}}\edtext{\edtext{\skt{घृतञ्च}}{\lemma{\skt{घृतञ्}}\Afootnote{H; \skt{घृत}\ A}}}{\lemma{\skt{च}}\Afootnote{H; \skt{स्य}\ A}}\skt{ }\edtext{\skt{गौरी\-म\-धु\-का\-र\-वि\-न्द\-लो\-ध्रा\-म्वु\-रा\-जा\-द\-न\-जी\-व\-केषु}}{\lemma{\skt{गौरी\appbrk{}म\appbrk{}धु\appbrk{}का\appbrk{}र\appbrk{}वि\appbrk{}न्द\appbrk{}लो\appbrk{}ध्रा\appbrk{}म्वु\appbrk{}रा\appbrk{}जा\appbrk{}द\appbrk{}न\appbrk{}जी\appbrk{}व\appbrk{}केषु}}\Afootnote{H; \skt{गौरी\appbrk{}म\appbrk{}धु\appbrk{}का\appbrk{}र\appbrk{}वि\appbrk{}न्द\appbrk{}रो\appbrk{}ध्रा\appbrk{}म्बु\appbrk{}रा\appbrk{}जा\appbrk{}द\appbrk{}न\appbrk{}गै\appbrk{}रि\appbrk{}केषु}\ A}}\skt{~। }\pend
\pstart\versepar \skt{}\edtext{\skt{तथार्षभे}}{\lemma{\skt{तथार्षभे}}\Afootnote{H; \skt{तथर्षभे}\ A}}\skt{ }\edtext{\skt{पद्म\-क\-शा\-रि\-वायां}}{\lemma{\skt{पद्म\appbrk{}क\appbrk{}शा\appbrk{}रि\appbrk{}वायां}}\Afootnote{H; \skt{पद्म\appbrk{}क\appbrk{}सा\appbrk{}रि\appbrk{}वासु}\ A}}\skt{ काको\-लि२\-मे\-दा\-कु\-मु\-दो\-त्प\-लेषु~॥}\pend

\pstart\versepar \ledsidenote{\passref{4.17.11}}\edtext{\skt{सच\-न्द\-ना\-या\-म्म\-धु\-श\-र्क्क\-रा\-या\-न्द्रा\-क्षा\-स्थि\-रा\-पृ\-ष्णि\-स\-म\-ङ्गया च}}{\lemma{\skt{द्राक्षा\appbrk{}स्थि\appbrk{}रा\appbrk{}पृ\appbrk{}ष्णि\appbrk{}स\appbrk{}म\appbrk{}ङ्गया च}}\Afootnote{H; \skt{द्राक्षा\appbrk{}स्थि\appbrk{}रा\appbrk{}पृ\appbrk{}श्नि\appbrk{}श\appbrk{}ता\appbrk{}ह्व\appbrk{}यासु}\ A}}\skt{~। }\pend
\pstart\versepar \skt{}\edtext{\skt{कल्की\-कृ\-ता\-सू\-द\-क\-म\-त्र\-गन्धा}}{\lemma{\skt{कल्की\appbrk{}कृ\appbrk{}ता\appbrk{}सू\appbrk{}द\appbrk{}क\appbrk{}म\appbrk{}त्र\appbrk{}गन्धा}}\Afootnote{H; \skt{कल्की\appbrk{}कृ\appbrk{}ता\appbrk{}सू\appbrk{}द\appbrk{}क\appbrk{}मत्र दत्त्वा}\ A}}\skt{ }\edtext{\skt{न्यग्रोधवर्ग्गस्य च}}{\lemma{\skt{न्यग्रोधवर्ग्गस्य च}}\Afootnote{H; \skt{न्यग्रोधवर्गस्य त}\ A}}\skt{ }\edtext{\skt{शारिवादेः}}{\lemma{\skt{शारिवादेः}}\Afootnote{H; \skt{था स्थिरादेः}\ A}}\skt{~॥}\pend

\pstart\versepar \ledsidenote{\passref{4.17.12}}\skt{गणस्य }\edtext{\edtext{\skt{विल्वा\-दि\-क\-प\-ञ्च\-मू\-ल्यो\-श्च\-तु\-र्ग्गुणं}}{\lemma{\skt{विल्वा\appbrk{}दि\appbrk{}क\appbrk{}प\appbrk{}ञ्च\appbrk{}मू\appbrk{}ल्योश्}}\Afootnote{H; \skt{बिल्वा\appbrk{}दि\appbrk{}क\appbrk{}प\appbrk{}ञ्च\appbrk{}मूल्या}\ A}}}{\lemma{\skt{चतुर्ग्गुणं}}\Afootnote{H; \skt{श्चतुर्गुणं}\ A}}\skt{ क्षीरमथापि तद्वत्~। }\pend
\pstart\versepar \skt{प्रस्थं विपक्वं परिषेचनेन }\edtext{\skt{पैत्तानि}}{\lemma{\skt{पैत्तानि}}\Afootnote{H; \skt{पैत्तीर्नि}\ A}}\skt{ हन्यात्तु }\edtext{\skt{विसप्प्पनाडी}}{\lemma{\skt{विसप्प्पनाडी}}\Afootnote{H; \skt{विसर्पनाडीः}\ A}}\skt{~॥}\pend

\pstart\versepar \ledsidenote{\passref{4.17.13}}\edtext{\skt{विस्फो\-ट\-दु\-ष्ट\-व्र\-ण\-शी\-र्ष\-रो\-गा\-न्पा\-क\-न्त\-था\-स्यास्य}}{\lemma{\skt{तथास्यास्य}}\Afootnote{H; \skt{तथा++आस्यस्य}\ A}}\skt{ }\edtext{\skt{तथैव}}{\lemma{\skt{तथैव}}\Afootnote{H; \skt{निहन्ति}\ A}}\skt{ पानात्~। }\pend
\pstart\versepar \skt{ग्रहार्द्दिते शोषिणि चापि वाले~। }\edtext{\skt{घृतं तु}}{\lemma{\skt{घृतं तु}}\Afootnote{H; \skt{घृतं हि}\ A}}\skt{ }\edtext{\edtext{\skt{गौरा\-दि\-क\-मे\-त\-दिष्टं}}{\lemma{\skt{गौरादिकम्}}\Afootnote{H; \skt{गौर्यादिक}\ A}}}{\lemma{\skt{एतदिष्टं}}\Afootnote{H; \skt{मेतदिष्टम्}\ A}}\skt{~॥}\pend

\pstart\versepar \ledsidenote{\passref{4.17.14}}\edtext{\skt{अथाश्वगन्धा}}{\lemma{\skt{अथाश्वगन्धा}}\Afootnote{H; \skt{अजा+अश्वगन्धा}\ A}}\skt{ }\edtext{\skt{सरलाथ}}{\lemma{\skt{सरलाथ}}\Afootnote{H; \skt{सरला स}\ A}}\skt{ काला }\edtext{\skt{सैके\-षि\-का\-वा\-प्य\-थ\-वा\-ज\-शृङ्गी}}{\lemma{\skt{सैके\appbrk{}षि\appbrk{}का\appbrk{}वा\appbrk{}प्य\appbrk{}थ\appbrk{}वा\appbrk{}ज\appbrk{}शृङ्गी}}\Afootnote{H; \skt{सैकैषिका चाप्य\appbrk{}थ\appbrk{}वा+\appbrk{}अ\appbrk{}ज\appbrk{}शृङ्गी}\ A}}\skt{~। }\pend
\pstart\versepar \skt{}\edtext{\skt{गोमूत्रपिष्टो}}{\lemma{\skt{गोमूत्रपिष्टो}}\Afootnote{A H\textsuperscript{pc}; \skt{गोमूत्रपि}\ H\textsuperscript{ac}}}\skt{ }\edtext{\skt{विहिमप्रदेहो}}{\lemma{\skt{विहिमप्रदेहो}}\Afootnote{H; \skt{विहितः प्रदेहो}\ A}}\skt{ हन्या\-द्वि\-स\-र्प्प\-ङ्क\-फजं }\edtext{\skt{सुशीघ्रं}}{\lemma{\skt{सुशीघ्रं}}\Afootnote{H; \skt{स शीघ्रम्}\ A}}\skt{~॥}\pend

\pstart\versepar \ledsidenote{\passref{4.17.15}}\edtext{\skt{काला\-नु\-सा\-र्या\-गु\-रु\-चो\-च\-गु\-ञ्जा\-र\-स्ना\-व\-चा\-शी\-त\-शि\-वे\-न्द्र\-पर्ण्ण्यः}}{\lemma{\skt{काला\appbrk{}नु\appbrk{}सा\appbrk{}र्या\appbrk{}गु\appbrk{}रु\appbrk{}चो\appbrk{}च\appbrk{}गु\appbrk{}ञ्जा\appbrk{}र\appbrk{}स्ना\appbrk{}व\appbrk{}चा\appbrk{}शी\appbrk{}त\appbrk{}शि\appbrk{}वे\appbrk{}न्द्र\appbrk{}पर्ण्ण्यः}}\Afootnote{H; \skt{काला\appbrk{}नु\appbrk{}सा\appbrk{}र्या\appbrk{}गु\appbrk{}रु\appbrk{}चो\appbrk{}च\appbrk{}गु\appbrk{}ञ्जा\appbrk{}रा\appbrk{}स्ना\appbrk{}व\appbrk{}चा\appbrk{}शी\appbrk{}त\appbrk{}शि\appbrk{}वे\appbrk{}द्र\appbrk{}पर्ण्यः}\ A}}\skt{~। }\pend
\pstart\versepar \skt{}\edtext{\skt{पालि\-न्दि\-म\-ज्जा\-स\-म\-ही\-क\-दम्वा}}{\lemma{\skt{पालि\appbrk{}न्दि\appbrk{}म\appbrk{}ज्जा\appbrk{}स\appbrk{}म\appbrk{}ही\appbrk{}क\appbrk{}दम्वा}}\Afootnote{H; \skt{पालि\appbrk{}न्दि\appbrk{}मु\appbrk{}ञ्जा\appbrk{}त\appbrk{}म\appbrk{}ही\appbrk{}क\appbrk{}दम्बा}\ A}}\skt{ }\edtext{\skt{हितो}}{\lemma{\skt{हितो}}\Afootnote{H; \skt{हिता}\ A}}\skt{ विसर्प्पेषु कफात्मकेषु~॥}\pend

\pstart\versepar \ledsidenote{\passref{4.17.16}}\skt{गणस्तु }\edtext{\skt{सर्व्वो}}{\lemma{\skt{सर्व्वो}}\Afootnote{H; \skt{योज्यो}\ A}}\skt{ }\edtext{\skt{वरुणः}}{\lemma{\skt{वरुणः}}\Afootnote{H; \skt{वरुण}\ A}}\skt{ }\edtext{\skt{प्रयोज्यः}}{\lemma{\skt{प्रयोज्यः}}\Afootnote{H; \skt{प्रवृत्तः}\ A}}\skt{ क्रियासु सर्व्वासु विचक्षणेन~। }\pend
\pstart\versepar \skt{}\edtext{\skt{सर्व्वेषु}}{\lemma{\skt{सर्व्वेषु}}\Afootnote{H; \emph{om.}\ A}}\skt{ }\edtext{\skt{संशोधनमेव}}{\lemma{\skt{एव}}\Afootnote{H; \emph{om.}\ A}}\skt{ }\edtext{\edtext{\skt{कुर्या\-द्र\-क्ता\-व\-से\-कञ्च}}{\lemma{\skt{कुर्याद्}}\Afootnote{H; \emph{om.}\ A}}}{\lemma{\skt{रक्तावसेकञ्च}}\Afootnote{H; \emph{om.}\ A}}\skt{ }\edtext{\skt{जलायुकाभिः}}{\lemma{\skt{जलायुकाभिः}}\Afootnote{H; \skt{शोणितमोक्षणं च श्रेष्ठं विसर्पेषु चिकित्सितं हि}\ A}}\skt{~॥}\pend

\pstart\versepar \ledsidenote{\passref{4.17.17}}\skt{सर्व्वाश्च पक्वान्परिशोध्य }\edtext{\skt{धीमान्व्रणं}}{\lemma{\skt{व्रणं}}\Afootnote{H; \skt{व्रण}\ A}}\skt{ क्रमे\-णो\-प\-च\-रे\-द्य\-थोक्तं~। }\pend
\pstart\versepar \skt{नाडी\-त्रि\-दो\-ष\-प्र\-भवा न सिध्ये }\edtext{\edtext{\skt{च्छेषाश्चतस्रः}}{\lemma{\skt{च्छेषाश्}}\Afootnote{H; \skt{च्छेषा}\ A}}}{\lemma{\skt{चतस्रः}}\Afootnote{H; \skt{श्चतस्रः}\ A}}\skt{ खलु यत्नसाध्याः~॥}\pend

\pstart\versepar \ledsidenote{\passref{4.17.18}}\skt{तत्रा\-नि\-लो\-त्था\-मु\-प\-नाह्य }\edtext{\edtext{\skt{पूर्व्वमशेषतः}}{\lemma{\skt{पूर्व्वम्}}\Afootnote{H; \skt{पूर्व}\ A}}}{\lemma{\skt{अशेषतः}}\Afootnote{H; \skt{मशेषतः}\ A}}\skt{ पूय\-ग\-ति\-म्वि\-दार्य~। }\pend
\pstart\versepar \skt{}\edtext{\edtext{\skt{फलै\-र\-पा\-मा\-र्ग्ग\-भवैः}}{\lemma{\skt{फलैर्}}\Afootnote{H; \skt{तिलैर्}\ A}}}{\lemma{\skt{अपामार्ग्गभवैः}}\Afootnote{H; \skt{अपा\appbrk{}मा\appbrk{}र्ग\appbrk{}फ\appbrk{}लैश्}\ A}}\skt{ }\edtext{\skt{सुपिष्टैः}}{\lemma{\skt{सुपिष्टैः}}\Afootnote{H; \skt{च पिष्ट्वा}\ A}}\skt{ ससैन्धवैः }\edtext{\skt{पूर्य}}{\lemma{\skt{पूर्य}}\Afootnote{H; \skt{बन्धन}\ A}}\skt{ }\edtext{\skt{हिताय}}{\lemma{\skt{हिताय}}\Afootnote{H; \skt{मत्र}\ A}}\skt{ }\edtext{\skt{वध्येत्}}{\lemma{\skt{वध्येत्}}\Afootnote{H; \skt{कुर्यात्}\ A}}\skt{~॥}\pend

\pstart\versepar \ledsidenote{\passref{4.17.19}}\skt{प्रक्षालने चापि सदा }\edtext{\skt{व्रणस्य}}{\lemma{\skt{व्रणस्य}}\Afootnote{H; \skt{वरणस्य}\ A}}\skt{ }\edtext{\skt{योज्य\-म्म\-ह\-द्य\-त्खलु}}{\lemma{\skt{महद्}}\Afootnote{H; \skt{महद्}\ A}}\skt{ पञ्चमूलं~। }\pend
\pstart\versepar \skt{}\edtext{\skt{हिंस्रा}}{\lemma{\skt{हिंस्रा}}\Afootnote{H; \skt{हिंस्रां}\ A}}\skt{ }\edtext{\skt{हरिद्रे}}{\lemma{\skt{हरिद्रे}}\Afootnote{H; \skt{हरिद्रां}\ A}}\skt{ }\edtext{\skt{कटुकावले च}}{\lemma{\skt{कटुकावले च}}\Afootnote{H; \skt{कटुकां बलां च}\ A}}\skt{ }\edtext{\skt{गोजिह्विका}}{\lemma{\skt{गोजिह्विका}}\Afootnote{H; \skt{गोजिह्विकां}\ A}}\skt{ }\edtext{\skt{चैव}}{\lemma{\skt{चैव}}\Afootnote{H; \skt{चापि}\ A}}\skt{ }\edtext{\skt{सविल्वमूलं}}{\lemma{\skt{सविल्वमूलं}}\Afootnote{H; \skt{सबिल्वमूलाम्}\ A}}\skt{~॥}\pend

\pstart\versepar \ledsidenote{\passref{4.17.20}}\edtext{\skt{सम्भृत्य}}{\lemma{\skt{सम्भृत्य}}\Afootnote{H; \skt{संहृत्य}\ A}}\skt{ }\edtext{\edtext{\skt{कुर्यात्त्रिवृतं}}{\lemma{\skt{कुर्यात्}}\Afootnote{H; \skt{तै}\ A}}}{\lemma{\skt{त्रिवृतं}}\Afootnote{H; \skt{लं विपचेद्}\ A}}\skt{ व्रणस्य संशोधनं }\edtext{\skt{शोध\-न\-रो\-प\-णज्ञः}}{\lemma{\skt{शोध\appbrk{}न\appbrk{}रो\appbrk{}प\appbrk{}णज्ञः}}\Afootnote{H; \skt{पूरणरोपणं च}\ A}}\skt{~। }\pend
\pstart\versepar \skt{}\edtext{\skt{पित्ता\-त्म\-कां\-न्प्रा\-गु\-प\-नाह्य}}{\lemma{\skt{पित्तात्मकांन्}}\Afootnote{H; \skt{पित्तात्मिकां}\ A}}\skt{ धीमा\-नु\-त्का\-रि\-काभिः सपयोघृताभिः~॥}\pend

\pstart\versepar \ledsidenote{\passref{4.17.21}}\edtext{\skt{निपा\-त्य\-श\-स्त्र\-न्ति\-ल\-ना\-ग\-द\-न्ती\-म\-ञ्जि\-ष्ठ\-कल्कैः}}{\lemma{\skt{तिल\appbrk{}ना\appbrk{}ग\appbrk{}द\appbrk{}न्ती\appbrk{}म\appbrk{}ञ्जि\appbrk{}ष्ठ\appbrk{}कल्कैः}}\Afootnote{H; \skt{तिल\appbrk{}ना\appbrk{}ग\appbrk{}द\appbrk{}न्ती\appbrk{}य\appbrk{}ष्ट्या\appbrk{}ह्व\appbrk{}कल्कैः}\ A}}\skt{ }\edtext{\skt{प्रतिपूरयीत}}{\lemma{\skt{प्रतिपूरयीत}}\Afootnote{H; \skt{परि\appbrk{}पू\appbrk{}र\appbrk{}ये\appbrk{}त्ताम्}\ A}}\skt{~। }\pend
\pstart\versepar \skt{प्रक्षालने चापि }\edtext{\skt{हरि\-द्र\-सो\-म\-नि\-म्वाः\-प्र\-योज्याः}}{\lemma{\skt{हरि\appbrk{}द्र\appbrk{}सो\appbrk{}म\appbrk{}नि\appbrk{}म्वाः\appbrk{}प्र\appbrk{}योज्याः}}\Afootnote{H; \skt{ससोमनिम्बा निशा प्रयोज्या}\ A}}\skt{ कुशलेन नित्यं~॥}\pend

\pstart\versepar \ledsidenote{\passref{4.17.22}}\skt{श्यामा\-त्रि\-भ\-ण्डी\-त्रि\-फ\-लासु सिद्धं }\edtext{\edtext{\skt{हरि\-द्र\-यो\-स्ति\-ल्व\-क\-वृ\-क्षके}}{\lemma{\skt{हरिद्रयोस्}}\Afootnote{H; \skt{हरिद्रयो}\ A}}}{\lemma{\skt{तिल्वकवृक्षके}}\Afootnote{H; \skt{रोध्रकवृक्षयो}\ A}}\skt{ }\edtext{\skt{च}}{\lemma{\skt{च}}\Afootnote{H; \skt{श्च}\ A}}\skt{~। }\pend
\pstart\versepar \skt{घृतं सदु\-ग्ध\-म्व्र\-ण\-त\-र्प्प\-णेन हन्याद्गतिं }\edtext{\skt{कोष्ठगतापि या}}{\lemma{\skt{कोष्ठगतापि या}}\Afootnote{H; \skt{कोष्ठगता+अपि या}\ A}}\skt{ स्यात्~॥}\pend

\pstart\versepar \ledsidenote{\passref{4.17.23}}\edtext{\skt{नाडी\-ङ्क\-फो\-त्था\-मु\-प\-ना\-ह\-यित्वा}}{\lemma{\skt{उपनाहयित्वा}}\Afootnote{H; \skt{उपनाह्य सम्यक्}\ A}}\skt{ कुल\-त्थ\-सि\-द्धा\-र्थ\-क\-श\-क्तु\-किण्वैः~। }\pend
\pstart\versepar \skt{}\edtext{\skt{मृद्वी\-कृ\-ता\-मे\-ष्य\-ग\-ति\-म्वि\-दित्वा}}{\lemma{\skt{मृद्वी\appbrk{}कृ\appbrk{}ता\appbrk{}मे\appbrk{}ष्य\appbrk{}ग\appbrk{}तिम्}}\Afootnote{H; \skt{मृदूकृतामेष्य गतिं}\ A}}\skt{ }\edtext{\skt{निपा\-त\-ये\-च्छ\-स्त्र\-म\-शे\-ष\-कारी}}{\lemma{\skt{छस्त्र\appbrk{}म\appbrk{}शे\appbrk{}ष\appbrk{}कारी}}\Afootnote{\skt{छस्त्र\appbrk{}म\appbrk{}शे\appbrk{}ष\appbrk{}कारी}\ A H}}\skt{~॥}\pend

\pstart\versepar \ledsidenote{\passref{4.17.24}}\edtext{\skt{दद्याद्व्रणे}}{\lemma{\skt{दद्याद्}}\Afootnote{H; \skt{दद्याद्}\ A}}\skt{ }\edtext{\skt{निम्वतिलां}}{\lemma{\skt{निम्वतिलां}}\Afootnote{H; \skt{निम्बतिलान्}\ A}}\skt{ }\edtext{\skt{सदन्तीं}}{\lemma{\skt{सदन्तीं}}\Afootnote{H; \skt{सदन्तीन्}\ A}}\skt{ }\edtext{\skt{सुराष्ट्रजां}}{\lemma{\skt{सुराष्ट्रजां}}\Afootnote{H; \skt{सुराष्ट्रजा}\ A}}\skt{ }\edtext{\skt{सैन्ध\-व\-स\-म्प्र\-युक्तां}}{\lemma{\skt{सैन्ध\appbrk{}व\appbrk{}स\appbrk{}म्प्र\appbrk{}युक्तां}}\Afootnote{H; \skt{सैन्ध\appbrk{}व\appbrk{}सं\appbrk{}प्र\appbrk{}यु\appbrk{}क्तान्}\ A}}\skt{~। }\pend
\pstart\versepar \skt{प्रक्षालने चापि कर\-ञ्ज\-नि\-म्व\-जा\-त्य\-क्ष\-पी\-लु\-स्व\-रसाः प्रयोज्याः~॥}\pend

\pstart\versepar \ledsidenote{\passref{4.17.25}}\skt{सुव\-र्च्चि\-का\-सै\-न्ध\-व\-चि\-त्र\-केषु }\edtext{\skt{निकुम्भनीलीत}}{\lemma{\skt{निकुम्भनीलीत}}\Afootnote{H; \skt{निकुम्भतालीत}\ A}}\skt{ }\edtext{\skt{लरूपिकासु}}{\lemma{\skt{लरूपिकासु}}\Afootnote{A H\textsuperscript{pc}; \skt{लरूकासु}\ H\textsuperscript{ac}}}\skt{~। }\pend
\pstart\versepar \skt{फले\-ष्व\-पा\-मा\-र्ग्ग\-भ\-वेषु चैव }\edtext{\skt{कुर्या\-त्स\-म\-मू\-त्रेषु}}{\lemma{\skt{सममूत्रेषु}}\Afootnote{H; \skt{समूत्रेषु}\ A}}\skt{ हिताय तैलं~॥}\pend

\pstart\versepar \ledsidenote{\passref{4.17.26}}\skt{नाडीन्तु शल्यप्रभवां विदार्य }\edtext{\skt{निर्घात्यशल्यं}}{\lemma{\skt{निर्घात्यशल्यं}}\Afootnote{H; \skt{निर्हृत्य शल्यं}\ A}}\skt{ }\edtext{\skt{प्रविदार्य}}{\lemma{\skt{प्रविदार्य}}\Afootnote{H; \skt{प्रविशोध्य}\ A}}\skt{ मार्ग्गं~। }\pend
\pstart\versepar \skt{}\edtext{\edtext{\skt{वन्धेद्व्रणं}}{\lemma{\skt{वन्धेद्}}\Afootnote{H; \skt{संशोधयेत्}\ A}}}{\lemma{\skt{व्रणं}}\Afootnote{H; \emph{om.}\ A}}\skt{ }\edtext{\skt{क्षौद्र\-घृ\-त\-प्र\-गाढं}}{\lemma{\skt{क्षौद्र\appbrk{}घृ\appbrk{}त\appbrk{}प्र\appbrk{}गाढं}}\Afootnote{H; \skt{क्षौद्र\appbrk{}घृ\appbrk{}त\appbrk{}प्र\appbrk{}गा\appbrk{}ढैस्}\ A}}\skt{ }\edtext{\skt{दत्वा}}{\lemma{\skt{दत्वा}}\Afootnote{H; \emph{om.}\ A}}\skt{ }\edtext{\edtext{\skt{तिला\-ञ्च्छो\-ध्य\-व\-रो\-प\-यीत}}{\lemma{\skt{तिलाञ्}}\Afootnote{H; \skt{तिलैस्}\ A}}}{\lemma{\skt{च्छोध्य\appbrk{}व\appbrk{}रो\appbrk{}प\appbrk{}यीत}}\Afootnote{H; \skt{ततो रोपणमाशु कुर्यात्}\ A}}\skt{~॥}\pend

\pstart\versepar \ledsidenote{\passref{4.17.27}}\skt{कुम्भी\-क\-ख\-र्ज्जू\-र\-क\-पि\-त्थ\-वि\-ल्व\-व\-न\-स्प\-ती\-नाञ्च }\edtext{\skt{शलाटुवर्ग्गे}}{\lemma{\skt{शलाटुवर्ग्गे}}\Afootnote{H; \skt{शलाटुवर्गैः}\ A}}\skt{~। }\pend
\pstart\versepar \skt{कृत्वा कषा\-य\-म्वि\-प\-चेत्तु }\edtext{\edtext{\skt{तैलमावाप्य}}{\lemma{\skt{तैलम्}}\Afootnote{H; \skt{तैल}\ A}}}{\lemma{\skt{आवाप्य}}\Afootnote{H; \skt{मावाप्य}\ A}}\skt{ }\edtext{\skt{मुस्ता\-श\-कु\-न\-प्रि\-यङ्गूः}}{\lemma{\skt{मुस्ता\appbrk{}श\appbrk{}कु\appbrk{}न\appbrk{}प्रि\appbrk{}यङ्गूः}}\Afootnote{H; \skt{मुस्ता\appbrk{}स\appbrk{}र\appbrk{}ला\appbrk{}प्रि\appbrk{}यङ्गूः}\ A}}\skt{~॥}\pend

\pstart\versepar \ledsidenote{\passref{4.17.28}}\edtext{\skt{सौग\-न्धि\-का\-मो\-च\-र\-सा\-हि\-पुष्पं}}{\lemma{\skt{सौग\appbrk{}न्धि\appbrk{}का\appbrk{}मो\appbrk{}च\appbrk{}र\appbrk{}सा\appbrk{}हि\appbrk{}पुष्पं}}\Afootnote{H; \skt{सुग\appbrk{}न्धि\appbrk{}का\appbrk{}मो\appbrk{}च\appbrk{}र\appbrk{}सा\appbrk{}हि\appbrk{}पु\appbrk{}षोअं}\ A}}\skt{ }\edtext{\skt{लोध्राणि}}{\lemma{\skt{लोध्राणि}}\Afootnote{H; \skt{रोध्रं वि}\ A}}\skt{ }\edtext{\skt{दत्वा}}{\lemma{\skt{दत्वा}}\Afootnote{H; \skt{दध्या}\ A}}\skt{ }\edtext{\skt{खलु}}{\lemma{\skt{खलु}}\Afootnote{H; \skt{दपि}\ A}}\skt{ धातकीञ्च~। }\pend
\pstart\versepar \skt{एतेन शल्यप्रभवा तु नाडी }\edtext{\skt{रोहेद्व्रणे वा}}{\lemma{\skt{व्रणे वा}}\Afootnote{H; \skt{व्रणो वा}\ A}}\skt{ }\edtext{\skt{सुखमाशु}}{\lemma{\skt{सुखमाशु}}\Afootnote{H; \skt{सुखभाशु}\ A}}\skt{ चैव~॥}\pend

\pstart\versepar \ledsidenote{\passref{4.17.29}}\edtext{\edtext{\edtext{\skt{कृश\-दु\-र्व्व\-ल\-भी\-रू\-णा\-न्ना\-डी\-म्म\-र्म्मा\-श्रि\-ता\-मपि}}{\lemma{\skt{नाडीम्}}\Afootnote{H; \skt{नाडी}\ A}}}{\lemma{\skt{मर्म्माश्रिताम्}}\Afootnote{H; \skt{मर्माश्रिता}\ A}}}{\lemma{\skt{अपि}}\Afootnote{H; \skt{च या}\ A}}\skt{~। }\pend
\pstart\versepar \skt{क्षारसूत्रेण तां च्छिन्द्यान्न तु शस्त्रेण }\edtext{\skt{दारयेत्}}{\lemma{\skt{दारयेत्}}\Afootnote{H; \skt{बुद्धिमान्}\ A}}\skt{~॥}\pend

\pstart\versepar \ledsidenote{\passref{4.17.30}}\edtext{\skt{एषण्या }\edtext{\skt{गमन्विष्य}}{\lemma{\skt{अन्विष्य}}\Afootnote{H; \skt{तिमन्विष्य}\ A}}}{\lemma{\skt{एषण्या गम्}}\Afootnote{H; \skt{एषण्या ग}\ A}}\skt{ }\edtext{\skt{क्षार\-सू\-त्रा\-नु\-शा\-रिणीं}}{\lemma{\skt{क्षार\appbrk{}सू\appbrk{}त्रा\appbrk{}नु\appbrk{}शा\appbrk{}रिणीं}}\Afootnote{H; \skt{आर\appbrk{}सू\appbrk{}त्रा\appbrk{}नु\appbrk{}सा\appbrk{}रि\appbrk{}णीम्}\ A}}\skt{~। }\pend
\pstart\versepar \skt{सूची\-न्नि\-द\-ध्या\-द्ग\-त्यन्ते }\edtext{\edtext{\skt{प्रोन्नाम्पाण्डु}}{\lemma{\skt{प्रोन्नाम्}}\Afootnote{H; \skt{तहोन्न}\ A}}}{\lemma{\skt{पाण्डु}}\Afootnote{H; \skt{म्या}\ A}}\skt{ }\edtext{\skt{विनिर्हरेत्}}{\lemma{\skt{विनिर्हरेत्}}\Afootnote{H; \skt{शु निर्हरेत्}\ A}}\skt{~॥}\pend

\pstart\versepar \ledsidenote{\passref{4.17.31}}\skt{सूत्रस्यान्तं समानीय }\edtext{\skt{गाढवन्धं}}{\lemma{\skt{गाढवन्धं}}\Afootnote{H; \skt{गाढं बन्धं}\ A}}\skt{ समाचरेत्~। }\pend
\pstart\versepar \skt{ततः क्षारवलम्वीक्ष्य }\edtext{\edtext{\edtext{\skt{सूत्र\-म\-न्य\-त्प्र\-यो\-ज\-येत्}}{\lemma{\skt{सूत्रम्}}\Afootnote{H; \skt{सूत्र}\ A}}}{\lemma{\skt{अन्यत्}}\Afootnote{H; \skt{मन्यत्}\ A}}}{\lemma{\skt{प्रयोजयेत्}}\Afootnote{H; \skt{प्रवेशयेत्}\ A}}\skt{~।}\pend

\pstart\versepar \ledsidenote{\passref{4.17.32}}\skt{क्षारा\-क्त\-म्म\-ति\-मा\-न्वैद्यो यावन्न च्छिद्यते गतिः~। }\pend
\pstart\versepar \skt{भगन्दरेऽप्येष विधिः }\edtext{\skt{कार्या}}{\lemma{\skt{कार्या}}\Afootnote{H; \skt{कार्यो}\ A}}\skt{ वैद्येन जानता~॥}\pend

\pstart\versepar \ledsidenote{\passref{4.17.33}}\edtext{\skt{अर्व्वुदादिषु}}{\lemma{\skt{अर्व्वुदादिषु}}\Afootnote{H; \emph{om.}\ A}}\skt{ }\edtext{\skt{चोत्क्षिप्य}}{\lemma{\skt{चोत्क्षिप्य}}\Afootnote{H; \emph{om.}\ A}}\skt{ }\edtext{\skt{मूले}}{\lemma{\skt{मूले}}\Afootnote{H; \emph{om.}\ A}}\skt{ }\edtext{\edtext{\skt{सूत्र\-न्नि\-धा\-प\-येत्}}{\lemma{\skt{सूत्रन्}}\Afootnote{H; \emph{om.}\ A}}}{\lemma{\skt{निधापयेत्}}\Afootnote{H; \emph{om.}\ A}}\skt{~। }\pend
\pstart\versepar \skt{}\edtext{\skt{सूची\-भि\-र्य\-व\-व\-क्त्रा\-भि\-रा\-चि\-तम्वा}}{\lemma{\skt{आचितम्वा}}\Afootnote{H; \skt{आचितान्वा}\ A}}\skt{ समन्ततः~।. }\pend
\pstart\versepar \skt{मूले }\edtext{\skt{सूत्रे च}}{\lemma{\skt{सूत्रे च}}\Afootnote{H; \skt{सूत्रेण}\ A}}\skt{ वध्नीयाच्च्छिन्ने }\edtext{\skt{चोपक्रमेद्व्रणं}}{\lemma{\skt{चोपक्रमेद्}}\Afootnote{H; \skt{चोपवरेद्}\ A}}\skt{~॥}\pend

\pstart\versepar \ledsidenote{\passref{4.17.34}}\skt{या द्विव्र\-णी\-येऽ\-भि\-हि\-तास्तु }\edtext{\edtext{\skt{वर्त्त्यस्ताः}}{\lemma{\skt{वर्त्त्यस्}}\Afootnote{H; \skt{वर्त्य}\ A}}}{\lemma{\skt{ताः}}\Afootnote{H; \skt{स्ताः}\ A}}\skt{ सर्व्वनाडीषु भिषग्विदध्यात्~। }\pend
\pstart\versepar \skt{}\edtext{\skt{घोण्टाफलं}}{\lemma{\skt{घोण्टाफलं}}\Afootnote{H; \skt{घोण्टाफल}\ A}}\skt{ }\edtext{\skt{त्वग्मदनं}}{\lemma{\skt{मदनं}}\Afootnote{H; \skt{लवणानि लाक्षापूगी}\ A}}\skt{ }\edtext{\skt{फलञ्च}}{\lemma{\skt{फलञ्च}}\Afootnote{H; \skt{फलं}\ A}}\skt{ }\edtext{\skt{पूगस्य च}}{\lemma{\skt{पूगस्य च}}\Afootnote{H; \skt{चा}\ A}}\skt{ }\edtext{\skt{त्वग्लवणञ्च}}{\lemma{\skt{त्वग्}}\Afootnote{H; \emph{om.}\ A}}\skt{ }\edtext{\skt{मुख्यं}}{\lemma{\skt{मुख्यं}}\Afootnote{H; \skt{पत्रम्}\ A}}\skt{~॥}\pend

\pstart\versepar \ledsidenote{\passref{4.17.35}}\edtext{\skt{स्नुहार्क्कदुग्धेन}}{\lemma{\skt{स्नुहार्क्कदुग्धेन}}\Afootnote{H; \skt{स्रुह्यर्कदुग्धेन}\ A}}\skt{ }\edtext{\skt{सहैष}}{\lemma{\skt{सहैष}}\Afootnote{H; \skt{तु}\ A}}\skt{ }\edtext{\skt{कल्को}}{\lemma{\skt{कल्को}}\Afootnote{H; \skt{कल्क}\ A}}\skt{ }\edtext{\skt{वर्त्तीकृता}}{\lemma{\skt{वर्त्तीकृता}}\Afootnote{H; \skt{एष वर्तीकृतो}\ A}}\skt{ }\edtext{\skt{हन्त्य}}{\lemma{\skt{हन्त्य}}\Afootnote{H; \skt{हन्त्य}\ A}}\skt{ चिरेण }\edtext{\skt{नाडीं}}{\lemma{\skt{नाडीं}}\Afootnote{H; \skt{नाडीः}\ A}}\skt{ }\pend
\pstart\versepar \skt{}\edtext{\skt{विभी\-त\-का\-म्रा\-स्थि\-व\-ट\-प्र\-वा\-ल\-ह\-रे\-णु\-का\-शं\-खि\-नि\-वी\-ज\-मिश्रा}}{\lemma{\skt{विभी\appbrk{}त\appbrk{}का\appbrk{}म्रा\appbrk{}स्थि\appbrk{}व\appbrk{}ट\appbrk{}प्र\appbrk{}वा\appbrk{}ल\appbrk{}ह\appbrk{}रे\appbrk{}णु\appbrk{}का\appbrk{}शं\appbrk{}खि\appbrk{}नि\appbrk{}वी\appbrk{}ज\appbrk{}मिश्रा}}\Afootnote{H; \skt{बिभी\appbrk{}त\appbrk{}का\appbrk{}म्रा\appbrk{}स्थि\appbrk{}व\appbrk{}ट\appbrk{}प्र\appbrk{}वाला हरे\appbrk{}णु\appbrk{}का\appbrk{}श\appbrk{}ङ्खि\appbrk{}नि\appbrk{}बी\appbrk{}ज\appbrk{}मस्यः}\ A}}\skt{~॥}\pend

\pstart \ledsidenote{\passref{4.17.36}}\edtext{\skt{वरा\-ह\-वि\-ट्सू\-क्ष्म\-म\-सी\-प्र\-देया}}{\lemma{\skt{वरा\appbrk{}ह\appbrk{}वि\appbrk{}ट्सू\appbrk{}क्ष्म\appbrk{}म\appbrk{}सी\appbrk{}प्र\appbrk{}देया}}\Afootnote{H; \skt{वाराहिकन्दश्च तथा प्रदेयो}\ A}}\skt{ नाडीषु तैलेन }\edtext{\skt{विमिश्रयित्वा}}{\lemma{\skt{विमिश्रयित्वा}}\Afootnote{H; \skt{च मिश्रयित्वा}\ A}}\skt{~।}\pend

\pstart\versepar \ledsidenote{\passref{4.17.38}}\skt{तथा च गोमांसमसीं हिताय }\edtext{\skt{कोष्ठाश्रिता}}{\lemma{\skt{कोष्ठाश्रिता}}\Afootnote{H; \skt{कोष्ठाश्रित}\ A}}\skt{ }\edtext{\skt{स्वादरतो}}{\lemma{\skt{स्वादरतो}}\Afootnote{H; \skt{स्यादरतो}\ A}}\skt{ दिशन्ति~।. }\pend
\pstart\versepar \skt{}\edtext{\skt{वर्त्ती\-कृ\-त\-म्मा\-क्षि\-क\-सं\-प्र\-यु\-क्त\-न्ना\-डी\-घ्न\-मुक्तं}}{\lemma{\skt{वर्त्तीकृतम्}}\Afootnote{H; \skt{वरीकृतं}\ A}}\skt{ लवणोत्तमम्वा~।}\pend

\pstart \ledsidenote{\passref{4.17.39}}\skt{दुष्टव्रणे }\edtext{\skt{यद्विहितं}}{\lemma{\skt{यद्}}\Afootnote{H; \skt{यद्}\ A}}\skt{ च }\edtext{\edtext{\edtext{\skt{तैल\-न्त\-त्से\-व्य\-मा\-न\-ङ्ग\-ति\-माशु}}{\lemma{\skt{तत्}}\Afootnote{H; \skt{तत्सर्वनाडीषु भिषग्विदध्यात्}\ A}}}{\lemma{\skt{सेव्यमानङ्}}\Afootnote{H; \skt{चूर्णीकृतैरथ विमि\appbrk{}श्रि\appbrk{}त\appbrk{}मे\appbrk{}भि\appbrk{}रेव तैलं प्रयु\appbrk{}क्त\appbrk{}म\appbrk{}चि\appbrk{}रेण}\ A}}}{\lemma{\skt{आशु}}\Afootnote{H; \skt{नि}\ A}}\skt{ हन्ति~॥}\pend

\pstart\versepar \ledsidenote{\passref{4.17.37}}\edtext{\skt{धत्तू\-र\-क\-म्म\-द\-न\-को\-द्र\-व\-जञ्च}}{\lemma{\skt{धत्तूरकम्}}\Afootnote{H; \skt{धत्तूरजं}\ A}}\skt{ वीजं कोशातकी शुकनसा मृगभोजिनी च~। }\pend
\pstart\versepar \skt{}\edtext{\edtext{\skt{साङ्को\-ल्ल\-वी\-ज\-कु\-सु\-मैस्तु}}{\lemma{\skt{साङ्को\appbrk{}ल्ल\appbrk{}वी\appbrk{}ज\appbrk{}कु\appbrk{}सु\appbrk{}मैस्}}\Afootnote{H; \skt{अङ्को\appbrk{}ट\appbrk{}बी\appbrk{}ज\appbrk{}कु\appbrk{}सुमं}\ A}}}{\lemma{\skt{तु}}\Afootnote{H; \skt{गति}\ A}}\skt{ }\edtext{\skt{कृतञ्च}}{\lemma{\skt{कृतञ्च}}\Afootnote{H; \skt{षु प्र}\ A}}\skt{ }\edtext{\skt{चूर्ण्णं}}{\lemma{\skt{चूर्ण्णं}}\Afootnote{H; \skt{योज्यं}\ A}}\skt{ लाक्षो\-द\-का\-हृ\-त\-म\-लासु }\edtext{\skt{गतीषु}}{\lemma{\skt{गतीषु}}\Afootnote{H; \emph{om.}\ A}}\skt{ }\edtext{\skt{देयं}}{\lemma{\skt{देयं}}\Afootnote{H; \skt{विकृत्य चूर्णम्}\ A}}\skt{~॥}\pend

\pstart\versepar \ledsidenote{\passref{4.17.40}}\skt{एष्वेव }\edtext{\skt{साधुषु}}{\lemma{\skt{साधुषु}}\Afootnote{H; \skt{मूत्रस}\ A}}\skt{ }\edtext{\skt{कृतेषु}}{\lemma{\skt{कृतेषु}}\Afootnote{H; \skt{हितेषु}\ A}}\skt{ }\edtext{\skt{विमिश्रयित्वा}}{\lemma{\skt{विमिश्रयित्वा}}\Afootnote{H; \skt{विधाय}\ A}}\skt{ तैलं }\edtext{\skt{कृत\-ङ्ग\-ति\-म\-पो\-हति}}{\lemma{\skt{कृतङ्}}\Afootnote{H; \skt{तत्साधितं}\ A}}\skt{ सप्तरात्रात्~। }\pend
\pstart\versepar \skt{}\edtext{\skt{पिण्डीकृतस्य तु}}{\lemma{\skt{पिण्डीकृतस्य तु}}\Afootnote{H; \skt{पिण्डीतकस्य तु}\ A}}\skt{ वरा\-ह\-वि\-भा\-वि\-तस्य मूलेषु कन्दशकलेषु च सौवहेषु~॥}\pend

\pstart\versepar \ledsidenote{\passref{4.17.41}}\skt{तैल\-ङ्कृ\-त\-ङ्ग\-ति\-म\-पो\-हति }\edtext{\edtext{\skt{शीघ्रमेव}}{\lemma{\skt{शीघ्रम्}}\Afootnote{H; \skt{शीघ्र}\ A}}}{\lemma{\skt{एव}}\Afootnote{H; \skt{मेतत्}\ A}}\skt{ कन्देषु चाम\-र\-व\-रा\-यु\-ध\-सा\-ह्व\-येषु~। }\pend
\pstart\versepar \skt{भल्ला\-त\-का\-र्क्क\-म\-रि\-चै\-र्ल्ल\-व\-णो\-त्त\-मेन सिद्धं }\edtext{\skt{विडं\-ग\-र\-ज\-नी\-द्व\-य\-चि\-त्र\-कै\-र्यत्}}{\lemma{\skt{यत्}}\Afootnote{H; \skt{च}\ A}}\skt{~॥}\pend

\pstart\versepar \ledsidenote{\passref{4.17.42}}\edtext{\skt{स्यामार्क्कवस्य च}}{\lemma{\skt{स्यामार्क्कवस्य च}}\Afootnote{H; \skt{स्यन्मार्कवस्य च}\ A}}\skt{ रसेन निहन्ति तैलं नाडीं }\edtext{\skt{कफा\-नि\-ल\-कृ\-ता\-न\-प\-ची\-न्व्र\-णाश्च}}{\lemma{\skt{कफा\appbrk{}नि\appbrk{}ल\appbrk{}कृ\appbrk{}ता\appbrk{}न\appbrk{}प\appbrk{}चीन्}}\Afootnote{H; \skt{कफा\appbrk{}नि\appbrk{}ल\appbrk{}कृ\appbrk{}ता\appbrk{}म\appbrk{}पचीं}\ A}}\skt{~। }\pend
\pstart\versepar \skt{स्तन्ये गते }\edtext{\edtext{\skt{विकृतिमाशु}}{\lemma{\skt{विकृतिम्}}\Afootnote{H; \skt{विकृति}\ A}}}{\lemma{\skt{आशु}}\Afootnote{H; \skt{माशु}\ A}}\skt{ भिषक्तु धात्रीं }\edtext{\skt{पीत्वा}}{\lemma{\skt{पीत्वा}}\Afootnote{H; \skt{पीतां}\ A}}\skt{ घृतं }\edtext{\skt{परिणतेहनि}}{\lemma{\skt{परिणतेहनि}}\Afootnote{H; \skt{परिणतेऽहनि}\ A}}\skt{ }\edtext{\skt{वामयीत}}{\lemma{\skt{वामयीत}}\Afootnote{H; \skt{वामयेत्तु}\ A}}\skt{~॥}\pend

\pstart\versepar \ledsidenote{\passref{4.17.43}}\skt{निम्वोदकेन मधु\-मा\-ग\-धि\-का\-यु\-तेन }\edtext{\skt{वान्तापरेऽहनि}}{\lemma{\skt{वान्तापरे}}\Afootnote{H; \skt{वान्तागते}\ A}}\skt{ च }\edtext{\skt{मुद्गरसासना}}{\lemma{\skt{मुद्गरसासना}}\Afootnote{H; \skt{मुद्गरसाशना}\ A}}\skt{ स्यात्~। }\pend
\pstart\versepar \skt{एव\-न्त्र्य\-ह\-ञ्च\-तु\-रहं षडहम्वमेद्वा सर्प्पिः }\edtext{\skt{पिवेद्त्रिफलया}}{\lemma{\skt{पिवेद्}}\Afootnote{H; \skt{पिबेत्}\ A}}\skt{ सह संयुतम्वा~॥}\pend

\pstart\versepar \ledsidenote{\passref{4.17.44}}\edtext{\edtext{\skt{भार्ग्गी\-म्व\-चा\-म\-ति\-विषां}}{\lemma{\skt{वचाम्}}\Afootnote{H; \emph{om.}\ A}}}{\lemma{\skt{अतिविषां}}\Afootnote{H; \emph{om.}\ A}}\skt{ }\edtext{\edtext{\edtext{\edtext{\skt{सुर\-दा\-रु\-पा\-ठा\-म्मु\-स्ता\-दि\-क\-म्म\-धु\-र\-सा\-ङ्क\-टु\-रो\-हि\-णीञ्च}}{\lemma{\skt{सुर\appbrk{}दा\appbrk{}रु\appbrk{}पा\appbrk{}ठाम्}}\Afootnote{H; \emph{om.}\ A}}}{\lemma{\skt{मुस्तादिकम्}}\Afootnote{H; \emph{om.}\ A}}}{\lemma{\skt{मधुरसाङ्}}\Afootnote{H; \emph{om.}\ A}}}{\lemma{\skt{कटुरोहिणीञ्च}}\Afootnote{H; \emph{om.}\ A}}\skt{~। }\pend
\pstart\versepar \skt{}\edtext{\skt{धात्री}}{\lemma{\skt{धात्री}}\Afootnote{H; \emph{om.}\ A}}\skt{ पिवेत्तु पयसः }\edtext{\skt{परिशोधनार्थ}}{\lemma{\skt{परिशोधनार्थ}}\Afootnote{H; \skt{परि\appbrk{}शो\appbrk{}ध\appbrk{}ना\appbrk{}र्थम्}\ A}}\skt{ }\edtext{\skt{मुस्तावचादि तु}}{\lemma{\skt{मुस्तावचादि तु}}\Afootnote{H; \skt{आरग्वधादिषु}\ A}}\skt{ }\edtext{\skt{वरम्मधुनाथ}}{\lemma{\skt{मधुनाथ}}\Afootnote{H; \skt{मधुना क}\ A}}\skt{ }\edtext{\skt{वापि}}{\lemma{\skt{वापि}}\Afootnote{H; \skt{षायम्}\ A}}\skt{~॥}\pend

\pstart \ledsidenote{\passref{4.17.45}}\edtext{\edtext{\skt{स्तन\-रो\-ग\-मु\-त्थि\-त\-म\-वेक्ष्य}}{\lemma{\skt{स्तनरोगम्}}\Afootnote{H; \skt{सामा\appbrk{}न्य\appbrk{}मे\appbrk{}त\appbrk{}दु\appbrk{}प\appbrk{}दि\appbrk{}ष्ट\appbrk{}मतो विशे\appbrk{}षा\appbrk{}द्दो\appbrk{}षा\appbrk{}न्प\appbrk{}यो\appbrk{}नि\appbrk{}प\appbrk{}ति\appbrk{}ता\appbrk{}न्श\appbrk{}म\appbrk{}ये\appbrk{}द्य\appbrk{}था\appbrk{}स्वम्~। रोगं स्तनो}\ A}}}{\lemma{\skt{उत्थितम्}}\Afootnote{H; \skt{त्थितम्}\ A}}\skt{ }\edtext{\skt{भिष\-ग्वि\-द\-ध्या\-द्य\-द्वि\-द्र\-धा\-व\-भि\-हि\-त\-म्ब\-हु\-धा\-भि\-धा\-नम्}}{\lemma{\skt{बहु\appbrk{}धा\appbrk{}भि\appbrk{}धा\appbrk{}नम्}}\Afootnote{H; \skt{बहुशो विधानम्}\ A}}\skt{~॥}\pend

\pstart \ledsidenote{\passref{4.17.47}}\edtext{\skt{आमे}}{\lemma{\skt{आमे}}\Afootnote{H; \skt{पक्वे तु दुग्धहरिणीः परिहृत्य नाडीः कृष्णं च चूचुकयुगं विदधीत शस्त्रम्~। आमे}\ A}}\skt{ }\edtext{\skt{विदह्यति}}{\lemma{\skt{विदह्यति}}\Afootnote{H; \skt{विदाहिनि}\ A}}\skt{ तथैव गते च }\edtext{\skt{पाकन्तस्या}}{\lemma{\skt{तस्या}}\Afootnote{H; \skt{धात्र्याः}\ A}}\skt{ स्तनौ सततमेव च }\edtext{\skt{निर्द्दुहीतेति}}{\lemma{\skt{निर्द्दुहीतेति}}\Afootnote{H; \skt{निर्दुहीत}\ A}}\skt{~॥}\pend

\pstart \ledsidenote{\passref{4.17.trailer}}\edtext{\skt{॥ चिकित्सा .. ओ}}{\lemma{\skt{॥ चिकित्सा .. ओ}}\Afootnote{H; \skt{इति सुश्रु\appbrk{}त\appbrk{}सं\appbrk{}हि\appbrk{}तायां चिकित्सास्थाने विस\appbrk{}र्प\appbrk{}ना\appbrk{}डी\appbrk{}स्त\appbrk{}न\appbrk{}रो\appbrk{}ग\appbrk{}चि\appbrk{}कि\appbrk{}त्सितं नाम सप्तदशोऽध्यायः}\ A}}\skt{~॥}\pend

\endnumbering
\beginnumbering
\eledsection*{Adhyāya 18}
\pstart \ledsidenote{\passref{4.18.1}}\skt{अथातो }\edtext{\skt{ग्रन्थ्य\-प\-च्य\-र्ब्बु\-द\-ग\-ल\-ग\-ण्डा\-ना\-ञ्चि\-कि\-त्सि\-त\-म्व्या\-ख्या\-स्यामः}}{\lemma{\skt{ग्रन्थ्य\appbrk{}प\appbrk{}च्य\appbrk{}र्ब्बु\appbrk{}द\appbrk{}ग\appbrk{}ल\appbrk{}ग\appbrk{}ण्डा\appbrk{}नाञ्}}\Afootnote{H; \skt{ग्रन्थ्य\appbrk{}प\appbrk{}च्य\appbrk{}र्बु\appbrk{}द\appbrk{}ग\appbrk{}ल\appbrk{}गण्ड}\ A}}\skt{~॥}\pend

\pstart \ledsidenote{\passref{4.18.3}}\edtext{\edtext{\skt{ग्रन्थिष्वथामेषु}}{\lemma{\skt{ग्रन्थिष्व्}}\Afootnote{H; \skt{ग्रन्थिष्}\ A}}}{\lemma{\skt{अथामेषु}}\Afootnote{H; \skt{वथामेषु}\ A}}\skt{ }\edtext{\skt{भिष\-ग्वि\-द\-ध्या\-च्छो\-फ\-क्रिया}}{\lemma{\skt{छोफक्रिया}}\Afootnote{H; \skt{छोफक्रियां}\ A}}\skt{ }\edtext{\skt{विस्तरशो}}{\lemma{\skt{विस्तरशो}}\Afootnote{H; \skt{विस्तरशोकार्याः पित्तार्बुदे कायविरेचनं च}\ A}}\skt{ }\edtext{\skt{विधिज्ञः}}{\lemma{\skt{विधिज्ञः}}\Afootnote{H; \skt{विघृष्य चोदु\appbrk{}म्ब\appbrk{}र\appbrk{}शा\appbrk{}क\appbrk{}गो\appbrk{}जी\appbrk{}प\appbrk{}त्रै\appbrk{}र्भृशं क्षौद्रयुतैः प्रलिम्पेत्}\ A}}\skt{~।}\pend

\pstart \ledsidenote{\passref{4.18.3.add}}\skt{रक्षेद्बलञ्चापि नरस्य नित्यं तद्र\-क्षि\-त\-म्व्या\-धि\-ब\-ल\-न्नि\-हन्ति~॥}\pend

\pstart \ledsidenote{\passref{4.18.4}}\skt{तैलं }\edtext{\skt{पिबेत्सर्प्पिरथ}}{\lemma{\skt{सर्प्पिरथ}}\Afootnote{H; \skt{सर्पिरथो}\ A}}\skt{ द्वयं वा दत्वा वसाम्वा }\edtext{\skt{त्रिवृ\-त\-म्वि\-द\-ध्यात्। अपी\-ह\-वा\-ते\-द\-श\-मू\-ल\-सिद्धं}}{\lemma{\skt{विदध्यात्। अपी\appbrk{}ह\appbrk{}वा\appbrk{}ते\appbrk{}द\appbrk{}श\appbrk{}मू\appbrk{}ल\appbrk{}सिद्धं}}\Afootnote{H; \skt{विदध्यात्। अपे\appbrk{}हि\appbrk{}वा\appbrk{}ता\appbrk{}द\appbrk{}श\appbrk{}मू\appbrk{}ल\appbrk{}सिद्धं}\ A}}\skt{ }\edtext{\skt{वैद्य\-श्च\-तुः\-स्ने\-ह\-म\-थापि}}{\lemma{\skt{वैद्य\appbrk{}श्च\appbrk{}तुः\appbrk{}स्ने\appbrk{}ह\appbrk{}म\appbrk{}थापि}}\Afootnote{H; \skt{वैद्य\appbrk{}श्च\appbrk{}तुः\appbrk{}स्ने\appbrk{}ह\appbrk{}मथो द्}\ A}}\skt{ }\edtext{\skt{दद्यात्}}{\lemma{\skt{दद्यात्}}\Afootnote{H; \skt{वयं वा}\ A}}\skt{~॥}\pend

\pstart \ledsidenote{\passref{4.18.5}}\edtext{\skt{हिंस्रास}}{\lemma{\skt{हिंस्रास}}\Afootnote{H; \skt{हिंस्रा+अथ}\ A}}\skt{ }\edtext{\skt{रोहिण्यमृताथ}}{\lemma{\skt{रोहिण्यमृताथ}}\Afootnote{H; \skt{रोहि\appbrk{}ण्य\appbrk{}मृ\appbrk{}ता+\appbrk{}अथ}\ A}}\skt{ भर्ग्गी श्योना\-क\-बि\-ल्वा\-गु\-रु\-कृ\-ष्ण\-गन्धा~। गोजी च पिष्टा सह तालपत्त्र्या ग्रन्थौ }\edtext{\skt{विधेया}}{\lemma{\skt{विधेया}}\Afootnote{H; \skt{विधेयो}\ A}}\skt{ }\edtext{\skt{निलजे}}{\lemma{\skt{निलजे}}\Afootnote{H; \skt{ऽनिलजे}\ A}}\skt{ प्रलेपः॥}\pend

\pstart \ledsidenote{\passref{4.18.6}}\edtext{\skt{स्वेदोपनाहा}}{\lemma{\skt{स्वेदोपनाहा}}\Afootnote{H; \skt{स्वेदोपनाहान्}\ A}}\skt{ }\edtext{\skt{विविधाश्च}}{\lemma{\skt{विविधाश्च}}\Afootnote{H; \skt{विविधांश्च}\ A}}\skt{ }\edtext{\skt{योज्याः}}{\lemma{\skt{योज्याः}}\Afootnote{H; \emph{om.}\ A}}\skt{ }\edtext{\skt{कुर्या\-त्प्र\-सि\-द्धा\-न\-प\-रांश्च}}{\lemma{\skt{प्रसि\appbrk{}द्धा\appbrk{}न\appbrk{}प\appbrk{}रांश्च}}\Afootnote{H; \skt{तथा प्रसि\appbrk{}द्धा\appbrk{}न\appbrk{}प\appbrk{}रांश्च}\ A}}\skt{ लेपान्। }\edtext{\skt{विदार्य}}{\lemma{\skt{विदार्य}}\Afootnote{H; \skt{विदार्य वा}\ A}}\skt{ }\edtext{\skt{पक्वं}}{\lemma{\skt{पक्वं}}\Afootnote{H; \skt{पक्}\ A}}\skt{ }\edtext{\skt{समुपोह्य}}{\lemma{\skt{समुपोह्य}}\Afootnote{H; \skt{वमपोह्य}\ A}}\skt{ पूयम्प्रक्षाल्य }\edtext{\skt{बिल्वा\-र्क्क\-ग\-ण\-दि\-तोयैः}}{\lemma{\skt{बिल्वा\appbrk{}र्क्क\appbrk{}ग\appbrk{}ण\appbrk{}दि\appbrk{}तोयैः}}\Afootnote{H; \skt{बिल्वा\appbrk{}र्क\appbrk{}न\appbrk{}रे\appbrk{}न्द्र\appbrk{}तोयैः}\ A}}\skt{~॥}\pend

\pstart \ledsidenote{\passref{4.18.7}}\skt{तिलैः सप\-ञ्चा\-ङ्गु\-ल\-प\-त्त्र\-मिश्रैः संशो\-ध\-ये\-त्सै\-न्ध\-व\-सं\-प्र\-युक्तैः। }\edtext{\edtext{\skt{शुद्ध\-म्व्र\-ण\-ञ्चा\-प्यु\-प\-रो\-प\-यीत}}{\lemma{\skt{चाप्य्}}\Afootnote{H; \skt{वा+अप्}\ A}}}{\lemma{\skt{उपरोपयीत}}\Afootnote{H; \skt{युपरोपयेत्तु}\ A}}\skt{ तैलेन }\edtext{\skt{रास्ना\-स\-र\-ला\-यु\-तेन॥}}{\lemma{\skt{रास्ना\appbrk{}स\appbrk{}र\appbrk{}ला\appbrk{}यु\appbrk{}तेन}}\Afootnote{H; \skt{रास्ना\appbrk{}स\appbrk{}र\appbrk{}ला\appbrk{}न्वि\appbrk{}तेन}\ A}}\pend

\pstart \ledsidenote{\passref{4.18.8}}\skt{विड\-ङ्ग\-य\-ष्टी\-म\-धु\-का\-मृ\-ताभिः }\edtext{\skt{सिद्धेन तं}}{\lemma{\skt{सिद्धेन तं}}\Afootnote{H; \skt{सिद्धेन वा}\ A}}\skt{ }\edtext{\skt{क्षीर\-वि\-मि\-श्रि\-तेन।}}{\lemma{\skt{क्षीर\appbrk{}वि\appbrk{}मि\appbrk{}श्रि\appbrk{}तेन}}\Afootnote{H; \skt{क्षीर\appbrk{}स\appbrk{}म\appbrk{}न्वि\appbrk{}तेन}\ A}}\skt{ }\edtext{\skt{जलायुकाः}}{\lemma{\skt{जलायुकाः}}\Afootnote{H; \skt{जलौकसः}\ A}}\skt{ पित्तकृते हितास्तु }\edtext{\skt{क्षीरो\-द\-का\-भ्या\-म्प\-रि\-षे\-च\-नञ्च॥}}{\lemma{\skt{परिषेचनञ्}}\Afootnote{H; \skt{परिबेचनं}\ A}}\pend

\pstart \ledsidenote{\passref{4.18.9}}\edtext{\skt{काकोलिवर्ग्गस्य तु}}{\lemma{\skt{काकोलिवर्ग्गस्य तु}}\Afootnote{H; \skt{काकोलिवर्गस्य च}\ A}}\skt{ शीतलानि पिबेत्कषायाणि सशर्क्कराणि~। द्राक्षा\-र\-से\-ने\-क्षु\-र\-सेन }\edtext{\skt{वापि}}{\lemma{\skt{वापि}}\Afootnote{H; \skt{वा+अपि}\ A}}\skt{ चूर्ण्णं }\edtext{\edtext{\skt{पिबेद्वापि}}{\lemma{\skt{पिबेद्}}\Afootnote{H; \skt{पिबेच्}\ A}}}{\lemma{\skt{वापि}}\Afootnote{H; \skt{चापि}\ A}}\skt{ हरीतकीनां॥}\pend

\pstart \ledsidenote{\passref{4.18.10}}\skt{मधू\-क\-ज\-म्ब्व\-र्ज्जु\-न\-वे\-त\-सानां त्वग्भिः }\edtext{\skt{प्रदे\-हा\-न\-व\-चा\-र\-येच्च।}}{\lemma{\skt{प्रदे\appbrk{}हा\appbrk{}न\appbrk{}व\appbrk{}चा\appbrk{}र\appbrk{}येच्च}}\Afootnote{H; \skt{प्रदे\appbrk{}हा\appbrk{}न\appbrk{}व\appbrk{}चा\appbrk{}र\appbrk{}येत}\ A}}\skt{ सशर्क्करैर्व्वा }\edtext{\edtext{\skt{तृण\-शू\-ल्य\-क\-न्दै\-र्द्दि\-ह्या\-द\-भीक्ष्णं}}{\lemma{\skt{तृण\appbrk{}शू\appbrk{}ल्य\appbrk{}क\appbrk{}न्दैर्}}\Afootnote{\skt{तृण\appbrk{}शू\appbrk{}न्य\appbrk{}क\appbrk{}न्दैर्}\ A; \skt{तृणशूल्यकन्दै}\ H}}}{\lemma{\skt{द्दिह्यादभीक्ष्णं}}\Afootnote{\skt{दिह्यादभीक्ष्णं}\ A; \skt{र्द्दिह्यादभीक्ष्णं}\ H}}\skt{ }\edtext{\skt{मुचिलिन्दजैर्व्वा॥}}{\lemma{\skt{मुचिलिन्दजैर्}}\Afootnote{H; \skt{मुचुलुन्दजैर्}\ A}}\pend

\pstart \ledsidenote{\passref{4.18.11}}\skt{विदार्य वा पक्वमपोह्य पूयं धावेत्कषायेण वनस्पतीनां। तिलैः सय\-ष्टी\-म\-धु\-कै\-र्व्वि\-शोध्य सर्प्पिः प्रयो\-ज्य\-म्म\-धु\-रै\-र्व्वि\-पक्वं॥}\pend

\pstart \ledsidenote{\passref{4.18.12}}\skt{हृतेषु दोषेषु }\edtext{\skt{यथानुपूर्व्वा}}{\lemma{\skt{यथानुपूर्व्वा}}\Afootnote{H; \skt{यथानुपूर्व्या}\ A}}\skt{ ग्रन्थौ }\edtext{\skt{भिष\-क्छ्ले\-ष्म\-स\-मु\-त्थि\-तेषु}}{\lemma{\skt{छ्लेष्म\appbrk{}स\appbrk{}मु\appbrk{}त्थि\appbrk{}तेषु}}\Afootnote{H; \skt{श्लेष्मसमुत्थिते तु}\ A}}\skt{~। स्विन्नस्य }\edtext{\edtext{\skt{विम्लापनमेव}}{\lemma{\skt{विम्लापनम्}}\Afootnote{H; \skt{विम्लापन}\ A}}}{\lemma{\skt{एव}}\Afootnote{H; \skt{मेव}\ A}}\skt{ }\edtext{\edtext{\skt{कुर्या\-दं\-गु\-ष्ठ\-वे\-णु\-दृ\-श\-दी\-सु\-तैस्तु॥}}{\lemma{\skt{कुर्या\appbrk{}दं\appbrk{}गु\appbrk{}ष्ठ\appbrk{}वे\appbrk{}णु\appbrk{}दृ\appbrk{}श\appbrk{}दी\appbrk{}सु\appbrk{}तैस्}}\Afootnote{H; \skt{कुर्या\appbrk{}द\appbrk{}ङ्गु\appbrk{}ष्ठ\appbrk{}लो\appbrk{}हो\appbrk{}प\appbrk{}ल\appbrk{}वे\appbrk{}णु\appbrk{}दण्डैः}\ A}}}{\lemma{\skt{तु}}\Afootnote{H; \emph{om.}\ A}}\pend

\pstart \ledsidenote{\passref{4.18.13}}\skt{विकं\-क\-ता\-र\-ग्व\-ध\-का\-क\-ण\-न्ती\-का\-का\-द\-नी\-ता\-प\-स\-वृ\-क्ष\-मूलैः। }\edtext{\skt{आले\-प\-ये\-दे\-नम}}{\lemma{\skt{आले\appbrk{}प\appbrk{}ये\appbrk{}दे\appbrk{}नम}}\Afootnote{H; \skt{आले\appbrk{}प\appbrk{}ये\appbrk{}त्पि\appbrk{}ण्डफ}\ A}}\skt{ }\edtext{\skt{लावु\-भा\-र्ग्गी\-क\-र\-ञ्ज\-का\-ला\-म\-द\-नैश्च}}{\lemma{\skt{लावु\appbrk{}भा\appbrk{}र्ग्गी\appbrk{}क\appbrk{}र\appbrk{}ञ्ज\appbrk{}का\appbrk{}ला\appbrk{}म\appbrk{}द\appbrk{}नैश्च}}\Afootnote{H; \skt{लार्क\appbrk{}भा\appbrk{}र्गी\appbrk{}क\appbrk{}र\appbrk{}ञ्ज\appbrk{}का\appbrk{}ला\appbrk{}म\appbrk{}द\appbrk{}नैश्च}\ A}}\skt{ विद्वान्॥}\pend

\pstart \ledsidenote{\passref{4.18.14}}\edtext{\skt{अमर्म्मजे}}{\lemma{\skt{अमर्म्मजे}}\Afootnote{H; \skt{अमर्मजा}\ A}}\skt{ }\edtext{\skt{चानुप}}{\lemma{\skt{चानुप}}\Afootnote{H; \skt{तं शम}\ A}}\skt{ }\edtext{\edtext{\skt{शीम्य्मा\-न\-स\-प\-क्व\-मे\-वा\-प\-ह\-रे\-द्वि\-दार्य।}}{\lemma{\skt{शीम्य्मा\appbrk{}न\appbrk{}स\appbrk{}प\appbrk{}क्वम्}}\Afootnote{H; \skt{मप्र\appbrk{}या\appbrk{}न्त\appbrk{}म\appbrk{}पक्व}\ A}}}{\lemma{\skt{एवापहरेद्}}\Afootnote{H; \skt{मेवापहरेद्}\ A}}\skt{ दहेत्स्थिते चासृजि सिद्धकर्म्मा सद्यः क्षतोक्तञ्च विधिम्विदध्यात्॥}\pend

\pstart \ledsidenote{\passref{4.18.15}}\skt{या मांसकन्द्यः कठिना }\edtext{\skt{बृहत्यस्तास्वप्य यं}}{\lemma{\skt{बृहत्यस्तास्वप्य यं}}\Afootnote{H; \skt{बृहत्यस्तास्}\ A}}\skt{ }\edtext{\skt{कर्म्म}}{\lemma{\skt{कर्म्म}}\Afootnote{H; \emph{om.}\ A}}\skt{ }\edtext{\skt{विधिः}}{\lemma{\skt{विधिः}}\Afootnote{H; \skt{वे}\ A}}\skt{ }\edtext{\skt{प्रयोज्यः।}}{\lemma{\skt{प्रयोज्यः}}\Afootnote{H; \skt{ष योज्यश्च}\ A}}\skt{ }\edtext{\skt{शस्त्रेण}}{\lemma{\skt{शस्त्रेण}}\Afootnote{H; \skt{विधिर्विधिज्ञैः। शस्त्रेण}\ A}}\skt{ }\edtext{\skt{चापाट्य}}{\lemma{\skt{चापाट्य}}\Afootnote{H; \skt{वा++आपाठ्य}\ A}}\skt{ }\edtext{\skt{विपक्वमाशु}}{\lemma{\skt{विपक्वमाशु}}\Afootnote{\skt{सुपक्वमाशु}\ A; \skt{विपक्वमाशु}\ H}}\skt{ प्रक्षा\-ल\-ये\-त्प\-थ्य\-तमैः कषायैः॥}\pend

\pstart \ledsidenote{\passref{4.18.16}}\skt{संशोधनैस्तञ्च }\edtext{\skt{विशोधयीत}}{\lemma{\skt{विशोधयीत}}\Afootnote{H; \skt{विशोधयेत्तु}\ A}}\skt{ }\edtext{\skt{क्षारोत्तमैः}}{\lemma{\skt{क्षारोत्तमैः}}\Afootnote{H; \skt{क्षारोत्तरैः}\ A}}\skt{ }\edtext{\skt{क्षौद्र\-घृ\-त\-प्र\-गाढैः।}}{\lemma{\skt{क्षौद्र\appbrk{}घृ\appbrk{}त\appbrk{}प्र\appbrk{}गाढैः}}\Afootnote{H; \skt{क्षौद्र\appbrk{}गु\appbrk{}ड\appbrk{}प्र\appbrk{}गाढैः}\ A}}\skt{ शुद्धे च तैलं त्वव\-चा\-र\-णी\-य\-म्वि\-ड\-ङ्ग\-पा\-ठा\-र\-ज\-नी\-वि\-पक्वं॥}\pend

\pstart \ledsidenote{\passref{4.18.17}}\skt{मेदःसमुत्थे तिलकल्कदिग्धं }\edtext{\skt{त्वचो\-प\-रि\-ष्टा\-द्द्वि\-गु\-ण\-म्प\-टान्तं।}}{\lemma{\skt{त्वचोपरिष्टाद्}}\Afootnote{H; \skt{दत्त्वोपरिष्टा}\ A}}\skt{ हुताशतप्तेन मुहुः }\edtext{\skt{प्रमृज्यालेन}}{\lemma{\skt{प्रमृज्यालेन}}\Afootnote{H; \skt{प्रमृज्याल्लोहेन}\ A}}\skt{ धीमा\-न\-द\-ह\-न्हि\-ताय॥}\pend

\pstart \ledsidenote{\passref{4.18.18}}\skt{प्रलिप्य दार्व्वीमथ लाक्षया वा प्रतप्तया }\edtext{\edtext{\skt{स्वेदनमस्य}}{\lemma{\skt{स्वेदनम्}}\Afootnote{H; \skt{स्वेदन}\ A}}}{\lemma{\skt{अस्य}}\Afootnote{H; \skt{मस्य}\ A}}\skt{ कार्यं। निपात्य वा शस्त्रमपोह्य मेदो दहे\-त्सु\-प\-क्व\-न्त्व\-थवा विदार्य॥}\pend

\pstart \ledsidenote{\passref{4.18.19}}\skt{प्रक्षाल्य }\edtext{\skt{मूत्रैश्च}}{\lemma{\skt{मूत्रैश्च}}\Afootnote{H; \skt{मूत्रेण}\ A}}\skt{ तिलैः सुपिष्टैः }\edtext{\skt{सुव\-र्च्चि\-का\-द्यै\-र्ह\-रि\-ता\-ल\-मिस्रैः।}}{\lemma{\skt{हरितालमिस्रैः}}\Afootnote{H; \skt{हरितालमिश्रैः}\ A}}\skt{ ससैन्धवैः क्षौद्र\-घृ\-त\-प्र\-गाढैः }\edtext{\edtext{\skt{क्षारो\-त्त\-रै\-रे\-न\-म\-भि\-प्र\-शोध्य॥}}{\lemma{\skt{एनम्}}\Afootnote{H; \skt{एन}\ A}}}{\lemma{\skt{अभिप्रशोध्य}}\Afootnote{H; \skt{मभिप्रशोध्य}\ A}}\pend

\pstart \ledsidenote{\passref{4.18.20}}\skt{तैलं }\edtext{\edtext{\skt{विद\-द्या\-द्द्वि\-क\-र\-ञ्ज\-गु\-ञ्जा\-व\-ङ्गा\-व\-ले\-खे\-ङ्गु\-दि\-मू\-त्र\-सिद्धैः।}}{\lemma{\skt{विदद्याद्}}\Afootnote{H; \skt{विदद्याद्}\ A}}}{\lemma{\skt{द्विक\appbrk{}र\appbrk{}ञ्ज\appbrk{}गु\appbrk{}ञ्जा\appbrk{}व\appbrk{}ङ्गा\appbrk{}व\appbrk{}ले\appbrk{}खे\appbrk{}ङ्गु\appbrk{}दि\appbrk{}मू\appbrk{}त्र\appbrk{}सिद्धैः}}\Afootnote{H; \skt{द्विक\appbrk{}र\appbrk{}ञ्ज\appbrk{}गु\appbrk{}ञ्जा\appbrk{}वं\appbrk{}शा\appbrk{}व\appbrk{}ले\appbrk{}खे\appbrk{}ङ्गु\appbrk{}द\appbrk{}मू\appbrk{}त्र\appbrk{}सि\appbrk{}द्धम्}\ A}}\skt{ }\edtext{\skt{जीमूतकी}}{\lemma{\skt{जीमूतकी}}\Afootnote{H; \skt{जीमूतकैः}\ A}}\skt{ }\edtext{\skt{कोश\-व\-ती\-फ\-लेषु}}{\lemma{\skt{कोश\appbrk{}व\appbrk{}ती\appbrk{}फ\appbrk{}लेषु}}\Afootnote{H; \skt{कोश\appbrk{}व\appbrk{}ती\appbrk{}फ\appbrk{}लैश्च}\ A}}\skt{ दन्ती द्रवन्ती त्रिवृतासु चैव॥}\pend

\pstart \ledsidenote{\passref{4.18.21}}\skt{सर्प्पिः कृतं }\edtext{\edtext{\skt{हन्त्यपचीं}}{\lemma{\skt{हन्त्य्}}\Afootnote{H; \skt{हन्त्}\ A}}}{\lemma{\skt{अपचीं}}\Afootnote{H; \skt{यपचीं}\ A}}\skt{ प्रवृद्धां द्विधा }\edtext{\skt{प्रवृ\-त्त\-न्त\-दु\-दा\-र\-वीर्यं।}}{\lemma{\skt{तदुदारवीर्यं}}\Afootnote{H; \skt{तदु\appbrk{}दा\appbrk{}र\appbrk{}वी\appbrk{}र्यम्}\ A}}\skt{ निर्ग्गु\-ण्डि\-जा\-ती\-ब\-रि\-हि\-ष्ठ\-युक्तं जीमूतकं माक्षि\-क\-सै\-न्ध\-वाढ्यं॥}\pend

\pstart \ledsidenote{\passref{4.18.22}}\skt{अभिप्रतप्तं वमनं }\edtext{\skt{प्रशस्तं}}{\lemma{\skt{प्रशस्तं}}\Afootnote{H; \skt{प्रगाढं}\ A}}\skt{ दुष्टा\-प\-ची\-षू\-त्त\-म\-मा\-दि\-शन्ति। }\edtext{\skt{निर्ग्गुण्डि}}{\lemma{\skt{निर्ग्गुण्डि}}\Afootnote{H; \skt{कैडर्य}\ A}}\skt{ बिम्बी\-क\-र\-वी\-र\-सिद्धं तैलं हितं }\edtext{\skt{शीर्षविरेचनं}}{\lemma{\skt{शीर्षविरेचनं}}\Afootnote{H; \skt{मूर्धविरेचनं}\ A}}\skt{ च॥}\pend

\pstart \ledsidenote{\passref{4.18.23}}\edtext{\skt{शाखोटकल्क}}{\lemma{\skt{शाखोटकल्क}}\Afootnote{H; \skt{शास्वोटकस्य}\ A}}\skt{ स्वरसेन सिद्धं तैलं हितं }\edtext{\skt{नस्य\-नि\-शे\-व\-नेन।}}{\lemma{\skt{नस्य\appbrk{}नि\appbrk{}शे\appbrk{}व\appbrk{}नेन}}\Afootnote{H; \skt{नस्य\appbrk{}वि\appbrk{}रे\appbrk{}च\appbrk{}नेषु}\ A}}\skt{ मधूकसारश्च हितो }\edtext{\skt{वपीडे}}{\lemma{\skt{वपीडे}}\Afootnote{H; \skt{ऽवपीडे}\ A}}\skt{ फलानि शिग्रोः खरमञ्जरेर्व्वा॥}\pend

\pstart \ledsidenote{\passref{4.18.24}}\skt{ग्रन्थी\-न\-म\-र्म्म\-प्र\-भ\-वा\-न\-प\-क्वा\-नु\-द्धृत्य चाग्निं विदधीत }\edtext{\skt{वैद्यः।}}{\lemma{\skt{वैद्यः}}\Afootnote{H; \skt{पश्चात्}\ A}}\skt{ क्षारेण }\edtext{\skt{चैनं}}{\lemma{\skt{चैनं}}\Afootnote{H; \skt{वा+अपि}\ A}}\skt{ }\edtext{\edtext{\skt{प्रति\-सा\-र\-यी\-त्मु\-हु\-र्म्म्हु\-र्ल्लि\-ख्यम}}{\lemma{\skt{प्रतिसारयीत्}}\Afootnote{H; \skt{प्रतिसारयेत्}\ A}}}{\lemma{\skt{मुहुर्म्म्हुर्ल्लिख्यम}}\Afootnote{H; \skt{तु संलिख्य शस्त्रेण य}\ A}}\skt{ }\edtext{\skt{थोयथावत्॥}}{\lemma{\skt{थोयथावत्}}\Afootnote{H; \skt{थोपदेशम्}\ A}}\pend

\pstart \ledsidenote{\passref{4.18.25}}\edtext{\skt{पार्ष्णिम्प्रति}}{\lemma{\skt{पार्ष्णिम्}}\Afootnote{H; \skt{पार्ष्णि}\ A}}\skt{ }\edtext{\skt{द्वा}}{\lemma{\skt{द्वा}}\Afootnote{H; \skt{द्वे}\ A}}\skt{ दश चाङ्गुलानि }\edtext{\skt{मात्वे\-न्द्र\-ब\-स्ति\-म्प\-रि\-वर्ज्ज्य}}{\lemma{\skt{मात्वेन्द्रबस्तिम्}}\Afootnote{H; \skt{मित्वेन्द्रबस्तिं}\ A}}\skt{ धीमान्। विदार्य मत्स्याण्डनिभानि वैद्यो }\edtext{\skt{जालानि}}{\lemma{\skt{जालानि}}\Afootnote{H; \skt{निष्कृष्य जाला}\ A}}\skt{ }\edtext{\skt{निःक्र\-म्य\-वि\-धि\-म्वि\-द्ध्यात्}}{\lemma{\skt{निःक्र\appbrk{}म्य\appbrk{}वि\appbrk{}धि\appbrk{}म्वि\appbrk{}द्ध्यात्}}\Afootnote{H; \skt{न्यनलं विदध्यात्}\ A}}\skt{~॥}\pend

\pstart \ledsidenote{\passref{4.18.27}}\skt{मणि\-ब\-न्धो\-प\-रि\-ष्टाद्वा }\edtext{\skt{कुर्या\-द्रे\-खा\-त्र\-य\-म्भि\-षक्। अङ्गुलान्तरितं}}{\lemma{\skt{भिषक्। अङ्गुलान्तरितं}}\Afootnote{H; \skt{भिषक्। अङ्गुल्यन्तरितं}\ A}}\skt{ सम्यगपचीनां निवृत्तये॥}\pend

\pstart \ledsidenote{\passref{4.18.28}}\skt{चूर्ण्णस्य काले }\edtext{\edtext{\skt{पर्च\-ला\-क\-का\-क\-गो\-धा\-हि\-कू\-र्म्म\-प्र\-भ\-वा\-स्तु\-मष्यः}}{\lemma{\skt{पर्च\appbrk{}ला\appbrk{}क\appbrk{}का\appbrk{}क\appbrk{}गो\appbrk{}धा\appbrk{}हि\appbrk{}कू\appbrk{}र्म्म\appbrk{}प्र\appbrk{}भ\appbrk{}वास्}}\Afootnote{H; \skt{पर्च\appbrk{}ला\appbrk{}क\appbrk{}का\appbrk{}क\appbrk{}गो\appbrk{}धा\appbrk{}हि\appbrk{}कू\appbrk{}र्म\appbrk{}प्र\appbrk{}भवां}\ A}}}{\lemma{\skt{तुमष्यः}}\Afootnote{H; \skt{मसीं तु}\ A}}\skt{~। दद्याच्च तैलेन सहेङ्गुदीनां }\edtext{\skt{तैलानि}}{\lemma{\skt{तैलानि}}\Afootnote{H; \skt{यद्}\ A}}\skt{ }\edtext{\skt{वक्ष्यामि च}}{\lemma{\skt{वक्ष्यामि च}}\Afootnote{H; \skt{वक्ष्यते}\ A}}\skt{ }\edtext{\skt{श्लीपदेषु}}{\lemma{\skt{श्लीपदेषु}}\Afootnote{H; \skt{श्लीपदिनां च तैलम्}\ A}}\skt{~॥}\pend

\pstart \ledsidenote{\passref{4.18.29}}\edtext{\edtext{\edtext{\skt{वैरे\-च\-न\-न्धू\-म\-मु\-पा\-द\-दीत}}{\lemma{\skt{वैरेचनन्}}\Afootnote{H; \skt{विरेचनं}\ A}}}{\lemma{\skt{धूमम्}}\Afootnote{H; \skt{धूम}\ A}}}{\lemma{\skt{उपाददीत}}\Afootnote{H; \skt{मुपाददीत}\ A}}\skt{ भवेच्च नित्यं यवमुद्गभोजी। }\edtext{\edtext{\skt{कर्क्का\-रु\-के\-र्व्वा\-रु\-क\-ना\-लि\-के\-र\-पि\-या\-ल\-प\-ञ्चा\-ङ्गु\-ल\-बी\-ज\-पूर्ण्णैः॥}}{\lemma{\skt{कर्क्कारुकेर्}}\Afootnote{H; \skt{कर्कारुकैर्}\ A}}}{\lemma{\skt{व्वारु\appbrk{}क\appbrk{}ना\appbrk{}लि\appbrk{}के\appbrk{}र\appbrk{}पि\appbrk{}या\appbrk{}ल\appbrk{}प\appbrk{}ञ्चा\appbrk{}ङ्गु\appbrk{}ल\appbrk{}बी\appbrk{}ज\appbrk{}पूर्ण्णैः}}\Afootnote{H; \skt{वारु\appbrk{}क\appbrk{}ना\appbrk{}सि\appbrk{}र्के\appbrk{}ल\appbrk{}प्रि\appbrk{}या\appbrk{}ल\appbrk{}प\appbrk{}ञ्चा\appbrk{}ङ्गु\appbrk{}ल\appbrk{}बी\appbrk{}ज\appbrk{}चूर्णैः}\ A}}\pend

\pstart \ledsidenote{\passref{4.18.30}}\skt{वातार्ब्बुदं क्षीर\-घृ\-ता\-म्बु\-सि\-द्धै\-रुष्णैः सतै\-लै\-रु\-प\-ना\-ह\-येत्तु। कुर्याच्च }\edtext{\edtext{\skt{मुख्या\-नु\-उ\-प\-ना\-ह\-ना\-र्थ\-म्भृ\-तेषु}}{\lemma{\skt{मुख्या\appbrk{}नु\appbrk{}उ\appbrk{}प\appbrk{}ना\appbrk{}ह\appbrk{}ना\appbrk{}र्थम्}}\Afootnote{H; \skt{मुख्या\appbrk{}न्यु\appbrk{}प\appbrk{}ना\appbrk{}हना}\ A}}}{\lemma{\skt{भृतेषु}}\Afootnote{H; \skt{नि सिद्धैश्च}\ A}}\skt{ }\edtext{\skt{मांसेषु च}}{\lemma{\skt{मांसेषु च}}\Afootnote{H; \skt{मांसैरथ}\ A}}\skt{ }\edtext{\skt{वेशवारां॥}}{\lemma{\skt{वेशवारां}}\Afootnote{H; \skt{वेसवारैः}\ A}}\pend

\pstart \ledsidenote{\passref{4.18.31}}\edtext{\skt{स्वेद\-म्वि\-द\-ध्या\-त्कु\-श\-लन्तु}}{\lemma{\skt{कुशलन्तु}}\Afootnote{H; \skt{कुशलस्तु}\ A}}\skt{ नाड्या }\edtext{\skt{शृङ्गेन}}{\lemma{\skt{शृङ्गेन}}\Afootnote{H; \skt{शृङ्गेण}\ A}}\skt{ रक्तम्बहुशो }\edtext{\skt{हरेच्च।}}{\lemma{\skt{हरेच्च}}\Afootnote{H; \skt{हरेच्च}\ A}}\skt{ वात\-घ्न\-नि\-र्यू\-ह\-प\-यो\-म्ल\-भागैः सिद्धं }\edtext{\skt{शताख्य}}{\lemma{\skt{शताख्य}}\Afootnote{H; \skt{शताख्यं}\ A}}\skt{ }\edtext{\skt{त्रिवृतम्पिबेद्वा॥}}{\lemma{\skt{पिबेद्वा}}\Afootnote{H; \skt{पिबेद्वा}\ A}}\pend

\pstart \ledsidenote{\passref{4.18.32}}\skt{स्वेदोपनाहा }\edtext{\skt{मृदुवस्तु}}{\lemma{\skt{मृदुवस्तु}}\Afootnote{H; \skt{मृदवस्तु}\ A}}\skt{ }\edtext{\skt{पथ्याः}}{\lemma{\skt{पथ्याः}}\Afootnote{H; \skt{कार्याः}\ A}}\skt{ पित्तार्ब्बुदे काय\-वि\-रे\-च\-नञ्च। विघृष्य }\edtext{\skt{चोदुम्बरि}}{\lemma{\skt{चोदुम्बरि}}\Afootnote{H; \skt{चोदुम्बर}\ A}}\skt{ }\edtext{\skt{काक\-गो\-जी\-प\-त्रै\-र्भ्भृशं}}{\lemma{\skt{काक\appbrk{}गो\appbrk{}जी\appbrk{}प\appbrk{}त्रै\appbrk{}र्भ्भृशं}}\Afootnote{H; \skt{शाक\appbrk{}गो\appbrk{}जी\appbrk{}प\appbrk{}त्रै\appbrk{}र्भृशं}\ A}}\skt{ क्षौद्रयुतैः प्रलिम्पेत्॥}\pend

\pstart\versepar \ledsidenote{\passref{4.18.33}}\edtext{\skt{सूक्ष्णीकृतैः}}{\lemma{\skt{सूक्ष्णीकृतैः}}\Afootnote{H; \skt{श्लक्ष्णीकृतैः}\ A}}\skt{ }\edtext{\skt{सर्ज्ज\-र\-स\-प्रि\-यं\-गू\-प\-त्त\-ङ्ग\-लो\-ध्रा\-ञ्ज\-न\-य\-ष्टि\-काह्वैः}}{\lemma{\skt{सर्ज्ज\appbrk{}र\appbrk{}स\appbrk{}प्रि\appbrk{}यं\appbrk{}गू\appbrk{}प\appbrk{}त्त\appbrk{}ङ्ग\appbrk{}लो\appbrk{}ध्रा\appbrk{}ञ्ज\appbrk{}न\appbrk{}य\appbrk{}ष्टि\appbrk{}काह्वैः}}\Afootnote{H; \skt{सर्ज\appbrk{}र\appbrk{}स\appbrk{}प्रि\appbrk{}य\appbrk{}ङ्गु\appbrk{}प\appbrk{}त्त\appbrk{}ङ्ग\appbrk{}रो\appbrk{}ध्रा\appbrk{}ञ्ज\appbrk{}न\appbrk{}य\appbrk{}ष्टि\appbrk{}काह्वैः}\ A}}\skt{~। }\pend
\pstart\versepar \skt{विस्राव्य }\edtext{\skt{चार\-ग्व\-ध\-गो\-जि\-सोमाः}}{\lemma{\skt{चार\appbrk{}ग्व\appbrk{}ध\appbrk{}गो\appbrk{}जि\appbrk{}सोमाः}}\Afootnote{H; \skt{चार\appbrk{}ग्व\appbrk{}धा\appbrk{}गो\appbrk{}जि\appbrk{}सोमाः}\ A}}\skt{ }\edtext{\skt{श्यामा व}}{\lemma{\skt{श्यामा व}}\Afootnote{H; \skt{श्यामा च}\ A}}\skt{ }\edtext{\skt{योज्याः}}{\lemma{\skt{योज्याः}}\Afootnote{H; \skt{योज्या}\ A}}\skt{ कुशलेन }\edtext{\skt{लेपः}}{\lemma{\skt{लेपः}}\Afootnote{H; \skt{लेपे}\ A}}\skt{~॥}\pend

\pstart\versepar \ledsidenote{\passref{4.18.34}}\edtext{\skt{श्यामा\-गि\-रि\-ह्वा\-र\-ज\-नी\-र\-सेषु}}{\lemma{\skt{श्यामा\appbrk{}गि\appbrk{}रि\appbrk{}ह्वा\appbrk{}र\appbrk{}ज\appbrk{}नी\appbrk{}र\appbrk{}सेषु}}\Afootnote{H; \skt{श्यामा\appbrk{}गि\appbrk{}रि\appbrk{}ह्वा\appbrk{}ञ्ज\appbrk{}न\appbrk{}की\appbrk{}र\appbrk{}सेषु}\ A}}\skt{ द्राक्षारसे सप्तलिकारसे च~॥ }\pend
\pstart\versepar \skt{घृतं }\edtext{\edtext{\skt{पचे\-त्क्ली\-त\-क\-जे\-च\-सिद्धं}}{\lemma{\skt{पचेत्}}\Afootnote{H; \skt{पिबेत्}\ A}}}{\lemma{\skt{क्लीत\appbrk{}क\appbrk{}जे\appbrk{}च\appbrk{}सिद्धं}}\Afootnote{H; \skt{क्लीत\appbrk{}क\appbrk{}सं\appbrk{}प्र\appbrk{}सिद्धं}\ A}}\skt{ पित्तार्बुदी }\edtext{\skt{तज्जठरे च}}{\lemma{\skt{तज्जठरे च}}\Afootnote{H\textsuperscript{pc}; \skt{तज्जठरी च}\ A; \skt{तज्जठरे}\ H\textsuperscript{ac}}}\skt{ जन्तुः~॥}\pend

\pstart\versepar \ledsidenote{\passref{4.18.35}}\edtext{\edtext{\skt{शुद्धश्च}}{\lemma{\skt{शुद्धश्}}\Afootnote{H; \skt{शुद्ध}\ A}}}{\lemma{\skt{च}}\Afootnote{H; \skt{स्य}\ A}}\skt{ जन्तोः कफजे }\edtext{\skt{र्बुदे तु}}{\lemma{\skt{र्बुदे तु}}\Afootnote{H; \skt{ऽर्बुदे तु}\ A}}\skt{ }\edtext{\skt{रक्तेऽवसिक्ते च}}{\lemma{\skt{ऽवसिक्ते च}}\Afootnote{H; \skt{ऽवसिक्ते तु}\ A}}\skt{ ततोऽर्बुदन्तत्~। }\pend
\pstart\versepar \skt{}\edtext{\edtext{\skt{कल्कै\-र्द्दि\-हे\-दू\-र्द्ध्व\-म\-धश्च}}{\lemma{\skt{कल्कैर्}}\Afootnote{H; \emph{om.}\ A}}}{\lemma{\skt{द्दिहे\appbrk{}दू\appbrk{}र्द्ध्व\appbrk{}म\appbrk{}धश्च}}\Afootnote{H; \emph{om.}\ A}}\skt{ }\edtext{\skt{यनि}}{\lemma{\skt{यनि}}\Afootnote{H; \emph{om.}\ A}}\skt{ द्रव्याणि }\edtext{\skt{दोषानपचिं}}{\lemma{\skt{दोषानपचिं}}\Afootnote{H; \skt{यान्यूर्ध्वमधश्च दोषान्ह}\ A}}\skt{ }\edtext{\skt{नुदन्ति}}{\lemma{\skt{नुदन्ति}}\Afootnote{H; \skt{रन्ति तैः कल्ककृतैः प्रदिह्यात्}\ A}}\skt{~॥}\pend

\pstart\versepar \ledsidenote{\passref{4.18.36}}\edtext{\skt{कपो\-त\-पा\-रा\-व\-त\-वि\-ड्वि\-मिश्रैः}}{\lemma{\skt{कपो\appbrk{}त\appbrk{}पा\appbrk{}रा\appbrk{}व\appbrk{}त\appbrk{}वि\appbrk{}ड्वि\appbrk{}मिश्रैः}}\Afootnote{H; \skt{कपो\appbrk{}त\appbrk{}पा\appbrk{}रा\appbrk{}व\appbrk{}त\appbrk{}वि\appbrk{}ङ्वि\appbrk{}मिश्रैः}\ A}}\skt{ }\edtext{\skt{सकांशनीली}}{\lemma{\skt{सकांशनीली}}\Afootnote{H; \skt{सकांस्यनीलैः}\ A}}\skt{ शुकलाङ्गलाख्यैः~। }\pend
\pstart\versepar \skt{}\edtext{\skt{मूलैस्तु}}{\lemma{\skt{मूलैस्तु}}\Afootnote{H; \skt{मूत्रैस्तु}\ A}}\skt{ }\edtext{\skt{काका\-द\-नि\-मू\-ल\-मिस्रैः}}{\lemma{\skt{काका\appbrk{}द\appbrk{}नि\appbrk{}मू\appbrk{}ल\appbrk{}मिस्रैः}}\Afootnote{H; \skt{काका\appbrk{}द\appbrk{}नि\appbrk{}मू\appbrk{}ल\appbrk{}मिश्रैः}\ A}}\skt{ क्षार\-प्र\-दि\-ग्धै\-र\-थवा प्रदिह्यात्~॥}\pend

\pstart\versepar \ledsidenote{\passref{4.18.37}}\edtext{\skt{निष्पा\-व\-पि\-ण्या\-क\-कु\-ल\-त्थ\-क\-ल्कै\-र्म्मां\-स\-प्र\-गा\-ढै\-र्द्द\-धि\-म\-र्द्दि\-तैस्तु}}{\lemma{\skt{द्दधिमर्द्दितैस्तु}}\Afootnote{H; \skt{दधिमस्तुयु}\ A}}\skt{~। }\pend
\pstart\versepar \skt{}\edtext{\skt{लेपं}}{\lemma{\skt{लेपं}}\Afootnote{H; \skt{क्तैः~। लेपं}\ A}}\skt{ }\edtext{\skt{विदध्यात्क्रिमयो}}{\lemma{\skt{क्रिमयो}}\Afootnote{H; \skt{कृमयो}\ A}}\skt{ }\edtext{\skt{यथात्र}}{\lemma{\skt{यथात्र}}\Afootnote{H; \skt{यथाऽत्र}\ A}}\skt{ मूर्च्छन्ति }\edtext{\edtext{\skt{मुञ्चन्त्यथ}}{\lemma{\skt{मुञ्चन्त्य्}}\Afootnote{H; \skt{मुञ्चन्त्}\ A}}}{\lemma{\skt{अथ}}\Afootnote{H; \skt{यथ}\ A}}\skt{ }\edtext{\skt{मक्षिका}}{\lemma{\skt{मक्षिका}}\Afootnote{H; \skt{मक्षिकाश्}\ A}}\skt{ }\edtext{\skt{वा}}{\lemma{\skt{वा}}\Afootnote{H; \skt{च}\ A}}\skt{~॥}\pend

\pstart\versepar \ledsidenote{\passref{4.18.38}}\edtext{\skt{अल्पावशिष्टं}}{\lemma{\skt{अल्पावशिष्टं}}\Afootnote{H; \skt{अल्पावशिष्टे}\ A}}\skt{ }\edtext{\edtext{\skt{क्रिमि\-भ\-क्षि\-त\-न्त\-ल्लि\-खे\-त्त\-तोऽ\-ग्नि\-म्वि\-द\-धीत}}{\lemma{\skt{क्रिमिभक्षितन्}}\Afootnote{H; \skt{कृमिभक्षिते}\ A}}}{\lemma{\skt{तल्लिखेत्ततो}}\Afootnote{H; \skt{च लिखेत्ततो}\ A}}\skt{ }\edtext{\skt{पश्चात्}}{\lemma{\skt{पश्चात्}}\Afootnote{H; \skt{पश्चात्}\ A}}\skt{~। }\pend
\pstart\versepar \skt{यदल्पमूलं }\edtext{\edtext{\skt{त्रपु\-शी\-स\-ता\-स्रै\-स्त\-द्वेष्ट्य}}{\lemma{\skt{त्रपु\appbrk{}शी\appbrk{}स\appbrk{}ता\appbrk{}स्रैस्}}\Afootnote{H; \skt{त्रपु\appbrk{}ता\appbrk{}म्र\appbrk{}सी\appbrk{}स\appbrk{}पट्टैः}\ A}}}{\lemma{\skt{तद्वेष्ट्य}}\Afootnote{H; \skt{समावेष्ट्य}\ A}}\skt{ }\edtext{\edtext{\skt{पट्टै\-र\-थ\-वा\-य\-सैर्व्वा}}{\lemma{\skt{पट्टैर्}}\Afootnote{H; \emph{om.}\ A}}}{\lemma{\skt{अथवायसैर्}}\Afootnote{H; \skt{तदायसैर्}\ A}}\skt{~॥}\pend

\pstart\versepar \ledsidenote{\passref{4.18.39}}\edtext{\skt{क्षारा\-ग्नि\-श\-स्त्रा\-ण्य\-व\-चा\-र\-यीत}}{\lemma{\skt{क्षारा\appbrk{}ग्नि\appbrk{}श\appbrk{}स्त्रा\appbrk{}ण्य\appbrk{}व\appbrk{}चा\appbrk{}र\appbrk{}यीत}}\Afootnote{H; \skt{क्षारा\appbrk{}ग्नि\appbrk{}श\appbrk{}स्त्रा\appbrk{}ण्य\appbrk{}स\appbrk{}कृ\appbrk{}द्वि\appbrk{}दध्या}\ A}}\skt{ }\edtext{\edtext{\skt{मुहुर्म्मुहुः}}{\lemma{\skt{मुहुर्}}\Afootnote{H; \emph{om.}\ A}}}{\lemma{\skt{म्मुहुः}}\Afootnote{H; \skt{त्}\ A}}\skt{ }\edtext{\skt{प्राण\-म\-हिं\-स\-मानः}}{\lemma{\skt{प्राण\appbrk{}म\appbrk{}हिं\appbrk{}स\appbrk{}मानः}}\Afootnote{H; \skt{प्राणा\appbrk{}न\appbrk{}हिं\appbrk{}स\appbrk{}न्भि\appbrk{}ष\appbrk{}ग्प्र\appbrk{}मत्तः}\ A}}\skt{~। }\pend
\pstart\versepar \skt{आस्फो\-त\-जा\-ती\-क\-र\-वी\-र\-पत्रैः }\edtext{\edtext{\skt{कषायमिष्टं}}{\lemma{\skt{कषायम्}}\Afootnote{H; \skt{कषाय}\ A}}}{\lemma{\skt{इष्टं}}\Afootnote{H; \skt{मिष्टं}\ A}}\skt{ व्रणशोधनार्थं~॥}\pend

\pstart\versepar \ledsidenote{\passref{4.18.40}}\skt{शुद्धे च तैलं विदधीत }\edtext{\skt{भार्ग्गी\-वि\-ड\-ङ्ग\-पा\-ठा\-त्रि\-फ\-ला\-सु\-सिद्धं}}{\lemma{\skt{भार्ग्गी\appbrk{}वि\appbrk{}ड\appbrk{}ङ्ग\appbrk{}पा\appbrk{}ठा\appbrk{}त्रि\appbrk{}फ\appbrk{}ला\appbrk{}सु\appbrk{}सिद्धं}}\Afootnote{H; \skt{भार्गी\appbrk{}वि\appbrk{}ड\appbrk{}ङ्ग\appbrk{}पा\appbrk{}ठा\appbrk{}त्रि\appbrk{}फ\appbrk{}ला\appbrk{}वि\appbrk{}प\appbrk{}क्वम्}\ A}}\skt{~। }\pend
\pstart\versepar \skt{यदृच्छया चोपगतानि पाकं }\edtext{\skt{क्रमेण}}{\lemma{\skt{क्रमेण}}\Afootnote{H; \skt{पाकक्रमेणो}\ A}}\skt{ }\edtext{\skt{चैव\-ओ\-प\-च\-रे\-द्वि\-धिज्ञः}}{\lemma{\skt{चैव\appbrk{}ओ\appbrk{}प\appbrk{}च\appbrk{}रे\appbrk{}द्वि\appbrk{}धिज्ञः}}\Afootnote{H; \skt{पचरेद्विधिज्ञः}\ A}}\skt{~॥}\pend

\pstart\versepar \ledsidenote{\passref{4.18.41}}\edtext{\edtext{\skt{मेदोऽर्ब्बुदं}}{\lemma{\skt{मेदो}}\Afootnote{H; \skt{मेदो}\ A}}}{\lemma{\skt{ऽर्ब्बुदं}}\Afootnote{H; \skt{र्बुदं}\ A}}\skt{ }\edtext{\skt{स्वेद्य}}{\lemma{\skt{स्वेद्य}}\Afootnote{H; \skt{स्विन्नमथो}\ A}}\skt{ विदार्य }\edtext{\skt{सौध्य}}{\lemma{\skt{सौध्य}}\Afootnote{H; \skt{विशोध्य}\ A}}\skt{ }\edtext{\skt{सीव्याद्यथोक्तं}}{\lemma{\skt{सीव्याद्यथोक्तं}}\Afootnote{H; \skt{सीव्येद्}\ A}}\skt{ }\edtext{\skt{गतरक्तमाशु}}{\lemma{\skt{गतरक्तमाशु}}\Afootnote{H; \skt{गतरक्तमाशु}\ A}}\skt{~। }\pend
\pstart\versepar \skt{ततो }\edtext{\skt{हरि\-द्रा\-गृ\-ह\-धू\-म\-लो\-ध्र\-प\-त्त\-ङ्ग\-चूर्ण्णैः}}{\lemma{\skt{हरि\appbrk{}द्रा\appbrk{}गृ\appbrk{}ह\appbrk{}धू\appbrk{}म\appbrk{}लो\appbrk{}ध्र\appbrk{}प\appbrk{}त्त\appbrk{}ङ्ग\appbrk{}चूर्ण्णैः}}\Afootnote{H; \skt{हरि\appbrk{}द्रा\appbrk{}गृ\appbrk{}ह\appbrk{}धू\appbrk{}म\appbrk{}रो\appbrk{}ध्र\appbrk{}प\appbrk{}त\appbrk{}ङ्ग\appbrk{}चूर्णैः}\ A}}\skt{ समनः शिलालैः~॥}\pend

\pstart\versepar \ledsidenote{\passref{4.18.42}}\edtext{\skt{व्रणम्प्रतिग्राह}}{\lemma{\skt{प्रतिग्राह}}\Afootnote{H; \skt{प्रतिग्राह्य}\ A}}\skt{ मधुप्रगाढैः करञ्जतैलं विदधीत शुद्धे~। }\pend
\pstart\versepar \skt{}\edtext{\skt{सशेषदोषाणि हि योऽर्ब्बुदानि}}{\lemma{\skt{सशेषदोषाणि हि यो}}\Afootnote{A H\textsuperscript{pc}; \skt{सशेषाणि हि यो}\ H\textsuperscript{ac}}}\skt{ करोति }\edtext{\skt{तान्याशु}}{\lemma{\skt{तान्याशु}}\Afootnote{H; \skt{तस्याशु}\ A}}\skt{ }\edtext{\edtext{\skt{पुनर्भ्भवन्ति}}{\lemma{\skt{पुनर्}}\Afootnote{H; \skt{पुन}\ A}}}{\lemma{\skt{भ्भवन्ति}}\Afootnote{H; \skt{र्भवन्ति}\ A}}\skt{~॥}\pend

\pstart\versepar \ledsidenote{\passref{4.18.43}}\skt{तस्मादशेषाणि समुद्धरेत्तु हन्युः सशेषाणि यथा }\edtext{\skt{विषाग्नी}}{\lemma{\skt{विषाग्नी}}\Afootnote{H; \skt{हि वह्निः}\ A}}\skt{~। }\pend
\pstart\versepar \skt{}\edtext{\edtext{\skt{स्विन्नेऽनिलोथं}}{\lemma{\skt{स्विन्ने}}\Afootnote{H; \skt{संस्वे}\ A}}}{\lemma{\skt{ऽनिलोथं}}\Afootnote{H; \skt{द्य गण्डं}\ A}}\skt{ }\edtext{\skt{गलगण्डमदौ}}{\lemma{\skt{गलगण्डमदौ}}\Afootnote{H; \skt{पवनोत्थमादौ}\ A}}\skt{ }\edtext{\skt{नाड्यानि}}{\lemma{\skt{नाड्यानि}}\Afootnote{H; \skt{नाड्याऽनि}\ A}}\skt{ लघ्नौ\-ष\-ध\-प\-त्र\-भङ्गैः~॥}\pend

\pstart\versepar \ledsidenote{\passref{4.18.44}}\edtext{\skt{अम्लेषु}}{\lemma{\skt{अम्लेषु}}\Afootnote{H\textsuperscript{pc}; \skt{अम्लैः स}\ A; \skt{अम्लैः प}\ H\textsuperscript{ac}}}\skt{ }\edtext{\skt{पत्त्रैषु}}{\lemma{\skt{पत्त्रैषु}}\Afootnote{H\textsuperscript{pc}; \skt{मूत्रैर्}\ A; \skt{त्त्रैषु}\ H\textsuperscript{ac}}}\skt{ }\edtext{\skt{पयः}}{\lemma{\skt{पयः}}\Afootnote{H; \emph{om.}\ A}}\skt{ }\edtext{\skt{सुचापि}}{\lemma{\skt{सुचापि}}\Afootnote{H; \emph{om.}\ A}}\skt{ }\edtext{\skt{स्वेदं}}{\lemma{\skt{स्वेदं}}\Afootnote{H; \emph{om.}\ A}}\skt{ }\edtext{\edtext{\skt{विद\-द्ध्या\-त्पि\-शि\-तेषु}}{\lemma{\skt{विदद्ध्यात्}}\Afootnote{H; \skt{विविधैः पयोभिरुष्णैः सतैलैः}\ A}}}{\lemma{\skt{पिशितेषु}}\Afootnote{H; \skt{पिशितैश्च}\ A}}\skt{ }\edtext{\skt{चापि}}{\lemma{\skt{चापि}}\Afootnote{H; \skt{विद्वान्}\ A}}\skt{~। }\pend
\pstart\versepar \skt{}\edtext{\edtext{\skt{विस्रा\-व\-ये\-त्स्वि\-न्न\-म\-त\-न्द्रि\-तस्तु}}{\lemma{\skt{स्विन्नम्}}\Afootnote{H; \skt{स्विन्न}\ A}}}{\lemma{\skt{अतन्द्रितस्तु}}\Afootnote{H; \skt{मतन्द्रियश्च}\ A}}\skt{ }\edtext{\skt{रूढं}}{\lemma{\skt{रूढं}}\Afootnote{H; \skt{शुद्धं}\ A}}\skt{ व्रणं }\edtext{\skt{चाप्यु\-प\-ना\-ह\-यीत}}{\lemma{\skt{उपनाहयीत}}\Afootnote{H; \skt{उपनाहयेत्तु}\ A}}\skt{~॥}\pend

\pstart\versepar \ledsidenote{\passref{4.18.45}}\edtext{\skt{शणा\-त\-सी\-मू\-ल\-के\-शि\-ग्रु\-सिद्ग्रैः}}{\lemma{\skt{शणा\appbrk{}त\appbrk{}सी\appbrk{}मू\appbrk{}ल\appbrk{}के\appbrk{}शि\appbrk{}ग्रु\appbrk{}सिद्ग्रैः}}\Afootnote{H; \skt{शणा\appbrk{}त\appbrk{}सी\appbrk{}मू\appbrk{}ल\appbrk{}क\appbrk{}शि\appbrk{}ग्रु\appbrk{}कि\appbrk{}ण्वप्}\ A}}\skt{ }\edtext{\skt{पिया\-ल\-म\-ज्जा\-स\-हि\-तै\-स्ति\-लैस्तु}}{\lemma{\skt{पिया\appbrk{}ल\appbrk{}म\appbrk{}ज्जा\appbrk{}स\appbrk{}हि\appbrk{}तैस्}}\Afootnote{H; \skt{रिया\appbrk{}ल\appbrk{}म\appbrk{}ज्जा\appbrk{}नु\appbrk{}यु\appbrk{}तैस्}\ A}}\skt{~। }\pend
\pstart\versepar \skt{}\edtext{\skt{काला\-मृ\-ता\-शि\-ग्रु\-पु\-न\-र्न्न\-वा\-र्क्क\-ग\-जा\-दि\-ना\-मा\-क\-र\-घा\-ट\-कुष्ठैः}}{\lemma{\skt{काला\appbrk{}मृ\appbrk{}ता\appbrk{}शि\appbrk{}ग्रु\appbrk{}पु\appbrk{}न\appbrk{}र्न्न\appbrk{}वा\appbrk{}र्क्क\appbrk{}ग\appbrk{}जा\appbrk{}दि\appbrk{}ना\appbrk{}मा\appbrk{}क\appbrk{}र\appbrk{}घा\appbrk{}ट\appbrk{}कुष्ठैः}}\Afootnote{H; \skt{काला\appbrk{}मृ\appbrk{}ता\appbrk{}शि\appbrk{}ग्रु\appbrk{}पु\appbrk{}न\appbrk{}र्न\appbrk{}वा\appbrk{}र्क\appbrk{}ग\appbrk{}जा\appbrk{}दि\appbrk{}ना\appbrk{}मा\appbrk{}क\appbrk{}र\appbrk{}हा\appbrk{}ट\appbrk{}कुष्ठैः}\ A}}\skt{~॥}\pend

\pstart \ledsidenote{\passref{4.18.46}}\skt{एकै\-षि\-का\-वृ\-क्ष\-क\-ति\-ल्व\-कैश्च सुराम्लपिष्टैः }\edtext{\skt{प्रदिहेदभीक्ष्णं}}{\lemma{\skt{प्रदिहेदभीक्ष्णं}}\Afootnote{H; \skt{अस\appbrk{}कृ\appbrk{}त्प्र\appbrk{}दि\appbrk{}ह्यात्}\ A}}\skt{~।}\pend

\pstart\versepar \ledsidenote{\passref{4.18.47}}\skt{तैल\-म्पि\-बे\-च्चा\-मृ\-त\-व\-ल्लि\-नि\-म्ब\-ह\-न्सा\-ह्व\-या\-वृ\-क्ष\-क\-पि\-प्प\-लीभिः~॥ }\pend
\pstart\versepar \skt{सिद्धं बलाभ्यां }\edtext{\skt{सहदेवदारु}}{\lemma{\skt{सहदेवदारु}}\Afootnote{H; \skt{च सदेवदारु}\ A}}\skt{ हिताय नित्यं गलगण्डरोगे~॥}\pend

\pstart\versepar \ledsidenote{\passref{4.18.48}}\skt{स्वेदोपनाहैः कफसम्भवन्तु संस्वेदय विस्रावणमेव कुर्यात्~॥ }\pend
\pstart\versepar \skt{}\edtext{\skt{ततोऽ\-ज\-ग\-न्धा\-ति\-वि\-षा\-वि\-श\-ल्या\-वि\-षा\-णि\-का\-ला\-वु\-शु\-का\-ञ्जा\-नाशा}}{\lemma{\skt{ऽज\appbrk{}ग\appbrk{}न्धा\appbrk{}ति\appbrk{}वि\appbrk{}षा\appbrk{}वि\appbrk{}श\appbrk{}ल्या\appbrk{}वि\appbrk{}षा\appbrk{}णि\appbrk{}का\appbrk{}ला\appbrk{}वु\appbrk{}शु\appbrk{}का\appbrk{}ञ्जा\appbrk{}नाशा}}\Afootnote{H\textsuperscript{pc}; \skt{ऽज\appbrk{}ग\appbrk{}न्धा\appbrk{}ति\appbrk{}वि\appbrk{}षा\appbrk{}वि\appbrk{}श\appbrk{}ल्या\appbrk{}वि\appbrk{}षा\appbrk{}णिका}\ A; \skt{ऽज\appbrk{}ग\appbrk{}न्धा\appbrk{}ति\appbrk{}वि\appbrk{}षा\appbrk{}वि\appbrk{}श\appbrk{}ल्या\appbrk{}वि\appbrk{}षा\appbrk{}णि\appbrk{}का\appbrk{}ला\appbrk{}वु\appbrk{}शुका}\ H\textsuperscript{ac}}}\skt{ }\edtext{\edtext{\skt{केचि\-त्पा\-ठ\-ह्व\-याभिः}}{\lemma{\skt{केचित्}}\Afootnote{H\textsuperscript{pc}; \skt{कुष्ठ}\ A; \emph{om.}\ H\textsuperscript{ac}}}}{\lemma{\skt{पाठह्वयाभिः}}\Afootnote{H\textsuperscript{pc}; \skt{शुकाह्वयाभिः}\ A; \skt{ह्वयाभिः}\ H\textsuperscript{ac}}}\skt{~।}\pend

\pstart\versepar \ledsidenote{\passref{4.18.49}}\edtext{\skt{पला\-श\-भ\-स्मो\-द\-क\-पे\-शि\-ता\-भि\-र्द्दि\-ह्या\-त्सु\-गु\-ञ्जा\-भि\-र\-शी\-त\-लाभिः}}{\lemma{\skt{पला\appbrk{}श\appbrk{}भ\appbrk{}स्मो\appbrk{}द\appbrk{}क\appbrk{}पे\appbrk{}शि\appbrk{}ता\appbrk{}भिर्}}\Afootnote{H\textsuperscript{pc}; \skt{पला\appbrk{}श\appbrk{}भ\appbrk{}स्मो\appbrk{}द\appbrk{}क\appbrk{}पे\appbrk{}षि\appbrk{}ता\appbrk{}भिर्}\ A; \skt{पला\appbrk{}श\appbrk{}भ\appbrk{}स्मो\appbrk{}द\appbrk{}क\appbrk{}पे\appbrk{}शि\appbrk{}भिर्}\ H\textsuperscript{ac}}}\skt{~॥ }\pend
\pstart\versepar \skt{}\edtext{\skt{दशा\-र्द्ध\-श\-ङ्खै\-र्ल्ल\-व\-णैश्च}}{\lemma{\skt{दशार्द्धशङ्खैर्}}\Afootnote{H; \skt{दशार्धसङ्ख्यैर्}\ A}}\skt{ युक्तं तैलं }\edtext{\skt{पिबे\-त्मा\-ग\-धि\-का\-दि\-सिद्धं}}{\lemma{\skt{पिबे\appbrk{}त्मा\appbrk{}ग\appbrk{}धि\appbrk{}का\appbrk{}दि\appbrk{}सिद्धं}}\Afootnote{H; \skt{पिबे\appbrk{}न्मा\appbrk{}ग\appbrk{}धि\appbrk{}का\appbrk{}दि\appbrk{}सि\appbrk{}द्धम्}\ A}}\skt{~।}\pend

\pstart\versepar \ledsidenote{\passref{4.18.50}}\skt{प्रच्छर्द्दनं }\edtext{\skt{शीर्ष\-वि\-रे\-च\-नञ्च}}{\lemma{\skt{शीर्ष\appbrk{}वि\appbrk{}रे\appbrk{}च\appbrk{}नञ्च}}\Afootnote{H; \skt{मूर्धविरेचनं च}\ A}}\skt{ }\edtext{\skt{धूपश्तु वैरेचनिको}}{\lemma{\skt{धूपश्तु वैर्}}\Afootnote{H; \skt{धूमश्च वैर्}\ A}}\skt{ हितस्तु~॥ }\pend
\pstart\versepar \skt{पाकक्रमो वा }\edtext{\skt{प्यनयो}}{\lemma{\skt{प्यनयो}}\Afootnote{H; \skt{ऽपि सदा}\ A}}\skt{ विधेयो वैद्येन }\edtext{\skt{पाकंगतयो}}{\lemma{\skt{पाकंगतयो}}\Afootnote{H; \skt{पाकङ्गतयोः}\ A}}\skt{ कथञ्चित्~।}\pend

\pstart\versepar \ledsidenote{\passref{4.18.51}}\edtext{\skt{कटु\-त्रि\-क\-क्षौ\-द्र\-यु\-ताश्च}}{\lemma{\skt{कटु\appbrk{}त्रि\appbrk{}क\appbrk{}क्षौ\appbrk{}द्र\appbrk{}यु\appbrk{}ताश्च}}\Afootnote{H; \skt{कटु\appbrk{}त्रि\appbrk{}क\appbrk{}क्षौ\appbrk{}द्र\appbrk{}युताः स}\ A}}\skt{ }\edtext{\skt{शस्ता}}{\lemma{\skt{शस्ता}}\Afootnote{H; \skt{मूत्रा}\ A}}\skt{ भक्ष्या यवान्नानि रसाश्च }\edtext{\skt{मौद्गाः}}{\lemma{\skt{मौद्गाः}}\Afootnote{A H\textsuperscript{pc}; \skt{मुद्गाः}\ H\textsuperscript{ac}}}\skt{~। }\pend
\pstart\versepar \skt{सशृंगवेराः सपटोलनिम्बा हिताय देया गलगण्डरोगे~।}\pend

\pstart\versepar \ledsidenote{\passref{4.18.52}}\skt{मेदःसमुत्थे तु यथोपदिष्टां विध्येत्सिरां स्निग्ध\-त\-नो\-र्न्न\-रस्य~॥ }\pend
\pstart\versepar \skt{श्यामा\-सु\-धा\-लो\-ह\-पु\-री\-ष\-द\-न्ती\-र\-सा\-ञ्ज\-नै\-श्चापि हितः }\edtext{\skt{प्रलेपः}}{\lemma{\skt{प्रलेपः}}\Afootnote{H; \skt{प्रदेहः}\ A}}\skt{~।}\pend

\pstart\versepar \ledsidenote{\passref{4.18.53}}\skt{मूत्रेण }\edtext{\skt{चालोड्य}}{\lemma{\skt{चालोड्य}}\Afootnote{H; \skt{वा+आलोड्य}\ A}}\skt{ हिताय }\edtext{\skt{कल्कं}}{\lemma{\skt{कल्कं}}\Afootnote{H; \skt{सारं}\ A}}\skt{ प्रातः }\edtext{\skt{पिबे\-त्सा\-र\-म\-ही\-रु\-हाणां}}{\lemma{\skt{सार\appbrk{}म\appbrk{}ही\appbrk{}रु\appbrk{}हाणां}}\Afootnote{H; \skt{साल\appbrk{}म\appbrk{}ही\appbrk{}रु\appbrk{}हा\appbrk{}णाम्}\ A}}\skt{~॥ }\pend
\pstart\versepar \skt{शस्त्रेण }\edtext{\skt{वाऽपाट्य}}{\lemma{\skt{ऽपाट्य}}\Afootnote{H; \skt{ऽऽआपाट्य}\ A}}\skt{ }\edtext{\skt{विदारिते च}}{\lemma{\skt{विदारिते च}}\Afootnote{H; \skt{विदार्य चैनं}\ A}}\skt{ मेदः समुद्धृत्य हिताय सीव्येत्~।}\pend

\pstart\versepar \ledsidenote{\passref{4.18.54}}\edtext{\edtext{\skt{मज्जा\-व\-सा\-क्ष\-द्र\-यु\-तै\-र्द्द\-हेद्वा}}{\lemma{\skt{मज्जा\appbrk{}व\appbrk{}सा\appbrk{}क्ष\appbrk{}द्र\appbrk{}यु\appbrk{}तैर्}}\Afootnote{H; \skt{मज्जा\appbrk{}ज्य\appbrk{}मे\appbrk{}दो\appbrk{}म\appbrk{}धु\appbrk{}भिर्}\ A}}}{\lemma{\skt{द्दहेद्वा}}\Afootnote{H; \skt{दहेद्वा}\ A}}\skt{ दग्धे च सर्प्पिर्म्मधु चावचार्यं~॥ }\pend
\pstart\versepar \skt{कासीसतुत्थे च ततोऽत्र देये चूर्ण्णीकृते }\edtext{\skt{रोचनया च}}{\lemma{\skt{रोचनया च}}\Afootnote{H; \skt{रोचनया स}\ A}}\skt{ }\edtext{\skt{युक्ते}}{\lemma{\skt{युक्ते}}\Afootnote{H; \skt{मेते}\ A}}\skt{~।}\pend

\pstart\versepar \ledsidenote{\passref{4.18.55}}\skt{तैलेन चाभ्यज्य हिताय दद्यात्सारोद्भवं गोमयजञ्च भस्म~। }\pend
\pstart\versepar \skt{}\edtext{\skt{हितञ्च}}{\lemma{\skt{हितञ्च}}\Afootnote{H; \skt{हितश्च}\ A}}\skt{ }\edtext{\skt{नित्य\-न्त्रि\-फ\-ला\-क\-षायं गा ढश्च}}{\lemma{\skt{त्रिफलाकषायं गा ढश्च}}\Afootnote{H; \skt{त्रिफलाकषायो गाढश्च}\ A}}\skt{ बन्धो }\edtext{\skt{यव\-भो\-ज\-न\-ञ्चेति}}{\lemma{\skt{चेति}}\Afootnote{H; \skt{च}\ A}}\skt{~॥}\pend

\pstart \ledsidenote{\passref{4.18.trailer}}\edtext{\skt{चिकित्सा}}{\lemma{\skt{चिकित्सा}}\Afootnote{H; \skt{इति सुश्रु\appbrk{}त\appbrk{}सं\appbrk{}हि\appbrk{}तायां चिकित्सास्थाने ग्रन्थ्य\appbrk{}प\appbrk{}च्य\appbrk{}र्बु\appbrk{}द\appbrk{}ग\appbrk{}ल\appbrk{}ग\appbrk{}ण्ड\appbrk{}चि\appbrk{}कि\appbrk{}त्सितं नामा\appbrk{}ष्टा\appbrk{}द\appbrk{}शोऽ\appbrk{}ध्यायः}\ A}}\skt{ १९~॥}\pend

\endnumbering
\beginnumbering
\eledsection*{Adhyāya 19}
\pstart \ledsidenote{\passref{4.19.1}}\skt{अथातो }\edtext{\skt{वृद्ध्यु\-प\-दं\-श\-श्ली\-प\-दानां}}{\lemma{\skt{वृद्ध्यु\appbrk{}प\appbrk{}दं\appbrk{}श\appbrk{}श्ली\appbrk{}प\appbrk{}दानां}}\Afootnote{H; \skt{वृद्ध्यु\appbrk{}प\appbrk{}दं\appbrk{}श\appbrk{}श्ली\appbrk{}पद}\ A}}\skt{ चिकित्सितं व्याख्यास्यामः~॥}\pend

\pstart\versepar \ledsidenote{\passref{4.19.3}}\edtext{\edtext{\skt{अन्त्रवृद्धिमृते}}{\lemma{\skt{अन्त्रवृद्धिम्}}\Afootnote{H; \skt{अन्त्रवृद्ध्}\ A}}\skt{ }\edtext{\skt{येऽन्ये}}{\lemma{\skt{ऽन्ये}}\Afootnote{H; \skt{षड्या}\ A}}}{\lemma{\skt{ऋते ये}}\Afootnote{H; \skt{या विना}\ A}}\skt{ }\edtext{\skt{वृद्धयस्तेषु}}{\lemma{\skt{तेषु}}\Afootnote{H; \skt{तासु}\ A}}\skt{ }\edtext{\skt{वर्जयेत्}}{\lemma{\skt{वर्जयेत्}}\Afootnote{H; \skt{वर्जयेत्~। अ}\ A}}\skt{~। }\pend
\pstart\versepar \skt{}\edtext{\skt{वेगाघातं}}{\lemma{\skt{वेगाघातं}}\Afootnote{H; \skt{श्वा}\ A}}\skt{ }\edtext{\skt{पृष्ठयानं}}{\lemma{\skt{पृष्ठयानं}}\Afootnote{H; \skt{दियानं}\ A}}\skt{ व्यायामं मैथुनं }\edtext{\skt{तथा}}{\lemma{\skt{तथा}}\Afootnote{H; \skt{वेगनिग्रहम्}\ A}}\skt{~॥}\pend

\pstart\versepar \ledsidenote{\passref{4.19.4}}\edtext{\skt{अनशनं}}{\lemma{\skt{अनशनं}}\Afootnote{H; \skt{अ}\ A}}\skt{ }\edtext{\edtext{\edtext{\skt{चात्य\-श\-न\-म\-ध्य\-श\-न\-मथो}}{\lemma{\skt{चात्यशनम्}}\Afootnote{H; \skt{त्यासनं}\ A}}}{\lemma{\skt{अध्यशनम्}}\Afootnote{H; \skt{चङ्क्रमणम्}\ A}}}{\lemma{\skt{अथो}}\Afootnote{H; \skt{उपवासं}\ A}}\skt{ }\edtext{\skt{गुरुः}}{\lemma{\skt{गुरुः}}\Afootnote{H; \skt{गुरूणि च}\ A}}\skt{~। }\pend
\pstart\versepar \skt{तत्रादितो वातवृद्धौ }\edtext{\edtext{\skt{त्रिवृ\-ता\-स्नि\-ग्ध\-मा\-तु\-रम्}}{\lemma{\skt{त्रिवृतास्निग्धम्}}\Afootnote{H; \skt{त्रैवृतस्निग्ध}\ A}}}{\lemma{\skt{आतुरम्}}\Afootnote{H; \skt{मातुरम्}\ A}}\skt{~॥}\pend

\pstart\versepar \ledsidenote{\passref{4.19.5}}\edtext{\edtext{\skt{स्निन्नञ्चैव}}{\lemma{\skt{स्निन्नञ्}}\Afootnote{H; \skt{स्विन्नं}\ A}}}{\lemma{\skt{चैव}}\Afootnote{H; \skt{चैनं}\ A}}\skt{ }\edtext{\skt{यथा\-न्या\-य\-म्पा\-य\-यीत}}{\lemma{\skt{पाययीत}}\Afootnote{H; \skt{पाययेत}\ A}}\skt{ विरेचनं~। }\pend
\pstart\versepar \skt{कोशा\-म्र\-ति\-ल्व\-कै\-र\-ण्ड\-फ\-ल\-तै\-लानि वा }\edtext{\skt{नरः}}{\lemma{\skt{नरः}}\Afootnote{H; \skt{नरम्}\ A}}\skt{~॥}\pend

\pstart\versepar \ledsidenote{\passref{4.19.6}}\skt{सक्षीरम्वा }\edtext{\skt{पिबेत्मासं}}{\lemma{\skt{पिबेत्}}\Afootnote{H; \skt{पिबेन्}\ A}}\skt{ }\edtext{\edtext{\skt{तैल\-मे\-र\-ण्ड\-स\-म्भवं}}{\lemma{\skt{तैलम्}}\Afootnote{H; \skt{तैल}\ A}}}{\lemma{\skt{एरण्डसम्भवं}}\Afootnote{H; \skt{मेर\appbrk{}ण्ड\appbrk{}सं\appbrk{}भ\appbrk{}वम्}\ A}}\skt{~। }\pend
\pstart\versepar \skt{ततः कालेऽनिलघ्नानां क्वाथैः कल्कैश्च बुद्धिमान्~॥}\pend

\pstart\versepar \ledsidenote{\passref{4.19.7}}\edtext{\skt{निरू\-ह\-ये\-न्नि\-रू\-ढञ्च}}{\lemma{\skt{निरूढञ्च}}\Afootnote{A H\textsuperscript{pc}; \skt{निरूढश्च}\ H\textsuperscript{ac}}}\skt{ भुक्तवन्तं }\edtext{\skt{रसोदनं}}{\lemma{\skt{रसोदनं}}\Afootnote{H; \skt{रसौदनम्}\ A}}\skt{~। }\pend
\pstart\versepar \skt{यष्टी\-म\-धु\-क\-सि\-द्धेन ततस्तैलेन योजयेत्~॥}\pend

\pstart\versepar \ledsidenote{\passref{4.19.8}}\edtext{\skt{स्वेदोपनाहौ}}{\lemma{\skt{स्वेदोपनाहौ}}\Afootnote{H; \skt{स्नेहोपनाहौ}\ A}}\skt{ कुर्याच्च प्रदे\-हा\-श्चा\-नि\-ला\-प\-हान्~। }\pend
\pstart\versepar \skt{}\edtext{\skt{विद\-ग्ध\-म्पा\-च\-यित्वा}}{\lemma{\skt{विदग्धम्}}\Afootnote{H; \skt{विदग्धां}\ A}}\skt{ वा सेव\-नी\-म्प\-रि\-व\-र्ज्ज\-येत्~॥}\pend

\pstart\versepar \ledsidenote{\passref{4.19.9}}\edtext{\skt{भिन्द्याद्भित्वा च}}{\lemma{\skt{भित्वा च}}\Afootnote{H; \skt{ततः}\ A}}\skt{ }\edtext{\skt{मतिमान्यथोक्तं}}{\lemma{\skt{मतिमान्}}\Afootnote{H; \skt{प्रभिन्नायां}\ A}}\skt{ क्रममाचरेत्~। }\pend
\pstart\versepar \skt{}\edtext{\edtext{\edtext{\skt{पित्त\-ग्र\-न्थि\-क्र\-म\-ङ्कु\-र्या\-द\-प\-क्व\-म्पि\-त्त\-स\-म्भवे}}{\lemma{\skt{पित्त\appbrk{}ग्र\appbrk{}न्थि\appbrk{}क्र\appbrk{}मङ्}}\Afootnote{H; \skt{पित्त\appbrk{}जा\appbrk{}या\appbrk{}म\appbrk{}प\appbrk{}क्वायां पित्तग्रन्थिक्रमो}\ A}}}{\lemma{\skt{कुर्यादपक्वम्}}\Afootnote{H; \emph{om.}\ A}}}{\lemma{\skt{पित्तसम्भवे}}\Afootnote{H; \skt{हितः}\ A}}\skt{~॥}\pend

\pstart\versepar \ledsidenote{\passref{4.19.10}}\edtext{\edtext{\skt{शस्त्रो\-प\-प\-न्न\-म्भित्वा}}{\lemma{\skt{शस्त्रोपपन्नम्}}\Afootnote{H; \skt{पक्वां}\ A}}}{\lemma{\skt{भित्वा}}\Afootnote{H; \skt{वा}\ A}}\skt{ }\edtext{\skt{वा}}{\lemma{\skt{वा}}\Afootnote{H; \skt{भेदयेद्भिन्नां}\ A}}\skt{ }\edtext{\edtext{\skt{शोध\-ये\-न्म\-धु\-स\-र्प्पिषा}}{\lemma{\skt{शोधयेन्}}\Afootnote{H; \skt{शोधयेत्}\ A}}}{\lemma{\skt{मधुसर्प्पिषा}}\Afootnote{H; \skt{क्षौद्रसर्पिषा}\ A}}\skt{~। }\pend
\pstart\versepar \skt{}\edtext{\skt{शुद्धे च}}{\lemma{\skt{शुद्धे च}}\Afootnote{H\textsuperscript{pc}; \skt{शुद्धा}\ A; \skt{शुद्धे चो}\ H\textsuperscript{ac}}}\skt{ }\edtext{\edtext{\skt{चोपशङ्कल्कं}}{\lemma{\skt{चोपशङ्}}\Afootnote{H\textsuperscript{pc}; \skt{यां च भि}\ A; \skt{पशङ्}\ H\textsuperscript{ac}}}}{\lemma{\skt{कल्कं}}\Afootnote{H; \skt{षग्}\ A}}\skt{ }\edtext{\skt{दद्यात्तैलन्तु}}{\lemma{\skt{तैलन्तु}}\Afootnote{H; \skt{तैलं}\ A}}\skt{ }\edtext{\skt{यद्धितं}}{\lemma{\skt{यद्धितं}}\Afootnote{H; \skt{कल्कं च रोपणम्}\ A}}\skt{~॥}\pend

\pstart\versepar \ledsidenote{\passref{4.19.11}}\edtext{\edtext{\edtext{\skt{मुहु\-र्मु\-हु\-र्ज्ज\-लौ\-काभिः}}{\lemma{\skt{मुहुर्}}\Afootnote{H; \emph{om.}\ A}}}{\lemma{\skt{मुहुर्}}\Afootnote{H; \skt{रक्}\ A}}}{\lemma{\skt{ज्जलौकाभिः}}\Afootnote{H; \skt{तजायां जलौकोभिः}\ A}}\skt{ शोणितं }\edtext{\skt{रक्तजे}}{\lemma{\skt{रक्तजे}}\Afootnote{H; \skt{निर्}\ A}}\skt{ हरेत्~। }\pend
\pstart\versepar \skt{}\edtext{\edtext{\skt{पिबे\-द्वि\-रे\-च\-न\-म्वापि}}{\lemma{\skt{पिबेद्}}\Afootnote{H; \skt{भिषक्~। पिबेद्}\ A}}}{\lemma{\skt{वापि}}\Afootnote{H; \skt{वाऽपि}\ A}}\skt{ शर्क्क\-रा\-क्षौ\-द्र\-सं\-युतं~॥}\pend

\pstart\versepar \ledsidenote{\passref{4.19.12}}\skt{पित्तग्रन्थिक्रमं कुर्यादामे पक्वे च }\edtext{\skt{रक्तजे}}{\lemma{\skt{रक्तजे}}\Afootnote{H; \skt{सर्व}\ A}}\skt{~। }\pend
\pstart\versepar \skt{}\edtext{\edtext{\edtext{\skt{श्लेष्म\-वृ\-द्धि\-मु\-ष्ण\-वी\-र्यै\-र्म्मू\-त्र\-पि\-ष्टैस्तु}}{\lemma{\skt{श्लेष्मवृद्धिम्}}\Afootnote{H; \skt{दा~। वृद्धिं}\ A}}}{\lemma{\skt{उष्णवीर्यैर्}}\Afootnote{H; \skt{कफा\appbrk{}त्मि\appbrk{}का\appbrk{}मु\appbrk{}ष्णैर्}\ A}}}{\lemma{\skt{तु}}\Afootnote{H; \skt{प्र}\ A}}\skt{ लेपयेत्~॥}\pend

\pstart\versepar \ledsidenote{\passref{4.19.13}}\skt{पीत\-दा\-रु\-क\-षा\-यञ्च }\edtext{\skt{पिबेत्मूत्रेण}}{\lemma{\skt{पिबेत्}}\Afootnote{H; \skt{पिबेन्}\ A}}\skt{ संयुतं~। }\pend
\pstart\versepar \skt{विम्लापनादृते }\edtext{\skt{चापि}}{\lemma{\skt{चापि}}\Afootnote{H; \skt{वाऽपि}\ A}}\skt{ }\edtext{\skt{श्लेष्मग्रन्थिक्रमे}}{\lemma{\skt{श्लेष्मग्रन्थिक्रमे}}\Afootnote{H; \skt{श्लेष्मग्रन्थिक्रमो}\ A}}\skt{ हितः~॥}\pend

\pstart\versepar \ledsidenote{\passref{4.19.14}}\edtext{\skt{पक्वभिन्ने}}{\lemma{\skt{पक्वभिन्ने}}\Afootnote{H; \skt{पक्वायां च}\ A}}\skt{ }\edtext{\skt{विशुद्ध्यर्थन्तैलं}}{\lemma{\skt{विशुद्ध्यर्थन्}}\Afootnote{H; \skt{विभिन्नायां}\ A}}\skt{ }\edtext{\edtext{\skt{देयम्विजानता}}{\lemma{\skt{देयम्}}\Afootnote{H; \skt{शोधन}\ A}}}{\lemma{\skt{विजानता}}\Afootnote{H; \skt{मिष्यते}\ A}}\skt{~। }\pend
\pstart\versepar \skt{सुम\-ना\-रु\-ष्क\-रा\-ङ्को\-ठ\-स\-प्त\-प\-र्ण्णेषु साधितं~॥}\pend

\pstart\versepar \ledsidenote{\passref{4.19.15}}\edtext{\skt{स्वेद्यम्मेदः}}{\lemma{\skt{स्वेद्यम्}}\Afootnote{H; \emph{om.}\ A}}\skt{ }\edtext{\edtext{\skt{समुत्थन्तु}}{\lemma{\skt{समुत्थन्}}\Afootnote{H; \skt{समुत्थां}\ A}}}{\lemma{\skt{तु}}\Afootnote{H; \skt{संस्वेद्य}\ A}}\skt{ लेप\-ये\-त्सु\-र\-सा\-दिना~। }\pend
\pstart\versepar \skt{शिरो\-वि\-रे\-क\-द्र\-व्यैर्व्वा सुखो\-ष्णै\-र्म्मू\-त्र\-सं\-युतैः~॥}\pend

\pstart\versepar \ledsidenote{\passref{4.19.16}}\edtext{\edtext{\skt{स्विन्नम्वावेष्ट}}{\lemma{\skt{स्विन्नम्}}\Afootnote{H; \skt{स्विन्नां}\ A}}}{\lemma{\skt{वावेष्ट}}\Afootnote{H; \skt{चावेष्ट्य}\ A}}\skt{ पट्टेन }\edtext{\skt{समाश्वास्य च}}{\lemma{\skt{समाश्वास्य च}}\Afootnote{H; \skt{समाश्वास्य तु}\ A}}\skt{ मानवं~। }\pend
\pstart\versepar \skt{}\edtext{\skt{रक्षमूलं}}{\lemma{\skt{रक्षमूलं}}\Afootnote{H; \skt{रक्षन्फले}\ A}}\skt{ सेवनीञ्च वृद्धिपत्त्रेण दारयेत्~॥}\pend

\pstart\versepar \ledsidenote{\passref{4.19.17}}\edtext{\skt{मेदः}}{\lemma{\skt{मेदः}}\Afootnote{H; \skt{मेदस्ततः}\ A}}\skt{ समुद्धृत्य }\edtext{\skt{ततो}}{\lemma{\skt{ततो}}\Afootnote{H; \emph{om.}\ A}}\skt{ दद्या\-त्का\-सी\-स\-सै\-न्धवे~। }\pend
\pstart\versepar \skt{बध्नीयाच्च यथोद्दिष्टं शुद्धे तैलञ्च दापयेत्~॥}\pend

\pstart\versepar \ledsidenote{\passref{4.19.18}}\skt{मनः\-शि\-ला\-ल\-ल\-वणैः सिद्ध\-मा\-रुः\-क\-रेषु च~। }\pend
\pstart\versepar \skt{}\edtext{\skt{स्वेद्य\-मू\-त्र\-नि\-मि\-त्तन्तु}}{\lemma{\skt{स्वेद्य\appbrk{}मू\appbrk{}त्र\appbrk{}नि\appbrk{}मि\appbrk{}त्तन्तु}}\Afootnote{H; \skt{मूत्रजां स्वेदयित्वा तु}\ A}}\skt{ वस्त्रपट्टेन वेष्टयेत्~॥}\pend

\pstart\versepar \ledsidenote{\passref{4.19.19}}\edtext{\edtext{\skt{सेव\-न्या\-म्पा\-र्श्व\-तोऽ\-ध\-स्ता\-द्वि\-ध्ये\-द्व्री\-हि\-मु\-खेन}}{\lemma{\skt{सेवन्याम्}}\Afootnote{H; \skt{सेवन्याः}\ A}}\skt{ वै}}{\lemma{\skt{व्रीहिमुखेन वै}}\Afootnote{H; \skt{व्रीहिमुखेन तु}\ A}}\skt{~। }\pend
\pstart\versepar \skt{}\edtext{\skt{अथा\-त्रो\-द\-र\-व\-न्ना\-डी\-न्दत्त्वा}}{\lemma{\skt{अथा\appbrk{}त्रो\appbrk{}द\appbrk{}र\appbrk{}वन्}}\Afootnote{H; \skt{अथात्र द्विमुखां}\ A}}\skt{ विस्रा\-व\-ये\-द्भि\-षक्~॥}\pend

\pstart\versepar \ledsidenote{\passref{4.19.20}}\edtext{\skt{मूत्र\-ना\-डी\-म\-थो\-द्धृत्य}}{\lemma{\skt{मूत्रनाडीम्}}\Afootnote{H; \skt{मूत्रं नाडीम्}\ A}}\skt{ }\edtext{\skt{स्थवि\-का\-ब\-न्ध\-मा\-च\-रेत्}}{\lemma{\skt{स्थविकाबन्धम्}}\Afootnote{H; \skt{स्थगिकाबन्धम्}\ A}}\skt{~। }\pend
\pstart\versepar \skt{}\edtext{\skt{शुद्धे च}}{\lemma{\skt{शुद्धे च}}\Afootnote{H; \skt{शुद्धायां}\ A}}\skt{ }\edtext{\skt{रोप\-ण\-न्द\-द्या\-द्व\-र्ज्ज\-ये\-द\-न्त्र\-है\-तुकं}}{\lemma{\skt{वर्ज्ज\appbrk{}ये\appbrk{}द\appbrk{}न्त्र\appbrk{}है\appbrk{}तुकं}}\Afootnote{H; \skt{वर्ज\appbrk{}ये\appbrk{}द\appbrk{}न्त्र\appbrk{}है\appbrk{}तु\appbrk{}कीम्}\ A}}\skt{~॥}\pend

\pstart\versepar \ledsidenote{\passref{4.19.21}}\edtext{\skt{अप्रा\-प्त\-फ\-ल\-कोसे वा}}{\lemma{\skt{अप्रा\appbrk{}प्त\appbrk{}फ\appbrk{}ल\appbrk{}कोसे वा}}\Afootnote{H; \skt{अप्रा\appbrk{}प्त\appbrk{}फ\appbrk{}ल\appbrk{}को\appbrk{}षायां}\ A}}\skt{ वातवृद्धिक्रमो हितः~। }\pend
\pstart\versepar \skt{}\edtext{\skt{दहेद्वा}}{\lemma{\skt{दहेद्वा}}\Afootnote{H; \skt{तत्र या वङ्खणस्था तां दहे}\ A}}\skt{ }\edtext{\skt{वह्निना}}{\lemma{\skt{वह्निना}}\Afootnote{H; \emph{om.}\ A}}\skt{ }\edtext{\edtext{\skt{गाढन्तस्य}}{\lemma{\skt{गाढन्}}\Afootnote{H; \emph{om.}\ A}}}{\lemma{\skt{तस्य}}\Afootnote{H; \skt{दर्धेन्दुवत्त्}\ A}}\skt{ }\edtext{\skt{मार्ग्गावरोधना}}{\lemma{\skt{मार्ग्गावरोधना}}\Afootnote{H; \skt{रया}\ A}}\skt{~॥}\pend

\pstart \ledsidenote{\passref{4.19.22}}\edtext{\skt{अङ्गुष्ठमध्ये}}{\lemma{\skt{अङ्गुष्ठमध्ये}}\Afootnote{H; \skt{सम्य\appbrk{}ङ्मा\appbrk{}र्गा\appbrk{}व\appbrk{}रो\appbrk{}धार्थं कोशप्राप्तां तु वर्जयेत्~। त्वचं भित्त्वाऽङ्गुष्ठमध्ये}\ A; \emph{om.}\ K}}\skt{ }\edtext{\edtext{\skt{त्वक्छित्वा}}{\lemma{\skt{त्वक्}}\Afootnote{K H; \emph{om.}\ A}}}{\lemma{\skt{छित्वा}}\Afootnote{K H; \emph{om.}\ A}}\skt{ }\edtext{\skt{दहे\-द्वा\-ङ्ग\-वि\-प\-र्यये}}{\lemma{\skt{वाङ्गविपर्यये}}\Afootnote{K H; \skt{चाङ्ग\appbrk{}वि\appbrk{}प\appbrk{}र्य\appbrk{}यात्}\ A}}\skt{~।}\pend

\pstart\versepar \ledsidenote{\passref{4.19.23}}\skt{अनेनैव }\edtext{\skt{विधाणेन}}{\lemma{\skt{विधाणेन}}\Afootnote{K; \skt{विधानेन}\ A H}}\skt{ वृद्धी }\edtext{\skt{वातकफात्मकौ}}{\lemma{\skt{वातकफात्मकौ}}\Afootnote{K H; \skt{वातकफात्मिके}\ A}}\skt{~। }\pend
\pstart\versepar \skt{प्रदहेत्प्रयतः किन्तु }\edtext{\edtext{\skt{स्नायु\-च्छे\-दोऽ\-धि\-क\-स्तयोः}}{\lemma{\skt{स्नायुच्छेदो}}\Afootnote{A; \skt{च्छेदो}\ K\textsuperscript{ac}; \skt{तस्मा च्छेदो}\ K\textsuperscript{pc}; \skt{स्नुषा च्छेदो}\ H}}}{\lemma{\skt{ऽधिकस्}}\Afootnote{A H; \skt{धिकस्}\ K}}\skt{~।}\pend

\pstart\versepar \ledsidenote{\passref{4.19.24}}\skt{शङ्खोपरि च कर्ण्णान्ते त्यक्त्वा }\edtext{\edtext{\skt{सेव\-नि\-मा\-द\-रात्}}{\lemma{\skt{सेवनिम्}}\Afootnote{K H; \skt{यत्नेन सेवनीम्}\ A}}}{\lemma{\skt{आदरात्}}\Afootnote{K H; \emph{om.}\ A}}\skt{~। }\pend
\pstart\versepar \skt{}\edtext{\edtext{\edtext{\skt{व्यत्या\-सा\-त्त\-त्सि\-रा\-म्वि\-ध्ये\-द\-न्त्र\-वृ\-द्धे\-र्नि\-वृ\-त्तये}}{\lemma{\skt{व्यत्यासात्}}\Afootnote{K; \skt{व्यत्यासाद्}\ A H\textsuperscript{ac}; \skt{व्यत्या गाद्}\ H\textsuperscript{pc}}}}{\lemma{\skt{तत्सिराम्}}\Afootnote{K H; \skt{वा सिरां}\ A}}}{\lemma{\skt{विध्ये\appbrk{}द\appbrk{}न्त्र\appbrk{}वृ\appbrk{}द्धेर्}}\Afootnote{K H; \skt{विध्येदन्त्रवृद्धि}\ A}}\skt{~॥}\pend

\pstart\versepar \ledsidenote{\passref{4.19.25}}\skt{उपदंशेषु साध्येषु }\edtext{\skt{सुस्नि\-ग्ध\-स्वि\-न्न\-दे\-हिनः}}{\lemma{\skt{सुस्नि\appbrk{}ग्ध\appbrk{}स्वि\appbrk{}न्न\appbrk{}दे\appbrk{}हिनः}}\Afootnote{K; \skt{स्निग्धस्विन्नस्य देहिनः}\ A; \skt{सुस्निग्धं स्विन्नदेहिनः}\ H}}\skt{~। }\pend
\pstart\versepar \skt{}\edtext{\skt{सिरा\-म्वि\-ध्ये\-न्मे\-ढ्र\-मध्ये}}{\lemma{\skt{विध्येन्}}\Afootnote{K H; \skt{विध्येन्}\ A}}\skt{ पातयेद्वा जलौकसः~।}\pend

\pstart\versepar \ledsidenote{\passref{4.19.26}}\skt{हरे\-दु\-भ\-य\-त\-श्चापि }\edtext{\edtext{\skt{दोषा\-ण\-त्य\-र्थ\-मु\-च्छ्रितां}}{\lemma{\skt{दोषाणत्यर्थम्}}\Afootnote{K; \skt{दोषानत्यर्थम्}\ A H}}}{\lemma{\skt{उच्छ्रितां}}\Afootnote{K; \skt{उच्छ्रितान्}\ A H}}\skt{~। }\pend
\pstart\versepar \skt{सद्यो }\edtext{\skt{निहृतदोषस्य}}{\lemma{\skt{निहृतदोषस्य}}\Afootnote{K H; \skt{ऽप\appbrk{}हृ\appbrk{}त\appbrk{}दो\appbrk{}षस्य}\ A}}\skt{ }\edtext{\skt{रुक्छो\-था\-वु\-प\-शा\-म्यतः}}{\lemma{\skt{रुक्छोथाव्}}\Afootnote{K H; \skt{रुक्शोफाव्}\ A}}\skt{~।}\pend

\pstart\versepar \ledsidenote{\passref{4.19.27}}\skt{यदि वा दुर्बलो जन्तुर्न वा प्राप्तं विरेचनं~। }\pend
\pstart\versepar \skt{}\edtext{\skt{निरू\-हे\-णा\-ह\-रे\-त्तस्य}}{\lemma{\skt{निरू\appbrk{}हे\appbrk{}णा\appbrk{}ह\appbrk{}रेत्}}\Afootnote{K H; \skt{निरूहेण हरेत्}\ A}}\skt{ }\edtext{\skt{दोषा\-न\-त्य\-र्थ\-मु\-च्छ्रितां}}{\lemma{\skt{उच्छ्रितां}}\Afootnote{K; \skt{उच्छ्रितान्}\ A H}}\skt{~।}\pend

\pstart\versepar \ledsidenote{\passref{4.19.28}}\edtext{\skt{प्रपौण्डरीकं}}{\lemma{\skt{प्रपौण्डरीकं}}\Afootnote{K H; \skt{प्रपौण्डरीक}\ A}}\skt{ }\edtext{\skt{मधुकं}}{\lemma{\skt{मधुकं}}\Afootnote{K H; \skt{यष्ट्}\ A}}\skt{ }\edtext{\skt{रास्ना\-कु\-पु\-न\-र्ण्णवैः}}{\lemma{\skt{रास्ना\appbrk{}कु\appbrk{}पु\appbrk{}न\appbrk{}र्ण्णवैः}}\Afootnote{\skt{याह्व\appbrk{}व\appbrk{}र्षा\appbrk{}भू\appbrk{}कु\appbrk{}ष्ठ\appbrk{}दा\appbrk{}रुभिः}\ A; \skt{रास्ना\appbrk{}कु\appbrk{}ष्ठ\appbrk{}पु\appbrk{}न\appbrk{}र्ण्णवैः}\ K; \skt{रास्ना\appbrk{}कु\appbrk{}ष्ठ\appbrk{}पु\appbrk{}न\appbrk{}र्न्नवैः}\ H}}\skt{~। }\pend
\pstart\versepar \skt{}\edtext{\skt{सर\-ला\-गु\-रु\-भ\-द्रा\-ख्यै\-र्वा\-तजं}}{\lemma{\skt{सर\appbrk{}ला\appbrk{}गु\appbrk{}रु\appbrk{}भ\appbrk{}द्रा\appbrk{}ख्यैर्}}\Afootnote{K H; \skt{सर\appbrk{}ला\appbrk{}गु\appbrk{}रु\appbrk{}रा\appbrk{}स्ना\appbrk{}भिर्}\ A}}\skt{ संप्रलेपयेत्~।}\pend

\pstart\versepar \ledsidenote{\passref{4.19.29}}\skt{निचु\-लै\-र\-ण्ड\-बी\-जानि यव\-गो\-धू\-म\-स\-क्तवः~। }\pend
\pstart\versepar \skt{एतैश्च वातजं }\edtext{\skt{स्निग्धं}}{\lemma{\skt{स्निग्धं}}\Afootnote{K H; \skt{स्निग्धैः}\ A}}\skt{ सुखोष्णैः सम्प्रलेपयेत्~।}\pend

\pstart\versepar \ledsidenote{\passref{4.19.30}}\edtext{\skt{प्रपौण्डरीक}}{\lemma{\skt{प्रपौण्डरीक}}\Afootnote{A K; \skt{प्रपौण्डरीकं}\ H}}\skt{ पूर्वैश्च द्रव्यैः सेकः प्रशस्यते~। }\pend
\pstart\versepar \skt{}\edtext{\skt{गैरि\-का\-ञ्ज\-न\-म\-ञ्जि\-ष्ठा\-म\-धु\-को\-शी\-र\-प\-द्मकैः}}{\lemma{\skt{गैरि\appbrk{}का\appbrk{}ञ्ज\appbrk{}न\appbrk{}म\appbrk{}ञ्जि\appbrk{}ष्ठा\appbrk{}म\appbrk{}धु\appbrk{}को\appbrk{}शी\appbrk{}र\appbrk{}प\appbrk{}द्मकैः}}\Afootnote{K H; \skt{गौरि\appbrk{}का\appbrk{}ञ्ज\appbrk{}न\appbrk{}य\appbrk{}ष्ट्या\appbrk{}ह्वा\appbrk{}रि\appbrk{}वो\appbrk{}शी\appbrk{}र\appbrk{}प\appbrk{}द्मकैः}\ A}}\skt{~।}\pend

\pstart\versepar \ledsidenote{\passref{4.19.31}}\skt{सचन्दनोत्पलैः स्निग्धैः }\edtext{\skt{पैत्तिकां}}{\lemma{\skt{पैत्तिकां}}\Afootnote{K; \skt{पैत्तिकं}\ A H}}\skt{ सम्प्रलेपयेत्~। }\pend
\pstart\versepar \skt{}\edtext{\skt{पद्मो\-त्प\-ल\-मृ\-णा\-लाद्यैः}}{\lemma{\skt{पद्मो\appbrk{}त्प\appbrk{}ल\appbrk{}मृ\appbrk{}णा\appbrk{}लाद्यैः}}\Afootnote{K H\textsuperscript{pc}; \skt{पद्मो\appbrk{}त्प\appbrk{}ल\appbrk{}मृ\appbrk{}णा\appbrk{}लैश्च}\ A; \skt{पद्मो\appbrk{}त्प\appbrk{}ल\appbrk{}मृ\appbrk{}णाद्यैः}\ H\textsuperscript{ac}}}\skt{ सस\-र्जा\-र्जु\-न\-वे\-तसैः~।}\pend

\pstart\versepar \ledsidenote{\passref{4.19.32}}\edtext{\skt{सर्पिष्कैः}}{\lemma{\skt{सर्पिष्कैः}}\Afootnote{K H; \skt{सर्पिःस्निग्धैः}\ A}}\skt{ समधुकैः पैत्तिकं सम्प्रलेपयेत्~। }\pend
\pstart\versepar \skt{}\edtext{\skt{सेचयेच्च}}{\lemma{\skt{सेचयेच्च}}\Afootnote{A K\textsuperscript{pc} H; \skt{सेचये}\ K\textsuperscript{ac}}}\skt{ }\edtext{\skt{घृत\-क्षी\-र\-श\-र्क\-रे\-क्षु\-म\-धू\-दकैः}}{\lemma{\skt{घृत\appbrk{}क्षी\appbrk{}र\appbrk{}श\appbrk{}र्क\appbrk{}रे\appbrk{}क्षु\appbrk{}म\appbrk{}धू\appbrk{}दकैः}}\Afootnote{A K; \skt{घृतक्षीरं शर्क्क\appbrk{}रे\appbrk{}क्षु\appbrk{}म\appbrk{}धू\appbrk{}दकैः}\ H}}\skt{~।}\pend

\pstart\versepar \ledsidenote{\passref{4.19.33}}\edtext{\skt{अथवापि}}{\lemma{\skt{अथवापि}}\Afootnote{K H; \skt{अथवाऽपि}\ A}}\skt{ सुशीतेन कषायेण वटादिना~॥ }\pend
\pstart\versepar \skt{}\edtext{\skt{साला\-ज\-क\-र्ण्णा\-श्व\-क\-र्ण्ण\-ध\-व\-त्वग्भिः}}{\lemma{\skt{साला\appbrk{}ज\appbrk{}क\appbrk{}र्ण्णा\appbrk{}श्व\appbrk{}क\appbrk{}र्ण्ण\appbrk{}ध\appbrk{}व\appbrk{}त्वग्भिः}}\Afootnote{K H; \skt{साला\appbrk{}श्व\appbrk{}क\appbrk{}र्णा\appbrk{}ज\appbrk{}क\appbrk{}र्ण\appbrk{}ध\appbrk{}व\appbrk{}त्वग्भिः}\ A}}\skt{ कफोत्थितं~।}\pend

\pstart\versepar \ledsidenote{\passref{4.19.34}}\edtext{\skt{सुर\-सा\-दि\-तै\-ल\-यु\-क्ताभिः}}{\lemma{\skt{सुर\appbrk{}सा\appbrk{}दि\appbrk{}तै\appbrk{}ल\appbrk{}यु\appbrk{}क्ताभिः}}\Afootnote{K; \skt{सुरापिष्टाभिर्}\ A; \skt{सुर\appbrk{}सा\appbrk{}दि\appbrk{}तै\appbrk{}ल\appbrk{}यु\appbrk{}क्ताभिः सुखो}\ H}}\skt{ }\edtext{\skt{सुखोष्णाभिः}}{\lemma{\skt{सुखोष्णाभिः}}\Afootnote{K; \skt{उष्णाभिः सतैलाभिः}\ A; \skt{ष्णाभिः}\ H}}\skt{ प्रलेपयेत्~। }\pend
\pstart\versepar \skt{आर\-ग्व\-धा\-दि\-क्वा\-थेन }\edtext{\skt{परिषेकञ्च}}{\lemma{\skt{परिषेकञ्च}}\Afootnote{K H; \emph{om.}\ A}}\skt{ }\edtext{\skt{दापयेत्}}{\lemma{\skt{दापयेत्}}\Afootnote{K H; \skt{त्~। रज\appbrk{}न्य\appbrk{}ति\appbrk{}वि\appbrk{}षा\appbrk{}मु\appbrk{}स्ता\appbrk{}सु\appbrk{}रा\appbrk{}सु\appbrk{}र\appbrk{}दा\appbrk{}रुभिः}\ A}}\skt{~।}\edtext{}{\lemma{\skt{प्रलेपयेत्~। आर\appbrk{}ग्व\appbrk{}धा\appbrk{}दि\appbrk{}क्वा\appbrk{}थेन}}\Afootnote{K H; \skt{प्रलेपये}\ A}}\pend

\pstart\versepar \ledsidenote{\passref{4.19.37}}\edtext{\skt{न च}}{\lemma{\skt{न च}}\Afootnote{K H; \skt{न}\ A}}\skt{ }\edtext{\skt{यान्ति}}{\lemma{\skt{यान्ति}}\Afootnote{K H; \skt{याति च}\ A}}\skt{ यथा पाकं प्रयतेत तथा भिषक्~। }\pend
\pstart\versepar \skt{}\edtext{\skt{विदग्धैर्हि}}{\lemma{\skt{विदग्धैर्हि}}\Afootnote{K H; \skt{विदग्धैस्तु}\ A}}\skt{ }\edtext{\skt{सिरा\-स्ना\-यु\-त्व\-ङ्मासैः}}{\lemma{\skt{सिरा\appbrk{}स्ना\appbrk{}यु\appbrk{}त्व\appbrk{}ङ्मासैः}}\Afootnote{K; \skt{सिरा\appbrk{}स्ना\appbrk{}यु\appbrk{}त्व\appbrk{}ङ्मांसैः}\ A H}}\skt{ क्षीयते ध्वजः~।}\pend

\pstart\versepar \ledsidenote{\passref{4.19.36}}\edtext{\skt{तस्मा\-त्सं\-शो\-ध\-ना\-ले\-प\-से\-क\-शो\-णि\-त\-मो\-क्षणैः}}{\lemma{\skt{तस्मात्}}\Afootnote{K H; \skt{एवं}\ A}}\skt{~। }\pend
\pstart\versepar \skt{}\edtext{\edtext{\skt{क्रिया\-यो\-गै\-रु\-प\-च\-रे\-त्प्रा\-क्स्था\-नो\-क्तै\-र्हि\-तै\-रपि}}{\lemma{\skt{क्रियायोगैर्}}\Afootnote{K H; \skt{प्रति\appbrk{}कु\appbrk{}र्या\appbrk{}त्क्रि\appbrk{}या\appbrk{}योगैः}\ A}}}{\lemma{\skt{उपचरेत्}}\Afootnote{K H; \emph{om.}\ A}}\skt{~।}\pend

\pstart\versepar \ledsidenote{\passref{4.19.38}}\edtext{\edtext{\skt{शस्त्रे\-णो\-प\-च\-रे\-द्वापि}}{\lemma{\skt{शस्त्रे\appbrk{}णो\appbrk{}प\appbrk{}च\appbrk{}रेद्}}\Afootnote{K H; \skt{शस्त्रे\appbrk{}णो\appbrk{}प\appbrk{}च\appbrk{}रेच्}\ A}}}{\lemma{\skt{वापि}}\Afootnote{K H; \skt{चापि}\ A}}\skt{ }\edtext{\edtext{\skt{पाक\-मा\-ग\-त\-माशु}}{\lemma{\skt{पाकम्}}\Afootnote{K H; \skt{पाक}\ A}}}{\lemma{\skt{आगतम्}}\Afootnote{K H; \skt{मागत}\ A}}\skt{ वै~। }\pend
\pstart\versepar \skt{}\edtext{\skt{तमपोह्य}}{\lemma{\skt{अपोह्य}}\Afootnote{K H; \skt{दा+अपोह्य}\ A}}\skt{ तिलैः सर्पिः क्षौद्रयुक्तैः प्रलेपयेत्~।}\edtext{}{\lemma{\skt{आशु वै~। तम्}}\Afootnote{K H; \skt{माशु वै। त}\ A}}\pend

\pstart\versepar \ledsidenote{\passref{4.19.39}}\skt{करवीरस्य पत्राणि }\edtext{\edtext{\skt{जात्या\-र\-ग्व\-ध\-यो\-रपि}}{\lemma{\skt{जात्या\appbrk{}र\appbrk{}ग्व\appbrk{}ध\appbrk{}योर्}}\Afootnote{K H; \skt{जात्यारग्वधयो}\ A}}}{\lemma{\skt{अपि}}\Afootnote{K H; \skt{स्तथा}\ A}}\skt{ }\pend
\pstart\versepar \skt{प्रक्षालने प्रयोज्यानि वैज\-य\-न्त्य\-र्क\-यो\-रपि~।}\pend

\pstart\versepar \ledsidenote{\passref{4.19.40}}\edtext{\skt{सौराष्ट्रीं}}{\lemma{\skt{सौराष्ट्रीं}}\Afootnote{A; \skt{सौराष्ट्रि ं}\ K\textsuperscript{ac}; \skt{सौराष्ट्रीं}\ K\textsuperscript{pc}; \skt{सौराष्ट्री}\ H}}\skt{ गैरिकं तुत्थं पुष्प\-का\-सी\-स\-सै\-न्धवं~। }\pend
\pstart\versepar \skt{}\edtext{\skt{लोध्रं}}{\lemma{\skt{लोध्रं}}\Afootnote{K H; \skt{रोध्रं}\ A}}\skt{ }\edtext{\skt{रसाञ्जनञ्चापि}}{\lemma{\skt{चापि}}\Afootnote{K H; \skt{दार्वीं}\ A}}\skt{ हरितालं }\edtext{\skt{मनःशिला}}{\lemma{\skt{मनःशिला}}\Afootnote{K H; \skt{मनःशिलाम्}\ A}}\skt{~।}\pend

\pstart\versepar \ledsidenote{\passref{4.19.41}}\edtext{\skt{हरेणुकैलाश्च}}{\lemma{\skt{हरेणुकैलाश्च}}\Afootnote{K H; \skt{हरेणुकैले च}\ A}}\skt{ तथा }\edtext{\skt{समान्याहृत्य}}{\lemma{\skt{समान्याहृत्य}}\Afootnote{K; \skt{सूक्ष्म}\ A; \skt{सामान्याहृत्य}\ H}}\skt{ }\edtext{\skt{चूर्ण्णयेत्}}{\lemma{\skt{चूर्ण्णयेत्}}\Afootnote{K H; \skt{चूर्णानि कारयेत्}\ A}}\skt{~। }\pend
\pstart\versepar \skt{}\edtext{\skt{तच्चूर्ण्णं}}{\lemma{\skt{तच्चूर्ण्णं}}\Afootnote{A H; \skt{तच्चूर्ण्णं}\ K}}\skt{ क्षौद्रसंयुक्तं उपदंशेषु पूजितं~॥}\pend

\pstart\versepar \ledsidenote{\passref{4.19.42}}\edtext{\skt{जम्ब्वा\-म्र\-सु\-म\-ना\-नि\-म्ब\-श्वे\-त\-का\-म्बो\-जि\-प\-ल्लवां}}{\lemma{\skt{जम्ब्वा\appbrk{}म्र\appbrk{}सु\appbrk{}म\appbrk{}ना\appbrk{}नि\appbrk{}म्ब\appbrk{}श्वे\appbrk{}त\appbrk{}का\appbrk{}म्बो\appbrk{}जि\appbrk{}प\appbrk{}ल्लवां}}\Afootnote{K; \skt{जम्ब्वा\appbrk{}म्र\appbrk{}सु\appbrk{}म\appbrk{}ना\appbrk{}नि\appbrk{}म्ब\appbrk{}श्वे\appbrk{}त\appbrk{}का\appbrk{}म्बो\appbrk{}जि\appbrk{}प\appbrk{}ल्लवाः}\ A; \skt{जम्ब्वा\appbrk{}म्र\appbrk{}सु\appbrk{}म\appbrk{}ना\appbrk{}नि\appbrk{}म्ब\appbrk{}श्वे\appbrk{}त\appbrk{}का\appbrk{}म्बो\appbrk{}जि\appbrk{}प\appbrk{}ल्ल\appbrk{}वान्}\ H}}\skt{~। }\pend
\pstart\versepar \skt{}\edtext{\skt{सल्ल\-की\-ब\-द\-री\-बि\-ल्व\-प\-ला\-श\-ति\-नि\-श\-त्वचः}}{\lemma{\skt{सल्ल\appbrk{}की\appbrk{}ब\appbrk{}द\appbrk{}री\appbrk{}बि\appbrk{}ल्व\appbrk{}प\appbrk{}ला\appbrk{}श\appbrk{}ति\appbrk{}नि\appbrk{}श\appbrk{}त्वचः}}\Afootnote{K H; \skt{शल्ल\appbrk{}की\appbrk{}ब\appbrk{}द\appbrk{}री\appbrk{}बि\appbrk{}ल्व\appbrk{}प\appbrk{}ला\appbrk{}श\appbrk{}ति\appbrk{}नि\appbrk{}श\appbrk{}त्वचः}\ A}}\skt{~।}\pend

\pstart\versepar \ledsidenote{\passref{4.19.43}}\skt{क्षीरिणां च त्वचो }\edtext{\skt{योज्या}}{\lemma{\skt{योज्या}}\Afootnote{K H; \skt{योज्याः}\ A}}\skt{ क्वाथे त्रिफलया सह~। }\pend
\pstart\versepar \skt{}\edtext{\skt{तेन}}{\lemma{\skt{तेन}}\Afootnote{A K\textsuperscript{pc} H; \emph{om.}\ K\textsuperscript{ac}}}\skt{ }\edtext{\skt{क्वाथेन}}{\lemma{\skt{क्वाथेन}}\Afootnote{A K\textsuperscript{pc} H; \emph{om.}\ K\textsuperscript{ac}}}\skt{ }\edtext{\skt{सततं}}{\lemma{\skt{सततं}}\Afootnote{K\textsuperscript{pc} H; \skt{नियतं}\ A; \skt{ततं}\ K\textsuperscript{ac}}}\skt{ }\edtext{\edtext{\skt{प्रण\-म्प्र\-क्षा\-ल\-ये\-द्भि\-षक्}}{\lemma{\skt{प्रणम्}}\Afootnote{K; \skt{व्रणं}\ A H}}}{\lemma{\skt{प्रक्षालयेद्}}\Afootnote{K H; \skt{प्राक्षालयेद्}\ A}}\skt{~।}\pend

\pstart\versepar \ledsidenote{\passref{4.19.44}}\skt{अस्मिन्नेव }\edtext{\skt{कषाये तु}}{\lemma{\skt{कषाये तु}}\Afootnote{A K H\textsuperscript{pc}; \skt{कषायेण}\ Kb; \skt{कषाये तु कषाये तु}\ H\textsuperscript{ac}}}\skt{ तैलं धीरो विपाचयेत्~। }\pend
\pstart\versepar \skt{}\edtext{\skt{गोजी\-वि\-ड\-ङ्ग\-म\-धुक}}{\lemma{\skt{गोजी\appbrk{}वि\appbrk{}ड\appbrk{}ङ्ग\appbrk{}म\appbrk{}धुक}}\Afootnote{K; \skt{गोजी\appbrk{}वि\appbrk{}ड\appbrk{}ङ्ग\appbrk{}य\appbrk{}ष्टीभिः सर्वगन्धैश्च संयुतम्}\ A; \skt{गोजी\appbrk{}वि\appbrk{}ड\appbrk{}ङ्ग\appbrk{}म\appbrk{}धु\appbrk{}क\appbrk{}स\appbrk{}र्व\appbrk{}ग\appbrk{}न्धै\appbrk{}स्स\appbrk{}म\appbrk{}न्वित}\ Kb; \skt{गोजी\appbrk{}वि\appbrk{}डं\appbrk{}ग\appbrk{}म\appbrk{}धु\appbrk{}क\appbrk{}स\appbrk{}र्व्व\appbrk{}गन्धैः समन्वितं}\ H}}\skt{ }\edtext{\skt{समन्वितं}}{\lemma{\skt{समन्वितं}}\Afootnote{K Kb; \emph{om.}\ A H}}\skt{~॥}\pend

\pstart \ledsidenote{\passref{4.19.45ab}}\edtext{\edtext{\skt{एत\-त्स\-र्वो\-प\-दं\-शेषु}}{\lemma{\skt{एतत्}}\Afootnote{A Kb H; \emph{om.}\ K}}}{\lemma{\skt{सर्वोपदंशेषु}}\Afootnote{A Kb H; \emph{om.}\ K}}\skt{ }\edtext{\skt{श्रेष्ठं}}{\lemma{\skt{श्रेष्ठं}}\Afootnote{A Kb H; \emph{om.}\ K}}\skt{ }\edtext{\skt{रोपणमिष्यते}}{\lemma{\skt{रोपणम्}}\Afootnote{A Kb H; \skt{म्}\ K}}\skt{~॥}\pend

\pstart \ledsidenote{\passref{4.19.48cd}}\skt{उपदंशद्वये }\edtext{\skt{शेषे}}{\lemma{\skt{शेषे}}\Afootnote{K Kb H; \skt{अप्येतां}\ A}}\skt{ प्रत्या\-ख्या\-या\-च\-रे\-त्क्रियां~।}\pend

\pstart \ledsidenote{\passref{4.19.49ab}}\edtext{\edtext{\skt{एतेषामेव}}{\lemma{\skt{एतेषाम्}}\Afootnote{K Kb H; \skt{तयोर्}\ A}}}{\lemma{\skt{एव}}\Afootnote{K Kb H; \skt{एव च या}\ A}}\skt{ }\edtext{\skt{योज्यन्तु}}{\lemma{\skt{योज्यन्तु}}\Afootnote{K Kb H; \skt{योग्या}\ A}}\skt{ वीक्ष्य दोषबलाबलं~॥}\pend

\pstart\versepar \ledsidenote{\passref{4.19.52}}\edtext{\skt{स्वेह\-स्वे\-दो\-प\-ना\-हांश्च}}{\lemma{\skt{स्वेह\appbrk{}स्वे\appbrk{}दो\appbrk{}प\appbrk{}ना\appbrk{}हांश्च}}\Afootnote{K Kb; \skt{स्नेह\appbrk{}स्वे\appbrk{}दो\appbrk{}प\appbrk{}पन्ने तु}\ A; \skt{स्नेह\appbrk{}स्वे\appbrk{}दो\appbrk{}प\appbrk{}ना\appbrk{}हाश्च}\ H}}\skt{ श्लीपदे }\edtext{\skt{निलजे}}{\lemma{\skt{निलजे}}\Afootnote{K Kb; \skt{अनिलने}\ A; \skt{ऽनिलजे}\ H}}\skt{ भिषक्~। }\pend
\pstart\versepar \skt{कृत्वा गुल्फोपरि }\edtext{\skt{शिरां}}{\lemma{\skt{शिरां}}\Afootnote{K Kb H; \skt{सिरां}\ A}}\skt{ }\edtext{\skt{व्यध\-ये\-च्च\-तु\-र\-ङ्गुले}}{\lemma{\skt{व्यधयेच्}}\Afootnote{K H; \skt{विध्येत्तु}\ A; \skt{व्याधयेच्}\ Kb}}\skt{~।}\pend

\pstart\versepar \ledsidenote{\passref{4.19.53}}\edtext{\skt{मुपाचरेत्}}{\lemma{\skt{मुपाचरेत्}}\Afootnote{\skt{समा\appbrk{}प्या\appbrk{}यि\appbrk{}त\appbrk{}देहं च बस्तिभिः समुपाचरेत्}\ A; \skt{मुपाचरेत्}\ K; \skt{समा\appbrk{}प्या\appbrk{}यि\appbrk{}त\appbrk{}दे\appbrk{}हश्च बस्तिभिः समुपाचरेत्}\ Kb H}}\skt{~। }\pend
\pstart\versepar \skt{}\edtext{\edtext{\skt{मास\-मे\-र\-ण्ड\-तै\-लञ्च}}{\lemma{\skt{मासम्}}\Afootnote{K Kb H; \skt{मास}\ A}}}{\lemma{\skt{एरण्डतैलञ्च}}\Afootnote{K Kb H; \skt{मेरण्डजं तैलं}\ A}}\skt{ }\edtext{\skt{पिबेन्मूत्रेण}}{\lemma{\skt{पिबेन्}}\Afootnote{K Kb H; \skt{पिबेन्}\ A}}\skt{ }\edtext{\skt{मानवः}}{\lemma{\skt{मानवः}}\Afootnote{K Kb H; \skt{संयुतम्}\ A}}\skt{~।}\pend

\pstart\versepar \ledsidenote{\passref{4.19.54}}\edtext{\skt{महौ\-ष\-ध\-वि\-प\-क्वेन}}{\lemma{\skt{महौ\appbrk{}ष\appbrk{}ध\appbrk{}वि\appbrk{}प\appbrk{}क्वेन}}\Afootnote{K Kb H; \emph{om.}\ A}}\skt{ }\edtext{\skt{पयसा}}{\lemma{\skt{पयसा}}\Afootnote{K Kb H; \skt{पयसौ}\ A}}\skt{ }\edtext{\edtext{\skt{चान्नमादिशेत्}}{\lemma{\skt{चान्नम्}}\Afootnote{K Kb H; \skt{दन\appbrk{}म\appbrk{}श्नी\appbrk{}या\appbrk{}न्नाग}\ A}}}{\lemma{\skt{आदिशेत्}}\Afootnote{K H; \skt{रक्वथितेन च}\ A; \skt{आदिशे}\ Kb}}\skt{~। }\pend
\pstart\versepar \skt{}\edtext{\skt{त्रिवृतं}}{\lemma{\skt{त्रिवृतं}}\Afootnote{K H; \skt{तैवृतं}\ A; \skt{त्~। तृवृतं}\ Kb}}\skt{ चोपयुञ्जीत }\edtext{\edtext{\edtext{\skt{हित\-श्चा\-त्रा\-ग्नि\-रु\-च्यते}}{\lemma{\skt{हितश्}}\Afootnote{K Kb H; \emph{om.}\ A}}}{\lemma{\skt{चात्राग्निर्}}\Afootnote{K Kb H; \emph{om.}\ A}}}{\lemma{\skt{उच्यते}}\Afootnote{K Kb H; \skt{शस्तो दाह\appbrk{}स्त\appbrk{}था+\appbrk{}अ\appbrk{}ग्निना}\ A}}\skt{~॥}\pend

\pstart\versepar \ledsidenote{\passref{4.19.55}}\skt{गुल्फस्याधः }\edtext{\skt{शिरां}}{\lemma{\skt{शिरां}}\Afootnote{K Kb H; \skt{सिरां}\ A}}\skt{ विध्येच्छ्लीपदे }\edtext{\skt{पित्तस}}{\lemma{\skt{पित्तस}}\Afootnote{\skt{पित्तसंभवे। पीतघ्नीं च क्रियां कुर्या\appbrk{}त्पि\appbrk{}त्ता\appbrk{}र्बु\appbrk{}द\appbrk{}वि\appbrk{}स\appbrk{}र्प\appbrk{}वत्}\ A; \skt{पित्तसष्ठवे~। पित्तघ्नाश्च क्रियाः …}\ K; \skt{पित्तसम्भवे पित्तघ्नाश्च क्रियाः कुर्या\appbrk{}त्पि\appbrk{}त्ता\appbrk{}र्बु\appbrk{}द\appbrk{}वि\appbrk{}स\appbrk{}र्प\appbrk{}वत्}\ Kb H}}\skt{~। }\pend
\pstart\versepar \skt{~।}\pend

\pstart\versepar \ledsidenote{\passref{4.19.56}}\edtext{\skt{तां}}{\lemma{\skt{तां}}\Afootnote{K; \skt{सिरां सुविदितां}\ A; \skt{सिरां सुविदिता}\ Kb H}}\skt{ विध्येदंगुष्ठे }\edtext{\skt{श्लेष्मश्लीपदे}}{\lemma{\skt{श्लेष्मश्लीपदे}}\Afootnote{K Kb H; \skt{श्लैष्मिके भिषक्}\ A}}\skt{~। }\pend
\pstart\versepar \skt{मधुयुक्तानि }\edtext{\skt{चाभीक्ष्णं}}{\lemma{\skt{चाभीक्ष्णं}}\Afootnote{A K H; \skt{चाभीक्ष्ण}\ Kb}}\skt{ कषायाणि }\edtext{\skt{पिबेन्नरः}}{\lemma{\skt{पिबेन्}}\Afootnote{A K Kb; \skt{पिबेत्}\ H}}\skt{~।}\pend

\pstart\versepar \ledsidenote{\passref{4.19.57}}\skt{पिबेद्वा }\edtext{\edtext{\skt{प्यभयाकल्कं}}{\lemma{\skt{प्य्}}\Afootnote{K Kb; \skt{अप्}\ A; \skt{ऽप्य्}\ H}}}{\lemma{\skt{अभयाकल्कं}}\Afootnote{K Kb H; \skt{यभयाकल्कं}\ A}}\skt{ }\edtext{\skt{मूत्रे\-णा\-न्य\-त\-मेन वै}}{\lemma{\skt{मूत्रे\appbrk{}णा\appbrk{}न्य\appbrk{}त\appbrk{}मेन वै}}\Afootnote{K Kb H; \skt{मूत्रे\appbrk{}णा\appbrk{}न्य\appbrk{}त\appbrk{}मेन}\ A}}\skt{~। }\pend
\pstart\versepar \skt{}\edtext{\edtext{\skt{पिबे\-दे\-द\-ङ्गु\-डू\-चीम्वा}}{\lemma{\skt{पिबेदेदङ्}}\Afootnote{K; \emph{om.}\ A; \skt{पिबेदिदं}\ Kb; \skt{पिबेदेनङ्}\ H}}}{\lemma{\skt{गुडूचीम्वा}}\Afootnote{K H; \emph{om.}\ A; \skt{गुडुचीम्वा}\ Kb}}\skt{ }\edtext{\edtext{\skt{नाग\-र\-म्भ\-द्र\-दारु}}{\lemma{\skt{नागरम्}}\Afootnote{K Kb H; \emph{om.}\ A}}\skt{ वा}}{\lemma{\skt{भद्रदारु वा}}\Afootnote{K Kb H; \skt{च। कटुकाममृतां शुण्ठीं विडङ्गं दारु चित्रकम्}\ A}}\skt{~।}\pend

\pstart\versepar \ledsidenote{\passref{4.19.58}}\edtext{\edtext{\skt{हितश्चालेपणं}}{\lemma{\skt{हितश्}}\Afootnote{K Kb H; \skt{हितं}\ A}}\skt{ चित्रको}}{\lemma{\skt{चालेपणं चित्रको}}\Afootnote{\skt{वा लेपने नित्यं भद्रदारु सचित्रकम्}\ A; \skt{चालेपणं वित्यं चित्रको}\ K; \skt{चालेपणन्नित्यं चित्रको}\ Kb; \skt{चाले\appbrk{}प\appbrk{}न\appbrk{}न्नि\appbrk{}त्य\appbrk{}ञ्चि\appbrk{}त्रको}\ H}}\skt{ }\edtext{\skt{देवदारु च}}{\lemma{\skt{देवदारु च}}\Afootnote{K Kb H; \emph{om.}\ A}}\skt{~। }\pend
\pstart\versepar \skt{}\edtext{\skt{वि}}{\lemma{\skt{वि}}\Afootnote{\skt{विड\appbrk{}ङ्ग\appbrk{}म\appbrk{}रि\appbrk{}चा\appbrk{}र्केषु नागरे चित्रकेअथवा}\ A H; \skt{विडि\appbrk{}ङ्ग\appbrk{}म\appbrk{}रि\appbrk{}चा\appbrk{}र्केषु नागरे चित्रके थवा~।}\ K; \skt{विड\appbrk{}ङ्ग\appbrk{}म\appbrk{}रि\appbrk{}चा\appbrk{}केषु नागरचित्रके थवा}\ Kb}}\pend

\pstart\versepar \ledsidenote{\passref{4.19.59}}\edtext{\skt{वणेषु च}}{\lemma{\skt{वणेषु च}}\Afootnote{K; \skt{भद्र\appbrk{}दा\appbrk{}र्वे\appbrk{}लु\appbrk{}काख्ये च सर्वेषु लवणेषु च}\ A; \skt{भद्र\appbrk{}दा\appbrk{}र्वे\appbrk{}ल\appbrk{}का\appbrk{}क्षश्च सर्वेषु लवणेषु च}\ Kb H}}\skt{~। }\pend
\pstart\versepar \skt{तैलं पक्वं }\edtext{\skt{पिबेद्वापि}}{\lemma{\skt{वापि}}\Afootnote{K Kb H; \skt{वा+अपि}\ A}}\skt{ यवान्नञ्च हितं सदा~।}\pend

\pstart\versepar \ledsidenote{\passref{4.19.60}}\edtext{\skt{पिबेत्सदा च}}{\lemma{\skt{सदा च}}\Afootnote{K Kb H; \skt{सर्षप}\ A}}\skt{ }\edtext{\skt{तैलेन}}{\lemma{\skt{तैलेन}}\Afootnote{K Kb H; \skt{तैलं वा}\ A}}\skt{ श्लीप\-दा\-ना\-न्नि\-वृ\-त्तये~। }\pend
\pstart\versepar \skt{}\edtext{\skt{पूती\-क\-र\-ञ्ज\-च्छ\-दजं}}{\lemma{\skt{पूती\appbrk{}क\appbrk{}र\appbrk{}ञ्ज\appbrk{}च्छ\appbrk{}दजं}}\Afootnote{K H; \skt{पूती\appbrk{}क\appbrk{}र\appbrk{}ञ्ज\appbrk{}प\appbrk{}त्राणां}\ A; \skt{पूती\appbrk{}क\appbrk{}र\appbrk{}ञ्जा\appbrk{}च्छ\appbrk{}दजं}\ Kb}}\skt{ }\edtext{\skt{रसञ्चापि}}{\lemma{\skt{चापि}}\Afootnote{K Kb H; \skt{वा+अपि}\ A}}\skt{ यथाबलं~।}\pend

\pstart\versepar \ledsidenote{\passref{4.19.61}}\skt{अनेनैव विधानेन पुत्रं जीवकजं~॥ }\pend
\pstart\versepar \skt{}\edtext{\skt{प्रयुञ्जीत}}{\lemma{\skt{प्रयुञ्जीत}}\Afootnote{\skt{रसम्प्रयुञ्जीत}\ A K Kb H}}\skt{ भिषक्प्राज्ञः }\edtext{\skt{कालसंगवित्}}{\lemma{\skt{कालसंगवित्}}\Afootnote{\skt{काल\appbrk{}सा\appbrk{}त्म्य\appbrk{}वि\appbrk{}भा\appbrk{}ग\appbrk{}वित्}\ A H; \skt{काल\appbrk{}सं\appbrk{}स्य\appbrk{}क्वि\appbrk{}भा\appbrk{}ग\appbrk{}वित्}\ K; \skt{काल\appbrk{}स\appbrk{}म्य\appbrk{}ग्वि\appbrk{}भा\appbrk{}ग\appbrk{}वित्}\ Kb}}\skt{~॥}\pend

\pstart\versepar \ledsidenote{\passref{4.19.63}}\edtext{\skt{कण्टकारिकां}}{\lemma{\skt{कण्टकारिकां}}\Afootnote{\skt{काकादनीं काकजङ्घां बृहतीं कण्टकारिकाम्}\ A; \skt{काकादनीः काक\appbrk{}जं\appbrk{}घा\appbrk{}म्बृ\appbrk{}हतीं कण्टकारिकां}\ K; \skt{काकादनी काकजंघा बृहतीं कण्ठकारिकां}\ Kb; \skt{काकादनीः काक\appbrk{}जं\appbrk{}घा\appbrk{}म्बृ\appbrk{}ह\appbrk{}ती\appbrk{}ङ्क\appbrk{}ण्ट\appbrk{}का\appbrk{}रि\appbrk{}कान्}\ H}}\skt{~। }\pend
\pstart\versepar \skt{}\edtext{\skt{कद\-म्ब\-पु\-ष्पी\-मा\-न्दारीं}}{\lemma{\skt{आन्दारीं}}\Afootnote{K Kb; \skt{मन्दारीं}\ A; \skt{मान्दारीं}\ H}}\skt{ }\edtext{\skt{था}}{\lemma{\skt{था}}\Afootnote{\skt{लम्बां शुकनसां तथा}\ A K H; \skt{लम्बा शुकनसान्तथा}\ Kb}}\skt{~।}\pend

\pstart\versepar \ledsidenote{\passref{4.19.64}}\skt{दग्ध्वा }\edtext{\skt{मूत्रेण तं}}{\lemma{\skt{मूत्रेण तं}}\Afootnote{K Kb H; \skt{मूत्रेण त}\ A}}\skt{ }\edtext{\skt{भस्म}}{\lemma{\skt{भस्म}}\Afootnote{K Kb H; \skt{द्भस्म}\ A}}\skt{ }\edtext{\skt{स्राव\-ये\-त्क्षा\-र\-क\-ल्क\-वत्}}{\lemma{\skt{क्षारकल्कवत्}}\Afootnote{K Kb H; \skt{क्षारकल्पवत्}\ A}}\skt{~। }\pend
\pstart\versepar \skt{तत्र दद्यात्प्रतीवापं }\edtext{\skt{काको\-दु\-म्ब\-रि\-का\-रसं}}{\lemma{\skt{काको\appbrk{}दु\appbrk{}म्ब\appbrk{}रि\appbrk{}का\appbrk{}रसं}}\Afootnote{A K H; \skt{काको\appbrk{}दु\appbrk{}रि\appbrk{}का\appbrk{}रसं}\ Kb}}\skt{~।}\pend

\pstart\versepar \ledsidenote{\passref{4.19.65}}\skt{मदनाच्च }\edtext{\skt{फलं}}{\lemma{\skt{फलं}}\Afootnote{K Kb H; \skt{फलात्}\ A}}\skt{ }\edtext{\skt{क्वाथं}}{\lemma{\skt{क्वाथं}}\Afootnote{A Kb; \skt{क्}\ K H}}\skt{ }\edtext{\skt{स्वरसञ्च}}{\lemma{\skt{स्वरसञ्च}}\Afootnote{Kb; \skt{शुकाख्यस्वरसं}\ A; \skt{वाथं स्वरसञ्च}\ K H}}\skt{ }\edtext{\skt{शुकाख्यया}}{\lemma{\skt{शुकाख्यया}}\Afootnote{K Kb H; \skt{तथा}\ A}}\skt{~। }\pend
\pstart\versepar \skt{एष क्षारस्तु पानीयः श्लीपदं }\edtext{\skt{हन्ति}}{\lemma{\skt{हन्ति}}\Afootnote{\skt{हन्ति सेवितः}\ A K Kb H}}\skt{~।}\pend

\pstart\versepar \ledsidenote{\passref{4.19.66}}\edtext{\skt{अपचिं}}{\lemma{\skt{अपचिं}}\Afootnote{K Kb; \skt{अपचीं}\ A H}}\skt{ गलगण्डञ्च ग्रह\-णी\-दो\-ष\-मेव च~। }\pend
\pstart\versepar \skt{}\edtext{\skt{अभ\-क्त\-रो\-च\-क\-ञ्चैव}}{\lemma{\skt{अभ\appbrk{}क्त\appbrk{}रो\appbrk{}च\appbrk{}कञ्}}\Afootnote{K Kb H; \skt{भक्तस्यानशनं}\ A}}\skt{ }\edtext{\skt{हन्या\-त्स\-र्व\-वि\-षाणि}}{\lemma{\skt{सर्वविषाणि}}\Afootnote{A K H; \skt{सर्वं विषाणि}\ Kb}}\skt{ च~।}\pend

\pstart\versepar \ledsidenote{\passref{4.19.67}}\edtext{\edtext{\skt{एष्वेव}}{\lemma{\skt{एष्व्}}\Afootnote{K Kb H; \skt{एष्}\ A}}}{\lemma{\skt{एव}}\Afootnote{K Kb H; \skt{वेव}\ A}}\skt{ }\edtext{\skt{सिद्धं}}{\lemma{\skt{सिद्धं}}\Afootnote{K Kb H; \skt{तैलं}\ A}}\skt{ }\edtext{\edtext{\skt{तैलञ्च}}{\lemma{\skt{तैलञ्}}\Afootnote{K Kb H; \skt{सं}\ A}}}{\lemma{\skt{च}}\Afootnote{K Kb H; \skt{सिद्धं}\ A}}\skt{ नस्याभ्यङ्गेषु }\edtext{\skt{सेवितं}}{\lemma{\skt{सेवितं}}\Afootnote{K Kb H; \skt{पूजितम्}\ A}}\skt{~। }\pend
\pstart\versepar \skt{}\edtext{\skt{एतानेवामयां}}{\lemma{\skt{एतानेवामयां}}\Afootnote{K Kb; \skt{एता\appbrk{}ने\appbrk{}वा\appbrk{}म\appbrk{}यान्}\ A H}}\skt{ हन्ति ये च दुष्टव्रणा नृणां~।}\pend

\pstart\versepar \ledsidenote{\passref{4.19.68}}\edtext{\skt{द्रन्ती}}{\lemma{\skt{द्रन्ती}}\Afootnote{K Kb; \emph{om.}\ A; \skt{दन्ती}\ H}}\skt{ }\edtext{\skt{द्रवन्ती}}{\lemma{\skt{द्रवन्ती}}\Afootnote{K Kb H; \skt{द्रवन्तीं}\ A}}\skt{ }\edtext{\skt{त्रिवृता}}{\lemma{\skt{त्रिवृता}}\Afootnote{K Kb; \skt{त्रिवृतां दन्तीं}\ A; \skt{तृवृता}\ H}}\skt{ }\edtext{\skt{नीली}}{\lemma{\skt{नीली}}\Afootnote{K Kb H; \skt{नीलीं}\ A}}\skt{ }\edtext{\skt{श्यामा}}{\lemma{\skt{श्यामा}}\Afootnote{K Kb H; \skt{श्यामां}\ A}}\skt{ तथैव च~। }\pend
\pstart\versepar \skt{}\edtext{\skt{सप्तलां}}{\lemma{\skt{सप्तलां}}\Afootnote{A K Kb; \skt{सप्तला}\ H}}\skt{ शङ्खिनीं चैव दग्ध्वा मूत्रेण गालयेत्~।}\pend

\pstart\versepar \ledsidenote{\passref{4.19.69}}\skt{दद्याच्च }\edtext{\skt{त्रिफलारस्तु}}{\lemma{\skt{त्रिफलारस्तु}}\Afootnote{\skt{त्रिफ\appbrk{}ला\appbrk{}क्वा\appbrk{}थ\appbrk{}मेष क्षारस्तु}\ A; \skt{त्रिफ\appbrk{}ला\appbrk{}क्वा\appbrk{}थ\appbrk{}मेष क्षारस्तु}\ K Kb H}}\skt{ साधितः~। }\pend
\pstart\versepar \skt{}\edtext{\skt{श्चाप्याशिषः}}{\lemma{\skt{चाप्य्}}\Afootnote{\skt{पूर्वैश्चाप्य्}\ A; \skt{पूर्वश्चाप्य्}\ K Kb H}}\skt{ समाः~।}\edtext{}{\lemma{\skt{साधितः~। श्}}\Afootnote{\skt{साधितः अधो गच्छति पीतस्तु}\ A Kb; \skt{साधितः अधो गच्छति पीतश्च}\ K H}}\pend

\pstart\versepar \ledsidenote{\passref{4.19.69.add1}}\skt{क्रियाभिः सम्यगारब्धा यदि नैवोपशाम्यति~। }\pend
\pstart\versepar \skt{स्वे स्वे भाग्ये }\edtext{\skt{समुदयित्वा}}{\lemma{\skt{समुदयित्वा}}\Afootnote{\skt{समुत्पाद्य च्छेदयित्वा}\ K Kb H}}\skt{ }\edtext{\skt{सिराम्व्यधेत्}}{\lemma{\skt{व्यधेत्}}\Afootnote{K H; \skt{व्यध्येत्}\ Kb}}\skt{~।}\pend

\pstart\versepar \ledsidenote{\passref{4.19.69.add2}}\skt{सिरामुखे च मतिमां दहे\-त्स\-म्य\-क्छ\-ला\-कया~॥ }\pend
\pstart\versepar \skt{कुर्या\-च्छी\-त\-क्रि\-या\-ञ्चापि व्रणोक्तञ्च विधिन्तथेति~॥}\pend

\pstart \ledsidenote{\passref{4.19.trailer}}\edtext{\skt{चि}}{\lemma{\skt{चि}}\Afootnote{K Kb; \skt{इति सुश्रु\appbrk{}त\appbrk{}सं\appbrk{}हि\appbrk{}तायां चिकित्सितस्थाने वृद्ध्यु\appbrk{}प\appbrk{}दं\appbrk{}श\appbrk{}श्लीपः चिकित्सितं नामै\appbrk{}को\appbrk{}न\appbrk{}विं\appbrk{}शोऽ\appbrk{}ध्यायः}\ A; \skt{चिकित्सा एको\appbrk{}न\appbrk{}विं\appbrk{}श\appbrk{}ति\appbrk{}त\appbrk{}मोऽ\appbrk{}ध्यायः}\ H}}\skt{ १९~॥}\pend

\endnumbering
\beginnumbering
\eledsection*{Adhyāya 20}
\pstart \ledsidenote{\passref{4.20.1}}\skt{अथातः }\edtext{\skt{क्षुद्ररोगाणा}}{\lemma{\skt{क्षुद्ररोगाणा}}\Afootnote{K; \skt{क्षुद्ररोग}\ A; \skt{क्षुद्ररोगाणाञ्}\ Kb; \skt{क्षुद्ररोगानां}\ H}}\skt{ }\edtext{\skt{चिकि\-त्सि\-त\-म्व्या\-ख्या\-स्यामः}}{\lemma{\skt{व्याख्यास्यामः}}\Afootnote{A Kb H; \skt{व्या व्याख्यास्यामः}\ K}}\skt{~॥}\pend

\pstart\versepar \ledsidenote{\passref{4.20.3}}\edtext{\skt{तत्रादित}}{\lemma{\skt{तत्रादित}}\Afootnote{K Kb H; \skt{तत्रा}\ A}}\skt{ }\edtext{\skt{एव}}{\lemma{\skt{एव}}\Afootnote{K Kb H; \emph{om.}\ A}}\skt{ }\edtext{\edtext{\skt{जग\-ञ्ज\-ला\-यु\-का\-भि\-र्ग्रा\-ह\-येत्}}{\lemma{\skt{जग\appbrk{}ञ्ज\appbrk{}ला\appbrk{}यु\appbrk{}का\appbrk{}भिर्}}\Afootnote{\skt{जगल्लिकामामां जलौकोभिर्}\ A; \skt{जग\appbrk{}ल्लि\appbrk{}का\appbrk{}म\appbrk{}प\appbrk{}क्वा\appbrk{}ञ्ज\appbrk{}ला\appbrk{}यु\appbrk{}का\appbrk{}भिर्}\ K Kb H\textsuperscript{pc}; \skt{जि गल्लि\appbrk{}का\appbrk{}म\appbrk{}प\appbrk{}क्वा\appbrk{}ञ्ज\appbrk{}ला\appbrk{}यु\appbrk{}का\appbrk{}भिर्}\ H\textsuperscript{ac}}}}{\lemma{\skt{ग्राहयेत्}}\Afootnote{K Kb H; \skt{उपाचरेत्। शुक्तिश्रुघ्}\ A}}\skt{~। }\pend
\pstart\versepar \skt{}\edtext{\edtext{\skt{कल्कै\-श्चा\-ले\-प\-येत्}}{\lemma{\skt{कल्कैश्}}\Afootnote{\skt{नीय\appbrk{}व\appbrk{}क्षा\appbrk{}र\appbrk{}क\appbrk{}ल्कैश्}\ A; \skt{यव\appbrk{}क्षा\appbrk{}र\appbrk{}सु\appbrk{}रा\appbrk{}ष्ट्रि\appbrk{}का\appbrk{}का\appbrk{}को\appbrk{}ली\appbrk{}क\appbrk{}ल्कैश्}\ K Kb H}}}{\lemma{\skt{चालेपयेत्}}\Afootnote{K Kb H; \skt{चाले\appbrk{}प\appbrk{}ये\appbrk{}द्भि\appbrk{}षक्}\ A}}\skt{~।}\pend

\pstart \ledsidenote{\passref{4.20.4}}\edtext{\skt{श्यामा\-पा\-ठा\-ला\-ङ्ग\-ली\-क\-ल्कैर्वा}}{\lemma{\skt{श्यामा\appbrk{}पा\appbrk{}ठा\appbrk{}ला\appbrk{}ङ्ग\appbrk{}ली\appbrk{}क\appbrk{}ल्कैर्वा}}\Afootnote{K Kb H; \skt{श्यामा\appbrk{}ला\appbrk{}ङ्ग\appbrk{}ल\appbrk{}कि\appbrk{}पा\appbrk{}ठा\appbrk{}क\appbrk{}ल्कैर्वा+}\ A}}\skt{ }\edtext{\skt{पक्वाम्वा}}{\lemma{\skt{पक्वाम्वा}}\Afootnote{K Kb; \skt{अपि विचक्षणः। पक्वां व्}\ A; \skt{विधा\appbrk{}ने\appbrk{}नो\appbrk{}प\appbrk{}च\appbrk{}रेत्~। पक्वम्वा}\ H}}\skt{ }\edtext{\skt{पाक\-वि\-धा\-ने\-नो\-प\-च\-रेत्}}{\lemma{\skt{पाक\appbrk{}वि\appbrk{}धा\appbrk{}ने\appbrk{}नो\appbrk{}प\appbrk{}च\appbrk{}रेत्}}\Afootnote{K Kb H; \skt{रणविधानेन यथोक्तेन प्रसाधयेत्}\ A}}\skt{~॥}\pend

\pstart \ledsidenote{\passref{4.20.5-6}}\edtext{\edtext{\skt{यव\-प्र\-ख्या\-म\-म्बा\-लजीं}}{\lemma{\skt{यवप्रख्याम्}}\Afootnote{K Kb H; \skt{अन्धालजीं यवप्रख्यां}\ A}}}{\lemma{\skt{अम्बालजीं}}\Afootnote{K Kb H; \emph{om.}\ A}}\skt{ }\edtext{\edtext{\skt{पन\-सि\-का\-म्पा\-षा\-ण\-ग\-र्द\-भ\-ङ्क\-च्छ\-पि\-काञ्च}}{\lemma{\skt{पनसिकाम्}}\Afootnote{K Kb H; \skt{पनसीं कच्छपीं तथा}\ A}}}{\lemma{\skt{कच्छपिकाञ्च}}\Afootnote{K Kb H; \skt{चैव पूर्वं}\ A}}\skt{ }\edtext{\skt{स्वेदयित्वा}}{\lemma{\skt{स्वेदयित्वा}}\Afootnote{K Kb H; \skt{स्वेदेन योजयेत्}\ A}}\skt{ }\edtext{\edtext{\skt{सुर\-दा\-रु\-कु\-ष्ठ\-ह\-रि\-ता\-ल\-म\-नः\-शि\-ला\-क\-ल्कै\-रा\-ले\-प\-येत्}}{\lemma{\skt{सुर\appbrk{}दा\appbrk{}रु\appbrk{}कु\appbrk{}ष्ठ\appbrk{}ह\appbrk{}रि\appbrk{}ता\appbrk{}ल\appbrk{}म\appbrk{}नः\appbrk{}शि\appbrk{}ला\appbrk{}क\appbrk{}ल्कैर्}}\Afootnote{K Kb H; \skt{मनः\appbrk{}शि\appbrk{}ला\appbrk{}ता\appbrk{}ल\appbrk{}कु\appbrk{}ष्ठ\appbrk{}दा\appbrk{}रु\appbrk{}कल्कैः}\ A}}}{\lemma{\skt{आलेपयेत्}}\Afootnote{K Kb H; \skt{प्रलेपयेत्}\ A}}\skt{~। }\edtext{\skt{पक्वां वा}}{\lemma{\skt{पक्वां वा}}\Afootnote{K Kb H; \skt{परि}\ A}}\skt{ }\edtext{\skt{पाक\-वि\-धा\-ने\-नो\-प\-च\-रेत्}}{\lemma{\skt{पाक\appbrk{}वि\appbrk{}धा\appbrk{}ने\appbrk{}नो\appbrk{}प\appbrk{}च\appbrk{}रेत्}}\Afootnote{K Kb H; \skt{पाकगतान्भित्त्वा व्रण\appbrk{}व\appbrk{}त्स\appbrk{}मु\appbrk{}पा\appbrk{}च\appbrk{}रेत्}\ A}}\skt{~॥}\pend

\pstart \ledsidenote{\passref{4.20.7-8}}\edtext{\edtext{\skt{विवृ\-ता\-मि\-न्द्र\-विद्धां}}{\lemma{\skt{विवृताम्}}\Afootnote{K Kb H; \skt{विवृता}\ A}}}{\lemma{\skt{इन्द्रविद्धां}}\Afootnote{K Kb H; \skt{मिन्द्रवृद्धां च}\ A}}\skt{ }\edtext{\skt{गर्दभिकां}}{\lemma{\skt{गर्दभिकां}}\Afootnote{K Kb H; \skt{गर्दभीं}\ A}}\skt{ }\edtext{\skt{जाल\-ग\-र्द\-भ\-मि\-रि\-वे\-ल्लिकां}}{\lemma{\skt{इरिवेल्लिकां}}\Afootnote{K Kb H; \skt{इरिवेल्लीं}\ A}}\skt{ }\edtext{\skt{गन्धनामां}}{\lemma{\skt{गन्धनामां}}\Afootnote{K H; \skt{गन्धनाम्नीं}\ A; \skt{गन्धनामा}\ Kb}}\skt{ }\edtext{\edtext{\skt{कक्षाम्विस्फोटांश्च}}{\lemma{\skt{कक्षाम्}}\Afootnote{A K H; \skt{कक्}\ Kb}}}{\lemma{\skt{विस्फोटांश्च}}\Afootnote{K\textsuperscript{pc}; \skt{विस्फोटकांस्तथा॥ पित्तजस्य विसर्पस्य क्रियया साध}\ A; \skt{विस्फोटा}\ K\textsuperscript{ac}; \skt{ष्याम्विस्फोटांश्च}\ Kb; \skt{विस्फोटका́श्च पित्तविसर्}\ H}}\skt{ }\edtext{\skt{पित्त\-वि\-स\-र्प\-वि\-धा\-ने\-नो\-प\-च\-रेत्}}{\lemma{\skt{पित्त\appbrk{}वि\appbrk{}स\appbrk{}र्प\appbrk{}वि\appbrk{}धा\appbrk{}ने\appbrk{}नो\appbrk{}प\appbrk{}च\appbrk{}रेत्}}\Afootnote{K\textsuperscript{pc} Kb; \skt{येद्भिषक्। रोपयेत्सर्पिषा}\ A; \skt{ंश्च पिर्प\appbrk{}वि\appbrk{}धा\appbrk{}ने\appbrk{}नो\appbrk{}प\appbrk{}च\appbrk{}रेत्}\ K\textsuperscript{ac}; \skt{प्पवि\appbrk{}धा\appbrk{}ने\appbrk{}नो\appbrk{}प\appbrk{}च\appbrk{}रेत्}\ H}}\skt{~। }\edtext{\skt{पक्वांश्च}}{\lemma{\skt{पक्वांश्च}}\Afootnote{K Kb; \skt{पक्वान्सिद्धेन}\ A; \skt{पक्वा́श्च}\ H}}\skt{ }\edtext{\skt{मधु\-रौ\-ष\-ध\-सि\-द्धेन}}{\lemma{\skt{मधु\appbrk{}रौ\appbrk{}ष\appbrk{}ध\appbrk{}सि\appbrk{}द्धेन}}\Afootnote{K Kb H; \skt{मधुरौषधैः}\ A}}\skt{ }\edtext{\skt{सर्पिषा}}{\lemma{\skt{सर्पिषा}}\Afootnote{K Kb H; \emph{om.}\ A}}\skt{ }\edtext{\skt{रोपयेत्}}{\lemma{\skt{रोपयेत्}}\Afootnote{K Kb H; \emph{om.}\ A}}\skt{~॥}\pend

\pstart \ledsidenote{\passref{4.20.9-10}}\edtext{\edtext{\skt{चिप्यान्तूष्णाम्बुना}}{\lemma{\skt{चिप्यान्}}\Afootnote{K Kb H; \skt{चिप्य}\ A}}}{\lemma{\skt{तूष्णाम्बुना}}\Afootnote{K Kb H; \skt{(प्प)मुष्णाम्बुना}\ A}}\skt{ }\edtext{\edtext{\skt{परि\-षे\-च्यो\-त्कृ\-त्या\-प\-न\-ये\-च्च\-क्र\-तै\-ले\-ना\-भ्यज्य}}{\lemma{\skt{परि\appbrk{}षे\appbrk{}च्यो\appbrk{}त्कृ\appbrk{}त्या\appbrk{}प\appbrk{}न\appbrk{}येच्}}\Afootnote{K Kb H; \skt{सिक्तमुत्कृत्य स्रावयेद्}\ A}}}{\lemma{\skt{चक्र\appbrk{}तै\appbrk{}ले\appbrk{}ना\appbrk{}भ्यज्य}}\Afootnote{K Kb; \skt{भिषक्। चक्रतैलेन चाभ्यज्य}\ A; \skt{चक्रतैलेनाभ्य}\ H}}\skt{ }\edtext{\skt{सर्ज\-र\-स\-र\-स\-चू\-र्ण्णे\-ना\-व\-चूर्ण्ण्य}}{\lemma{\skt{सर्ज\appbrk{}र\appbrk{}स\appbrk{}र\appbrk{}स\appbrk{}चू\appbrk{}र्ण्णे\appbrk{}ना\appbrk{}व\appbrk{}चूर्ण्ण्य}}\Afootnote{K H; \skt{सर्जचूर्णेन चूर्णयेत्}\ A; \skt{सर्ज\appbrk{}स\appbrk{}चू\appbrk{}र्ण्णे\appbrk{}ना\appbrk{}चूर्ण्ण्य}\ Kb}}\skt{ }\edtext{\edtext{\edtext{\skt{बन्धे\-नो\-प\-च\-रे\-दे\-व\-म\-श\-क्य\-म\-ग्निना}}{\lemma{\skt{बन्धे\appbrk{}नो\appbrk{}प\appbrk{}च\appbrk{}रे\appbrk{}दे\appbrk{}वम्}}\Afootnote{K Kb H; \skt{बन्धे\appbrk{}नो\appbrk{}प\appbrk{}च\appbrk{}रे\appbrk{}च्चैन}\ A}}}{\lemma{\skt{अशक्यम्}}\Afootnote{K Kb H; \skt{मशक्यं}\ A}}}{\lemma{\skt{अग्निना}}\Afootnote{K Kb H; \skt{चाग्निना}\ A}}\skt{ दहेत्~। मधु\-रौ\-ष\-ध\-सि\-द्धेन }\edtext{\skt{च}}{\lemma{\skt{च}}\Afootnote{K Kb H; \skt{तत}\ A}}\skt{ }\edtext{\skt{तैले\-नो\-प\-च\-रेत्}}{\lemma{\skt{तैले\appbrk{}नो\appbrk{}प\appbrk{}च\appbrk{}रेत्}}\Afootnote{K Kb H; \skt{स्तैलेन रोपयेत्}\ A}}\skt{~॥}\pend

\pstart \ledsidenote{\passref{4.20.11}}\edtext{\skt{अनुशयीं}}{\lemma{\skt{अनुशयीं}}\Afootnote{K H; \skt{कुनखे विधिरप्येष कार्यो हि भिषजा भवेत्। उपा\appbrk{}च\appbrk{}रे\appbrk{}द\appbrk{}नु\appbrk{}शयीं}\ A; \skt{अनुशयी}\ Kb}}\skt{ श्लेष्म }\edtext{\skt{विद्र\-धि\-व\-दु\-प\-च\-रेत्}}{\lemma{\skt{विद्र\appbrk{}धि\appbrk{}व\appbrk{}दु\appbrk{}प\appbrk{}च\appbrk{}रेत्}}\Afootnote{K Kb H; \skt{विद्र\appbrk{}धि\appbrk{}व\appbrk{}द्भि\appbrk{}षक्}\ A}}\skt{~।}\pend

\pstart \ledsidenote{\passref{4.20.12-13}}\edtext{\skt{विदा\-रि\-का\-म\-भ्यज्य}}{\lemma{\skt{अभ्यज्य}}\Afootnote{K Kb H; \skt{समभ्यज्य}\ A}}\skt{ }\edtext{\skt{परिस्वेद्य}}{\lemma{\skt{परिस्वेद्य}}\Afootnote{K Kb H; \skt{स्विन्नां}\ A}}\skt{ }\edtext{\skt{विम्लापयेत्}}{\lemma{\skt{विम्लापयेत्}}\Afootnote{K Kb H; \skt{विम्लाप्य लेपयेत्}\ A}}\skt{~। }\edtext{\edtext{\skt{बिल्व\-मू\-ल\-ना\-ग\-वृ\-त्ति\-का\-व\-र्षा\-भू\-क\-ल्कै\-श्चा\-ले\-प\-येत्}}{\lemma{\skt{बिल्व\appbrk{}मू\appbrk{}ल\appbrk{}ना\appbrk{}ग\appbrk{}वृ\appbrk{}त्ति\appbrk{}का\appbrk{}व\appbrk{}र्षा\appbrk{}भू\appbrk{}क\appbrk{}ल्कैश्}}\Afootnote{K H\textsuperscript{ac}; \skt{नग\appbrk{}वृ\appbrk{}त्ति\appbrk{}क\appbrk{}व\appbrk{}र्षा\appbrk{}भू\appbrk{}बि\appbrk{}ल्व\appbrk{}मूलैः}\ A; \skt{बिल्व\appbrk{}मू\appbrk{}ल\appbrk{}ना\appbrk{}ग\appbrk{}द\appbrk{}न्ति\appbrk{}का\appbrk{}व\appbrk{}र्षा\appbrk{}भू\appbrk{}क\appbrk{}ल्कैश्}\ Kb H\textsuperscript{pc}}}}{\lemma{\skt{चालेपयेत्}}\Afootnote{K Kb H; \skt{सुपेषितैः}\ A}}\skt{~। }\edtext{\edtext{\skt{व्रण\-भा\-व\-ण\-मा\-पन्नं}}{\lemma{\skt{व्रणभावणम्}}\Afootnote{K H\textsuperscript{ac}; \skt{व्रण\appbrk{}भा\appbrk{}व\appbrk{}ग\appbrk{}तायां वा कृत्वा सं}\ A; \skt{व्रण\appbrk{}भा\appbrk{}व\appbrk{}ण\appbrk{}मा\appbrk{}पन्न वा}\ Kb; \skt{व्रण\appbrk{}भा\appbrk{}व\appbrk{}ण\appbrk{}मा\appbrk{}प\appbrk{}न्नम्वा}\ H\textsuperscript{pc}}}}{\lemma{\skt{आपन्नं}}\Afootnote{H\textsuperscript{ac}; \skt{शोधनक्रियाम्। रोपणार्थं}\ A; \skt{आपन्ना ं वा}\ K\textsuperscript{ac}; \skt{आपन्नं}\ K\textsuperscript{pc}; \emph{om.}\ Kb H\textsuperscript{pc}}}\skt{ }\edtext{\skt{वा मधु\-र\-क\-षा\-यौ\-ष\-ध\-सि\-द्धेन}}{\lemma{\skt{वा मधु\appbrk{}र\appbrk{}क\appbrk{}षा\appbrk{}यौ\appbrk{}ष\appbrk{}ध\appbrk{}सि\appbrk{}द्धेन}}\Afootnote{K\textsuperscript{pc} H\textsuperscript{ac}; \skt{हितं तैलं कषायमधुरैः शृतम्}\ A; \skt{मधु\appbrk{}र\appbrk{}क\appbrk{}षा\appbrk{}यौ\appbrk{}ष\appbrk{}ध\appbrk{}सि\appbrk{}द्धेन}\ K\textsuperscript{ac} Kb H\textsuperscript{pc}}}\skt{ }\edtext{\skt{तैलेन}}{\lemma{\skt{तैलेन}}\Afootnote{K Kb H; \emph{om.}\ A}}\skt{ }\edtext{\skt{रोपयेत्}}{\lemma{\skt{रोपयेत्}}\Afootnote{K Kb H; \emph{om.}\ A}}\skt{~॥}\pend

\pstart\versepar \ledsidenote{\passref{4.20.17-18}}\edtext{\edtext{\skt{शर्क\-रा\-र्बु\-द\-म्मे\-दो\-र्बु\-द\-व\-दु\-प\-च\-रेत्}}{\lemma{\skt{शर्करार्बुदम्}}\Afootnote{K H; \emph{om.}\ A; \skt{शर्करार्बुद}\ Kb}}}{\lemma{\skt{मेदो\appbrk{}र्बु\appbrk{}द\appbrk{}व\appbrk{}दु\appbrk{}प\appbrk{}च\appbrk{}रेत्}}\Afootnote{K Kb; \skt{मेदोऽ\appbrk{}र्बु\appbrk{}द\appbrk{}वि\appbrk{}धा\appbrk{}नेन साध\appbrk{}ये\appbrk{}च्छ\appbrk{}र्क\appbrk{}रा\appbrk{}र्बु\appbrk{}दम्}\ A; \skt{मेदोऽ\appbrk{}र्ब्बु\appbrk{}द\appbrk{}व\appbrk{}दु\appbrk{}प\appbrk{}च\appbrk{}रेत्}\ H}}\skt{~॥ }\pend
\pstart\versepar \skt{}\edtext{\skt{कच्छू\-वि\-च\-र्चि\-का\-र\-क\-शासु}}{\lemma{\skt{कच्छू\appbrk{}वि\appbrk{}च\appbrk{}र्चि\appbrk{}का\appbrk{}र\appbrk{}क\appbrk{}शासु}}\Afootnote{K H; \skt{कच्छूं विचर्चिकां पामां}\ A; \skt{कच्छ\appbrk{}वि\appbrk{}च\appbrk{}र्चि\appbrk{}का\appbrk{}र\appbrk{}क\appbrk{}शासु}\ Kb}}\skt{ }\edtext{\skt{कुष्ठ\-व\-त्प्र\-ती\-कारो}}{\lemma{\skt{प्रतीकारो}}\Afootnote{K Kb H; \skt{समुपाचरेत्॥ लेपश्}\ A}}\skt{ }\edtext{\skt{विशेषतस्तु}}{\lemma{\skt{विशेषतस्तु}}\Afootnote{K Kb H; \skt{च शस्यते सिक्थ}\ A}}\skt{ }\edtext{\edtext{\skt{शता\-ह्वा\-स\-र्ष\-प\-क\-ल्कै\-रा\-ले\-प\-येत्}}{\lemma{\skt{शता\appbrk{}ह्वा\appbrk{}स\appbrk{}र्ष\appbrk{}प\appbrk{}क\appbrk{}ल्कैर्}}\Afootnote{K Kb H; \skt{शता\appbrk{}ह्वा\appbrk{}गौ\appbrk{}र\appbrk{}सर्ष}\ A}}}{\lemma{\skt{आलेपयेत्}}\Afootnote{K Kb H; \skt{पैः}\ A}}\skt{~। }\edtext{\skt{सिद्धा\-र्थ\-क\-दा\-रु\-ह\-रि\-द्रा\-क\-ल्कैर्वा}}{\lemma{\skt{सिद्धा\appbrk{}र्थ\appbrk{}क\appbrk{}दा\appbrk{}रु\appbrk{}ह\appbrk{}रि\appbrk{}द्रा\appbrk{}क\appbrk{}ल्कैर्वा}}\Afootnote{K Kb H; \skt{वचादार्वी}\ A}}\skt{~। }\edtext{\skt{सरलातैलेन वा}}{\lemma{\skt{सरलातैलेन वा}}\Afootnote{K Kb H; \skt{सर्षपैर्वा}\ A}}\skt{~। }\edtext{\skt{कटु\-कौ\-ष\-ध\-सि\-द्धे\-ना\-भ्य\-ञ्ज्यात्}}{\lemma{\skt{कटु\appbrk{}कौ\appbrk{}ष\appbrk{}ध\appbrk{}सि\appbrk{}द्धे\appbrk{}ना\appbrk{}भ्य\appbrk{}ञ्ज्यात्}}\Afootnote{K Kb H; \skt{तैलं वा}\ A}}\skt{~। }\edtext{\skt{नक्त\-मा\-ल\-तै\-लेन}}{\lemma{\skt{नक्त\appbrk{}मा\appbrk{}ल\appbrk{}तै\appbrk{}लेन}}\Afootnote{K Kb H; \skt{नक्तमालजम्}\ A}}\skt{ }\edtext{\skt{वा}}{\lemma{\skt{वा}}\Afootnote{K Kb H; \emph{om.}\ A}}\skt{~।}\pend

\pstart\versepar \ledsidenote{\passref{4.20.19}}\edtext{\skt{पक्वांश्च}}{\lemma{\skt{पक्वांश्च}}\Afootnote{K Kb H; \skt{सा}\ A}}\skt{ }\edtext{\skt{पक्व\-वि\-धा\-ने\-नो\-प\-च\-रेत्}}{\lemma{\skt{पक्व\appbrk{}वि\appbrk{}धा\appbrk{}ने\appbrk{}नो\appbrk{}प\appbrk{}च\appbrk{}रेत्}}\Afootnote{K Kb H; \skt{रतै\appbrk{}ल\appbrk{}म\appbrk{}था\appbrk{}भ्यङ्गे कुर्वीत कटुकैः शृतम्}\ A}}\skt{~॥ }\pend
\pstart\versepar \skt{}\edtext{\skt{पाददार्यान्तु}}{\lemma{\skt{पाददार्यान्तु}}\Afootnote{K Kb H; \skt{पाददार्यां}\ A}}\skt{ }\edtext{\skt{सिराम्विध्येत्}}{\lemma{\skt{विध्येत्}}\Afootnote{K Kb H; \skt{विद्ध्वा स्वेदाभ्यङ्गौ प्रयोजयेत्}\ A}}\skt{~।}\pend

\pstart \ledsidenote{\passref{4.20.20}}\edtext{\skt{पादौ}}{\lemma{\skt{पादौ}}\Afootnote{K Kb H; \emph{om.}\ A}}\skt{ }\edtext{\skt{स्नेह\-स्वे\-दो\-प\-पन्ना}}{\lemma{\skt{स्नेह\appbrk{}स्वे\appbrk{}दो\appbrk{}प\appbrk{}पन्ना}}\Afootnote{K Kb H; \emph{om.}\ A}}\skt{ }\edtext{\edtext{\edtext{\skt{मधू\-च्छि\-ष्ट\-घृ\-त\-व\-सा\-म\-ज्जा\-भि\-र्य\-व\-क्षा\-र\-मि\-श्रै\-रा\-ले\-प\-येत्}}{\lemma{\skt{मधू\appbrk{}च्छि\appbrk{}ष्ट\appbrk{}घृ\appbrk{}त\appbrk{}व\appbrk{}सा\appbrk{}म\appbrk{}ज्जा\appbrk{}भिर्}}\Afootnote{K H; \skt{मधू\appbrk{}च्छि\appbrk{}ष्ट\appbrk{}व\appbrk{}सा\appbrk{}म\appbrk{}ज्ज\appbrk{}स\appbrk{}र्ज\appbrk{}चू\appbrk{}र्ण\appbrk{}घृतैः कृतः}\ A; \skt{मधूच्छिष्टा घृत\appbrk{}व\appbrk{}सा\appbrk{}म\appbrk{}ज्जाभि}\ Kb}}}{\lemma{\skt{यव\appbrk{}क्षा\appbrk{}र\appbrk{}मि\appbrk{}श्रैर्}}\Afootnote{K Kb H; \skt{यवा\appbrk{}ह्व\appbrk{}गै\appbrk{}रि\appbrk{}को\appbrk{}न्मिश्रैः}\ A}}}{\lemma{\skt{आलेपयेत्}}\Afootnote{K Kb H; \skt{पादलेपः प्रशस्यते}\ A}}\skt{~॥}\pend

\pstart \ledsidenote{\passref{4.20.21-22}}\edtext{\skt{अलसे}}{\lemma{\skt{अलसे}}\Afootnote{K Kb H; \skt{पादौ सिक्त्वा++\appbrk{}आ\appbrk{}र\appbrk{}ना\appbrk{}लेन लेपनं ह्यलसे}\ A}}\skt{ }\edtext{\skt{त्वाम्ल\-का\-ञ्जि\-क\-सिक्तौ}}{\lemma{\skt{त्वाम्ल\appbrk{}का\appbrk{}ञ्जि\appbrk{}क\appbrk{}सिक्तौ}}\Afootnote{K Kb H; \emph{om.}\ A}}\skt{ }\edtext{\skt{चरणौ}}{\lemma{\skt{चरणौ}}\Afootnote{K Kb H; \emph{om.}\ A}}\skt{ }\edtext{\edtext{\skt{पटो\-ल\-पि\-चु\-म\-र्द\-ति\-ल\-का\-सी\-स\-रो\-च\-ना\-क\-ल्कै\-रा\-ले\-प\-येत्}}{\lemma{\skt{पटो\appbrk{}ल\appbrk{}पि\appbrk{}चु\appbrk{}म\appbrk{}र्द\appbrk{}ति\appbrk{}ल\appbrk{}का\appbrk{}सी\appbrk{}स\appbrk{}रो\appbrk{}च\appbrk{}ना\appbrk{}क\appbrk{}ल्कैर्}}\Afootnote{K Kb H; \skt{हितम्। कल्की\appbrk{}कृ\appbrk{}तै\appbrk{}र्नि\appbrk{}म्ब\appbrk{}ति\appbrk{}ल\appbrk{}का\appbrk{}सीसा}\ A}}}{\lemma{\skt{आलेपयेत्}}\Afootnote{K Kb H; \skt{लैः ससैन्धवैः॥ लाक्षारसोऽभया वा+अपि कार्यं स्याद्र\appbrk{}क्त\appbrk{}मो\appbrk{}क्ष\appbrk{}णम्। सिद्धं रसे कण्टकार्यास्तैलं वा सार्षपं हितम्}\ A}}\skt{~।}\pend

\pstart \ledsidenote{\passref{4.20.23}}\edtext{\skt{निदि\-ग्धि\-का\-र\-स\-सि\-द्धेन च}}{\lemma{\skt{निदि\appbrk{}ग्धि\appbrk{}का\appbrk{}र\appbrk{}स\appbrk{}सि\appbrk{}द्धेन च}}\Afootnote{K Kb H; \emph{om.}\ A}}\skt{ }\edtext{\skt{सर्ष\-प\-तै\-ले\-ना\-भ्यज्य}}{\lemma{\skt{सर्ष\appbrk{}प\appbrk{}तै\appbrk{}ले\appbrk{}ना\appbrk{}भ्यज्य}}\Afootnote{K Kb H; \emph{om.}\ A}}\skt{ }\edtext{\skt{कासी\-स\-रो\-च\-ना\-म\-नः\-शि\-ला\-चूर्ण्णैः}}{\lemma{\skt{कासी\appbrk{}स\appbrk{}रो\appbrk{}च\appbrk{}ना\appbrk{}म\appbrk{}नः\appbrk{}शि\appbrk{}ला\appbrk{}चूर्ण्णैः}}\Afootnote{K Kb H; \skt{कासी\appbrk{}स\appbrk{}रो\appbrk{}च\appbrk{}न\appbrk{}शि\appbrk{}ला\appbrk{}चू\appbrk{}र्णैर्वा}\ A}}\skt{ }\edtext{\skt{प्रतिसारयेत्}}{\lemma{\skt{प्रतिसारयेत्}}\Afootnote{K Kb H; \skt{प्रतिसार}\ A}}\skt{~॥ }\edtext{\edtext{\skt{कदरमुद्धृत्य}}{\lemma{\skt{कदरम्}}\Afootnote{K Kb H; \skt{णम्}\ A}}}{\lemma{\skt{उद्धृत्य}}\Afootnote{K Kb H; \skt{उत्कृत्य}\ A}}\skt{ }\edtext{\edtext{\skt{दहेत्स्नेहेन}}{\lemma{\skt{दहेत्}}\Afootnote{K Kb H; \skt{दग्ध्वा}\ A}}}{\lemma{\skt{स्नेहेन}}\Afootnote{K Kb H; \skt{स्नेहेन जये\appbrk{}त्क\appbrk{}द\appbrk{}र\appbrk{}सं\appbrk{}ज्ञ\appbrk{}कम्}\ A}}\skt{~॥}\pend

\pstart \ledsidenote{\passref{4.20.24}}\edtext{\skt{इन्द्रलुप्ते तु}}{\lemma{\skt{इन्द्रलुप्ते तु}}\Afootnote{K Kb H; \skt{इन्द्रलुप्ते}\ A}}\skt{ }\edtext{\skt{शिरसि}}{\lemma{\skt{शिरसि}}\Afootnote{K Kb H; \emph{om.}\ A}}\skt{ }\edtext{\skt{सिराम्विधा तु }\edtext{\skt{न्दम\-नः\-शि\-ला\-म\-रि\-च\-क\-ल्कै\-रा\-ले\-प\-येत्}}{\lemma{\skt{आलेपयेत्}}\Afootnote{K Kb H; \skt{नस्य मोक्षयेत्। कल्कैः सम\appbrk{}रि\appbrk{}चै\appbrk{}र्दि\appbrk{}ह्या\appbrk{}च्छि\appbrk{}ला\appbrk{}का\appbrk{}सी\appbrk{}स\appbrk{}तु\appbrk{}त्थकैः}\ A}}}{\lemma{\skt{विधा तु न्दम\appbrk{}नः\appbrk{}शि\appbrk{}ला\appbrk{}म\appbrk{}रि\appbrk{}च\appbrk{}क\appbrk{}ल्कैर्}}\Afootnote{K; \skt{मूर्ध्नि स्निग्धस्विन्}\ A; \skt{विध्वा मनः\appbrk{}शि\appbrk{}ला\appbrk{}म\appbrk{}रि\appbrk{}च\appbrk{}क\appbrk{}ल्कैर्}\ Kb H\textsuperscript{ac}; \skt{विद्ध्वा तुन्द\appbrk{}का\appbrk{}सी\appbrk{}स\appbrk{}म\appbrk{}रि\appbrk{}च\appbrk{}क\appbrk{}ल्कैर्}\ H\textsuperscript{pc}}}\skt{~।}\pend

\pstart \ledsidenote{\passref{4.20.25-26}}\edtext{\skt{भद्र\-दा\-रु\-क\-ल्कैर्वा}}{\lemma{\skt{भद्र\appbrk{}दा\appbrk{}रु\appbrk{}क\appbrk{}ल्कैर्वा}}\Afootnote{\skt{कट\appbrk{}न्न\appbrk{}ट\appbrk{}दा\appbrk{}रु\appbrk{}क\appbrk{}ल्कैर्}\ A; \skt{कुट\appbrk{}न्न\appbrk{}ट\appbrk{}भ\appbrk{}द्र\appbrk{}दा\appbrk{}रु\appbrk{}क\appbrk{}ल्कैर्वा}\ K H; \skt{कुटं\appbrk{}न\appbrk{}टा\appbrk{}भ\appbrk{}द्र\appbrk{}दा\appbrk{}रु\appbrk{}क\appbrk{}ल्कैर्वा}\ Kb}}\skt{ }\edtext{\skt{जाती\-चि\-त्र\-क\-क\-र\-वी\-र\-न\-क्त\-मा\-ल\-क\-ल्क\-सि\-द्धेन}}{\lemma{\skt{जाती\appbrk{}चि\appbrk{}त्र\appbrk{}क\appbrk{}क\appbrk{}र\appbrk{}वी\appbrk{}र\appbrk{}न\appbrk{}क्त\appbrk{}मा\appbrk{}ल\appbrk{}क\appbrk{}ल्क\appbrk{}सि\appbrk{}द्धेन}}\Afootnote{K Kb H; \emph{om.}\ A}}\skt{ }\edtext{\skt{तैलेनाभ्यञ्ज्यात्}}{\lemma{\skt{तैलेनाभ्यञ्ज्यात्}}\Afootnote{K Kb; \emph{om.}\ A; \skt{तैलेनाभ्यज्यात्}\ H}}\skt{~। }\edtext{\edtext{\skt{रसा\-य\-न\-वि\-धा\-न\-ञ्चा\-से\-वेन}}{\lemma{\skt{रसा\appbrk{}य\appbrk{}न\appbrk{}वि\appbrk{}धा\appbrk{}नञ्}}\Afootnote{K Kb H; \emph{om.}\ A}}}{\lemma{\skt{चासेवेन}}\Afootnote{\skt{लेप}\ A; \skt{चासेन्धवेन}\ K; \skt{चासेत्ववेन}\ Kb; \skt{सेवेत}\ H}}\skt{~। }\edtext{\skt{अवगाढपदं}}{\lemma{\skt{अवगाढपदं}}\Afootnote{K Kb H; \skt{नं वा प्रशस्यते}\ A}}\skt{ }\edtext{\skt{प्रच्छयित्वा}}{\lemma{\skt{प्रच्छयित्वा}}\Afootnote{K Kb H; \skt{प्रच्छ\appbrk{}यि\appbrk{}त्वा+\appbrk{}अ\appbrk{}व\appbrk{}गाढं वा}\ A}}\skt{ }\edtext{\edtext{\skt{गुञ्जा\-फ\-ल\-क\-ल्कै\-र्भूयो}}{\lemma{\skt{गुञ्जा\appbrk{}फ\appbrk{}ल\appbrk{}क\appbrk{}ल्कैर्}}\Afootnote{K H; \skt{गुञ्जाकल्कैर्}\ A; \skt{गुंजाफलकल्कै}\ Kb}}}{\lemma{\skt{भूयो}}\Afootnote{K Kb; \emph{om.}\ A; \skt{ब्भूयो}\ H}}\skt{ }\edtext{\skt{भूयः}}{\lemma{\skt{भूयः}}\Afootnote{K Kb H; \emph{om.}\ A}}\skt{ }\edtext{\skt{प्रदिह्यात्}}{\lemma{\skt{प्रदिह्यात्}}\Afootnote{K Kb H; \skt{मुहुर्मुहुः॥ लेप\appbrk{}ये\appbrk{}दु\appbrk{}प\appbrk{}शा\appbrk{}न्त्यर्थं कुर्याद्वा+अपि रसायनम्। माल\appbrk{}ती\appbrk{}क\appbrk{}र\appbrk{}वी\appbrk{}रा\appbrk{}ग्नि\appbrk{}न\appbrk{}क्त\appbrk{}मा\appbrk{}ल\appbrk{}वि\appbrk{}पा\appbrk{}चि\appbrk{}तम्}\ A}}\skt{~॥}\pend

\pstart \ledsidenote{\passref{4.20.27-29}}\edtext{\edtext{\skt{शोणि\-त\-म\-व\-सिच्य}}{\lemma{\skt{शोणितम्}}\Afootnote{\skt{तैलमभ्यञ्जने शस्त\appbrk{}मि\appbrk{}न्द्र\appbrk{}लु\appbrk{}प्ता\appbrk{}पहं परम्। अरूंषिकां हृते रक्ते सेच\appbrk{}ये\appbrk{}न्नि\appbrk{}म्ब\appbrk{}वा\appbrk{}रिणा}\ A; \skt{अरूषिका तु जलौकोभिः शोणितम्}\ K H; \skt{अरूंषिका तु जलूकाभिः शोणितम्}\ Kb}}}{\lemma{\skt{अवसिच्य}}\Afootnote{K Kb H; \skt{दिह्या}\ A}}\skt{ }\edtext{\skt{निम्बोदकेन}}{\lemma{\skt{निम्बोदकेन}}\Afootnote{K Kb H; \skt{त्सैन्धवयुक्तेन}\ A}}\skt{ }\edtext{\skt{परि\-सि\-च्या\-श्व\-वि\-ड्र\-सेन}}{\lemma{\skt{परि\appbrk{}सि\appbrk{}च्या\appbrk{}श्व\appbrk{}वि\appbrk{}ड्र\appbrk{}सेन}}\Afootnote{K Kb H; \skt{वाजि\appbrk{}वि\appbrk{}ष्ठा\appbrk{}र\appbrk{}सेन}\ A}}\skt{ }\edtext{\skt{लव\-ण\-प्र\-गा\-ढे\-ना\-ले\-प\-येत्}}{\lemma{\skt{लव\appbrk{}ण\appbrk{}प्र\appbrk{}गा\appbrk{}ढे\appbrk{}ना\appbrk{}ले\appbrk{}प\appbrk{}येत्}}\Afootnote{K H; \emph{om.}\ A; \skt{लव\appbrk{}ण\appbrk{}प्र\appbrk{}गा\appbrk{}ढे\appbrk{}ण\appbrk{}ना\appbrk{}ले\appbrk{}प\appbrk{}येत्}\ Kb}}\skt{~। }\edtext{\skt{पटो\-ल\-पि\-चु\-म\-र्द\-ह\-रि\-द्रा\-क\-ल्कैर्वा}}{\lemma{\skt{पटो\appbrk{}ल\appbrk{}पि\appbrk{}चु\appbrk{}म\appbrk{}र्द\appbrk{}ह\appbrk{}रि\appbrk{}द्रा\appbrk{}क\appbrk{}ल्कैर्वा}}\Afootnote{K Kb H; \skt{तु। हरि\appbrk{}ता\appbrk{}ल\appbrk{}नि\appbrk{}शा\appbrk{}नि\appbrk{}म्ब\appbrk{}क\appbrk{}ल्कैर्वा}\ A}}\skt{~। }\edtext{\skt{मधु\-को\-त्प\-ल\-ति\-लै\-र\-ण्ड\-मा\-र्क\-व\-प\-त्रैर्वा}}{\lemma{\skt{मधु\appbrk{}को\appbrk{}त्प\appbrk{}ल\appbrk{}ति\appbrk{}लै\appbrk{}र\appbrk{}ण्ड\appbrk{}मा\appbrk{}र्क\appbrk{}व\appbrk{}प\appbrk{}त्रैर्वा}}\Afootnote{K Kb H; \skt{सपटोलजैः॥ यष्टी\appbrk{}नी\appbrk{}लो\appbrk{}त्प\appbrk{}लै\appbrk{}र\appbrk{}ण्ड\appbrk{}मा\appbrk{}र्क\appbrk{}वैर्वा}\ A}}\skt{~। }\edtext{\skt{इन्द्रलुप्तोक्तेन}}{\lemma{\skt{इन्द्रलुप्तोक्तेन}}\Afootnote{K Kb H; \skt{प्रलेपयेत्। इन्द्रलुप्तापहं}\ A}}\skt{ }\edtext{\skt{तैलेनाभ्यञ्ज्यात्}}{\lemma{\skt{तैलेनाभ्यञ्ज्यात्}}\Afootnote{K\textsuperscript{pc} Kb; \skt{तैलमभ्यङ्गे च प्रशस्यते। सिरां}\ A; \skt{तैलेभ्यंज्यात्}\ K\textsuperscript{ac}; \skt{तैलेनाभ्यज्यात्}\ H}}\skt{~॥ दारुणके }\edtext{\skt{तु}}{\lemma{\skt{तु}}\Afootnote{K Kb H; \skt{विद्ध्वा}\ A}}\skt{ स्निग्धस्विन्नस्य }\edtext{\skt{ललाटे}}{\lemma{\skt{ललाटे}}\Afootnote{K Kb H; \skt{मूर्ध}\ A}}\skt{ }\edtext{\edtext{\skt{सिराम्व्यधेत्}}{\lemma{\skt{सिराम्}}\Afootnote{K Kb H; \skt{नि}\ A}}}{\lemma{\skt{व्यधेत्}}\Afootnote{K Kb H; \emph{om.}\ A}}\skt{~।}\pend

\pstart\versepar \ledsidenote{\passref{4.20.30}}\edtext{\skt{अवपीडशिरो}}{\lemma{\skt{अवपीडशिरो}}\Afootnote{K Kb H; \skt{अवपीडं शिरो}\ A}}\skt{ }\edtext{\edtext{\skt{बस्त्य\-भ्य\-गा\-श्चा\-से\-वेत}}{\lemma{\skt{बस्त्यभ्यगाश्}}\Afootnote{\skt{बस्तिमभ्यङ्गं}\ A; \skt{बस्त्यभ्यंगांश्}\ K H; \skt{बस्त्यभ्यंगाश्}\ Kb}}}{\lemma{\skt{चासेवेत}}\Afootnote{K H; \skt{च प्रयोजयेत्। क्षालने}\ A; \skt{चासेवेत्}\ Kb}}\skt{~। }\pend
\pstart\versepar \skt{}\edtext{\edtext{\skt{कोद्र\-व\-तृ\-ण\-क्ष\-रौ\-द\-क\-प्र\-क्षा\-ल\-नञ्च}}{\lemma{\skt{कोद्र\appbrk{}व\appbrk{}तृ\appbrk{}ण\appbrk{}क्ष\appbrk{}रौ\appbrk{}द\appbrk{}क\appbrk{}प्र\appbrk{}क्षा\appbrk{}ल\appbrk{}नञ्}}\Afootnote{K; \skt{कोद्र\appbrk{}व\appbrk{}तृ\appbrk{}ण\appbrk{}क्षा\appbrk{}र\appbrk{}तोयं प्रशस्यते}\ A; \skt{कोद्र\appbrk{}व\appbrk{}तृ\appbrk{}ण\appbrk{}क्षा\appbrk{}रौ\appbrk{}प्र\appbrk{}क\appbrk{}प्र\appbrk{}क्षा\appbrk{}ल\appbrk{}नञ्}\ Kb; \skt{कोद्र\appbrk{}व\appbrk{}तृ\appbrk{}ण\appbrk{}क्षा\appbrk{}रौ\appbrk{}द\appbrk{}क\appbrk{}प्र\appbrk{}क्षा\appbrk{}ल\appbrk{}नञ्}\ H}}}{\lemma{\skt{च}}\Afootnote{K H; \emph{om.}\ A; \skt{च~॥ {\symfont ❈}}\ Kb}}\skt{~॥}\pend

\pstart\versepar \ledsidenote{\passref{4.20.31}}\edtext{\edtext{\skt{पलि\-त\-प्र\-का\-र\-म\-न्यत्र}}{\lemma{\skt{पलि\appbrk{}त\appbrk{}प्र\appbrk{}का\appbrk{}रम्}}\Afootnote{K H; \skt{उपरिष्टात्प्}\ A}}}{\lemma{\skt{अन्यत्र}}\Afootnote{K H; \skt{र}\ A}}\skt{ }\edtext{\skt{वक्ष्यामः}}{\lemma{\skt{वक्ष्यामः}}\Afootnote{K H; \skt{वक्ष्यामि विधिं पलि\appbrk{}त\appbrk{}ना\appbrk{}श\appbrk{}नम्}\ A}}\skt{~। }\pend
\pstart\versepar \skt{}\edtext{\skt{मसूरिकां}}{\lemma{\skt{मसूरिकां}}\Afootnote{K H; \skt{मसूरि}\ A}}\skt{ }\edtext{\edtext{\skt{कफ\-पि\-त्त\-वि\-स\-र्प\-कु\-ष्ठ\-ह\-र\-प्र\-ले\-प\-ना\-दि\-क्रि\-या\-मु\-प\-दि\-शन्ति}}{\lemma{\skt{कफ\appbrk{}पि\appbrk{}त्त\appbrk{}वि\appbrk{}स\appbrk{}र्प\appbrk{}कु\appbrk{}ष्ठ\appbrk{}ह\appbrk{}र\appbrk{}प्र\appbrk{}ले\appbrk{}प\appbrk{}ना\appbrk{}दि\appbrk{}क्रि\appbrk{}याम्}}\Afootnote{K H; \skt{कायां कुष्ठ\appbrk{}घ्न\appbrk{}ले\appbrk{}प\appbrk{}ना\appbrk{}दि\appbrk{}क्रिया}\ A}}}{\lemma{\skt{उपदिशन्ति}}\Afootnote{K H; \skt{हिता}\ A}}\skt{~॥}\pend

\pstart \ledsidenote{\passref{4.20.32cd-33ab}}\edtext{\skt{जतु\-म\-णि\-म\-श\-क\-च\-र्म\-की\-ल\-ति\-ल\-का\-ल\-कांस्तु}}{\lemma{\skt{जतु\appbrk{}म\appbrk{}णि\appbrk{}म\appbrk{}श\appbrk{}क\appbrk{}च\appbrk{}र्म\appbrk{}की\appbrk{}ल\appbrk{}ति\appbrk{}ल\appbrk{}का\appbrk{}ल\appbrk{}कांस्तु}}\Afootnote{K H; \skt{जतुमणिं समुत्कृत्य मषकं तिलकालकम्}\ A}}\skt{ }\edtext{\skt{शस्त्रेणोत्कृत्य}}{\lemma{\skt{शस्त्रेणोत्कृत्य}}\Afootnote{K H; \emph{om.}\ A}}\skt{ क्षारेण }\edtext{\skt{प्रतिसारयेत्~। अग्निना वा}}{\lemma{\skt{प्रतिसारयेत्~। अग्निना वा}}\Afootnote{K H; \skt{प्र}\ A}}\skt{ }\edtext{\skt{दहेत्}}{\lemma{\skt{दहेत्}}\Afootnote{K H; \skt{दहेद्युक्त्या वह्निना वा शनैः शनैः}\ A}}\skt{~॥}\pend

\pstart \ledsidenote{\passref{4.20.33cd}}\edtext{\skt{नच्छ\-नी\-लि\-का\-व्य\-ङ्गेषु}}{\lemma{\skt{नच्छ\appbrk{}नी\appbrk{}लि\appbrk{}का\appbrk{}व्य\appbrk{}ङ्गेषु}}\Afootnote{K H; \skt{न्यच्छे व्यङ्गे सिरामोक्षो नीलिकायां च श}\ A}}\skt{ }\edtext{\skt{यथाभ्याशं}}{\lemma{\skt{यथाभ्याशं}}\Afootnote{K H; \skt{स्यते}\ A}}\skt{ }\edtext{\skt{ललाटे}}{\lemma{\skt{ललाटे}}\Afootnote{K H; \emph{om.}\ A}}\skt{ }\edtext{\edtext{\skt{सिरा\-म्वि\-स्रा\-व\-येत्}}{\lemma{\skt{सिराम्}}\Afootnote{K H; \emph{om.}\ A}}}{\lemma{\skt{विस्रावयेत्}}\Afootnote{K H; \emph{om.}\ A}}\skt{~।}\pend

\pstart \ledsidenote{\passref{4.20.34}}\edtext{\skt{क्षीर\-वृ\-क्ष\-क\-ल्कैश्च}}{\lemma{\skt{क्षीर\appbrk{}वृ\appbrk{}क्ष\appbrk{}क\appbrk{}ल्कैश्च}}\Afootnote{K H; \skt{यथान्यायं यथाभ्यासं लाला\appbrk{}ट्या\appbrk{}दि\appbrk{}सि\appbrk{}रा\appbrk{}व्यधः। घृष्ट्वा दिह्यात्त्वचं पिष्ट्वा क्षीरिणां}\ A}}\skt{ }\edtext{\edtext{\skt{क्षीर\-पि\-ष्टै\-र\-व\-घृ\-ष्या\-ले\-प\-येत्}}{\lemma{\skt{क्षीरपिष्टैर्}}\Afootnote{K H; \skt{क्षीर}\ A}}}{\lemma{\skt{अव\appbrk{}घृ\appbrk{}ष्या\appbrk{}ले\appbrk{}प\appbrk{}येत्}}\Afootnote{K H; \skt{संयुताम्}\ A}}\skt{~।}\pend

\pstart \ledsidenote{\passref{4.20.35}}\edtext{\skt{बला\-ति\-ब\-ला\-म\-धु\-क\-ह\-रि\-द्रा\-क\-ल्कैर्वा}}{\lemma{\skt{बला\appbrk{}ति\appbrk{}ब\appbrk{}ला\appbrk{}म\appbrk{}धु\appbrk{}क\appbrk{}ह\appbrk{}रि\appbrk{}द्रा\appbrk{}क\appbrk{}ल्कैर्वा}}\Afootnote{K H; \skt{बल\appbrk{}ति\appbrk{}ब\appbrk{}ल\appbrk{}यष्ट्या}\ A}}\skt{~। }\edtext{\edtext{\skt{सुव\-र्ण्ण\-गै\-रि\-का\-गु\-रु\-प\-य\-स्या\-क\-ल्कैर्वा}}{\lemma{\skt{सुव\appbrk{}र्ण्ण\appbrk{}गै\appbrk{}रि\appbrk{}का\appbrk{}गु\appbrk{}रु\appbrk{}प\appbrk{}य\appbrk{}स्या\appbrk{}क\appbrk{}ल्कैर्}}\Afootnote{K H; \skt{ह्वरजनीर्वा प्रलेपनम्। पय\appbrk{}स्या\appbrk{}गु\appbrk{}रु\appbrk{}का\appbrk{}ली\appbrk{}य\appbrk{}ले\appbrk{}पनं वा सगैरिकम्}\ A}}}{\lemma{\skt{वा}}\Afootnote{K H; \emph{om.}\ A}}\skt{~।}\pend

\pstart \ledsidenote{\passref{4.20.36}}\edtext{\skt{तरु\-ण\-क\-पि\-त्थ\-रा\-जा\-द\-न\-क\-ल्कैर्वा}}{\lemma{\skt{तरु\appbrk{}ण\appbrk{}क\appbrk{}पि\appbrk{}त्थ\appbrk{}रा\appbrk{}जा\appbrk{}द\appbrk{}न\appbrk{}क\appbrk{}ल्कैर्वा}}\Afootnote{K H; \skt{क्षौद्राज्ययुक्तया लिम्पेद्दंष्ट्रया शूकरस्य च। कपि\appbrk{}त्थ\appbrk{}रा\appbrk{}जा\appbrk{}द\appbrk{}नयोः कल्}\ A}}\skt{~। }\edtext{\skt{सूकरदंष्ट्रया वा}}{\lemma{\skt{सूकरदंष्ट्रया वा}}\Afootnote{K H; \skt{कं वा}\ A}}\skt{ }\edtext{\skt{मधु\-घृ\-त\-यु\-क्तया}}{\lemma{\skt{मधु\appbrk{}घृ\appbrk{}त\appbrk{}यु\appbrk{}क्तया}}\Afootnote{K H; \skt{हितमुच्यते}\ A}}\skt{~।}\pend

\pstart \ledsidenote{\passref{4.20.37}}\edtext{\edtext{\skt{यूनामपि}}{\lemma{\skt{यूनाम्}}\Afootnote{K H; \emph{om.}\ A}}}{\lemma{\skt{अपि}}\Afootnote{K H; \emph{om.}\ A}}\skt{ }\edtext{\skt{युवानपि }\edtext{\skt{अकास्वेष}}{\lemma{\skt{एष}}\Afootnote{K H; \skt{वेष}\ A}}}{\lemma{\skt{युवानपि अकास्व्}}\Afootnote{\skt{यौवने पिडकास्}\ A; \skt{युवा\appbrk{}न\appbrk{}पि\appbrk{}ट\appbrk{}कास्व्}\ K H}}\skt{ }\edtext{\skt{एव}}{\lemma{\skt{एव}}\Afootnote{K H; \emph{om.}\ A}}\skt{ }\edtext{\skt{विधिः}}{\lemma{\skt{विधिः}}\Afootnote{K H; \emph{om.}\ A}}\skt{~। }\edtext{\skt{विशेषतस्तु}}{\lemma{\skt{विशेषतस्तु}}\Afootnote{K H; \skt{विशेषा}\ A}}\skt{ }\edtext{\edtext{\edtext{\skt{प्रच्छ\-र्द\-न\-मि\-ष्ट\-मा\-ले\-पश्च}}{\lemma{\skt{प्रच्छर्दनम्}}\Afootnote{K; \skt{च्छर्दनं}\ A; \skt{प्रच्छनम्}\ H}}}{\lemma{\skt{इष्टम्}}\Afootnote{K H; \skt{हि}\ A}}}{\lemma{\skt{आलेपश्च}}\Afootnote{K H; \skt{तम्। लेप}\ A}}\skt{ }\edtext{\skt{सिद्धा\-र्थ\-क\-व\-चा\-लो\-ध्र\-सै\-न्ध\-व\-क\-ल्कैर्वा}}{\lemma{\skt{सिद्धा\appbrk{}र्थ\appbrk{}क\appbrk{}व\appbrk{}चा\appbrk{}लो\appbrk{}ध्र\appbrk{}सै\appbrk{}न्ध\appbrk{}व\appbrk{}क\appbrk{}ल्कैर्वा}}\Afootnote{K H; \skt{नं च वचा\appbrk{}रो\appbrk{}ध्र\appbrk{}सै\appbrk{}न्धवैः सर्षपान्वितैः}\ A}}\skt{~॥}\pend

\pstart\versepar \ledsidenote{\passref{4.20.38.add}}\skt{शाल्म\-ली\-क\-ण्ट\-का\-म्पिष्ट्वा गव्येन पयसा सह~। }\pend
\pstart\versepar \skt{विव\-र्ण्ण\-म्मु\-ख\-मा\-लि\-म्पे\-त्स\-व\-र्ण्ण\-ञ्जा\-यते ततः~॥}\pend

\pstart \ledsidenote{\passref{4.20.38cd-39ab}}\skt{पद्मिनीकण्टके }\edtext{\skt{तु}}{\lemma{\skt{तु}}\Afootnote{K H; \skt{रोगे छर्दयेन्}\ A}}\skt{ }\edtext{\skt{निम्बो\-द\-क\-वा\-न्तस्य}}{\lemma{\skt{निम्बो\appbrk{}द\appbrk{}क\appbrk{}वा\appbrk{}न्तस्य}}\Afootnote{K; \skt{निम्बवा}\ A; \skt{निम्बोदक}\ H}}\skt{ }\edtext{\skt{निम्बोदकसर्पिः}}{\lemma{\skt{निम्बोदकसर्पिः}}\Afootnote{K; \skt{रिणा॥ तेनैव सिद्धं}\ A; \skt{सर्प्पिः}\ H}}\skt{ }\edtext{\edtext{\edtext{\skt{क्षौद्र\-म्पा\-तु\-म्प्र\-य\-च्छेत्}}{\lemma{\skt{क्षौद्रम्}}\Afootnote{K; \skt{सक्षौद्रं}\ A H}}}{\lemma{\skt{पातुम्}}\Afootnote{K H; \skt{सर्पिःपानं}\ A}}}{\lemma{\skt{प्रयच्छेत्}}\Afootnote{K H; \skt{प्रदापयेत्}\ A}}\skt{~।}\pend

\pstart \ledsidenote{\passref{4.20.39cd}}\edtext{\skt{निम्बा\-र\-ग्व\-ध\-क\-षायं}}{\lemma{\skt{निम्बा\appbrk{}र\appbrk{}ग्व\appbrk{}ध\appbrk{}क\appbrk{}षायं}}\Afootnote{K H; \skt{निम्बारग्वधयोः कल्को हित}\ A}}\skt{ }\edtext{\skt{चोत्सादने}}{\lemma{\skt{चोत्सादने}}\Afootnote{K\textsuperscript{pc} H\textsuperscript{ac}; \skt{उत्सादने भवेत्}\ A; \skt{त्सादने}\ K\textsuperscript{ac}; \skt{चोत्साद नार्थं}\ H\textsuperscript{pc}}}\skt{~॥}\pend

\pstart \ledsidenote{\passref{4.20.40-41}}\edtext{\skt{परिवर्त्तिकां तु}}{\lemma{\skt{परिवर्त्तिकां तु}}\Afootnote{K H; \skt{परिवृ}\ A}}\skt{ }\edtext{\skt{सर्पिषाभ्यज्य}}{\lemma{\skt{सर्पिषाभ्यज्य}}\Afootnote{K H; \skt{त्तिं घृताभ्य}\ A}}\skt{ }\edtext{\skt{स्वेदयेत्}}{\lemma{\skt{स्वेदयेत्}}\Afootnote{K H; \emph{om.}\ A}}\skt{~। }\edtext{\edtext{\skt{वातघ्नैस्ततः}}{\lemma{\skt{वातघ्नैस्}}\Afootnote{K H; \skt{क्तां}\ A}}}{\lemma{\skt{ततः}}\Afootnote{K H; \skt{सु}\ A}}\skt{ }\edtext{\edtext{\skt{स्निन्ना\-मु\-प\-ना\-ह\-येत्}}{\lemma{\skt{स्निन्नाम्}}\Afootnote{K H; \skt{स्विन्नाम्}\ A}}}{\lemma{\skt{उपनाहयेत्}}\Afootnote{K H; \skt{उपनाहये}\ A}}\skt{~। }\edtext{\skt{साल्वलेन}}{\lemma{\skt{साल्वलेन}}\Afootnote{K H; \emph{om.}\ A}}\skt{ }\edtext{\skt{वेशवारेण वा}}{\lemma{\skt{वेशवारेण वा}}\Afootnote{K H; \emph{om.}\ A}}\skt{ }\edtext{\edtext{\skt{त्रिरा\-त्र\-म्प\-ञ्च\-रा\-त्रम्वा}}{\lemma{\skt{त्रिरात्रम्}}\Afootnote{K H; \emph{om.}\ A}}}{\lemma{\skt{पञ्चरात्रम्वा}}\Afootnote{K H; \skt{त्}\ A}}\skt{ ततः }\edtext{\skt{स्वेदयित्वाभ्यज्य}}{\lemma{\skt{स्वेदयित्वाभ्यज्य}}\Afootnote{K H; \skt{ऽभ्यज्य}\ A}}\skt{ }\edtext{\skt{शनैः}}{\lemma{\skt{शनैः}}\Afootnote{K H; \emph{om.}\ A}}\skt{ }\edtext{\skt{शनैश्चर्म}}{\lemma{\skt{चर्म}}\Afootnote{K H; \skt{चर्}\ A}}\skt{ }\edtext{\skt{समानयं}}{\lemma{\skt{समानयं}}\Afootnote{K H; \emph{om.}\ A}}\skt{ }\edtext{\edtext{\skt{मनि\-प\-थ\-म\-व\-पी\-ड\-येत्}}{\lemma{\skt{मनिपथम्}}\Afootnote{K H; \skt{म}\ A}}}{\lemma{\skt{अवपीडयेत्}}\Afootnote{K H; \skt{चान\appbrk{}ये\appbrk{}त्पी\appbrk{}ड\appbrk{}ये\appbrk{}न्म\appbrk{}णिम्॥ प्रविष्टे च मणौ चर्म स्वेद\appbrk{}ये\appbrk{}दु\appbrk{}प\appbrk{}ना\appbrk{}हनैः। त्रिरात्रं पञ्चरात्रं वा वातघ्नैः साल्वणादिभिः}\ A}}\skt{~।}\pend

\pstart \ledsidenote{\passref{4.20.42}}\edtext{\skt{परिवेष्टितञ्च}}{\lemma{\skt{परिवेष्टितञ्च}}\Afootnote{K H; \emph{om.}\ A}}\skt{ }\edtext{\edtext{\skt{पुन\-रु\-प\-ना\-ह\-येत्}}{\lemma{\skt{पुनर्}}\Afootnote{K H; \skt{द}\ A}}}{\lemma{\skt{उपनाहयेत्}}\Afootnote{K H; \skt{द्याद्वातहरा}\ A}}\skt{~। }\edtext{\skt{नित्यञ्च}}{\lemma{\skt{नित्यञ्च}}\Afootnote{K H; \skt{न्ब}\ A}}\skt{ }\edtext{\edtext{\edtext{\skt{स्निग्ध\-म\-न्न\-मा\-ह\-रेत्}}{\lemma{\skt{स्निग्धम्}}\Afootnote{K H; \skt{स्तीन्स्निग्धान्य्}\ A}}}{\lemma{\skt{अन्नम्}}\Afootnote{K H; \skt{अन्ना}\ A}}\skt{~॥ }\edtext{\edtext{\skt{अव\-पा\-टि\-क\-या\-म\-प्येवं}}{\lemma{\skt{अप्य्}}\Afootnote{K H; \skt{जये}\ A}}}{\lemma{\skt{एवं}}\Afootnote{K H; \skt{देवं}\ A}}}{\lemma{\skt{आहरेत्~॥ अव\appbrk{}पा\appbrk{}टि\appbrk{}क\appbrk{}याम्}}\Afootnote{K; \skt{नि भोजयेत्~। वपाटिकां}\ A; \skt{आहरेत्~॥ अव\appbrk{}पा\appbrk{}टि\appbrk{}का\appbrk{}याम्}\ H}}\skt{ }\edtext{\skt{यथादोषं}}{\lemma{\skt{यथादोषं}}\Afootnote{K H; \skt{यथादोषं चिकित्सकः}\ A}}\skt{~॥}\pend

\pstart \ledsidenote{\passref{4.20.43}}\edtext{\skt{निरुद्धप्रकाशे तु}}{\lemma{\skt{निरुद्धप्रकाशे तु}}\Afootnote{K H; \skt{निरुद्धप्रकशे}\ A}}\skt{ }\edtext{\skt{सुवर्ण्णादीनां}}{\lemma{\skt{सुवर्ण्णादीनां}}\Afootnote{K H; \emph{om.}\ A}}\skt{ }\edtext{\skt{नाडी\-मु\-भ\-य\-तो\-मुखीं}}{\lemma{\skt{उभयतोमुखीं}}\Afootnote{K H; \skt{लौहीमुभयतो मुखीम्}\ A}}\skt{ }\edtext{\skt{दारुवीम्वा}}{\lemma{\skt{दारुवीम्वा}}\Afootnote{K; \skt{दारवीं वा}\ A; \skt{दार्व्वीम्वा}\ H}}\skt{ }\edtext{\skt{जतुवृतां}}{\lemma{\skt{जतुवृतां}}\Afootnote{K H; \skt{जतुकृतां}\ A}}\skt{ }\edtext{\skt{सर्पिषाभ्यज्य}}{\lemma{\skt{सर्पिषाभ्यज्य}}\Afootnote{K H; \skt{घृताभ्य}\ A}}\skt{ }\edtext{\skt{स्रोतसि}}{\lemma{\skt{स्रोतसि}}\Afootnote{K H; \skt{क्तां प्रवेशयेत्}\ A}}\skt{~।}\pend

\pstart \ledsidenote{\passref{4.20.44}}\edtext{\skt{प्रणिधाय}}{\lemma{\skt{प्रणिधाय}}\Afootnote{K H; \skt{परिषेके वसामज्ज}\ A}}\skt{ }\edtext{\skt{शिशु\-मा\-र\-व\-सा\-व\-रा\-ह\-म\-ज्जाभिः}}{\lemma{\skt{शिशु\appbrk{}मा\appbrk{}र\appbrk{}व\appbrk{}सा\appbrk{}व\appbrk{}रा\appbrk{}ह\appbrk{}म\appbrk{}ज्जाभिः}}\Afootnote{H; \skt{शिशुमारव}\ A; \skt{शिशि मार\appbrk{}व\appbrk{}सा\appbrk{}व\appbrk{}रा\appbrk{}ह\appbrk{}म\appbrk{}ज्जाभिः}\ K\textsuperscript{ac}; \skt{शिशु\appbrk{}मा\appbrk{}र\appbrk{}व\appbrk{}सा\appbrk{}व\appbrk{}रा\appbrk{}ह\appbrk{}म\appbrk{}ज्जाभिः}\ K\textsuperscript{pc}}}\skt{ }\edtext{\edtext{\skt{परि\-षे\-च\-ये\-च्च\-क्र\-तै\-लेन}}{\lemma{\skt{परिषेचयेच्}}\Afootnote{K\textsuperscript{pc} H; \skt{राहयोः}\ A; \skt{परिषेयेच्}\ K\textsuperscript{ac}}}}{\lemma{\skt{चक्रतैलेन}}\Afootnote{K H; \skt{चक्रतैलं तथा}\ A}}\skt{ }\edtext{\skt{वात्}}{\lemma{\skt{वात्}}\Afootnote{K; \skt{योज्यं वात\appbrk{}घ्न\appbrk{}द्र\appbrk{}व्य\appbrk{}सं\appbrk{}यु\appbrk{}तम्}\ A; \skt{वा}\ H}}\skt{~। }\edtext{\skt{सुर\-भि\-व\-च\-दे\-व\-दा\-रु\-सि\-द्धेन}}{\lemma{\skt{सुर\appbrk{}भि\appbrk{}व\appbrk{}च\appbrk{}दे\appbrk{}व\appbrk{}दा\appbrk{}रु\appbrk{}सि\appbrk{}द्धेन}}\Afootnote{K; \emph{om.}\ A; \skt{सुर\appbrk{}भि\appbrk{}व\appbrk{}चा\appbrk{}दे\appbrk{}व\appbrk{}दा\appbrk{}रु\appbrk{}सि\appbrk{}द्धेन}\ H}}\skt{~।}\pend

\pstart \ledsidenote{\passref{4.20.45}}\edtext{\edtext{\skt{त्र्यहात्त्र्यहाच्च}}{\lemma{\skt{त्र्यहात्}}\Afootnote{K H; \skt{त्र्यहा}\ A}}}{\lemma{\skt{त्र्यहाच्च}}\Afootnote{K H; \skt{त्त्र्यहात्}\ A}}\skt{ }\edtext{\skt{स्थूलतरीं}}{\lemma{\skt{स्थूलतरीं}}\Afootnote{K H; \skt{स्थूलतरां}\ A}}\skt{ }\edtext{\edtext{\edtext{\skt{नाडी\-म्प्र\-णि\-द\-ध्या\-दे\-व\-मस्य}}{\lemma{\skt{नाडीम्}}\Afootnote{K H; \skt{सम्यङ्नाडीं}\ A}}}{\lemma{\skt{प्रणि\appbrk{}द\appbrk{}ध्या\appbrk{}दे\appbrk{}वम्}}\Afootnote{K H; \skt{प्रवेश}\ A}}}{\lemma{\skt{अस्य}}\Afootnote{K H; \skt{येत्}\ A}}\skt{ }\edtext{\skt{स्रोतोद्वारं}}{\lemma{\skt{स्रोतोद्वारं}}\Afootnote{K H; \skt{स्रोतो वि}\ A}}\skt{ }\edtext{\skt{वर्द्धयेत्}}{\lemma{\skt{वर्द्धयेत्}}\Afootnote{K H; \skt{वर्धयेदेवं}\ A}}\skt{~। }\edtext{\edtext{\edtext{\skt{स्निग्ध\-ञ्चा\-हा\-र\-मु\-प\-से\-वेत}}{\lemma{\skt{स्निग्धञ्}}\Afootnote{K H; \skt{स्निग्}\ A}}}{\lemma{\skt{चाहारम्}}\Afootnote{K H; \emph{om.}\ A}}}{\lemma{\skt{उपसेवेत}}\Afootnote{K H; \skt{धमन्नं च भोजयेत्}\ A}}\skt{~।}\pend

\pstart \ledsidenote{\passref{4.20.46ab}}\edtext{\skt{शस्त्रेण वा}}{\lemma{\skt{शस्त्रेण वा}}\Afootnote{K H; \skt{भित्त्वा वा}\ A}}\skt{ सेवनीं }\edtext{\skt{विमु\-च्या\-व\-दा\-र\-येत्}}{\lemma{\skt{विमु\appbrk{}च्या\appbrk{}व\appbrk{}दा\appbrk{}र\appbrk{}येत्}}\Afootnote{K H; \skt{मुक्त्}\ A}}\skt{~। }\edtext{\skt{ततः}}{\lemma{\skt{ततः}}\Afootnote{K H; \skt{वा}\ A}}\skt{ }\edtext{\skt{सद्य}}{\lemma{\skt{सद्य}}\Afootnote{K; \skt{सद्यः}\ A H}}\skt{ }\edtext{\skt{क्षत\-व\-दु\-पा\-च\-रेत्}}{\lemma{\skt{क्षत\appbrk{}व\appbrk{}दु\appbrk{}पा\appbrk{}च\appbrk{}रेत्}}\Afootnote{K H; \skt{क्षत\appbrk{}व\appbrk{}दा\appbrk{}च\appbrk{}रेत्}\ A}}\skt{~।}\pend

\pstart \ledsidenote{\passref{4.20.46cd-47ab}}\edtext{\skt{अग्निरोहिणीं}}{\lemma{\skt{अग्निरोहिणीं}}\Afootnote{K H; \emph{om.}\ A}}\skt{ }\edtext{\skt{सन्निरुद्धगुदं}}{\lemma{\skt{सन्निरुद्धगुदं}}\Afootnote{K; \skt{सन्निरुद्धगुदं रोगं}\ A; \skt{सन्निरुद्धगुद}\ H}}\skt{ }\edtext{\skt{वल्मीकञ्च}}{\lemma{\skt{वल्मीकञ्च}}\Afootnote{K H; \skt{वल्मीकं व}\ A}}\skt{ }\edtext{\skt{प्रत्या\-ख्या\-यो\-प\-च\-रेत्}}{\lemma{\skt{प्रत्या\appbrk{}ख्या\appbrk{}यो\appbrk{}प\appbrk{}च\appbrk{}रेत्}}\Afootnote{K H; \skt{ह्निरोहिणीम्॥ प्रत्याख्याय यथायोगं चिकि\appbrk{}त्सि\appbrk{}त\appbrk{}म\appbrk{}था\appbrk{}च\appbrk{}रेत्}\ A}}\skt{~।}\pend

\pstart \ledsidenote{\passref{4.20.47cd}}\skt{विसर्पोक्तेन }\edtext{\skt{विधा\-ने\-ना\-ग्नि\-रो\-हिणीं}}{\lemma{\skt{विधा\appbrk{}ने\appbrk{}ना\appbrk{}ग्नि\appbrk{}रो\appbrk{}हिणीं}}\Afootnote{K H; \skt{विधिना साध\appbrk{}ये\appbrk{}द\appbrk{}ग्नि\appbrk{}रो\appbrk{}हि\appbrk{}णीम्}\ A}}\skt{ }\edtext{\skt{चिकित्सेत्}}{\lemma{\skt{चिकित्सेत्}}\Afootnote{K H; \emph{om.}\ A}}\skt{~।}\pend

\pstart \ledsidenote{\passref{4.20.48ab}}\edtext{\skt{सन्निरुद्धगुदं}}{\lemma{\skt{सन्निरुद्धगुदं}}\Afootnote{K H; \skt{सन्निरुद्धगुदे}\ A}}\skt{ }\edtext{\skt{निरु\-द्ध\-प्र\-क\-श\-वत्}}{\lemma{\skt{निरु\appbrk{}द्ध\appbrk{}प्र\appbrk{}क\appbrk{}श\appbrk{}वत्}}\Afootnote{K H; \skt{योज्या निरु\appbrk{}द्ध\appbrk{}प्र\appbrk{}क\appbrk{}श\appbrk{}क्रिया}\ A}}\skt{~।}\pend

\pstart \ledsidenote{\passref{4.20.48cd-49}}\edtext{\skt{वल्मीकं}}{\lemma{\skt{वल्मीकं}}\Afootnote{K H; \emph{om.}\ A}}\skt{ }\edtext{\edtext{\edtext{\skt{शस्त्रे\-णो\-त्कृ\-त्या\-ग्नि\-क्षा\-रा\-भ्या\-मु\-प\-च\-रे\-त्ततः}}{\lemma{\skt{शस्त्रे\appbrk{}णो\appbrk{}त्कृ\appbrk{}त्या\appbrk{}ग्नि\appbrk{}क्षा\appbrk{}रा\appbrk{}भ्याम्}}\Afootnote{K H; \skt{शस्त्रेणोत्कृत्य वल्मीकं क्षाराग्निभ्यां}\ A}}}{\lemma{\skt{उपचरेत्}}\Afootnote{K H; \emph{om.}\ A}}}{\lemma{\skt{ततः}}\Afootnote{K H; \emph{om.}\ A}}\skt{ }\edtext{\skt{संशोध्य}}{\lemma{\skt{संशोध्य}}\Afootnote{K H; \skt{प्र}\ A}}\skt{ }\edtext{\skt{रोप\-ये\-द\-र्बु\-द\-वि\-धा\-णेन}}{\lemma{\skt{रोप\appbrk{}ये\appbrk{}द\appbrk{}र्बु\appbrk{}द\appbrk{}वि\appbrk{}धा\appbrk{}णेन}}\Afootnote{K; \skt{साधयेत्॥ विधा\appbrk{}ने\appbrk{}ना\appbrk{}र्बु\appbrk{}दो\appbrk{}क्तेन शोधयित्वा च रोपयेत्। वल्मीकं तु भवेद्यस्य नाति\appbrk{}वृ\appbrk{}द्ध\appbrk{}म\appbrk{}म\appbrk{}र्म\appbrk{}जम्}\ A; \skt{रोप\appbrk{}ये\appbrk{}द\appbrk{}र्ब्बु\appbrk{}द\appbrk{}वि\appbrk{}धा\appbrk{}नेन}\ H}}\skt{~॥}\pend

\pstart \ledsidenote{\passref{4.20.57-58ab}}\edtext{\skt{अहिपूतने}}{\lemma{\skt{अहिपूतने}}\Afootnote{K H; \emph{om.}\ A}}\skt{ धात्र्याः स्तन्यं }\edtext{\skt{शोधेत}}{\lemma{\skt{शोधेत}}\Afootnote{K; \skt{शोधयित्वा बाले साध्या+\appbrk{}अ\appbrk{}हि\appbrk{}पू\appbrk{}तना। पटो}\ A; \skt{शोध्येत~। बा}\ H}}\skt{~। }\edtext{\skt{बालस्य}}{\lemma{\skt{बालस्य}}\Afootnote{K; \skt{लपत्र}\ A; \skt{लस्य}\ H}}\skt{ }\edtext{\skt{त्रिफ\-ला\-र\-सां\-ज\-न\-प\-टोलैः}}{\lemma{\skt{त्रिफ\appbrk{}ला\appbrk{}र\appbrk{}सां\appbrk{}ज\appbrk{}न\appbrk{}प\appbrk{}टोलैः}}\Afootnote{K H; \skt{त्रिफ\appbrk{}ला\appbrk{}र\appbrk{}सा\appbrk{}ञ्जन}\ A}}\skt{ }\edtext{\skt{सिद्धं}}{\lemma{\skt{सिद्धं}}\Afootnote{K H; \emph{om.}\ A}}\skt{ }\edtext{\skt{सर्पिः}}{\lemma{\skt{सर्पिः}}\Afootnote{K H; \skt{वि}\ A}}\skt{ }\edtext{\edtext{\skt{पानमिष्टं}}{\lemma{\skt{पानम्}}\Afootnote{K H; \skt{पाचितम्॥ पीतं घृतं नाशयति कृच्छ्रा\appbrk{}म\appbrk{}प्य\appbrk{}हि\appbrk{}पू\appbrk{}त\appbrk{}नाम्}\ A}}}{\lemma{\skt{इष्टं}}\Afootnote{K H; \emph{om.}\ A}}\skt{~।}\pend

\pstart \ledsidenote{\passref{4.20.58cd}}\edtext{\edtext{\skt{त्रिफ\-ला\-ख\-दि\-र\-ब\-द\-री\-क\-षा\-य\-म्व्र\-ण\-शो\-धनं}}{\lemma{\skt{त्रिफ\appbrk{}ला\appbrk{}ख\appbrk{}दि\appbrk{}र\appbrk{}ब\appbrk{}द\appbrk{}री\appbrk{}क\appbrk{}षा\appbrk{}यम्}}\Afootnote{K H; \skt{त्रिफ\appbrk{}ला\appbrk{}को\appbrk{}ल\appbrk{}ख\appbrk{}दि\appbrk{}र\appbrk{}क\appbrk{}षायं}\ A}}}{\lemma{\skt{व्रणशोधनं}}\Afootnote{K H; \skt{व्रणरोपणम्}\ A}}\skt{~।}\pend

\pstart \ledsidenote{\passref{4.20.59}}\edtext{\edtext{\skt{कासी\-स\-रो\-च\-ना\-तु\-त्थ\-ह\-रि\-ता\-ल\-र\-सा\-ञ्ज\-नै\-र\-म्ल\-पिष्टैः}}{\lemma{\skt{कासी\appbrk{}स\appbrk{}रो\appbrk{}च\appbrk{}ना\appbrk{}तु\appbrk{}त्थ\appbrk{}ह\appbrk{}रि\appbrk{}ता\appbrk{}ल\appbrk{}र\appbrk{}सा\appbrk{}ञ्ज\appbrk{}नैर्}}\Afootnote{A K; \skt{कासी\appbrk{}स\appbrk{}रो\appbrk{}च\appbrk{}ना\appbrk{}तु\appbrk{}त्था\appbrk{}ह\appbrk{}रि\appbrk{}ता\appbrk{}र\appbrk{}सा\appbrk{}ञ्ज\appbrk{}नैर्}\ H}}}{\lemma{\skt{अम्लपिष्टैः}}\Afootnote{K H; \skt{लेपोऽम्लपिष्टो}\ A}}\skt{ }\edtext{\skt{प्रलेपनं}}{\lemma{\skt{प्रलेपनं}}\Afootnote{K H; \emph{om.}\ A}}\skt{~। }\edtext{\skt{बदरचूर्ण्णैर्वा}}{\lemma{\skt{बदरचूर्ण्णैर्वा}}\Afootnote{K H; \skt{बदरीत्वग्वा}\ A}}\skt{ }\edtext{\skt{सैन्धवमिश्रैः}}{\lemma{\skt{सैन्धवमिश्रैः}}\Afootnote{K H; \skt{सैन्धवसंयुता}\ A}}\skt{~।}\pend

\pstart \ledsidenote{\passref{4.20.60}}\edtext{\skt{चूर्ण्णकाले तु}}{\lemma{\skt{चूर्ण्णकाले तु}}\Afootnote{K H; \emph{om.}\ A}}\skt{ }\edtext{\skt{चूर्ण्णः}}{\lemma{\skt{चूर्ण्णः}}\Afootnote{K H; \emph{om.}\ A}}\skt{ }\edtext{\skt{कपालतुत्थजः}}{\lemma{\skt{कपालतुत्थजः}}\Afootnote{K H; \skt{कपालतुत्थजं}\ A}}\skt{~॥ }\edtext{\edtext{\skt{वृष\-ण\-क\-च्छ्वा\-म्पा\-मा\-व\-दु\-प\-चारो}}{\lemma{\skt{वृषणकच्छ्वाम्}}\Afootnote{K H; \skt{चूर्णं चूर्णकाले प्रयोजयेत्। चिकि\appbrk{}त्से\appbrk{}न्मु\appbrk{}ष्क\appbrk{}कच्छूं}\ A}}}{\lemma{\skt{पामा\appbrk{}व\appbrk{}दु\appbrk{}प\appbrk{}चारो}}\Afootnote{K H; \skt{चाप्य्}\ A}}\skt{ }\edtext{\edtext{\skt{हिपूतनवच्च}}{\lemma{\skt{हिपूतनवच्}}\Afootnote{K; \skt{अहि\appbrk{}पू\appbrk{}त\appbrk{}न\appbrk{}पा\appbrk{}म\appbrk{}वत्}\ A; \skt{ऽहि\appbrk{}पू\appbrk{}त\appbrk{}न\appbrk{}वच्}\ H}}}{\lemma{\skt{च}}\Afootnote{K H; \emph{om.}\ A}}\skt{~॥}\pend

\pstart \ledsidenote{\passref{4.20.61}}\edtext{\skt{गुदभ्रंशे तु}}{\lemma{\skt{गुदभ्रंशे तु}}\Afootnote{K H; \skt{गुदभ्रंशे}\ A}}\skt{ गुदं }\edtext{\skt{स्वेद्य\-स्ने\-हे\-ना\-भ्यज्य}}{\lemma{\skt{स्वेद्य\appbrk{}स्ने\appbrk{}हे\appbrk{}ना\appbrk{}भ्यज्य}}\Afootnote{K H; \skt{स्विन्नं स्नेहाभ्यक्तं}\ A}}\skt{ }\edtext{\skt{प्रवेश्य}}{\lemma{\skt{प्रवेश्य}}\Afootnote{K H; \skt{प्रवे}\ A}}\skt{ }\edtext{\skt{प्रविष्टे च}}{\lemma{\skt{प्रविष्टे च}}\Afootnote{K H; \skt{श}\ A}}\skt{ }\edtext{\skt{स्वेदं}}{\lemma{\skt{स्वेदं}}\Afootnote{K H; \emph{om.}\ A}}\skt{ }\edtext{\skt{कुर्या\-दु\-च्चा\-र\-नि\-र्ग\-म\-नार्थं}}{\lemma{\skt{कुर्या\appbrk{}दु\appbrk{}च्चा\appbrk{}र\appbrk{}नि\appbrk{}र्ग\appbrk{}म\appbrk{}नार्थं}}\Afootnote{K H; \skt{येत्। कार\appbrk{}ये\appbrk{}द्गो\appbrk{}फ\appbrk{}णा\appbrk{}बन्धं}\ A}}\skt{~। }\edtext{\skt{छिद्रेण}}{\lemma{\skt{छिद्रेण}}\Afootnote{K H; \skt{मध्यच्छिद्रेण}\ A}}\skt{ }\edtext{\skt{चास्य}}{\lemma{\skt{चास्य}}\Afootnote{K H; \emph{om.}\ A}}\skt{ चर्मणा }\edtext{\skt{कौपीनबन्धः}}{\lemma{\skt{कौपीनबन्धः}}\Afootnote{K H; \emph{om.}\ A}}\skt{~।}\pend

\pstart \ledsidenote{\passref{4.20.62}}\edtext{\skt{तैलं}}{\lemma{\skt{तैलं}}\Afootnote{K H; \emph{om.}\ A}}\skt{ }\edtext{\skt{चास्य}}{\lemma{\skt{चास्य}}\Afootnote{K H; \emph{om.}\ A}}\skt{ }\edtext{\edtext{\skt{पाना\-भ्य\-ञ्ज\-ना\-र्थ\-म्वि\-प\-चेत्}}{\lemma{\skt{पाना\appbrk{}भ्य\appbrk{}ञ्ज\appbrk{}ना\appbrk{}र्थम्}}\Afootnote{K H; \skt{विनिर्गमार्थं}\ A}}}{\lemma{\skt{विपचेत्}}\Afootnote{K H; \skt{वायोश्च स्}\ A}}\skt{~। }\edtext{\skt{मूषिकं}}{\lemma{\skt{मूषिकं}}\Afootnote{K H; \skt{वेदये}\ A}}\skt{ }\edtext{\skt{निरंत्रं}}{\lemma{\skt{निरंत्रं}}\Afootnote{K H; \skt{च्}\ A}}\skt{ }\edtext{\skt{कल्पयित्वा}}{\lemma{\skt{कल्पयित्वा}}\Afootnote{K H; \skt{च मुहु}\ A}}\skt{ }\edtext{\edtext{\skt{महा\-प\-ञ्च\-मू\-ल\-क्षी\-र\-म्वि\-प\-चेत्}}{\lemma{\skt{महा\appbrk{}प\appbrk{}ञ्च\appbrk{}मू\appbrk{}ल\appbrk{}क्षी\appbrk{}रम्}}\Afootnote{K H; \skt{र्मुहुः। क्षीरे महत्पञ्चमूलं मूषिकां}\ A}}}{\lemma{\skt{विपचेत्}}\Afootnote{K H; \skt{चान्त्रवर्जिताम्}\ A}}\skt{~।}\pend

\pstart \ledsidenote{\passref{4.20.63}}\edtext{\edtext{\skt{ततस्तस्मिं}}{\lemma{\skt{ततस्}}\Afootnote{K H; \skt{पक्त्वा}\ A}}}{\lemma{\skt{तस्मिं}}\Afootnote{K H; \skt{तस्मिन्}\ A}}\skt{ }\edtext{\edtext{\edtext{\edtext{\edtext{\skt{तैल\-म्वा\-त\-ह\-र\-क\-ल्क\-प्र\-ति\-वा\-प\-म्वि\-प\-चे\-दे\-त\-त्कृ\-च्छ्र\-मपि}}{\lemma{\skt{तैलम्}}\Afootnote{K H; \skt{पचेतैलं}\ A}}}{\lemma{\skt{वात\appbrk{}ह\appbrk{}र\appbrk{}क\appbrk{}ल्क\appbrk{}प्र\appbrk{}ति\appbrk{}वा\appbrk{}पम्}}\Afootnote{K H; \skt{वातघ्}\ A}}}{\lemma{\skt{विपचेदेतत्}}\Afootnote{K H; \skt{नौषध}\ A}}}{\lemma{\skt{कृच्छ्रम्}}\Afootnote{K H; \skt{सं}\ A}}\skt{ गुदत्रं संसाधयतेति}}{\lemma{\skt{अपि गुदत्रं संसाधयतेति}}\Afootnote{K H; \skt{युतम्। गुदभ्रंशमिदं कृच्छ्रं पाना\appbrk{}भ्य\appbrk{}ङ्गा\appbrk{}त्प्र\appbrk{}सा\appbrk{}ध\appbrk{}येत्}\ A}}\skt{~॥}\pend

\pstart \ledsidenote{\passref{4.20.trailer}}\edtext{\skt{चि}}{\lemma{\skt{चि}}\Afootnote{K; \skt{इति सुश्रु\appbrk{}त\appbrk{}सं\appbrk{}हि\appbrk{}तायां चिकित्सास्थाने क्षुद्र\appbrk{}रो\appbrk{}ग\appbrk{}चि\appbrk{}कि\appbrk{}त्सितं नाम विंशोऽध्यायः}\ A; \skt{चिकित्सा}\ Kb H}}\skt{ २०~॥ ०~॥}\pend

\endnumbering
